\documentclass[11pt]{article}
\usepackage[a4paper,margin=25mm]{geometry}
\usepackage{amsmath,amssymb,amsthm,mathtools}
\usepackage[hidelinks,hypertexnames=false]{hyperref}
\usepackage{graphicx}
\usepackage{mathrsfs}
\graphicspath{{figures/}{figures_029/}{figures_033/}{figures_finite_residue/}{figures_source_spectrum/}{figures_global_residue/}{figures_es_total/}{figures_prime_collision/}{./}}
\newtheorem{proposition}{Proposition}
\newtheorem{theorem}{Theorem}
\providecommand{\W}{\mathcal W}
\newcommand{\C}{\mathbb C}
\newcommand{\R}{\mathbb R}
\newcommand{\diag}{\operatorname{diag}}
\newcommand{\ord}{\operatorname{ord}}
\newcommand{\rank}{\operatorname{rank}}
\newcommand{\tr}{\operatorname{tr}}
\title{The ES--Fable inverse correspondence in the original weighted conductor}
\author{Split-Zero programme: coordinate calculation}
\date{20 September 2026}

\providecommand{\Spec}{\operatorname{Spec}}
\providecommand{\Disc}{\operatorname{Disc}}
\providecommand{\Tr}{\operatorname{Tr}}
\providecommand{\Span}{\operatorname{span}}

\providecommand{\im}{\operatorname{im}}
\providecommand{\Div}{\operatorname{div}}
\providecommand{\Tr}{\operatorname{Tr}}
\providecommand{\Hom}{\operatorname{Hom}}
\providecommand{\Res}{\operatorname{Res}}
\begin{document}
\maketitle

This calculation responds to equation (11) of the supplied audit: the
inverse-exterior bound for the matrix in WCF6 \cite{WCF}. It constructs
the marked map from the four ES--Fable locales, follows it through the
original quartet and degree-eight conductor, and determines where the
receiving finite linear map can lose rank. All four labels are retained.
The conclusions concern this precise map. They neither assert RH nor
classify all nonlinear maps in the programme.
The original three-dimensional polynomial example was announced by
Levent Alp\"oge with credit to Fable; Tao's exposition \cite{Tao}
reproduces its formula. The additional four-Time constructions used here
are credited to the ES programme source \cite{ES}, rather than attributed
to that original three-point fibre.

The continuation solves the entire inverse problem for the displayed
four-dimensional map: every finite fibre, its complete omitted image,
its degree-eight inverse domain and its complete nonproper-value locus.
It gives the eight signed states over the literal arithmetic quartet,
their common limiting omitted target, and their escape in a fixed
original Gamma norm. For the earlier seven-point target it evaluates
the complete singular spectrum, all exterior inverse asymptotes, and
the exact volume distortion of explicitly transported coordinate flags.
The original conductor intertwines these maps throughout; its coefficients
are never chosen to fit the new polynomial.

The signed-state continuation gives a nonzero exact determinant for the
four sign differences on the literal quartet, the full rank-three mean
map and its explicit kernel line, and the entire eight-state evaluation
kernel. Retaining labels and evaluation gives an invertible total map
through two copies of the original conductor. Its signed-root symmetry
has order $192$; all monodromy generators, its deck transformations,
the complete mixing block and its ranks are calculated explicitly.

\section{Four points, including the coordinates hidden by local symmetry}
The source \cite{ES}, equations \texttt{four-Time-quarter-locales},
\texttt{times-weight-and-unit-vectors}, and
\texttt{four-Time-quarter-coordinate-matrix}, defines
\[
 \mathcal T=(\mathrm M,\mathrm N,\mathrm T,\mathrm E),\qquad
 V=\{v\in\R^4:\mathbf1^Tv=0\},\qquad
 \mu_a=\frac2{\sqrt3}(e_a-\tfrac14\mathbf1),\qquad
 \mathfrak q_a=-\frac13\mu_a.
 \tag{FC1}
\]
Here $e_a$ are the four separate coordinate vectors; no two are identified.
Consequently
\[
 \langle\mathfrak q_a,\mathfrak q_b\rangle
 =\begin{cases}1/9&a=b,\\-1/27&a\ne b,\end{cases}
 \qquad
 Q_{ab}:=\tfrac34\langle\mathfrak q_b,\mu_a\rangle
 =\frac1{12}-\frac13\delta_{ab}.
 \tag{FC2}
\]
The matrix $Q$ has four distinct columns, although each column has its
own local altitude $-1/4$. Direct multiplication gives
\[
 Q\mathbf1=0,\qquad Q|_{\mathbf1^\perp}=-\tfrac13 I,
 \qquad \rank Q=3.
 \tag{FC3}
\]
Thus a specific noninvertible map is already present: simultaneous
altitude readout. Its entire kernel is the constant vector. This is the
source's tetrahedral projection, not a singularity of WCF6.

The affine map which recovers the four indicators is
\[
 b:V_{\C}\longrightarrow\{c\in\C^4:\mathbf1^Tc=1\},\qquad
 b(v)=\tfrac14\mathbf1-\tfrac{3\sqrt3}{2}v,
 \quad b(\mathfrak q_a)=e_a.
 \tag{FC4}
\]
Its inverse is $c\mapsto-2(c-\mathbf1/4)/(3\sqrt3)$.
Both formulas follow by substitution in FC1. They commute with all
simultaneous permutations of the four labels. Retaining the affine term
in FC4 recovers the constant direction that FC3 loses.

\section{Exact transport into the original quartet}
Retain the WCF simple quartet and its full unit, with
$0<\delta<1/2$ and $\gamma>2$. Fix the explicit marking
\[
 \begin{array}{c|cccc}
 a&\mathrm M&\mathrm N&\mathrm T&\mathrm E\\\hline
 \rho_a&\frac12+\delta+i\gamma&\frac12+\delta-i\gamma&
 \frac12-\delta+i\gamma&\frac12-\delta-i\gamma.
 \end{array}
 \tag{FC5}
\]
This is a declared marking, not an inference that the four Time operators
are the arithmetic dilation operator. All other markings are obtained
by simultaneous permutations; their effects on the Segre relation are
calculated below. Put $h(S)=\prod_a(S-\rho_a)$ and
\[
 \ell_a(S)=\prod_{b\ne a}\frac{S-\rho_b}{\rho_a-\rho_b}.
 \tag{FC6}
\]
The exact algebra map is
\[
 \Theta:\operatorname{Map}(\{\mathfrak q_a\},\C)
 \longrightarrow\C[S]/(h),\qquad
 f\longmapsto\sum_a f(\mathfrak q_a)[\ell_a].
 \tag{FC7}
\]
Evaluation at the four $\rho_a$ is its inverse. Indeed
$\ell_a(\rho_b)=\delta_{ab}$; a polynomial of degree at most three
vanishing at all four distinct points is zero. The same argument applied
to remainders modulo $h$ proves
$[\ell_a\ell_b]=\delta_{ab}[\ell_a]$ and
$\sum_a[\ell_a]=1$, so FC7 is an algebra isomorphism, not merely a
vector-space bijection.

For the full original Taylor unit $u_a=(g/h)(\rho_a)\ne0$, the
additional source-unit map is $c\mapsto\diag(u_a)c$, with inverse
$c\mapsto\diag(u_a^{-1})c$. Its determinant is $\prod_a u_a$.
No unit has been replaced by one in the preceding isomorphism or in
the period map below. The unit multiplication is a separate linear
isomorphism, not an unital algebra homomorphism.

The original exponential observation of a coefficient vector is
\[
 \mathcal E(c)(z)=\sum_a c_a e^{\rho_a z},\qquad
 \bigl(\mathcal E(c)^{(n)}(0)\bigr)_{n=0}^3
 =V_\rho c,\quad (V_\rho)_{na}=\rho_a^n.
 \tag{FC8}
\]
The Vandermonde determinant, with precisely the order FC5, is
\[
 \det V_\rho=\prod_{a<b}(\rho_b-\rho_a)
 =-64\delta^2\gamma^2(\delta^2+\gamma^2)\ne0.
 \tag{FC9}
\]
The product formula follows by subtracting columns and induction in the
number of columns; its evaluation uses the six differences
$-2i\gamma,-2\delta,-2\delta-2i\gamma,
-2\delta+2i\gamma,-2\delta,-2i\gamma$.
In particular the four transported points cannot collide under the full
first-four-jet observation. With the unit included the determinant is
FC9 times $\prod_a u_a$.

\section{The degree-eight map is the original conductor map}
Write $t=(t_1,t_2,t_3,t_4)$ in the FC5 order, and retain
\[
 Q_0(t)=t_1t_4-t_2t_3,\qquad
 t(z)=(e^{\rho_1z},e^{\rho_2z},e^{\rho_3z},e^{\rho_4z}).
 \tag{FC10}
\]
Let $\mathcal S_n=\C[t_1,t_2,t_3,t_4]_n$. Substitution along this
curve gives the linear map
\[
 \mathcal E_n:\mathcal S_n\to
 \operatorname{span}\{e^{\beta_{n,u,w}z}:0\le u,w\le n\},
 \quad P\mapsto P(t(z)),\qquad
 \beta_{n,u,w}=\tfrac n2+(2u-n)\delta+i(2w-n)\gamma.
 \tag{FC11}
\]
Its exact kernel is
\[
 \ker\mathcal E_n=Q_0\mathcal S_{n-2},\qquad
 \dim\operatorname{im}\mathcal E_n=(n+1)^2,
 \tag{FC12}
\]
where $\mathcal S_{n-2}=0$ for $n<2$.
To prove this, a monomial $t_1^{k_1}\cdots t_4^{k_4}$ is sent to the
exponential indexed by $u=k_1+k_2,w=k_1+k_3$. Every pair $0\le u,w\le n$
occurs: choose
$k_1=\max(0,u+w-n)$, $k_2=u-k_1$, $k_3=w-k_1$,
$k_4=n-u-w+k_1$. Replacing $t_1t_4$ by $t_2t_3$ reduces any monomial to
this unique representative. Thus the quotient by $Q_0$ has those
$(n+1)^2$ monomial representatives. Their exponents are distinct because
equality of real parts forces $u=u'$ and equality of imaginary parts
forces $w=w'$. Differentiating a linear relation in these exponentials
through order $(n+1)^2-1$ gives a Vandermonde matrix with nonzero
determinant. Hence they are linearly independent and FC12 follows.
At $n=8$, the map has rank $81$ and kernel dimension
$\binom{11}{3}-81=84$.

For the original admitted period let $H$ be its entire invertible period
matrix, including its Gamma factors, branch powers and Taylor unit.
Retain $D_5=\diag(\zeta,\zeta^2,\zeta^3,\zeta^4)$,
$\zeta=e^{2\pi i/5}$, and the precise original polynomials
\[
 Q_j(X)=Q_0(H^{-1}D_5^{-j}X),\qquad
 A_H(X)=\prod_{j=1}^4Q_j(X),\qquad
 E_A(z)=A_H(Ht(z)).
 \tag{FC13}
\]
These are NC5 and DFX4--5 \cite{NC,Cumulative}; no alternative conductor
is being selected. Expanding
$A_H(H(ZW,Z,W,1)^T)=\sum_{u,w=0}^8 a_{uw}Z^uW^w$ gives exactly
\[
 E_A(z)=\sum_{u,w=0}^8a_{uw}e^{\beta_{8,u,w}z},\qquad
 \mu_j=\sum_{u,w=0}^8a_{uw}\beta_{8,u,w}^{\,j}.
 \tag{FC14}
\]
This follows because
$t(z)=e^{(1/2-\delta-i\gamma)z}
(e^{2\delta z}e^{2i\gamma z},e^{2\delta z},e^{2i\gamma z},1)^T$
and $A_H$ is homogeneous of degree eight.

There is also a literal point map into the original quadric:
\[
 \mathfrak q_a\xmapsto{b}e_a\xmapsto{H}He_a,
 \qquad Q_0(e_a)=Q_0(H^{-1}He_a)=0.
\tag{FC15}
\]
The canonical 488-page reader \cite{ES488} also supplies an explicit
four-point polynomial collision in four coordinates, theorem
\texttt{explicit-four-time-Jacobian-collision}.
Its ordered source columns form the matrix
\[
 K=\begin{pmatrix}
 0&i&1/\sqrt2&1/\sqrt2\\
 0&-1&-1-\sqrt2 i&1-\sqrt2 i\\
 i/2&-3i&2\sqrt2+6i&-2\sqrt2+6i\\
 0&13&-34-19\sqrt2 i&34-19\sqrt2 i
 \end{pmatrix},\qquad \det K=77\sqrt2 i/2.
 \tag{FC15a}
\]
Expansion along the first column reduces this determinant to $i/2$
times a three-by-three determinant. Subtracting its second column from
its third gives the column $(0,2,68)^T$; the remaining expansion is
$30\sqrt2+94/\sqrt2=77\sqrt2$, proving the value. Hence the explicit
linear map $x\mapsto HK^{-1}x$, with
$K^{-1}=\operatorname{adj}(K)/(77\sqrt2 i/2)$, sends those four
original collision coordinates to exactly the same four $He_a$ in
FC15. The labelled map from these four source points to the Time locales
is $Ke_a\mapsto\mathfrak q_a$. This supplies a second, fully explicit
source for FC15 without asserting that two different ambient spaces or
their dynamics are identical. The companion independent derivation
records all four polynomial evaluations and their exact certificate.
Put $M_j=H^{-1}D_5^{-j}H$. The exact conductor value at this point is
\[
 A_H(He_a)=\prod_{j=1}^4
 \bigl((M_j)_{1a}(M_j)_{4a}-(M_j)_{2a}(M_j)_{3a}\bigr).
 \tag{FC16}
\]
Its vanishing means that $He_a$ lies on another quadric in the orbit.
It is not the same evaluation as any $\mu_v$: the latter is the
derivative at $z=0$ on the path $Ht(z)$, whose base point is $H\mathbf1$.
Equations FC13--16 give both maps and their different evaluation points
explicitly, rather than discarding the relation between the objects.

More precisely, the four incidences are four coordinates of the
\emph{original} symbol, not an auxiliary symbol. With
$P(t)=A_H(Ht)=\sum_{|m|=8}c_mt^m$, collection in FC11 gives
\[
 a_{u,w}=\sum_{\substack{|m|=8\\m_1+m_2=u,\ m_1+m_3=w}}c_m,
 \quad
 (A_H(He_1),A_H(He_2),A_H(He_3),A_H(He_4))
 =(a_{88},a_{80},a_{08},a_{00}).
 \tag{FC16a}
\]
Indeed each of those four corner pairs admits only the corresponding
pure eighth power. Define
$\mathcal D_\rho=\sum_j\rho_jt_j\partial_{t_j}$. The chain rule,
iterated on each polynomial $\mathcal D_\rho^nP$, proves
\[
 \mu_n=(\mathcal D_\rho^nP)(\mathbf1)
 =\sum_{j=1}^4A_H(He_j)(8\rho_j)^n
  +\sum_{(u,w)\notin\{(8,8),(8,0),(0,8),(0,0)\}}
      a_{uw}\beta_{8,u,w}^{\,n}.
 \tag{FC16b}
\]
Thus all four points enter the exact moment test, with the remaining
77 collected terms explicitly retained. Computing those four incidences
alone does not evaluate the remaining sum. For example, the test
polynomial $t_1^7t_2$ has zero at all four $e_j$ and value one at
$\mathbf1$; this establishes independence of the evaluation functionals,
without substituting that polynomial for the original conductor.

On the original five-orbit domain the quadrics $Q_0(H^{-1}X),Q_1,\ldots,Q_4$
are nonassociate irreducibles. The polynomial $Q_0(t)$ is irreducible:
as a degree-one polynomial in $t_1$ it is primitive over
$\C[t_2,t_3,t_4]$, with coprime coefficients $t_4$ and $-t_2t_3$,
and is irreducible over that coefficient field. An invertible linear
substitution preserves irreducibility. If the degree-eight polynomial
$A_H(Ht)$ belonged to $(Q_0)$, primality would force $Q_0$ to divide
one factor $Q_j(Ht)$, making those degree-two polynomials associates.
This contradicts the original orbit domain. Thus FC12 proves
\[
 E_A\not\equiv0,\qquad
 v=\ord_0E_A\le80,\qquad\mu_v\ne0.
 \tag{FC17}
\]
The order bound follows again from the $81$ distinct exponents and their
first $81$ derivative matrix. It is not a numerical lower bound for the
first nonzero derivative.
More precisely, if $r_A$ is the number of nonzero original collected
coefficients, the same derivative argument gives $v\le r_A-1$.
In particular, vanishing of all four incidences in FC16a leaves
$r_A\le77$ and gives $v\le76$; it does not annihilate the original
nonzero symbol. This support statement evaluates the effect of all four
point-incidence equations on the vanishing-order bound.

\section{All tetrahedral permutations, with the lost equation restored}
The full $S_4$ action on the four labels takes the line $\C Q_0$ to
the three lines represented by
\[
 Q_{14|23}=t_1t_4-t_2t_3,\quad
 Q_{12|34}=t_1t_2-t_3t_4,\quad
 Q_{13|24}=t_1t_3-t_2t_4.
 \tag{FC18}
\]
These correspond to the three partitions into two unordered pairs.
The stabilizer of the first line has order eight (independent exchanges
inside the pairs and exchange of the pairs); it is the dihedral group
of the rectangle. A sign can change when the pairs are exchanged and
is retained by these equations. The three lines are distinct, since
their monomial supports differ. On the unchanged original curve,
\[
 Q_{14|23}(t(z))=0,\quad
 Q_{12|34}(t(z))=2e^z\sinh(2\delta z),\quad
 Q_{13|24}(t(z))=2ie^z\sin(2\gamma z).
 \tag{FC19}
\]
All identities follow by multiplying the original exponentials. The
last two have nonzero first derivatives $4\delta$ and $4i\gamma$ at zero.
Thus equal-looking tetrahedral presentations must carry their pairing
equation as well as their point labels. To transport the whole diagram
under a permutation matrix $P$, use $t'=Pt$, $H'=HP^{-1}$ and
$Q'_0(t')=Q_0(P^{-1}t')$. Then $H't'=Ht$ and
$Q'_0(H'^{-1}X)=Q_0(H^{-1}X)$ exactly; the original conductor is unchanged.
Permuting only labels while silently fixing a different pairing would
instead change its defining equation.

\section{The singularity test in equation (11)}
Retain $d=D+1$, $s\in\{1,a\}$ with $a=4\ell+1\ge9$,
$b=s/2$, and $N=D+v$. WCF2--6 specify the original Gamma polynomial
norms and their orthonormal frames; write
\[
 \rho_n^{(s)}=\sqrt{n!/(b)_n},\qquad
 w(t)=-i\arctan t,\quad
 g(t)=t^{-v}e^{-4w(t)}E_A(w(t)),\quad
 g_0=(-i)^v\mu_v/v!.
 \tag{FC20}
\]
The source and target of the unscaled map are exactly
\[
 T_A:\mathcal P_{D+v}/\mathcal P_{v-1}\to\mathcal P_D,
 \qquad T_AP(S')=\sum_{u,w}a_{uw}P(S'+\beta_{8,u,w}),
 \tag{FC21}
\]
with centers $a/2$ and $a/2-4$ and the WCF2 quotient/target norms.
Its WCF6 matrix and determinant are
\[
 A_{D,s}=\diag(\rho_0^{(s)},\ldots,\rho_D^{(s)})^{-1}
 T_D(g)\diag(\rho_v^{(s)},\ldots,\rho_{D+v}^{(s)}),
 \tag{FC22}
\]
\[
 \det A_{D,s}=\left(\frac{(-i)^v\mu_v}{v!}\right)^{D+1}
 \prod_{j=0}^D\frac{\rho_{j+v}^{(s)}}{\rho_j^{(s)}}\ne0.
 \tag{FC23}
\]
Indeed Taylor expansion gives
$T_AS^n=\sum_{j=v}^n\binom nj\mu_j S'^{n-j}$.
It annihilates precisely $\mathcal P_{v-1}$ and reduces each remaining
leading degree by $v$ with nonzero factor $\binom nv\mu_v$.
Successive elimination proves bijectivity in FC21. In the stated frames
FC22 is triangular; its diagonal proves FC23. Since $b>0$, all displayed
weights are finite and positive. Nothing about the four labelled points
can defeat FC23 while retaining the stated orbit domain and actual
vanishing order.

For completeness, every inverse coefficient is explicitly reconstructed:
\[
 h_0=g_0^{-1},\qquad
 h_n=-g_0^{-1}\sum_{j=1}^ng_jh_{n-j},\qquad
 A_{D,s}^{-1}=\diag(\rho_v^{(s)},\ldots,\rho_{D+v}^{(s)})^{-1}
 T_D(h)\diag(\rho_0^{(s)},\ldots,\rho_D^{(s)}).
 \tag{FC24}
\]
Multiplication of the two truncated Toeplitz matrices gives the identity
by the convolution recurrence, proving this inverse with all factors.

There is a precise rank-drop locus for a \emph{fixed} old quotient.
Keep $v_0$ and $N=D+v_0$ fixed while varying coefficients and shifts
within the conductor formula. It defines
$\mathcal P_N/\mathcal P_{v_0-1}\to\mathcal P_D$ exactly when
$\mu_0=\cdots=\mu_{v_0-1}=0$. Necessity follows by applying the map to
$1,S,\ldots,S^{v_0-1}$ in that order; sufficiency follows from its Taylor
formula. At such a point of finite actual order $v\ge v_0$,
\[
 \rank T_A=\max(D+1-(v-v_0),0),\qquad
 \ker T_A=\mathcal P_{\min(v-1,N)}/\mathcal P_{v_0-1}.
 \tag{FC25}
\]
To prove this, degrees below $v$ vanish and the images of
$S^v,\ldots,S^N$ have different degrees and nonzero leading coefficients;
they span $\mathcal P_{N-v}$ when $N\ge v$. If $E_A=0$ identically the
rank is zero. If actual order is smaller than $v_0$, the old quotient
arrow is undefined, not singular. Restoring the actual kernel and actual
target degree gives FC21 again. Thus a zero of the old leading moment
is an order jump, and the full kernel correction must be retained.

The audit's inverse-exterior expression (11) has singular factors only
when the retained leading moment vanishes or a displayed weight ceases
to be defined. The latter never occurs for the admitted positive $s$;
the former contradicts using the actual order $v$. FC25 states exactly
what a genuine fixed-quotient degeneration would be. No original-period
order-jump point is supplied merely by the four Fable locale coordinates.
The point-conductor incidence test FC16 remains an explicit separate
polynomial calculation in the original $H$, with all four columns retained.

\section{Four explicit inputs to the very operator in equation (11)}
For the programme cutoffs $D\ge3$, the four labels have the following
literal receiving vectors in the WCF source, without introducing another
conductor:
\[
 x_a=[S^v\ell_a(S)]\in\mathcal P_{D+v}/\mathcal P_{v-1},\qquad
 y_a(S')=\sum_{u,w=0}^8 a_{uw}
 (S'+\beta_{8,u,w})^v\ell_a(S'+\beta_{8,u,w}).
 \tag{FC26}
\]
Thus $\mathfrak q_a\mapsto x_a\mapsto y_a=T_Ax_a$ is an explicit
map into equation (11)'s actual operator; the four-dimensional collision
coordinates $Ke_a$ map to the same $x_a$ by first applying $K^{-1}$.
All four output polynomials lie in $\mathcal P_3$.
With the source label basis in FC5 and target monomial basis
$(1,S',S'^2,S'^3)$, their four-column determinant is
\[
 \det\bigl([y_1]\ [y_2]\ [y_3]\ [y_4]\bigr)
 =\frac{\mu_v^4\displaystyle\prod_{j=0}^3\binom{v+j}{v}}
 {-64\delta^2\gamma^2(\delta^2+\gamma^2)}\ne0.
 \tag{FC27}
\]
Here the brackets on the left denote coefficient columns, not an
orthonormal change of the WCF metric. To prove the formula, evaluation
of a cubic at the four roots has matrix $V_\rho^T$. Therefore the
coefficient matrix of the four $\ell_a$ is $(V_\rho^T)^{-1}$.
On the source monomials $S^v,\ldots,S^{v+3}$, Taylor expansion in
FC21 gives a triangular output matrix with diagonal
$\mu_v\binom{v+j}{v}$. Multiplying these two determinants proves
FC27. It also proves all four outputs are linearly independent. The
source norm is the original quotient norm of the classes in FC26;
choosing their minimum-norm representatives only subtracts elements of
$\mathcal P_{v-1}$, which $T_A$ kills exactly. The target norm is the
unchanged original Gamma norm at center $a/2-4$. Thus no modification
of the norms is used to establish the determinant or the images.

The full four-vector Gram determinants can also be evaluated, so the
norm transport need not be left implicit. Set
$\Delta_\rho=-64\delta^2\gamma^2(\delta^2+\gamma^2)$ and retain
\[
 dm_s(y)=\frac{(2\pi)^{s/2}2^{s/2}}{4\pi\Gamma(s/2)}
 |\Gamma(s/4+iy/2)|^2dy,\qquad
 h_{n,s}=(2\pi)^{s/2}n!(s/2)_n.
 \tag{FC28}
\]
Let $G_x=(\langle x_a,x_b\rangle_{s,a/2,\mathrm{quot}})_{ab}$
and $G_y=(\langle y_a,y_b\rangle_{s,a/2-4})_{ab}$. Then
\[
 \det G_x=\frac{\prod_{j=0}^3 h_{v+j,s}}{|\Delta_\rho|^2},\qquad
 \det G_y=\frac{|\mu_v|^8
       \left(\prod_{j=0}^3\binom{v+j}{v}\right)^2
       \prod_{j=0}^3h_{j,s}}{|\Delta_\rho|^2},
 \tag{FC29}
\]
and in particular
\[
 \frac{\det G_y}{\det G_x}
 =\left|\frac{\mu_v}{v!}\right|^8
       \prod_{j=0}^3\frac{(\rho_{v+j}^{(s)})^2}
                              {(\rho_j^{(s)})^2}
 =|\det A_{3,s}|^2>0.
 \tag{FC30}
\]
For the proof, the original orthonormal polynomial of degree $n$ in
$S$ has leading coefficient $i^{-n}/\sqrt{h_{n,s}}$, as in WCF2.
Expanding $S^v,\ldots,S^{v+3}$ in this frame and removing the
degrees below $v$ gives a triangular matrix with diagonal
$i^{v+j}\sqrt{h_{v+j,s}}$. Its Gram determinant is
$\prod_{j=0}^3h_{v+j,s}$. Multiplication by $(V_\rho^T)^{-1}$
therefore proves the first formula in FC29. In the target, the
monomial Gram determinant is $\prod_{j=0}^3h_{j,s}$ by the same
leading-coefficient argument; FC27 proves the second formula.
Dividing the two equalities, substituting $h_{n,s}$ and the binomial
coefficients, proves FC30. Both Gram determinants retain the mass
$(2\pi)^{2s}$; it cancels only in their displayed ratio. This is the
original quotient metric on the four vectors, not a Euclidean norm
imposed on their coefficient columns.

These formulas locate the noninvertibility that survives transport:
FC3 loses the constant direction when reading the four locales together,
and the collision polynomial forgets which marked source was chosen.
FC26 instead retains that source identity through the original
conductor; FC27 excludes a four-point collision or a kernel there.
An attempt to factor FC26's four outputs through a single common image
would make all four columns equal, contradicting FC27. This proves the
failure of that specific factorization with the exact source and target
maps in hand, rather than asserting that the programmes are unrelated.

\section{Transport of the complete nonlinear collision}
In this section and the next two sections, $a$ is the first ES coordinate.
It is separate from the fixed WCF amplification integer, which is now
written $a_{\mathrm{amp}}=4\ell+1\ge9$. Thus the original source and
target centers remain $a_{\mathrm{amp}}/2$ and
$a_{\mathrm{amp}}/2-4$, and $s\in\{1,a_{\mathrm{amp}}\}$.
This dictionary changes no coordinate value or measure.
Here the receiving map is constructed, rather than inferred from a
similarity of determinants. In the original source coordinates
$q=(a,y,z,w)$, the ES reader's four-dimensional polynomial is
\[
\begin{aligned}
 P_0(q)&=a^3z+2a^2y-ia,\\
 P_2(q)&=-a^3y^2z-2ia^2yz+a^2w-2a^2y^3-10iay^2+3az+y,\\
 P_3(q)&=2a^3y^3z+6ia^2y^2z+2a^2wy+4a^2y^4
             +2iaw-4iay^3-2ayz+2iz+7y^2,\\
 P_4(q)&=2a^3y^4z+8ia^2y^3z+a^2wy^2+4a^2y^5
             +2iawy+7iay^4-10ay^2z-4iyz-w-3y^3.
\end{aligned}\tag{FC31}
\]
Write $P=(P_0,P_2,P_3,P_4)$ and $h=(0,0,-1,0)^T$.
Substitution in FC31 gives $P(Ke_a)=h$ for the four columns of FC15a,
and $h=2iKe_{\mathrm M}$. Direct differentiation and determinant
expansion give $\det DP=-2$ identically; the accompanying
\texttt{independent/verify\_four\_labels.py} performs that exact
polynomial calculation and all four evaluations, with rational
coefficients in $\mathbb Q(i,\sqrt2)$.

Let $\mathcal U=\mathcal P_{v+3}/\mathcal P_{v-1}$ and
$\mathcal W=\mathcal P_3$, as subspaces of FC21, and use the original
quotient and target Gamma norms. Define the isomorphisms
\[
 X:\C^4\to\mathcal U,\quad Xe_a=x_a,\qquad
 Y:\C^4\to\mathcal W,\quad Ye_a=y_a,\qquad
 \Phi=XK^{-1},\quad\Psi=YK^{-1}.
 \tag{FC32}
\]
FC27 proves they are isomorphisms and $T_A\Phi=\Psi$.
Consequently the following are explicit polynomial self-maps of the
\emph{actual} original four-dimensional source and target spaces:
\[
 P_{\mathcal U}=\Phi P\Phi^{-1},\qquad
 P_{\mathcal W}=\Psi P\Psi^{-1},\qquad
 T_AP_{\mathcal U}=P_{\mathcal W}T_A.
 \tag{FC33}
\]
The last equality follows by substituting $T_A\Phi=\Psi$ and
$\Phi^{-1}=\Psi^{-1}T_A$ on $\mathcal U$. Both polynomial maps have
constant Jacobian determinant $-2$: the chain rule cancels the two
constant change-of-coordinate determinants. Their norms remain those
of FC28--30. The pulled-back forms on $\C^4$ are, explicitly,
$(K^{-1})^*G_xK^{-1}$ and $(K^{-1})^*G_yK^{-1}$.

In particular the four distinct original target polynomials satisfy
\[
 P_{\mathcal W}(y_{\mathrm M})=
 P_{\mathcal W}(y_{\mathrm N})=
 P_{\mathcal W}(y_{\mathrm T})=
 P_{\mathcal W}(y_{\mathrm E})=2i\,y_{\mathrm M}.
 \tag{FC34}
\]
This is a transported noninjective nonlinear map on the conductor
target. The original linear conductor is the invertible intertwiner
in FC33; it has not acquired a kernel and none of its coefficients
has been fitted to the new map.

No inverse is hidden in these definitions: put
$\ell_a(S)=\sum_{r=0}^3\lambda_{ra}S^r$, where its coefficients are
given by the product FC6. In the target monomial basis the matrix of
$Y$ is exactly
\[
 Y_{ka}=\sum_{r=k}^3\lambda_{ra}\binom{v+r}{k}\mu_{v+r-k},
 \quad 0\le k\le3,\qquad
 P_{\mathcal W}(p)=YK^{-1}P(KY^{-1}[p]).
 \tag{FC35}
\]
Here $[p]$ is the coefficient column of $p$; the last output column is
read in the same monomial basis. The inverses are adjugates divided
by their nonzero determinants in FC15a and FC27. Formula FC35,
together with FC31, is the complete polynomial map in the original
target coordinates, with all original moments retained.

\section{The full fibre contains seven points, with all four labels retained}
The four displayed source points of FC15a are not the entire polynomial
fibre. This does not identify or remove any of the four Time locales.
We now solve every equation $P(q)=h$ over $\C$.

When $a=0$, FC31 reads
\[
 P(q)=(0,y,2iz+7y^2,-4iyz-w-3y^3).
\]
Thus $y=0,z=i/2,w=0$, the single point $q_{\mathrm M}=Ke_{\mathrm M}$.
When $a\ne0$, the equations $P_0=P_2=0$ solve uniquely as
\[
 z=\frac{i-2ay}{a^2},\qquad
 w=\frac{i(7a^2y^2-3iay-3)}{a^3}.
 \tag{FC36}
\]
Substitution into $P_3+1$ and $P_4$ gives exactly
\[
 P_3+1=-\frac{3a^2y^2-a^2+6iay-4}{a^2},\qquad
 P_4=-\frac{(ay+i)(2a^2y^2+4iay-3)}{a^3}.
 \tag{FC37}
\]
Set $t=ay$. Hence
\[
 (t+i)(2t^2+4it-3)=0,\qquad a^2=3t^2+6it-4.
 \tag{FC38}
\]
There are exactly three possible $t$: $-i$ and $-i\pm1/\sqrt2$.
For $t=-i$ the second equation gives $a^2=-1$; for the other two it
gives $a^2=1/2$. In each case both square roots are nonzero, $y=t/a$
is determined, and FC36 determines $z,w$. Thus there are exactly six
points with $a\ne0$ and the one already found with $a=0$.
All implications used above are equivalences, so no elimination root
is extraneous and none is omitted.

The exact involution
\[
 J(a,y,z,w)=(-a,-y,z,-w)
\]
satisfies, term by term in FC31,
\[
 P\circ J=J_{\mathrm{out}}\circ P,\qquad
 J_{\mathrm{out}}(u_0,u_2,u_3,u_4)=(-u_0,-u_2,u_3,-u_4).
 \tag{FC39}
\]
Since $J_{\mathrm{out}}h=h$, the full fibre is
\[
 P^{-1}(h)=\{q_{\mathrm M},q_{\mathrm N},q_{\mathrm T},q_{\mathrm E},
              Jq_{\mathrm N},Jq_{\mathrm T},Jq_{\mathrm E}\}.
 \tag{FC40}
\]
The three additional points, in the unchanged $(a,y,z,w)$ coordinates,
are
\[
\begin{aligned}
 Jq_{\mathrm N}&=(-i,1,-3i,-13),\\
 Jq_{\mathrm T}&=(-1/\sqrt2,1+\sqrt2i,2\sqrt2+6i,34+19\sqrt2i),\\
 Jq_{\mathrm E}&=(-1/\sqrt2,-1+\sqrt2i,-2\sqrt2+6i,-34+19\sqrt2i).
\end{aligned}\tag{FC41}
\]
Their first and second coordinates distinguish all seven, and FC36--38
prove their exhaustiveness.

This entire fibre transports to the original WCF target. Since
$Jq=-q+2z e_z$ and $e_z=-2i q_{\mathrm M}$, its seven points are
\[
\begin{gathered}
 y_{\mathrm M},\ y_{\mathrm N},\ y_{\mathrm T},\ y_{\mathrm E},\quad
 -12y_{\mathrm M}-y_{\mathrm N},\\
 (24-8\sqrt2i)y_{\mathrm M}-y_{\mathrm T},\quad
 (24+8\sqrt2i)y_{\mathrm M}-y_{\mathrm E}.
\end{gathered}\tag{FC42}
\]
They are precisely $P_{\mathcal W}^{-1}(2iy_{\mathrm M})$, by the
bijective conjugacy FC33. Every one is a locally invertible point of
$P_{\mathcal W}$ because its Jacobian determinant is $-2$.
The failure is a global inverse at the common target, not a zero of
the local derivative or of the linear conductor determinant.

\section{The explicit escaping branch and its original norm}
There is also a rational symmetry which leaves every output unchanged.
On the exact domain $a\ne0$, set
\[
 \Sigma(a,y,z,w)=\left(-a,\ y+\frac{2i}{a},\
 \frac{6i}{a^2}-z,\
 w-\frac{14iy^2}{a}+\frac{28y}{a^2}+\frac{40i}{a^3}\right).
 \tag{FC43}
\]
Direct substitution gives $\Sigma^2=I$ and $P\Sigma=P$ on this
domain. One can verify the latter identity with fewer terms in the
reader's original factor coordinates:
\[
\begin{aligned}
 b&=i+ay,&c&=-i+2ay+a^2z,\\
 d&=-iy-a(iz+2y^2)-a^2yz,&e&=2z-7iy^2+aw,\\
 f&=iw+3iy^3-4yz+a(6iy^2z+wy+4y^4)+2a^2y^3z.
\end{aligned}\tag{FC44}
\]
Substitution of FC43 into these formulas changes
$(a,b,c,d,e,f)$ to $(-a,-b,-c,-d,-e,-f)$.
The polynomial outputs are $(ac,ae+bd,af+be,bf)$, so all four
remain unchanged. The first four solved coordinates also show the
sign operation is an involution; alternatively it follows by inserting
FC43 twice. For $a=0$, the chart fixes $b=i$. The negative of the
$\mathrm M$ factor state has $b=-i$ and lies outside this chart.
The other three factor-state negatives are present among FC41.
This records the orientation data without treating an opposite
factorization as an identical labelled point.

The exact factor count explains why four labels coexist with seven
points in this affine fibre. A marked factorization of
$H(U,V)=UV(U-V)(U+V)$ chooses one of its four distinct linear factors
as $L=aU+bV$, and places the remaining factors in
$C=cU^3+dU^2V+eUV^2+fV^3$. The resultant is
\[
 R(L,C)=C(-b,a)=a^3f-a^2be+ab^2d-b^3c.
 \tag{FC44a}
\]
For each marked root it is nonzero, because none of the remaining
three factors vanishes on that root. Replacing $(L,C)$ by
$(\lambda L,\lambda^{-1}C)$ changes this resultant by $\lambda^2$.
Thus $R=1$ has exactly two scalar choices for each of the four
marked roots, eight factor states in all. Their signs are distinct.
All have $ad+bc=1$, the $U^3V$ coefficient of $H$.

On $a\ne0$, FC44 has the exact inverse
\[
 y=(b-i)/a,\qquad z=(c+i-2ay)/a^2,\qquad
 w=(e-2z+7iy^2)/a.
 \tag{FC44b}
\]
The incidence equation $ad+bc=1$ then determines $d$ uniquely, and
FC44a with $R=1$ determines $f$ uniquely. This proves that all six
factor states with $a\ne0$ occur in the chart. When $a=0$, those two
incidence equations force $bc=1$ and $-b^3c=1$, hence $b^2=-1$.
The chart FC44 has $b=i$ at $a=0$; its remaining coordinates are
recovered by $y=id$, $z=(e+7iy^2)/2$ and
$w=(f-3iy^3+4yz)/i$. It contains the entire $b=i$ component and none
of the $b=-i$ component. At the particular product $H$, the omitted
state is $(a,b,c,d,e,f)=(0,-i,i,0,-i,0)$. This proves precisely that
the chosen chart retains seven of the eight signed factor states,
still over the same four marked roots. No Time label has been added
or discarded.

Take real $\epsilon>0$ and the two explicit source curves
\[
 q_\epsilon=(\epsilon,0,i/2,0),\qquad
 \widehat q_\epsilon=\Sigma q_\epsilon
 =(-\epsilon,2i/\epsilon,6i/\epsilon^2-i/2,40i/\epsilon^3).
 \tag{FC45}
\]
The latter escapes to infinity whereas the former tends to
$q_{\mathrm M}$. Their identical images are
\[
 P(q_\epsilon)=P(\widehat q_\epsilon)
 =(-i\epsilon+i\epsilon^3/2,\ 3i\epsilon/2,\ -1,\ 0)
 \longrightarrow h.
 \tag{FC46}
\]
Thus $h$ is an explicit nonproper value of the nonlinear map: a
sequence escaping every compact subset has convergent images.
No numerical search is needed.

The same statement holds in the original target norm, with exact
constants. Put $z_\epsilon=\Psi\widehat q_\epsilon$ and
\[
 R=\Psi e_w=
 \frac6{77}y_{\mathrm M}+\frac1{77}y_{\mathrm N}
 +\frac{-3-\sqrt2i}{154}y_{\mathrm T}
 +\frac{3-\sqrt2i}{154}y_{\mathrm E}.
 \tag{FC47}
\]
The displayed coefficients are the last column of $K^{-1}$, verified
by multiplying by FC15a. They are not all zero, so $R\ne0$ by FC27.
Multiplying FC45 by $\epsilon^3\Psi$ gives the norm-convergent limit
\[
 \epsilon^3 z_\epsilon\longrightarrow40iR,\qquad
 \epsilon^3\|z_\epsilon\|_{s,a_{\mathrm{amp}}/2-4}
 \longrightarrow40\|R\|_{s,a_{\mathrm{amp}}/2-4}>0.
 \tag{FC48}
\]
The limit uses only a fixed finite-dimensional linear map and the
original positive Gamma norm. Simultaneously
\[
 P_{\mathcal W}(z_\epsilon)
 =\Psi(-i\epsilon+i\epsilon^3/2,3i\epsilon/2,-1,0)
 \longrightarrow2iy_{\mathrm M}.
 \tag{FC49}
\]
Equations FC42 and FC45--49 therefore identify the full finite fibre
and a branch escaping to infinity above its common value in the
actual conductor target. They do not make that introduced nonlinear
map a linear conductor or an arithmetic dilation action. The exact
relation to the original conductor is the intertwining square FC33,
which preserves its inverse bound.

\section{Compatibility with the original translation action}
The nonlinear transport can be tested against an actual symmetry of the
original conductor. It is not enough to have a conjugate polynomial
on its target. The original finite-shift formula FC21 satisfies
\[
 T_A\partial_S=\partial_{S'}T_A,\qquad
 T_A(P(S+t))=(T_AP)(S'+t).
 \tag{FC50}
\]
Both identities follow by differentiating, respectively translating,
each term of the unchanged finite sum. Differentiation preserves
$\mathcal P_{v-1}$ and therefore descends on $\mathcal U$.
In the source basis $[S^v],\ldots,[S^{v+3}]$ and target basis
$1,S',S'^2,S'^3$, its matrices are
\[
 N_v=\begin{pmatrix}0&v+1&0&0\\0&0&v+2&0\\0&0&0&v+3\\0&0&0&0\end{pmatrix},
 \qquad
 N_0=\begin{pmatrix}0&1&0&0\\0&0&2&0\\0&0&0&3\\0&0&0&0\end{pmatrix}.
 \tag{FC51}
\]
Let $B_v$ have entries
$(B_v)_{kj}=\binom{v+j}{k}\mu_{v+j-k}$ for $k\le j$ and zero
otherwise. FC35 gives $Y=B_v(V_\rho^T)^{-1}$, and FC50 gives
$N_0B_v=B_vN_v$. Thus the complete original translation matrix
in the Fable coordinates is exactly
\[
 M=\Psi^{-1}N_0\Psi
   =K V_\rho^T N_v(V_\rho^T)^{-1}K^{-1},\qquad M^4=0,
 \quad M^3\ne0.
 \tag{FC52}
\]
No moment has been set to one in this formula: their cancellation is
the proved identity $N_0B_v=B_vN_v$. They remain in $\Psi$ and in
the native metrics.

The quotient correction has an explicit form in the labelled basis.
Writing $R_v=V_\rho^T N_v(V_\rho^T)^{-1}$ gives
\[
 (R_v)_{ba}=\ell_a'(\rho_b)
          +v\frac{\delta_{ba}-\ell_a(0)}{\rho_b}.
 \tag{FC53}
\]
Indeed, modulo $\mathcal P_{v-1}$,
\[
 \partial_S(S^vf(S))
 =S^v\left(f'(S)+v\frac{f(S)-f(0)}S\right).
\]
The removed term is exactly $vf(0)S^{v-1}$, and the quotient in
parentheses is a polynomial. Evaluating at the four nonzero $\rho_b$
proves FC53. It retains the term that would be lost by treating the
quotient as an ordinary cubic space before transporting its action.

The introduced nonlinear map does not preserve this original
translation flow. From FC31,
\[
 P(0)=0,\qquad J_0:=DP(0)=\diag(-i,1,2i,-1).
 \tag{FC54}
\]
All four eigenvalues of $J_0$ are distinct. A matrix commuting with
this diagonal matrix must be diagonal, since its $(j,k)$ entry is
multiplied in the commutator by the nonzero difference of the two
eigenvalues whenever $j\ne k$. If $M$ commuted with $J_0$, it
would consequently be diagonal; $M^4=0$ would then force $M=0$,
contradicting $M^3\ne0$. Hence $J_0M-MJ_0\ne0$.

More strongly, $P e^{tM}=e^{tM}P$ fails as a polynomial identity
for every $t\ne0$ in $\C$. Otherwise differentiation at zero
would make $J_0$ commute with $R=e^{tM}$. Since $(R-I)^4=0$,
\[
 M=t^{-1}\left((R-I)-\tfrac12(R-I)^2+\tfrac13(R-I)^3\right).
\]
This identity follows by substituting
$R=I+tM+t^2M^2/2+t^3M^3/6$ and collecting powers of $M$.
It would force $J_0$ to commute with $M$, the contradiction just
proved. Conjugating by $\Psi$ proves the same failure for
$P_{\mathcal W}$ and the original target translations $e^{tN_0}$.

The entire correction is explicit, rather than an unspecified failure:
\[
 \Delta(q):=DP(q)Mq-MP(q),\qquad
 P(e^{tM}q)-e^{tM}P(q)
 =\int_0^t e^{(t-s)M}\Delta(e^{sM}q)\,ds.
 \tag{FC55}
\]
For complex $t$, integrate on the straight segment from zero to $t$.
Differentiating $e^{(t-s)M}P(e^{sM}q)$ proves FC55 by the
fundamental theorem of calculus. Every exponential truncates after
its cubic term, so the integrand and integral are finite explicit
polynomials. The linear part of $\Delta$ is precisely
$(J_0M-MJ_0)q\ne0$. The original-target correction is
$\Psi$ times the right side of FC55, with $q=\Psi^{-1}p$;
its norm is evaluated in the unchanged form
$(K^{-1})^*G_yK^{-1}$.

Thus FC33 is a complete nonlinear intertwining with the original
conductor, FC42 is its complete seven-point collision, FC48--49
is its exact nonproper branch, and FC55 is its nonzero correction
for the conductor's original translation action. None of these
calculations identifies this additional nonlinear map with the
arithmetic dilation action or makes the inverse-exterior bound false.

% BEGIN ESCAPE_DIFFERENTIAL
\section{Every exterior inverse along the escaping branch, in the original metric}
The original branch FC45 yields an exact pointwise differential calculation,
not only an unbounded sequence. Throughout this section $\epsilon>0$, the
quartet, original conductor, $a_{\mathrm{amp}}$, and $s$ are fixed. Let
$\widehat q_\epsilon$ be FC45, $z_\epsilon=\Psi\widehat q_\epsilon$,
and retain the full positive matrix
\[
 G=(K^{-1})^*G_yK^{-1},\qquad
 \|q\|_G^2=q^*Gq=\|\Psi q\|_{s,a_{\mathrm{amp}}/2-4}^2.
 \tag{FC56}
\]
No entry of this form is replaced by a Euclidean coordinate norm.
Its entries can be obtained by a finite formula in the original moments,
so FC56 does not leave an unspecified metric to be chosen. Write
$F=YK^{-1}$ in the target monomial basis, $c=a_{\mathrm{amp}}/2-4$,
$b=s/2$, and $M_s=(2\pi)^{s/2}$. Then
\[
 G=F^*\mathcal H F,\qquad
 \mathcal H_{jk}=\sum_{p=0}^j\sum_{q=0}^k
 \binom jp\binom kq c^{j+k-p-q}(-i)^p i^q m_{p+q},
 \quad 0\le j,k\le3,
 \tag{FC56a}
\]
where $m_0=M_s$, $m_2=M_sb$, $m_4=M_s(3b^2+2b)$,
$m_6=M_s(15b^3+30b^2+16b)$, and the odd moments vanish.
To prove this formula, expand
$(c-iy)^j(c+iy)^k$ in the defining measure FC28. The measure is
even. Integrating the WCF generating function in WCF3, cited in
\cite{WCF}, leaves only its constant orthogonal-polynomial coefficient,
so its moment generating function at $u=\arctan t$ is
$\int e^{uy}\,dm_s(y)=M_s(\cos u)^{-b}$ near zero. The coefficients
through order six give exactly the stated moments by differentiation.
The identity can equivalently be read coefficientwise through degree
six, using orthogonality, so no exchange of an infinite series outside
its convergence domain is needed. FC35 specifies every entry of $F$
in $\mu_v,\ldots,\mu_{v+3}$ and the unchanged quartet.
Differentiating the polynomial FC31 and substituting FC45 gives
\[
 J_\epsilon=DP(\widehat q_\epsilon)=
 \begin{pmatrix}
 9i-3i\epsilon^2/2&2\epsilon^2&-\epsilon^3&0\\
 -7i/2-30i/\epsilon^2&\epsilon^2-27&-3\epsilon&\epsilon^2\\
 -2/\epsilon+104/\epsilon^3&-i\epsilon-56i/\epsilon&-2i&2i\epsilon\\
 -4i/\epsilon^2+64i/\epsilon^4&2+28/\epsilon^2&0&-1
 \end{pmatrix},\qquad \det J_\epsilon=-2.
 \tag{FC57}
\]
Thus $D P_{\mathcal W}(z_\epsilon)=\Psi J_\epsilon\Psi^{-1}$.
The inverse matrix, including all its lower-order terms, is
\[
\renewcommand{\arraystretch}{1.25}
\begin{pmatrix}
 \frac{3i\epsilon^2}{2}-i&-\frac{3i\epsilon^4}{2}+2i\epsilon^2&
 \frac{3\epsilon^5}{2}-\frac{5\epsilon^3}{2}&\frac{3i\epsilon^6}{2}-3i\epsilon^4\\
 -\frac32+\frac2{\epsilon^2}&\frac{3\epsilon^2}{2}-3&
 \frac{3i\epsilon^3}{2}-\frac{7i\epsilon}{2}&-\frac{3\epsilon^4}{2}+4\epsilon^2\\
 \frac{9\epsilon}{4}-\frac{18}{\epsilon}+\frac{12}{\epsilon^3}&
 -\frac{9\epsilon^3}{4}+\frac{39\epsilon}{2}-\frac{24}{\epsilon}&
 -\frac{9i\epsilon^4}{4}+\frac{81i\epsilon^2}{4}-\frac{59i}{2}&
 \frac{9\epsilon^5}{4}-21\epsilon^3+35\epsilon\\
 3-\frac{138}{\epsilon^2}+\frac{120}{\epsilon^4}&
 -3\epsilon^2+140-\frac{212}{\epsilon^2}&
 -3i\epsilon^3+141i\epsilon-\frac{258i}{\epsilon}&
 3\epsilon^4-142\epsilon^2+303
\end{pmatrix}.
\tag{FC58}
\]
This is $J_\epsilon^{-1}$: multiplying FC57 by FC58 gives the identity
entry by entry. It is also $-\operatorname{adj}(J_\epsilon)/2$.
These two exact checks, rather than numerical singular-value sampling,
are recorded by \texttt{derive\_escape\_differential.py}.

For exterior powers use increasing coordinate subsets in the order
$1=a,2=y,3=z,4=w$. Let $E^{(r)}_{I,J}$ send the coordinate wedge
$e_J$ to $e_I$ and annihilate the other coordinate wedges. Expanding
the minors of FC57 proves all four finite matrix limits
\[
\begin{array}{c|c|c|c}
 r&m_r&\displaystyle\lim_{\epsilon\downarrow0}
       \epsilon^{m_r}\wedge^rJ_\epsilon&
 \displaystyle\lim_{\epsilon\downarrow0}
       \epsilon^{m_r}\wedge^rJ_\epsilon^{-1}\\\hline
 1&4&64iE^{(1)}_{4,1}&120E^{(1)}_{4,1}\\
 2&5&-672E^{(2)}_{34,12}&336E^{(2)}_{34,12}\\
 3&4&240E^{(3)}_{234,123}&32iE^{(3)}_{234,123}\\
 4&0&-2&-1/2
\end{array}.
\tag{FC59}
\]
For explicit verification, each exterior matrix has entries
$\det(J_{I,J})$ or $\det((J^{-1})_{I,J})$. Their minimum Laurent
exponents are respectively $-4,-5,-4,0$; the coefficients at those
exponents are exactly the displayed single-entry matrices. Every
remaining monomial has larger exponent. This proves the limits and
the absence of any more singular omitted entry. The accompanying
script expands every minor and checks the limits symbolically over
$\mathbb Q(i)(\epsilon)$.

Here are the exact constants in the original norm. For a principal
submatrix write $G_{I,I}$; inverses below are of the full matrix $G$.
Set
\[
 \tau_1=\sqrt{G_{44}(G^{-1})_{11}},\qquad
 \tau_2=\sqrt{\det G_{34,34}\det(G^{-1})_{12,12}}
       =\frac{\det G_{34,34}}{\sqrt{\det G}}>0.
 \tag{FC60}
\]
Both constants are strictly positive because $G$ is positive definite.
For the original wedge metric, the squared norm of $e_I$ is
$\det G_{I,I}$ and that of its coordinate dual functional $e_J^*$ is
$\det(G^{-1})_{J,J}$. The norm of a rank-one map $x\lambda$ is
$\|x\|\|\lambda\|$: the inequality is Cauchy--Schwarz for the
Riesz representative of $\lambda$, and equality is attained at that
representative. It follows that the rank-one norms in FC59 are
$64\tau_1$, $672\tau_2$, and $240\tau_1$.
For the last equality use
\[
 \det G_{234,234}=\det G\,(G^{-1})_{11},\quad
 \det(G^{-1})_{123,123}=G_{44}/\det G.
\]
These complementary-minor formulas follow from the cofactor inverse
formula for three-element minors. For the two-element formula in
FC60, block-eliminate the positive principal block $G_{34,34}$.
The complementary block of $G^{-1}$ is the inverse Schur complement,
and $\det G$ is the product of that complement's determinant with
$\det G_{34,34}$, giving the claimed identity without discarding
any cross term.

Continuity of the norm of a finite matrix with fixed positive domain
and codomain forms now proves
\[
\begin{array}{c|c|c}
r&\displaystyle\lim_{\epsilon\downarrow0}\epsilon^{m_r}
 \|\wedge^r DP_{\mathcal W}(z_\epsilon)\|&
\displaystyle\lim_{\epsilon\downarrow0}\epsilon^{m_r}
 \|\wedge^r DP_{\mathcal W}(z_\epsilon)^{-1}\|\\\hline
1&64\tau_1&120\tau_1\\
2&672\tau_2&336\tau_2\\
3&240\tau_1&32\tau_1\\
4&2&1/2
\end{array}.
\tag{FC61}
\]
Every norm here is in the original target Gamma metric and its induced
exterior metric. The pulled-back matrix $G$ is only its exact coordinate
expression; the maps $\Psi$ and $K$ have not been declared unitary.

Let $\sigma_1(\epsilon)\ge\cdots\ge\sigma_4(\epsilon)>0$
be the four singular values of this derivative in that metric. The
largest singular value of $\wedge^r J$ is the product of its largest
$r$ singular values: diagonalizing the two positive singular-value
factors gives the wedge basis with all $r$-fold products. Taking
successive ratios of FC61 therefore proves the full spectrum
\[
\begin{aligned}
 \sigma_1(\epsilon)&\sim64\tau_1\epsilon^{-4},&
 \sigma_2(\epsilon)&\sim\frac{21\tau_2}{2\tau_1}\epsilon^{-1},\\
 \sigma_3(\epsilon)&\sim\frac{5\tau_1}{14\tau_2}\epsilon,&
 \sigma_4(\epsilon)&\sim\frac{1}{120\tau_1}\epsilon^4.
\end{aligned}\tag{FC62}
\]
The product of the four displayed constants is $2$, agreeing with
the exact determinant modulus. In particular the condition number
is asymptotic to $7680\tau_1^2\epsilon^{-8}$. The degeneration is
at infinity along a proved branch; each finite point still has
nonzero Jacobian determinant.

There is also an exact inverse derivative identity in the direction
of the image curve. Differentiating FC45--46 gives
\[
 J_\epsilon
 \begin{pmatrix}-1\\-2i\epsilon^{-2}\\-12i\epsilon^{-3}\\-120i\epsilon^{-4}\end{pmatrix}
 =\begin{pmatrix}-i+3i\epsilon^2/2\\3i/2\\0\\0\end{pmatrix}.
 \tag{FC63}
\]
This identifies a specific bounded target velocity whose inverse
grows as $\epsilon^{-4}$, in addition to the full singular spectrum.
If $p_\epsilon=P_{\mathcal W}(z_\epsilon)$, put
$v_0=\Psi(-i,3i/2,0,0)^T\ne0$ and $p_*=2iy_{\mathrm M}$.
FC46 implies $(p_\epsilon-p_*)/\epsilon\to v_0$. Combining
this limit with FC48 and FC61 gives the exact distance laws
\[
 \|p_\epsilon-p_*\|^3\|z_\epsilon\|
       \longrightarrow40\|\Psi e_w\|\|v_0\|^3,
 \qquad
 \|p_\epsilon-p_*\|^4
 \|DP_{\mathcal W}(z_\epsilon)^{-1}\|
       \longrightarrow120\tau_1\|v_0\|^4.
 \tag{FC64}
\]
All norms and all constants in these two limits are the same original
Gamma norms of FC56. Injectivity of $\Psi$ proves their positivity.

\section{The exact return to audit equation (11)}
For an invertible map $B$ of a four-dimensional Hilbert space,
its singular values prove
\[
 \|\wedge^r B^{-1}\|=
       \frac{\|\wedge^{4-r}B\|}{|\det B|}.
 \tag{FC65}
\]
Applying this to FC57 gives precisely the inverse constants in FC61,
including $\|\wedge^2 B^{-1}\|=\|\wedge^2B\|/2$.
Thus the escaping branch realizes highly anisotropic distortion while
satisfying the same exterior identity used to derive audit (11).
Its forward bound is not fixed along the branch: it grows with the
explicit poles in FC59. No contradiction with that identity occurs.

The exact relation to the original conductor is stronger than applying
the identity to an unrelated matrix. Put $L=T_A|_{\mathcal U}$,
$u_\epsilon=\Phi\widehat q_\epsilon$, and
$B_{U,\epsilon}=DP_{\mathcal U}(u_\epsilon)$,
$B_{W,\epsilon}=DP_{\mathcal W}(z_\epsilon)$. Differentiating FC33
at the very same source point proves
\[
 LB_{U,\epsilon}=B_{W,\epsilon}L,\qquad
 B_{W,\epsilon}^{-1}=L B_{U,\epsilon}^{-1}L^{-1}.
 \tag{FC66}
\]
This identity retains the original source quotient metric on $\mathcal U$
and target metric on $\mathcal W$. In particular, for every $0\le r\le4$,
\[
 \frac{\|\wedge^r B_{U,\epsilon}^{-1}\|}
 {\|\wedge^r L\|\,\|\wedge^r L^{-1}\|}
 \le\|\wedge^r B_{W,\epsilon}^{-1}\|
 \le\|\wedge^r L\|\,\|\wedge^r L^{-1}\|
       \|\wedge^r B_{U,\epsilon}^{-1}\|.
 \tag{FC67}
\]
The upper bound follows from composition of exterior maps; the lower
bound follows by solving FC66 for $B_{U,\epsilon}^{-1}$. On the four-plane,
$L$ is exactly $A_{3,s}$ in the WCF orthonormal frames, so the actual
bound in audit (11) applies to the unchanged factor
$\|\wedge^rL^{-1}\|$. For fixed conductor this factor is independent
of $\epsilon$. All four pole orders therefore survive transport through
the original conductor in both directions. This supplies a quantitative
return to the audited map, without assigning to it the nonlinear map's
variable derivative or removing a kernel without its quotient.

For completeness, define the source pullback form
$G_U=(K^{-1})^*G_xK^{-1}$. Replacing $G$ in FC60 by $G_U$ proves
the exact source constants in FC61--64 as well, since the Laurent
matrices FC57--59 are unchanged. Both source and target degenerations
are now evaluated in their original norms. The constant determinant
and the nonproper inverse are compatible; the first does not control
the separate singular values, which FC62 evaluates completely.

\section{The singular flags and the exact volumes of fixed Fable flags}
The compound limits determine the directions of this distortion as well
as its size. Let $I_k=\{1,\ldots,k\}$ and
$J_k=\{5-k,\ldots,4\}$. For $k=1,2,3$, the limiting right
singular $k$-plane in the $G$ coordinates and the limiting left plane are
\[
 \mathcal R_k=G^{-1}\operatorname{span}\{e_j:j\in I_k\},\qquad
 \mathcal L_k=\operatorname{span}\{e_j:j\in J_k\}.
 \tag{FC68}
\]
To prove this, the rank-one compound limit in FC59 has image line
$\C e_{J_k}$ and covector proportional to $e_{I_k}^*$.
The Riesz vector of the coordinate covector $e_j^*$ in the original
form is $G^{-1}e_j$, because
$\langle G^{-1}e_j,x\rangle_G=e_j^*x$.
Taking wedges gives the right singular line
$\C\bigwedge_{j\in I_k}(G^{-1}e_j)$. The nonzero singular value
of each limiting rank-one map is simple; its other singular values
are zero. For convergent matrices with this property, their leading
left and right singular lines converge: a convergent subsequence of
unit maximizing vectors must attain the limiting positive norm and
therefore belongs to that unique line. Compactness of the unit sphere
and uniqueness give convergence of the full family of lines. The
Pl\"ucker line of a decomposable nonzero wedge determines its plane
by $x\wedge\omega=0$, so the singular $k$-planes have the stated
limits. Applying $\Psi$ gives these planes in the literal original
target, with its unchanged Gamma metric.

In particular, the greatest right-singular direction tends to
$\C G^{-1}e_1$, and the greatest left-singular direction tends to
$\C e_4$. Applying the inverse compound limit in FC59 shows that
the least right-singular direction tends to $\C e_4$, and its
least left-singular direction tends to $\C G^{-1}e_1$.
These are directions, not a declaration that the coordinate vectors
are orthogonal or of unit length.

There is also a complete calculation for fixed, explicitly transported
subspaces, as in the audit's endpoint Gram formula (12). Let
\[
 E_k=\Psi\operatorname{span}\{e_1,\ldots,e_k\}\subset\mathcal W,
 \quad 1\le k\le4,
 \qquad (d_1,d_2,d_3,d_4)=(64,672,240,2).
\]
The ratio of the output Gram determinant of this particular frame under
$B_{W,\epsilon}$ to its input Gram determinant is exactly
\[
 \mathcal V_k(\epsilon)=
 \frac{\det\bigl((J_\epsilon^*GJ_\epsilon)_{I_k,I_k}\bigr)}
      {\det G_{I_k,I_k}},\qquad
 \epsilon^{2m_k}\mathcal V_k(\epsilon)
 \longrightarrow d_k^2\frac{\det G_{J_k,J_k}}{\det G_{I_k,I_k}}>0.
 \tag{FC69}
\]
Indeed the numerator is the squared original wedge norm of
$(\wedge^kJ_\epsilon)e_{I_k}$. FC59 takes this wedge, after
the displayed power of $\epsilon$, to a scalar of modulus $d_k$
times $e_{J_k}$. Taking its squared norm proves FC69, while the
denominator is the exact squared norm of the input wedge. Hence
\[
 \log\mathcal V_k(\epsilon)=2m_k\log(1/\epsilon)
   +\log\left(d_k^2\frac{\det G_{J_k,J_k}}
                                {\det G_{I_k,I_k}}\right)+o(1).
 \tag{FC70}
\]
For $k=4$, this ratio is identically $4$, not merely asymptotic to it.
For $k=1,2,3$ it diverges with the specified rates, even though the
full determinant stays fixed. This computes actual subspaces and their
exact original metric constants; it is not an upper enclosure for an
unspecified endpoint. The $E_k$ are the transported Fable coordinate
flags defined here, and no equality with a different programme endpoint
is inferred. Their source partners $\Phi\operatorname{span}(e_1,\ldots,e_k)$
are carried to them by the original $L$ in FC66. Replacing $G$ by
$G_U$ gives their full source volume formulas by the same proof.

The difference between the actual coordinate-flag constant and the
optimizing singular-flag constant is itself exact. For $k=1,2,3$,
FC59--60 and FC69 give
\[
 \lim_{\epsilon\downarrow0}
 \frac{\|\wedge^k B_{W,\epsilon}\|^2}{\mathcal V_k(\epsilon)}
 =\det G_{I_k,I_k}\det(G^{-1})_{I_k,I_k}\ge1,
 \tag{FC71}
\]
with equality precisely when $G_{I_k,I_k^c}=0$. To prove this
without omitting the cross terms, write $G$ in these blocks as
$\left(\begin{smallmatrix}A&C\\C^*&D\end{smallmatrix}\right)$.
The corresponding inverse block is $(A-CD^{-1}C^*)^{-1}$.
The matrix $A^{-1/2}(A-CD^{-1}C^*)A^{-1/2}$ is positive and at
most the identity, since $CD^{-1}C^*$ is positive semidefinite.
Its determinant is at most one, which proves the inequality.
Equality makes all its eigenvalues one, hence $CD^{-1}C^*=0$.
Positive definiteness of $D^{-1}$ then gives $C=0$; the reverse
implication follows by substitution. The limit equality follows
by dividing the rank-one operator constant
$d_k^2\det G_{J_k,J_k}\det(G^{-1})_{I_k,I_k}$ by the fixed-frame
constant in FC69. Thus the exact preserved cross Gram determines
the slack in replacing these particular flag volumes by the operator
bound used in audit (11).
% END ESCAPE_DIFFERENTIAL

% BEGIN GLOBAL_FIBRE
\providecommand{\cP}{\mathcal P}
% Full independent global proof reviewed and integrated by root.
The following complete inverse calculation extends the canonical ES source
\cite{ES488}, whose original polynomial and four marked representatives
are retained. It classifies every target, then specializes the result
to the literal four roots used by the original conductor.

\section{The original polynomial and the factor coordinates}
Fix $i^2=-1$. The original coordinates are $(a,y,z,w)\in\C^4$,
and the ordered output coordinates are $(u_0,u_2,u_3,u_4)$. Define
\begin{align}
 P_0={}&a^3z+2a^2y-ia,\nonumber\\
 P_2={}&-a^3y^2z-2ia^2yz+a^2w-2a^2y^3-10iay^2+3az+y,\nonumber\\
 P_3={}&2a^3y^3z+6ia^2y^2z+2a^2wy+4a^2y^4+2iaw
       -4iay^3-2ayz+2iz+7y^2,\nonumber\\
 P_4={}&2a^3y^4z+8ia^2y^3z+a^2wy^2+4a^2y^5+2iawy
       +7iay^4-10ay^2z-4iyz-w-3y^3.
 \tag{GF1}
\end{align}
The factor-coordinate map $\Phi_{\mathrm{fac}}(a,y,z,w)=(a,b,c,d,e,f)$ is
\begin{align}
 b&=i+ay,&c&=-i+2ay+a^2z,\nonumber\\
 d&=-iy-a(iz+2y^2)-a^2yz,&e&=2z-7iy^2+aw,\nonumber\\
 f&=iw+3iy^3-4yz+a(6iy^2z+wy+4y^4)+2a^2y^3z.
 &&\tag{GF2}
\end{align}
Let
\[
 L=aU+bV,\qquad C=cU^3+dU^2V+eUV^2+fV^3.
\]
Their original incidence and resultant equations are
\begin{equation}
 ad+bc=1,\qquad
 R(L,C)=C(-b,a)=a^3f-a^2be+ab^2d-b^3c=1.\tag{GF3}
\end{equation}
Direct substitution of GF2 proves both identities. Multiplication
of the two factors gives the complete coefficient map
\begin{equation}
 LC=(ac)U^4+U^3V+(ae+bd)U^2V^2+(af+be)UV^3+(bf)V^4,
 \qquad P=(ac,ae+bd,af+be,bf).\tag{GF4}
\end{equation}
Substitution of GF2 in GF4 gives precisely GF1.

On $a\ne0$, every factor state satisfying GF3 has the unique inverse
\begin{equation}
 y=(b-i)/a,\qquad z=(c+i-2ay)/a^2,\qquad
 w=(e-2z+7iy^2)/a.\tag{GF5}
\end{equation}
These formulas recover $a,b,c,e$. The first equation GF3 uniquely
determines $d$, and its second uniquely determines $f$ because
$a^3\ne0$. Thus GF5 is a full inverse on that open set.
At $a=0$, GF3 forces $bc=1$ and $-b^3c=1$, hence $b=\pm i$.
The original chart GF2 has $b=i$, $c=-i$, and contains all such
states: its inverse there is
\[
 y=id,\qquad z=(e+7iy^2)/2,\qquad
 w=(f-3iy^3+4yz)/i.
\]
It contains none of the $a=0,b=-i$ states. This description of the
chart is used explicitly in every fibre count below.

\section{Exact inverse over every target}
For $u=(u_0,u_2,u_3,u_4)$ define
\begin{equation}
 H_u(U,V)=u_0U^4+U^3V+u_2U^2V^2+u_3UV^3+u_4V^4,
 \qquad h_u(x)=H_u(x,1).\tag{GF6}
\end{equation}
The fixed $U^3V$ coefficient is part of the original map, not a
choice made during this calculation.

\begin{theorem}[Complete fibre, including every multiple-root case]
Every finite simple root $r$ of $h_u$ contributes exactly two
distinct points of $P^{-1}(u)$. Their complete coordinates are
\begin{align}
 a^2&=1/h_u'(r),\qquad a\in\{\text{both square roots}\},\nonumber\\
 y&=-r-i/a,\nonumber\\
 z&=u_0/a^3+2r/a+3i/a^2,\nonumber\\
 w&=7ir^2/a+(u_2-17r+u_0r^2)/a^2-13i/a^3-2u_0/a^4.
 \tag{GF7}
\end{align}
A finite multiple root contributes no point. When $u_0=0$ there
is exactly one additional point, namely
\begin{equation}
 q_\infty(u)=
 \left(0,u_2,\frac{i(7u_2^2-u_3)}2,
                  11u_2^3-2u_2u_3-u_4\right).\tag{GF8}
\end{equation}
There are no other points. In particular, writing $m(u)$ for the
number of distinct finite simple roots,
\begin{equation}
 \#P^{-1}(u)=2m(u)+\mathbf1_{\{u_0=0\}}.\tag{GF9}
\end{equation}
\end{theorem}
\begin{proof}
A source point with $a\ne0$ gives the marked finite root
$r=-b/a$ of $H_u=LC$. The quotient factor is necessarily
\begin{align*}
 b&=-ar,&c&=u_0/a,&d&=(1+u_0r)/a,\\
 e&=(u_2+r+u_0r^2)/a,&
 f&=(u_3+u_2r+r^2+u_0r^3)/a.
\end{align*}
These formulas follow by comparing successive coefficients in
$H_u=a(U-rV)C$. The final constant coefficient is equivalent to
$h_u(r)=0$. Evaluating $C$ at $(-b,a)=(ar,a)$ gives
\[
 R(L,C)=a^3 C(r,1)=a^2 h_u'(r).
\]
Thus GF3 holds exactly when $a^2h_u'(r)=1$. A multiple root is
excluded by this equation, while each simple root supplies exactly
two nonzero values of $a$. The other incidence equation follows
without a further constraint:
$ad+bc=1+u_0r-u_0r=1$. The complete inverse GF5 now gives GF7;
for example substitution in $w=(e-2z+7iy^2)/a$ yields all four
displayed terms of its last line. These operations are reversible,
so they prove both sufficiency and exhaustiveness on $a\ne0$.
Distinct roots give different $r=-b/a$, and the two signs of $a$
give distinct source points.

If $a=0$ then $u_0=ac=0$. Conversely, GF1 on this hyperplane reads
\[
 P(0,y,z,w)=(0,y,2iz+7y^2,-4iyz-w-3y^3).
\]
Solving successively gives exactly GF8. In projective factor terms,
$u_0=0$ makes $[U:V]=[1:0]$ a simple root: $H_u=V(U^3+\cdots)$
and the leading coefficient of the parenthesis is exactly one.
Its two resultant-one factor scalings have $b=\pm i$, but GF2
retains only $b=i$. This is precisely the single point just solved,
regardless of any multiple finite roots. Hence there are no omitted
chart exceptions, and GF9 follows.
\end{proof}

For clarity the exhaustive multiplicity counts are
\[
\begin{array}{c|c|c}
\text{target chart}&\text{finite root multiplicities}&\#P^{-1}(u)\\\hline
u_0\ne0&1+1+1+1&8\\
u_0\ne0&2+1+1&4\\
u_0\ne0&3+1&2\\
u_0\ne0&2+2\text{ or }4&0\\\hline
u_0=0&1+1+1&7\\
u_0=0&2+1&3\\
u_0=0&3&1
\end{array}
\]
Every row follows directly from GF9. The polynomial always has
degree four or three, because the cubic coefficient is fixed at one.

\section{A short complete Jacobian proof}
On $a\ne0$, use parameters $(a,r,A,B)$, where $A=u_0$ and $B=u_2$.
The two equations $h_u(r)=0$, $a^2h_u'(r)=1$ solve as
\begin{equation}
 u_3=a^{-2}-4Ar^3-3r^2-2Br,
 \qquad u_4=3Ar^4+2r^3+Br^2-r/a^2.\tag{GF10}
\end{equation}
GF7 gives the source parameterization
\[
 q(a,r,A,B)=\left(a,-r-i/a,A/a^3+2r/a+3i/a^2,
 7ir^2/a+(B-17r+Ar^2)/a^2-13i/a^3-2A/a^4\right).
\]
Its Jacobian determinant in this parameter order is $-a^{-5}$:
after its first row, the remaining triangular minor has diagonal
$-1,a^{-3},a^{-2}$. For the target parameterization
$(A,B,u_3,u_4)$, reorder the first two parameter columns past the
last two; this permutation has positive sign. The remaining
two-by-two determinant is
\[
 \det\begin{pmatrix}
 -2/a^3&-(12Ar^2+6r+2B)\\
 2r/a^3&r(12Ar^2+6r+2B)-a^{-2}
 \end{pmatrix}=2a^{-5}.
\]
The chain rule therefore proves
\begin{equation}
 \det DP=-2\quad\text{on }a\ne0.\tag{GF11}
\end{equation}
Every original source point with $a\ne0$ has these parameters by
GF5--7, so this identity holds on that entire open set. Both sides
of GF11 are polynomials in the original coordinates; their
difference vanishes on a nonempty dense open set and is consequently
the zero polynomial. This proves GF11 everywhere, including the
$a=0$ chart, with its original constant and sign.

\section{The exact omitted image}
The only quartic multiplicity patterns with no simple root are
$2+2$ and $4$. Equivalently a quartic with $u_0\ne0$ and no simple
root is the square of a quadratic times its leading coefficient.
Because its cubic coefficient is one, it must have the form
\[
 H_u(U,V)=u_0\left(U^2+\frac{1}{2u_0}UV+u_3V^2\right)^2.
\]
Indeed, write the monic quadratic as $U^2+pUV+qV^2$; comparison
of cubic and linear coefficients gives $p=1/(2u_0)$ and $q=u_3$.
The remaining coefficients then give the exact image complement
\begin{equation}
 \C^4\setminus P(\C^4)
 =\{4u_0u_2-1-8u_0^2u_3=0,\quad u_4-u_0u_3^2=0\}.
 \tag{GF12}
\end{equation}
The first equation excludes $u_0=0$. Conversely the two equations
give exactly the displayed square. GF9 therefore proves both
inclusions in GF12, not only exclusion of the square targets.

\section{Discriminant and the genuine eight-sheet domain}
Write $A=u_0,B=u_2,C=u_3,D=u_4$. Define the original binary-quartic
discriminant polynomial by
\begin{align}
 \mathfrak D(u)={}&256A^3D^3-128A^2B^2D^2+144A^2BC^2D
 -27A^2C^4-192A^2CD^2\nonumber\\
 &+16AB^4D-4AB^3C^2-80AB^2CD+18ABC^3
 +144ABD^2-6AC^2D\nonumber\\
 &-4B^3D+B^2C^2+18BCD-4C^3-27D^2.
 \tag{GF13}
\end{align}
For $A\ne0$, the Sylvester determinant of $h_u$ and $h_u'$,
divided by $A$, expands to GF13. If $r_1,\ldots,r_4$ are its
roots counted with multiplicity, that same determinant gives
\begin{equation}
 \mathfrak D=A^6\prod_{j<k}(r_j-r_k)^2.\tag{GF14}
\end{equation}
One can verify this identity first for distinct roots: the
resultant is $A^3\prod_jh_u'(r_j)$, and
$\prod_jh_u'(r_j)=A^4\prod_{j<k}(r_j-r_k)^2$ because
$(-1)^{4\cdot3/2}=1$. Polynomial continuity gives the identity
also at multiple roots. At $A=0$, GF13 becomes
\[
 B^2C^2-4C^3-4B^3D+18BCD-27D^2,
\]
the discriminant of the monic cubic $x^3+Bx^2+Cx+D$. The same
resultant calculation in degree three proves that it vanishes
exactly at a repeated finite root. Since the projective root at
infinity is simple when $A=0$, GF13 vanishes exactly when $H_u$
has any multiple projective root, on both charts.

Let
\begin{equation}
 \mathcal B=\{u\in\C^4:u_0\mathfrak D(u)\ne0\},\qquad
 R=\C[u_0,u_2,u_3,u_4,(u_0\mathfrak D)^{-1}].\tag{GF15}
\end{equation}
Over this precise domain the map $P$ is an eight-sheet finite
unramified cover. Its complete coordinate algebra is
\begin{equation}
 S=R[r,a]/(h_u(r),\ a^2h_u'(r)-1).\tag{GF16}
\end{equation}
Here is a direct proof of each assertion. Since $u_0$ is a unit,
the first relation is a monic relation of degree four after division
by $u_0$, and $R[r]/(h_u)$ is free of rank four over $R$ with basis
$1,r,r^2,r^3$. The resultant of $h_u$ and $h_u'$ is a unit in $R$.
The adjugate of their Sylvester matrix therefore gives a Bezout
identity over $R$, so $h_u'(r)$ is a unit in this rank-four algebra.
The second relation is consequently the monic relation
$a^2=h_u'(r)^{-1}$. Thus $S$ is free of rank eight over $R$, with
basis $1,r,r^2,r^3,a,ar,ar^2,ar^3$.

The original chart inverse GF7 has entries in $S$: $a$ is a unit
with inverse $a h_u'(r)$. In the other direction, for an original
source point over $\mathcal B$ the identity $u_0=ac$ makes $a$
a unit, and $r=-b/a$ is regular. GF3--7 prove these two maps are
mutually inverse. Thus GF16 is the coordinate algebra of the
\emph{actual} preimage $P^{-1}(\mathcal B)$, rather than a new
eight-point model.

The derivative matrix of the two defining equations of GF16 with
respect to $(r,a)$ has determinant
$2a\,h_u'(r)^2$, a unit. This proves that each of its eight simple
solutions gives one local holomorphic inverse, equivalently that
the finite free map is unramified. Explicitly the four simple roots
vary holomorphically in a sufficiently small target neighbourhood,
and each nonzero $1/h_u'(r)$ has its two separate holomorphic
square roots there. GF7 gives eight disjoint local inverse charts.
There are exactly eight by GF9. These statements require neither
a chosen global square-root branch nor identification of its two signs.

\section{The complete nonproper-value locus}
A target $u$ is called a nonproper value here if there is a sequence
$q_n\in\C^4$ escaping every compact set with $P(q_n)\to u$.
For this polynomial map that is the local failure of properness
near $u$. The exact locus is
\begin{equation}
 \boxed{\quad\mathcal N(P)=\{u_0=0\}\ \cup\ \{\mathfrak D(u)=0\}.\quad}
 \tag{GF17}
\end{equation}

First take $u\in\mathcal B$ and a convergent sequence of targets
$u_n\to u$. On a sufficiently small closed neighbourhood of $u$,
$u_0$ and $\mathfrak D$ are bounded away from zero. Every finite
root is bounded: after dividing $h$ by $u_0$, any root has modulus
at most one plus the maximum modulus of its nonleading
coefficients. This follows by comparing the leading term with
the finite geometric sum of the remaining terms outside that disk.
Moreover $|h_{u_n}'(r_n)|$ is bounded away from zero uniformly over
all roots. Otherwise a convergent subsequence of the bounded
$r_n$ would give a root $r$ of $h_u$ with $h_u'(r)=0$,
contradicting $\mathfrak D(u)\ne0$. The derivatives also have a
uniform upper bound, by the root and coefficient bounds. Hence
$a^2=1/h'(r)$ bounds both $a$ and $a^{-1}$; GF7 then bounds every
original source coordinate. There is no $a=0$ solution on this
neighbourhood. Thus no sequence above $u_n$ can escape, proving
$\mathcal N(P)\subseteq\C^4\setminus\mathcal B$.

For the reverse inclusion at $u_0=0$, let GF8 be
$q_0=(0,y_0,z_0,w_0)$ and set
$q_\epsilon=(\epsilon,y_0,z_0,w_0)$ for real $\epsilon>0$.
On $a\ne0$ the exact factor-sign involution is
\begin{equation}
 \Sigma(a,y,z,w)=\left(-a,y+\frac{2i}{a},\frac{6i}{a^2}-z,
 w-\frac{14iy^2}{a}+\frac{28y}{a^2}+\frac{40i}{a^3}\right).
 \tag{GF18}
\end{equation}
To verify its invariance, substitute GF18 in GF2: $a,b,c,e$
all change sign. GF3 then forces $d$ and $f$ also to change sign
because $a\ne0$. Both factors in GF4 change sign, so
$P\Sigma=P$. The full inverse GF5 proves $\Sigma^2=I$ as well.
Now the second coordinate of $\Sigma q_\epsilon$ is
$y_0+2i/\epsilon$ and escapes, while
$P\Sigma q_\epsilon=Pq_\epsilon\to Pq_0=u$.
Every point of the entire hyperplane $u_0=0$ is consequently
nonproper, including its discriminant intersections.

Finally let $u_0\ne0$ and $\mathfrak D(u)=0$. Choose any multiple
root $r_0$ of $h_u$ and define
\[
 h_\epsilon(x)=h_u(x)+\epsilon(x-r_0),\qquad
 u_\epsilon=(u_0,u_2,u_3+\epsilon,u_4-\epsilon r_0).
\]
These preserve the original fixed cubic coefficient.
The root $r_0$ persists and has derivative exactly $\epsilon$.
For real $\epsilon>0$ choose $a=\epsilon^{-1/2}$ and apply the
complete inverse GF7 to this simple root. It gives a genuine
source point $q_\epsilon$ with $Pq_\epsilon=u_\epsilon\to u$,
while its first coordinate tends to infinity. This argument works
for every multiplicity of $r_0$ and does not require the other roots
to be simple. It proves the remaining inclusion in GF17.

\section{The literal arithmetic quartet as a marked target quartic}
Retain the original four distinct roots, with both original signs:
\[
 \rho_{\varepsilon,\eta}=\tfrac12+\varepsilon\delta+i\eta\gamma,
 \qquad\varepsilon,\eta\in\{1,-1\},\quad
 0<\delta<\tfrac12,\quad\gamma>2.
\]
Their sum is exactly two. Therefore the quartic with these literal
roots and the required original cubic coefficient one is
\begin{equation}
 H_\rho(U,V)=-\tfrac12\prod_{\varepsilon,\eta}
                    (U-\rho_{\varepsilon,\eta}V).\tag{GF19}
\end{equation}
The factor $-1/2$ is forced by the cubic coefficient and is retained
in every derivative, discriminant, and inverse below. With
$x=S-1/2$, its affine polynomial is
\[
 h_\rho(S)=-\tfrac12\bigl(x^4+2(\gamma^2-\delta^2)x^2
                              +(\delta^2+\gamma^2)^2\bigr).
\]
Expanding gives the exact target
\begin{equation}
 u_\rho(\delta)=\left(-\tfrac12,\
 \delta^2-\gamma^2-\tfrac34,\
 \gamma^2-\delta^2+\tfrac14,\
 -\tfrac1{32}-\frac{\gamma^2-\delta^2}{4}
                     -\frac{(\delta^2+\gamma^2)^2}{2}\right).
 \tag{GF20}
\end{equation}
At the marked root $\rho_{\varepsilon,\eta}$,
\begin{equation}
 h_\rho'(\rho_{\varepsilon,\eta})
       =4\delta\gamma(\varepsilon\gamma-i\eta\delta),\qquad
 a_{\varepsilon,\eta,\pm}^{2}
       =\frac1{4\delta\gamma(\varepsilon\gamma-i\eta\delta)}.
 \tag{GF21}
\end{equation}
To verify the derivative, the three differences from that root are
$2i\eta\gamma$, $2\varepsilon\delta$, and
$2\varepsilon\delta+2i\eta\gamma$; multiply them and the retained
leading coefficient $-1/2$. This gives precisely GF21. Both square
roots are present at every one of the four marked roots.

In the order $(+,+),(+,-),(-,+),(-,-)$ the original Vandermonde is
$\Delta_\rho=-64\delta^2\gamma^2(\delta^2+\gamma^2)$. Hence
\begin{equation}
 \mathfrak D(u_\rho)=\Delta_\rho^2/64
                   =64\delta^4\gamma^4(\delta^2+\gamma^2)^2>0.
 \tag{GF22}
\end{equation}
The denominator $64$ is the sixth power of the retained factor
$2$. Thus $u_\rho\in\mathcal B$, and GF7 and GF21 give all
eight original chart states over it. Unlike the earlier target
$UV(U-V)(U+V)$, it has no root at infinity and loses no sign
partner from the source chart.

For fixed $\gamma$ and $\delta\downarrow0$, GF19 tends to
\begin{equation}
 H_0=-\tfrac12\left((U-V/2)^2+\gamma^2V^2\right)^2.\tag{GF23}
\end{equation}
This target is on the omitted square surface GF12, so its complete
source fibre is empty. The eight nearby states all escape; they
do not become fewer finite chart points. Indeed GF21 gives
\[
 |a|=\bigl(4\delta\gamma\sqrt{\gamma^2+\delta^2}\bigr)^{-1/2},
 \qquad \sqrt\delta\,|a|\longrightarrow\frac1{2\gamma}.
\]
For each continuous choice of its sign branch,
\begin{equation}
 \sqrt\delta\,a_{\varepsilon,\eta,\pm}\longrightarrow
 \kappa_{\varepsilon,\pm},\qquad
 \kappa_{\varepsilon,\pm}^2=\frac1{4\varepsilon\gamma^2}.
 \tag{GF24}
\end{equation}
GF7 then gives
$y\to-(1/2+i\eta\gamma)$ and $z,w\to0$. Consequently, in
the original source coordinates,
\begin{equation}
 \sqrt\delta\,q_{\varepsilon,\eta,\pm}(\delta)
       \longrightarrow\kappa_{\varepsilon,\pm}e_a,
 \qquad e_a=(1,0,0,0)^{\mathsf T}.\tag{GF25}
\end{equation}
Each claim follows from the displayed inverse formulas; no moving
period or limiting identification of distinct root coordinates is
used to infer these limits.

\section{All eight escaping branches in one fixed original conductor norm}
Fix once and for all an admitted reference quartet with parameter
$\delta_*>0$, its original conductor, an original permitted source
order $s$, and its original target Gamma norm on $W=\cP_3$.
Let $\Psi_*:\C^4\to W$ be the exact invertible receiving map
constructed from that reference conductor in FC26--35
above. Its entries include the reference period and every
reference moment; none is varied in this section. Define
\[
 K=\begin{pmatrix}
 0&i&1/\sqrt2&1/\sqrt2\\
 0&-1&-1-\sqrt2 i&1-\sqrt2 i\\
 i/2&-3i&2\sqrt2+6i&-2\sqrt2+6i\\
 0&13&-34-19\sqrt2 i&34-19\sqrt2 i
 \end{pmatrix},\quad \det K=77\sqrt2 i/2.
\]
For complete coordinates, let $V_*=(\rho_{*,j}^n)_{0\le n\le3,1\le j\le4}$
be the reference root Vandermonde, and let $\mu_j$ and $v$ be
the original reference conductor moments and their actual first
nonzero order. In the target basis $(1,S',S'^2,S'^3)$ the complete
matrix of this fixed map is
\[
 [\Psi_*]=B_{v,*}(V_*^{\mathsf T})^{-1}K^{-1},\qquad
 (B_{v,*})_{kj}=\begin{cases}
 \binom{v+j}{k}\mu_{v+j-k},&0\le k\le j\le3,\\
 0,&k>j.
 \end{cases}
\]
The first determinant is
$\det B_{v,*}=\mu_v^4\prod_{j=0}^3\binom{v+j}{v}\ne0$;
the reference Vandermonde is nonzero and $K$ has the determinant
displayed above. This proves the asserted invertibility with all
original factors retained. The target norm in this monomial frame
is the original one, not the Euclidean norm of the column:
\[
 \|p\|_{\Gamma,*}^2=
 \int_{\mathbb R}|p(c_*+iy)|^2
 \frac{(2\pi)^{s/2}2^{s/2}}{4\pi\Gamma(s/2)}
 |\Gamma(s/4+iy/2)|^2\,dy,\qquad
 c_*=a_{\mathrm{amp},*}/2-4,
\]
where $s\in\{1,a_{\mathrm{amp},*}\}$ is the fixed original
reference source order. This integral defines a positive definite
norm on $\cP_3$: its density is positive on the real axis and a
nonzero polynomial cannot vanish on that entire axis. Finiteness is
the original WCF2 orthogonal-polynomial norm identity
$h_{n,s}=(2\pi)^{s/2}n!(s/2)_n$, applied to the finite expansion
of the polynomial in that frame, with no change of measure.
With these exact maps and norms, define
\[
 P_{W,*}=\Psi_*P\Psi_*^{-1},\qquad
 p_\delta=\Psi_*u_\rho(\delta),\qquad
 Z_{\varepsilon,\eta,\pm}(\delta)
              =\Psi_*q_{\varepsilon,\eta,\pm}(\delta).
\]
The variable $\delta$ in these formulas parametrizes an auxiliary
target path inside this \emph{fixed} operator and its \emph{fixed}
norm. It does not move the reference arithmetic period. Exact
conjugacy gives all eight fibre identities
\[
 P_{W,*}(Z_{\varepsilon,\eta,\pm}(\delta))=p_\delta.
\]
GF25 and continuity of the fixed linear map give the actual norm
asymptote, for every one of the eight branches,
\begin{equation}
 \sqrt\delta\,\|Z_{\varepsilon,\eta,\pm}(\delta)\|_{\Gamma,*}
       \longrightarrow\frac{\|\Psi_*e_a\|_{\Gamma,*}}{2\gamma}>0.
 \tag{GF26}
\end{equation}
The strict inequality follows because $\Psi_*$ is injective and
the original Gamma norm is positive definite.

The target displacement in its original coordinates is exactly
\begin{equation}
 u_\rho(\delta)-u_\rho(0)
 =\delta^2 d_\gamma-\tfrac12\delta^4e_w,
 \quad d_\gamma=(0,1,-1,\tfrac14-\gamma^2)^{\mathsf T},
 \quad e_w=(0,0,0,1)^{\mathsf T}.\tag{GF27}
\end{equation}
Subtract GF20 at zero to verify every coefficient. Therefore
$\delta^{-2}\|p_\delta-p_0\|_{\Gamma,*}
\to\|\Psi_*d_\gamma\|_{\Gamma,*}>0$.
Combining this with GF26 proves the complete original-norm rate
\begin{equation}
 \boxed{\quad
 \|p_\delta-p_0\|_{\Gamma,*}^{1/4}
 \|Z_{\varepsilon,\eta,\pm}(\delta)\|_{\Gamma,*}
 \longrightarrow
 \frac{\|\Psi_*d_\gamma\|_{\Gamma,*}^{1/4}
       \|\Psi_*e_a\|_{\Gamma,*}}{2\gamma}>0.
 \quad}\tag{GF28}
\end{equation}
The limiting point $p_0$ has no finite preimage under $P_{W,*}$,
because GF12 and the exact conjugacy transfer that assertion as
well. For completeness, the entire nonproper-value locus of this
fixed nonlinear receiving map is precisely
$\Psi_*(\{u_0\mathfrak D(u)=0\})$: invertible linear maps preserve
boundedness, escaping sequences, and convergence. This statement
concerns $P_{W,*}$; the original linear conductor remains the
invertible intertwiner defining $\Psi_*$.

GF26--28 avoid using a singular limit of a moving Vandermonde or a
moving conductor norm. Such a limit would be a separate calculation:
the construction of $\Psi_\delta$ itself assumes distinct original
roots, and its original Vandermonde vanishes at $\delta=0$.
Everything asserted here has instead been proved with the reference
operator, its full original coefficients, and its original norm held
fixed. No assertion about RH follows from introducing this nonlinear
map or this auxiliary target path.
% END GLOBAL_FIBRE

% BEGIN SIGNED_STATE_EVALUATION
\section{The eight-state evaluation and its exact sign decomposition}
The eight inverse states from GF7 give a second, intrinsic four-vector
frame in the original conductor target. This construction uses the actual
opposite-sign pairs in that fibre, rather than selecting four unrelated
vectors. Write $A=u_0$, $B=u_2$, $C=u_3$, $D=u_4$ and
$h(r)=Ar^4+r^3+Br^2+Cr+D$, retaining $A\ne0$.
For each of its four distinct roots $r_j$ choose one of the two values
$a_j^2 h'(r_j)=1$, and write $q_{j,+}=q(r_j,a_j)$ and
$q_{j,-}=q(r_j,-a_j)$ using the full inverse GF7. Put $d_j=h'(r_j)$.
Direct addition and subtraction of GF7 give
\[
 e_j:=\frac{q_{j,+}+q_{j,-}}2=
 \begin{pmatrix}0\\-r_j\\3i d_j\\
 (B-17r_j+Ar_j^2)d_j-2A d_j^2\end{pmatrix},\qquad
 o_j:=\frac{q_{j,+}-q_{j,-}}2=
 a_j\begin{pmatrix}1\\-i d_j\\A d_j^2+2r_jd_j\\
 7ir_j^2d_j-13i d_j^2\end{pmatrix}.
 \tag{SE1}
\]
For example $a_j^{-1}=a_jd_j$, $a_j^{-2}=d_j$,
$a_j^{-3}=a_jd_j^2$, and $a_j^{-4}=d_j^2$; these identities
retain every constant and sign in that addition. Thus SE1 is an exact
reconstruction $q_{j,\pm}=e_j\pm o_j$.

Let $\mathcal F_u=\C^{\{(j,+),(j,-):1\le j\le4\}}$ with its
eight separately marked basis vectors $f_{j,\pm}$, and let
$\mathsf S f_{j,\pm}=f_{j,\mp}$. The two explicit embeddings are
\[
 \iota_+e_j^{\rm coord}=(f_{j,+}+f_{j,-})/2,\qquad
 \iota_-e_j^{\rm coord}=(f_{j,+}-f_{j,-})/2.
 \tag{SE2}
\]
Their images are the $+1$ and $-1$ eigenspaces of $\mathsf S$;
every vector of $\mathcal F_u$ has a unique expression
$\iota_+\alpha+\iota_-\beta$. The factors $1/2$ in SE2 are the
explicit coefficients of these maps, not an unrecorded metric change.
Define the source-coordinate evaluation map by
$\mathcal E f_{j,\pm}=q_{j,\pm}$ and matrices
$E=(e_1\ e_2\ e_3\ e_4)$, $O=(o_1\ o_2\ o_3\ o_4)$.
Then the whole eight-state map is exactly
\[
 \mathcal E(\iota_+\alpha+\iota_-\beta)=E\alpha+O\beta.
 \tag{SE3}
\]

\section{The sign differences are a basis on the literal quartet}
The rank of these two pieces can be proved in the original arithmetic
quartet, without dropping the phases in $a_j$. For this purpose form
the four polynomials
\[
 v(r)=(1,-ih'(r),A h'(r)^2+2rh'(r),
                         7ir^2h'(r)-13ih'(r)^2)^T.
\]
Reduce each polynomial modulo the original $h$ by polynomial division;
this changes no evaluation at a root. Write the remainders as
$v(r)=V(1,r,r^2,r^3)^T$ modulo $h$. Their full coefficient matrix is
\[
\begin{aligned}
 V_{1,*}&=(1,0,0,0),\\
 V_{2,*}&=(-iC,-2iB,-3i,-4iA),\\
 V_{3,*}&=(AC^2-9D,\ 4ABC-8AD-7C,\
             -16A^2D+4AB^2-2AC-5B,\ -8A^2C+4AB-3),\\
 V_{4,*}&=(-13iC^2+20iD/A,\ -52iBC+76iD+20iC/A,\
             208iAD-52iB^2+5iC+20iB/A,\
             104iAC-66iB+20i/A).
\end{aligned}\tag{SE4}
\]
This is polynomial division by a quartic with its retained leading
coefficient $A$; the denominators in the last row have not been removed.
For the literal $u_\rho$ in GF20, substitution and expansion of the
three-by-three determinant remaining after the first row give
\[
 \det V=448\delta^2\gamma^2
                   (2\delta^2\gamma^2-\delta^2+\gamma^2)>0.
 \tag{SE5}
\]
The inequality follows since the last factor is
$\gamma^2+\delta^2(2\gamma^2-1)>0$ for $\gamma>2$.
The exact polynomial remainder and determinant calculations in SE4--5
are reproduced by \texttt{derive\_signed\_evaluation.py}; they use
rational function arithmetic, not a sampled determinant sign.

Write $\Delta=\det(r_j^{n})_{0\le n\le3,1\le j\le4}$ in the
same original marking, and retain
\[
 \chi=A^2\Delta\prod_{j=1}^4a_j\in\{1,-1\}.
 \tag{SE6}
\]
Indeed $\prod_jh'(r_j)=A^4\Delta^2$ by multiplying all pair
differences; the six sign reversals have positive product. Squaring
SE6 and using $a_j^2=1/h'(r_j)$ proves $\chi^2=1$. No particular
choice of its sign has been imposed. Evaluating SE4 at the roots gives
$O=V(r_j^n)_{n,j}\diag(a_1,a_2,a_3,a_4)$, hence
\[
 \det O=\frac{\chi}{A^2}\det V
   =1792\chi\delta^2\gamma^2
                   (2\delta^2\gamma^2-\delta^2+\gamma^2)\ne0
 \tag{SE7}
\]
on the literal quartet, where $A=-1/2$ is retained. Thus the four
opposite-sign differences are an actual basis in the original source
coordinates for every admitted quartet, for either choice of each sign.

The pair means have rank exactly three. Their first coordinate is zero.
For a direct lower bound on rank, let $E_{\mathrm{pol}}$ be the coefficient
matrix of the four polynomials in the first column expression of SE1,
again reduced modulo the original $h$. The minor with rows $y,z,w$
and columns corresponding to $r,r^2,r^3$ is
\[
 6i\bigl(64A^3D-16A^2B^2-22A^2C+89AB-30\bigr)
 =6i\bigl(16\delta^2\gamma^2-35\delta^2+35\gamma^2\bigr)\ne0.
 \tag{SE8}
\]
The equality follows by substituting GF20. Positivity of the last
real factor follows from
$16\delta^2\gamma^2+35(\gamma^2-\delta^2)>0$ under the original
bounds $0<\delta<1/2$, $\gamma>2$. Evaluation at four distinct
roots multiplies $E_{\mathrm{pol}}$ by the invertible Vandermonde,
so its rank and the rank of $E$ are both exactly three.

The remaining even kernel line is also explicit. The full coefficient
matrix of the means has rows
\[
 (0,0,0,0),\quad(0,-1,0,0),\quad(3iC,6iB,9i,12iA),\quad
 (w_0,w_1,w_2,w_3),
\]
where
\[
\begin{aligned}
 w_0&=-2AC^2+BC+71D,&w_1&=-8ABC+12AD+2B^2+54C,\\
 w_2&=32A^2D-8AB^2+AC+40B,&w_3&=16A^2C-6AB+20.
\end{aligned}\tag{SE8a}
\]
For $W=(r_j^n)_{n,j}$ define
\[
 t=(3w_3-4Aw_2,\ 0,\ 4Aw_0-Cw_3,\ Cw_2-3w_0)^T,
 \qquad \lambda=W^{-1}t.
 \tag{SE8b}
\]
The first coordinate of $t$ is nonzero by SE8. Substitution in
the three nonzero rows of SE8a gives zero: the last two equations
are the cross-product cancellation between $(C,3,4A)$ and
$(w_0,w_2,w_3)$; the second coordinate of $t$ is zero. Consequently
\[
 \ker E=\C\lambda,\qquad
 \lambda_j=
 \frac{t_0(C+Br_j+r_j^2+Ar_j^3)+t_2(1+Ar_j)+t_3A}{h'(r_j)}.
 \tag{SE8c}
\]
To verify the last formula, the $j$th row of $W^{-1}$ is the
coefficient row of $h(r)/((r-r_j)h'(r_j))$. Polynomial division
gives the four coefficients
$h'(r_j)^{-1}(C+Br_j+r_j^2+Ar_j^3,B+r_j+Ar_j^2,1+Ar_j,A)$.
Pairing with $t$, whose $r$ coefficient is zero, gives SE8c.
Thus the one even direction lost by evaluation is calculated
in the original labelled coordinates, not inferred to be the
constant direction from the separate tetrahedral altitude map.

\section{The entire eight-state kernel and the original-metric return}
Because $O$ is invertible on the literal quartet, SE3 gives the complete
kernel, with no unspecified missing sign or root state:
\[
 \ker\mathcal E=
 \{\iota_+\alpha-\iota_-O^{-1}E\alpha:\alpha\in\C^4\},
 \qquad
 \mathcal E\,\iota_-O^{-1}=I_{\C^4}.
 \tag{SE9}
\]
The inverse here is exactly $\operatorname{adj}(O)/\det O$, with
SE1 and SE7 specifying numerator and nonzero denominator. The kernel
therefore has dimension four. Its intersection with the even subspace
has dimension one by SE8; its intersection with the odd subspace is zero
by SE7. These assertions follow by solving $E\alpha+O\beta=0$;
they do not identify the even kernel with a different tetrahedral
constant direction by analogy.

The receiving map to the actual conductor target is
$\mathcal E_W=\Psi\mathcal E$, and its source partner is
$\mathcal E_U=\Phi\mathcal E$, with the exact original maps FC32.
Thus
\[
 T_A\mathcal E_U=\mathcal E_W,\quad
 \ker\mathcal E_U=\ker\mathcal E_W=\ker\mathcal E,\quad
 \mathcal E_W\,\iota_-O^{-1}\Psi^{-1}=I_{\mathcal W}.
 \tag{SE10}
\]
The last identity constructs a complete section into the sign-odd
four-plane of the actual eight-state space. Changing which partner
is named $+$ changes the relevant column of $O$ and of $\iota_-$
by the same sign; it leaves this section unchanged. Permuting the
four labels likewise permutes these two maps inversely, so no root
label is implicitly forgotten.

The four original target vectors $\Psi o_j$ have the full Gram
$O^*GO$, with $G$ from FC56--56a. Consequently
\[
 \det(O^*GO)=1792^2\delta^4\gamma^4
       (2\delta^2\gamma^2-\delta^2+\gamma^2)^2\det G>0.
 \tag{SE11}
\]
This retains the entire original Gamma measure and all its cross
terms. Their source Gram is $O^*G_UO$; therefore the exact ratio is
\[
 \frac{\det(O^*GO)}{\det(O^*G_UO)}
 =\frac{\det G}{\det G_U}=|\det A_{3,s}|^2.
 \tag{SE12}
\]
The first equality cancels the same nonzero explicitly calculated
$|\det O|^2$ on both sides, while the second is FC30 with the
unchanged $K$ factors. This is a basis constructed from the actual
eight signed inverse states and transported through the original
linear conductor. Its full sign-odd part is invertible; the rank-three
loss in the pair means and the four-dimensional graph kernel SE9
record exactly what the sign-forgetting evaluation removes.

\section{An invertible total state map retaining both labels and evaluation}
The graph kernel can be resolved by retaining the original root-label
vector, rather than discarding the nonlinear evaluation. Define
\[
 \mathcal J_U(\iota_+\alpha+\iota_-\beta)
       =(X\alpha,\ \Phi(E\alpha+O\beta))\in\mathcal U\oplus\mathcal U,
 \qquad
 \mathcal J_W(\iota_+\alpha+\iota_-\beta)
       =(Y\alpha,\ \Psi(E\alpha+O\beta))\in\mathcal W\oplus\mathcal W.
 \tag{SE13}
\]
Both are isomorphisms of the complete eight-dimensional state space.
For the target, the entire inverse is
\[
 \mathcal J_W^{-1}(p_1,p_2)=
 \iota_+Y^{-1}p_1+
 \iota_-O^{-1}(\Psi^{-1}p_2-EY^{-1}p_1).
 \tag{SE14}
\]
Substitution in SE13 proves each inverse composition. Replacing
$Y,\Psi$ by $X,\Phi$ proves the source inverse. Moreover the
original conductor itself supplies the exact map between these
total state spaces:
\[
 (T_A|_{\mathcal U}\oplus T_A|_{\mathcal U})\,\mathcal J_U
      =\mathcal J_W.
 \tag{SE15}
\]
This is direct substitution of $T_AX=Y$ and $T_A\Phi=\Psi$.
In particular, the two original sign states of root $j$ become
$(y_j,\Psi q_{j,+})$ and $(y_j,\Psi q_{j,-})$ in the same
target direct sum. The root coordinate and the nonlinear input
are both retained.

The exact action of signed root permutations on this object includes
the terms that evaluation alone would lose. For a permutation $\pi$
and signs $\sigma_j\in\{1,-1\}$, let
$g f_{j,t}=f_{\pi(j),\sigma_jt}$, let $\Pi e_j=e_{\pi(j)}$,
and let $R e_j=\sigma_j e_{\pi(j)}$. Thus
$g\iota_+=\iota_+\Pi$ and $g\iota_-=\iota_-R$.
In coordinates $(\alpha,z)$ for the target pair $(Y\alpha,\Psi z)$,
SE13--14 give the full transported action
\[
 \widehat g=
 \begin{pmatrix}
  \Pi&0\\ E\Pi-ORO^{-1}E&ORO^{-1}
 \end{pmatrix}.
 \tag{SE16}
\]
Indeed $z=E\alpha+O\beta$ gives
$z'=E\Pi\alpha+OR\beta
=ORO^{-1}z+(E\Pi-ORO^{-1}E)\alpha$.
The off-diagonal block is therefore evaluated exactly, with all
coefficients given in SE1 and SE7. Since SE16 is
$\mathcal J_W g\mathcal J_W^{-1}$ in the displayed coordinates,
it obeys the group law for every pair of signed permutations.
The actual monodromy subgroup is determined in the signed-root
continuation; this formula retains its full action as well.
The source uses the identical block matrix in coordinates
$(X\alpha,\Phi z)$, so SE15 intertwines every such action.

Every nontrivial even sign change has a nonzero lower-left block,
and its rank is exact. Put $F=O^{-1}E$. The image of $E$ is the
entire first-coordinate-zero hyperplane by SE8, and the first row
of $O$ is $(a_1,a_2,a_3,a_4)$. Hence
$\operatorname{im}F=\{x:\sum_j a_jx_j=0\}$.
Each $a_j$ is nonzero. Projection of this hyperplane onto any
proper subset of coordinates is surjective: choose one omitted
coordinate to solve the single displayed equation. For a diagonal
sign action $D_J$ negating the subset $J$, the lower-left block
of SE16 is $O(I-D_J)F$. Therefore
\[
 \operatorname{rank}\bigl(E-OD_JO^{-1}E\bigr)
 =\begin{cases}2,&|J|=2,\\3,&|J|=4.\end{cases}
 \tag{SE18}
\]
For two coordinates this is the just-proved surjective projection;
for all four it is twice a rank-three matrix. These are all
nonidentity elements of the even-sign kernel. Thus the extra
label coordinate in SE13 retains a necessary, quantitatively
specified mixing term. SE16 supplies the full action including it.

Put the original orthogonal direct-sum Gamma norms on the two target
copies and the original quotient direct-sum norms on the two source
copies. In the coordinates just used their Grams are
$\diag(G_y,G)$ and $\diag(G_x,G_U)$, respectively. This choice
retains each original mass and cross Gram within its copy; it does
not identify either $\mathcal J$ with an isometry from a previously
assigned norm on the abstract finite state space. For $L=T_A|_U$
and every $0\le r\le8$, the actual inverse-exterior norm is
\[
 \left\|\wedge^r(L\oplus L)^{-1}\right\|
 =\max_{\substack{p+q=r\\0\le p,q\le4}}
          \|\wedge^pL^{-1}\|\,\|\wedge^qL^{-1}\|.
 \tag{SE17}
\]
To prove it, the orthogonal direct sum decomposes its $r$th exterior
space into the orthogonal summands
$\wedge^p\mathcal W\otimes\wedge^q\mathcal W$, $p+q=r$.
The map on each summand is the tensor product of the displayed
two exterior inverses; in singular bases its norm is the product
of their largest singular values. The norm of a block diagonal
map on an orthogonal direct sum is the maximum of block norms,
proving SE17. Each factor is the original four-plane $A_{3,s}$
inverse appearing in audit (11). This proves the full eight-state
return to that bound without deleting the rank-three mean part,
the four root labels, either factor sign, or the mixing block.
% END SIGNED_STATE_EVALUATION

% BEGIN SIGNED_MONODROMY
\providecommand{\Disc}{\operatorname{Disc}}
\providecommand{\sgn}{\operatorname{sgn}}
\providecommand{\Aut}{\operatorname{Aut}}
% Root-reviewed complete independent signed-root monodromy proof.
The next calculation determines the symmetry of the actual cover GF16,
with the original ES map and four marked labels retained \cite{ES488}.

\section{The actual cover and its eight labelled states}
For $u=(u_0,u_2,u_3,u_4)$ put
\[
 H_u(U,V)=u_0U^4+U^3V+u_2U^2V^2+u_3UV^3+u_4V^4,
 \qquad h_u(x)=H_u(x,1),\qquad A=u_0.
\]
The original cubic coefficient remains one throughout. Let
\[
 \mathcal B=\{u:A\Disc(H_u)\ne0\},\qquad
 R=\C[u_0,u_2,u_3,u_4,(A\Disc H_u)^{-1}].
\]
Over this domain the original source is isomorphic to the finite
cover with coordinate algebra
\begin{equation}
 S=R[r,a]/(h_u(r),\ a^2h_u'(r)-1).\tag{SM1}
\end{equation}
In the unchanged original source coordinates its inverse is
\begin{align}
 y&=-r-i/a,\nonumber\\
 z&=A/a^3+2r/a+3i/a^2,\nonumber\\
 w&=7ir^2/a+(u_2-17r+Ar^2)/a^2-13i/a^3-2A/a^4.
 \tag{SM2}
\end{align}
These formulas are obtained by factoring
$H_u=(aU+bV)(cU^3+dU^2V+eUV^2+fV^3)$, setting $r=-b/a$,
and retaining the resultant equation
$C(-b,a)=a^2h_u'(r)=1$. Thus SM1 describes the actual inverse
states; its two signs are two distinct original source points.

For completeness, $R[r]/(h_u)$ is free of rank four because $A$
is a unit and $A^{-1}h_u$ is monic. The discriminant condition
and the Sylvester Bezout identity make $h_u'(r)$ a unit. The
second equation of SM1 is therefore the monic relation
$a^2=h_u'(r)^{-1}$, and $S$ is free of rank eight over $R$.
The Jacobian of the two equations with respect to $(r,a)$ has
determinant $2a h_u'(r)^2$, a unit. In a sufficiently small target
neighbourhood the four simple roots and their two nonzero square
roots therefore give eight separate holomorphic inverse charts.
This provides the local path-lifting property used below.

Both the base and the actual cover are path connected. The base
is the complement in $\C^4$ of the zero set of the nonzero
polynomial $A\Disc H_u$. The actual source is the complement
of the zero set of its polynomial pullback under the original map
$P$. This pullback is nonzero, since the eight states of any point
in $\mathcal B$ exist by SM1--2. To connect two points in either
complement, restrict its defining polynomial to the complex line
through them. Its value at the endpoints is nonzero, so the
restricted polynomial is not identically zero and has only
finitely many zeros. A path in the complex line parameter from
zero to one, with small detours around these finitely many points,
connects the endpoints within the complement. This proves the
asserted connectedness without an unproved topological
classification of the base.

\section{The exact upper bound from the retained product invariant}
At a base point order the four roots as $r_1,r_2,r_3,r_4$ and
choose one of the two actual values $a_j$ at each root. Denote
the eight states by $(j,+)$ and $(j,-)$, corresponding respectively
to $a_j$ and $-a_j$ in SM2. During continuation the four roots
remain distinct, and their scalings remain nonzero. Put
\[
 \Delta=\prod_{j<k}(r_k-r_j),\qquad
 \chi=A^2\Delta\prod_{j=1}^4a_j.
\]
The complete derivative product is
\[
 \prod_{j=1}^4h_u'(r_j)
 =A^4\prod_{j=1}^4\prod_{k\ne j}(r_j-r_k)
 =A^4\Delta^2.
\]
There are six pairs, so the sign from reversing each pair is
$(-1)^6=1$. Consequently
\begin{equation}
 \chi^2=1.\tag{SM3}
\end{equation}
The continuous value $\chi$ is therefore constant along every
ordered lift of a path. Its initial value may be $1$ or $-1$;
neither value has been discarded.

A loop in $\mathcal B$ induces a root permutation $\pi\in S_4$
and signs $\sigma_j\in\{1,-1\}$ such that
$r_j\mapsto r_{\pi(j)}$ and $a_j\mapsto\sigma_j a_{\pi(j)}$.
The negative state continues as the negative of the positive one,
by uniqueness in the equation $a^2h'(r)=1$. The exact action is
therefore
\begin{equation}
 (j,\epsilon)\longmapsto(\pi(j),\sigma_j\epsilon),
 \qquad \epsilon\in\{1,-1\}.\tag{SM4}
\end{equation}
The final Vandermonde is $\sgn(\pi)\Delta$. Applying the
constancy of SM3 gives the necessary relation
\begin{equation}
 \prod_{j=1}^4\sigma_j=\sgn(\pi).\tag{SM5}
\end{equation}
Thus the monodromy is contained in the explicitly defined group
\begin{equation}
 G=\{(\sigma,\pi)\in\{\pm1\}^4\rtimes S_4:
                       \prod_j\sigma_j=\sgn(\pi)\}.
 \tag{SM6}
\end{equation}
For each permutation exactly eight sign choices satisfy SM5, so
$|G|=8\cdot24=192$. As signed permutation matrices its elements
have determinant $(\prod_j\sigma_j)\sgn(\pi)=1$.

\section{Explicit local exchanges realizing the full group}
Use the real roots
\[
 (r_{\mathrm T},r_{\mathrm E},r_{\mathrm M},r_{\mathrm N})
 =(7,9,11,13),\qquad A=-1/40.
\]
Their sum is forty; hence the cubic coefficient of
$A\prod(x-r_j)$ is exactly one. The base target is
\begin{equation}
 u_*^{\rm real}=(-1/40,-59/4,95,-9009/40)\in\mathcal B.
 \tag{SM7}
\end{equation}
For the next calculation index the real order as
$s_1=7,s_2=9,s_3=11,s_4=13$; these indices refer respectively
to the unchanged labels $\mathrm T,\mathrm E,\mathrm M,\mathrm N$.
The four derivatives are
\[
 (h'(s_1),h'(s_2),h'(s_3),h'(s_4))
       =(6/5,-2/5,2/5,-6/5).
\]
Choose the exact base scalings
\begin{equation}
 (a_{\mathrm T},a_{\mathrm E},a_{\mathrm M},a_{\mathrm N})
 =\left(\sqrt{5/6},-i\sqrt{5/2},-\sqrt{5/2},i\sqrt{5/6}\right),
 \tag{SM8}
\end{equation}
where the real square roots displayed are positive. Each satisfies
the original equation in SM1.

For each adjacent pair $(s_j,s_{j+1})$, let
$c=(s_j+s_{j+1})/2$ and $d=1$. Keeping the other two roots fixed,
take the explicit path
\begin{equation}
 s_j(\theta)=c-d e^{i\theta},\qquad
 s_{j+1}(\theta)=c+d e^{i\theta},\qquad0\le\theta\le\pi.
 \tag{SM9}
\end{equation}
The pair sum remains $2c$, so the total root sum and the original
coefficient $A=-1/40$ are unchanged. The two moving roots remain
distinct. The closed disk of radius one about $c$ contains no
other root, so no collision with a fixed root occurs. At the
endpoint the pair is interchanged. The unordered roots and all
polynomial coefficients consequently return to the same target;
SM9 is a loop in the actual base $\mathcal B$.

We compute its signs without assuming a braid-group presentation.
For a fixed root $s_k$ outside the pair disk, the moving contribution
to $h'(s_k)$ is
\[
 (s_k-s_j(\theta))(s_k-s_{j+1}(\theta))
       =(s_k-c)^2-d^2e^{2i\theta}.
\]
Its values lie in the disk with positive real centre $(s_k-c)^2$
and strictly smaller radius $d^2$. This disk is in the right
half-plane and has a single continuous logarithm. After the loop
that logarithm returns to its initial value. Therefore neither
scaling at a fixed root changes sign.

For a moving root, its difference from the other moving root turns
by exactly $\pi$. Each of its two differences from an outside
root lies in a disk not containing zero, with both endpoint values
real and of the same sign. A logarithm on the corresponding left
or right half-plane shows that the net imaginary part of each of
those logarithm differences is zero. The argument change of the
complete moving derivative is thus exactly $\pi$. Its inverse
square root changes phase by $-\pi/2$, with its magnitude taking
the new root's required magnitude. By the explicit choices SM8,
the resulting permutation is the four-cycle
\begin{equation}
 \beta_j=(j+\quad(j+1)+\quad j-\quad(j+1)-).
 \tag{SM10}
\end{equation}
Equivalently, in the four retained Time labels the three generators
are
\begin{equation}
 \begin{split}
 \beta_1&=(\mathrm T+\ \mathrm E+\ \mathrm T-\ \mathrm E-),\\
 \beta_2&=(\mathrm E+\ \mathrm M+\ \mathrm E-\ \mathrm M-),\\
 \beta_3&=(\mathrm M+\ \mathrm N+\ \mathrm M-\ \mathrm N-).
 \end{split}\tag{SM11}
\end{equation}
Every sign and every fixed state has been determined by the path.

The square $\beta_j^2$ flips exactly the two adjacent root signs.
These three pair flips generate all even sign patterns: as vectors
over $\mathbb F_2$ they are
$(1,1,0,0)$, $(0,1,1,0)$ and $(0,0,1,1)$, which are independent
and span the dimension-three subspace with coordinate sum zero.
The projections of the three $\beta_j$ are the three adjacent
transpositions, and they generate all of $S_4$ by successively
moving entries to any prescribed positions. Thus the group
generated by these actual loops contains all eight elements in
the sign kernel and projects onto all 24 permutations. It has
at least $8\cdot24=192$ elements. The universal upper bound SM6
has exactly 192 elements, proving equality:
\begin{equation}
 \boxed{\operatorname{Mon}(P^{-1}(\mathcal B)/\mathcal B)=G.}
 \tag{SM12}
\end{equation}
This proof establishes the entire group without an unproved claim
that every base loop is a braid word.

\section{Transport from the literal original quartet}
Start with the original marking
\[
 (\rho_{\mathrm M},\rho_{\mathrm N},\rho_{\mathrm T},\rho_{\mathrm E})
 =\left(\tfrac12+\delta+i\gamma,\tfrac12+\delta-i\gamma,
        \tfrac12-\delta+i\gamma,\tfrac12-\delta-i\gamma\right),
 \qquad0<\delta<\tfrac12,\quad\gamma>2.
\]
Its target is $H_\rho=-\frac12\prod(U-\rho_jV)$, with the
original cubic coefficient one. The following two paths give a
fully explicit base transport to SM7. First, for $0\le t\le1$,
put
\begin{equation}
 \begin{array}{ll}
 r_{\mathrm M}(t)=\tfrac12+(1-t)(\delta+i\gamma)+t,
 &r_{\mathrm N}(t)=\tfrac12+(1-t)(\delta-i\gamma)+3t,\\
 r_{\mathrm T}(t)=\tfrac12+(1-t)(-\delta+i\gamma)-3t,
 &r_{\mathrm E}(t)=\tfrac12+(1-t)(-\delta-i\gamma)-t.
 \end{array}\tag{SM13}
\end{equation}
Their sum is always two, so $A=-1/2$ throughout this segment.
Roots with different original imaginary signs have different
imaginary parts for $t<1$. Roots with the same original imaginary
sign have real difference $2\delta(1-t)+4t$, which is positive.
At $t=1$ the four real roots are $(3/2,7/2,-5/2,-1/2)$ and
are distinct. Thus the entire segment remains in $\mathcal B$.

Next translate all four roots by $19t/2$, $0\le t\le1$, and
put $A(t)=-1/(2+38t)$. Their differences remain fixed, their sum
is $2+38t\ne0$, and the retained original cubic coefficient is
again one. At its endpoint the marked roots are
$(11,13,7,9)$, namely SM7 in the original Time order. This
second segment also lies entirely in $\mathcal B$.

The actual signs transport by the original equations, not by
discarding a branch. An exact formula along either differentiable
segment is
\begin{equation}
 a_j(t)=a_j(0)\exp\left[-\frac12\int_0^t
 \left(\frac{A'(s)}{A(s)}+
 \sum_{k\ne j}\frac{r_j'(s)-r_k'(s)}{r_j(s)-r_k(s)}\right)ds\right].
 \tag{SM14}
\end{equation}
The integrand has no singularity on these paths. Differentiating
shows that $a_j(t)^2A(t)\prod_{k\ne j}(r_j(t)-r_k(t))$ is
constant, and its initial value is one. Hence SM14 is precisely
the lifted inverse branch. Both initial signs lift separately.
Using SM8 as the signs at the real endpoint defines their exact
transported signs at the original quartet; any other initial
sign convention conjugates SM11 by a diagonal sign permutation.
It does not change SM12, because the determinant-one subgroup
is invariant under that conjugation.

Conjugating each loop SM9 by the explicitly constructed path
SM13 and its translation segment gives based loops at the literal
quartet with the same marked-root action. This transport concerns
the auxiliary quartic targets of the fixed nonlinear map. It does
not assert that intermediate roots are arithmetic zeta zeros or
that an original conductor period is moved along these paths.

\section{Why this group is not the Weyl group of type $D_4$}
With composition acting first by the rightmost generator, the
product $\beta_1\beta_2\beta_3$ is the eight-cycle
\begin{equation}
 (\mathrm T+\ \mathrm E+\ \mathrm M+\ \mathrm N+\
   \mathrm T-\ \mathrm E-\ \mathrm M-\ \mathrm N-).
 \tag{SM15}
\end{equation}
This follows by applying the three four-cycles SM11 to each
state in order. It has order eight.

For comparison, the reflections in the roots
$\{\pm e_j\pm e_k:j\ne k\}$ generate precisely the signed
permutations with an even total number of sign changes. Indeed
the reflection in $e_j-e_k$ swaps those coordinates, and the
reflection in $e_j+e_k$ swaps them with both signs changed.
Their compositions generate every permutation and every pair
sign flip, and each such reflection has even total sign product.
This proves the usual signed-permutation description of the
Weyl group $W(D_4)$ directly:
\[
 W(D_4)=\{(\sigma,\pi):\prod_j\sigma_j=1\}.
\]
This group has no element of order eight. An underlying
permutation of four coordinates can have cycle lengths only
$1,2,3,4$. A signed cycle of length $m$ has order $m$ if its
sign product is positive and $2m$ if its sign product is
negative, as follows by applying the cycle $m$ times to each
coordinate. With no length-four cycle, the available orders
combine to divisors of four or six. With a length-four cycle,
the even total sign condition makes its sign product positive,
so its order is four. Thus order eight never occurs in $W(D_4)$,
whereas SM15 occurs in $G$. The groups are therefore not even
abstractly isomorphic, despite their equal order 192.

\section{Exact deck groups}
Let $\tau$ denote simultaneous reversal of all four root signs,
$\tau(j,\epsilon)=(j,-\epsilon)$. The centralizer of $G$ in
the full permutation group of the eight inverse states is
\begin{equation}
 C_{S_8}(G)=\{1,\tau\}.\tag{SM16}
\end{equation}
Here is a complete proof. Fix one state, say $(1,+)$. Its
stabilizer in $G$ also fixes $(1,-)$. No state at another root
is fixed by that stabilizer: a pair sign flip involving that
root and one of the other two roots different from root 1
belongs to the stabilizer and moves that state. Hence the
stabilizer has exactly the two fixed states $(1,+),(1,-)$.
Any permutation $c$ commuting with $G$ must send $(1,+)$ to
one of these fixed states. The group $G$ is transitive on all
eight states (permutations move root labels and pair flips
change either chosen sign), so the image of this one state
determines $c$ everywhere through $c(gx)=g c(x)$. The first
choice gives the identity and the second gives $\tau$.
Both commute with every signed permutation, proving SM16.

A deck transformation of the actual connected cover restricts
on the base fibre to a permutation commuting with monodromy:
apply the deck transformation to a lifted loop and use uniqueness
of path lifting in the local charts of SM1. Its restriction is
injective as a function of the deck transformation. Indeed a
transformation fixing one fibre point fixes the lift of every
path starting there; connectedness, or the explicit path-lifting
description, reaches the entire cover. Thus SM16 is an upper
bound of two for the actual deck group.

Both transformations are realized. The nontrivial one is the
original rational factor-sign involution, in unchanged source
coordinates,
\begin{equation}
 \Sigma(a,y,z,w)=\left(-a,\ y+\frac{2i}{a},\
 \frac{6i}{a^2}-z,\
 w-\frac{14iy^2}{a}+\frac{28y}{a^2}+\frac{40i}{a^3}\right).
 \tag{SM17}
\end{equation}
To check it directly in the inverse model, keep the same $r,u$
in SM2 and replace $a$ by $-a$. Its new $y$ is $y+2i/a$ and
its new $z$ is $6i/a^2-z$. Expanding the two expressions for
$w$ gives exactly SM17. It follows from this same operation
that $\Sigma^2=1$ and $P\Sigma=P$ on $a\ne0$. Since $u_0=ac$
is nonzero on $\mathcal B$, every point of its preimage has
$a\ne0$. SM17 is regular on the entire cover and realizes
$\tau$. The exact cover deck group is consequently
\begin{equation}
 \Aut_{\mathcal B}(P^{-1}(\mathcal B))=\{1,\Sigma\}\cong C_2.
 \tag{SM18}
\end{equation}

The same two elements are all rational birational deck
transformations of the original polynomial map. To justify the
passage from rational maps, write $K_0=\C(u_0,u_2,u_3,u_4)$
and $L_0=\C(a,y,z,w)$ with the embedding defined by $P$.
SM1 makes $[L_0:K_0]=8$. The actual ring $S$ is also the
localization
\[
 \C[a,y,z,w,(P_0\Disc(H_{P(q)}))^{-1}],
\]
so it is integrally closed: a localization of a polynomial
unique-factorization ring has this property, since a reduced
fraction integral over it has denominator dividing a power of
its relatively prime numerator and therefore has unit denominator.
Since $S$ is finite integral over $R$, it is the full integral
closure of $R$ in $L_0$. In fact any element of $L_0$ integral
over $R$ is integral over $S$ by the same monic polynomial,
and integrally closedness puts it in $S$.
Every $K_0$-automorphism of $L_0$ preserves this integral closure,
and hence extends to a regular deck transformation on the actual
cover. SM18 therefore bounds it by those same two elements;
SM17 realizes both. Thus
\begin{equation}
 \Aut_{K_0}(L_0)=\{1,\Sigma\}.\tag{SM19}
\end{equation}
This statement is about birational automorphisms fixing the
target field, not arbitrary rational maps of the source.

Finally the only global polynomial deck automorphism $D$ of $\C^4$
satisfying $P\circ D=P$ is the identity. Such an automorphism would
restrict to SM19. The nonidentity element cannot extend
polynomially across $a=0$, because its second coordinate is
$(ay+2i)/a$, whose numerator is not divisible by $a$.
Equality on the dense open set $a\ne0$ forces equality as
rational functions, so there is no alternative polynomial
extension. This proves the assertion for the whole original
source, while retaining the nontrivial regular deck
transformation on the precise eight-sheet domain. This does not
classify polynomial automorphisms satisfying the different equation
$P\circ D=D\circ P$; that centralizer is not asserted to be trivial.

\section{Receiving the group in the fixed original conductor target}
For the fixed original receiving isomorphism
$\Psi:\C^4\to W=\cP_3$, with all original moments and Gamma
norms retained, set $P_W=\Psi P\Psi^{-1}$. Its eight-sheet
base is exactly $\Psi(\mathcal B)$. Sending a target path
$u(t)$ to $\Psi u(t)$ and each source state to its image under
$\Psi$ gives a bijection of all path lifts, with inverse
$\Psi^{-1}$. Thus its monodromy is exactly SM12 on the eight
transported marked states, and its nontrivial rational deck
transformation is exactly $\Psi\Sigma\Psi^{-1}$. The cover
deck group is $C_2$ and its global polynomial deck group is
trivial by the same conjugacy.

These are statements about the constructed nonlinear cover on
the original target. They do not enlarge the symmetry group
of a fixed Segre pairing or assert that the original conductor
itself is a nonlinear map. The original linear conductor
remains the exact invertible intertwiner proved in
FC26--35 above. All four Time labels and both signs over each
are present in the monodromy action just calculated.
% END SIGNED_MONODROMY

% BEGIN SIGNED_FRAME_CONNECTION
The four signed inverse states are those of the original ES--Fable
map \cite{ES488}, with their exact inverse coordinates proved in
GF7 and GF19--28 above. This note differentiates their actual odd frame in
the unchanged coordinates $(a,y,z,w)$. It then transports the
result through one fixed original conductor receiving map. No
moving conductor, diagonal replacement of the frame, or change
of its original Gamma norm is used.

\section{The original odd frame and its derivative}
Fix $\gamma>2$ and let $0<\delta<1/2$. Order the labels as
$(\mathrm M,\mathrm N,\mathrm T,\mathrm E)$, with corresponding
$(\varepsilon,\eta)=(1,1),(1,-1),(-1,1),(-1,-1)$. Retain
\begin{equation}
 r_j=\tfrac12+\varepsilon_j\delta+i\eta_j\gamma,
 \qquad A=-\tfrac12,\qquad
 d_j=h_\rho'(r_j)=4\delta\gamma(\varepsilon_j\gamma-i\eta_j\delta),
 \qquad a_j^2=d_j^{-1}.\tag{SF1}
\end{equation}
Each of the two sign choices is retained; take either continuous
choice on the displayed positive interval when differentiating.
The exact half-difference of its two original source points is
\begin{equation}
 o_j=a_j\,v_j,\qquad
 v_j=\begin{pmatrix}
 1\\-i d_j\\-d_j^2/2+2r_jd_j\\7ir_j^2d_j-13id_j^2
 \end{pmatrix},\qquad
 N=(v_1\ v_2\ v_3\ v_4),\quad
 O=N\diag(a_1,a_2,a_3,a_4).\tag{SF2}
\end{equation}
For example these expressions follow from the complete inverse
by using $a^{-1}=ad$, $a^{-2}=d$, $a^{-3}=ad^2$,
$a^{-4}=d^2$ and subtracting the values at $a$ and $-a$.
All constants, including $A=-1/2$, remain in SF2.

The determinant and branch logarithmic derivatives are
\begin{align}
 \det N&=-28672\delta^4\gamma^4(\delta^2+\gamma^2)
                         (\gamma^2+(2\gamma^2-1)\delta^2),\nonumber\\
 \frac{a_j'}{a_j}&=-\frac1{2\delta}
       +\frac{i\eta_j}{2(\varepsilon_j\gamma-i\eta_j\delta)}.
 \tag{SF3}
\end{align}
The determinant is obtained from the four explicit columns SF2
by expansion; the independent coefficient identity and exact
polynomial numerator certificate are supplied below. Every factor
displayed in this determinant is nonzero on the original interval.
The second equation follows by differentiating $a_j^2d_j=1$;
it is independent of the retained choice of sign.
Writing $L=\diag(a_j'/a_j)$, ordinary product differentiation gives
the exact original-coordinate connection
\begin{equation}
 \Omega_\delta=O'O^{-1}=(N'+NL)N^{-1}.\tag{SF4}
\end{equation}
In particular the square roots and their signs cancel from the
connection. This is an identity of the original matrices, not an
identification of the signed source points with each other.

\section{Every entry of the rational matrix}
For readable formulas only, abbreviate
$t=\delta^2$, $g=\gamma^2$, $Q=t+g$ and $F=g+(2g-1)t$.
These abbreviations do not change any coordinate of the four-vector.
The complete matrix is
\begin{equation}
 \boxed{\quad\Omega_\delta=
             \frac{B(\delta^2,\gamma^2)}{448\delta QF}.\quad}
 \tag{SF5}
\end{equation}
Here all sixteen entries of the polynomial matrix $B(t,g)$ are
as follows, in the original row and column order $(a,y,z,w)$:
\begin{align*}
 B_{11}&=-224(5t^2g-2t^2+tg^2+3tg+g^2),&
 B_{12}&=-28i,&B_{13}&=28,&B_{14}&=0,\\
 B_{21}&=2880it gQ^2,&&&&\\
 B_{22}&=4(280t^2g-140t^2+56tg^2-82tg+7t+28g^2-7g),&&&&\\
 B_{23}&=-8i(104tg-7t+7g),& B_{24}&=16F,&&&\\
 B_{31}&=64t gQ^2(84tg-56Q+39),&&&&\\
 B_{32}&=-4i(224t^2g-84t^2+448tg^2-86tg-7t+28g^2+7g),&&&&\\
 B_{33}&=8(252t^2g-112t^2+252tg^2+20tg-7t+28g^2+7g),&&&&\\
 B_{34}&=16iF,&&&&\\
 B_{41}&=112it gQ^2(1248tg-1012t+68g+123),&&&&\\
 B_{42}&=-7\bigl(1120t^3g-560t^3+2240t^2g^2-2336t^2g+112t^2&&&&\\
 &\hspace{15mm}+1120tg^3-1440tg^2+254tg+7t+560g^3+112g^2-7g\bigr),&&&&\\
 B_{43}&=14i(832t^2g-28t^2+832tg^2-64tg-7t-140g^2+7g),&&&&\\
 B_{44}&=28F(28Q-1).&&&&
\end{align*}
These expressions include the complete polynomial numerators;
no coefficient has been fitted at a sample quartet.

To verify SF5 from the original definition, insert the four
columns SF2 and the four entries of $L$ from SF3. Direct
multiplication yields the exact matrix identity
\begin{equation}
 B(\delta^2,\gamma^2)N
 =448\delta(\delta^2+\gamma^2)
       (\gamma^2+(2\gamma^2-1)\delta^2)(N'+NL).\tag{SF6}
\end{equation}
Every entry is a rational expression with a nonzero denominator
on the original domain. Multiplication by those denominators
reduces SF6 to sixteen polynomial identities over $\mathbb Q(i)$.
The supplied certificate constructs these polynomials from the
displayed expressions and verifies that every coefficient is zero.
The nonzero determinant SF3 then implies SF5. This verification
also proves the determinant in SF3 directly from SF2, so the
calculation does not require an inferred frame rank.

The exact trace provides a separate scalar check:
\begin{equation}
 \tr\Omega_\delta=\frac2\delta+
        \frac{2\delta(2\gamma^2-1)}{\gamma^2+(2\gamma^2-1)\delta^2}.
 \tag{SF7}
\end{equation}
It follows either by adding the diagonal entries SF5 or by
differentiating the actual determinant. In fact, with
$\Delta_\rho=-64\delta^2\gamma^2(\delta^2+\gamma^2)$,
the preserved sign invariant
$\chi=A^2\Delta_\rho\prod_ja_j\in\{1,-1\}$ gives
\begin{equation}
 \det O=1792\chi\delta^2\gamma^2
                    (\gamma^2+(2\gamma^2-1)\delta^2).
 \tag{SF8}
\end{equation}
Indeed $A^2=1/4$, so multiplication of SF3 by
$\prod a_j=4\chi/\Delta_\rho$ gives SF8 with its full sign.
Its logarithmic derivative is precisely SF7.

\section{The exact pole and retained Laurent remainder}
The pole at $\delta=0$ is simple. Its residue, still in the
unchanged original coordinates and with $g=\gamma^2$, is
\begin{equation}
 R_\gamma=\lim_{\delta\to0}\delta\Omega_\delta
 =\begin{pmatrix}
 -\tfrac12&-\dfrac{i}{16g^2}&\dfrac1{16g^2}&0\\[5pt]
 0&\dfrac{4g-1}{16g}&-\dfrac{i}{8g}&\dfrac1{28g}\\[5pt]
 0&-\dfrac{i(4g+1)}{16g}&\dfrac{4g+1}{8g}&\dfrac{i}{28g}\\[5pt]
 0&-\dfrac{7(80g^2+16g-1)}{64g}&\dfrac{7i(1-20g)}{32g}
                                   &\dfrac{28g-1}{16g}
 \end{pmatrix}.
 \tag{SF9}
\end{equation}
It is $B(0,g)/(448g^2)$ by SF5. Define the polynomial matrix
\[
 C(t,g)=\frac{g^2B(t,g)-(t+g)(g+(2g-1)t)B(0,g)}{t}.
\]
The numerator vanishes at $t=0$, entry by entry, so this is
exact polynomial division. The complete remainder identity is
\begin{equation}
 \Omega_\delta=\frac{R_\gamma}{\delta}
       +\delta\,\frac{C(\delta^2,\gamma^2)}
        {448\gamma^4(\delta^2+\gamma^2)
                        (\gamma^2+(2\gamma^2-1)\delta^2)}.
 \tag{SF10}
\end{equation}
The second term is analytic and of order $\delta$ near zero.
No higher terms have been suppressed: its exact rational
numerator and denominator remain specified by SF5 and SF10.
The entry $R_{11}=-1/2$ proves that this pole does not cancel.

For completeness the entire finite complex-$\delta$ pole set
of SF5 is
\begin{equation}
 0,\quad\pm i\gamma,\quad
 \pm\frac{i\gamma}{\sqrt{2\gamma^2-1}}.\tag{SF11}
\end{equation}
All five poles are simple when $\gamma>2$. The denominator in
SF5 has these distinct simple zeros, and its $(1,2)$ numerator
is the nonzero constant $-28i$, so none can be removed from the
matrix. At the pair with $Q=0$ the common residue is
$B(-g,g)/(1792g^2(g-1))$; at the pair with $F=0$ it is
$-(2g-1)B(-g/(2g-1),g)/(1792g^2(g-1))$.
These formulas follow by differentiating the scalar denominator
$448\delta QF$ at its simple zeros. The original positive
interval contains none of these poles except its excluded
endpoint zero.

As a rational matrix function, SF5 has pole order three at
infinity, with exact leading term
\[
 \lim_{\delta\to\infty}\delta^{-3}\Omega_\delta
 =\frac{g}{2g-1}\bigl(4(3g-2)E_{31}
                         +i(312g-253)E_{41}\bigr).
\]
Here $E_{jk}$ is the matrix with its only nonzero entry, equal
to one, in position $(j,k)$. The limit follows from the highest
powers in SF5; its $(3,1)$ entry is nonzero for $g>4$.
If the connection is regarded as the one-form
$\Omega_\delta\,d\delta$, its order at infinity is five,
because $d\delta=-\zeta^{-2}d\zeta$ for $\zeta=1/\delta$.
This statement of complex pole orders does not extend the
original real arithmetic parameter domain.

\section{Residue identities without replacing the original matrix}
Direct expansion of SF9 gives
\begin{equation}
 \det(\lambda I-R_\gamma)
   =(\lambda+\tfrac12)(\lambda-\tfrac12)^2(\lambda-\tfrac32),
 \qquad
 (2R_\gamma+I)(2R_\gamma-I)(2R_\gamma-3I)=0.
 \tag{SF12}
\end{equation}
The characteristic determinant and each entry of the latter
identity are verified as polynomial numerator identities in
$g$ by the exact certificate. The annihilating polynomial has
three distinct roots; its Lagrange interpolation polynomials
in $R_\gamma$ give complementary spectral idempotents.
This proves semisimplicity as a property of the displayed
matrix, rather than assuming a diagonal frame.

Two eigenvectors needed for the native metric statement are
particularly explicit:
\begin{equation}
 e_a=(1,0,0,0)^{\mathsf T},\qquad
 x_+=(0,1,i,35\gamma^2-7/4)^{\mathsf T},\qquad
 R_\gamma e_a=-\tfrac12e_a,\quad
 R_\gamma x_+=\tfrac32x_+.
 \tag{SF13}
\end{equation}
The first follows from the first column of SF9; substitution
in its other three rows proves the second. No coordinates of
the operator have been changed to use these identities.

\section{Exact fixed-conductor transport and Hermitian defect}
Fix one admitted reference conductor and its exact receiving
isomorphism $\Psi_*:\C^4\to W=\cP_3$, as constructed in
FC26--35 and GF19--28 above. Hold it fixed as $\delta$ varies here.
In the original target monomial basis let $H_{\Gamma,*}$ be
the full original Gram matrix of the norm
\[
 \|p\|_{\Gamma,*}^2=\int_{\mathbb R}|p(c_*+iy)|^2
 \frac{(2\pi)^{s/2}2^{s/2}}{4\pi\Gamma(s/2)}
 |\Gamma(s/4+iy/2)|^2\,dy,
 \qquad c_*=a_{\rm amp,*}/2-4.
\]
The source order $s\in\{1,a_{\rm amp,*}\}$ and every mass
in this expression remain fixed. Put
$G=\Psi_*^*H_{\Gamma,*}\Psi_*>0$. This is an exact pullback
of the original norm, not a replacement by a coefficient norm.

The transported frame and connection are
\begin{equation}
 O_W=\Psi_*O,\qquad
 \Omega_W=O_W'O_W^{-1}=\Psi_*\Omega_\delta\Psi_*^{-1},
 \qquad R_W=\Psi_*R_\gamma\Psi_*^{-1}.
 \tag{SF14}
\end{equation}
The identities follow by differentiation because the reference
map is fixed. In particular every coefficient of the original
receiving map remains in these formulas; no derivative of a
moving reference period has been silently omitted.

The entire original-coordinate Hermitian defect is
\begin{equation}
 \mathscr H_\delta=G\Omega_\delta+\Omega_\delta^*G
    =\frac{G B(\delta^2,g)+B(\delta^2,g)^*G}
                 {448\delta(\delta^2+g)(g+(2g-1)\delta^2)}
 \quad(\delta>0).
 \tag{SF15}
\end{equation}
The denominator is real and positive on the original interval,
so taking the conjugate transpose gives exactly this expression.
For the actual moving-frame Gram matrix $K_\delta=O^*GO$,
ordinary differentiation yields
\begin{equation}
 K_\delta'=O^*\mathscr H_\delta O.
 \tag{SF16}
\end{equation}
Its native determinant is also exact:
\[
 \det K_\delta=1792^2\delta^4\gamma^4
       (\gamma^2+(2\gamma^2-1)\delta^2)^2\det G,
 \qquad
 \partial_\delta\log\det K_\delta
  =\frac4\delta+
    \frac{4\delta(2\gamma^2-1)}{\gamma^2+(2\gamma^2-1)\delta^2}.
\]
The first equality is $\det K=|\det O|^2\det G$ with SF8;
the second is its derivative. Thus the native metric and its
complete mass are retained even where a scalar trace is used.

SF10 proves the operator-norm limits
\begin{equation}
 \delta\Omega_\delta\longrightarrow R_\gamma,
 \qquad
 \delta\mathscr H_\delta\longrightarrow
       G R_\gamma+R_\gamma^*G,
 \qquad
 \delta\|\Omega_W\|_{\Gamma,*}
       \longrightarrow\|R_\gamma\|_G\ge\tfrac32.
 \tag{SF17}
\end{equation}
Continuity of matrix norms in these fixed finite-dimensional
positive inner products gives the limits; the lower bound
follows by applying SF13 to $x_+$. Both signs of the leading
Hermitian defect are present for \emph{every} such fixed
original metric, as the exact identities
\begin{equation}
 e_a^*(G R_\gamma+R_\gamma^*G)e_a=-\|e_a\|_G^2<0,
 \qquad
 x_+^*(G R_\gamma+R_\gamma^*G)x_+=3\|x_+\|_G^2>0
 \tag{SF18}
\end{equation}
show. They follow from the two eigenvector equalities in SF13
and conjugate symmetry; they do not assert that $R_\gamma$
is self-adjoint or normal in this original metric. Equivalently
the physical test vectors are $\Psi_*e_a$ and $\Psi_*x_+$,
and the physical Hermitian defect pulls back exactly to SF15.
The full remainder in all these limits remains the explicit
rational remainder SF10, not an unspecified discarded term.

\section{Exact verification and source use}
The accompanying script
\texttt{derive\_signed\_frame\_connection.py} differentiates
the four labelled columns SF2, forms SF4, and saves their full
rational expressions. The independent display certificate
\texttt{verify\_signed\_frame\_connection.py} reconstructs
the sixteen entries in SF5, checks all sixteen entries of
SF6 by their polynomial numerators over $\mathbb Q(i)$,
checks the determinant SF3, both identities SF12, both vectors
SF13, the trace SF7, and the exact divisibility used in SF10.
No decimal approximation or sample parameter is used.
The captured output is
\texttt{verify\_signed\_frame\_connection\_output.txt}.
% END SIGNED_FRAME_CONNECTION

% BEGIN COEFFICIENT_CONNECTION
\section{The signed-frame connection in all four target coordinates}
The full inverse GF7, proved from the original ES factor construction
\cite{ES488,ES}, determines a differential equation without changing its
coordinates. Let $u=(A,B,C,D)$ and retain
$h(x)=Ax^4+x^3+Bx^2+Cx+D$. Work on the explicitly defined open set
\[
 \mathcal U=\{u:A\,\operatorname{Disc}(h)\det V(u)\ne0\},
 \qquad V\text{ is the full matrix in SE4}.
 \tag{SC1}
\]
Every admitted literal quartet belongs to this set by GF21 and SE5.
For a local ordering of its four roots, put $W=(r_j^n)_{n=0}^3$.
With all choices $a_j^2h'(r_j)=1$ retained, its odd frame is
$O=VW\operatorname{diag}(a_j)$ as already proved in SE7.

In the original ordered polynomial coefficient basis $(1,x,x^2,x^3)$,
define the following four-by-four matrices:
\[
 \mathsf C=\begin{pmatrix}
 0&0&0&-D/A\\1&0&0&-C/A\\0&1&0&-B/A\\0&0&1&-1/A
 \end{pmatrix},\quad
 \mathsf F=4A\mathsf C^3+3\mathsf C^2+2B\mathsf C+CI_4,
 \quad
 \mathsf J=\begin{pmatrix}0&0&0&0\\1&0&0&0\\0&2&0&0\\0&0&3&0\end{pmatrix}.
 \tag{SC2}
\]
Here the scalar $C$ and the multiplication matrix $\mathsf C$ are
different objects. Multiplication by $x$ modulo the original polynomial
$h$ has precisely the displayed matrix: its last column follows from
$Ax^4=-x^3-Bx^2-Cx-D$. Evaluation proves
$W^T\mathsf C=\operatorname{diag}(r_j)W^T$. Consequently
\[
 \det\mathsf F=\prod_jh'(r_j)
   =A^4\prod_{j<k}(r_k-r_j)^2
   =\operatorname{Disc}(h)/A^2\ne0.
 \tag{SC3}
\]
The middle equality multiplies the twelve ordered differences in
$h'(r_j)=A\prod_{k\ne j}(r_j-r_k)$; its six sign reversals have
positive product. The discriminant here is the retained degree-four
discriminant $A^6\prod_{j<k}(r_k-r_j)^2$.

For $\nu=A,B,C,D$ respectively let
\[
 f_\nu(x)=(x^4,x^2,x,1)_\nu,\qquad
 \mathsf R_\nu=-f_\nu(\mathsf C)\mathsf F^{-1},\qquad
 \mathsf T_\nu=-\tfrac12\mathsf F^{-1}
 \bigl[f'_\nu(\mathsf C)
 +(12A\mathsf C^2+6\mathsf C+2BI_4)\mathsf R_\nu\bigr].
 \tag{SC4}
\]
All entries are specified rational functions of the original four
coordinates. No roots or square-root branches need be selected to
evaluate these matrices. The exact connection matrix is
\[
 \boxed{\quad
 \Omega_\nu=(\partial_\nu V)V^{-1}
       +V\bigl(\mathsf J\mathsf R_\nu^T+
                           \mathsf T_\nu^T\bigr)V^{-1}.
 \quad} \tag{SC5}
\]
This formula retains the complete SE4 matrix, its derivatives, the
original leading coefficient $A$, and every indicated inverse and
transpose. It is a finite rational matrix formula, rather than an
unspecified existence assertion.

Here is a direct proof of SC5. Differentiating $h(r_j)=0$ and
$a_j^2h'(r_j)=1$ gives, respectively,
\[
 \partial_\nu r_j=-\frac{f_\nu(r_j)}{h'(r_j)},\qquad
 \frac{\partial_\nu a_j}{a_j}
 =-\frac{f'_\nu(r_j)+h''(r_j)\partial_\nu r_j}{2h'(r_j)}.
 \tag{SC6}
\]
Evaluation of any rational function of $\mathsf C$ is legitimate here
when its denominator is invertible. SC3 verifies this for SC4.
The evaluation identity in SC2 implies that the two quantities in SC6
are the evaluations of $\mathsf R_\nu$ and $\mathsf T_\nu$.
More precisely,
\[
 W\operatorname{diag}(\partial_\nu r_j)W^{-1}
       =\mathsf R_\nu^T,\quad
 W\operatorname{diag}(\partial_\nu a_j/a_j)W^{-1}
       =\mathsf T_\nu^T,\quad
 (\partial_\nu W)W^{-1}=\mathsf J\mathsf R_\nu^T.
 \tag{SC7}
\]
The last identity follows by differentiating each column
$(1,r_j,r_j^2,r_j^3)^T$; its derivative with respect to $r_j$ is
$\mathsf J$ times that column. Applying the product rule to
$O=VW\operatorname{diag}(a_j)$ proves
\[
 \partial_\nu O=\Omega_\nu O. \tag{SC8}
\]
Changing the root marking or the four square-root choices multiplies
$O$ on the right by a constant signed permutation matrix on a connected
overlap. Hence SC8, and therefore SC5, is independent of all those
choices while retaining their action on the actual eight states.

Write $\Omega=\sum_\nu\Omega_\nu\,du_\nu$.
Since $\Omega=dO\,O^{-1}$ on every root chart, differentiation gives
\[
 d\Omega=\Omega\wedge\Omega,\qquad
 \partial_\mu\Omega_\nu-\partial_\nu\Omega_\mu
          =[\Omega_\mu,\Omega_\nu]. \tag{SC9}
\]
Indeed $d(O^{-1})=-O^{-1}(dO)O^{-1}$, so the product rule for the
one-form $dO$ gives $d(dO\,O^{-1})=dO\wedge O^{-1}dO\,O^{-1}$.
Thus $d-\Omega$ is flat, with the actual $O$ as a local fundamental
matrix. For a loop in $\mathcal U$ whose signed continuation is $R$,
its based transport is exactly $O(u_b)R O(u_b)^{-1}$: each continued
column is the specified sign of the specified endpoint column. This
statement identifies the transport of every such loop; it does not
replace the separately proved SM path calculations.

On the literal quartet with fixed $\gamma$, substitution of GF20 gives
\[
 \Omega_\delta=2\delta\Omega_B-2\delta\Omega_C+
  \delta(\tfrac12-2\delta^2-2\gamma^2)\Omega_D,
 \qquad O'=\Omega_\delta O. \tag{SC10}
\]
The coefficient of $\Omega_A$ is zero because $A=-1/2$ is retained.
SC10 is also independently checked against all sixteen entries of the
explicit quartet matrix SF5 by
\texttt{verify\_coefficient\_connection.py}. The certificate multiplies
out exact numerator identities; it uses no numeric root samples.

Fix a receiving original period and its maps $\Phi_*,\Psi_*$ and
$L\Phi_*=\Psi_*$ from FC31--35. The star denotes this fixed reference,
not a moving family of arithmetic periods. Define
\[
 \Omega^U_\nu=\Phi_*\Omega_\nu\Phi_*^{-1},\qquad
 \Omega^W_\nu=\Psi_*\Omega_\nu\Psi_*^{-1}.
 \quad L\Omega^U_\nu=\Omega^W_\nu L. \tag{SC11}
\]
The displayed intertwining follows by substituting $L\Phi_*=\Psi_*$
on each side. Hence both differential equations and their full path
transports pass through the actual fixed conductor. Let
$G=\Psi_*^*\langle\ ,\ \rangle_{\Gamma,W}\Psi_*$ and $G_U$ be its
source counterpart, retaining FC56a and all original Gamma masses
\cite{WCF}. For a real path and $z'=\Omega_tz$, differentiation of
the exact original norm gives
\[
 \frac{d}{dt}\|\Psi_*z\|_{\Gamma,W}^2
       =z^*(\Omega_t^*G+G\Omega_t)z. \tag{SC12}
\]
This is the full metric defect. Flatness SC9 does not remove it.
% END COEFFICIENT_CONNECTION

% BEGIN TOTAL_CONNECTION
\section{The connection of the complete eight-state object}
Fix the original receiving $X_*,Y_*,\Phi_*,\Psi_*$ of FC31--35 at
one admitted reference period. The auxiliary quartic is now varied
through GF20, while this original conductor and its norms stay fixed.
Keep the ordered roots
$r_j=\tfrac12+\epsilon_j\delta+i\eta_j\gamma$ in the original
$(\mathrm M,\mathrm N,\mathrm T,\mathrm E)$ marking, with
$(\epsilon_j,\eta_j)=(1,1),(1,-1),(-1,1),(-1,-1)$.
All eight states are still $q_{j,\pm}=e_j\pm o_j$ of SE1.
Define the complete coordinate map and its inverse by
\[
 \mathsf F_\delta=
 \begin{pmatrix}I_4&0\\ E_\delta&O_\delta\end{pmatrix},\qquad
 \mathsf F_\delta^{-1}=
 \begin{pmatrix}I_4&0\\-O_\delta^{-1}E_\delta&O_\delta^{-1}\end{pmatrix}.
 \tag{TC1}
\]
Here $\mathsf F_\delta(\alpha,\beta)=(\alpha,z)$ with
$z=E_\delta\alpha+O_\delta\beta$. SE7 proves invertibility on the
whole admitted interval. The physical maps are
$\mathsf J_U=\operatorname{diag}(X_*,\Phi_*)\mathsf F_\delta$ and
$\mathsf J_W=\operatorname{diag}(Y_*,\Psi_*)\mathsf F_\delta$.
They satisfy $(L\oplus L)\mathsf J_U=\mathsf J_W$ by FC32.

Differentiate TC1, retaining both its blocks. The exact differential
equation for a constant marked state $(\alpha,\beta)$ is
\[
 \frac{d}{d\delta}\begin{pmatrix}\alpha\\z\end{pmatrix}
 =\mathcal B_\delta\begin{pmatrix}\alpha\\z\end{pmatrix},\qquad
 \mathcal B_\delta=\mathsf F_\delta'\mathsf F_\delta^{-1}
 =\begin{pmatrix}0&0\\
 E_\delta'-\Omega_\delta E_\delta&\Omega_\delta\end{pmatrix}.
 \tag{TC2}
\]
For clarity every entry of the lower-left block is specified by the
full SF matrix and the following original columns and derivatives:
\[
 \begin{aligned}
 d_j&=4\delta\gamma(\epsilon_j\gamma-i\eta_j\delta),&
 d_j'&=4\gamma(\epsilon_j\gamma-2i\eta_j\delta),\\
 e_j&=(0,-r_j,3id_j,
       (B-17r_j+Ar_j^2)d_j-2Ad_j^2)^T,& A&=-\tfrac12,\\
 e_j'&=(0,-\epsilon_j,3id_j',
 [2\delta-17\epsilon_j+2Ar_j\epsilon_j]d_j
 +(B-17r_j+Ar_j^2)d_j'-4Ad_jd_j')^T.
 \end{aligned} \tag{TC3}
\]
The $j$th lower-left column is exactly $e_j'-\Omega_\delta e_j$.
No term from the moving means has been omitted.

For any two points joined by a continued branch in the admitted
domain, the full transport, with its inverse obtained by swapping
the endpoints, is
\[
 \mathcal T_{2,1}=\mathsf F_{\delta_2}\mathsf F_{\delta_1}^{-1}
 =\begin{pmatrix}I_4&0\\
 E_{\delta_2}-O_{\delta_2}O_{\delta_1}^{-1}E_{\delta_1}&
 O_{\delta_2}O_{\delta_1}^{-1}\end{pmatrix}.
 \tag{TC4}
\]
Multiplication proves $\mathcal T_{3,2}\mathcal T_{2,1}
=\mathcal T_{3,1}$; differentiation at the first endpoint gives TC2.
Thus the lower-left term is part of the proved transport, including
when the initial evaluation $z$ is zero.

Put $r_+=\tfrac12+i\gamma$, $r_-=\tfrac12-i\gamma$, and retain
the exact residue $R_\gamma=\lim_{\delta\to0}\delta\Omega_\delta$
calculated in SF. Substitution in TC3 gives
\[
 E_0=e_y\ell_\gamma,\quad e_y=(0,1,0,0)^T,\quad
 \ell_\gamma=(-r_+,-r_-,-r_+,-r_-).
 \tag{TC5}
\]
Indeed every $d_j$ vanishes at $\delta=0$, while the two repeated
root values stay as displayed. Since $E_\delta$ is polynomial in
$\delta$, TC2 and the explicit simple pole of SF give
\[
 \mathcal B_\delta=\frac{\mathcal R_\gamma}{\delta}+O(1),\qquad
 \mathcal R_\gamma=
 \begin{pmatrix}0&0\\-R_\gamma E_0&R_\gamma\end{pmatrix}.
 \tag{TC6}
\]
The $O(1)$ refers to every entry in these unchanged coordinates.
The lower-left residue has rank exactly one. In fact SF proves
$\det R_\gamma=-3/16\ne0$, and $E_0$ has rank one because neither
$e_y$ nor $\ell_\gamma$ is zero. This rank statement therefore
retains the actual means and all four labels.

Direct block determinants and multiplication give
\[
 \det(\lambda I_8-\mathcal R_\gamma)
 =\lambda^4(\lambda+\tfrac12)
                  (\lambda-\tfrac12)^2(\lambda-\tfrac32),
 \qquad
 \mathcal R_\gamma(2\mathcal R_\gamma+I_8)
 (2\mathcal R_\gamma-I_8)(2\mathcal R_\gamma-3I_8)=0.
 \tag{TC7}
\]
For the product identity, powers of the lower block have the form
$\mathcal R_\gamma^n=[\,0\ 0;\,-R_\gamma^nE_0\ R_\gamma^n\,]$
for $n\ge1$. Substitution of the SF cubic annihilator proves TC7.
Its four roots are distinct, so the residue is semisimple: for
completeness the interpolation polynomials
$p_\lambda(t)=\prod_{\mu\ne\lambda}(t-\mu)/(\lambda-\mu)$,
with $\lambda,\mu\in\{0,-1/2,1/2,3/2\}$, give mutually
annihilating idempotents $p_\lambda(\mathcal R_\gamma)$ summing
to $I_8$, with image in the indicated eigenspace. This is a proof
about the displayed original matrix, without replacing its coordinates.

\section{A complete boundary loop and its exact original metric defect}
Choose any admitted real $\delta_b\in(0,1/2)$, and continue
$\delta(t)=\delta_b e^{2\pi it}$ for $0\le t\le1$, with the same
real $\gamma>2$. All quartic roots remain distinct: their possible
collisions occur at $\delta=0$ or $\delta=\pm i\gamma$. The other
zeros of SE5 lie at
$\delta=\pm i\gamma/\sqrt{2\gamma^2-1}$, whose moduli exceed
$1/\sqrt2$. Thus this entire loop belongs to SC1.
The factors
$\epsilon_j\gamma-i\eta_j\delta(t)$ trace circles of radius
$\delta_b$ about $\epsilon_j\gamma$; they have winding zero around
zero. Therefore each
$d_j=4\gamma\delta(\epsilon_j\gamma-i\eta_j\delta)$ has winding
one. Its inverse square root $a_j$ changes sign. The root marking
and every mean $e_j$ return unchanged. It follows, without a
resonance assumption, that the full based transport is
\[
 \mathcal T_{\rm loop}=
 \begin{pmatrix}I_4&0\\2E_{\delta_b}&-I_4\end{pmatrix},\qquad
 (\mathcal T_{\rm loop})^2=I_8.
 \tag{TC8}
\]
On the eight marked inverse states this is precisely the simultaneous
swap of all four signs, the central deck element already identified
in SM19. Its lower-left block has rank three by SE8. The rank-one
boundary residue in TC6 and the rank-three finite based holonomy in
TC8 are both calculated from the same original means, at their stated
different parameter values.

The fixed original target metric on the two retained coordinates is
$\mathbb G=\operatorname{diag}(G_Y,G)$, with
$G_Y=Y_*^*\langle\ ,\ \rangle_{\Gamma,W}Y_*$ and
$G=\Psi_*^*\langle\ ,\ \rangle_{\Gamma,W}\Psi_*$ from FC56a.
For $\zeta'=\mathcal B_\delta\zeta$ its complete defect is
\[
 \frac{d}{d\delta}(\zeta^*\mathbb G\zeta)
 =\zeta^*(\mathcal B_\delta^*\mathbb G+
                              \mathbb G\mathcal B_\delta)\zeta,
 \quad
 \delta(\mathcal B_\delta^*\mathbb G+
                 \mathbb G\mathcal B_\delta)
 \longrightarrow
 \begin{pmatrix}
 0&-E_0^*R_\gamma^*G\\-GR_\gamma E_0&R_\gamma^*G+GR_\gamma
 \end{pmatrix}. \tag{TC9}
\]
For the original vectors $e_a=(1,0,0,0)^T$ and
$x_+=(0,1,i,35\gamma^2-7/4)^T$, SF proves
$R_\gamma e_a=-e_a/2$ and $R_\gamma x_+=3x_+/2$.
Consequently TC9 on $(0,e_a)$ equals $-\|e_a\|_G^2<0$, and
on $(0,x_+)$ equals $3\|x_+\|_G^2>0$. In particular the original
metric defect has both signs for all sufficiently small positive
$\delta$; no positive fixed original metric makes it vanish there.
Continuity of finite matrix norms also gives the exact limits
\[
 \delta\|\mathcal B_\delta\|_{\mathbb G}
       \longrightarrow\|\mathcal R_\gamma\|_{\mathbb G}\ge\tfrac32,
 \qquad
 \delta\|\Omega_\delta\|_G\longrightarrow\|R_\gamma\|_G
       \ge\tfrac32. \tag{TC10}
\]
The lower bounds follow by evaluating the operators on the displayed
$3/2$ eigenvector, with its unchanged positive norm.

Finally conjugate TC2--10 by
$\operatorname{diag}(Y_*,\Psi_*)$ or
$\operatorname{diag}(X_*,\Phi_*)$ to obtain their exact physical
target and source maps. Since the original conductor is fixed,
\[
 (L\oplus L)\mathcal B^U_\delta
      =\mathcal B^W_\delta(L\oplus L),\qquad
 (L\oplus L)\mathcal T^U_{2,1}
      =\mathcal T^W_{2,1}(L\oplus L). \tag{TC11}
\]
These equalities follow by cancelling the stated fixed isomorphisms
on the two sides. They preserve the original inverse-exterior
constants of SE17 and audit (11). The new singularity is the proved
boundary connection of the nonlinear inverse states carried through
that conductor; the fixed original linear map retains its original
nonzero determinant.
% END TOTAL_CONNECTION

% BEGIN FACTOR_COMPLETION
\section{Completing the missing factor sign with exact chart morphisms}
The web continuation \cite{WebFable}, Sections 1--2, proposes the
two-chart factor completion. Here its chart inverses are made explicit
at the boundary, so the completed object and the original polynomial
are related by exact morphisms on their stated domains.
Keep the original factor coordinates, with
$L=aU+bV$ and $Q=cU^3+dU^2V+eUV^2+fV^3$, and define
\[
 \mathfrak X=\{(a,b,c,d,e,f)\in\C^6:
 ad+bc=1,\ -cb^3+dab^2-ea^2b+fa^3=1\},
 \quad\pi=(ac,ae+bd,af+be,bf). \tag{FX1}
\]
For $\sigma\in\{i,-i\}$ define $\varphi_\sigma:\C^4\to\mathfrak X$
by the complete formulas
\[
 \begin{aligned}
 b&=\sigma+ay,&c&=-\sigma+2ay+a^2z,\\
 d&=-\sigma y-2ay^2-\sigma az-a^2yz,&
 e&=-7\sigma y^2+2z+aw,\\
 f&=2a^2y^3z+awy+4ay^4+6\sigma ay^2z+\sigma w
                                      +3\sigma y^3-4yz.
 \end{aligned} \tag{FX2}
\]
The first coordinate is the unchanged $a$.
Substitution using $\sigma^2=-1$ gives both constraints in FX1.
The original map is exactly $P=\pi\varphi_i$; the other chart
map is $P_{-i}=\pi\varphi_{-i}$.

The variety $\mathfrak X$ is smooth of complex dimension four.
Indeed scaling $L$ alone and $Q$ alone gives two ambient tangent
vectors on which the derivatives of the two constraint functions
are the columns of
$\left(\begin{smallmatrix}1&1\\3&1\end{smallmatrix}\right)$,
whose determinant is $-2$. Thus the two constraint gradients are
independent at every point. If a tangent vector is killed by
$d\pi$, then $\dot LQ+L\dot Q=0$, since the omitted cubic
coefficient has derivative zero by the first constraint. The
resultant-one constraint makes $L$ and $Q$ coprime. Hence $L$
divides $\dot L$, and homogeneity forces $\dot L=tL$;
then $\dot Q=-tQ$. The derivative of the resultant in this
direction is $2t$, which must vanish. Therefore $t=0$, and
$d\pi$ is an isomorphism between its four-dimensional tangent
spaces at every point.

At $a=0$, the constraints imply $bc=1$ and $-cb^3=1$,
so $b=\pm i$ and $c=-b$, with $d,e,f$ arbitrary. For a fixed
$\sigma$ set
\[
 U_\sigma=\mathfrak X\setminus\{a=0,b=-\sigma\}.
 \tag{FX3}
\]
On its part $a\ne0$ the inverse to FX2 is
\[
 y=(b-\sigma)/a,\quad
 z=(c+\sigma-2ay)/a^2,\quad
 w=(e+7\sigma y^2-2z)/a. \tag{FX4}
\]
Eliminating $d$ with the first constraint and $f$ with the second
shows these recover all six factor coefficients.
To prove regularity at $a=0,b=\sigma$, work in the explicit open
set $b(b+\sigma)\ne0$. The following formulas have only units
in their denominators and coincide with FX4:
\[
 \begin{aligned}
 y&=\frac{2db^2-aeb+a^2f}{b+\sigma},\\
 z&=\frac{7y^2-10\sigma ay^3-4a^2y^4-eb+af}{2b^3},\\
 w&=\frac{eb(2\sigma y+ay^2)-11\sigma y^3-17ay^4
                                      +7\sigma a^2y^5-f}{b^3}.
 \end{aligned} \tag{FX5}
\]
For a direct verification, eliminate $bc$ from the resultant to
get $b^2+1=a(2db^2-aeb+a^2f)$. Factoring its left side as
$(b-\sigma)(b+\sigma)$ proves the first formula and
$b=\sigma+ay$. Next $bc=1-ad$ gives
$c+\sigma-2ay=-a(d+\sigma y+2ay^2)/b$.
Substitution of $b=\sigma+ay$ into the preceding resultant
identity gives
\[
 2b^2(d+\sigma y+2ay^2)
 =a[-7y^2+10\sigma ay^3+4a^2y^4+eb-af].
 \tag{FX6}
\]
This proves the second formula, including its extra factor of $a$.
Finally expand $b^3=-\sigma-3ay+3\sigma a^2y^2+a^3y^3$;
substituting the second formula in $e+7\sigma y^2-2z$ proves
the third formula, with its factor of $a$. Thus no division by
zero or unproved divisibility is needed at the boundary.
At $a=0$ the formulas give
$y=\sigma d$, $z=(e+7\sigma y^2)/2$, and
$w=(f-3\sigma y^3+4yz)/\sigma$; substitution in FX2 recovers
$d,e,f$ there as well. FX4--5 cover $U_\sigma$ and agree on
their overlap. They therefore construct the full regular inverse.
In particular both $U_i,U_{-i}$ are actual affine four-space
charts, they cover $\mathfrak X$, and their overlap is exactly
$a\ne0$, isomorphic to $\C^*\times\C^3$.

The transition from the negative to the positive chart is
\[
 \varphi_i^{-1}\varphi_{-i}(a,Y,Z,W)
 =\left(a,Y-\frac{2i}{a},Z+\frac{6i}{a^2},
 W+\frac{14iY^2}{a}+\frac{28Y}{a^2}-\frac{40i}{a^3}\right).
 \tag{FX7}
\]
It follows by applying FX4 to FX2; its differential is triangular
with all diagonal entries one, so its determinant is exactly one.
The inverse transition is obtained by replacing $i$ with $-i$
in FX7. Their composition is the identity, as the same inverse
chart construction also proves. Multiplication gives
$P_i\varphi_i^{-1}\varphi_{-i}=P_{-i}$ on this overlap.

Every simple projective root of the original homogeneous quartic
now contributes both resultant-one factor signs. A multiple root
contributes none. At the finite roots this is GF7. At infinity,
$A=0$ and the fixed cubic coefficient is one, so infinity is
simple and its two signs are precisely $b=\pm i$ just constructed.
Consequently
\[
 \#\pi^{-1}(u)=2\#\{\text{simple projective roots of }H_u\},
 \qquad S_\pi=\{u:\operatorname{Disc}(H_u)=0\}.
 \tag{FX8}
\]
Here $S_\pi$ denotes finite target limits of factors escaping every
compact subset of the affine six-coordinate variety FX1.
For the nonproper inclusion, a multiple projective root must be
finite: infinity is simple whenever it is a root. Keep that root
$r$ and deform $h$ to $h+\varepsilon^2(x-r)$. Its selected
derivative is $\varepsilon^2$, so its factor coefficient
$a=\varepsilon^{-1}$ diverges, while all target coefficients
converge. This proves every discriminant point is in $S_\pi$.

Conversely fix a squarefree target. On a sufficiently small target
neighborhood its four projective roots lie in four disjoint coordinate
disks and are simple. Each root determines a factor line, locally
represented by a nonzero linear polynomial $\ell$. Polynomial
division gives $H=\ell q$ with bounded continuous coefficients,
and $R=\operatorname{Res}(\ell,q)\ne0$ throughout a smaller
neighborhood. Every factor pair over this selected root has the
unique form $(s\ell,s^{-1}q)$ with $s^2R=1$, since the
resultant of degrees one and three scales by $s^{3-1}$.
Thus both $s$ and $s^{-1}$ are bounded there, and every coefficient
of every lift is bounded. Four such disks cover all lifts. A
convergent sequence of targets in this neighborhood therefore
cannot have an escaping factor lift; this proves the reverse
inclusion and all of FX8. The fixed cubic coefficient ensures
the first constraint for every such factorization.

The original map has the additional nonproper hyperplane $A=0$
because it sees only $U_i$; FX7 exhibits its missing finite
representatives in $U_{-i}$. The selected-root escape RC1 persists
in both charts as a divergent factor coefficient. These are
explicitly different degenerations of the same completed object.
No identification of $\mathfrak X$ with the original four-space
or with the arithmetic source quotient is used.
% END FACTOR_COMPLETION

% BEGIN SELECTED_ROOT_COLLISION
\section{The second boundary: a selected root collides}
The incoming web continuation \cite{WebFable}, Section 4, provides
the following branch. Its polynomial $\widehat P_i$ is exactly the
original $P$ of FC31, under the sole coordinate dictionary $t=w$:
both use the identical five factor expressions in GF2, and multiplying
$(ac,ae+bd,af+be,bf)$ gives each of the four original components.
Thus no new polynomial is introduced. For $0<\varepsilon<1/2$ put
\[
 q_\varepsilon=(\varepsilon^{-1},-i\varepsilon,
          3i\varepsilon^2,\varepsilon^2-13i\varepsilon^3)^T.
 \quad P(q_\varepsilon)=(0,1,\varepsilon^2,0)^T.
 \tag{RC1}
\]
The displayed image follows by substitution in those four components.
Its selected root is $r=0$ of $h_\varepsilon(x)=x^3+x^2+
\varepsilon^2x$, so $h_\varepsilon'(0)=\varepsilon^2$ and
$a^2h_\varepsilon'(0)=1$ with the original $a=\varepsilon^{-1}$.
The discriminant is $\varepsilon^4(1-4\varepsilon^2)$, nonzero
on this interval. At its endpoint the selected zero root becomes
double. For fixed $\Psi_*$ and its original positive Gram $G$,
\[
 \varepsilon\|\Psi_*q_\varepsilon\|_{\Gamma,W}
       \longrightarrow\sqrt{G_{11}}>0,\qquad
 \frac{\|\Psi_*(P(q_\varepsilon)-(0,1,0,0)^T)\|_{\Gamma,W}}
      {\varepsilon^2}=\sqrt{G_{33}}.
 \tag{RC2}
\]
The first follows from $\varepsilon q_\varepsilon\to e_a$ and
continuity of this fixed norm; the second is exact for every
$\varepsilon$. In particular the square root of the displayed
output distance times the source norm tends to
$G_{33}^{1/4}G_{11}^{1/2}$. The factor coefficient $a$ itself
diverges. This escape therefore persists in the full factor variety,
not just its $b=i+ay$ chart.

The complete derivative and inverse in the original coordinate order
are
\[
 J_\varepsilon=DP(q_\varepsilon)=
 \begin{pmatrix}
 4i&2\varepsilon^{-2}&\varepsilon^{-3}&0\\
 2\varepsilon-14i\varepsilon^2&-13&2\varepsilon^{-1}&\varepsilon^{-2}\\
 -2\varepsilon^3-2i\varepsilon^2&2&0&0\\
 -i\varepsilon^4&\varepsilon^2&0&0
 \end{pmatrix},\quad\det J_\varepsilon=-2, \tag{RC3}
\]
\[
 J_\varepsilon^{-1}=
 \begin{pmatrix}
 0&0&-1/(2\varepsilon^3)&\varepsilon^{-5}\\
 0&0&-i/(2\varepsilon)&\varepsilon^{-2}+i\varepsilon^{-3}\\
 \varepsilon^3&0&3i&-2/\varepsilon-6i/\varepsilon^2\\
 -2\varepsilon^4&\varepsilon^2&1-39i\varepsilon/2&
                  17+39i/\varepsilon-2/\varepsilon^2
 \end{pmatrix}. \tag{RC4}
\]
Differentiation of the original components gives RC3; multiplying
RC3 and RC4 in either order gives $I_4$. The velocity identity is
$J_\varepsilon q_\varepsilon'=2\varepsilon e_z$, where
$e_z=(0,0,1,0)^T$. The complete exact checks are reproduced by
\texttt{derive\_selected\_root\_collision.py}, including every
compound matrix, rather than only the scalar determinant.

Use increasing multi-indices for each exterior coordinate and let
$E_{I,J}$ send the coefficient of $e_J$ to $e_I$. Direct expansion
of all minors in RC3 and RC4 gives the full leading matrices:
\[
 \begin{array}{c|c|c|c|c}
 r& m_r&\lim\varepsilon^{m_r}\Lambda^rJ_\varepsilon&
        n_r&\lim\varepsilon^{n_r}\Lambda^rJ_\varepsilon^{-1}\\\hline
 1&3&E_{1,3}&5&E_{1,4}\\
 2&5&E_{12,34}&5&-E_{12,34}/2\\
 3&5&2E_{123,234}&3&-E_{124,234}/2\\
 4&0&-2&0&-1/2
 \end{array} \tag{RC5}
\]
These are entrywise limits, with all entries outside the displayed
positions zero. For example the $(12,34)$ minor in RC3 is exactly
$\varepsilon^{-5}$, and the $(123,234)$ minor is exactly
$2\varepsilon^{-5}$. The finite Laurent entries in RC3--4 determine
every other minor by the determinant formula, and the certificate
lists all leading positions. Thus RC5 is an identity calculation,
not an inference from a sampled singular-value plot.

Let $G$ remain the complete fixed original target Gram from FC56a,
and let $G^{-1}$ denote its actual inverse. Define positive constants
\[
 \alpha=\sqrt{G_{11}(G^{-1})_{33}},\qquad
 \beta=\sqrt{\det G_{12,12}\det(G^{-1})_{34,34}},\qquad
 \xi=\sqrt{G_{11}(G^{-1})_{44}}. \tag{RC6}
\]
Every principal minor here is positive. In the induced exterior
norm the coordinate operator $cE_{I,J}$ has exact norm
$|c|\sqrt{\det G_{I,I}\det(G^{-1})_{J,J}}$: it is the vector
$ce_I$ times the coordinate functional dual to $e_J$, whose
dual squared norm is the indicated inverse Gram minor.
Continuity of that finite-dimensional operator norm in RC5, and
the complementary-minor identities
\[
 \det G_{123,123}\det(G^{-1})_{234,234}
       =G_{11}(G^{-1})_{44},\quad
 \det G_{124,124}\det(G^{-1})_{234,234}
       =G_{11}(G^{-1})_{33}, \tag{RC7}
\]
give the forward exterior constants $(\alpha,\beta,2\xi,2)$
with powers $(3,5,5,0)$, and inverse constants
$(\xi,\beta/2,\alpha/2,1/2)$ with powers $(5,5,3,0)$.
To verify RC7, use
$\det G_{I,I}=\det G\det(G^{-1})_{I^c,I^c}$ and its inverse
version, obtained by the corresponding cofactor identities in
$G^{-1}=\operatorname{adj}(G)/\det G$.

Let $\sigma_1\ge\cdots\ge\sigma_4>0$ be the singular values of
the literal operator $\Psi_*J_\varepsilon\Psi_*^{-1}$ in the
original Gamma norm. The norm of its $r$th exterior power is
$\sigma_1\cdots\sigma_r$: in orthonormal singular bases the
exterior eigenvalues are all $r$-fold products, whose greatest
one uses the $r$ greatest factors. Taking successive quotients
of the four proved exterior asymptotes yields
\[
 \boxed{
 \sigma_1\sim\alpha\varepsilon^{-3},\quad
 \sigma_2\sim(\beta/\alpha)\varepsilon^{-2},\quad
 \sigma_3\sim2\xi/\beta,\quad
 \sigma_4\sim\varepsilon^5/\xi.}
 \qquad
 \frac{\sigma_1}{\sigma_4}\sim\alpha\xi\varepsilon^{-8}.
 \tag{RC8}
\]
This extends the scales stated in the incoming web calculation to
all exact constants in the unchanged original metric. The product
is exactly $2$, and the squared volume Jacobian is exactly $4$
for every $\varepsilon$, because both sides use the same $G$.
Each squared singular value is the corresponding eigenvalue of
the relative pullback metric $G^{-1}J_\varepsilon^*GJ_\varepsilon$,
so RC8 also evaluates all four metric-degeneration constants.

The leading right and left exterior directions in ranks one,
two and three are determined without discarding the original
Gram. For RC5 position $(I,J)$ they are respectively
$G^{-1}\operatorname{span}\{e_j:j\in J\}$ and
$\operatorname{span}\{e_i:i\in I\}$, carried physically by
$\Psi_*$. Indeed the Riesz vector of the coefficient functional
$e_J^*$ is $(\Lambda^rG)^{-1}e_J$, and its decomposable factors
are the indicated $G^{-1}e_j$; the range is the stated $e_I$.
The limiting rank-one map is nonzero, so its largest singular
line is simple, and convergence of the finite Hermitian Gram
matrices forces the maximizing lines to converge to these
directions. This supplies the native singular flags for RC8.

Finally set $J_U=\Phi_*J_\varepsilon\Phi_*^{-1}$ and
$J_W=\Psi_*J_\varepsilon\Psi_*^{-1}$. Differentiating the already
proved nonlinear conductor intertwining gives
\[
 LJ_U=J_WL,\qquad
 J_W^{-1}=LJ_U^{-1}L^{-1},\quad
 \frac{\|\Lambda^rJ_U^{-1}\|}
      {\|\Lambda^rL\|\|\Lambda^rL^{-1}\|}
 \le\|\Lambda^rJ_W^{-1}\|
 \le\|\Lambda^rL\|\|\Lambda^rL^{-1}\|
                       \|\Lambda^rJ_U^{-1}\|. \tag{RC9}
\]
The lower bound follows by interchanging source and target in the
upper bound. Every norm is its actual original source or target
norm, and the same calculation RC6--8 with $G_U$ gives the exact
source constants. The original inverse-exterior factor in audit
(11) is thus retained explicitly in this second degeneration.
% END SELECTED_ROOT_COLLISION

% BEGIN CORNER_ACTION
\section{The original corner receiver and its exact multiplication residual}
The diagonal flow in the web continuation \cite{WebFable}, Section 6,
uses labelled corner coefficients. Its exact arithmetic receiver is
constructed here; the earlier cubic receiver $\Psi$ is retained through
an explicit comparison map. Fix the original admitted period with
$0<\delta_*<1/2$, $\gamma_*>2$, and let $k\equiv1\pmod4$,
$k\ge9$, $j=k-8$, $q_-=(j+1)^2$, $D=q_--1$. The original lower
grid and polynomial are
\[
 \lambda_{ab}=j/2+(2a-j)\delta_*+i(2b-j)\gamma_*,\quad
 \chi_j(S')=\prod_{a,b=0}^j(S'-\lambda_{ab}),\quad
 E_j=\C[S']/\bigl(\chi_j\bigr). \tag{CA1}
\]
Distinct pairs give distinct roots: their real and imaginary
differences force $a$ and $b$ separately equal. Order the four corners
as $(j,j),(j,0),(0,j),(0,0)$ with retained labels
$(\mathrm M,\mathrm N,\mathrm T,\mathrm E)$, and denote them by
$\lambda_l=j\rho_l$. Define the actual polynomials
\[
 p_l(S')=\frac{\chi_j(S')}
 {(S'-\lambda_l)\chi_j'(\lambda_l)},\qquad
 J_c:\C^4\longrightarrow\mathcal P_D,\quad
 J_c z=\sum_l z_l p_l(S'). \tag{CA2}
\]
Evaluation at the four marked corners is a left inverse, since
$p_l(\lambda_m)=\delta_{lm}$; evaluation at every other grid
point is zero. Thus $J_c$ is injective both in the polynomial
target and after reduction $I_D:\mathcal P_D\to E_j$.
The latter reduction is itself an isomorphism, because the
degree of $\chi_j$ is $D+1$. All its coefficients and all
cardinal denominators remain in CA2.

Retain the original finite-shift conductor, its actual vanishing
order $v$, and every moment $\mu_n$. In the source quotient
$U_D=\mathcal P_{v+D}/\mathcal P_{v-1}$ and target
$W_D=\mathcal P_D$, its monomial coefficient matrix is
\[
 (B_{D,v})_{rt}=\begin{cases}
 \binom{v+t}{r}\mu_{v+t-r},&0\le r\le t\le D,\\
 0,&r>t.
 \end{cases}\qquad
 b_r=(B_{D,v})_{rr}=\binom{v+r}{v}\mu_v\ne0. \tag{CA3}
\]
Indeed applying the finite-shift formula to $S^{v+t}$ and
using $\mu_0=\cdots=\mu_{v-1}=0$ gives exactly CA3.
For a complete finite expression for its inverse put
$D_b=\operatorname{diag}(b_0,\ldots,b_D)$ and
$N=D_b^{-1}B_{D,v}-I$; then $N^{D+1}=0$ and
\[
 B_{D,v}^{-1}=\left(\sum_{m=0}^D(-N)^m\right)D_b^{-1}.
 \tag{CA4}
\]
Multiplication of the finite geometric sum proves both inverse
identities. This formula is in the original monomial coordinates;
it is not a replacement of the physical norms. Let $T_D:U_D\to W_D$
be the original map, and define the source injection
$K_c=T_D^{-1}J_c$, using CA4 for every coefficient. Then
$T_DK_c=J_c$, and $K_c$ is injective.

On the indicated four-dimensional images the full nonlinear maps
and their exact comparison with the preceding cubic receiver are
\[
 P_c^W=J_cPJ_c^{-1},\quad P_c^U=K_cPK_c^{-1},\quad
 T_DP_c^U=P_c^WT_D,\qquad
 F_W=J_c\Psi_*^{-1}:W\to\operatorname{im}J_c,\quad
 F_U=K_c\Phi_*^{-1}:U\to\operatorname{im}K_c,
 \quad T_DF_U=F_WL. \tag{CA5}
\]
Every inverse here is on its stated image. The last equality is
obtained from $L\Phi_*=\Psi_*$; in particular this is a proved
relation between the two receivers, with every original moment
in the source comparison.
For the actual original target Gamma Gram $H_{D,s}$, the corner
metric is the complete positive matrix
$H_c=J_c^*H_{D,s}J_c$, $s\in\{1,a_{\rm amp,*}\}$.
Its positivity follows from injectivity and the original positive
norm \cite{WCF}; no column is declared orthonormal.
The source uses $K_c^*H_{U_D,s}K_c$ with its original quotient
metric. All previously proved fixed-metric singular formulas apply
with these explicit matrices, and CA5 preserves their differentiated
conductor relation.

Let $M_j$ be multiplication by $S'$ on the actual quotient $E_j$,
and put $B_c=\operatorname{diag}(\lambda_\mathrm M,
\lambda_\mathrm N,\lambda_\mathrm T,\lambda_\mathrm E)$.
The entire polynomial residual is
\[
 S'J_cz-J_cB_cz
 =\chi_j(S')\sum_l\frac{z_l}{\chi_j'(\lambda_l)},\qquad
 M_j I_DJ_c=I_DJ_c B_c. \tag{CA6}
\]
The first identity follows separately from each CA2 column and
has degree $D+1$, outside $W_D$ in general. Reduction by the
actual $\chi_j$ gives the second identity. Thus this is exactly
the arithmetic quotient action represented by $B_c$, with its
full polynomial residual displayed. It is not asserted to be the
unreduced multiplication operator on the earlier cubic space.
Exponentiating the finite-dimensional equality gives
$e^{tM_j}I_DJ_c=I_DJ_c e^{tB_c}$ for every complex $t$, term
by term in the convergent matrix series.

\section{Polynomial lift of this original action and finite complex escape}
For any constant $4\times4$ matrix $B$, the original constant
Jacobian $\det DP=-2$ gives the complete polynomial vector field
\[
 X_B(x)=-\tfrac12\operatorname{adj}(DP(x))\,BP(x),\qquad
 DP(x)X_B(x)=BP(x). \tag{CA7}
\]
The adjugate entry $(r,t)$ is explicitly
$(-1)^{r+t}\det DP(x)_{\widehat t,\widehat r}$, so CA7 is a
finite polynomial formula in the fully displayed original $P$.
The adjugate identity proves CA7, and invertibility of $DP$
proves uniqueness. Its exact correction to the constant linear
field is
\[
 X_B(x)-Bx=-\tfrac12\operatorname{adj}(DP(x))
                         [BP(x)-DP(x)Bx]. \tag{CA8}
\]
The possible zero corrections can also be classified completely.
If $DP(x)Bx=BP(x)$, comparison of linear terms forces $B$ to
commute with $DP(0)=\operatorname{diag}(-i,1,2i,-1)$.
These four eigenvalues are distinct, so $B$ is diagonal.
Writing its diagonal $(b_a,b_y,b_z,b_w)$, the monomials
$a^2y,a^3z$ in $P_0$ force $b_y=-b_a$, $b_z=-2b_a$;
the monomial $a^2w$ in $P_2$ forces $b_w=-3b_a$.
Conversely every monomial of each displayed component has
that component's weight for $(1,-1,-2,-3)$. Thus
\[
 \{B:DP(x)Bx=BP(x)\text{ for all }x\}
       =\C\operatorname{diag}(1,-1,-2,-3). \tag{CA9}
\]
For the corner action $B_c$, its trace is $2j$. A matrix
similar to a multiple of that displayed weight matrix would
have multiplier $-2j/5$, a real number, and therefore real
eigenvalues. The four eigenvalues $j\rho_l$ are nonreal.
Hence $B_c$ is not in that similarity class; CA7 supplies its
actual nonlinear lift, rather than an assumed linear symmetry.

Retain the CA1--2 ordering and put $\lambda_3=\lambda_\mathrm T$,
$\lambda_4=\lambda_\mathrm E$,
$\omega=3\lambda_3-2\lambda_4
=j(1/2-\delta_*+5i\gamma_*)\ne0$.
The following is an actual target orbit in the corner receiver:
\[
 u(t)=(0,0,-3e^{\lambda_3t},e^{\lambda_4t})^T,\qquad
 u'=B_cu,\qquad
 t_*=-\log(4)/\omega. \tag{CA10}
\]
Its original quotient image is $I_DJ_cu(t)=
e^{tM_j}I_DJ_cu(0)$ by CA6. The auxiliary factor cubic is
$h_t(x)=x^3-3p(t)x+q(t)$, with
$p(t)=e^{\lambda_3t}$ and $q(t)=e^{\lambda_4t}$.
Its discriminant is exactly
\[
 27e^{2\lambda_4t}(4e^{\omega t}-1), \tag{CA11}
\]
by the determinant definition or direct expansion of the cubic
resultant. This has a simple zero at $t_*$, because its
derivative there is $27e^{2\lambda_4t_*}\omega\ne0$.
Along $t=st_*$, $0\le s<1$, it is nonzero, since
$4e^{\omega st_*}=4^{1-s}>1$.
At $t_*$ set $r_*=q(t_*)/(2p(t_*))$. The equality
$q(t_*)^2=4p(t_*)^3$ gives $r_*^2=p(t_*)\ne0$,
and $h_{t_*}(x)=(x-r_*)^2(x+2r_*)$.

The exact leading local root behavior and its constants follow
without a root-identification guess. Put $t=t_*+s^2$ and
$x=r_*+s v$. Dividing $h_{t_*+s^2}(r_*+sv)$ by $s^2$
extends analytically to $s=0$, with value
$3r_*v^2-\omega r_*^3$ there. Its two roots
$v_\pm(0)=\pm r_*\sqrt{\omega/3}$ are nonzero and simple.
The analytic implicit-function theorem therefore gives two
actual root branches with
\[
 r_\pm(t)=r_*\pm r_*\sqrt{\omega/3}(t-t_*)^{1/2}
                   +O(t-t_*),\qquad
 h_t'(r_\pm(t))=
   c_\pm(t-t_*)^{1/2}+O(t-t_*),\quad
 c_\pm=\pm6r_*^2\sqrt{\omega/3}\ne0.
 \tag{CA12}
\]
The chosen square root fixes the branch and its sign; both are
retained. The four inverse factor states on these two roots
have
\[
 a(t)=\kappa(t-t_*)^{-1/4}(1+O((t-t_*)^{1/2})),
 \quad\kappa^2=c_\pm^{-1},\qquad
 y(t)\to-r_*,\ z(t)\to0,\ w(t)\to0. \tag{CA13}
\]
Indeed GF7 gives $a^2h_t'(r)=1$ and all three remaining
coordinates explicitly; here its coefficient $A=0$ and
$B=0$ throughout the orbit. Each colliding branch extends
backwards to $t=0$ along the stated segment, because the
cubic roots remain simple and finite there and their
nonzero derivatives provide continuous square-root lifts.
On every compact subsegment the resulting original source
coordinates are bounded. Differentiating $P(x(t))=u(t)$
then gives $x'=X_{B_c}(x)$ by CA7 and the nonzero Jacobian.
Thus these are actual trajectories of the polynomial field.

In the original corner metric their exact escape constant is
\[
 |t-t_*|^{1/4}\|J_c x(t)\|_{\Gamma,W_D}
           \longrightarrow |\kappa|\sqrt{(H_c)_{11}}>0.
 \tag{CA14}
\]
The same statement with $K_c$ uses its full original source
Gram. It follows by multiplying CA13 by $(t-t_*)^{1/4}$
on a chosen branch and using continuity of the fixed norms.
The target $J_cu(t)$ and its arithmetic quotient remain
finite at this complex time. A trajectory with unbounded
source norm has no holomorphic extension across that time,
so the polynomial lift is not complete in complex time.
No assertion of real-time escape follows from this calculation.
Through CA5 the two nonlinear vector fields, their trajectories,
and this singularity pass between the original high-degree
source and target by $T_D$ exactly.

Finally for any retained positive metric $H_c$ put
$G_P(x)=DP(x)^*H_cDP(x)$. The underlying real flow $\phi_t$
of $X_{B_c}$ satisfies $P\phi_t=e^{tB_c}P$ wherever defined.
Differentiating in the initial condition gives
$DP(\phi_t(x))D\phi_t(x)=e^{tB_c}DP(x)$.
Pulling back the metric and differentiating at real $t=0$
therefore proves the full Lie-derivative identity
\[
 \mathcal L_{X_{B_c}}G_P
     =DP^*(B_c^*H_c+H_cB_c)DP. \tag{CA15}
\]
It includes the derivative of the moving metric and both
vector-field derivatives. For every real center $c_0$, the
generalized characteristic polynomial of
$(\mathcal L_{X_{B_c}}G_P-2c_0G_P,G_P)$ equals that of
$(B_c^*H_c+H_cB_c-2c_0H_c,H_c)$: take the determinant of
the congruence and divide by $\det G_P=|\det DP|^2\det H_c$.
Thus the original metric defect is transported with its full
spectrum; the nonlinear escape does not set it to zero.
% END CORNER_ACTION

% BEGIN ACTUAL_MARKED_OBSERVATION
\section{The four original upper corners in the actual observation}
This continuation uses the actual invariant observation
$\Lambda_k:E_k\to B_k$, rather than selecting a lift into $B_k$.
Its original period and all its unit factors are those in FC13.
Retain $j=k-8$, $k\equiv1\pmod4$, $k\ge9$ and the two complete
grid polynomials in CA1, with $k$ and $j$ respectively. Define
$J_+:\C^4\to E_k$ and $J_-:\C^4\to E_j$ by the original
four corner cardinal classes in the order
$(\mathrm M,\mathrm N,\mathrm T,\mathrm E)$.
Their roots are $k\rho_l$ and $j\rho_l$. Write $p_l^+,p_l^-$
for their unique representatives of degrees less than $(k+1)^2$
and $(j+1)^2$. Every denominator is its original
$\chi_k'(k\rho_l)$ or $\chi_j'(j\rho_l)$.

Let $C:E_k\to E_j$ be the original conductor quotient map and
$\bar C:B_k\to E_j$ the map with $\bar C\Lambda_k=C$.
This is an actual factorization: the conductor-multiple invariant
polynomials give $\ker\Lambda_k\subset\ker C$, as proved in
\cite[OB3]{Observed}. The formula
$\bar C(\Lambda_kx)=Cx$ is independent of the choice of lift
by that kernel inclusion. Put
\[
 D_\partial=\diag(a_{88},a_{80},a_{08},a_{00}),\qquad
 O_+=\Lambda_kJ_+,\qquad
 CJ_+=J_-D_\partial,\quad
 \bar C O_+=J_-D_\partial. \tag{MO1}
\]
To prove the third identity, evaluate the finite-shift conductor
at every lower grid root. A shifted lower root reaches an upper
corner only when the lower root is that same corner and the
shift is the matching one among $(8,8),(8,0),(0,8),(0,0)$.
The values are exactly the corresponding original coefficients,
so equality of classes follows by interpolation. This proves
all four columns, including their actual phases.

Before reduction, the entire residual is retained:
\[
 T_Ap_l^+(S')=d_l p_l^-(S')+\chi_j(S')Z_l(S'),\qquad
 Z_l(S')=\frac{\sum_{u,w=0}^8a_{uw}
 p_l^+(S'+\beta_{8,u,w})-d_l p_l^-(S')}{\chi_j(S')},
 \tag{MO2}
\]
where $d_l=(D_\partial)_{ll}$.
The numerator vanishes at every distinct lower grid root by the
preceding evaluation proof, so $Z_l$ is a polynomial, not a
meromorphic substitution. Its degree is at most
$(k+1)^2-(j+1)^2-v-1=16k-49-v$; a negative bound means the
zero polynomial. MO2 gives every coefficient by exact finite
division of its displayed numerator. Thus MO1 has not silently
deleted a polynomial relation from the original WCF norm.

On the explicit original large-period domain of
\cite[OC50--74]{CornerNonzero}, every $d_l$ is nonzero. In the
original notation that domain is $\gamma\ge1000\delta$ and
$|\varpi|\ge R_{\rm corner}^{\rm explicit}$, where all constants
are given in OC53, OC64 and OC67--71; the exact lower bound is
$|d_l|\ge(m_{\rm corner}/2)^4$. This is an already evaluated
period-domain result, with the unit phase retained, rather than
an assumed coefficient sign. MO1 then proves that $O_+$ is
injective: if $O_+c=0$, injectivity of $J_-$ and invertibility
of $D_\partial$ give $c=0$. At every other original period,
the exact partial-support treatment AO1--23 below applies and
does not identify $d_l=0$ with $\Lambda_ke_l=0$.

Let $A_k$ and $A_j$ be the original multiplication operators
on $E_k,E_j$. They act on these corners by
\[
 A_kJ_+=J_+ \mathsf B_k,\quad A_jJ_-=J_-\mathsf B_j,\qquad
 \mathsf B_k=\diag(k\rho_l),\quad \mathsf B_j=\diag(j\rho_l).
 \tag{MO3}
\]
These are identities of the actual quotient operators by the
cardinal formula CA6. On the four-dimensional image of $O_+$,
the actual observed action is
$\Theta_+=O_+\mathsf B_kO_+^{-1}$: directly
$\Lambda_kA_kJ_+=\Theta_+\Lambda_kJ_+$.
No action on the full one-step quotient $B_k$ is inferred.
Its conductor defect is completely explicit:
\[
 (\bar C\Theta_+-A_j\bar C)O_+
       =8J_-D_\partial\diag(\rho_l),\qquad
 \bar C\Theta_+=\frac{k}{j}A_j\bar C
                 \quad\text{on }\operatorname{im}O_+.
 \tag{MO4}
\]
This follows from MO1--3 and the fact that the two displayed
diagonal matrices commute. The factor $k/j$ is retained;
neither original arithmetic action has been replaced by the other.

\section{The exact original minimum correction on these four corners}
Fix $N\ge(k+1)^2-1$ and the original polynomial Gram $H_N$
of $m_k(y)dy=w_h^{*k}(y)dy$ at $S=k/2+iy$, with
$w_h(y)=|(2\xi/h)(1/2+iy)|^2/(2\pi)$ and the original
quartet factor $h(S)=\prod_l(S-\rho_l)$.
Let $J_N:\mathcal P_N\to E_k$ be reduction and retain
\[
 G_N=(J_NH_N^{-1}J_N^*)^{-1},\quad
 Q_N=(\Lambda_kG_N^{-1}\Lambda_k^*)^{-1},\quad
 T_N=(\bar C Q_N^{-1}\bar C^*)^{-1},\quad
 R_N=Q_N^{-1}\bar C^*T_N. \tag{MO5}
\]
These are the original attained metrics of \cite[EQ1--5]{EightMinimum}.
Their inverses exist: the reductions and observation are onto their
defined images, and $C$ is onto because its triangular polynomial
map drops degree by $v\le80$ while the grid-degree difference is
$16k-48\ge96$. The actual density is positive almost everywhere;
its isolated zero set cannot annihilate a nonzero polynomial, so
the finite Grams are positive. The moments and original total mass
in $H_N$ are not replaced by a Gamma surrogate.

The difference between the original observation and its minimum
lift through the conductor is the exact matrix
\[
 \Delta_N=O_+-R_NJ_-D_\partial,\qquad
 \bar C\Delta_N=0. \tag{MO6}
\]
Since $R_N^*Q_N=T_N\bar C$, the two summands are orthogonal
in the original $Q_N$ metric. Thus the complete four-corner
observed Gram is
\[
 H_+=O_+^*Q_NO_+
 =D_\partial^*J_-^*T_NJ_-D_\partial+
                   \Delta_N^*Q_N\Delta_N. \tag{MO7}
\]
No assertion is made that the second term vanishes.
There is a separate exact correction from the original source
quotient: put $M_N=G_N^{-1}\Lambda_k^*Q_N$ and
$Z_N=(I-M_N\Lambda_k)J_+$. Then
$\Lambda_kM_N=I$, $M_N^*G_NM_N=Q_N$,
$\Lambda_kZ_N=0$, and direct orthogonal decomposition gives
\[
 J_+^*G_NJ_+=H_++Z_N^*G_NZ_N. \tag{MO8}
\]
Equations MO7 and MO8 retain the two different kernels and both
original positive residual energies. They also show precisely
which norms accompany an original observed trajectory.

\section{The literal Fable marking and its arithmetic generator}
The original four collision points are $q_l=Ke_l$, with the
complete invertible matrix $K$ in FC15a. Therefore the receiver
which sends each named collision point to its corresponding
original observed corner is
\[
 F_+=J_+K^{-1},\quad F_B=O_+K^{-1},\quad
 F_-=J_-D_\partial K^{-1},\qquad
 \Lambda_kF_+=F_B,\quad\bar CF_B=F_-.
 \tag{MO9}
\]
In particular $F_Bq_l=\Lambda_k[p_l^+]$, with all four labels
retained individually. On the nonzero-corner domain all three
maps are injections onto their indicated four-dimensional images.
Their nonlinear conjugates satisfy, exactly on these images,
$\Lambda_k(F_+PF_+^{-1})=(F_BPF_B^{-1})\Lambda_k$ and
$\bar C(F_BPF_B^{-1})=(F_-PF_-^{-1})\bar C$.
The polynomial relation before this second equality is precisely
MO2; it is not asserted to have zero original polynomial norm.

The arithmetic matrix in the unchanged Fable coordinates is
\[
 \widehat B_k=K\diag(k\rho_l)K^{-1},\qquad
 \Theta_+F_B=F_B\widehat B_k,
 \qquad \widehat G_N=K^{-*}H_+K^{-1}>0.
 \tag{MO10}
\]
This formula gives every coefficient through the displayed FC15a
$K$ and its exact adjugate inverse; all original signs and phases
remain. Replacing $\widehat B_k$ by $\diag(k\rho_l)$ in the
polynomial coordinates would change this marking. CA's different
corner-coefficient receiver remains related by the explicit $K$.
The unique polynomial lift of MO10 is
$X_{\widehat B_k}=-\operatorname{adj}(DP)\widehat B_kP/2$,
by the original determinant $-2$. Its image through $F_B$ is
therefore the exact polynomial lift of the actual observed
arithmetic action restricted to the four original corners.

\section{Every point of the original common fibre along its actual orbit}
Let $h_0=(0,0,-1,0)^T=2iKe_{\mathrm M}$ and
$\lambda=k\rho_\mathrm M$.
MO10 proves $\widehat B_k h_0=\lambda h_0$.
Set $b(t)=e^{\lambda t/2}$; this entire function never vanishes.
The complete original fibre over $e^{\lambda t}h_0$ consists of
the following seven trajectories:
\[
 \begin{aligned}
 q_\infty(t)&=(0,0,ib(t)^2/2,0)^T,\\
 q_{0,\epsilon}(t)&=
 (\epsilon i/b(t),-\epsilon b(t),-3ib(t)^2,
                       13\epsilon b(t)^3)^T,
                         &&\epsilon\in\{1,-1\},\\
 q_{\eta,\epsilon}(t)&=
 \bigl(\epsilon/(\sqrt2 b(t)),
  -(\eta+i\epsilon\sqrt2)b(t),
  (2\eta\epsilon\sqrt2+6i)b(t)^2,
  (-34\eta-19i\epsilon\sqrt2)b(t)^3\bigr)^T,
                         &&\eta,\epsilon\in\{1,-1\}.
 \end{aligned} \tag{MO11}
\]
GF7 gives them from the roots $0,\pm b(t)$ with derivatives
$-b(t)^2,2b(t)^2,2b(t)^2$, and GF8 gives the infinity point.
Substituting in the full original $P$ proves
$P(q(t))=e^{\lambda t}h_0$. Conversely that cubic has precisely
these three simple finite roots and the one chart infinity root,
so GF's full fibre classification proves completeness of the list.
The original four named points occur at $t=0$ as
$q_\infty,q_{0,+},q_{+,+},q_{-,+}$ in order M,N,T,E.

Differentiating their exact image equality and using $DP^{-1}$
shows that every trajectory solves $q'=X_{\widehat B_k}(q)$.
All seven are entire functions of complex time, so none escapes
at a finite complex time on this particular original orbit.
After the target period $2\pi i/\lambda$, $b$ changes sign:
$q_\infty$ returns, $q_{0,\epsilon}$ is exchanged with
$q_{0,-\epsilon}$, and $q_{\eta,\epsilon}$ with
$q_{-\eta,-\epsilon}$. This computes the original fibre's
continuation without merging any two of its marked states.

\section{An exact marked arithmetic orbit with five escaping branches}
The same original generator nevertheless has an explicit finite
boundary crossing. Put
\[
 u_*=(0,1,0,0)^T,\qquad
 u(t)=e^{t\widehat B_k}u_*
     =K\diag(e^{k\rho_l t})K^{-1}u_*.
 \tag{MO12}
\]
This is a complete formula in the fixed original period. Direct
multiplication of FC15a and its inverse gives
\[
 \begin{aligned}
 \nu_A:=A'(0)&=\frac{k}{154}
       [-68\gamma+i(25\sqrt2\gamma-136\delta)]\ne0,\\
 \nu_D:=D'(0)&=\frac{k}{77}
       [-884\delta+(1631\sqrt2+442i)\gamma]\ne0.
 \end{aligned} \tag{MO13}
\]
Nonvanishing follows from the nonzero real part of the first
and the positive imaginary part of the second. These identities
are checked exactly over $\mathbb Q(i,\sqrt2)$ by
\texttt{derive\_marked\_arithmetic\_generator.py}.
At $u_*$ the quartic is $h_*(x)=x^2(x+1)$.
The full degree-four discriminant polynomial, retaining its
fixed cubic coefficient, has gradient $(0,0,0,-4)$ there.
This follows by differentiation of its resultant determinant
(and is reproduced by \texttt{verify\_marked\_fibre\_flow.py}).
Thus the discriminant along MO12 has derivative $-4\nu_D\ne0$.

There is a completely specified punctured time disk containing
only eight-state fibres. For an explicit radius form the finite
polynomial
\[
 \mathcal D(z_1,z_2,z_3,z_4)=
 \operatorname{Disc}_x\bigl(A(z)x^4+x^3+B(z)x^2+C(z)x+D(z)\bigr)
 =\sum_\mu c_\mu z^\mu,
 \quad (A(z),B(z),C(z),D(z))^T
       =K\diag(z_l)K^{-1}u_*.
 \tag{MO14}
\]
This specifies every coefficient by the finite resultant and the
literal original matrix. Write $A(z)=\sum_l\alpha_lz_l$ and
$\theta_\mu=k\sum_l\mu_l\rho_l$, and retain the finite sums
\[
 M_D=\sum_\mu|c_\mu|\,|\theta_\mu|^2e^{|\theta_\mu|},\quad
 M_A=\sum_l|\alpha_l|\,|k\rho_l|^2e^{|k\rho_l|},\quad
 r_0=\min\{1,4|\nu_D|/M_D,|\nu_A|/M_A\}>0.
 \tag{MO15}
\]
Both sums are nonzero since the respective first derivatives
are nonzero. Taylor's integral remainder for each exponential
gives errors at most $M_D|t|^2/2$ and $M_A|t|^2/2$ for
$|t|\le1$. Their constant terms vanish. It follows that,
for $0<|t|\le r_0$, the discriminant and $A(t)$ have moduli
at least $2|\nu_D||t|$ and $|\nu_A||t|/2$ respectively.
GF16 therefore gives exactly eight inverse states throughout
this punctured disk. The radius does not require a numerical
search for zeros or a discarded higher Taylor term.

Four states select the two roots colliding at zero. Substitute
$t=s^2$ and $x=sv$ in the full $h_{u(t)}(x)$ and divide by
$s^2$. The analytic extension at zero is $v^2+\nu_D$,
whose two roots $v_\pm=\pm\sqrt{-\nu_D}$ are simple.
Analytic implicit inversion gives
\[
 r_\pm(t)=v_\pm t^{1/2}+O(t),\quad
 h_{u(t)}'(r_\pm(t))=2v_\pm t^{1/2}+O(t),\quad
 t^{1/4}q_{\pm,\epsilon}(t)\longrightarrow\kappa_{\pm,\epsilon}e_a,
 \quad\kappa_{\pm,\epsilon}^2=(2v_\pm)^{-1}.
 \tag{MO16}
\]
Both values of each square root are retained on a chosen time
sector. The final limit follows from all four coordinates of
GF7: $a\sim\kappa t^{-1/4}$, while $y,z,w\to0$.

The projective root at infinity is simple at $t=0$ and supplies
two more factor states. In the negative chart its finite limiting
coordinates are $(a,Y,Z,W)=(0,1,-7i/2,11)$.
The smooth local inverse of $\pi$ from FX1--8 gives its analytic
lift, and $A=ac$ with $c(0)=i$ proves
$a'(0)=-i\nu_A$. Applying the complete transition FX7 gives
\[
 t^3q_{\rm lost}(t)\longrightarrow
                     -40\nu_A^{-3}e_w. \tag{MO17}
\]
Indeed the third-order term is $-40i/a^3$, with
$a=-i\nu_A t+O(t^2)$; its sign and full coefficient are as
displayed. The other terms of FX7 have orders at most two.
The positive infinity state is finite. The remaining root
at $-1$ is simple and supplies two finite states. Their three
limits in the original positive chart are exactly
\[
 (0,1,7i/2,11)^T,\qquad
 (\epsilon,1-i\epsilon,-2\epsilon+3i,18-6i\epsilon)^T,
                       \quad\epsilon\in\{1,-1\}. \tag{MO18}
\]
Substitution in $P$ verifies their common image $u_*$, and
the nonzero constant Jacobian gives their local analytic lifts.
The fibre count eight on the punctured disk and three at its
center proves that MO16--18 exhaust every branch: precisely
four intrinsic root-collision escapes, one missing-chart escape,
and three finite states.

Every inverse branch satisfies the actual polynomial lifted
equation $q'=X_{\widehat B_k}(q)$ by differentiating MO12.
For example start at $t=-r_0/2$ and continue in a sector toward
zero; the elapsed complex time is finite, with no intermediate
discriminant or $A$ zero by MO15. The original observed target
$F_Bu(t)$ follows $\Theta_+$ and stays finite. In the complete
original $Q_N$ norm the five source escape constants are
\[
 |t|^{1/4}\|F_Bq_{\pm,\epsilon}(t)\|_{Q_N}
       \longrightarrow|\kappa_{\pm,\epsilon}|
                             \sqrt{(\widehat G_N)_{11}},\qquad
 |t|^3\|F_Bq_{\rm lost}(t)\|_{Q_N}
       \longrightarrow\frac{40}{|\nu_A|^3}
                             \sqrt{(\widehat G_N)_{44}}.
 \tag{MO19}
\]
Positivity follows from MO7--10. These constants keep both
the original observation and its full minimum residual.
At other periods the actual stack of AO gives the corresponding
nonzero constants whenever any original corner is observable;
it preserves both $K^{-1}e_a$ and $K^{-1}e_w$ and every invisible
kernel, rather than assigning positivity to a vanished column.

At a partially invisible period, the visible stack by itself is
not declared to carry an induced nonlinear $P$. The full relation
has an exact receiver with the invisible labels retained. Let
$I$ and $Z=I^c$ be AO's actual active and invisible sets and let
$\mathcal R_I$ denote the image of its $m$-stack. Define
\[
 \mathcal F_{\rm full}:\C^4\longrightarrow
              \mathcal R_I\oplus\C^Z,\qquad
 q\longmapsto\bigl(Z_mK^{-1}q,(K^{-1}q)_Z\bigr),\qquad
 \mathcal F_{\rm full}^{-1}(y,c_Z)
       =K\bigl(\mathcal R_{\rm AO}y,c_Z\bigr).
 \tag{MO20}
\]
Here $\mathcal R_{\rm AO}$ is the complete active-coordinate
inverse AO7, and the parentheses on the right place the two
vectors in their original label positions, without reordering
them. For $I$ empty the visible component is zero and this is
just $q\mapsto K^{-1}q$ on the four retained invisible labels.
AO7 verifies both inverse identities in every case. Consequently
$\mathcal F_{\rm full}P\mathcal F_{\rm full}^{-1}$ and the
transport of each inverse branch are exact maps on this complete
receiver, and AO11 proves its arithmetic action is
$\mathcal S|_{\mathcal R_I}\oplus\diag(k\rho_l)_{l\in Z}$.
Equip it with the original full corner norm by assigning
$\|(y,c_Z)\|^2=(\mathcal R_{\rm AO}y,c_Z)^*
(J_+^*G_NJ_+)(\mathcal R_{\rm AO}y,c_Z)$.
This retains all cross terms; minimizing over $c_Z$ is exactly
the Schur formula AO13. Thus no invisible direction is deleted
in the completed nonlinear or arithmetic correspondence.

Finally the exact metric action identity CA15 applies with
$B=\widehat B_k$ and $H=\widehat G_N$; its proof used only the
complete polynomial inverse Jacobian. Together with MO2, MO4,
MO7 and MO8 this records all original polynomial, action and
metric residuals in the actual observed realization. The fixed
original linear conductor still has the inverse in FC24.
% END ACTUAL_MARKED_OBSERVATION

% BEGIN ACTUAL_CORNER_OBSERVATION


This calculation uses the original observation $\Lambda_k$, the
original upper corner idempotents, the full original output Gram,
and the actual four conductor corner coefficients. The coefficient
test and conductor lift are the explicit maps in
\cite[O1--3 and O10]{Observe}. The four labelled Fable source
columns are those of \cite{ES488}, and their original-conductor
receivers are retained from the FC/CA sections above. This note does not
replace the upper corner classes by their minimum-norm lifts.
It proves the restriction and its iterated recovery even when
some actual observation columns vanish.

\section{The actual invariant corner subspace}
Fix an original admitted simple quartet
\[
 (\rho_{\mathrm M},\rho_{\mathrm N},\rho_{\mathrm T},\rho_{\mathrm E})
 =\left(\tfrac12+\delta+i\gamma,\tfrac12+\delta-i\gamma,
        \tfrac12-\delta+i\gamma,\tfrac12-\delta-i\gamma\right),
 \quad 0<\delta<\tfrac12,\quad\gamma>2.
\]
At the original degree $k\ge9$ let $e_l\in E_k$ be the four
upper-grid root idempotents with indices
$(k,k),(k,0),(0,k),(0,0)$ in this order. Put
\begin{equation}
 J_+:\C^4\longrightarrow E_k,\quad J_+c=\sum_lc_le_l,
 \qquad D_k=\diag(\lambda_l),\quad\lambda_l=k\rho_l.
 \tag{AO1}
\end{equation}
The original arithmetic action satisfies exactly
$A_kJ_+=J_+D_k$, since each $e_l$ is its actual root idempotent.
All four $\lambda_l$ are distinct. Hence the four-dimensional
image of $J_+$ is invariant under that action.

Let $\mathcal B_k$ be the original output space, with its
complete positive Hermitian Gram matrix $Q$. Set
\begin{equation}
 b_l=\Lambda_k e_l,\qquad
 B=(b_{\mathrm M}\ b_{\mathrm N}\ b_{\mathrm T}\ b_{\mathrm E})
       =\Lambda_kJ_+,
 \qquad I=\{l:b_l\ne0\},\quad m=|I|.
 \tag{AO2}
\end{equation}
The letters $B$ and $Q$ in this note always denote these
actual observation data. No column is normalized. Inner
products are conjugate linear in the first argument.
On the corner subspace, for every integer $n\ge0$,
\begin{equation}
 \Lambda_k A_k^nJ_+c=B D_k^n c.\tag{AO3}
\end{equation}
This is obtained by iterating AO1. It does not require the
one-step kernel of $\Lambda_k$ in the entire $E_k$ to be invariant,
and makes no assertion that the whole observation quotient
inherits $A_k$.

\section{Full stacked Gram and exact recovery with at most four observations}
For $p\ge1$ retain the unscaled stack
\begin{equation}
 Z_p c=(BD_k^t c)_{t=0}^{p-1}\in\mathcal B_k^{\oplus p},
 \qquad Q_p=I_p\otimes Q.
 \tag{AO4}
\end{equation}
Its full original Gram has entries
\begin{equation}
 (K_p)_{lr}=(Z_p^*Q_pZ_p)_{lr}
   =b_l^*Qb_r\sum_{t=0}^{p-1}\overline{\lambda_l}^{\,t}\lambda_r^t.
 \tag{AO5}
\end{equation}
These are exact finite sums with all original root scales and
off-diagonal output correlations retained. If $l\notin I$,
the corresponding row and column are zero.

Suppose first $m>0$. For each $l\in I$ define the actual spectral
separating polynomial and actual output functional
\begin{equation}
 p_l(S)=\prod_{r\in I\setminus\{l\}}
              \frac{S-\lambda_r}{\lambda_l-\lambda_r}
        =\sum_{t=0}^{m-1}p_{l,t}S^t,
 \qquad
 \vartheta_l(y)=\frac{b_l^*Qy}{b_l^*Qb_l}.
 \tag{AO6}
\end{equation}
All denominators are nonzero because the original roots are
distinct and $Q>0$. The exact inverse on the active coordinates is
\begin{equation}
 c_l=\sum_{t=0}^{m-1}p_{l,t}\vartheta_l(y_t),
 \qquad (y_t)_{t=0}^{m-1}=Z_m c,\quad l\in I.
 \tag{AO7}
\end{equation}
To prove it, substitute $y_t=\sum_rb_r\lambda_r^tc_r$.
The sum over $t$ becomes
$\sum_{r\in I}p_l(\lambda_r)\vartheta_l(b_r)c_r=c_l$,
since $p_l(\lambda_r)=\delta_{lr}$ and
$\vartheta_l(b_l)=1$. Thus
\begin{equation}
 \ker Z_m=\operatorname{span}\{e_l^{\rm label}:l\notin I\},
 \qquad K_{m,I}>0.
 \tag{AO8}
\end{equation}
Here $e_l^{\rm label}$ is a coordinate vector of $\C^4$, and
$K_{m,I}$ is the active principal submatrix of AO5.
In particular even linearly dependent or collinear nonzero
$b_l$ are recovered by their original, distinct arithmetic
eigenvalues. At most $m\le4$ observations are needed.
If $m=0$, every stack is identically zero, and AO8 is to be
read as the full four-dimensional kernel; no inverse is assigned.

The polynomial
\begin{equation}
 \chi_I(S)=\prod_{l\in I}(S-\lambda_l)
       =S^m+\sum_{t=0}^{m-1}c_tS^t\tag{AO9}
\end{equation}
is the minimal polynomial of the actual active corner action.
Indeed an annihilating polynomial must vanish at each of its
$m$ distinct eigenvalues; divisibility by their product proves
minimality. The degree-$m-1$ separation in AO6 uses precisely
these roots, without the other $(k+1)^2-4$ grid roots.
For the particular columns the least number of observations is
exactly
\[
 p_{\min}=\min\{1\le p\le m:\rank K_{p,I}=m\}.
\]
This is a finite criterion using the full original matrices.
The upper bound $m$ is sharp in general: when all active columns
are nonzero multiples of one fixed vector, a stack of $p<m$
observations has rank at most $p$, so cannot recover $m$
coordinates. This last example states sharpness of the proved
general bound, not a replacement of the programme's columns.

\section{Exact finite action on the recovered observation space}
Let $\mathcal R_I=\operatorname{im}Z_m$ for $m>0$. Define on the
ambient stack the literal block companion operator
\begin{equation}
 \mathcal S(y_0,\ldots,y_{m-1})
    =(y_1,\ldots,y_{m-1},-\sum_{t=0}^{m-1}c_ty_t),
 \tag{AO10}
\end{equation}
with the $m=1$ formula interpreted as $y_0\mapsto-c_0y_0$.
For actual corner data, AO9 gives
\begin{equation}
 \mathcal S Z_m=Z_mD_k.\tag{AO11}
\end{equation}
The first $m-1$ components are shifts of the original observed
iterates; the last follows from
$\lambda_l^m+\sum_tc_t\lambda_l^t=0$ for every active label.
Consequently $\mathcal R_I$ is invariant, and the restriction
of $\mathcal S$ to it is exactly conjugate to the active
diagonal action $D_{k,I}$ by AO7. Its exponential satisfies
$e^{z\mathcal S}Z_m=Z_m e^{zD_k}$ for every complex $z$ by
the convergent finite-dimensional matrix series. This is a
constructed action on the actual stacked data, not an action
postulated on the whole one-step quotient.

The original inherited stack metric on active coordinates is
$K_{m,I}$. Its complete infinitesimal Hermitian defect is
\begin{equation}
 D_{k,I}^*K_{m,I}+K_{m,I}D_{k,I},\qquad
 [D_{k,I}^*K_{m,I}+K_{m,I}D_{k,I}]_{lr}
       =(\overline{\lambda_l}+\lambda_r)(K_m)_{lr}.
 \tag{AO12}
\end{equation}
Subtracting the original central term $kK_{m,I}$ leaves
$(\overline{\lambda_l}+\lambda_r-k)(K_m)_{lr}$.
On the single active eigenvector of label $l$ its value is
$2k\varepsilon_l\delta\,(K_m)_{ll}$, where
$\varepsilon_l$ is the original real-part sign. This proves
both signs whenever the active set contains both real-part
signs, with the entire original Gram still present.

For completeness, if $H=J_+^*G_{E_k}J_+>0$ is the original
corner Gram of a specified physical source metric, its quotient
by the invisible coordinate subspace has the exact Schur Gram
\begin{equation}
 H_{\rm quot}=H_{II}-H_{IZ}H_{ZZ}^{-1}H_{ZI},
 \qquad Z=I^c.
 \tag{AO13}
\end{equation}
With $Z$ empty this means $H_{II}$; with $I$ empty no inverse
is asserted. Complete the square in
$(c_I,c_Z)^*H(c_I,c_Z)$: its minimum occurs at
$c_Z=-H_{ZZ}^{-1}H_{ZI}c_I$ and gives AO13. This minimum
represents only the norm of the recovered corner quotient;
the receiver AO1 still consists of the original upper
idempotents. The exact squared inverse norm from the actual
stack to this quotient is
\begin{equation}
 \sup_{c_I\ne0}\frac{c_I^*H_{\rm quot}c_I}
                         {c_I^*K_{m,I}c_I}
 =\lambda_{\max}
     (K_{m,I}^{-1/2}H_{\rm quot}K_{m,I}^{-1/2}).\tag{AO14}
\end{equation}
This equality follows by substituting $K_{m,I}^{1/2}c_I$ in
the Rayleigh quotient. It retains both original metrics and
does not assign a small or degree-uniform value to this
finite expression.

\section{What the four actual conductor coefficients prove}
Set $j=k-8$. Let $J_-:\C^4\to E_j$ insert the four lower
corner idempotents in their same marked order. The original
conductor coefficient map has the exact corner restriction
\begin{equation}
 C J_+=J_-D_{\rm edge},\qquad
 D_{\rm edge}=\diag(d_{\mathrm M},d_{\mathrm N},d_{\mathrm T},d_{\mathrm E})
      =\diag(a_{88},a_{80},a_{08},a_{00}).\tag{AO15}
\end{equation}
This is the corner term of the original finite convolution,
with no other shift omitted: an extreme upper index can be
formed only from the corresponding extreme lower index
and the extreme degree-eight coefficient. Equivalently the
four entries equal the original point evaluations
$A_H(He_l)$, proved in FC16a of the FC/CA sections above.

Each nonzero $d_l$ certifies that $b_l\ne0$. To prove this
directly from the original observation, use the conductor
lift \cite[O3]{Observe}. Multiply the actual degree-eight
conductor biform by the lower corner monomial with index
$(j,j),(j,0),(0,j)$ or $(0,0)$. Its coefficient at the matching
upper corner is exactly $d_l$; there is no other contribution
to that extreme coefficient. The original five-term invariant
lift in O3 realizes this biform in $\Theta_k(S_k^G)$.
The original coefficient criterion O1 then gives
\begin{equation}
 I_{\rm edge}:=\{l:d_l\ne0\}\ \subseteq\ I.\tag{AO16}
\end{equation}
If $d_l=0$, this particular coefficient test is zero, but the
full column vanishes only if every original invariant
coefficient test in O1 is zero. Thus $I$ in AO2 is the exact
active set; $I_{\rm edge}$ is a proved subset, not an assumed
equality. On any locus where all four actual coefficients
are nonzero, $I$ consists of all four labels.

The arithmetic shift across AO15 is also exact:
\begin{equation}
 C A_kJ_+-A_j C J_+
       =J_-D_{\rm edge}\diag(8\rho_l).\tag{AO17}
\end{equation}
Indeed use AO1 and its lower counterpart
$A_jJ_-=J_-\diag(j\rho_l)$; the two diagonal matrices commute,
and $k-j=8$. This retains the original additive shift and
does not identify the upper and lower arithmetic actions.

\section{The labelled Fable receiver and all previously computed escapes}
Let the four original Fable collision columns be $q_l=Ke_l^{\rm label}$,
where
\[
 K=\begin{pmatrix}
 0&i&1/\sqrt2&1/\sqrt2\\
 0&-1&-1-\sqrt2i&1-\sqrt2i\\
 i/2&-3i&2\sqrt2+6i&-2\sqrt2+6i\\
 0&13&-34-19\sqrt2i&34-19\sqrt2i
 \end{pmatrix},\qquad \det K=77\sqrt2i/2.
\]
The label-preserving actual upper observation and its stack are
\begin{equation}
 q\longmapsto BK^{-1}q,\qquad
 \mathcal F_m(q)=Z_mK^{-1}q.\tag{AO18}
\end{equation}
At the four marked points these give respectively $b_l$ and
$(\lambda_l^t b_l)_{t=0}^{m-1}$. These formulas retain each
marked source identity; the different raw-coordinate insertion
$q\mapsto Bq$ does not have that property and is not substituted.
Their exact action satisfies
$\mathcal S\mathcal F_m=\mathcal F_m K D_kK^{-1}$ by AO11.

For an arbitrary original direction $v\in\C^4$, AO8 gives
the complete visibility test and its exact norm:
\begin{equation}
 \mathcal F_m(v)=0\iff (K^{-1}v)_l=0\text{ for every }l\in I,
 \qquad
 \|\mathcal F_m(v)\|_{Q_m}^2
    =(K^{-1}v)_I^*K_{m,I}(K^{-1}v)_I.
 \tag{AO19}
\end{equation}
In contrast a one-step observation can cancel nonzero columns;
AO19 is proved for the complete finite stack. The exact
one-step visibility test is simply $BK^{-1}v\ne0$, with
its full vector sum retained.

The two leading directions of the established Fable escapes
have particularly useful coefficients. With the original raw
coordinate vectors $e_a=(1,0,0,0)^{\mathsf T}$ and
$e_w=(0,0,0,1)^{\mathsf T}$, direct multiplication by $K$ gives
\begin{align}
 c_a:=K^{-1}e_a
 &=\left(-4\sqrt2-\frac{180i}{77},\ -\frac{30i}{77},\
          \frac{47\sqrt2-64i}{154},\
          \frac{47\sqrt2+64i}{154}\right)^{\mathsf T},\nonumber\\
 c_w:=K^{-1}e_w
 &=\left(\frac6{77},\frac1{77},\
          \frac{-3-\sqrt2i}{154},\
          \frac{3-\sqrt2i}{154}\right)^{\mathsf T}.
 \tag{AO20}
\end{align}
Every one of these eight entries is nonzero. Therefore either
direction is visible in the actual $m$-stack as soon as even
one actual corner column is nonzero. Its exact positive norm
is given by AO19. If $I$ is empty, both and every other
corner direction are invisible to all iterates.

In particular for the original sign-partner path
\[
 \widehat q_\epsilon=(-\epsilon,2i/\epsilon,
                     6i/\epsilon^2-i/2,40i/\epsilon^3),
 \qquad\epsilon>0,
\]
the literal vector limit $\epsilon^3\widehat q_\epsilon\to40ie_w$
and the fixed linearity of AO18 prove
\begin{equation}
 \epsilon^3\|\mathcal F_m(\widehat q_\epsilon)\|_{Q_m}
 \longrightarrow40\sqrt{(c_w)_I^*K_{m,I}(c_w)_I}>0
 \quad(I\ne\varnothing).
 \tag{AO21}
\end{equation}
For all eight auxiliary quartet-limit branches of GF25 above,
with the original receiver held fixed, the established coordinate
limit is $\sqrt\tau\,q_{\varepsilon,\eta,\pm}(\tau)
\to\kappa_{\varepsilon,\pm}e_a$, where
$|\kappa_{\varepsilon,\pm}|=1/(2\gamma_{\rm aux})$.
Here $\tau$ is the auxiliary quartet separation, not a change
in the fixed original arithmetic $\delta$ of AO1. Their exact
observed escape constants are
\begin{equation}
 \sqrt\tau\,\|\mathcal F_m(q_{\varepsilon,\eta,\pm}(\tau))\|_{Q_m}
 \longrightarrow\frac{\sqrt{(c_a)_I^*K_{m,I}(c_a)_I}}
                         {2\gamma_{\rm aux}}>0.
 \tag{AO22}
\end{equation}
The selected-root collision trajectories of CA13 above
have $(t-t_*)^{1/4}q(t)\to\kappa e_a$ on the retained branch;
the identical argument gives constant
$|\kappa|\sqrt{(c_a)_I^*K_{m,I}(c_a)_I}$ with the original
quarter-power rate. Every limit here follows from the actual
source coordinates and a fixed finite linear map, so no
estimate of an unproved one-step quotient action is involved.

These conclusions also have a direct conductor counterpart.
For any prescribed original positive lower-corner Gram
$H_->0$, the marked conductor image is
$J_-D_{\rm edge}K^{-1}q$. Its exact squared norm in a raw
direction $v$ is
\begin{equation}
 (K^{-1}v)^*D_{\rm edge}^*H_-D_{\rm edge}(K^{-1}v).
 \tag{AO23}
\end{equation}
For $v=e_a$ or $e_w$ this is strictly positive if and only
if at least one of the four explicit original $d_l$ is
nonzero, because AO20 has no zero entry and $H_->0$.
If all four $d_l$ vanish, the entire marked corner conductor
image is zero; the actual observation may still be nonzero,
and its definitive test remains AO2 and AO19. If even one
$d_l$ is nonzero, AO16 certifies that both Fable escape
directions survive the original finite observation stack.
On the explicit domain where all four $d_l$ are nonzero,
all four marked coordinates are recovered, with the exact
inverse AO7 and full original norm cost AO14.

\section{Exact coefficient certificate and provenance}
The accompanying certificate
\texttt{verify\_actual\_corner\_observation.py} verifies
the two complete columns AO20 over $\mathbb Q(i,\sqrt2)$,
their nonzero entries, and the actual quartet Vandermonde
determinant $-64k^6\delta^2\gamma^2(\delta^2+\gamma^2)$.
The Gram, recovery, and action assertions are proved above
for arbitrary actual columns and arbitrary positive original
$Q$; no sampled Gram substitutes for those proofs.
% END ACTUAL_CORNER_OBSERVATION

% BEGIN EXACT_MARKED_SPECIALIZATION
\section{The exact two-leg specialization}
Original polynomial and marked points: The Clankers \cite{ES488}.


This note constructs two explicit target paths for the original
ES--Fable polynomial. Their concatenated specialization has
eight, then seven, then three finite points in the original
source chart. It is not the single arithmetic trajectory in
MO12 of the preceding sections: both paths and all their signs are
specified here. In particular its second-leg escape exponent
differs from that transverse arithmetic trajectory.

\section{The original map and its exact fibre coordinates}
Use the original source coordinates $(a,y,z,w)$. Its complete
factor-coordinate formulas and polynomial map are
\begin{align}
 b&=i+ay,&c&=-i+2ay+a^2z,\nonumber\\
 d&=-iy-a(iz+2y^2)-a^2yz,&e&=2z-7iy^2+aw,\nonumber\\
 f&=iw+3iy^3-4yz+a(6iy^2z+wy+4y^4)+2a^2y^3z,&&\nonumber\\
 P(a,y,z,w)&=(ac,ae+bd,af+be,bf).
 &&\tag{VA1}
\end{align}
Here $L=aU+bV$, $C=cU^3+dU^2V+eUV^2+fV^3$ satisfy
$ad+bc=1$ and $C(-b,a)=1$, directly by substitution. Thus
their product at target $u=(A,B,C_1,D)$ is
$H_u=AU^4+U^3V+BU^2V^2+C_1UV^3+DV^4$.
The symbol $C_1$ is the third target coefficient, distinct
from the cubic factor denoted $C$. No source coordinate or
polynomial coefficient has been changed in VA1.

For a simple finite root $r$ of
$h_u(x)=Ax^4+x^3+Bx^2+C_1x+D$, the two inverse states are
\begin{align}
 a^2&=1/h_u'(r),&b&=-ar,&c&=A/a,&d&=(1+Ar)/a,\nonumber\\
 e&=(B+r+Ar^2)/a,& f&=(C_1+Br+r^2+Ar^3)/a,&&\nonumber\\
 y&=-r-i/a,&z&=A/a^3+2r/a+3i/a^2,&&\nonumber\\
 w&=7ir^2/a+(B-17r+Ar^2)/a^2-13i/a^3-2A/a^4.
 &&\tag{VA2}
\end{align}
Both nonzero square roots for $a$ are retained. To prove
completeness, divide $H_u$ by $a(U-rV)$ and compare successive
coefficients to obtain the six factor entries in VA2.
The resultant is $C(-b,a)=a^2h_u'(r)$, which must be one;
thus a multiple root gives no state. The original chart inverse
$y=(b-i)/a$, $z=(c+i-2ay)/a^2$,
$w=(e-2z+7iy^2)/a$ gives the remaining entries of VA2.
All operations are reversible for $a\ne0$.
At $A=0$ the root at infinity is simple, because the original
cubic coefficient is one. Its positive-chart state is exactly
\[
 q_\infty(A=0,B,C_1,D)
 =\left(0,B,\frac{i(7B^2-C_1)}2,11B^3-2BC_1-D\right).
\]
Indeed VA1 at $a=0$ gives
$P=(0,y,2iz+7y^2,-4iyz-w-3y^3)$, which solves to this formula.
The opposite factor sign has $b=-i$ and is outside the original
chart, since $a=0$ forces $b=i$ in VA1. Therefore every
simple finite root contributes two source states and infinity
contributes one in this chart.

\section{First leg: eight states specialize to the original seven}
For $0<\epsilon<1/2$ put $m_\epsilon=1-\epsilon^2/2$ and
$d_\epsilon=\sqrt{4-9\epsilon^2/4}>0$. The first target path,
its factorization and its four distinct marked roots are
\begin{align}
 u^{(1)}_\epsilon
   &=(-i\epsilon+i\epsilon^3/2,3i\epsilon/2,-1,0)
                \longrightarrow h_0=(0,0,-1,0),\nonumber\\
 h^{(1)}_\epsilon(x)
   &=(\epsilon x+i)x
             \left(-im_\epsilon x^2+\frac\epsilon2 x+i\right),\nonumber\\
 r_{\mathrm M}&=-i/\epsilon,&r_{\mathrm N}&=0,&
 r_{\mathrm T}&=\frac{d_\epsilon-i\epsilon/2}{2m_\epsilon},&
 r_{\mathrm E}&=\frac{-d_\epsilon-i\epsilon/2}{2m_\epsilon},\nonumber\\
 \operatorname{Disc} H_{u^{(1)}_\epsilon}&=4-9\epsilon^2/4>0.
 &&\tag{VA3}
\end{align}
The factorization follows by multiplying its displayed factors.
The quadratic formula gives the last two roots. The full
quartic discriminant is obtained from the resultant of this
polynomial and its derivative; substitution and expansion
gives exactly the last line of VA3. Its leading coefficient
$-i\epsilon m_\epsilon$ is nonzero on the stated interval.
Consequently all four roots are finite and distinct and VA2
gives exactly eight source states, with no root collision
anywhere before the endpoint. Their full inverse coordinates
are VA2 evaluated at the four explicit roots VA3 and both
values of $a^2=1/(h^{(1)}_\epsilon)'(r)$.

For the M root these two original-coordinate states simplify
completely to
\[
 q_{\mathrm M,+}(\epsilon)=(\epsilon,0,i/2,0),\qquad
 q_{\mathrm M,-}(\epsilon)
 =(-\epsilon,2i/\epsilon,6i/\epsilon^2-i/2,40i/\epsilon^3).
\]
Substitution in VA1 proves both have image VA3. Equivalently
their selected derivatives are $\epsilon^{-2}$ and their
factor scalings are $a=\pm\epsilon$. Their factor states
are negatives of each other. The positive one tends to
$(0,i,-i,0,i,0)$; the negative one tends to
$(0,-i,i,0,-i,0)$, a finite point in the factor completion
excluded from the original chart. Only its chart coordinates
escape. For the three remaining roots, VA3 gives limits
$0,1,-1$ and nonzero derivative limits $-1,2,2$.
Every expression in VA2 then has a finite limit on each
retained square-root branch. Thus exactly one of the eight
source states escapes, and seven remain finite at $h_0$.

\section{The complete signed marking at the junction}
The four original displayed collision points are
\[
\begin{aligned}
 q_{\mathrm M}&=(0,0,i/2,0),&q_{\mathrm N}&=(i,-1,-3i,13),\\
 q_{\mathrm T}&=(1/\sqrt2,-1-\sqrt2i,2\sqrt2+6i,-34-19\sqrt2i),\\
 q_{\mathrm E}&=(1/\sqrt2,1-\sqrt2i,-2\sqrt2+6i,34-19\sqrt2i).
\end{aligned}
\]
Define the original linear involution
$J(a,y,z,w)=(-a,-y,z,-w)$. The exact marked factor-root
associations and factor states at $h_0$ are
\begin{equation}
\begin{array}{c|c|c|c}
\text{label}&\text{root}&\text{positive factor state }(a,b,c,d,e,f)
                         &\text{two source-chart states}\\\hline
\mathrm M&\infty&(0,i,-i,0,i,0)&q_{\mathrm M},\ \text{opposite sign omitted}\\
\mathrm N&0&(i,0,0,-i,0,i)&q_{\mathrm N},\ Jq_{\mathrm N}\\
\mathrm T&1&(1/\sqrt2,-1/\sqrt2,0,\sqrt2,\sqrt2,0)
                       &q_{\mathrm T},\ Jq_{\mathrm E}\\
\mathrm E&-1&(1/\sqrt2,1/\sqrt2,0,\sqrt2,-\sqrt2,0)
                       &q_{\mathrm E},\ Jq_{\mathrm T}
\end{array}\tag{VA4}
\end{equation}
The negative factor state in each row is the negative of
every one of its six displayed entries. Multiplying $L$ and
$C$ in every row gives $UV(U-V)(U+V)$, with resultant one.
Substituting the original source states in VA1 recovers these
rows and their negatives. This proves all associations in
VA4. In particular $J$ interchanges the T/E factor-root
marking; $Jq_{\mathrm T}$ is not the T negative-sign state.
The missing M negative factor state is the one displayed
at the end of the preceding section. Thus all eight signed
factor states are described, and exactly seven of them
belong to the original source chart.

\section{Second leg: seven states specialize to three finite states}
Take $0\le s<1$ and put $\tau=1-s\in(0,1]$. The second
target path is
$u^{(2)}(s)=(0,s,s-1,0)=(0,1-\tau,-\tau,0)$.
It starts at $h_0$ when $s=0$ and tends to
$u_*=(0,1,0,0)$ when $s\uparrow1$. Its complete polynomial
and all seven inverse states are
\begin{align}
 h^{(2)}_\tau(x)&=x(x-\tau)(x+1),&
 \operatorname{Disc} H_{u^{(2)}}&=\tau^2(1+\tau)^2>0,\nonumber\\
 q_{\mathrm M}(\tau)&=
 \left(0,1-\tau,\frac{i(7(1-\tau)^2+\tau)}2,
                      11(1-\tau)^3+2\tau(1-\tau)\right),\nonumber\\
 q_{\mathrm N,\sigma}(\tau)&=
 \left(\frac{\sigma i}{\sqrt\tau},-\sigma\sqrt\tau,-3i\tau,
                     -\tau+\tau^2+13\sigma\tau^{3/2}\right),\nonumber\\
 q_{\mathrm T,\sigma}(\tau)&=
 \left(\frac{\sigma}{\sqrt{\tau(1+\tau)}},
       -\tau-i\sigma\sqrt{\tau(1+\tau)},\right.\nonumber\\
 &\hspace{12mm}2\sigma\tau^{3/2}\sqrt{1+\tau}+3i\tau(1+\tau),
 \left.\tau(1+\tau)(1-18\tau)
      -i\sigma\tau^{3/2}\sqrt{1+\tau}(13+6\tau)\right),\nonumber\\
 q_{\mathrm E,\sigma}(\tau)&=
 \left(\frac{\sigma}{\sqrt{1+\tau}},1-i\sigma\sqrt{1+\tau},
       -2\sigma\sqrt{1+\tau}+3i(1+\tau),\right.\nonumber\\
 &\hspace{12mm}\left.(18-\tau)(1+\tau)
               -i\sigma\sqrt{1+\tau}(6+13\tau)\right),
 \qquad\sigma\in\{1,-1\}.\tag{VA5}
\end{align}
All real square roots in VA5 are the positive ones; $\sigma$
retains both actual signs. The finite roots are precisely
$0,\tau,-1$ and their derivatives are $-\tau,\tau(1+\tau),1+\tau$.
They are pairwise distinct for every $0<\tau\le1$, and
their squared pair differences give the discriminant in
VA5. Substitution of these roots and both scalings into
VA2 gives every displayed finite-root formula. The infinity
formula after VA2 gives $q_{\mathrm M}(\tau)$.
Thus VA5 is the complete fibre and has no other root
coalescence before its endpoint.

At $\tau=1$ these formulas agree exactly with VA4:
the N positive and negative states are $q_{\mathrm N}$ and
$Jq_{\mathrm N}$; the T states are $q_{\mathrm T}$ and
$Jq_{\mathrm E}$; the E states are $q_{\mathrm E}$ and
$Jq_{\mathrm T}$. At $\tau=0$ the N and T roots collide.
Their four original source states escape, since their
$a$ coordinates diverge. The three remaining finite limits are
\[
 (0,1,7i/2,11),\qquad
 (\sigma,1-i\sigma,-2\sigma+3i,18-6i\sigma),
 \qquad\sigma\in\{1,-1\}.
\]
The complete fibre rule proved after VA2 gives exactly these
three states at the endpoint: the double root zero gives
none, the simple root $-1$ gives two, and infinity gives
the single retained chart state. The N and T factor
coefficients $a$ themselves diverge, so these four escapes
persist in the two-chart factor completion. The opposite
infinity sign was already outside the original chart all
along this second leg and is not a further escape on it.

\section{Exact fixed original-Gram transfer of the two legs}
Let $F:\C^4\to V$ be any one of the already constructed
fixed injective original receiving maps, with its unchanged
physical Hermitian inner product. Put $G=F^*H_VF>0$.
For example $F$ may be the fixed original conductor map
$\Psi_*$, or the actual injective marked observer $F_B$
on its proved nonzero-corner domain. Its full Gram, including
off-diagonal terms and original mass, is retained below.
Write $e_a=(1,0,0,0)^{\mathsf T}$,
$e_y=(0,1,0,0)^{\mathsf T}$,
$e_z=(0,0,1,0)^{\mathsf T}$ and
$e_w=(0,0,0,1)^{\mathsf T}$. The exact limiting vectors
and norm constants are
\begin{align}
 \epsilon^3q_{\mathrm M,-}(\epsilon)&\longrightarrow40ie_w,&
 \epsilon^3\|Fq_{\mathrm M,-}(\epsilon)\|
                       &\longrightarrow40\sqrt{G_{44}},\nonumber\\
 \sqrt\tau\,q_{\mathrm N,\sigma}(\tau)&\longrightarrow\sigma i e_a,&
 \sqrt\tau\,q_{\mathrm T,\sigma}(\tau)&\longrightarrow\sigma e_a,\nonumber\\
 \sqrt\tau\,\|Fq_{\mathrm N,\sigma}(\tau)\|
    &\longrightarrow\sqrt{G_{11}},&
 \sqrt\tau\,\|Fq_{\mathrm T,\sigma}(\tau)\|
    &\longrightarrow\sqrt{G_{11}},\nonumber\\
 \|F(u^{(2)}(s)-u_*)\|
   &=\tau\sqrt{(e_y+e_z)^*G(e_y+e_z)},&&\nonumber\\
 \|F(u^{(2)}(s)-u_*)\|^{1/2}\|Fq_{l,\sigma}(\tau)\|
   &\longrightarrow
     \bigl((e_y+e_z)^*G(e_y+e_z)\bigr)^{1/4}\sqrt{G_{11}},
 &&l\in\{\mathrm N,\mathrm T\}.
 \tag{VA6}
\end{align}
The vector limits are immediate from the complete coordinates
above. Applying the fixed linear $F$ and continuity of its
original norm proves each limit, with positive constants
because $F$ is injective. The target-distance equality is
exact at every $\tau$, since
$u^{(2)}(s)-u_*=-\tau(e_y+e_z)$.
In particular its coefficient is
$G_{22}+G_{33}+2\Re G_{23}$ under the square root;
the off-diagonal metric term has not been deleted.

These statements prove an exact two-leg specialization
$8\to7\to3$, including its missing factor sign and its
four original root-collision escapes. They do not claim a
single flow trajectory traverses those two paths. The
literal arithmetic orbit through $h_0$ has seven entire
finite trajectories, as proved in MO11 of the preceding sections.
The separate transverse arithmetic orbit MO12--19 has
four quarter-power escapes, one cubic chart escape and
three finite limits at $u_*$. The second path VA5 is
tangent to the discriminant: its discriminant vanishes
quadratically in $\tau$ and its four escape rates are
$\tau^{-1/2}$, exactly as VA6 proves. These are compatible
specializations of the same original polynomial, with
their different paths and rates retained explicitly.

The exact certificate
\texttt{verify\_exact\_marked\_specialization.py} verifies
the two first-leg M states, both discriminants, and all
seven original-coordinate second-leg formulas. For the
latter it uses $t=\sqrt\tau$ and $b=\sqrt{1+\tau}$ and
reduces every polynomial numerator modulo $b^2-t^2-1$
over $\mathbb Q(i)$; no numerical sample replaces an
identity. Its captured output is
\texttt{verify\_exact\_marked\_specialization\_output.txt}.
% END EXACT_MARKED_SPECIALIZATION

% BEGIN RATIONAL_NATIVE_RECEIVER
\section{The rational test receiver with the full original Gamma metric}
The received continuation \cite[Equations (33)--(36)]{ResidueIntake}
proposes a rational-function receiver of the signed object. Here its
entire coefficient matrix, inverse, original metric, graph kernel and
spectral comparison are calculated. No positivity claim from that
intake is used in this calculation. Its full-Weil form below is the
actual zero sum, including all multiplicities.

Retain the original $a_{\rm amp}$, $s\in\{1,a_{\rm amp}\}$,
$v$, $c_*=a_{\rm amp}/2-4$, $\kappa=c_*-1/2$,
the original measure $dm_s$ in FC28, and the exact maps
$\Phi,\Psi$ in FC32. For the test pole set
\[
 \omega=\sigma+i\tau,\quad\sigma>1,\quad
 h=2\sigma-1>0,\quad b_\omega=\bar\omega-1,\quad
 d=\sigma-\tfrac12=h/2.
 \tag{RW1}
\]
The positive square root $\sqrt h$ is retained. The symbol $s$ in
$dm_s$ is the original Gamma parameter; the complex variable of
the rational functions will be $z$ throughout this section.
Define, for $0\le j\le n$,
\[
 N_j^{(n)}(z)=\sqrt h\,(z+b_\omega)^j(\omega-z)^{n-j},\qquad
 \varphi_j^\omega(z)=\frac{N_j^{(n)}(z)}{(\omega-z)^{n+1}}
 =\frac{\sqrt h}{\omega-z}
           \left(\frac{z+b_\omega}{\omega-z}\right)^j.
 \tag{RW2}
\]
In particular this defines the same rational function for every
$n\ge j$; it does not change the denominator convention of a
fixed finite numerator space.

Let $R_\omega^{(n)}$ take a column $d_0,\ldots,d_n$ to the
coefficients of $\sum_jd_jN_j^{(n)}$ in $(1,z,\ldots,z^n)$.
Every entry and every inverse entry is the following finite sum:
\[
 \begin{aligned}
 (R_\omega^{(n)})_{\ell j}
  &=\sqrt h\sum_{r=0}^j
    \binom jr b_\omega^{j-r}
    \binom{n-j}{\ell-r}\omega^{n-j-\ell+r}(-1)^{\ell-r},\\
 ((R_\omega^{(n)})^{-1})_{j\ell}
  &=h^{-n-1/2}\sum_{r=0}^{\ell}
     \binom\ell r\omega^r(-b_\omega)^{\ell-r}
                         \binom{n-\ell}{j-r}.
 \end{aligned}\tag{RW3}
\]
In the first line only indices with $0\le\ell-r\le n-j$
are included, so there is no negative power hidden behind a
vanishing binomial coefficient. Binomial coefficients outside
their nonnegative range in the second line are zero.
Expansion of RW2 proves the first formula. For the inverse put
$u=(z+b_\omega)/(\omega-z)$, so
$z=(\omega u-b_\omega)/(1+u)$ and
$\omega-z=h/(1+u)$. For an arbitrary numerator $p$, the polynomial
\[
 h^{-n-1/2}(1+u)^n
                   p\left(\frac{\omega u-b_\omega}{1+u}\right)
\]
has coefficient column $d$. Expanding each monomial of $p$
gives exactly RW3. Substitution back proves both inverse identities.

The determinant, with the original coefficient order, is
\[
 \det R_\omega^{(n)}=h^{(n+1)^2/2},\qquad
 R_\omega^{(3)}=\sqrt h
 \begin{pmatrix}
 \omega^3&b_\omega\omega^2&b_\omega^2\omega&b_\omega^3\\
 -3\omega^2&\omega^2-2b_\omega\omega&2b_\omega\omega-b_\omega^2&3b_\omega^2\\
 3\omega&b_\omega-2\omega&\omega-2b_\omega&3b_\omega\\
 -1&1&-1&1
 \end{pmatrix}.
 \tag{RW4}
\]
For a direct determinant proof use $t=\omega-z$. The $j$th
polynomial is $\sqrt h(h-t)^jt^{n-j}$. Reversing the column
order makes its coefficient matrix triangular with diagonal
$\sqrt h h^j$. Reversal contributes
$(-1)^{n(n+1)/2}$, and changing $t$ back to $z$ contributes the
same sign, because the substitution matrix has diagonal $(-1)^\ell$.
Their product is positive and the remaining power is
$(n+1)/2+n(n+1)/2=(n+1)^2/2$. This proves RW4 without dropping
the square-root factor or changing a coefficient basis.

\section{Exact conductor return and the native Toeplitz matrix}
On $\mathcal P_n$ let
\[
 (C_\kappa)_{r\ell}=\binom\ell r\kappa^{\ell-r}\quad(r\le\ell),
 \qquad C_\kappa^{-1}=C_{-\kappa},\quad\det C_\kappa=1.
 \tag{RW5}
\]
This is precisely $p(S')\mapsto p(z+\kappa)$, by the binomial
formula; composition of translations proves the inverse.
For the four-dimensional receiver take $n=3$ and write $R=R_\omega^{(3)}$.
With the original monomial realization of $\Psi$, define
\[
 B_\omega=R^{-1}C_\kappa\Psi,\quad
 B_\omega^{-1}=\Psi^{-1}C_{-\kappa}R,\qquad
 B_{\omega,U}=R^{-1}C_\kappa T_A\Phi=B_\omega.
 \tag{RW6}
\]
Thus an original target state $p=\Psi q$ has the exact rational
image $p(z+\kappa)/(\omega-z)^4=\sum_j(B_\omega q)_j\varphi_j^\omega(z)$.
Equation RW6 is the original conductor commutation, including
all the moments in $\Psi$ and $T_A$; it is not an identification
of $P$ with $T_A$.

Let $\mathcal H$ be the original target monomial Gram in FC56a,
and $G=\Psi^*\mathcal H\Psi$. In rational coefficient coordinates
the original physical target norm has the exact positive Gram
\[
 H_\omega=B_\omega^{-*}GB_\omega^{-1}
       =R^*C_{-\kappa}^*\mathcal HC_{-\kappa}R.
 \tag{RW7}
\]
To evaluate it, put $S'=c_*+iy$ in the unchanged original
measure. Since $S'-\kappa=1/2+iy$, the actual numerator
polynomial is $N_j^{(3)}(1/2+iy)$. Its two factors are
$d+i(y-\tau)$ and $d-i(y-\tau)$, which are conjugate.
Consequently all sixteen entries are
\[
 (H_\omega)_{ij}=g_{j-i},\qquad
 g_r=h\int_{\R}(d+i(y-\tau))^{3+r}
                    (d-i(y-\tau))^{3-r}\,dm_s(y),
 \quad -3\le r\le3,\quad g_{-r}=\bar g_r.
 \tag{RW8}
\]
Thus the original metric itself is Toeplitz in this basis.
This is a proved identity, not a unitary-metric substitution.
The original center moved through the stated translation;
the measure, its total mass and every correlation remain.

Here is a complete finite evaluation of RW8. Write
$\mu_j(\tau)=\int(y-\tau)^jdm_s(y)$ and retain
\[
 \mu_j(\tau)=\sum_{\ell=0}^j\binom j\ell(-\tau)^{j-\ell}m_\ell,
 \quad m_0=M_s=(2\pi)^{s/2},\quad m_2=M_s\beta,
 \quad m_4=M_s(3\beta^2+2\beta),\quad
 m_6=M_s(15\beta^3+30\beta^2+16\beta),\quad\beta=s/2,
 \tag{RW9}
\]
where all odd $m_\ell$ vanish. These are exactly the proved
original moments of FC56a, with their cited DLMF provenance
\cite{WCF}. In terms of these full shifted moments,
\[
 g_r=h\sum_{p=0}^{3+r}\sum_{q=0}^{3-r}
 \binom{3+r}p\binom{3-r}q d^{6-p-q}i^p(-i)^q\mu_{p+q}(\tau).
 \tag{RW10}
\]
Every sum is finite and includes all odd shifted moments.
It is obtained by expanding the two factors in RW8, so it
specifies every entry without a quadrature or asymptotic error.
For example $g_0=h(d^6m_0+3d^4\mu_2+3d^2\mu_4+\mu_6)$.

The full determinant is also explicit:
\[
 \det H_\omega
 =h^{16}\prod_{j=0}^3\bigl[M_s j!(s/2)_j\bigr]
 =12h^{16}M_s^4\beta^3(\beta+1)^2(\beta+2)>0.
 \tag{RW11}
\]
Indeed RW7, RW4 and $\det C_{-\kappa}=1$ give
$\det H_\omega=h^{16}\det\mathcal H$; the exact monomial
Gram determinant is the product of the original monic
orthogonal squared norms in FC28--29. This proves independence
of $\tau$ for the determinant, not bounded conditioning as
$|\tau|$ grows. Each entry in RW8 still retains its full height
dependence.

\section{The same exact comparison for the full eight-state object}
Define the actual full-Weil form of the rational functions by
\[
 (T_3(\omega))_{ij}
 =\sum_{\rho}m_\rho\overline{\varphi_i^\omega(1-\bar\rho)}
                                     \varphi_j^\omega(\rho).
 \tag{RW12}
\]
Every sum includes the complete nontrivial zero multiset.
For fixed $\omega,i,j$ each summand is $O(|\rho|^{-2})$:
both pole factors have size comparable to $|\Im\rho|$, and
the finite powers of their ratio are bounded for large height.
For completeness the convergence needed here also follows
directly from the classical theta integral used in the residue
calculation below. Its Gaussian bound gives, for $|z|\le R$,
$|\xi(z)|\le\exp(CR\log(R+2))$ with a fixed constant $C$:
bound the power of $x\ge1$ by $x^{\lceil(R+1)/2\rceil}$,
integrate against $e^{-\pi x}$ and use $k!\le k^k$.
Jensen's formula centered at $\xi(0)=1/2$, on a circle between
radii $2R$ and $3R$ avoiding zeros, then gives
$N(R)\log2\le C'(R+1)\log(R+2)$ for the zero count with
multiplicity in $|z|\le R$. Splitting into annuli
$2^j\le|\rho|<2^{j+1}$ proves
$\sum_\rho m_\rho|\rho|^{-2}<\infty$, because
$\sum_j(j+1)2^{-j}<\infty$. This is the same absolute
convergence underlying the classical canonical product
\cite[25.2.12]{DLMFXi}. Reflection of the zero multiset
proves it Hermitian by changing $\rho$ to $1-\bar\rho$.
The exact identity
$\overline{U_\omega(1-\bar\rho)}=U_\omega(\rho)^{-1}$,
where $U_\omega(z)=(z+b_\omega)/(\omega-z)$, also proves
\[
 (T_3(\omega))_{ij}=c_{j-i},\qquad
 c_r=h\sum_\rho m_\rho
       \frac{(\rho+b_\omega)^{r-1}}{(\omega-\rho)^{r+1}},
 \quad c_{-r}=\bar c_r.
 \tag{RW13}
\]
Neither denominator vanishes in the open critical strip, since
$\sigma>1$. This verifies the exact form needed for the
receiving calculation; finite arithmetic evaluation and tail
bounds of $c_r$ belong to the separately retained full-Weil
calculation, not to a claim that these sums have been numerically
evaluated here.

Retain both original label and evaluation coordinates of SE/TC.
Let $E,O$ be their full mean and difference matrices, and let
$G_Y=Y^*\langle\ ,\ \rangle_{\Gamma,W}Y>0$.
The complete original target direct-sum Gram and Weil form are
\[
 \mathbb H=
 \begin{pmatrix}I&0\\E&O\end{pmatrix}^{\!*}
 \begin{pmatrix}G_Y&0\\0&G\end{pmatrix}
 \begin{pmatrix}I&0\\E&O\end{pmatrix},\qquad
 \mathbb W=[E\ O]^*B_\omega^*T_3(\omega)B_\omega[E\ O].
 \tag{RW14}
\]
The original $O$ is invertible by SE5 on its stated quartet
domain. The exact simultaneous change and inverse are
\[
 S_\omega=\begin{pmatrix}I&0\\-O^{-1}E&O^{-1}B_\omega^{-1}\end{pmatrix},
 \qquad
 S_\omega^{-1}=\begin{pmatrix}I&0\\B_\omega E&B_\omega O\end{pmatrix}.
 \tag{RW15}
\]
Multiplying either way proves the inverse; multiplying by the
two actual receiving matrices proves simultaneously
\[
 S_\omega^*\mathbb W S_\omega=\diag(0_4,T_3(\omega)),\qquad
 S_\omega^*\mathbb H S_\omega=\diag(G_Y,H_\omega).
 \tag{RW16}
\]
Thus no positive physical norm has been discarded to obtain
the Weil congruence. In the full original native metric the
generalized eigenvalues are four zeros together with those
of $H_\omega^{-1/2}T_3(\omega)H_\omega^{-1/2}$.
To prove this, apply the positive square root of each diagonal
block in RW16 to the equation $\mathbb Wv=\lambda\mathbb Hv$.
Both transformations are invertible; the resulting equation is
the direct sum of a zero Hermitian matrix and the displayed
Hermitian matrix. If $T_3$ itself has a kernel, the total
nullity is $4+\dim\ker T_3$, and the complete original kernel is
\[
 \left\{\left(\alpha,-O^{-1}E\alpha+
          O^{-1}B_\omega^{-1}d\right):\ d\in\ker T_3\right\}.
 \tag{RW17}
\]
In particular the four-dimensional evaluation graph is retained
exactly, without claiming that it exhausts every possible Weil
kernel. Positive and negative inertia are those of the actual
four-dimensional $T_3$. This statement is an exact consequence
of two constructed forms, not a supposition about their signs.

For every $d$ the canonical lift with zero retained-label
coordinate is the explicit original state
$(0,O^{-1}B_\omega^{-1}d)$. Its physical target polynomial is
$C_{-\kappa}Rd$, and its original conductor source is
\[
 T_A^{-1}C_{-\kappa}Rd,
 \tag{RW18}
\]
where the full original inverse FC24 supplies every coefficient
and its source norm. Hence any evaluated test vector is carried
back through the actual conductor, rather than only through a
formal Hermitian congruence. For degrees $n>3$, RW2--5 use the
actual $\mathcal P_n$ conductor restriction in FC21 and FC24;
the original eight-state space itself stays eight-dimensional.

The actual original source direct-sum norm is retained as well.
Let $G_X=X^*\langle\ ,\ \rangle_{\Gamma,U}X$ and
$G_U=\Phi^*\langle\ ,\ \rangle_{\Gamma,U}\Phi$, where the
bracket is the original quotient norm. Its full Gram and
transported rational coefficient Gram are
\[
 \mathbb H_U=
 \begin{pmatrix}I&0\\E&O\end{pmatrix}^{\!*}
 \begin{pmatrix}G_X&0\\0&G_U\end{pmatrix}
 \begin{pmatrix}I&0\\E&O\end{pmatrix},\qquad
 H_{U,\omega}=B_\omega^{-*}G_UB_\omega^{-1}.
 \tag{RW19}
\]
The source state represented by $(\alpha,\beta)$ is
$(X\alpha,\Phi(E\alpha+O\beta))$. Applying $T_A$ on
each component gives exactly the target state
$(Y\alpha,\Psi(E\alpha+O\beta))$, because
$T_AX=Y$ and $T_A\Phi=\Psi$.
Consequently the same explicit matrix RW15 proves
\[
 S_\omega^*\mathbb H_U S_\omega=\diag(G_X,H_{U,\omega}),
 \qquad
 S_\omega^*\mathbb W S_\omega=\diag(0_4,T_3).
 \tag{RW20}
\]
No mixed source--target direct-sum metric is being substituted.
The actual source generalized eigenvalues are four zeros and
those of $H_{U,\omega}^{-1/2}T_3H_{U,\omega}^{-1/2}$, by
the same invertible square-root argument used for RW16.
Its kernel is still exactly RW17. The complete entries and
all source height constants of $H_{U,\omega}$ are evaluated
in NH26 onward using the original FC24 inverse and quotient
Gram. Thus the two original sides of the conductor have
separate fully specified metrics and the same exact Weil receiver.
% END RATIONAL_NATIVE_RECEIVER

% BEGIN RATIONAL_NATIVE_HEIGHT
\section{Exact target and source height spectra}

This calculation evaluates the entire high-height spectrum of the
original native Gram in RW7--11 of
\texttt{RATIONAL\_NATIVE\_RECEIVER\_BODY.tex}. Its polynomial
basis, original Gamma measure, center translation, phases, and total
mass remain exactly those of that receiver. The original Gamma
moments are those proved in FC56a and WCF \cite{WCF}, with the
human source of the orthogonal-polynomial dictionary recorded there.
No assertion about the sign of the separate full-Weil form is used.

\section{The exact original coefficient map and Gram}
Fix $\sigma>1$, let $\omega=\sigma+i\tau$ with $\tau\in\R$,
and retain
\[
 h=2\sigma-1,\qquad d=\sigma-\tfrac12=h/2,\qquad
 b=s/2>0,\qquad M=(2\pi)^{s/2}.
 \tag{NH1}
\]
Here $s$ is the fixed original Gamma parameter. The original measure is
\[
 dm_s(y)=\frac{M2^{s/2}}{4\pi\Gamma(s/2)}
               |\Gamma(s/4+iy/2)|^2dy.
 \tag{NH2}
\]
It has the original moments
\[
 m_0=M,\quad m_2=Mb,\quad m_4=M(3b^2+2b),\quad
 m_6=M(15b^3+30b^2+16b),\qquad m_1=m_3=m_5=0.
 \tag{NH3}
\]
The exact polynomials in the rational receiver are
\[
 N_j(z)=\sqrt h\,(z+\bar\omega-1)^j(\omega-z)^{3-j},
 \qquad 0\le j\le3.
 \tag{NH4}
\]
For the original physical target center $c_*$ put
$\kappa=c_*-1/2$. Define the actual numerator map
\[
 \mathcal M_\omega:\C^4\longrightarrow\mathcal P_3,
 \qquad \mathcal M_\omega c(S')=\sum_{j=0}^3c_jN_j(S'-\kappa).
 \tag{NH5}
\]
The target has its unchanged norm
$\|p\|_\Gamma^2=\int|p(c_*+iy)|^2dm_s(y)$.
In NH5, the domain is the displayed coefficient space with its
ordinary Euclidean norm when singular values are computed below.
This declaration identifies the operator whose singular values are
being evaluated; it changes no original target norm.

At $S'=c_*+iy$, the polynomial in NH4 is
$\sqrt h(d+i(y-\tau))^j(d-i(y-\tau))^{3-j}$.
Let $Q(\tau)$ be its coefficient matrix in $(1,y,y^2,y^3)$ after
factoring out the displayed constant $\sqrt h$ from the formula.
Its complete entries, still retaining all phases, are
\[
\begin{pmatrix}
d^3+3id^2\tau-3d\tau^2-i\tau^3&d^3+id^2\tau+d\tau^2+i\tau^3&
d^3-id^2\tau+d\tau^2-i\tau^3&d^3-3id^2\tau-3d\tau^2+i\tau^3\\
-3id^2+6d\tau+3i\tau^2&-id^2-2d\tau-3i\tau^2&
id^2-2d\tau+3i\tau^2&3id^2+6d\tau-3i\tau^2\\
-3d-3i\tau&d+3i\tau&d-3i\tau&-3d+3i\tau\\
i&-i&i&-i
\end{pmatrix}.
\tag{NH6}
\]
The Gram of the displayed $y$-monomials is exactly
\[
 K=M\begin{pmatrix}
 1&0&b&0\\0&b&0&3b^2+2b\\
 b&0&3b^2+2b&0\\0&3b^2+2b&0&15b^3+30b^2+16b
 \end{pmatrix}.
 \tag{NH7}
\]
It is positive definite: every nonzero polynomial has positive
integral of its squared modulus against the strictly positive
density NH2. Direct computation also gives the positive leading
principal determinants
\[
 M,\quad M^2b,\quad2M^3b^2(b+1),\quad
 12M^4b^3(b+1)^2(b+2).
 \tag{NH8}
\]
Consequently the original coefficient Gram is the exact matrix
\[
 H_\omega=hQ(\tau)^*KQ(\tau),\qquad
 (H_\omega)_{ij}=h\int(d+i(y-\tau))^{3+j-i}
                      (d-i(y-\tau))^{3-j+i}dm_s(y).
 \tag{NH9}
\]
The second identity follows by conjugating the $i$th polynomial
before multiplying by the $j$th. It proves the native Toeplitz
formula and keeps every shifted odd moment as well as the even ones.

The determinant of NH6 is $-64d^6$, for every $\tau$.
One proof is the translation factorization in NH11 below and
direct expansion at $\tau=0$. Hence
\[
 \det H_\omega
 =h^4(64d^6)^2\det K
 =12h^{16}M^4b^3(b+1)^2(b+2)>0.
 \tag{NH10}
\]
The coefficient map $\sqrt h Q$ in the $y$-basis has determinant
$-h^8$. The receiver matrix $R_\omega^{(3)}$ in the original
$z$-monomial basis has determinant $+h^8$: substitution
$z=1/2+iy$ has determinant $i^{0+1+2+3}=-1$. Thus NH10
agrees with RW11 with its original sign conventions, not only
with its determinant modulus.

\section{All compound leading matrices}
Use row and column indices $0,1,2,3$, preserving the ordered
coefficient basis. Define the full translation matrix
\[
 T(\tau)_{k\ell}=\begin{cases}
 \binom\ell k(-\tau)^{\ell-k},&k\le\ell,\\0,&k>\ell.
 \end{cases}
 \qquad Q(\tau)=T(\tau)Q(0).
 \tag{NH11}
\]
Indeed substituting $y-\tau$ into each degree-$\ell$ monomial
gives exactly the indicated coefficient column. This identity
proves every entry and keeps the actual shift rather than discarding
it in the measure.

For increasing $r$-subsets $I,J\subset\{0,1,2,3\}$, the
$(I,J)$ entry of $\wedge^rT$ is
\[
 (-\tau)^{\sum J-\sum I}
       \det\bigl(\tbinom j i\bigr)_{i\in I,j\in J},
 \tag{NH12}
\]
whenever the coefficient determinant is nonzero. This follows
by factoring the powers out of each row and column of the
minor. Its maximal possible degree is
$m_r=r(4-r)$, attained only at
$I_r=\{0,\ldots,r-1\}$ and $J_r=\{4-r,\ldots,3\}$.
The binomial minors at those positions are one: in degrees
one and four this is immediate, in degree two the matrix is
$\left(\begin{smallmatrix}1&1\\2&3\end{smallmatrix}\right)$,
and in degree three it is
$\left(\begin{smallmatrix}1&1&1\\1&2&3\\0&1&3\end{smallmatrix}\right)$.
Each determinant is one. All other nonzero minors have strictly
smaller degree, so there is no omitted term of a higher order.

Here are convenient unchanged input vectors:
\[
 v_1=(1,-1,1,-1)^T,\quad v_2=(3,-1,-1,3)^T,\quad
 v_3=(1,1,-1,-1)^T,\quad v_4=(1,3,3,1)^T.
 \tag{NH13}
\]
They have Gram $\diag(4,20,4,20)$ in coefficient Euclidean
coordinates. Multiplying the unique leading row of each compound
in NH12 by $\wedge^rQ(0)$ gives the following full limits:
\[
\begin{array}{c|c|c}
 r&m_r&\lim_{\tau\to\pm\infty}\tau^{-m_r}\wedge^rQ(\tau)\\\hline
 1&3&-i\,e_0v_1^{\mathsf T}\\
 2&4&id\,e_{01}(v_1\wedge v_2)^{\mathsf T}\\
 3&3&-2d^3\,e_{012}(v_1\wedge v_2\wedge v_3)^{\mathsf T}\\
 4&0&-64d^6
\end{array}.
\tag{NH14}
\]
For direct coefficient verification, the three nonzero row
vectors are respectively
\[
 (-i,i,-i,i),\qquad
 (2id,-4id,6id,2id,-4id,2id),\qquad
 (-8d^3,24d^3,-24d^3,8d^3).
\]
The compound subsets use increasing lexicographic order. These
rows follow by minors of $Q(0)$; all other rows vanish in the
limit by NH12. The limit with $\tau^{-m_r}$ retains the sign
for negative $\tau$; its squared Gram limits use $|\tau|^{2m_r}$
and therefore coincide on both ends.

\section{Full eigenvalue and native singular-value asymptotics}
Taking exterior powers of the exact Gram NH9 gives
\[
 \wedge^rH_\omega
 =h^r(\wedge^rQ)^*(\wedge^rK)(\wedge^rQ).
 \tag{NH15}
\]
The squared $K$-lengths of $e_0,e_{01},e_{012}$ are the first
three determinants NH8. Thus the full leading matrices in
NH15 are rank one for $r=1,2,3$. If $w_2=v_1\wedge v_2$
and $w_3=v_1\wedge v_2\wedge v_3$, they are
\[
\begin{aligned}
 |\tau|^{-6}H_\omega&\longrightarrow hM v_1v_1^{\mathsf T},\\
 |\tau|^{-8}\wedge^2H_\omega&\longrightarrow
              h^2d^2M^2b\,w_2w_2^{\mathsf T},\\
 |\tau|^{-6}\wedge^3H_\omega&\longrightarrow
              8h^3d^6M^3b^2(b+1)\,w_3w_3^{\mathsf T}.
\end{aligned}\tag{NH16}
\]
Their nonzero eigenvalues are their traces. Orthogonality in
NH13 gives $\|v_1\|^2=4$, $\|w_2\|^2=80$ and
$\|w_3\|^2=320$. Substituting the retained identity $d=h/2$
therefore yields all four largest compound-eigenvalue constants
\[
 C_1=4hM,\qquad C_2=20h^4M^2b,\qquad
 C_3=40h^9M^3b^2(b+1),\qquad
 C_4=12h^{16}M^4b^3(b+1)^2(b+2).
 \tag{NH17}
\]
Here $C_4$ is the exact determinant NH10. Every constant is
strictly positive, with the entire original mass $M$ retained.

Let $\lambda_1\ge\lambda_2\ge\lambda_3\ge\lambda_4>0$
be the eigenvalues of the displayed coefficient matrix $H_\omega$
with respect to coefficient Euclidean coordinates. A positive
matrix has an orthonormal eigenbasis, whose coordinate wedges
diagonalize its compounds with all products of eigenvalues.
Consequently $\|\wedge^rH_\omega\|=\lambda_1\cdots\lambda_r$.
Continuity of the norm in NH16 and successive ratios of the
constants NH17 prove the complete spectrum
\[
\boxed{\begin{aligned}
 \lambda_1&\sim4hM|\tau|^6,\\
 \lambda_2&\sim5h^3Mb|\tau|^2,\\
 \lambda_3&\sim2h^5Mb(b+1)|\tau|^{-2},\\
 \lambda_4&\sim\tfrac3{10}h^7Mb(b+1)(b+2)|\tau|^{-6}.
\end{aligned}}\tag{NH18}
\]
The product of the four constants is precisely NH10. This
also gives the exact coefficient condition-number asymptote
\[
 \frac{\lambda_1}{\lambda_4}
 \sim\frac{40}{3h^6b(b+1)(b+2)}|\tau|^{12}.
 \tag{NH19}
\]

Because the Gram of the typed numerator operator NH5 is
$H_\omega$, its singular values are $\sqrt{\lambda_j}$.
Thus they are, in decreasing order,
\[
\boxed{\begin{aligned}
 s_1(\mathcal M_\omega)&\sim2\sqrt{hM}|\tau|^3,\\
 s_2(\mathcal M_\omega)&\sim\sqrt{5h^3Mb}|\tau|,\\
 s_3(\mathcal M_\omega)&\sim\sqrt{2h^5Mb(b+1)}|\tau|^{-1},\\
 s_4(\mathcal M_\omega)&\sim\sqrt{\tfrac3{10}h^7Mb(b+1)(b+2)}|\tau|^{-3}.
\end{aligned}}\tag{NH20}
\]
Its inverse uses the original native space as domain and the
coefficient Euclidean space as codomain, and has the reciprocal
singular values in reverse order. All inverse exterior norms
are determined with constants by
\[
 \|\wedge^r\mathcal M_\omega^{-1}\|
 \sim\sqrt{\frac{C_{4-r}}{C_4}}\,
            |\tau|^{m_{4-r}},\qquad 0\le r\le4,
 \quad C_0=1, m_0=m_4=0.
 \tag{NH21}
\]
This follows either by complementary products of the four
singular values or by the exact identity
$\|\wedge^r\mathcal M_\omega^{-1}\|
=\|\wedge^{4-r}\mathcal M_\omega\|/\sqrt{\det H_\omega}$.

\section{Original input and output singular flags}
The limiting coefficient eigenlines of $H_\omega$ in descending
eigenvalue order are exactly
\[
 \C(v_1/2),\qquad\C(v_2/\sqrt{20}),\qquad
 \C(v_3/2),\qquad\C(v_4/\sqrt{20}).
 \tag{NH22}
\]
There is no assertion about individual phases of eigenvectors.
To prove the statement, the rank-one limits of NH16 show that
the leading one-, two-, and three-dimensional coefficient
eigenflags converge to
$\operatorname{span}(v_1)$, $\operatorname{span}(v_1,v_2)$ and
$\operatorname{span}(v_1,v_2,v_3)$. More explicitly the normalized compound
matrices converge to a rank-one orthogonal projection onto
$v_1$, $w_2$ or $w_3$. Their top eigenspace projections converge:
any unit eigenvector subsequential limit solves the limiting
eigenvector equation and lies in its unique one-dimensional
range. A nonzero decomposable wedge determines its plane.
The eigenvalue ratios in NH18 tend to infinity, so these are
the unique leading flags for sufficiently large $|\tau|$.
Taking orthogonal successive differences yields NH22, using
the exact orthogonality in NH13.

The native output singular flags of NH5 converge to the full
polynomial degree filtration $\mathcal P_0\subset\mathcal P_1
\subset\mathcal P_2\subset\mathcal P_3$ in its original
Gamma metric. Indeed each rank-one forward compound limit
NH14 has output range $e_0$, $e_{01}$ or $e_{012}$; the
same positive-matrix argument applied using the original
$K$ metric proves these output flags. The limiting individual
native output lines have the following original unit
representatives on $S'=c_*+iy$:
\[
 \frac1{\sqrt M},\qquad
 \frac{y}{\sqrt{Mb}},\qquad
 \frac{y^2-b}{\sqrt{2Mb(b+1)}},\qquad
 \frac{y^3-(3b+2)y}{\sqrt{6Mb(b+1)(b+2)}}.
 \tag{NH23}
\]
As polynomials in the original physical coordinate, each $y$
in NH23 means $(S'-c_*)/i$; the phases and center are thereby
specified. Orthogonality and the displayed squared norms
follow by substituting the exact moments NH3, so these
expressions are the original metric's unit representatives,
not a replacement metric. They retain its total mass $M$.

\section{Exact meaning for the original conductor receiver}
Retain the original map $\Psi:\C^4\to\mathcal P_3$,
its Gram $G=\Psi^*\langle\ ,\ \rangle_\Gamma\Psi$, and
the original receiver coefficient map
$B_\omega=R_\omega^{-1}C_\kappa\Psi$ from RW6. The exact
identity NH5 gives
\[
 \mathcal M_\omega B_\omega=\Psi,\qquad
 B_\omega^*H_\omega B_\omega=G.
 \tag{NH24}
\]
Thus $B_\omega$ from the original metric $G$ to the original
coefficient metric $H_\omega$ is an exact isometry. Its
singular values with those two specified native norms are
all one. The height-dependent eigenvalues NH18 compare the
actual coefficient Gram to the Euclidean coordinate metric;
the numerator singular values NH20 refer to the explicitly
typed map NH5. These identities specify their mathematical
meaning without dropping the original norm or inferring a
degeneration of the fixed linear conductor.

Finally let $L=T_A|_U$ be the original invertible four-plane
conductor from its original quotient space to the target.
The complete original source numerator receiver is
$\mathcal M_{U,\omega}=L^{-1}\mathcal M_\omega$. For every
$0\le r\le4$, composition gives
\[
 \frac{\|\wedge^r\mathcal M_\omega\|}{\|\wedge^rL\|}
 \le\|\wedge^r\mathcal M_{U,\omega}\|
 \le\|\wedge^rL^{-1}\|\|\wedge^r\mathcal M_\omega\|.
 \tag{NH25}
\]
The left inequality is the right one applied to
$\mathcal M_\omega=L\mathcal M_{U,\omega}$.
The original WCF inverse-exterior bound applies to the fixed
factor $\|\wedge^rL^{-1}\|$ here; all target numerator
constants have already been evaluated in NH17--21. No
substitution of the separate full-Weil form for this original
positive Gram is made.

The independent exact certificate
\texttt{derive\_rational\_native\_height.py} constructs NH6
directly from NH4, verifies the full factorization NH11,
computes every compound and its leading coefficient,
checks rank one and every squared native compound constant,
and verifies the Gram of the four polynomials NH23 using
NH3. Its complete machine-readable output is
\texttt{rational\_native\_height\_exact.json}.

\section{Every original source Gram entry and its evaluated spectrum}
The source side of NH25 can be evaluated completely. Retain
the actual conductor order $v$, its original moments $\mu_n$,
$\mu_v\ne0$, its centers $c_0=c_*+4$, $c_1=c_*$, and the
exact function $g$ of WCF5. No conductor coefficient varies
with $\tau$. To give the inverse coefficients explicitly, put
\[
 \nu_{v+j}=\sum_{r=0}^j\binom{v+j}{v+r}(-4)^{j-r}\mu_{v+r},
 \qquad 0\le j\le3.
 \tag{NH26}
\]
These are precisely the original centered moments
$\sum_{u,w}a_{uw}(\beta_{8,u,w}-4)^{v+j}$; the omitted
lower terms vanish by the actual order $v$. Expanding the
unchanged $w(t)=-i\arctan t$ gives every needed coefficient:
\[
\begin{aligned}
 g_0&=\frac{(-i)^v\mu_v}{v!},&
 g_1&=\frac{(-i)^{v+1}\nu_{v+1}}{(v+1)!},\\
 g_2&=\frac{(-i)^{v+2}\nu_{v+2}}{(v+2)!}
                           -\frac v3g_0,&
 g_3&=\frac{(-i)^{v+3}\nu_{v+3}}{(v+3)!}
                           -\frac{v+1}{3}g_1.
\end{aligned}\tag{NH27}
\]
Indeed $g(t)=t^{-v}\sum_{n\ge v}\nu_nw(t)^n/n!$,
and $w(t)^n=(-i)^nt^n(1-nt^2/3+O(t^4))$; only the
four displayed moments contribute through degree three.
Thus the four exact FC24 inverse coefficients are
\[
 \xi_0=g_0^{-1},\quad \xi_1=-g_1g_0^{-2},\quad
 \xi_2=g_1^2g_0^{-3}-g_2g_0^{-2},\quad
 \xi_3=-g_1^3g_0^{-4}+2g_1g_2g_0^{-3}-g_3g_0^{-2}.
 \tag{NH28}
\]
Multiplication by $g_0+g_1t+g_2t^2+g_3t^3$ gives one
through degree three, proving these are the original inverse
coefficients. In particular
$\xi_0=i^vv!/\mu_v$ and
$|\xi_0|^2=(v!)^2/|\mu_v|^2$; its original phase is retained.

Let $\phi_n(S)=p_{n,s}((S-c_0)/i)/\sqrt{Mn!(b)_n}$ be
the original orthonormal polynomials at the original source
center, where
\[
 p_{0,s}(y)=1,\quad p_{1,s}(y)=y,\qquad
 p_{n+1,s}(y)=yp_{n,s}(y)-n(n-1+b)p_{n-1,s}(y).
 \tag{NH29}
\]
Their classes $[\phi_v],\ldots,[\phi_{v+3}]$ form the
original quotient orthonormal frame. Each representative is
orthogonal to the full removed kernel $\mathcal P_{v-1}$;
hence every combination in this frame is its class's unique
minimum-norm representative. This follows directly from
Pythagoras upon adding any vector of $\mathcal P_{v-1}$.
No nonminimal monomial representative is used below.

Put $\eta_j(S')=((S'-c_*)/i)^j$ and let
$u_j=L^{-1}\eta_j$ in this original quotient, $0\le j\le3$.
Define the positive original weights
\[
 k_j=\rho_{v+j}^{-2}=\frac{(b)_{v+j}}{(v+j)!},
 \qquad 0\le j\le3.
 \tag{NH30}
\]
The complete coefficient matrix of the minimum representatives
$u_0,\ldots,u_3$ in the frame NH29 is
\[
 U=\sqrt M\diag(\sqrt{k_0},\sqrt{k_1},\sqrt{k_2},\sqrt{k_3})F,
 \qquad
 F=\begin{pmatrix}
 \xi_0&\xi_1&b\xi_0+2\xi_2&(3b+2)\xi_1+6\xi_3\\
 0&\xi_0&2\xi_1&(3b+2)\xi_0+6\xi_2\\
 0&0&2\xi_0&6\xi_1\\
 0&0&0&6\xi_0
 \end{pmatrix}.
 \tag{NH31}
\]
This is an evaluated matrix of actual inverse polynomials:
$u_j=\sum_{r=0}^jU_{rj}[\phi_{v+r}]$ with $\phi$ given
by NH29 and every $\xi$ given by NH26--28.

To prove NH31 directly from FC24, the original target
orthonormal frame expands the four monomials by
\[
 \begin{pmatrix}
 \sqrt M&0&b\sqrt M&0\\
 0&\sqrt{Mb}&0&(3b+2)\sqrt{Mb}\\
 0&0&\sqrt{2Mb(b+1)}&0\\
 0&0&0&\sqrt{6Mb(b+1)(b+2)}
 \end{pmatrix}.
 \tag{NH32}
\]
These are exactly $1$, $y$, $y^2=p_2+b$ and
$y^3=p_3+(3b+2)y$, with their original squared norms.
FC24 has entry $\xi_{j-r}\rho_j/\rho_{v+r}$ for $r\le j$.
The identity
$\rho_j\sqrt{Mj!(b)_j}=\sqrt M\,j!$
therefore proves NH31 by multiplication, retaining all phases
in the inverse coefficients and every factorial.

For fully explicit entries put
$A_2=b\xi_0+2\xi_2$, $A_3=(3b+2)\xi_1+6\xi_3$,
$B_3=(3b+2)\xi_0+6\xi_2$. The original source Gram
$K_U=(\langle u_j,u_\ell\rangle)$ has these ten entries;
the remaining six are their Hermitian conjugates:
\[
\begin{aligned}
 (K_U)_{00}/M&=k_0|\xi_0|^2,&
 (K_U)_{01}/M&=k_0\bar\xi_0\xi_1,\\
 (K_U)_{02}/M&=k_0\bar\xi_0A_2,&
 (K_U)_{03}/M&=k_0\bar\xi_0A_3,\\
 (K_U)_{11}/M&=k_0|\xi_1|^2+k_1|\xi_0|^2,\\
 (K_U)_{12}/M&=k_0\bar\xi_1A_2+2k_1\bar\xi_0\xi_1,\\
 (K_U)_{13}/M&=k_0\bar\xi_1A_3+k_1\bar\xi_0B_3,\\
 (K_U)_{22}/M&=k_0|A_2|^2+4k_1|\xi_1|^2+4k_2|\xi_0|^2,\\
 (K_U)_{23}/M&=k_0\bar A_2A_3+2k_1\bar\xi_1B_3
                                       +12k_2\bar\xi_0\xi_1,\\
 (K_U)_{33}/M&=k_0|A_3|^2+k_1|B_3|^2
                                +36k_2|\xi_1|^2+36k_3|\xi_0|^2.
\end{aligned}\tag{NH33}
\]
The proof is the exact product $K_U=U^*U$
$=MF^*\diag(k_j)F$ in the original quotient orthonormal
frame. It evaluates every source Gram entry in finite original
conductor moments and the original Gamma weights, rather
than leaving an unknown positive metric.

All four leading principal minors are explicitly
\[
\begin{aligned}
 D_{U,1}&=Mk_0|\xi_0|^2,\\
 D_{U,2}&=M^2k_0k_1|\xi_0|^4,\\
 D_{U,3}&=4M^3k_0k_1k_2|\xi_0|^6,\\
 D_{U,4}&=144M^4k_0k_1k_2k_3|\xi_0|^8.
\end{aligned}\tag{NH34}
\]
For the first $r$ columns of NH31, every row below the
$r$th is zero. Their Gram determinant is the squared
modulus of their upper triangular $r$-by-$r$ determinant.
The diagonal of $F$ is $(\xi_0,\xi_0,2\xi_0,6\xi_0)$,
which proves each factor in NH34. Positivity is strict
because $M,k_j>0$ and $\xi_0\ne0$. The coefficients
$\xi_1,\xi_2,\xi_3$ remain in NH33; their cancellation
from these particular determinants follows from the proved
triangular matrix, not from setting them to zero.

In the source orthonormal frame the exact numerator map is
\[
 \mathcal M_{U,\omega}=L^{-1}\mathcal M_\omega,
 \qquad [\mathcal M_{U,\omega}]=\sqrt h\,UQ(\tau),
 \quad H_{U,\omega}=hQ(\tau)^*K_UQ(\tau).
 \tag{NH35}
\]
To record all leading phases as well, let $f_0,\ldots,f_3$
be the fixed original source frame $[\phi_v],\ldots,
[\phi_{v+3}]$. Combining NH14 with NH31 proves
\[
\begin{array}{c|c|c}
 r&m_r&\lim\tau^{-m_r}\wedge^r\mathcal M_{U,\omega}\\\hline
 1&3&-i\sqrt{hMk_0}\,\xi_0 f_0v_1^{\mathsf T}\\
 2&4&ihdM\sqrt{k_0k_1}\,\xi_0^2
                           f_{01}(v_1\wedge v_2)^{\mathsf T}\\
 3&3&-4h^{3/2}d^3M^{3/2}\sqrt{k_0k_1k_2}\,\xi_0^3
                           f_{012}(v_1\wedge v_2\wedge v_3)^{\mathsf T}\\
 4&0&-12h^8M^2\sqrt{k_0k_1k_2k_3}\,\xi_0^4.
\end{array}\tag{NH36}
\]
Each is a full matrix limit in the unchanged original
source metric, because the displayed $f$ frame is the
actual quotient orthonormal frame. The last value is the
exact determinant for every $\tau$.

Taking squared norms of NH36 gives the four exact leading
compound constants of the source Gram:
\[
\begin{aligned}
 C_{U,1}&=4hMk_0|\xi_0|^2,\\
 C_{U,2}&=20h^4M^2k_0k_1|\xi_0|^4,\\
 C_{U,3}&=80h^9M^3k_0k_1k_2|\xi_0|^6,\\
 C_{U,4}&=144h^{16}M^4k_0k_1k_2k_3|\xi_0|^8.
\end{aligned}\tag{NH37}
\]
Their exponents are $6,8,6,0$. Successive ratios, by the
same proved singular-product identity as NH18, now give
every source Gram eigenvalue in decreasing order:
\[
\boxed{\begin{aligned}
 \lambda_{U,1}&\sim4hMk_0\frac{(v!)^2}{|\mu_v|^2}|\tau|^6,\\
 \lambda_{U,2}&\sim5h^3Mk_1\frac{(v!)^2}{|\mu_v|^2}|\tau|^2,\\
 \lambda_{U,3}&\sim4h^5Mk_2\frac{(v!)^2}{|\mu_v|^2}|\tau|^{-2},\\
 \lambda_{U,4}&\sim\tfrac95h^7Mk_3\frac{(v!)^2}{|\mu_v|^2}|\tau|^{-6}.
\end{aligned}}\tag{NH38}
\]
The source receiver singular values are the positive square
roots of these expressions, because its domain is again
the displayed Euclidean coefficient space and its codomain
has the original quotient norm. Its inverse exterior norms
are fully evaluated by
\[
 \|\wedge^r\mathcal M_{U,\omega}^{-1}\|
 \sim\sqrt{C_{U,4-r}/C_{U,4}}\,
                         |\tau|^{m_{4-r}},
 \quad C_{U,0}=1,\quad 0\le r\le4.
 \tag{NH39}
\]
The source Gram condition number is consequently
\[
 \frac{\lambda_{U,1}}{\lambda_{U,4}}
 \sim\frac{20}{9h^6}
 \frac{(v+1)(v+2)(v+3)}{(b+v)(b+v+1)(b+v+2)}|\tau|^{12}.
 \tag{NH40}
\]
This follows from NH30, with the complete Gamma quotient
weight ratio retained.

The input eigenlines are still the four lines NH22: the
input covectors in NH36 are unchanged and the positive
source constants only multiply their rank-one Gram
limits. The native output singular flags are now exactly
the limiting quotient-degree flags
\[
 \operatorname{span}([\phi_v])\ \subset\
 \operatorname{span}([\phi_v],[\phi_{v+1}])\ \subset\
 \operatorname{span}([\phi_v],[\phi_{v+1}],[\phi_{v+2}])
 \ \subset\ U.
 \tag{NH41}
\]
The individual limiting native output lines are
$\C[\phi_{v+j}]$, $j=0,1,2,3$, in that order. The proof
is the same rank-one compound projection argument used
for NH22--23, now applied to the actual orthonormal
quotient frame in NH36. Formula NH29 specifies every
polynomial, its original center, phase and squared norm.

Finally the source determinant gives a direct exact
comparison with the original conductor:
\[
 \frac{\det H_{U,\omega}}{\det H_\omega}
 =\frac{144M^4\prod_{j=0}^3 k_j|\xi_0|^8}
        {12M^4b^3(b+1)^2(b+2)}
 =|\det A_{3,s}^{-1}|^2.
 \tag{NH42}
\]
To verify the last equality from FC24,
$\det A_{3,s}^{-1}=\xi_0^4\prod_{j=0}^3\rho_j/\rho_{v+j}$
and $\prod_{j=0}^3\rho_j^2=12/[b^3(b+1)^2(b+2)]$.
Both sides of NH42 are independent of the test-pole height,
while the full source and target spectra above retain
their explicitly calculated height dependence. This
completes the source calculation in NH25 rather than
leaving it as an enclosure with unspecified constants.
% END RATIONAL_NATIVE_HEIGHT

% BEGIN ARITHMETIC_RESIDUES
\section{Arithmetic residues and complete confluent jets}

The original marked polynomial is credited to The Clankers \cite{ES488}. This continuation independently checks equations (1)--(17) of the incoming
calculation \cite{ResidueIntake}. It retains the original quartic
coefficient $-1/2$, both factor signs, the projective root at
infinity and the original polynomial coordinates. It corrects
the incoming unnumbered zeta determinant following (16).
No numerical zero calculation is used.

\section{The complete original factors and arithmetic residue map}
In original coordinates $q=(a,y,z,w)$ the polynomial is
\begin{align}
b&=i+ay,&c&=-i+2ay+a^2z,\nonumber\\
d&=-iy-a(iz+2y^2)-a^2yz,&e&=2z-7iy^2+aw,\nonumber\\
f&=iw+3iy^3-4yz+a(6iy^2z+wy+4y^4)+2a^2y^3z,\nonumber\\
P(q)&=(ac,ae+bd,af+be,bf). \tag{AR1}
\end{align}
Substitution proves $ad+bc=1$ and $C(-b,a)=1$, where
$L=aU+bV$ and $C=cU^3+dU^2V+eUV^2+fV^3$.
The full factor space and its exact marked-root morphism are
\begin{align}
{\cal E}&=\{ad+bc=1,\ a^3f-a^2be+ab^2d-b^3c=1\},\nonumber\\
{\cal P}&=(ac,ae+bd,af+be,bf),\qquad
{\cal E}\longrightarrow\{(u,[U:V]):H_u(U,V)=0\},\nonumber\\
(a,b,c,d,e,f)&\longmapsto({\cal P},[-b:a]),\qquad
H_u=AU^4+U^3V+BU^2V^2+C_1UV^3+DV^4.
\tag{AR2}
\end{align}
The first equation ensures $(a,b)\ne(0,0)$. Simultaneously
negating all six factors leaves the target and marked root fixed.
For a finite simple root $r$ of
$h_u(s)=As^4+s^3+Bs^2+C_1s+D$, division by $a(U-rV)$ gives
\[
b=-ar,\quad c=A/a,\quad d=(1+Ar)/a,\quad
e=(B+r+Ar^2)/a,\quad f=(C_1+Br+r^2+Ar^3)/a.
\]
The resultant becomes $a^2h'_u(r)=1$. This proves that both
square roots give exactly the two full factor states over a
simple finite root, while a multiple root gives none.
For the original completed function \cite{DLMFXi},
\begin{equation}
r=-y-i/a,\quad
d(\xi\circ r)=\xi'(r)(ia^{-2}da-dy),\qquad
\operatorname{Res}_r\frac{\xi(s)\,ds}{h_u(s)}
=\frac{\xi(r)}{h'_u(r)}=a^2\xi(r).
\tag{AR3}
\end{equation}
Thus each pair of factor signs gives one arithmetic value
and one residue, with no sign-dependent zero test.

For the other projective chart, put $t=V/U$, $g_u(t)=H_u(1,t)$
and $t_0=-a/b$. Division by $b(V-t_0U)$ gives
\begin{equation}
-b^2g'_u(t_0)=1,\qquad
\frac{\xi(s)\,ds}{h_u(s)}
=-\frac{t^2\xi(1/t)}{g_u(t)}\,dt.
\tag{AR4}
\end{equation}
Indeed $C(-b,a)=(-b)^3C(1,t_0)=-b^2g'_u(t_0)$.
For $t_0\ne0$ the residue is
$b^2t_0^2\xi(1/t_0)=a^2\xi(r)$, proving exact equality
of the chart formulas. At $A=0$ the root $t_0=0$ is simple
because $g'_u(0)=1$. Its two full factor states are
\begin{equation}
(a,b,c,d,e,f)=
(0,\sigma i,-\sigma i,-\sigma iB,-\sigma iC_1,-\sigma iD),
\qquad \sigma=\pm1.
\tag{AR5}
\end{equation}
Only $\sigma=1$ occurs in AR1, since $a=0$ forces $b=i$.
Its original coordinates are
$(0,B,i(7B^2-C_1)/2,11B^3-2BC_1-D)$.

\section{The exact finite remainder and infinity residue}
Let $n=\deg h_u$, equal to four or three, and let $\alpha_n$
be its actual leading coefficient, respectively $A$ or $1$.
For each entire function $F$ there is exactly one polynomial
$R_F$ of degree less than $n$ such that
\begin{equation}
F=R_F+h_uG_F,\quad G_F\text{ entire},\qquad
R_F^{(j)}(r)=F^{(j)}(r)\quad(0\le j<m_r),\qquad
{\cal A}_u=\C[s]/(h_u).
\tag{AR6}
\end{equation}
Here $m_r$ is the root multiplicity. For a complete existence
proof, the ideals $(s-r)^{m_r}$ are pairwise coprime; the
Euclidean algorithm gives their Bezout identities and the
polynomial Chinese remainder construction. Prescribing the
displayed Taylor data and reducing modulo $h_u$ gives a
representative of degree less than $n$. Its quotient in AR6
has only removable singularities. Uniqueness follows because
a polynomial of degree less than $n$ divisible by $h_u$ is
zero. Hence $j_uF=[R_F]$ is an algebra morphism.
In each local Taylor basis its multiplication matrix is
triangular with $F(r)$ on the diagonal $m_r$ times, so
\begin{equation}
\det M_{j_uF}=\prod_rF(r)^{m_r},\qquad
j_uF\text{ is a unit exactly when every }F(r)\ne0.
\tag{AR7}
\end{equation}
This includes nonreduced root schemes, without asserting
existence of a finite factor state at their multiple roots.

Write $h_u=(s-r)^{m_r}k_r(s)$, $k_r(r)\ne0$. The full
principal part and the infinity residue are
\begin{equation}
\left[\frac{F}{h_u}\right]_{\mathrm{pp},r}
=\sum_{j=1}^{m_r}\frac{(F/k_r)^{(m_r-j)}(r)}{(m_r-j)!}
 (s-r)^{-j},\qquad
\operatorname{Res}_\infty\frac{F\,ds}{h_u}
=-\frac{[s^{n-1}]R_F}{\alpha_n}.
\tag{AR8}
\end{equation}
The first formula is Taylor expansion. For the second,
$F/h_u=R_F/h_u+G_F$; an entire $G_F$ contributes no
$s^{-1}$ coefficient at infinity. The rational first term
has coefficient $[s^{n-1}]R_F/\alpha_n$, and $t=1/s$
gives the minus sign. Equivalently this residue is minus
the sum of all finite residues. It has the same value at
both infinity signs AR5.

There is no holomorphic point value $\xi(\infty)$.
Indeed $\xi(2m)=m(2m-1)\pi^{-m}(m-1)!\zeta(2m)$ and
$\zeta(2m)>1$ by its positive Dirichlet series. The factorial
grows faster than any fixed power, so $\xi$ is not a polynomial.
An entire function meromorphic at infinity would have a
polynomial growth bound, and Cauchy's coefficient formula would
make it a polynomial. Thus $\xi(1/t)$ is essential at $0$.
AR8 supplies a differential residue, not an invented value
or local zero test at infinity.

At $h_0=s^3-s$, the completion and reflection give
$\xi(0)=\xi(1)=1/2$ and $\xi(-1)=\xi(2)=\pi/6$.
Evaluation at the three roots proves
\begin{align}
R_\xi&=\tfrac12+\frac{\pi-3}{12}(s^2-s),&
R_\xi^{-1}&=2+(3/\pi-1)(s^2-s)\pmod{s^3-s},\nonumber\\
\det M_{R_\xi}&=\pi/24,&
(\operatorname{Res}_0,\operatorname{Res}_1,\operatorname{Res}_{-1},\operatorname{Res}_\infty)\frac{\xi\,ds}{s^3-s}
&=(-\tfrac12,\tfrac14,\tfrac{\pi}{12},-\tfrac{\pi-3}{12}).
\tag{AR9}
\end{align}
Each value occurs at both full factor signs; only one of the
infinity signs is in the original chart. All are nonzero.

\section{The literal quartet and every completion factor}
Fix $0<\delta<1/2$, $\gamma>0$, and retain
$r_{\epsilon,\eta}=1/2+\epsilon\delta+i\eta\gamma$,
with original order $(+,+),(+,-),(-,+),(-,-)$.
Put $\rho=r_{+,+}$, $x=s-1/2$, $k=\delta^2-\gamma^2$,
$l=2\delta\gamma$, $H_\delta=(x^2-k)^2+l^2$ and
$h_\rho=-H_\delta/2$. The last factor is forced by the
original cubic coefficient one. Write
$Z=\xi(\rho)=U+iV$, $\alpha=V/l$, $\beta=U-k\alpha$.
Reflection and conjugation give the four values
$(Z,\overline Z,\overline Z,Z)$, so
\begin{equation}
R_\xi=\alpha x^2+\beta=U+\alpha(x^2-k),\qquad
\det M_{R_\xi}=|Z|^4,\qquad
R_\xi^{-1}=\frac{U-\alpha(x^2-k)}{|Z|^2}\quad(Z\ne0).
\tag{AR10}
\end{equation}
Evaluation proves the remainder by AR6. The identity
$(x^2-k)^2=-l^2$ in the quotient proves the inverse.
For both original sign states at every root,
\begin{equation}
h'_\rho(r_{\epsilon,\eta})
=4\delta\gamma(\epsilon\gamma-i\eta\delta),\quad
a_{\epsilon,\eta,\pm}^2
=\frac{1}{4\delta\gamma(\epsilon\gamma-i\eta\delta)},\quad
\operatorname{Res}_r\frac{\xi\,ds}{h_\rho}=a^2\xi(r),\quad
\operatorname{Res}_r\frac{2\xi\,ds}{H_\delta}=-a^2\xi(r).
\tag{AR11}
\end{equation}
Multiplication of the three root differences and $-1/2$
proves the derivative, retaining the original factor and sign.

On the open critical strip the exact completion is
$\xi(s)=\Phi(s)\zeta(s)$ with
$\Phi(s)=s(s-1)\pi^{-s/2}\Gamma(s/2)/2$.
It is holomorphic and nonzero there: the polynomial factors
do not vanish, and Gamma has no poles or zeros in the right
half-plane. Thus its jet is a unit, including at the doubled
quartet below, and the complete multiplication morphism is
\begin{equation}
M_{j_u\xi}=M_{j_u\Phi}M_{j_u\zeta},\qquad
M_{j_u\zeta}=M_{j_u(1/\Phi)}M_{j_u\xi}.
\tag{AR12}
\end{equation}
This is an invertible map in the same coefficient algebra,
not a replacement of the original coordinates or norm.
It is not a global unit assertion: $\Phi(0)=-1$,
$\Phi(1)=0$, $\Phi'(1)=1/2$, and $\Phi$ has poles at the
negative even integers. The zeta pole and trivial zeros
compensate these factors in the entire completion.
AR12 therefore is not applied to the root $1$ of AR9 or to
the projective root at infinity.

The exact zeta determinant is
\begin{align}
\det M_{j_{H_\delta}\zeta}
&=|\zeta(\rho)\zeta(1-\rho)|^2
=\frac{|Z|^4}{|\Phi(\rho)\Phi(1-\rho)|^2}\nonumber\\
&=\left|\frac{\Phi(\rho)}{\Phi(1-\rho)}\right|^2
                         |\zeta(\rho)|^4,\nonumber\\
\frac{\Phi(\rho)}{\Phi(1-\rho)}
&=\pi^{1/2-\rho}\frac{\Gamma(\rho/2)}{\Gamma((1-\rho)/2)},\nonumber\\
\left|\frac{\Phi(\rho)}{\Phi(1-\rho)}\right|^2
&=\pi^{-2\delta}
\frac{|\Gamma(1/4+\delta/2+i\gamma/2)|^2}
     {|\Gamma(1/4-\delta/2+i\gamma/2)|^2}.
\tag{AR13}
\end{align}
Indeed the four values are
$\zeta(\rho),\overline{\zeta(\rho)},
\overline{\zeta(1-\rho)},\zeta(1-\rho)$.
Apply AR7 and then $\xi(\rho)=\xi(1-\rho)$, dividing only
by the nonzero completion factors, never by a possible zero
of zeta. The polynomial factors cancel since
$(1-s)(-s)=s(s-1)$, leaving the displayed Gamma and $\pi$
factors. The incoming $|\zeta(\rho)|^4$ formula omits this
generally nonunit-modulus ratio. It is not an identity:
as $\rho\to1$ near the real axis inside the strip, the ratio
tends to zero while $\xi(\rho)\to1/2$. Sufficiently nearby
nonreal points have nonzero zeta values and ratio modulus
different from one.

\section{Confluence in the unchanged coefficient coordinates}
Write $H_\delta=s^4+c_3s^3+c_2s^2+c_1s+c_0$, where
\[
c_3=-2,\quad c_2=\tfrac32+2(\gamma^2-\delta^2),\quad
c_1=-\tfrac12-2(\gamma^2-\delta^2),\quad
c_0=\tfrac1{16}+\tfrac{\gamma^2-\delta^2}{2}
                              +(\delta^2+\gamma^2)^2.
\]
In the original basis $(1,s,s^2,s^3)$, multiplication by
$s$ and by the completed remainder are exactly
\begin{equation}
C_\delta=\begin{pmatrix}
0&0&0&-c_0\\1&0&0&-c_1\\0&1&0&-c_2\\0&0&1&2
\end{pmatrix},\qquad
[M_{j\xi}]=\alpha(C_\delta-\tfrac12I)^2+\beta I .
\tag{AR14}
\end{equation}
The last column is the defining polynomial relation; the
others multiply the first three original monomials.
Let $r_\eta=1/2+i\eta\gamma$. Retain the exact paired
evaluation morphism
\[
{\cal V}_\delta p=
\left(\frac{p(r_++\delta)+p(r_+-\delta)}2,\
\frac{p(r_++\delta)-p(r_+-\delta)}{2\delta},\
\frac{p(r_-+\delta)+p(r_--\delta)}2,\
\frac{p(r_-+\delta)-p(r_--\delta)}{2\delta}\right)^{\mathsf T}.
\]
Its original-coordinate rows and determinant are
\begin{equation}
(1,r_\eta,r_\eta^2+\delta^2,r_\eta^3+3r_\eta\delta^2),\quad
(0,1,2r_\eta,3r_\eta^2+\delta^2),\qquad
\det{\cal V}_\delta=16\gamma^2(\gamma^2+\delta^2).
\tag{AR15}
\end{equation}
Each pair recovers both original evaluations by $u\pm\delta v$.
The complete inverse in original coefficient coordinates, for
$\delta>0$, is the polynomial
\[
{\cal V}_\delta^{-1}(u_+,v_+,u_-,v_-)(s)
=\sum_{\eta=\pm1}\sum_{\epsilon=\pm1}
(u_\eta+\epsilon\delta v_\eta)
\frac{H_\delta(s)}
 {(s-r_{\epsilon,\eta})H'_\delta(r_{\epsilon,\eta})}.
\]
Every quotient in this sum is a polynomial of degree three;
evaluation at each marked root proves both inverse compositions.
Its product law is $(u,v)(u',v')=(uu'+\delta^2vv',uv'+vu')$,
the algebra $\C[t]/(t^2-\delta^2)$. The Vandermonde and the
two explicit pair changes give the determinant; direct
expansion gives the same expression. It remains nonzero
at $\delta=0$. No root value or factor sign has been discarded.
At zero the map is
$ {\cal V}_0p=(p(r_+),p'(r_+),p(r_-),p'(r_-))$.
Put $d_\eta=r_\eta-r_{-\eta}$ and
\[
L_\eta(s)=(s-r_{-\eta})/d_\eta,\quad
A_\eta(s)=(1-2(s-r_\eta)/d_\eta)L_\eta(s)^2,\quad
B_\eta(s)=(s-r_\eta)L_\eta(s)^2.
\]
Direct evaluation and differentiation prove its full inverse
\begin{equation}
{\cal V}_0^{-1}(u_+,v_+,u_-,v_-)
=u_+A_++v_+B_++u_-A_-+v_-B_-.
\tag{AR16}
\end{equation}
For any holomorphic $F$, its local multiplication block is
$\begin{psmallmatrix}F(r_\eta)&0\\F'(r_\eta)&F(r_\eta)\end{psmallmatrix}$.
Conjugation by the explicit $ {\cal V}_0$ gives the matrix
back in the original coefficient basis.
Put $X(\gamma)=\xi(1/2+i\gamma)\in\mathbb R$.
Differentiation gives $\xi'(r_\eta)=-i\eta X'$;
Taylor expansion in real $\delta$ gives $U\to X$,
$V/\delta\to-X'$. Hence the full limiting data are
\begin{align}
H_0&=((s-\tfrac12)^2+\gamma^2)^2,\quad
T=(s-\tfrac12)^2+\gamma^2,\quad T^2=0\text{ in }{\cal A}_0,\nonumber\\
R_{\xi,0}&=X-\frac{X'}{2\gamma}T,\qquad
{\cal V}_0M_{j\xi}{\cal V}_0^{-1}
=\diag\left(\begin{pmatrix}X&0\\-iX'&X\end{pmatrix},
            \begin{pmatrix}X&0\\iX'&X\end{pmatrix}\right),\nonumber\\
\det M_{j\xi}&=X^4,\qquad
R_{\xi,0}^{-1}=X^{-1}+\frac{X'}{2\gamma X^2}T\quad(X\ne0).
\tag{AR17}
\end{align}
The same identities follow from AR16 using
$T(r_\eta)=0$, $T'(r_\eta)=2i\eta\gamma$.
If $X'\ne0$, there are two Jordan blocks of size two with
eigenvalue $X$. If $X'=0$, the operator is exactly $XI$.
At $X=0$, $X'\ne0$, its rank is two and its square is zero;
at $X=X'=0$, it is zero. These exhaustive alternatives follow
from the matrices, without any assumed zero multiplicity.
Thus $H_0$ divides $\xi$ exactly when $X=X'=0$.

\section{Completion jets and the converging residues}
Let $\phi_\eta=\Phi(r_\eta)\ne0$, $z_\eta=\zeta(r_\eta)$.
The complete original completion data are
\begin{align}
L_\Phi(s)&=\frac{\Phi'(s)}{\Phi(s)}
=\frac1s+\frac1{s-1}-\frac{\log\pi}{2}
                  +\frac12\frac{\Gamma'(s/2)}{\Gamma(s/2)},\nonumber\\
{\cal V}_0M_\Phi{\cal V}_0^{-1}\big|_\eta
&=\begin{pmatrix}\phi_\eta&0\\\phi_\eta L_\Phi(r_\eta)&\phi_\eta\end{pmatrix},
\quad
({\cal V}_0M_\Phi{\cal V}_0^{-1})^{-1}\big|_\eta
=\begin{pmatrix}\phi_\eta^{-1}&0\\
-L_\Phi(r_\eta)\phi_\eta^{-1}&\phi_\eta^{-1}\end{pmatrix},\nonumber\\
z_\eta&=X/\phi_\eta,\quad
z_\eta'=\frac{-i\eta X'-XL_\Phi(r_\eta)}{\phi_\eta},\qquad
\det M_{j_{H_0}\zeta}=|z_+|^4=\frac{X^4}{|\phi_+|^4}.
\tag{AX1}
\end{align}
Differentiate the original product defining $\Phi$ to prove its
logarithmic derivative. Multiplication of the triangular matrices
proves their inverse. Conjugation gives
$\phi_-=\overline{\phi_+}$ and the determinant. Thus the bare
$|\zeta(\rho)|^4$ formula holds at the doubled critical-line
limit, but needs AR13 at the distinct off-line quartet.
The exact unit map preserves the vanishing order even at $X=0$.

At $r_\eta$, the original polynomial $h_{\rho,0}=-H_0/2$
has complete arithmetic principal part
\begin{equation}
\left[\frac{\xi(s)}{h_{\rho,0}(s)}\right]_{\mathrm{pp},r_\eta}
=\frac{X}{2\gamma^2}(s-r_\eta)^{-2}
+\frac{i\eta(X-\gamma X')}{2\gamma^3}(s-r_\eta)^{-1}.
\tag{AX2}
\end{equation}
Expand $-2\xi(s)/(s-r_{-\eta})^2$ at $r_\eta$.
Its constant and linear terms are
$-2X/d_\eta^2$ and
$-2\xi'(r_\eta)/d_\eta^2+4X/d_\eta^3$.
Substitution $d_\eta=2i\eta\gamma$ proves every sign in AX2.

For $\delta>0$ denote the two finite-root residues at fixed
$\eta$ by $ {\cal R}_{\epsilon,\eta}$; each still occurs twice
among the signed states. AR11 and Taylor expansion give
\begin{align}
{\cal R}_{\epsilon,\eta}(\delta)
&=\frac{\epsilon X}{4\delta\gamma^2}
  +\frac{i\eta(X-\gamma X')}{4\gamma^3}+O(\delta),\nonumber\\
\sum_{\epsilon=\pm1}{\cal R}_{\epsilon,\eta}(\delta)
&\longrightarrow\frac{i\eta(X-\gamma X')}{2\gamma^3},&
\sum_{\epsilon=\pm1}\epsilon\delta\,{\cal R}_{\epsilon,\eta}(\delta)
&\longrightarrow\frac{X}{2\gamma^2}.
\tag{AX3}
\end{align}
The remainder is a convergent local Taylor remainder for fixed
$\gamma>0$. The exact rational identity
\[
\sum_{\epsilon=\pm1}\frac{{\cal R}_{\epsilon,\eta}}
 {s-r_\eta-\epsilon\delta}
=\frac{(s-r_\eta)({\cal R}_{+,\eta}+{\cal R}_{-,\eta})
       +\delta({\cal R}_{+,\eta}-{\cal R}_{-,\eta})}
       {(s-r_\eta)^2-\delta^2}
\]
then gives precisely AX2, locally uniformly away from $r_\eta$.
At $X=0$, $X'\ne0$, the individual residues tend to
$-i\eta X'/(4\gamma^2)$ while their factors $a$ diverge.
No finite full factor state exists over the doubled roots:
$a^2h'(r)=1$ excludes every finite root and $A=-1/2$ excludes
an infinity root. AR6, AR16 and AX2 nevertheless give the exact
jet and residue limit of these escaping factors.

\section{The whole explicit escape family: an exact arithmetic exclusion}
For every real $\varepsilon>0$, the original source
$(\varepsilon,0,i/2,0)$ has the exact target polynomial
\begin{equation}
h_\varepsilon(s)=i\varepsilon s(s+i/\varepsilon)
 [(\varepsilon^2/2-1)s^2-i\varepsilon s/2+1],\qquad
r_\pm=\frac{i\varepsilon\pm\sqrt{16-9\varepsilon^2}}
                 {2(\varepsilon^2-2)}.
\tag{AX4}
\end{equation}
Multiplication gives target
$(-i\varepsilon+i\varepsilon^3/2,3i\varepsilon/2,-1,0)$.
The other two labels are $0,-i/\varepsilon$.
At $\varepsilon^2=2$ the last factor is linear, its root is
$-2i/\varepsilon$, and the fourth projective root is infinity.
At $\varepsilon^2=16/9$ the two quadratic roots coincide
and are purely imaginary; for larger squares they are also
purely imaginary. When $0<\varepsilon^2<16/9$, their real
parts are opposite. The positive part equals
$R=\sqrt{16-9\varepsilon^2}/[2(2-\varepsilon^2)]$.
Squaring positive denominators proves $0<R<1$ exactly when
$7/4<\varepsilon^2<16/9$. There
$|\Im r_\pm|=\varepsilon/[2(2-\varepsilon^2)]<3$.
This covers all parameters, including changes of degree
and multiplicity.

Here is the complete classical boundary argument used for
these labels. For $\sigma>1$, the absolutely convergent Euler
logarithm gives, when $t\ne0$,
\[
3\log|\zeta(\sigma)|+4\log|\zeta(\sigma+it)|
 +\log|\zeta(\sigma+2it)|
=\sum_{p,m\ge1}
\frac{3+4\cos(mt\log p)+\cos(2mt\log p)}{mp^{m\sigma}}\ge0,
\]
since the numerator is $2(1+\cos(mt\log p))^2$.
If $\zeta$ had a zero of order $m_0\ge1$ at $1+it$,
the product obtained by exponentiating would be
$O((\sigma-1)^{4m_0-3})\to0$, using the simple pole
at $1$ and boundedness at $1+2it$. This contradicts its
lower bound one. Thus zeta has no zeros on $\Re s=1$;
the Euler product gives nonvanishing for $\Re s>1$.
The original completion factors are nonzero there away
from $1$, where $\xi(1)=1/2$. Reflection gives the same
for $\xi$ on $\Re s\le0$. This is the classical
Hadamard--de la Vallée Poussin line-one argument, reproduced
here and not claimed as a new theorem.

The precise theta input is the original series and modular
transformation in \cite{DLMFTheta}. At $z=0,\tau=ix$ these give,
with $\psi(x)=\sum_{n\ge1}e^{-\pi n^2x}$,
$1+2\psi(x)=x^{-1/2}(1+2\psi(1/x))$.
For $\Re s>1$, termwise integration gives
$\int_0^\infty\psi(x)x^{s/2-1}dx
=\Gamma(s/2)\pi^{-s/2}\zeta(s)$.
Split the integral at one and substitute $x=1/y$ in its
lower part. The elementary contribution is
$1/(s-1)-1/s=1/[s(s-1)]$.
Multiplying by $s(s-1)/2$ and using the completion gives
\[
\xi(s)=\tfrac12+\frac{s(s-1)}2
\int_1^\infty(x^{s/2}+x^{(1-s)/2})\psi(x)\,\frac{dx}{x}.
\]
The integral and all its $s$ derivatives converge locally
uniformly on the entire plane by exponential decay, so
analytic continuation proves this displayed identity for all $s$.
On $0\le\Re s\le1$, $|\Im s|\le3$, the product
$|s(s-1)|\le10$, the powers have combined modulus at most
$2x^{1/2}$, and $n^2\ge1+3(n-1)$ for positive integers.
Using $x^{-1/2}\le1$ for $x\ge1$ gives the full integral bound
\begin{equation}
\left|\xi(s)-\tfrac12\right|
\le\frac{10e^{-\pi}}{\pi(1-e^{-3\pi})}
<\frac{8000}{47994}<\frac13,\qquad \Re\xi(s)>\frac16.
\tag{AX5}
\end{equation}
The strict rational bound uses only $\pi>3$, $e^3>20$ and
$e^9>8000$. No tail or Gamma factor was discarded.
The classification AX4, the boundary argument and AX5 prove
\begin{equation}
\xi(r)\ne0\text{ at every finite root of }h_\varepsilon
\quad(\varepsilon>0),\qquad
j_{h_\varepsilon}\xi\text{ is a unit in }\C[s]/(h_\varepsilon).
\tag{AX6}
\end{equation}
AR7 retains root multiplicity at the double root. At
$\varepsilon^2=2$ this algebra has degree three; the infinity
residue is AR8, not a fictitious value $\xi(\infty)$.

\section{The precise obstruction to extending the selected residue}
The global infinity residue AR8 is not the continuous extension
of the selected finite-root residue AR3. This failure already
occurs at the original finite source point $q_{\mathrm M}$,
not only at the missing sign chart. For complex $a\ne0$ take
$q_a=(a,0,i/2,0)$, so $q_a\to q_{\mathrm M}=(0,0,i/2,0)$
and the selected root is $-i/a$. For real $\varepsilon>0$,
the completion, reflection and the exact Gamma identity
$|\Gamma(1/2+it)|^2=\pi/\cosh(\pi t)$ give the equality
below for all $\varepsilon>0$, and its upper bound for
$0<\varepsilon\le1$:
\begin{equation}
\left|\varepsilon^2\xi(-i/\varepsilon)\right|
=\frac{\sqrt{1+\varepsilon^2}}
       {2\sqrt{\cosh(\pi/(2\varepsilon))}}
 |\zeta(1+i/\varepsilon)|
\le\sqrt{\frac{1+\varepsilon^2}{2}}\,
 [\log(1+\varepsilon^{-1})+4]e^{-\pi/(4\varepsilon)}
\longrightarrow0 .
\tag{AX7}
\end{equation}
For completeness, the elementary bound used here follows by
summation against the integer counting function: for $\Re s>0$,
$s\ne1$, and integer $N\ge1$,
\[
\zeta(s)=\sum_{n=1}^{N}n^{-s}
       +\frac{N^{1-s}}{s-1}
       -s\int_N^\infty\{x\}x^{-s-1}\,dx .
\]
It follows first for $\Re s>1$ by integrating the tail counting
function, and then throughout the stated half-plane by uniform
convergence of the integral on compacta. At $s=1+iT$, $T\ge1$,
$N=\lceil T\rceil$, its three absolute bounds are
$1+\log N$, $1/T$, and $|1+iT|/N$, whose sum is less than
$\log(T+1)+4$. Since $\cosh t\ge e^t/2$, AX7 follows.
The Gamma identity follows directly from its reflection
formula at $1/2+it$, retaining the conjugate factor.

In contrast, take $a_n=-i/(2n)$, with $n$ a positive integer.
The same original source path has selected root $2n$ and
exact selected residue
\begin{equation}
a_n^2\xi(2n)
=-\frac{2n-1}{4n}\,(n-1)!\,\pi^{-n}\zeta(2n),
\qquad |a_n^2\xi(2n)|\longrightarrow\infty.
\tag{AX8}
\end{equation}
This follows by inserting $\Gamma(n)=(n-1)!$ in the original
completion. The lower bound $\zeta(2n)>1$ and the ratio
$n/\pi$ of successive factorial-over-$\pi$ factors prove
divergence. Both $q_{a_n}$ and $q_{1/(2n)}$ tend to the
same original point $q_{\mathrm M}$, but their selected
residues have incompatible limits. Hence AR3 has no continuous
extension there. The same conclusion holds at the negative
factor state in the full completion: its residue is identical
for each $a$, and its six factors have the opposite finite limit.
At the limiting target $h_0=s^3-s$, AR9 instead gives the
global residue $\operatorname{Res}_\infty=(3-\pi)/12\ne0$.
It therefore cannot repair this noncontinuity by being assigned
as the missing selected-root value. The exact relation that
does hold is the global residue identity AR8 and the complete
finite remainder map AR6, with all marked signs retained.
Every fixed injective receiving map transports these two
sequences and their unequal arithmetic limits unchanged;
no change of the original Gamma norm removes the obstruction.

The exact certificate
\texttt{verify\_arithmetic\_residues\_and\_jets.py}
checks the original factors, both residue charts, the quartet
matrix and inverse, the confluent evaluation and Hermite inverse,
the principal-part constants and the original escape polynomial.
It treats transcendental values and derivatives as formal scalars.
% END ARITHMETIC_RESIDUES

% BEGIN VISUAL_ATLAS
\clearpage

\section{Continuous positivity and the exact original native return}
The five isolated cubic pole certificates in WP26 of \cite{WeilCertificates}
are strengthened here to one continuous real interval. The
proof uses the actual finite-height theorem of Platt and Trudgian
and the unconditional zero-count bound of Hasanalizade, Shen
and Wong. Every unverified zero is retained in the tail bound;
no location of such a zero is assumed. The original factor
$h=3$, the cubic coefficient coordinates and both original
conductor norms are preserved by explicit receiving maps.

\section{The complete form and the two human theorems used}
For every real $\tau$ set
\begin{equation}
\omega_\tau=2+i\tau,\qquad h=3,\qquad
U_\tau(s)=\frac{s+1-i\tau}{2+i\tau-s},\qquad
\varphi_j^\tau(s)=\frac{\sqrt3}{2+i\tau-s}U_\tau(s)^j
\quad(0\le j\le3).
\tag{CP1}
\end{equation}
For $p=(p_0,p_1,p_2,p_3)^{\mathsf T}\in\C^4$ write
$F_p^\tau=\sum_{j=0}^3p_j\varphi_j^\tau$.
Use the full reflected Weil form
\begin{equation}
\W(F,G)=\sum_\rho m_\rho
       \overline{F(1-\bar\rho)}G(\rho),\qquad
T_3(\omega_\tau)=\bigl(\W(\varphi_i^\tau,\varphi_j^\tau)\bigr)_{i,j=0}^3 .
\tag{CP2}
\end{equation}
Here $\rho$ ranges over all distinct nontrivial zeta zeros,
and $m_\rho$ is its actual multiplicity. Absolute convergence
is proved below. Thus $p^*T_3p=\W(F_p^\tau,F_p^\tau)$.

The first human input is \cite[Theorem 1]{PlattTrudgian}, using only
the strictly smaller height
\begin{equation}
T_0=3\cdot10^{12}:
\qquad \Re\rho=\tfrac12
\quad\hbox{for every nontrivial zero with }|\Im\rho|\le T_0.
\tag{CP3}
\end{equation}
The published theorem reaches $3\,000\,175\,332\,800$.
The negative-height assertion in CP3 follows by conjugation.
The second input is \cite[Corollary 1.2]{HSW}, with its exact
original constants and counting convention:
\begin{equation}
N(T)=\sum_{0<\Im\rho\le T}m_\rho,\quad
|N(T)-f(T)|\le E(T)\quad(T\ge e),\qquad
f(T)=\frac{T}{2\pi}\log\frac{T}{2\pi e},\quad
E(T)=.1038\log T+.2573\log\log T+9.3675.
\tag{CP4}
\end{equation}
These are previously proved unconditional theorems, not
additional assumptions about the zeros being tested.

\section{Four separated actual zero intervals at every pole height}
For every $A\ge100$, CP4 implies
\begin{equation}
N(A+100)-N(A)>13\log A-52>0.
\tag{CP5}
\end{equation}
Here are all constants in the deduction. Since
$\log\log T\le\log T$ for $T\ge e$,
$E(T)\le\log T+10$. Also
$f'(T)=(2\pi)^{-1}\log(T/(2\pi))$ increases. Hence
\[
N(A+100)-N(A)
\ge\frac{100}{2\pi}\log\frac{A}{2\pi}
                 -\log A-\log(A+100)-20 .
\]
The elementary bounds $\pi<22/7$, $\log(2\pi)<2$ and
$\log2<1$ show that the right side is strictly larger than
$15(\log A-2)-2\log A-22=13\log A-52$.
Here $\log A-2>0$, so the inequality in its multiplier is
in the correct direction. Finally $\log100>4$:
$e<3$ gives $e^4<81<100$. For the other elementary bound,
$e^2>7>2\pi$ follows from the first five positive exponential
terms, whose sum is exactly seven. This proves CP5 including
its domain and endpoints.

Set
\begin{equation}
T_*=\frac{99}{100}T_0=2\,970\,000\,000\,000 .
\tag{CP6}
\end{equation}
For $|\tau|\le T_*$, let $a=|\tau|$ and choose the least
actual zero ordinate in each interval
\begin{equation}
I_j(a)=(a+100+200j,\ a+200+200j]\quad(j=0,1,2,3).
\tag{CP7}
\end{equation}
CP5 proves that each interval contains a zero. Such a least
ordinate exists because there are finitely many zeros in
each bounded interval. Every interval is strictly below
$T_0$, since $T_*+800<T_0$, so CP3 proves that these selected
zeros lie on the critical line. No simplicity is needed.
Write their positive ordinates as $\gamma_j$. For $\tau\ge0$
set $t_j=\gamma_j$; for $\tau<0$ set $t_j=-\gamma_j$.
These are four actual critical zeros in either case.
They are distinct, and their relative heights satisfy
\begin{equation}
100+200j<|t_j-\tau|\le200+200j\le800,\qquad
|t_j-t_i|\ge100(2|j-i|-1)\quad(i\ne j).
\tag{CP8}
\end{equation}
The intervals supply these inequalities directly; no continuity
of the selected zero labels as a function of $\tau$ is needed.
The resulting matrix bound will hold for every real $\tau$
in CP6, with the same numerical constant.

\section{An exact rational Vandermonde inverse bound}
For a critical zero of height $t$, put $z(t)=U_\tau(1/2+it)$.
Direct substitution into CP1 gives
\begin{equation}
z(t)=\frac{3/2+i(t-\tau)}{3/2-i(t-\tau)},\quad |z(t)|=1,\qquad
|z(t)-z(u)|
=\frac{3|t-u|}
 {\sqrt{9/4+(t-\tau)^2}\sqrt{9/4+(u-\tau)^2}}.
\tag{CP9}
\end{equation}
Retain the exact constants
\begin{equation}
L=800^2+\frac94=\frac{2560009}{4},\qquad
w_*=\frac3L,\qquad d_*=\frac{300}{L}.
\tag{CP10}
\end{equation}
For $z_j=z(t_j)$, each of the four row weights
$w_j=3/[9/4+(t_j-\tau)^2]$ is at least $w_*$, and
CP8--9 give
$|z_j-z_i|\ge d_*(2|j-i|-1)$.

Let $V=(z_j^\ell)_{0\le j,\ell\le3}$.
Its inverse has column $j$ equal to the coefficient vector of
the exact Lagrange polynomial
\begin{equation}
\ell_j(z)=\frac{\prod_{i\ne j}(z-z_i)}
                     {\prod_{i\ne j}(z_j-z_i)}.
\tag{CP11}
\end{equation}
Evaluation at the four distinct $z_i$ proves both inverse
compositions. Since $|z_i|=1$, the four numerator coefficient
moduli are at most $1,3,3,1$, so their squared Euclidean norm
is at most twenty. The three denominator separations are
at least $(d_*,3d_*,5d_*)$ for $j=0,3$ and
$(d_*,d_*,3d_*)$ for $j=1,2$. Therefore
\begin{equation}
\|V^{-1}\|_2^2\le\|V^{-1}\|_{\mathrm F}^2
\le20d_*^{-6}\left(\frac2{225}+\frac2{9}\right)
=\frac{208}{45}d_*^{-6}.
\tag{CP12}
\end{equation}
This bound retains all four separated intervals; it does not
replace their product by a determinant divided by a trace.

Taking one copy of each selected zero gives a positive Gram
$H_{\mathrm{sel},\tau}=V^*\diag(w_0,w_1,w_2,w_3)V$.
Because $p=V^{-1}Vp$, CP12 proves for every coefficient vector
\begin{equation}
p^*H_{\mathrm{sel},\tau}p\ge b_*\|p\|_2^2,\qquad
b_*=\frac{45}{208}w_*d_*^6
=\frac{135\cdot1200^6}{52\cdot2560009^7}>10^{-26}.
\tag{CP13}
\end{equation}
The last inequality is exact integer arithmetic: the difference
$135\cdot1200^6\cdot10^{26}-52\cdot2560009^7$ is the positive
integer
$2839912978528311935972297436813173952969845612$.
Multiplicities can only increase the contribution of these
verified zeros.

\section{Every unverified zero and the complete reflected tail}
Write a point in the closed strip as $s=\beta+it$,
$0\le\beta\le1$, and put $x=t-\tau$. For $|x|\ge10$,
the exact factors in CP1 imply
\[
|U_\tau(s)|^2
=\frac{(\beta+1)^2+x^2}{(2-\beta)^2+x^2}
\le1+\frac3{x^2}<\left(\frac{51}{50}\right)^2,\qquad
|U_\tau'(s)|=\frac3{|\omega_\tau-s|^2}\le\frac3{x^2}.
\]
Differentiate each of the four original functions without
changing its degree or the factor $\sqrt3$. For $0\le j\le3$,
\begin{equation}
|(\varphi_j^\tau)'(s)|
\le\frac{\sqrt3}{x^2}
\left[\left(\frac{51}{50}\right)^3+
 \frac9{|x|}\left(\frac{51}{50}\right)^2\right]
<\frac4{x^2},\qquad
|(F_p^\tau)'(s)|\le\frac8{x^2}\|p\|_2 .
\tag{CP14}
\end{equation}
For the strict numerical inequality it is enough to use
$\sqrt3<7/4$ and $|x|\ge10$, giving the rational upper
coefficient $3.495744<4$.
The last inequality is Cauchy--Schwarz on the four derivatives.

The reflected zeros $\rho=\beta+it$ and
$\rho^\vee=1-\bar\rho=1-\beta+it$ have equal multiplicity.
Integrating CP14 on their horizontal segment, whose length
is at most one and whose imaginary part stays $t$, gives
\[
|F_p^\tau(\rho)-F_p^\tau(\rho^\vee)|
\le\frac8{|t-\tau|^2}\|p\|_2.
\]
Every noncritical reflected pair of multiplicity $m$ consequently
has its exact contribution bounded below by
\begin{align}
2m\Re(\overline{F_p^\tau(\rho^\vee)}F_p^\tau(\rho))
&=m\bigl(|F_p^\tau(\rho)|^2+|F_p^\tau(\rho^\vee)|^2
 -|F_p^\tau(\rho)-F_p^\tau(\rho^\vee)|^2\bigr)\nonumber\\
&\ge-\frac{64m}{|t-\tau|^4}\|p\|_2^2.
\tag{CP15}
\end{align}
Critical zeros, whether verified or not, contribute nonnegatively.
For each positive-height noncritical pair, its two members
contribute $2m$ to $N$. The conjugate negative-height pair
has the same multiplicity and relative height $-t-\tau$.
Therefore the total possible negative contribution from all
unverified zeros is bounded by
\[
32\|p\|_2^2
\sum_{\Im\rho>T_0}m_\rho
\left[\frac1{|\Im\rho-\tau|^4}
      +\frac1{|\Im\rho+\tau|^4}\right].
\]
For $|\tau|\le T_*$ and $t>T_0$,
$|t\pm\tau|\ge t/100>T_0/100>10$. Thus CP14 applies,
and this last expression is at most
$64\cdot100^4\|p\|_2^2\sum_{\Im\rho>T_0}m_\rho/(\Im\rho)^4$.

The zero count CP4 also gives $N(T)\le T\log T$ for $T\ge100$.
In fact $f(T)<T\log T/6$, and
$E(T)\le\log T+10<T$, while $\log T>4$.
The last strict bound follows at $100$ and persists because
$T-\log T-10$ increases. Stieltjes integration by parts now gives
\begin{align}
\sum_{\Im\rho>T_0}\frac{m_\rho}{(\Im\rho)^4}
&=-\frac{N(T_0)}{T_0^4}
     +4\int_{T_0}^\infty\frac{N(t)}{t^5}\,dt\nonumber\\
&\le\frac4{T_0^3}
              \left(\frac{\log T_0}{3}+\frac19\right).
\tag{CP16}
\end{align}
The boundary term at infinity is zero by the same count bound.
Every term above the verified height has therefore been retained
in a convergent upper estimate.
For completeness CP4 also proves absolute convergence of CP2:
for fixed $\tau$, each $\varphi_j^\tau(\rho)=O(|\Im\rho|^{-1})$,
uniformly across the strip at large height, and
$\sum_{\Im\rho>R}m_\rho/(\Im\rho)^2<\infty$ follows by the
same Stieltjes calculation. Conjugate negative heights satisfy
the same bound. Hence the reflection grouping used above is valid.

Combining CP15--16, the whole unknown negative tail is at most
$e_*\|p\|_2^2$, where the following constants are exact:
\begin{equation}
e_*=\frac{256\cdot100^4}{T_0^3}\frac{91}{9}
=\frac{23296}{243\cdot10^{28}}
<\frac{24}{25}\,10^{-26},\qquad
\mu_*=b_*-e_*>4\cdot10^{-28}.
\tag{CP17}
\end{equation}
Here $\log T_0<30$. One elementary exact check is that the
finite sum $\sum_{k=0}^{30}30^k/k!$ exceeds $3\cdot10^{12}$,
so $e^{30}>T_0$. The rational comparison for $e_*$ is
$23296/243<96$, equivalent to $23296<23328$.
The last lower bound follows from CP13:
$10^{-26}-(24/25)10^{-26}=4\cdot10^{-28}$.

Every remaining zero below $T_0$ is critical by CP3 and
contributes nonnegatively. Adding CP13 and subtracting this
complete tail proves the continuous, unconditional result
\begin{equation}
\boxed{\quad
T_3(2+i\tau)\succeq\mu_* I_4
\succeq4\cdot10^{-28}I_4
\quad\hbox{for every real }|\tau|\le2\,970\,000\,000\,000.
\quad}
\tag{CP18}
\end{equation}
The same coefficient bound holds for all four coordinates
and all relative phases simultaneously. The selection of four
verified zeros proves one positive contribution; the rest
of the complete form is controlled, not omitted.

\section{Return to the unchanged original conductor and its Gamma norm}
Retain any fixed admitted reference conductor from FC21--32
of \cite{Fable029}, its actual symbol order $v$, its original shifts
and coefficients, and its exact source quotient and target:
\[
T_A:\mathcal P_{v+3}/\mathcal P_{v-1}\longrightarrow\mathcal P_3,
\qquad T_Ap(S')=\sum_{u,w=0}^8 a_{uw}p(S'+\beta_{8,u,w}).
\]
The reference roots are denoted $r_{\ell,*}$ and remain fixed
as the pole height $\tau$ varies. With the original marked
collision matrix $K$, let
\[
\ell_{\ell,*}(S)=
\prod_{b\ne\ell}\frac{S-r_{b,*}}{r_{\ell,*}-r_{b,*}},\qquad
y_\ell(S')=\sum_{u,w=0}^8a_{uw}
 (S'+\beta_{8,u,w})^v\ell_{\ell,*}(S'+\beta_{8,u,w}).
\]
Writing $Y$ for their columns in the original target monomials,
the retained map is $\Psi=YK^{-1}$. Its source partner is
$\Phi_{\rm FC}$ with $T_A\Phi_{\rm FC}=\Psi$.
All these original maps are the proved isomorphisms on the
actual quotient; the reference periods, moments and norms do
not depend on $\tau$.

Let $c_*=a_{\rm amp}/2-4$ and $\kappa=c_*-1/2$.
The translation matrix and rational numerator matrix are
\begin{align}
(C_\kappa)_{r\ell}&=\binom{\ell}{r}\kappa^{\ell-r}
 \quad(0\le r\le\ell\le3),\qquad
C_\kappa^{-1}=C_{-\kappa},\nonumber\\
R_\tau&=\sqrt3
\begin{pmatrix}
\omega^3&b\omega^2&b^2\omega&b^3\\
-3\omega^2&\omega^2-2b\omega&2b\omega-b^2&3b^2\\
3\omega&b-2\omega&\omega-2b&3b\\
-1&1&-1&1
\end{pmatrix},
\quad \omega=2+i\tau,\quad b=1-i\tau,\quad \det R_\tau=3^8.
\tag{CP19}
\end{align}
The $j$th column is exactly the coefficient vector of
$\sqrt3(z+b)^j(\omega-z)^{3-j}$.
For clarity its full inverse, also fixing every sign, is
\[
(R_\tau^{-1})_{j\ell}
=3^{-7/2}\sum_{r=0}^{\ell}
\binom{\ell}{r}\omega^r(-b)^{\ell-r}\binom{3-\ell}{j-r}.
\]
Indeed $u=(z+b)/(\omega-z)$ gives
$z=(\omega u-b)/(1+u)$ and $\omega-z=3/(1+u)$;
expanding $3^{-7/2}(1+u)^3p((\omega u-b)/(1+u))$
proves the inverse formula by substitution.
The determinant follows by changing to $t=\omega-z$:
the columns $\sqrt3(3-t)^jt^{3-j}$ become triangular
after reversal, and the two reversal signs cancel.

The exact original receiver and its inverse are consequently
\begin{equation}
B_\tau=R_\tau^{-1}C_\kappa\Psi,\qquad
B_\tau^{-1}=\Psi^{-1}C_{-\kappa}R_\tau,\qquad
F_{B_\tau q}^\tau(z)=\frac{(\Psi q)(z+\kappa)}{(\omega-z)^4}.
\tag{CP20}
\end{equation}
Multiplying the numerator coefficients proves the last identity.
It also proves the original conductor return
$B_\tau=R_\tau^{-1}C_\kappa T_A\Phi_{\rm FC}$.
Thus CP18 gives the actual, continuously parameterized form
on the original Fable coordinates:
\begin{equation}
\W_{\tau,\rm original}
=B_\tau^*T_3(2+i\tau)B_\tau
\succeq\mu_* B_\tau^*B_\tau
\quad(|\tau|\le T_*).
\tag{CP21}
\end{equation}
This is a strict positive form because $B_\tau$ is invertible.
The coefficient norm on the right is not substituted for
the original physical norm; its exact comparison is next.

For the original positive parameter $s_0\in\{1,a_{\rm amp}\}$,
retain the full target measure
\[
dm_{s_0}(y)=
\frac{(2\pi)^{s_0/2}2^{s_0/2}}{4\pi\Gamma(s_0/2)}
 \left|\Gamma(s_0/4+iy/2)\right|^2dy,
\quad
\|p\|_{\Gamma,W}^2=\int_{\R}|p(c_*+iy)|^2\,dm_{s_0}(y).
\]
Let $\mathcal H$ be this Gram in $(1,S',S'^2,S'^3)$ and
$G=\Psi^*\mathcal H\Psi$. In rational coefficients the same
norm has the exact Gram
\begin{equation}
H_\tau=B_\tau^{-*}GB_\tau^{-1}
=R_\tau^*C_{-\kappa}^*\mathcal HC_{-\kappa}R_\tau,\qquad
(H_\tau)_{ij}
=3\int_{\R}
 (3/2+i(y-\tau))^{3+j-i}
 (3/2-i(y-\tau))^{3+i-j}\,dm_{s_0}(y).
\tag{CP22}
\end{equation}
Here $H_\tau$ denotes the full native Gram; the selected-zero
Gram $H_{\mathrm{sel},\tau}$ was used only in CP13.
To prove CP22, evaluate its numerator at
$S'=c_*+iy$: the translated argument is exactly $1/2+iy$.
The two factors are conjugate, with exponents displayed
after multiplying conjugate column $i$ and column $j$.
This retains the original center, total mass and every cross term.

Put $M=(2\pi)^{s_0/2}$ and $\beta_0=s_0/2$. The original
Gamma moments are
$m_0=M$, $m_2=M\beta_0$,
$m_4=M(3\beta_0^2+2\beta_0)$,
$m_6=M(15\beta_0^3+30\beta_0^2+16\beta_0)$, with odd
moments zero, as proved for this unchanged measure in
FC56a/RW9 \cite{Fable029}. Hence every diagonal entry of CP22 is
the following explicit positive polynomial:
\begin{align}
g_0(\tau)=3\bigg[&
\frac{729}{64}m_0+\frac{243}{16}(m_2+\tau^2m_0)
+\frac{27}{4}(m_4+6\tau^2m_2+\tau^4m_0)\nonumber\\
&+m_6+15\tau^2m_4+15\tau^4m_2+\tau^6m_0\bigg],
\qquad \operatorname{tr}H_\tau=4g_0(\tau).
\tag{CP23}
\end{align}
Expansion of $3(9/4+(y-\tau)^2)^3$ proves this formula.
Since a positive Hermitian matrix has largest eigenvalue
at most its trace, CP18 yields the exact native lower bound
\begin{equation}
q^*\W_{\tau,\rm original}q
\ge\frac{\mu_*}{4g_0(\tau)}\,q^*Gq
\quad(|\tau|\le T_*).
\tag{CP24}
\end{equation}
Every coefficient of $g_0$ as a polynomial in $\tau^2$ is
positive. Thus $\mu_*/[4g_0(T_*)]$ is also a single explicit
positive lower constant for the full continuous interval
in the unchanged original target norm.

\section{All eight signed states and the exact remaining kernel}
In the original eight-state construction, retain the actual
four columns of pair means and half differences
$E=(e_1,\ldots,e_4)$, $O=(o_1,\ldots,o_4)$, with
$e_j=(q_{j,+}+q_{j,-})/2$ and
$o_j=(q_{j,+}-q_{j,-})/2$. These are the complete original
source points over the literal quartet, including both signs,
as defined and proved in SE1--7 of \cite{Fable029}. Their determinant
is
$1792\chi\delta^2\gamma^2(2\delta^2\gamma^2-\delta^2+\gamma^2)$,
where $\chi=\pm1$. It is nonzero on the original admitted
domain $0<\delta<1/2$, $\gamma>2$. Thus $O^{-1}$ is the
exact original inverse on this domain, with no extra state
or sign omitted. Write a full signed vector as $(\alpha,\beta)$,
whose actual evaluation is $z=E\alpha+O\beta$.

Let $G_Y=Y^*\langle\, ,\,\rangle_{\Gamma,W}Y$ be the
original target label Gram. The complete target norm and
arithmetic form are
\begin{align}
\mathbb H_W&=
\begin{pmatrix}I&0\\E&O\end{pmatrix}^{*}
\diag(G_Y,G)
\begin{pmatrix}I&0\\E&O\end{pmatrix},\nonumber\\
\mathbb W_\tau&=[E\ O]^*B_\tau^*T_3(2+i\tau)B_\tau[E\ O],
\qquad
S_\tau=\begin{pmatrix}I&0\\-O^{-1}E&O^{-1}B_\tau^{-1}\end{pmatrix}.
\tag{CP25}
\end{align}
Direct multiplication proves
\[
\begin{pmatrix}I&0\\E&O\end{pmatrix}S_\tau
=\diag(I,B_\tau^{-1}),\qquad
[E\ O]S_\tau=[0\ B_\tau^{-1}].
\]
Consequently CP18 proves the full simultaneous comparison
\begin{equation}
S_\tau^*\mathbb W_\tau S_\tau
=\diag(0_4,T_3(2+i\tau))
\succeq\diag(0_4,\mu_*I_4),\qquad
S_\tau^*\mathbb H_W S_\tau=\diag(G_Y,H_\tau).
\tag{CP26}
\end{equation}
The inverse of $S_\tau$ is
$\begin{psmallmatrix}I&0\\B_\tau E&B_\tau O\end{psmallmatrix}$,
verified by multiplying either way. The complete kernel, rank,
and target quotient coercivity are therefore exactly
\begin{equation}
\ker\mathbb W_\tau
=\{(\alpha,-O^{-1}E\alpha):\alpha\in\C^4\},\qquad
\operatorname{rank}\mathbb W_\tau=4,\qquad
\mathbb W_\tau|_{\rm quotient}
\succeq\frac{\mu_*}{4g_0(\tau)}
                         \mathbb H_W|_{\rm quotient}.
\tag{CP27}
\end{equation}
To specify the quotient norm without an informal deletion,
the full coordinate pair is $(\alpha,z)$ and its squared
norm is $\alpha^*G_Y\alpha+z^*Gz$.
For fixed $z$ the unique minimum over the kernel coset is
at $\alpha=0$. Its section is
$(0,O^{-1}z)$. In rational coordinates $d=B_\tau z$ it is
$(0,O^{-1}B_\tau^{-1}d)$ and has squared norm $d^*H_\tau d$.
This proves exactly the quotient and the constant in CP27.
There are no additional arithmetic null directions in this
continuous pole range, because CP18 makes $T_3$ positive definite.

The source norm remains distinct. Let
$G_X=X^*\langle\, ,\,\rangle_{\Gamma,U}X$ and
$G_U=\Phi_{\rm FC}^*\langle\, ,\,\rangle_{\Gamma,U}\Phi_{\rm FC}$,
with the original source quotient norm, and put
$H_{U,\tau}=B_\tau^{-*}G_UB_\tau^{-1}$.
For the actual source pair
$(X\alpha,\Phi_{\rm FC}z)$ the same multiplication gives
\begin{equation}
S_\tau^*\mathbb H_U S_\tau=\diag(G_X,H_{U,\tau}),\qquad
\mathbb W_\tau|_{\rm source\ quotient}
\succeq
\frac{\mu_*}{\operatorname{tr}H_{U,\tau}}\,
                 \mathbb H_U|_{\rm quotient}.
\tag{CP28}
\end{equation}
The source quotient again minimizes at $\alpha=0$;
$H_{U,\tau}$ is its exact positive Gram, with every original
inverse coefficient specified by FC24. The actual conductor
maps this pair to $(Y\alpha,\Psi z)$ because
$T_AX=Y$ and $T_A\Phi_{\rm FC}=\Psi$. Thus CP25--28 preserve
both original metrics and their precise morphism.

There is also a uniform source constant with every original
inverse weight explicitly present. Write the retained conductor
symbol as
$g_A(t)=t^{-v}e^{-4w(t)}E_A(w(t))=\sum_{n\ge0}c_nt^n$,
where $w(t)=-i\arctan t$ and
$c_0=(-i)^v\mu_v/v!\ne0$ is its actual leading moment.
The four inverse coefficients are exactly
\[
\eta_0=c_0^{-1},\quad \eta_1=-c_1/c_0^2,\quad
\eta_2=(c_1^2-c_0c_2)/c_0^3,\quad
\eta_3=(-c_1^3+2c_0c_1c_2-c_0^2c_3)/c_0^4 .
\]
Their convolution with $c_0,c_1,c_2,c_3$ gives $(1,0,0,0)$.
With the original weights $\rho_n=\sqrt{n!/(s_0/2)_n}$,
the actual inverse matrix in the original orthonormal source
and target frames has entry
$\eta_{j-i}\rho_j/\rho_{v+i}$ for $0\le i\le j\le3$,
and zero below the diagonal, by FC24/WCF6.
Thus its complete squared Frobenius norm and uniform
source comparison are
\begin{equation}
J_A=\sum_{0\le i\le j\le3}
 |\eta_{j-i}|^2\frac{\rho_j^2}{\rho_{v+i}^2}>0,\qquad
H_{U,\tau}\preceq J_AH_\tau,\qquad
\mathbb W_\tau|_{\rm source\ quotient}
\succeq\frac{\mu_*}{4J_Ag_0(T_*)}
                  \mathbb H_U|_{\rm quotient}
\quad(|\tau|\le T_*).
\tag{CP29}
\end{equation}
Indeed a matrix has operator norm at most its Frobenius norm,
so the literal source inverse gives
$\|T_A^{-1}p\|_{\Gamma,U}^2\le J_A\|p\|_{\Gamma,W}^2$.
Substitute $p=C_{-\kappa}R_\tau d$ to obtain the matrix
inequality; CP23--24 prove the last bound. Orthogonal
coordinates are used only to calculate the norm of this
unchanged physical inverse: all masses and quotient weights
are those of the original source and target.

These results establish a continuous positive cubic test
range and its exact full eight-state kernel. The proof uses
the original full-Weil sum and the original cubic receiver
throughout. It makes no assertion about degrees above three
or about poles outside the explicit interval CP6.

The certificate
\texttt{verify\_continuous\_pole\_cubic\_positivity.py}
checks all rational constant comparisons by integer arithmetic,
the coefficient-norm and reflected-pair multiplicity factors,
the exact chord identity, and the native diagonal-moment formula.
Its output is retained separately. The human theorems CP3--4
are cited mathematical inputs, not assertions recertified by
that short algebra script.

\section{Complete original source and target coercivity}
The analytic inputs to this note are the proved full-zero-sum
inequalities WP14 and WP23, and the certified five-pole matrices
WP26, in \cite{WeilCertificates}. The receiving maps and their two distinct
original metrics are RW1--20 \cite{Fable029}; their exact native
matrices are NH1--42 \cite{Fable029}. We use those results at the
stated poles, retain every zero multiplicity in the Weil form,
and calculate the entire source and target coercivity on the
four-dimensional quotient of the full eight-state space.
No assertion about untested degrees or poles is used.

\section{The full signed space, its exact kernel and both minimum lifts}
Keep the four marked roots $r_j$ of the literal quartet polynomial
$h_\rho(r)=Ar^4+r^3+Br^2+Cr+D$, where
\[
 A=-\tfrac12,\quad B=\delta^2-\gamma^2-\tfrac34,\quad
 C=\gamma^2-\delta^2+\tfrac14,\quad
 D=-\tfrac1{32}-\tfrac{\gamma^2-\delta^2}{4}
                         -\tfrac{(\delta^2+\gamma^2)^2}{2},
 \quad0<\delta<\tfrac12,\quad\gamma>2.
\]
At each root retain both signs of $a_j^2h_\rho'(r_j)=1$,
put $d_j=h_\rho'(r_j)$, and use the actual mean and difference
columns of the inverse fibre:
\[
 e_j=\begin{pmatrix}0\\-r_j\\3id_j\\
 (B-17r_j+Ar_j^2)d_j-2Ad_j^2\end{pmatrix},\qquad
 o_j=a_j\begin{pmatrix}1\\-id_j\\Ad_j^2+2r_jd_j\\
 7ir_j^2d_j-13id_j^2\end{pmatrix},\quad E=(e_j),\ O=(o_j).
 \tag{NW1}
\]
These are SE1, proved from the original ES polynomial
\cite{ES488,Fable029}. The complete determinant proved in SE7 is
\[
 \det O=1792\chi\delta^2\gamma^2
 (2\delta^2\gamma^2-\delta^2+\gamma^2)\ne0,
 \qquad \chi=A^2\det(r_j^n)_{n,j}\prod_ja_j\in\{1,-1\}.
\]
Thus the inverse in every formula below is the explicit
$O^{-1}=\operatorname{adj}(O)/\det O$; no sign partner is
discarded.

Let $f_{j,+},f_{j,-}$ be the eight separate basis states.
Set $\iota_+c_j=(f_{j,+}+f_{j,-})/2$ and
$\iota_-c_j=(f_{j,+}-f_{j,-})/2$. Every state is uniquely
$v=\iota_+\alpha+\iota_-\beta$; define its complete evaluation
$z=E\alpha+O\beta$. The original source and target states
are, respectively,
\[
 (X\alpha,\Phi z)\in U\oplus U,\qquad
 (Y\alpha,\Psi z)\in W\oplus W,
 \quad T_AX=Y,\quad T_A\Phi=\Psi.
 \tag{NW2}
\]
All four maps are the existing exact conductor maps of
FC26--35, with the original periods and moments. Put
$G_X=X^*\langle\ ,\ \rangle_U X$,
$G_Y=Y^*\langle\ ,\ \rangle_W Y$,
$G_U=\Phi^*\langle\ ,\ \rangle_U\Phi$ and
$G=\Psi^*\langle\ ,\ \rangle_W\Psi$.
Each direct sum in NW2 has its original orthogonal direct-sum
norm, so the exact squared norms are
\[
 \|v\|_{8,U}^2=\alpha^*G_X\alpha+z^*G_Uz,
 \qquad
 \|v\|_{8,W}^2=\alpha^*G_Y\alpha+z^*Gz.
 \tag{NW3}
\]
No summand uses the norm from the other side of the conductor.

For the rational receiver retain
$\omega=2+i\tau$, $\tau\in\{0,14,21,100,1000\}$,
$h=3$, $d=h/2=3/2$, $c_*=a_{\rm amp}/2-4$ and
$\kappa=c_*-1/2$. Let $R_\omega$ be the complete coefficient
matrix of
$N_j(s)=\sqrt3(s+\bar\omega-1)^j(\omega-s)^{3-j}$,
$0\le j\le3$, and let $C_\kappa$ act by
$p(S')\mapsto p(s+\kappa)$. Its entries are
$(C_\kappa)_{r\ell}=\binom\ell r\kappa^{\ell-r}$ for
$r\le\ell$; its inverse is $C_{-\kappa}$. Define the exact
rational coefficient column and receiver by
\[
 d_v=B_\omega z,\qquad B_\omega=R_\omega^{-1}C_\kappa\Psi,
 \qquad
 F_v(s)=\sum_{j=0}^3(d_v)_j\frac{\sqrt3}{\omega-s}
                     \left(\frac{s+\bar\omega-1}{\omega-s}\right)^j.
 \tag{NW4}
\]
RW3 gives every coefficient of $R_\omega$ and its inverse;
$\det R_\omega=3^8\ne0$. Thus $B_\omega$ is an actual
isomorphism. The full Weil form is exactly
\[
 \mathbb W_\omega(v,v)=\W(F_v,F_v)=d_v^*T_3(\omega)d_v.
 \tag{NW5}
\]
This is the full zero sum of WP3 and RW12, including all
nontrivial zero multiplicities, not a four-zero truncation.

At all five poles $T_3(\omega)$ is strictly positive by
WP26, and at $\tau=0$ WP23 also supplies an analytic positive
lower bound. Therefore the full kernel in the eight-state
space is precisely
\[
 \mathcal N=\{\iota_+\alpha-\iota_-O^{-1}E\alpha:
                              \alpha\in\C^4\}.
 \tag{NW6}
\]
Its four explicit independent generators, in the original
eight marked basis states, are
\[
 n_j=\frac{f_{j,+}+f_{j,-}}2
       -\frac12\sum_{\ell=1}^4(O^{-1}E)_{\ell j}
                                      (f_{\ell,+}-f_{\ell,-}).
 \tag{NW7}
\]
Independence follows from their even components $c_j$.
Their evaluation is zero, so they are annihilated by NW5.
Conversely positivity forces $d_v=0$, and the invertible
$B_\omega,O$ then give exactly NW6. This proves both
inclusions and the exact nullity four.

For every coefficient column $d$ there is a unique vector
orthogonal to $\mathcal N$ in either of the two original
norms and having that column:
\[
 \ell_\omega(d)=\iota_-O^{-1}B_\omega^{-1}d.
 \tag{NW8}
\]
To prove both minimization assertions rather than just a
section identity, every vector with this column is
$\ell_\omega(d)+\iota_+\alpha-\iota_-O^{-1}E\alpha$.
In NW3 its evaluation is the fixed $B_\omega^{-1}d$,
and the remaining squared norm is $\alpha^*G_X\alpha$
or $\alpha^*G_Y\alpha$, strictly minimized at $\alpha=0$.
The cross term vanishes by NW2, because the first component
of the section and the second component of the null vector
are both zero. Consequently the exact quotient norms are
\[
 \|[v]\|_{Q,W}^2=d_v^*H_\omega d_v=z^*Gz,
 \qquad
 \|[v]\|_{Q,U}^2=d_v^*H_{U,\omega}d_v=z^*G_Uz.
 \tag{NW9}
\]
The actual target polynomial of NW8 is
$C_{-\kappa}R_\omega d$; its actual source is
$T_A^{-1}C_{-\kappa}R_\omega d$, with FC24 coefficients
evaluated below. Thus both minimum lifts are fully typed
through the original conductor.

\section{Evaluated original target and source constants}
Retain $b=s/2>0$, $M=(2\pi)^{s/2}$ and the original
measure
\[
 dm_s(y)=\frac{M2^{s/2}}{4\pi\Gamma(s/2)}
                 |\Gamma(s/4+iy/2)|^2dy.
 \tag{NW10}
\]
Its moments are $M$, $Mb$, $M(3b^2+2b)$ and
$M(15b^3+30b^2+16b)$ in degrees $0,2,4,6$, with
odd moments zero \cite{WCF,Fable029}. For the target, all diagonal
entries of $H_\omega$ are
$3\int((3/2)^2+(y-\tau)^2)^3dm_s(y)$. Expanding with
these unchanged moments gives the following exact trace:
\[
\begin{aligned}
 t_W(\tau):=\tr H_\omega=12M\big[&\tau^6+(15b+27/4)\tau^4
 +(45b^2+141b/2+243/16)\tau^2\\
 &+15b^3+201b^2/4+715b/16+729/64\big]>0.
\end{aligned}\tag{NW11}
\]
Its exact determinant is
\[
 D_W:=\det H_\omega=12\cdot3^{16}M^4b^3(b+1)^2(b+2)>0.
 \tag{NW12}
\]
No test-pole height is omitted from NW11; the height cancels
from NW12 by the exact coefficient determinant proved in
RW11 and NH10.

Here is an equally explicit source trace, including every
original inverse coefficient. Let $v$ be the actual order
of the original conductor and $\mu_n=E_A^{(n)}(0)$,
$\mu_v\ne0$. Retain
\[
\begin{gathered}
 \nu_{v+j}=\sum_{r=0}^j\binom{v+j}{v+r}(-4)^{j-r}\mu_{v+r},
 \quad0\le j\le3,\\
 g_0=(-i)^v\mu_v/v!,\quad
 g_1=(-i)^{v+1}\nu_{v+1}/(v+1)!,\\
 g_2=(-i)^{v+2}\nu_{v+2}/(v+2)!-(v/3)g_0,\quad
 g_3=(-i)^{v+3}\nu_{v+3}/(v+3)!-((v+1)/3)g_1,\\
 \xi_0=g_0^{-1},\quad \xi_1=-g_1g_0^{-2},\quad
 \xi_2=g_1^2g_0^{-3}-g_2g_0^{-2},\quad
 \xi_3=-g_1^3g_0^{-4}+2g_1g_2g_0^{-3}-g_3g_0^{-2},\\
 k_j=(b)_{v+j}/(v+j)!,\quad j=0,1,2,3.
\end{gathered}\tag{NW13}
\]
These are the finite FC24 inverse coefficients, including
the original center separation $4$ and the original
$w(t)=-i\arctan t$ phase. Their derivation and multiplication
check are NH26--32. Write the full matrix
\[
 F=\begin{pmatrix}
 \xi_0&\xi_1&b\xi_0+2\xi_2&(3b+2)\xi_1+6\xi_3\\
 0&\xi_0&2\xi_1&(3b+2)\xi_0+6\xi_2\\
 0&0&2\xi_0&6\xi_1\\0&0&0&6\xi_0
 \end{pmatrix}.
 \tag{NW14}
\]
In the original minimum quotient orthonormal frame
$[\phi_v],\ldots,[\phi_{v+3}]$ at center $c_0=c_*+4$,
the inverse images of $1,y,y^2,y^3$, where
$y=(S'-c_*)/i$, have matrix
$\sqrt M\diag(\sqrt{k_j})F$. The actual polynomials are
specified by
$\phi_n(S)=p_{n,s}((S-c_0)/i)/\sqrt{Mn!(b)_n}$ and
$p_{n+1,s}=yp_{n,s}-n(n-1+b)p_{n-1,s}$,
$p_{0,s}=1,p_{1,s}=y$. These representatives contain no
component in the removed kernel $\mathcal P_{v-1}$, so
their Gram $K_U=MF^*\diag(k_j)F$ is the original minimum
quotient Gram. Its entries and all leading minors are
evaluated in NH33--34.

For complete trace evaluation define the following explicit
real symmetric polynomial matrix, with $d=3/2$:
\[
 Z(\tau)=4\begin{pmatrix}
 (d^2+\tau^2)^3&-3\tau(d^2+\tau^2)^2&-(d^2-3\tau^2)(d^2+\tau^2)&\tau(d^2-\tau^2)\\
 -3\tau(d^2+\tau^2)^2&(d^2+\tau^2)(5d^2+9\tau^2)&-\tau(7d^2+9\tau^2)&-(d^2-3\tau^2)\\
 -(d^2-3\tau^2)(d^2+\tau^2)&-\tau(7d^2+9\tau^2)&5d^2+9\tau^2&-3\tau\\
 \tau(d^2-\tau^2)&-(d^2-3\tau^2)&-3\tau&1
 \end{pmatrix}.
 \tag{NW15}
\]
Every entry comes from $Z=Q(\tau)Q(\tau)^*$, where the
column $j$ of $Q$ is the full coefficient vector of
$(d+i(y-\tau))^j(d-i(y-\tau))^{3-j}$. NH6 gives all
sixteen entries of $Q$; expanding their pairwise products
proves NW15 directly. The exact source trace and determinant
are therefore the finite original-moment expressions
\[
 \boxed{\quad t_U(\tau)=3M\sum_{r=0}^3k_r
               \sum_{p,q=0}^3F_{rp}\overline{F_{rq}}Z_{pq}(\tau),
 \qquad
 D_U=144\cdot3^{16}M^4k_0k_1k_2k_3|\xi_0|^8.\quad}
 \tag{NW16}
\]
All symbols on the right are evaluated in NW13--15;
there is no unspecified matrix norm. Indeed
$H_{U,\omega}=3Q^*K_UQ$, so its trace is the first
expression by multiplication. Its determinant is the
squared modulus of the exact triangular source coefficient
matrix, times $|\det(\sqrt3Q)|^2=3^{16}$, giving the
second expression. Both are strictly positive. Formula
NW16 retains every cross term and every inverse moment;
it does not replace $K_U$ by its diagonal.

\section{Five explicit unconditional native coercivity bounds}
The provider's analytic full-form inequality WP23 is
\[
 T_3(2)\succeq a_*I,\qquad
 a_*=\frac1{2\cdot200^4\cdot30^{12}}>0.
 \tag{NW17}
\]
Its proof includes the unverified high zeros by the signed
tail estimate WP20--22, and depends on the stated theorems
of Platt--Trudgian and Hasanalizade--Shen--Wong through
WP4--5. It is unconditional. The five exact-ball matrices
in WP26 give stronger numerical coefficient bounds that
can be extracted without rerunning the zeta calculation.
The following are strict outward rational enclosures of
the provider's stored real balls:
\[
\begin{array}{c|c|c|c}
 \tau&c_0<q_\tau&\det T_3(2+i\tau)>D_\tau&
       a_\tau=27D_\tau/(4q_\tau)^3\\\hline
 0&7/50&468/10^{17}&3159/4390400000000000\\
 14&38/25&967/10^8&26109/22475571200\\
 21&173/100&1299/10^5&175365/165686944\\
 100&251/100&9&3796875/15813251\\
 1000&127/25&556&58640625/32774128
\end{array}
\tag{NW18}
\]
All first four principal minors in each stored matrix were
proved positive by WP26, hence the four eigenvalues are
positive. Their sum is $4c_0<4q_\tau=:S_\tau$. If
$\lambda$ is the smallest, the product of the other three
is at most $(S_\tau/3)^3$ by the arithmetic--geometric
mean inequality. Thus
\[
 a_\tau I\preceq T_3(2+i\tau)\preceq S_\tau I.
 \tag{NW19}
\]
The strict inequalities from the balls imply these weak
matrix inequalities with the displayed exact rational
constants. At zero, both $a_*$ and $a_0$ are available;
$a_0>a_*$ is also an exact rational comparison.

Every positive four-by-four Hermitian matrix $H$ satisfies
\[
 H\preceq(\tr H)I,\qquad
 H\succeq\frac{27\det H}{(\tr H)^3}I.
 \tag{NW20}
\]
The first follows by summing its positive eigenvalues;
the second follows by the same product-of-three argument
just used. Apply NW20 separately to the actual $H_\omega$
and $H_{U,\omega}$, with their explicitly evaluated
traces and determinants. Combining with NW5, NW9 and
NW19 proves the complete two-sided quotient estimates:
\[
\boxed{\begin{aligned}
 \frac{a_\tau}{t_W(\tau)}\|[v]\|_{Q,W}^2
 &\le\mathbb W_{2+i\tau}(v,v)
 \le\frac{S_\tau t_W(\tau)^3}{27D_W}\|[v]\|_{Q,W}^2,\\
 \frac{a_\tau}{t_U(\tau)}\|[v]\|_{Q,U}^2
 &\le\mathbb W_{2+i\tau}(v,v)
 \le\frac{S_\tau t_U(\tau)^3}{27D_U}\|[v]\|_{Q,U}^2,
 \qquad\tau\in\{0,14,21,100,1000\}.
\end{aligned}}\tag{NW21}
\]
Every constant on the right and left is the finite expression
NW11--16 and NW18. Thus these are quantitative inequalities
in each original native norm, with no coefficient norm left
unconverted. At $\tau=0$ the entirely analytic lower
constants $a_*/t_W(0)$ and $a_*/t_U(0)$ follow from
NW17 by exactly the same proof, independently of the
matrix ball computation.

For completeness, a sharper direct native expression is
also exact at each certified pole:
\[
 \mathbb W_\omega(v,v)\ge
 \frac{\det T_3(\omega)}{\tr(\operatorname{adj}(T_3(\omega))H_J)}
                         \|[v]\|_{Q,J}^2,
 \qquad J=W\text{ or }U.
 \tag{NW22}
\]
The denominator is strictly positive. To prove it, let
$A=H_J^{1/2}T_3^{-1}H_J^{1/2}>0$. Then
$\lambda_{\max}(A)\le\tr A$
$=\tr(T_3^{-1}H_J)$; inversion gives NW22.
Its numerator and denominator are finite determinant and
cofactor expressions in the actual certified Toeplitz
entries and the full original matrices, so NW22 can be
enclosed directly without introducing a spectral norm
parameter. NW21 already supplies explicit rational and
original-moment bounds without requiring this refinement.

\section{The cubic Gaussian counterpart in both original norms}
The provider also proves the complete Gaussian inequality
WP14. To retain its exact meaning, put
\[
 \mathcal A(s)=\tfrac12s(s-1)\pi^{-s/2}\Gamma(s/2),\quad
 \phi_*(x)=(4\pi^2x^4-6\pi x^2)e^{-\pi x^2},\quad
 D_x=-x\,d/dx.
 \tag{NW23}
\]
For $r(s)=\sum_{j=0}^3r_js^j$, its actual test function
$r(D_x)\phi_*$ has Mellin transform $\mathcal A(s)r(s)$
on $\Re s>-2$, by WP1--2 including their boundary-term
proof. WP14 gives, for this full cubic source space,
\[
 \W(\mathcal A r,\mathcal A r)\ge e^{-300}\sum_j|r_j|^2.
 \tag{NW24}
\]
For our full eight-state evaluation, its exact polynomial
is $r(s)=(\Psi z)(s+\kappa)$ and its source is therefore
$[\Psi z](D_x+\kappa)\phi_*$. This retains the original
center translation and every entry of $\Psi$, as in WP15--16.

The coefficient conversion from $r(s)$ to its value on
$s=1/2+iy$ is exactly
\[
 D_G=\begin{pmatrix}
 1&1/2&1/4&1/8\\0&i&i&3i/4\\0&0&-1&-3/2\\0&0&0&-i
 \end{pmatrix},\qquad
 Z_G=D_GD_G^*=
 \begin{pmatrix}
 85/64&-27i/32&-7/16&i/8\\
 27i/32&41/16&-17i/8&-3/4\\
 -7/16&17i/8&13/4&-3i/2\\
 -i/8&-3/4&3i/2&1
 \end{pmatrix}.
 \tag{NW25}
\]
The original target monomial Gram in $y$ is
$K=(m_{p+q})_{0\le p,q\le3}$ with the moments NW10;
the source one is $K_U=MF^*\diag(k_j)F$ from NW14.
The Gaussian coefficient Grams in the two original
arithmetic norms are exactly $J_W=D_G^*KD_G$ and
$J_U=D_G^*K_UD_G$. Their evaluated traces are
\[
\boxed{\begin{aligned}
 t_{G,W}&=M(85/64+339b/16+141b^2/4+15b^3),\\
 t_{G,U}&=M\sum_{r=0}^3k_r\sum_{p,q=0}^3
                         F_{rp}\overline{F_{rq}}(Z_G)_{pq}.
\end{aligned}}\tag{NW26}
\]
For the first identity, integrate
$\sum_{j=0}^3|1/2+iy|^{2j}$ using the full moments NW10;
for the second, multiply the explicitly displayed matrices.
Both traces are strictly positive, and no Gaussian
$L^2$ norm has been substituted for the original conductor
norm. The two original centers have been transported
through the exact $s=S'-\kappa$ dictionary.

The Gaussian full eight-state form has exactly the same
four-dimensional kernel NW6: its coefficient receiver
$C_\kappa\Psi$ is invertible and NW24 is strictly positive.
For a Gaussian coefficient column $r$, the actual
orthogonal minimum lift in either original norm is
\[
 \ell_G(r)=\iota_-O^{-1}\Psi^{-1}C_{-\kappa}r.
 \tag{NW27}
\]
Its target is the polynomial $r(S'-\kappa)$ and its
source is $T_A^{-1}r(S'-\kappa)$, with inverse coefficients
NW13--14. The same minimization argument as NW8 proves
that its quotient squared norms are $r^*J_Wr$ and
$r^*J_Ur$. Combining NW24 with the trace upper bounds
therefore gives the two unconditional original-native
Gaussian inequalities on every full signed state:
\[
\boxed{\quad
 \W(\mathcal A r,\mathcal A r)
 \ge\frac{e^{-300}}{t_{G,W}}\|[v]\|_{Q,W}^2,
 \qquad
 \W(\mathcal A r,\mathcal A r)
 \ge\frac{e^{-300}}{t_{G,U}}\|[v]\|_{Q,U}^2.
 \quad}\tag{NW28}
\]
These apply to the same original quotient spaces as NW21,
through the displayed Gaussian coefficient receiver; the
rational and Gaussian coefficient columns are not
identified with one another.

\section{Exact source use and independent checks}
The file \texttt{verify\_native\_weil\_coercivity.py}
reads the existing provider certificate as immutable data.
Using exact rational arithmetic on each displayed ball's
decimal midpoint and radius, it proves the outward
enclosures NW18 and their rational $a_\tau$. It does not
rerun Arb, compute zeta zeros, or repeat the provider's
Taylor-series calculation. Separately, it symbolically
verifies every entry of NW15 and NW25, the fully evaluated
traces NW11 and NW26, and both source trace identities
from the complete FC24 matrix NW14. Its receipt is
\texttt{native\_weil\_coercivity\_exact.json}.
The provider's exact file identities are
\[
\begin{array}{ll}
 \texttt{WEIL\_POSITIVITY\_AND\_ARITHMETIC\_CERTIFICATES.tex}:&
 \texttt{ffe3f0cd195a3a2652f3d47540cd4ce9bfd8cb730cab29223d730b145a2ba73f},\\
 \texttt{ARITHMETIC\_BALL\_CERTIFICATES.json}:&
 \texttt{9c3ae36dc2f03fa42b2d54c8f5a30802fcc0f49b4a2bd9d99bc598b6151a12f7}.
\end{array}
\]
The four-zero proof uses Platt--Trudgian's verified height,
Hasanalizade--Shen--Wong's explicit zero count, Askey--Roy's
Gamma remainder, and Johansson's Arb arithmetic exactly
as cited in WP. The full-zero-sum coefficient identity
comes from MP, retaining the primary genus-one-product
input cited there to Lagarias. This note imports those
proved bounded results; it does not represent this
independent receiver calculation as a rereading of every
primary paper.

\section{Certified arithmetic eigenlines in the original target norm}
The five-pole calculation in \cite[WP24--27]{WeilCertificates}
encloses the complete Weil coefficient matrices. Here the receiving
calculation encloses all four generalized eigenvalues in the
unchanged original target norm, constructs their exact eigenlines
in the original eight signed coordinates, and proves that no
additional kernel is hidden at any of these poles.

Fix the original permitted Gamma parameter $s=1$, so
$b=1/2$, $M=\sqrt{2\pi}$, and retain the full original target
center $c_*$ through $S'=z+c_*-1/2$. At
\[
 \omega=2+i\tau,\qquad \tau\in\{0,14,21,100,1000\},\qquad
 h=3,\quad d=3/2,
 \tag{NV1}
\]
let $Q(\tau)$ be exactly NH6 and $K$ exactly NH7. The actual
positive target Gram and arithmetic Gram are
\[
 H_\omega=3Q(\tau)^*KQ(\tau),\quad
 T_\omega=(c_{j-i}(\omega))_{0\le i,j\le3},\quad
 \det H_\omega=12\cdot3^{16}M^4b^3(b+1)^2(b+2).
 \tag{NV2}
\]
Thus every original correlation and the original total mass remain.
The entries $c_j$ are the full zero sums RW12--13. Their
finite evaluation uses the proved exact generating identity
\[
 m_\omega(u)=\frac{\omega+(1-\bar\omega)u}{1+u},\qquad
 X_\omega(u)=\xi(m_\omega(u)),\qquad
 \frac{X_\omega'(u)}{X_\omega(u)m_\omega'(u)}
 =\frac{\xi'}{\xi}(m_\omega(u)).
 \tag{NV3}
\]
For $j=1,2,3$, $c_j$ is its coefficient of $u^j$; the
diagonal is twice the real part of its constant term.
The exact derivation, including the zero sum and all
Hadamard constants, is \cite[MP6--15]{MovableWeil}.
Both divisors in NV3 have nonzero constant term:
$\xi(\omega)\ne0$ on $\Re\omega>1$ and
$m_\omega'(0)=-3$. Formal series division is therefore
legitimate. The script retains eight coefficients before
differentiation, more than needed for every displayed entry.
No truncation of a zero, prime or Gamma sum is used to
replace these completed-function values.

Here is the complete real-root certification argument,
so a floating-point eigenvalue is not being used as a sign
test. For real $x$ set
\[
 A_\omega(x)=T_\omega-xH_\omega,\qquad
 \Delta_0(x)=1,\quad
 \Delta_j(x)=\det A_\omega(x)_{\{0,\ldots,j-1\}},
 \quad1\le j\le4.
 \tag{NV4}
\]
At each rational endpoint used below the certificate proves
all four $\Delta_j$ real and nonzero, with strict signs.
Gaussian elimination without a row exchange gives
$A=L\,\diag(p_1,\ldots,p_4)L^*$, where $L$ is invertible
unit lower triangular and $p_j=\Delta_j/\Delta_{j-1}$.
For completeness, at the first step $p_1=A_{00}\ne0$;
subtracting its column multiples yields its Schur complement.
Its first $r$ leading determinant is
$\Delta_{r+1}/p_1$, by the block determinant identity.
Induction proves every pivot and the factorization.
Hence the number of negative eigenvalues of $A$ equals
the number of sign changes in
$(1,\Delta_1,\ldots,\Delta_4)$, since multiplication by
$L^*$ is a bijection preserving the maximal dimensions
of positive and negative quadratic subspaces.

The positive matrix $H_\omega$ has an invertible positive
square root. Congruence by $H_\omega^{-1/2}$ identifies
$A_\omega(x)$ with $H_\omega^{-1/2}T_\omega H_\omega^{-1/2}-xI$.
Its negative inertia is therefore precisely the number
of original-native generalized eigenvalues less than $x$.
The certificate provides four disjoint rational intervals
$(l_j,u_j)$ with
\[
 0<l_1<u_1<l_2<u_2<l_3<u_3<l_4<u_4,\qquad
 \operatorname{neg}A_\omega(l_j)=j-1,\quad
 \operatorname{neg}A_\omega(u_j)=j.
 \tag{NV5}
\]
Every interval consequently contains exactly one eigenvalue
with multiplicity. There are four in total, so all have
been enclosed, are positive and are simple.
The exact rational endpoints, all forty endpoint inertias,
all $160$ leading-minor balls, and every original matrix
entry are retained in
\texttt{ORIGINAL\_NATIVE\_WEIL\_CERTIFICATES.json}.
Each interval has relative width less than $10^{-60}$.

The following rounded numbers only identify the certified
intervals; the rational endpoints and strict signs are
the evidence:
\[
\begin{array}{c|rrrr}
 \tau&\lambda_1&\lambda_2&\lambda_3&\lambda_4\\\hline
 0&8.14745\,10^{-11}&1.31906\,10^{-7}&1.90271\,10^{-5}&1.59784\,10^{-3}\\
 14&8.28289\,10^{-11}&3.32645\,10^{-7}&2.28517\,10^{-3}&1.07200\,10^{4}\\
 21&2.49471\,10^{-10}&8.59801\,10^{-7}&3.53539\,10^{-2}&1.19546\,10^{5}\\
 100&8.86137\,10^{-14}&5.05459\,10^{-7}&4.05208\,10^{1}&3.46066\,10^{8}\\
 1000&1.54757\,10^{-19}&3.26961\,10^{-8}&4.69894\,10^{3}&1.63176\,10^{15}
\end{array}.
\tag{NV6}
\]
The calculation uses python-flint 0.9.0, Arb at $768$ bits,
and one computational worker. The original interval-arithmetic
source is Johansson \cite{JohanssonArb}; exact algorithms,
script hash and complete enclosures are supplied. Every
determinant's imaginary ball contains zero, consistent
with the already proved exact Hermitian symmetry. A ball
containing zero in its real part never passes a sign test.

The eigenlines are also explicitly constructed, rather than
left as unnamed vectors. Let $\lambda_j$ be the unique
real root of $\det A_\omega(x)$ in the certified interval
NV5. For each one, the certificate selects an index
$k_j\in\{0,1,2,3\}$ and proves the principal cofactor
$(\operatorname{adj}A_\omega(x))_{k_jk_j}$ nonzero
throughout that entire interval. It evaluates the same
cofactor by interval arithmetic with an enclosure of
the whole rational interval. Consequently
\[
 d_j=\operatorname{adj}(T_\omega-\lambda_jH_\omega)e_{k_j}
 \ne0,\qquad
 T_\omega d_j=\lambda_jH_\omega d_j.
 \tag{NV7}
\]
Every coordinate of $d_j$ is the specified signed
three-by-three determinant, hence a polynomial of degree
at most three in the exactly defined $\lambda_j$.
The adjugate identity proves the eigenvector equation;
the certified diagonal cofactor proves nonvanishing.
Simplicity proves that its line is the complete eigenspace.

Retain the eight signed vectors and the exact factors $1/2$
in SE2. The original eight-state eigenvector and its two
physical components are therefore
\[
 v_j=\iota_-O^{-1}B_\omega^{-1}d_j
   =\frac12\sum_{\ell=1}^4
     (O^{-1}B_\omega^{-1}d_j)_\ell(f_{\ell,+}-f_{\ell,-}),
 \quad
 \mathcal J_Wv_j=(0,C_{-\kappa}R_\omega d_j),\quad
 \mathcal J_Uv_j=(0,T_A^{-1}C_{-\kappa}R_\omega d_j).
 \tag{NV8}
\]
No sign state is identified with another in this map.
RW15--20 show that $v_j$ is an eigenvector of the
actual target-native pencil
$\mathbb Wv=\lambda\mathbb H_Wv$.
The four remaining eigenvectors have eigenvalue zero and are
\[
 g_\ell=\iota_+e_\ell-\iota_-O^{-1}Ee_\ell,\qquad
 1\le\ell\le4.
 \tag{NV9}
\]
Indeed their evaluations vanish and their label coordinates
are independent; NV5 excludes any further kernel.
For $j\ne k$, Hermitian symmetry in NV7 gives
$(\lambda_j-\lambda_k)d_k^*H_\omega d_j=0$, proving
orthogonality in the original target norm. All $g_\ell$
are orthogonal to all $v_j$ in that same norm by RW16.
The four graph vectors retain their full label Gram $G_Y$,
so no Euclidean label norm is silently substituted.

The numerical eigenvalues in NV6 concern the actual target
metric with the stated permitted $s=1$. The original
source metric is separately and completely evaluated in
NH26--42 and the native coercivity calculation. Its
conductor moments are not assigned numerical values here.
The symbolic exact source vector in NV8 nevertheless exists
for every original conductor on its established domain,
using the actual FC24 inverse. Thus each certified
arithmetic eigenline has a fully constructed return
through the original Fable and conductor maps.


\section{The finite completion and the exact residue return}
This note proves the finite completion proposed in EA1--28 of
\cite{ExceptionalAtlas}, retaining the original reciprocal factor
$t=1/a$, both infinity signs and all derivative constants.
The original polynomial and marked states remain attributed
to \cite{ES488}; the ordinary-fibre calculations and existing
nonfinite factor completion are recorded in \cite{Fable029}.
The finite completion below is a different, explicitly related
object. All schemes and algebras in this note are over $\C$.

\section{The original map and two completed root charts}
The original source coordinates are $(a,y,z,w)$. Define
\begin{align}
b&=i+ay,&c&=-i+2ay+a^2z,\nonumber\\
d&=-iy-a(iz+2y^2)-a^2yz,&e&=2z-7iy^2+aw,\nonumber\\
f&=iw+3iy^3-4yz+a(6iy^2z+wy+4y^4)+2a^2y^3z,\nonumber\\
P&=(ac,ae+bd,af+be,bf)=(A,B,C,D).
\tag{FS1}
\end{align}
Direct multiplication gives $ad+bc=1$ and
$a^3f-a^2be+ab^2d-b^3c=1$. Thus
$(aU+bV)(cU^3+dU^2V+eUV^2+fV^3)=H_u(U,V)$, where
\[
H_u(U,V)=AU^4+U^3V+BU^2V^2+CUV^3+DV^4,\quad
h_u(r)=Ar^4+r^3+Br^2+Cr+D,\quad
k_u(v)=A+v+Bv^2+Cv^3+Dv^4 .
\]
The coefficient of $U^3V$ remains one throughout.
For $a\ne0$, division by the linear factor gives
$r=-b/a=-y-i/a$ and $a^2h_u'(r)=1$.
Keeping the reciprocal factor as the actual new coordinate,
define
\begin{equation}
X_f=\Spec\frac{\C[A,B,C,D,r,t]}{(h_u(r),\,t^2-h_u'(r))},
\qquad
X_\infty=\Spec\frac{\C[A,B,C,D,v,\theta]}
 {(k_u(v),\,\theta^2+k_u'(v))}.
\tag{FS2}
\end{equation}
On $r\ne0$ and $v\ne0$ the exact gluing maps and inverses are
\begin{equation}
v=r^{-1},\quad\theta=-t/r,\qquad
r=v^{-1},\quad t=-\theta/v .
\tag{FS3}
\end{equation}
At a root, differentiating $k(v)=v^4h(1/v)$ gives
$k'(v)=-v^2h'(1/v)$. Hence FS3 preserves both equations
and signs, and direct substitution proves the inverse
compositions. Glue these charts to form $X$ with projection
$\pi:X\to\mathbb A^4_u$. No coordinate $t$ or $\theta$ is
inverted in defining the completed charts.

The equations on the finite chart solve as
\begin{equation}
C=t^2-4Ar^3-3r^2-2Br,\qquad
D=3Ar^4+2r^3+Br^2-rt^2.
\tag{FS4}
\end{equation}
Substitution in both directions proves
$X_f=\mathbb A^4_{t,r,A,B}$, including $A=0$ and $t=0$.
At $v=0$, the other chart has exactly
\begin{equation}
A=0,\quad\theta=\pm i,\qquad
D_\pm=\{v=0,\ \theta=\mp i\}\simeq\mathbb A^3_{B,C,D}.
\tag{FS5}
\end{equation}
Thus $D_+$ is the included original sign $\theta=-i$ and
$D_-$ is the missing original sign $\theta=i$.
Eliminating $A$ on $X_\infty$ leaves the equation
$1+\theta^2+2Bv+3Cv^2+4Dv^3=0$.
At $v=0$ its derivative in $\theta$ is $2\theta\ne0$.
Therefore $X$ is smooth there, and it is smooth elsewhere
by FS3--4.

\section{Five explicit finite free descriptions over the target}
For $p\in\{0,1,-1,2\}$ write $c_p=H_u(p,1)=h_u(p)$.
The following identity is literal coefficient arithmetic:
\begin{equation}
1=\tfrac12h_u(0)-\tfrac12h_u(1)-\tfrac16h_u(-1)
                       +\tfrac16h_u(2)-2A .
\tag{FS6}
\end{equation}
It is verified by comparing each coefficient of $A,B,C,D$
and the fixed cubic term. Consequently the five principal
opens $D(A),D(c_0),D(c_1),D(c_{-1}),D(c_2)$ cover the
entire characteristic-zero target.

On $D(c_p)$ use the determinant-one substitution
$(U,V)=(pU'-V',U')$ and put
$h_p(R)=H_u(pR-1,R)$. Its leading coefficient is precisely
the unit $c_p$. The complete root/sign transitions are
\begin{align}
R&=(p-r)^{-1},&T&=t/(p-r)
                    &&\text{on the finite overlap},\nonumber\\
R&=-v/(1-pv),&T&=\theta/(1-pv)
                    &&\text{on the infinity overlap},\nonumber\\
r&=p-R^{-1},&t&=T/R&& (R\ne0),\nonumber\\
v&=R/(pR-1),&\theta&=T/(1-pR)&& (pR-1\ne0).
\tag{FS7}
\end{align}
Every denominator is an explicitly stated unit on the relevant
overlap. On $D(c_p)$ the marked root cannot equal $p$;
thus $p-r$ is a unit on its finite root scheme. More
algebraically $h_u(p)-h_u(r)$ is divisible by $p-r$,
and $h_u(r)=0$, so the unit $c_p$ proves this assertion.
Likewise the infinity denominator $1-pv$ cannot vanish
there. The inverse opens cover the new chart because $R$
and $pR-1$ generate the unit ideal.

The derivative identity is
$h_p'(R)=R^2h_u'(p-R^{-1})$ at a root, obtained by
differentiating $h_p(R)=R^4h_u(p-R^{-1})$.
Thus $T^2=h_p'(R)$ on both overlaps; at $v=0$, $R=0$
and $T=\theta$, so the infinity sign is preserved exactly.
FS7 and its inverse formulas give an isomorphism of the
whole preimage over $D(c_p)$ with
\begin{equation}
\Spec\,\mathcal O_{D(c_p)}[R,T]/
                 (h_p(R),\,T^2-h_p'(R)).
\tag{FS8}
\end{equation}
The same description on $D(A)$ is the original finite
chart FS2 with the additional unit $A$; all marked roots
there are finite. In either case, writing $\lambda$ for
the retained leading coefficient, division by the unit
$\lambda$ makes only the root relation monic.
The sign relation is still exactly $T^2=h_p'(R)$,
not the derivative of the divided polynomial.
The successive monic relations give the free module basis
\begin{equation}
1,R,R^2,R^3,\ T,TR,TR^2,TR^3 .
\tag{FS9}
\end{equation}
Indeed the root quotient is free on the first four monomials;
the monic quadratic relation in $T$ makes its extension
free on $1,T$. This proves on all five opens, without
any restriction on their discriminants, that
\begin{equation}
\pi:X\longrightarrow\mathbb A^4_u
\text{ is finite locally free of rank eight.}
\tag{FS10}
\end{equation}

The finite chart is irreducible. It is dense in $X$:
on the infinity chart the remaining equation is monic
in $\theta$, so its coordinate algebra is free over
$\C[v,B,C,D]$ and multiplication by $v$ is injective.
Thus no irreducible component is supported on $v=0$;
every such point is also smooth by FS5.
It follows that $X$ is irreducible and smooth, hence normal.
Since FS10 has rank eight and is dominant, its function
field is a degree-eight extension of $\C(A,B,C,D)$.
Finally, a finite normal integral scheme over an affine
base is the integral closure of that base in its function
field: any element integral over the base is also integral
over its finite coordinate ring, so normality includes it.
This proves that $X$ is the finite normalization in the
actual original degree-eight function-field extension.

\section{The original open immersion, including its infinity sheet}
The full resultant-one factor space is
$\mathcal F=\{ad+bc=1,\ a^3f-a^2be+ab^2d-b^3c=1\}$.
Its exact identification with the unramified part of $X$
is given on the finite chart by
\begin{align}
a&=t^{-1},&b&=-r/t,&c&=At,&d&=(1+Ar)t,\nonumber\\
e&=(B+r+Ar^2)t,&f&=(C+Br+r^2+Ar^3)t&& (t\ne0),
\tag{FS11}
\end{align}
and on the infinity chart by
\begin{align}
a&=-v/\theta,&b&=\theta^{-1},&
c&=\theta(1+Bv+Cv^2+Dv^3),\nonumber\\
d&=\theta(B+Cv+Dv^2),&
e&=\theta(C+Dv),&f&=\theta D
\qquad(\theta\ne0).
\tag{FS12}
\end{align}
Multiplication of the factors gives $H_u$ in either formula.
The resultant is $t^{-2}h'_u(r)=1$ in FS11 and
$-\theta^{-2}k'_u(v)=1$ in FS12.
Their inverse is $(r,t)=(-b/a,1/a)$ when $a\ne0$ and
$(v,\theta)=(-a/b,1/b)$ when $b\ne0$.
Since $ad+bc=1$, these inverse charts cover $\mathcal F$.
Direct substitution and FS3 prove every inverse composition.
The excluded locus is the reduced divisor
$R_0=\{t=0\}\simeq\mathbb A^3_{r,A,B}$; on the infinity
chart it is $\theta=0$, which never meets $v=0$.
Thus $\mathcal F\simeq X\setminus R_0$ exactly.

The original source maps to $X$ via FS1 and these inverse
factor charts. For $a\ne0$ its complete inverse is
\begin{equation}
q(t,r,A,B)=\left(t^{-1},-r-it,\,
At^3+2rt+3it^2,\,
7ir^2t+(B-17r+Ar^2)t^2-13it^3-2At^4\right).
\tag{FS13}
\end{equation}
All terms come from the successive original factor equations;
substitution verifies $Pq=\pi$.
At $a=0$ its factors have $b=i$, hence $\theta=-i$.
It has exactly the original state
$q_\infty=(0,B,i(7B^2-C)/2,11B^3-2BC-D)$ over
every $A=0$ target.

For an explicit inverse on an open neighborhood of that
entire infinity sheet, use $b(b+i)\ne0$ in $\mathcal F$
and set
\begin{align}
Y&=\frac{a^2f-abe+2b^2d}{b+i},\nonumber\\
Z&=\frac{af-be+7Y^2-10iaY^3-4a^2Y^4}{2b^3},\nonumber\\
W&=\frac{a^2eY^3+7ia^2Y^5+3iaeY^2-17aY^4
                       -2eY-f-11iY^3}{b^3}.
\tag{FS14}
\end{align}
No factor $a$ is divided out in FS14. The original two
factor equations give
$b^2+1=a(a^2f-abe+2b^2d)$, so $b-i=aY$.
Substitution of FS1 gives the three polynomial identities
\[
a^2f-abe+2b^2d=(b+i)y,\qquad
af-be+7y^2-10iay^3-4a^2y^4=2b^3z,
\]
\[
a^2ey^3+7ia^2y^5+3iaey^2-17ay^4-2ey-f-11iy^3=b^3w.
\]
These identities follow by expanding the displayed factor
polynomials and prove that FS14 is the inverse on the
dense subopen $a\ne0$. Both it and FS1 are regular on
the whole stated open set, which is reduced as an open
subset of the smooth irreducible $X$. Equality of the
compositions there follows from equality on that dense
open. This proves the inverse also at $a=0$.
The opens $a\ne0$ and $b(b+i)\ne0$ cover precisely the
original image in $\mathcal F$: at $a=0$ the resultant
forces $b=\pm i$, and the omitted case is $b=-i$.
We have therefore proved, with complete coordinate maps,
\begin{equation}
\mathbb A^4_q\simeq X\setminus(R_0\cup D_-),\qquad
R_0\cap D_-=\varnothing,\qquad
R_0\simeq D_-\simeq\mathbb A^3 .
\tag{FS15}
\end{equation}
It is an open immersion, rather than an assertion that
$t=1/a$ is an invertible change at its new boundary.

\section{The different, branch divisor and full trace discriminant}
In the original column order $(t,r,A,B)$ and target order
$(A,B,C,D)$, differentiating FS4 and FS13 gives
\begin{equation}
\det D\pi=-2t^3,\qquad
\det Dq=t^3\quad(t\ne0),\qquad
\det DP=-2.
\tag{FS16}
\end{equation}
For the first determinant, the remaining $(t,r)$ minor
after the $(A,B)$ rows is
$\begin{psmallmatrix}2t&-\kappa\\-2rt&r\kappa-t^2\end{psmallmatrix}$,
where $\kappa=12Ar^2+6r+2B$; it has determinant $-2t^3$.
For the second, the triangular pivots of the source
parameterization are $-t^{-2},-1,t^3,t^2$, whose product
is $t^3$. The chain rule proves the last identity where
$a\ne0$, and polynomial equality extends it to all original
source points.

The reduced ramification support is $R_0$; the actual
different (or ramification Cartier divisor) is
\begin{equation}
\operatorname{Fitt}_0\Omega_{X/\mathbb A^4}
=(t^3)\quad\hbox{on }X_f,\qquad
\operatorname{Ram}(\pi)=3R_0 .
\tag{FS17}
\end{equation}
Indeed both source and target are smooth of dimension four,
so their differential presentation has determinant FS16.
Alternatively the relations in FS2 have relative Jacobian
in variables $(r,t)$ equal to
$\begin{psmallmatrix}h'(r)&0\\-h''(r)&2t\end{psmallmatrix}$,
whose determinant is $2t^3$ on $X_f$.
At an infinity point, the relative derivative matrix of
$(k(v),\theta^2+k'(v))$ in $(v,\theta)$ has determinant
$2\theta k'(0)=2\theta\ne0$.
Thus no ramification component is hidden at infinity.
The changes FS7 transport the determinant by units and
give FS17 globally.

Define the exact binary quartic discriminant $\mathfrak D$
by its polynomial extension from
$\mathfrak D=A^6\prod_{i<j}(r_j-r_i)^2$ for $A\ne0$.
At $A=0$ its specialization is exactly
\begin{equation}
\mathfrak D(0,B,C,D)
=B^2C^2-4C^3-4B^3D+18BCD-27D^2.
\tag{FS18}
\end{equation}
One way to verify both the polynomial extension and this
constant is the determinant of the Sylvester matrix of
$h,h'$ divided by $A$ on $A\ne0$, then expansion and
evaluation at zero. FS18 is the ordinary discriminant of
the monic cubic; the projective infinity root is simple.
Thus the reduced branch locus is exactly
$\{\mathfrak D=0\}$, the image of $R_0$.
This also identifies the locus where every root chart of
the completed cover is etale; no $A=0$ component is added.

In a free root/sign chart with polynomial leading coefficient
$\lambda$ and basis FS9, let $M_R$ be multiplication by $R$
in the degree-four root algebra. Define
$T_{ij}=\Tr(M_R^{i+j})$ and
$T'_{ij}=\Tr(M_{h_p'}M_R^{i+j})$, $0\le i,j\le3$.
The full bilinear trace Gram in the unaltered basis is
\begin{equation}
\left(\Tr_{X/u}(e_ie_j)\right)_{i,j=1}^{8}
=\begin{pmatrix}2T&0\\0&2T'\end{pmatrix},\qquad
\det\left(\Tr_{X/u}(e_ie_j)\right)
=256\,\frac{\mathfrak D^3}{\lambda^{14}} .
\tag{FS19}
\end{equation}
The algebra is the direct sum of its root module and $T$
times that module. Multiplication by a root-algebra element
acts identically on these two summands, so its trace is twice
the root trace. Odd multiplication has trace zero (also by
conjugating with the sign involution), and the product of
two odd factors is exactly $h_p'$ times their root product.
This proves the block formula even on nonreduced fibres.

For distinct roots, write
$\Delta=\prod_{i<j}(R_j-R_i)$ and
$V=(R_j^i)_{0\le i\le3,1\le j\le4}$.
Then $T=VV^{\mathsf T}$ and
$T'=V\diag(h_p'(R_j))V^{\mathsf T}$, so
$\det T=\Delta^2$, $\det T'=\lambda^4\Delta^4$:
the six pair-sign reversals give
$\prod_j h_p'(R_j)=\lambda^4\Delta^2$.
The determinant of the full blocks is consequently
$2^8\lambda^4\Delta^6
=256(\lambda^6\Delta^2)^3/\lambda^{14}$, as asserted.
Both sides are regular coefficient functions on this free
chart; equality on the dense distinct-root locus proves
equality on every fibre.

The binary discriminant remains unchanged under the
determinant-one substitutions used above. Explicitly,
$c_p=A\prod_j(p-r_j)$ and
$R_j-R_i=(r_j-r_i)/[(p-r_j)(p-r_i)]$, so the factors
in $c_p^6\prod_{i<j}(R_j-R_i)^2$ cancel exactly to
$A^6\prod_{i<j}(r_j-r_i)^2$. The identity extends to all
coefficient points. Thus FS19 on $D(A)$ has denominator
$A^{14}$, and on $D(c_p)$ it has denominator $c_p^{14}$.
Since those denominators are units only on their explicitly
stated opens, the complete global discriminant ideal is
\begin{equation}
\operatorname{disc}(\pi)=(\mathfrak D^3).
\tag{FS20}
\end{equation}
In particular $A^{-14}$ in the first basis is not an
additional ramification divisor. At $A=0$ choose any
$p\in\{0,1,-1,2\}$ with $c_p\ne0$, which FS6 guarantees;
FS19 then evaluates the exact trace determinant with its
full constant $256$ and no singular denominator.

\section{Every nonreduced fibre in the original reciprocal coordinate}
Fix a target $u$. At a finite root $r_0$ of exact multiplicity
$m$ put $\epsilon=r-r_0$ and factor the actual polynomial as
$h_u(r_0+\epsilon)=\epsilon^m g(\epsilon)$ with $g(0)\ne0$.
The root algebra local at $r_0$ is exactly
$\C[\epsilon]/(\epsilon^m)$: all other root factors are
units there. In this algebra
$h'_u(r_0+\epsilon)=m g(0)\epsilon^{m-1}$.
Thus the completed local factor, with its original $t$,
is
\begin{equation}
\mathcal A_{m,c}
=\C[\epsilon,t]/(\epsilon^m,\ t^2-c\epsilon^{m-1}),
\qquad c=m g(0)\ne0 .
\tag{FS21}
\end{equation}
No leading coefficient has been absorbed into $t$.
The two successive monic relations give the exact basis
and multiplication
\begin{align}
\mathcal B_m&=(1,\epsilon,\ldots,\epsilon^{m-1},
              t,t\epsilon,\ldots,t\epsilon^{m-1}),\nonumber\\
\epsilon^i\epsilon^j&=\epsilon^{i+j},&
\epsilon^i(t\epsilon^j)&=t\epsilon^{i+j},&
(t\epsilon^i)(t\epsilon^j)&=c\epsilon^{m-1+i+j},
\tag{FS22}
\end{align}
with every exponent at least $m$ interpreted as zero.
These formulas specify every structure constant and prove
the length $2m$, including all nilpotents.
For $m=1$ the algebra is
$\C[t]/(t^2-c)=\C\oplus\C$, its two points being
$t=\pm\sqrt c$; these are the two original factor signs.
For $m=2$, the exact isomorphism without changing $t$ is
\begin{equation}
\mathcal A_{2,c}\simeq\C[t]/(t^4),\qquad
\epsilon\longmapsto t^2/c,\qquad
t\longmapsto t .
\tag{FS23}
\end{equation}
The inverse sends the class of $t$ to the actual reciprocal
factor. Its square and fourth power satisfy the two original
relations, proving both compositions.
For $m=3$ and $m=4$ the full original algebras remain
$\C[\epsilon,t]/(\epsilon^3,t^2-c\epsilon^2)$ and
$\C[\epsilon,t]/(\epsilon^4,t^2-c\epsilon^3)$.
If one wishes only to compare their abstract isomorphism
types, choose an actual square root $\lambda^2=c$ and send
the new symbol $s$ to $t/\lambda$; the inverse sends $t$
to $\lambda s$. This records the precise comparison to
the coefficient-one presentation and both choices of
$\lambda$; it does not replace the retained coordinate.

For completeness the local maximal-ideal filtrations are
also exact. For $m=2$ they are those of $(t)$ in
$\C[t]/t^4$, with dimensions $4,3,2,1,0$.
For $m=3,4$, put $\mathfrak n=(\epsilon,t)$. Direct
multiplication using FS22 gives
\begin{equation}
\begin{array}{c|c|c|c}
m&(\dim\mathcal A,\dim\mathfrak n,\dim\mathfrak n^2,\ldots)
 &\text{associated graded dimensions}&\operatorname{Soc}(\mathcal A)\\\hline
2&(4,3,2,1,0)&(1,1,1,1)&\C\,t\epsilon\\
3&(6,5,3,1,0)&(1,2,2,1)&\C\,t\epsilon^2\\
4&(8,7,5,3,1,0)&(1,2,2,2,1)&\C\,t\epsilon^3
\end{array}
\tag{FS24}
\end{equation}
For example at $m=4$ the bases of the successive nonzero
powers from the second onward are
$(\epsilon^2,\epsilon^3,t\epsilon,t\epsilon^2,t\epsilon^3)$,
$(\epsilon^3,t\epsilon^2,t\epsilon^3)$ and $(t\epsilon^3)$.
At $m=3$ they are $(\epsilon^2,t\epsilon,t\epsilon^2)$
and $(t\epsilon^2)$. The socle statements follow by
annihilating an arbitrary FS22 basis combination first by
$\epsilon$ and then by $t$.
Since multiplication by a nilpotent has trace zero,
for $m\ge2$ the local trace pairing is exactly
\begin{equation}
\Tr_{\mathcal A_{m,c}}(uv)=2m\,u(0,0)v(0,0),\qquad
\operatorname{rad}\Tr=\mathfrak n,\qquad
\operatorname{rank}\Tr=1 .
\tag{FS25}
\end{equation}
To see the trace identity directly, write a local element
as its constant plus a nilpotent in $\mathfrak n$; its
multiplication operator has that constant on all $2m$
diagonal entries after ordering any compatible filtration.
The reduced $m=1$ factor has its two nondegenerate scalar
trace pairings. Thus FS25 also specifies the complete
degenerate trace forms, not just their vanishing determinants.

The constants in FS21 have the following original-coordinate
expressions, with no unspecified unit:
\begin{equation}
c=\frac{h_u^{(m)}(r_0)}{(m-1)!},\qquad
c_2=12Ar_0^2+6r_0+2B,\quad
c_3=12Ar_0+3,\quad c_4=4A .
\tag{FS26}
\end{equation}
The expression for $m=1$ is just $c_1=h_u'(r_0)$.
Exact multiplicity means the stated constant is nonzero.
For example a cubic triple root has $c_3=3$.
The quadruple-root algebra is not $\C[t]/t^8$:
FS24 shows its embedding dimension is two, and FS22 gives
$t^4=c_4^2\epsilon^6=0$. Its length eight is retained by
the two generators and their full multiplication, not by
declaring the reciprocal factor to have nilpotency order eight.

The entire fibre is the product of these local factors,
by the Chinese remainder theorem in its root algebra.
When $A=0$ the finite-root part has total length six,
and the two infinity points FS5 contribute two reduced
copies of $\C$. The complete list is therefore
\begin{equation}
\begin{array}{c|c|c|c|c}
\text{degree}&\text{root multiplicities}&\text{full completed algebra}
 &\#\text{original points}&\operatorname{rank}\Tr\\\hline
4&1+1+1+1&\C^8&8&8\\
4&2+1+1&\mathcal A_{2,c}\times\C^4&4&5\\
4&2+2&\mathcal A_{2,c}\times\mathcal A_{2,c'}&0&2\\
4&3+1&\mathcal A_{3,c}\times\C^2&2&3\\
4&4&\mathcal A_{4,c}&0&1\\\hline
3&1+1+1&\C^6\times\C^2_\infty&7&8\\
3&2+1&\mathcal A_{2,c}\times\C^2\times\C^2_\infty&3&5\\
3&3&\mathcal A_{3,3}\times\C^2_\infty&1&3
\end{array}
\tag{FS27}
\end{equation}
Every algebra has length eight. Each displayed local
coefficient is the actual FS26 coefficient at that root;
in the $2+2$ case $c$ and $c'$ refer to its two different
roots and are both retained. The original open removes
the entire local factor of each multiple root, because
its $t$ is nilpotent and its principal open $D(t)$ is
empty. Of the two infinity factors it removes only
$\theta=i$. This proves every original count in FS27.
The trace ranks follow by summing FS25 and the simple
scalar trace ranks. Thus the table records both the
geometric points and their nonreduced local schemes.

\section{The omitted surface, ramification parametrization and deck map}
A quartic has no simple root precisely in the two
patterns $2+2$ and $4$. Its leading coefficient is
nonzero. With the unchanged cubic coefficient one,
being a scalar times a quadratic square is equivalent to
\begin{equation}
4AB-1-8A^2C=0,\qquad D-AC^2=0,\qquad
h_u(s)=A\left(s^2+\frac{s}{2A}+C\right)^2 .
\tag{FS28}
\end{equation}
Indeed write the monic quadratic as $s^2+ps+q$;
comparison of its cubic and linear coefficients gives
$p=1/(2A)$, $q=C$, and the other two give FS28.
The maps
$(A,C)\mapsto(A,(4A)^{-1}+2AC,C,AC^2)$ and
$(A,B,C,D)\mapsto(A,C)$ are inverse, with $A\ne0$.
This proves that the original omitted image is exactly
the smooth surface $\C^*\times\C$.
Its quadratic has a double root exactly when
$C=1/(16A^2)$. Thus its full quadruple curve is
\begin{equation}
B=\frac3{8A},\qquad C=\frac1{16A^2},\qquad
D=\frac1{256A^3},\qquad r_0=-\frac1{4A},\qquad A\ne0 .
\tag{FS29}
\end{equation}
Outside this curve FS28 has two length-four completed
points; on the curve it has the full two-generator
length-eight algebra in FS27.

On $R_0$ the projection is exactly FS4 at $t=0$.
Its differential in coordinates $(A,B,r)$ has its first
two rows equal to $(1,0,0)$ and $(0,1,0)$; the third
column in full target coordinates is
$(0,0,-\kappa,r\kappa)^{\mathsf T}$,
$\kappa=12Ar^2+6r+2B$.
It therefore has rank three exactly at roots of
multiplicity two, and rank two at roots of multiplicity
at least three. This proves the precise rank drop,
including cubic triples, without substituting the name
of the discriminant for its parametrization.

Both signs extend on the whole completed algebra as the
regular deck involution
\begin{equation}
(r,t)\mapsto(r,-t),\qquad
(v,\theta)\mapsto(v,-\theta).
\tag{FS30}
\end{equation}
FS3 makes the two formulas agree. Its fixed scheme is
exactly $R_0$, since characteristic zero makes $2$ a unit;
it exchanges $D_+$ and $D_-$. Its invariant algebra
in FS9 is precisely the degree-four root algebra, because
every element is uniquely $f(R)+Tg(R)$.
Thus its quotient is the universal projective root scheme
$\{H_u(U,V)=0\}\subset\mathbb A^4_u\times\mathbb P^1$,
including every nonreduced root scheme. The odd part
continues to multiply by the exact rule $T^2=h_p'(R)$;
it is not discarded at a fixed point.

Finally the exact nonproper-value locus of the original
affine map is
\begin{equation}
\mathcal N(P)=\pi(R_0\cup D_-)
            =\{\mathfrak D=0\}\cup\{A=0\}.
\tag{FS31}
\end{equation}
To prove the first containment, a finite morphism has
compact inverse images of compact complex sets: its
affine coordinate generators satisfy finitely many
monic relations, whose coefficient bounds bound every
root over a compact target. If a target lies outside
the images of the two removed closed divisors, a small
compact neighborhood has inverse image disjoint from
them. The regular inverse FS13--14 is bounded on that
compact preimage, excluding an escaping original
sequence there.

For the reverse containment, at any finite repeated root
fix $(r,A,B)$ and take $t=\varepsilon\ne0$ in FS4.
The target converges to the prescribed repeated-root
target, while FS13 has $a=1/\varepsilon\to\infty$.
For any $A=0$ target let
$q_\varepsilon=(\varepsilon,B,i(7B^2-C)/2,11B^3-2BC-D)$.
It tends to the included original point. The simultaneous
factor sign in the original coordinates for $a\ne0$ is
\[
\Sigma(a,y,z,w)=
(-a,\ y+2i/a,\ 6i/a^2-z,\
 w-14iy^2/a+28y/a^2+40i/a^3).
\]
Substitution into FS1 negates every factor and therefore
preserves $P$. The second coordinate of
$\Sigma q_\varepsilon$ escapes as $2i/\varepsilon$,
while its image tends to the stated $A=0$ target.
This proves every point of the second divisor image is
nonproper, including its discriminant intersections.

The exact verifier
\texttt{verify\_finite\_signed\_completion.py} checks the
five-open identity, all coordinate-transition signs,
both Jacobians, the original-chart inverse identities,
the trace-discriminant constants and every local
algebra filtration and trace pairing in FS21--27.
No numerical fibre sampling replaces those algebra identities.

The received exceptional atlas \cite{ExceptionalAtlas}, EA29--36, specifies
constant-term unfoldings and claims their original-metric singular
asymptotes. This note verifies those claims, derives the complete
forward and inverse compound limits with their complex phases,
evaluates every multiplicity-dependent constant, and determines
the original singular flags. The polynomial is the retained
four-dimensional ES construction \cite{ES488}; its existing complete
inverse is GF7. Every metric below is the fixed original Gamma
metric transported by the existing \(\Psi\), with no change to
the period or the conductor as the unfolding parameter varies.

\section{The entire original differential and its inverse}
Use the original source coordinates \((a,y,z,w)\), and output
order \((A,B,C,D)\). The polynomial is
\[
\begin{aligned}
P_0&=a^3z+2a^2y-ia,\\
P_2&=-a^3y^2z-2ia^2yz+a^2w-2a^2y^3-10iay^2+3az+y,\\
P_3&=2a^3y^3z+6ia^2y^2z+2a^2wy+4a^2y^4
       +2iaw-4iay^3-2ayz+2iz+7y^2,\\
P_4&=2a^3y^4z+8ia^2y^3z+a^2wy^2+4a^2y^5
       +2iawy+7iay^4-10ay^2z-4iyz-w-3y^3.
\end{aligned}\tag{FD1}
\]
For \(t\ne0\), put
\[
\begin{aligned}
q(t,r,A,B)=\bigl(&t^{-1},-r-it,At^3+2rt+3it^2,\\
 &7ir^2t+(B-17r+Ar^2)t^2-13it^3-2At^4\bigr)^T,\\
\pi(t,r,A,B)=\bigl(&A,B,t^2-4Ar^3-3r^2-2Br,\,
                      3Ar^4+2r^3+Br^2-rt^2\bigr)^T.
\end{aligned}\tag{FD2}
\]
Substitution into FD1 proves \(Pq=\pi\). This is the complete
finite simple-root inverse, with \(t=1/a\), rather than a
replacement source coordinate at \(t=0\). With
\[
 h_u(s)=As^4+s^3+Bs^2+Cs+D,\quad
 \ell_r=(r^4,r^2,r,1),\quad \ell'_r=(4r^3,2r,1,0),\quad
 \kappa=12Ar^2+6r+2B,
 \tag{FD3}
\]
its equations are \(h_u(r)=0\), \(t^2=h_u'(r)\).
Differentiating these two equations in all four output directions
gives exactly
\[
 dr=-\ell_r\,du/t^2,\qquad
 dt=\ell'_r\,du/(2t)-\kappa\ell_r\,du/(2t^3).
 \tag{FD4}
\]
Here and throughout, \(\ell_r\) is the complex-linear row
functional with the displayed coefficients, not their conjugates.
The differential of the original coordinate vector in FD2 is
\[
\begin{aligned}
 da&=-dt/t^2,&dy&=-dr-i\,dt,\\
 dz&=t^3dA+2t\,dr+(3At^2+2r+6it)dt,\\
 dw&=(r^2t^2-2t^4)dA+t^2dB
       +[14irt+(-17+2Ar)t^2]dr\\
   &\hspace{13mm}
       +[7ir^2+2(B-17r+Ar^2)t-39it^2-8At^3]dt.
\end{aligned}\tag{FD5}
\]
FD4 inserted into FD5 is the full \(4\)-by-\(4\) matrix
\(J^{-1}=DP(q)^{-1}\), in the original coordinate orders.
In particular it includes every \(dA,dB,dC,dD\) term.
An equivalent explicit outer-product formula is
\[
 J^{-1}=q_t\left(\frac{\ell'_r}{2t}
                    -\frac{\kappa\ell_r}{2t^3}\right)
          -\frac{q_r\ell_r}{t^2}+q_Ae_A^*+q_Be_B^*,
 \tag{FD6}
\]
where every column \(q_t,q_r,q_A,q_B\) is specified by FD5.
Thus FD6 has no unevaluated matrix inverse.

For completeness, the entire forward Jacobian also has a compact
column expression retaining all lower terms. Put
\[
\begin{aligned}
v(r)&=(1,-r^2,-2r^3,2r^4)^T,&
w(r)&=(0,1,-2r,r^2)^T,\\
R_0&=(4i,3ir^2,-22ir^3,15ir^4)^T,\\
R_1&=(3A,-Ar^2+2B-16r,
                  -10Ar^3-4Br+32r^2,8Ar^4+2Br^2-16r^3)^T,\\
R_2&=(0,-2iAr-14i,-8iAr^2-2iB+22ir,
                                10iAr^3+2iBr-8ir^2)^T,\\
R_3&=(0,-2A,4Ar-2,-2Ar^2+2r)^T,\\
S_0&=(0,2Ar-13,8Ar^2+2B+32r,
                              -10Ar^3-2Br-19r^2)^T.
\end{aligned}
\]
Writing \(e_4=(0,0,0,1)^T\), its four columns are exactly
\[
\begin{aligned}
J_{\bullet a}&=\frac{2r}{t}v+R_0+tR_1+t^2R_2+t^3R_3-it^4e_4,\\
J_{\bullet y}&=\frac2{t^2}v+\frac{14ir}{t}w+S_0+t^2e_4,\\
J_{\bullet z}&=\frac1{t^3}v+\frac2t w,\qquad
J_{\bullet w}=\frac1{t^2}w .
\end{aligned}\tag{FD7}
\]
These columns follow by differentiating FD1 before substituting
FD2. The independent certificate verifies both products of FD7
with FD4--5 to be the identity over
\(\mathbb Q(i)(t,r,A,B)\).

The two intermediate Jacobians have determinants
\[
 \det Dq=t^3,\qquad \det D\pi=-2t^3,\qquad \det J=-2.
 \tag{FD8}
\]
For the first formula, after the first two rows the diagonal
entries of the remaining triangular minor are \(t^3,t^2\);
the first two pivots are \(-t^{-2},-1\). For the second,
move the two coefficient columns \(A,B\) to the front, an
even permutation. The residual \(t,r\) block is
\(\left(\begin{smallmatrix}2t&-\kappa\\-2rt&r\kappa-t^2\end{smallmatrix}\right)\),
whose determinant is \(-2t^3\). The chain rule proves the
last equality for every \(t\ne0\). These parameter values
cover the original \(a\ne0\) chart by \(t=1/a\),
\(r=-y-i/a\), \(A=P_0\), \(B=P_2\). Polynomial continuation
therefore gives \(\det DP=-2\) everywhere in the original
affine source. No finite critical point is created here.

\section{The exact complex Puiseux branches}
Fix an actual target polynomial \(h_0(s)=As^4+s^3+Bs^2+Cs+D\)
with a finite root \(r_0\) of exact multiplicity \(m=2,3,4\).
For \(x=s-r_0\), its original coefficients give exactly
\[
 h_0(r_0+x)=c_2x^2+c_3x^3+c_4x^4,\qquad
 c_2=6Ar_0^2+3r_0+B,\quad c_3=4Ar_0+1,\quad c_4=A.
 \tag{FD9}
\]
The first nonzero coefficient is \(c_m\).
For \(m=4\), in particular, \(c_3=0\) forces
\(A=-1/(4r_0)\) and \(r_0\ne0\); these are consequences
of the retained cubic coefficient one, not free normalizations.
The specified unfolding is exactly
\[
 h_\varepsilon(s)=h_0(s)-\varepsilon,\qquad
 u_\varepsilon=u_0-\varepsilon e_D.
 \tag{FD10}
\]
Keep the actual complex \(c_m\ne0\). The analytic map
\[
 f(x)=x\left(\frac{h_0(r_0+x)}{c_mx^m}\right)^{1/m}
 \tag{FD11}
\]
uses the unique local root with value one at \(x=0\).
Its derivative is one there, so it has a local analytic inverse
\(g(w)=w+O(w^2)\). This follows, for example, by the complex
inverse-function theorem; its nonzero derivative is explicit.
Let \(\zeta_j=e^{2\pi ij/m}\), \(0\le j<m\), and use the
Puiseux parameter \(u\) with \(\varepsilon=c_mu^{2m}\).
Then all \(m\) local roots and both original factor signs are
given by the exact analytic branches
\[
\begin{aligned}
 r_j(u)&=r_0+g(\zeta_j u^2),\\
 t_{j,\pm}(u)&=\theta_{j,\pm}u^{m-1}a_j(u),\qquad
 \theta_{j,\pm}^2=mc_m\zeta_j^{m-1},\qquad
 \theta_{j,-}=-\theta_{j,+},
\end{aligned}\tag{FD12}
\]
where \(a_j(u)\) is the analytic square root with value one of
\[
 \frac{h_0'(r_j(u))}
                  {mc_m\zeta_j^{m-1}u^{2m-2}}.
\]
The quotient has analytic extension with value one. In particular
\[
\begin{aligned}
 r_j-r_0&=\zeta_ju^2+O(u^4),\\
 t_{j,\pm}&=\theta_{j,\pm}u^{m-1}(1+O(u^2)),\\
 \kappa_j=h_0''(r_j)&=K_j u^{2m-4}(1+O(u^2)),\qquad
 K_j=m(m-1)c_m\zeta_j^{m-2},\\
 t_{j,\pm}^2/\kappa_j&=\frac{\zeta_j}{m-1}u^2(1+O(u^2)).
\end{aligned}\tag{FD13}
\]
These formulas hold as complex \(u\to0\), on either sign
branch, with every phase in \(\theta_{j,\pm},K_j,\zeta_j\)
retained. For sufficiently small nonzero \(u\), both \(t\)
and \(\kappa\) are nonzero. All resulting
\(q(t_{j,\pm}(u),r_j(u),A,B)\) are actual original source
points, by FD2 and FD10--12. No claim about a different
target path is used.

\section{All phased forward and inverse exterior limits}
Use increasing coordinate subsets of \(\{1,2,3,4\}\), in
the order \(a,y,z,w\) in the domain and \(A,B,C,D\) in the
codomain. A superscript star on a coordinate wedge \(e_I^*\)
in the following matrices means its coordinate covector;
it is not a metric adjoint or a conjugation of the root.
Let \(v_0=v(r_0)\), \(w_0=w(r_0)\),
\(\ell_0=(r_0^4,r_0^2,r_0,1)\), and
\(\ell'_0=(4r_0^3,2r_0,1,0)\). Define
\[
\begin{aligned}
 b_3&=(1,-r_0,r_0^2,-r_0^4)^T
       &&\text{in the wedge order }123,124,134,234,\\
 \eta_3&=(r_0^4,r_0^3,-r_0^2/2,-1/2)
       &&\text{in that same wedge order}.
\end{aligned}
\]
Along each branch FD12 the complete matrix limits are
\[
\boxed{\begin{array}{c|c|c}
 r&\text{forward limit}&\text{inverse limit}\\\hline
 1&t^3J\longrightarrow v_0e_3^*
       &(t^5/\kappa)J^{-1}\longrightarrow\tfrac12e_1\ell_0\\
 2&t^5\wedge^2J\longrightarrow(v_0\wedge w_0)e_{34}^*
       &t^5\wedge^2J^{-1}\longrightarrow
                    -\tfrac12e_{12}(\ell'_0\wedge\ell_0)\\
 3&(t^5/\kappa)\wedge^3J\longrightarrow b_3e_{234}^*
       &t^3\wedge^3J^{-1}\longrightarrow e_{124}\eta_3\\
 4&\wedge^4J=-2&\wedge^4J^{-1}=-\tfrac12
\end{array}.}\tag{FD14}
\]
Every entry and sign can also be read as the following
explicit wedge coefficient rows:
\[
\begin{aligned}
 v_0\wedge w_0&=(1,-2r_0,r_0^2,4r_0^3,-3r_0^4,2r_0^5)^T,\\
 \ell'_0\wedge\ell_0&=(2r_0^5,3r_0^4,4r_0^3,r_0^2,2r_0,1).
\end{aligned}\tag{FD15}
\]
For \(r=1\) on the forward side, FD7 gives the limit
entry by entry. For \(r=2\), the only order \(t^{-5}\)
minor uses columns \(z,w\). The potential order \(t^{-5}\)
term from columns \(y,z\) vanishes because both leading
columns are multiples of \(v(r)\); all remaining terms
have smaller pole. This proves the full second-compound
limit, not only one of its entries.

For the inverse first-compound limit, its first row in
FD4--5 is exactly
\[
 \frac{\kappa}{2t^5}\ell_r-\frac1{2t^3}\ell'_r.
 \tag{FD16}
\]
Multiplication by \(t^5/\kappa\), followed by FD13,
gives \(\ell_0/2\). The other rows tend to zero under
this factor: their terms contain positive powers of \(t\),
or such powers times \(t^2/\kappa\), after the multiplication.
For example the \(y\)-row becomes
\(t^3\ell_r/\kappa-it^4\ell'_r/(2\kappa)
+i t^2\ell_r/2\), which tends to zero. The \(z,w\)
rows follow directly from their bounded coefficient
polynomials in FD5 and the same relation
\(t^2/\kappa=O(u^2)\). This proves the inverse limit
also when \(\kappa\to0\), as happens for \(m=3,4\).

All remaining limits follow from the exact complementary
minor identity, with its original determinant:
\[
 \det(J^{-1})_{I,J}
 =(-1)^{\sum I+\sum J}\frac{\det J_{J^c,I^c}}{\det J}
 =-\tfrac12(-1)^{\sum I+\sum J}\det J_{J^c,I^c}.
 \tag{FD17}
\]
This identity follows from the adjugate formula, or by
expanding the determinant on complementary coordinate
subsets. For the third forward compound use the same
identity for \(J^{-1}\) and the first inverse limit.
It gives the column \((1,-r_0,r_0^2,-r_0^4)^T\)
at column \(234\). For the second inverse compound
it gives only row \(12\), with the row in FD15 multiplied
by \(-1/2\). For the third inverse compound it gives
only row \(124\), equal to \(\eta_3\). Thus the signs
and coordinate complements in FD14 are fully determined.

One can state these limits directly in the complex Puiseux
parameter without hiding any phase. For a fixed \(j,\pm\)
write \(\theta=\theta_{j,\pm}\), \(K=K_j\). The forward
scalings and coefficients are
\[
\begin{aligned}
 u^{3(m-1)}J&\longrightarrow\theta^{-3}v_0e_3^*,\\
 u^{5(m-1)}\wedge^2J&\longrightarrow
                     \theta^{-5}(v_0\wedge w_0)e_{34}^*,\\
 u^{3m-1}\wedge^3J&\longrightarrow K\theta^{-5}b_3e_{234}^*,
\end{aligned}\tag{FD18}
\]
and the inverse limits are respectively
\[
\begin{aligned}
 u^{3m-1}J^{-1}&\longrightarrow
                         \frac{K}{2\theta^5}e_1\ell_0,\\
 u^{5(m-1)}\wedge^2J^{-1}&\longrightarrow
                         -\frac1{2\theta^5}e_{12}
                                      (\ell'_0\wedge\ell_0),\\
 u^{3(m-1)}\wedge^3J^{-1}&\longrightarrow
                         \theta^{-3}e_{124}\eta_3.
\end{aligned}\tag{FD19}
\]
These give the phase for every one of the \(2m\) signed
branches, including the effect of \(t\mapsto-t\).

\section{The original metric and all constants}
Fix the original receiving map
\[
 \Psi=B_v(V_\rho^{\mathsf T})^{-1}K_{\rm ES}^{-1},\qquad
 (B_v)_{kj}=\begin{cases}
 \binom{v+j}{k}\mu_{v+j-k},&k\le j,\\0,&k>j,
 \end{cases}\quad 0\le k,j\le3,
 \tag{FD20}
\]
with its original fixed conductor order, moments,
quartet Vandermonde and collision matrix. Its determinant
is nonzero by the already proved FC27 and the original
\(\mu_v\ne0\). The norm on the target polynomial space
is unchanged:
\[
 \|p\|_\Gamma^2=\int_\R|p(c_*+iy)|^2dm_s(y),\qquad
 dm_s(y)=\frac{(2\pi)^{s/2}2^{s/2}}{4\pi\Gamma(s/2)}
                |\Gamma(s/4+iy/2)|^2dy .
 \tag{FD21}
\]
Let \(G=\Psi^*\mathcal H_\Gamma\Psi\) be its full positive
pullback. In the original monomial frame every entry is
explicit:
\[
 (\mathcal H_\Gamma)_{jk}
 =\sum_{p=0}^j\sum_{q=0}^k
 \binom jp\binom kq c_*^{j+k-p-q}(-i)^pi^q m_{p+q},
 \tag{FD22}
\]
where \(b=s/2\), \(M=(2\pi)^{s/2}\),
\[
 m_0=M,\quad m_2=Mb,\quad m_4=M(3b^2+2b),\quad
 m_6=M(15b^3+30b^2+16b),\qquad m_1=m_3=m_5=0.
\]
These are the original WCF Gamma moments \cite{WCF}.
All period coefficients, phases, center, mass and cross
terms remain in FD20--22. The fixed \(\Psi\) conjugates
the polynomial \(P\) to \(P_W=\Psi P\Psi^{-1}\); therefore
the derivative \(J\) uses the same \(G\) on its domain
and codomain. No moving-period limit is taken.

Define the three strictly positive original constants
\[
\begin{aligned}
 C_1&=\sqrt{v_0^*Gv_0}\sqrt{(G^{-1})_{33}},\\
 C_2&=\sqrt{\det([v_0\ w_0]^*G[v_0\ w_0])}
                       \sqrt{\det(G^{-1})_{34,34}},\\
 C_3&=\sqrt{G_{11}}\sqrt{\ell_0G^{-1}\ell_0^*}.
\end{aligned}\tag{FD23}
\]
Positivity follows from \(v_{0,1}=1\), independence
of \(v_0,w_0\), \(\ell_{0,4}=1\), and positive definiteness
of \(G\). The norm of a rank-one exterior map
\(x\lambda\) is \(\|x\|\|\lambda\|\): Cauchy--Schwarz
gives the upper bound and the Riesz representative of
\(\lambda\) attains it. The Gram determinant formula
for exterior norms therefore gives \(C_1,C_2\) for the
first two forward limits and \(C_3/2\) for the first
inverse limit in FD14.

For an invertible endomorphism of this fixed four-dimensional
Hermitian space, products of singular values prove the exact
identity
\[
 \|\wedge^rJ^{-1}\|_G
            =\frac{\|\wedge^{4-r}J\|_G}{|\det J|}
            =\tfrac12\|\wedge^{4-r}J\|_G .
 \tag{FD24}
\]
Indeed the inverse singular values are the reciprocals
in reverse order; canceling their product proves the
formula. Consequently the complete exterior asymptotes
in the original norm are
\[
\begin{array}{c|c|c}
 r&\|\wedge^rJ\|_G&\|\wedge^rJ^{-1}\|_G\\\hline
 1&\sim C_1|t|^{-3}&\sim C_3|\kappa|/(2|t|^5)\\
 2&\sim C_2|t|^{-5}&\sim C_2/(2|t|^5)\\
 3&\sim C_3|\kappa|/|t|^5&\sim C_1/(2|t|^3)\\
 4&2&1/2
\end{array}.\tag{FD25}
\]
The limits are norm limits of the full finite matrices,
not estimates from sampled numerical singular values.

Successive ratios of forward exterior norms give every
singular value in decreasing order:
\[
\boxed{\quad
 \sigma_1\sim C_1|t|^{-3},\qquad
 \sigma_2\sim\frac{C_2}{C_1}|t|^{-2},\qquad
 \sigma_3\sim\frac{C_3}{C_2}|\kappa|,\qquad
 \sigma_4\sim\frac{2|t|^5}{C_3|\kappa|}.
 \quad}\tag{FD26}
\]
For completeness \(\|\wedge^rJ\|\) is the product of the
largest \(r\) singular values: diagonalizing \(J^*J\) in
the fixed metric and taking wedges gives all \(r\)-fold
products. This proves the ratio step and verifies EA36.

The exact coefficient \(c_m\) gives
\[
 |t|\sim\sqrt m\,|c_m|^{1/(2m)}
                       |\varepsilon|^{(m-1)/(2m)},\qquad
 |\kappa|\sim m(m-1)|c_m|^{2/m}
                       |\varepsilon|^{(m-2)/m}.
 \tag{FD27}
\]
Substituting, with no amplitude set to one, yields
\[
\boxed{\begin{aligned}
 \sigma_1&\sim\frac{C_1}{m^{3/2}|c_m|^{3/(2m)}}
                  |\varepsilon|^{-3(m-1)/(2m)},\\
 \sigma_2&\sim\frac{C_2}{mC_1|c_m|^{1/m}}
                  |\varepsilon|^{-(m-1)/m},\\
 \sigma_3&\sim\frac{m(m-1)C_3|c_m|^{2/m}}{C_2}
                  |\varepsilon|^{(m-2)/m},\\
 \sigma_4&\sim\frac{2m^{3/2}|c_m|^{1/(2m)}}{(m-1)C_3}
                  |\varepsilon|^{(3m-1)/(2m)}.
\end{aligned}}\tag{FD28}
\]
All four constants multiply to exactly \(2\), agreeing
with FD8, and every phase has already been recorded in
the stronger complex matrix limits FD18--19. The powers are
\[
\begin{array}{c|rrrr}
m&\sigma_1&\sigma_2&\sigma_3&\sigma_4\\\hline
2&-3/4&-1/2&0&5/4\\
3&-1&-2/3&1/3&4/3\\
4&-9/8&-3/4&1/2&11/8
\end{array}.
\]
They belong to the specified constant-term path FD10.

\section{Original singular flags and exact fixed-plane volumes}
Every adjacent singular-value ratio tends to infinity
by FD28. Thus the leading singular flags are unique for
sufficiently small nonzero parameter. The leading right
one-, two-, and three-planes converge, in the original
metric, to
\[
 G^{-1}\Span(e_3),\qquad
 G^{-1}\Span(e_3,e_4),\qquad
 G^{-1}\Span(e_2,e_3,e_4).
 \tag{FD29}
\]
The corresponding leading left planes converge to
\[
 \Span(v_0),\qquad \Span(v_0,w_0),\qquad \ker\ell_0.
 \tag{FD30}
\]
For the proof, divide each full forward compound by its
norm and apply FD14. Its limit is a unit-norm rank-one
map, so its products with its adjoint tend to rank-one
orthogonal projections. The top eigenspace projections
converge: any subsequential limit of a unit top eigenvector
satisfies the limiting eigenvector equation and lies in
the unique limiting range. The right Riesz representative
of the coordinate wedge functional \(e_J^*\) is
\(\bigwedge_{j\in J}G^{-1}e_j\). The left third compound
vector \(b_3\) is decomposable and represents \(\ker\ell_0\),
as follows by taking its wedge coefficients or the
complementary-minor formula FD17. This proves both flags.
In particular the smallest right singular line tends to
\(\C e_1\), and the smallest left line to
\(\C G^{-1}\ell_0^*\), by taking their original-metric
orthogonal complements. Their actual polynomial-space
images are obtained by the fixed \(\Psi\).

The exact volume distortion on the fixed coordinate
input flags \(\Span(e_3)\), \(\Span(e_3,e_4)\), and
\(\Span(e_2,e_3,e_4)\) is respectively
\[
\begin{aligned}
 &\sim |t|^{-6}\frac{v_0^*Gv_0}{G_{33}},\\
 &\sim |t|^{-10}
              \frac{\det([v_0\ w_0]^*G[v_0\ w_0])}{\det G_{34,34}},\\
 &\sim |\kappa|^2|t|^{-10}
              \frac{\ell_0G^{-1}\ell_0^*}{(G^{-1})_{11}}.
\end{aligned}\tag{FD31}
\]
Each is a squared output Gram determinant divided by
the squared input volume in the original metric. For
the first two apply the compound limit to the displayed
coordinate wedge. For the third, the complementary
Gram identity gives
\(\|b_3\|_{\wedge^3G}^2
=\det G\,\ell_0G^{-1}\ell_0^*\) and
\(\det G_{234,234}=\det G(G^{-1})_{11}\).
This proves the last constant including its cross terms.
In degree four the squared-volume ratio is exactly four
at every finite point.

\section{Escape and the exact conductor relation}
Formula FD2 and the actual branches FD12 imply that
\(a=1/t\) is the sole divergent coordinate; the other
three tend to \((-r_0,0,0)\). Thus
\[
 |\varepsilon|^{(m-1)/(2m)}\|\Psi q\|_\Gamma
 \longrightarrow\frac{\sqrt{G_{11}}}
                         {\sqrt m\,|c_m|^{1/(2m)}}>0.
 \tag{FD32}
\]
The original factor is \(L=a(U-rV)\), so its actual
coefficient \(a\) diverges as well. This is intrinsic
factor-scale escape at the selected finite repeated root;
adding a missing \(a=0\) chart cannot supply a finite
limit of that coordinate. The image displacement is
exactly \(-\varepsilon\Psi e_4\). Consequently, if
\(\alpha_m=(m-1)/(2m)\),
\[
 \|\Psi(u_\varepsilon-u_0)\|_\Gamma^{\alpha_m}
 \|\Psi q\|_\Gamma
 \longrightarrow
 \frac{G_{44}^{\alpha_m/2}\sqrt{G_{11}}}
                         {\sqrt m\,|c_m|^{1/(2m)}}.
 \tag{FD33}
\]

The fixed original source map \(\Phi\) and conductor
restriction \(L=T_A|_U\) obey \(L\Phi=\Psi\).
For \(P_U=\Phi P\Phi^{-1}\), \(P_W=\Psi P\Psi^{-1}\),
differentiation at the matching points gives exactly
\[
 L\,DP_U(\Phi q)=DP_W(\Psi q)L.
 \tag{FD34}
\]
Both maps use their original respective norms; the
source-side formulas follow by replacing \(G\) by
the actual fixed \(G_U=\Phi^*\mathcal H_U\Phi\), with
no change to the differential matrices or phases.
In particular, exterior inverse transport is the exact
identity
\[
 \wedge^r DP_W^{-1}
   =(\wedge^rL)(\wedge^rDP_U^{-1})(\wedge^rL^{-1}).
 \tag{FD35}
\]
This is the original four-plane inverse from FC24, not
a variable conductor matrix or a claim that its leading
moment vanished. The finite-root degenerations above
are completely located in the specified nonlinear
derivative along its actual escaping branches.

\section{Exact verification and source use}
The independent script
\texttt{verify\_finite\_root\_native\_differential.py}
starts from FD1, differentiates it before substitution,
checks \(Pq=\pi\), both intermediate determinants,
\(JDq=D\pi\), both inverse products for FD4--5,
and every forward and inverse compound over
\(\mathbb Q(i)(t,r,A,B)\). Its full matrices and checks
are recorded in
\texttt{finite\_root\_native\_differential\_exact.json}.
No sampled period, artificial Gram or numerical
singular-value diagnostic is used in any proof above.
The Puiseux phases and the passage to the original metric
are proved explicitly in FD9--33.

\section{Residue duality on the complete signed-root algebra}
The finite completion FS1--25 retains nilpotents at a multiple
root. Here a nilpotent means a nonzero algebra element whose
positive power is zero. We compute a bilinear pairing that
detects every such element, the exact map from it to ordinary
trace, and both maps in the original conductor spaces.
The general relation between a coefficient-residue functional,
its Hankel matrix and the inverse Bezout matrix is classical;
we use Jouve and Rodriguez-Villegas \cite[Section 2, especially
Lemma \texttt{Q-lemma} and Theorem \texttt{thBsharp}]{JouveRV}.
All matrices and constants for the original quartic, its signed
extension and its conductor return are proved below.

Let $u=(A,B,C,D)$, with $A\ne0$, and retain
\begin{equation}
 h=Ar^4+r^3+Br^2+Cr+D,\quad
 E_u=\C[r]/(h),\quad
 S_u=E_u[t]/(t^2-h'),\quad
 \mathbf e=(1,r,r^2,r^3,t,tr,tr^2,tr^3).
 \tag{RD1}
\end{equation}
Division by $h$ has a unique remainder of degree at most three
because its leading coefficient is a unit. Thus $E_u$ has the
first four displayed basis vectors, and division by the monic
polynomial in $t$ proves that every element of $S_u$ has a
unique expression $R+tS$, with $R,S\in E_u$. Define
\begin{equation}
 \lambda_u(F)=\sum_{h(r_0)=0}\operatorname{Res}_{r=r_0}
                 \frac{F(r)\,dr}{h(r)}
     =\frac{[r^3](F\bmod h)}{A},\qquad
 \Lambda_u(R+tS)=\lambda_u(S),\qquad
 \mathcal B_u(F,G)=\Lambda_u(FG).
 \tag{RD2}
\end{equation}
The residue formula includes every multiple root. Subtracting
the remainder adds a polynomial differential with no finite
residue. For the remainder, its coefficient of $r^{-1}$ at
infinity is its cubic coefficient divided by $A$; the sum
of the finite residues is that coefficient. This proves RD2
directly and also proves that the definition descends to $E_u$.
The functional is complex-linear; the pairing is complex
bilinear and symmetric. No conjugation is used in RD2.

In the displayed order the two full Gram matrices and inverse are
\begin{align}
 H_u&=\begin{pmatrix}
 0&0&0&A^{-1}\\
 0&0&A^{-1}&-A^{-2}\\
 0&A^{-1}&-A^{-2}&A^{-3}-BA^{-2}\\
 A^{-1}&-A^{-2}&A^{-3}-BA^{-2}&-A^{-4}+2BA^{-3}-CA^{-2}
 \end{pmatrix},\nonumber\\
 J_u&=\begin{pmatrix}C&B&1&A\\B&1&A&0\\1&A&0&0\\A&0&0&0\end{pmatrix},
 &\mathcal B_u&\longleftrightarrow
 B_u=\begin{pmatrix}0&H_u\\H_u&0\end{pmatrix},\qquad
 B_u^{-1}=\begin{pmatrix}0&J_u\\J_u&0\end{pmatrix}.
 \tag{RD3}
\end{align}
To prove every entry, put $l_j=\lambda_u(r^j)$. RD2 gives
$l_0=l_1=l_2=0$, $l_3=A^{-1}$; multiplying $h=0$ by
$1,r,r^2$ gives successively
$l_4=-A^{-2}$, $l_5=A^{-3}-BA^{-2}$ and
$l_6=-A^{-4}+2BA^{-3}-CA^{-2}$.
Consequently $(H_u)_{ij}=l_{i+j}$ for $0\le i,j\le3$.
The product $(R+tS)(P+tQ)$ has odd coefficient $RQ+SP$,
which proves the block matrix. Direct row multiplication
gives $H_uJ_u=J_uH_u=I_4$. Equivalently, $J_u$ is the
coefficient matrix of $(h(x)-h(y))/(x-y)$; in the sign
convention of \cite[Section 2.3]{JouveRV} it is $B(h,1)$.
Expansion along the reversed diagonal gives
\begin{equation}
 \det H_u=A^{-4},\qquad \det B_u=A^{-8}.
 \tag{RD4}
\end{equation}
In particular $B_u$ is invertible on the entire $A\ne0$
chart, including its discriminant. Explicit dual vectors
are $t\sum_j(J_u)_{ji}r^j$ for $r^i$, and
$\sum_j(J_u)_{ji}r^j$ for $tr^i$. Pairing against those
eight vectors recovers the eight original coefficients.

\section{The exact passage from residue to trace}
For $F\in S_u$ let $M_F$ denote multiplication by $F$.
Define $L=M_{h'}$ on $E_u$. Multiplication by $t$ and by
$2t^3$ in the order RD1 are
\begin{equation}
 M_t=\begin{pmatrix}0&L\\I_4&0\end{pmatrix},\qquad
 M_{2t^3}=2\begin{pmatrix}0&L^2\\L&0\end{pmatrix}.
 \tag{RD5}
\end{equation}
Every trace, including on a nonreduced fibre, satisfies
\begin{equation}
 \operatorname{Tr}_{S_u}(M_F)=\Lambda_u(2t^3F),\qquad
 T_u:=\bigl(\operatorname{Tr}M_{e_i e_j}\bigr)_{ij}
    =B_uM_{2t^3}
    =2\begin{pmatrix}H_uL&0\\0&H_uL^2\end{pmatrix}.
 \tag{RD6}
\end{equation}
Here is a proof valid on each fibre without a distinct-root
assumption. On the length-$m$ algebra at $r_0$, multiplication
by $R$ is triangular in powers of $r-r_0$, and its diagonal
is $R(r_0)$ repeated $m$ times. Its trace is $mR(r_0)$.
Writing $h=(r-r_0)^m g$ shows that the residue of $h'R/h$
is also $mR(r_0)$. Addition over the roots gives
$\operatorname{Tr}_{E_u}M_R=\lambda_u(h'R)$.
On $E_u\oplus tE_u$ the diagonal blocks of $M_{R+tS}$
are both $M_R$, so its trace is $2\lambda_u(h'R)$.
The odd coefficient of $2t^3(R+tS)$ is $2h'R$.
This proves the functional identity and then the Gram identity.

The matrix of $M_r$ on $E_u$ is completely explicit:
\begin{equation}
 Z=\begin{pmatrix}
 0&0&0&-D/A\\1&0&0&-C/A\\0&1&0&-B/A\\0&0&1&-1/A
 \end{pmatrix},\qquad
 L=4AZ^3+3Z^2+2BZ+CI_4 .
 \tag{RD7}
\end{equation}
Thus RD5--6 contain no unevaluated quotient multiplication.
With the usual quartic discriminant retaining the factor $A^6$,
$\det L=\prod h'(r_j)=\operatorname{Disc}(h)/A^2$:
at four distinct roots each unordered difference contributes
its square and the six sign changes have positive product.
Both sides are polynomials in the coefficients with powers of
$A$ inverted, so this identity holds at repeated roots too.
Taking determinants in RD5--6 now proves
\begin{equation}
 \det M_{2t^3}=2^8(\det L)^3,\qquad
 \det T_u=\frac{256\operatorname{Disc}(h)^3}{A^{14}}.
 \tag{RD8}
\end{equation}
The determinant loss is exactly multiplication by $2t^3$.
The residue identification $S_u\to S_u^*$ in RD3 is an
isomorphism; its composition with $M_{2t^3}$ is the trace
identification. This specifies the actual morphism and its
kernel, rather than merely comparing the names of two forms.

\section{Every local nilpotent and every lost trace direction}
At a root of multiplicity $m\in\{2,3,4\}$ put
$\epsilon=r-r_0$ and write the complete polynomial
$h(r_0+\epsilon)=\epsilon^m g(\epsilon)$, $g_0=g(0)\ne0$.
The local signed algebra and its residue functional are
\begin{equation}
 S_{u,r_0}=\C[\epsilon,t]/
   (\epsilon^m,\ t^2-mg_0\epsilon^{m-1}),\qquad
 \Lambda_{r_0}(R+tS)=[\epsilon^{m-1}]\frac{S}{g(\epsilon)}.
 \tag{RD9}
\end{equation}
The first relation gives the local root algebra by the
polynomial Chinese remainder construction; differentiating
$\epsilon^m g$ and reducing modulo $\epsilon^m$ gives
the second relation with its exact coefficient $mg_0$.
The last equation is the definition of the residue at $r_0$.
All coefficients of $g$ needed in the truncated inverse
remain in that functional.

Write $g^{-1}=\sum_{j=0}^{m-1}b_j\epsilon^j$ modulo
$\epsilon^m$. These coefficients are finite exact expressions:
$b_0=g_0^{-1}$ and
$b_j=-g_0^{-1}\sum_{k=1}^j g_k b_{j-k}$.
In the order $(1,\epsilon,\ldots,\epsilon^{m-1},
t,t\epsilon,\ldots,t\epsilon^{m-1})$ the Gram matrix is
\begin{equation}
 B_{r_0}=\begin{pmatrix}0&H_{r_0}\\H_{r_0}&0\end{pmatrix},
 \quad (H_{r_0})_{ij}=
 \begin{cases}b_{m-1-i-j}&i+j\le m-1,\\0&i+j>m-1,\end{cases}
 \quad \det B_{r_0}=(-1)^m g_0^{-2m}.
 \tag{RD10}
\end{equation}
The determinant follows from the anti-triangular matrix:
$\det H_{r_0}=(-1)^{m(m-1)/2}g_0^{-m}$; swapping its two
blocks adds $(-1)^m$. A more explicit proof of invertibility
gives every coefficient of $F=R+tS$:
\begin{equation}
 R_j=\Lambda_{r_0}(F\,t\epsilon^{m-1-j}g),\qquad
 S_j=\Lambda_{r_0}(F\,\epsilon^{m-1-j}g),
 \quad 0\le j<m .
 \tag{RD11}
\end{equation}
Indeed the first odd coefficient is $R\epsilon^{m-1-j}g$,
and the second is $S\epsilon^{m-1-j}g$. Dividing by the
unchanged $g$ in RD9 extracts precisely the stated coefficient.
Thus every nonzero nilpotent has an explicitly specified
test vector with nonzero pairing.

On the other hand $t^4=0$ and
\begin{equation}
 2t^3F=2mg_0 R_0\,t\epsilon^{m-1},\qquad
 \operatorname{Tr}M_F=2mR_0,\qquad
 \ker M_{2t^3}=(\epsilon,t),\quad
 \operatorname{im}M_{2t^3}=\C\,t\epsilon^{m-1}.
 \tag{RD12}
\end{equation}
Here $t^4=m^2g_0^2\epsilon^{2m-2}=0$ because $m\ge2$.
Multiplication by $t^3=mg_0t\epsilon^{m-1}$ kills all
terms except $R_0$. The image vector has
$\Lambda_{r_0}(t\epsilon^{m-1})=g_0^{-1}$, proving
the trace formula and both displayed spaces. The trace
pairing is exactly $2mR_0P_0$ and has rank one on this
length-$2m$ algebra. For a simple root $h'(r_0)\ne0$,
the two signs are distinct, $t$ is a unit, and the trace
pairing has rank two. Therefore for a fibre with $n_1$
simple roots and $n_{\ge2}$ multiple roots the full trace
rank is $2n_1+n_{\ge2}$, while residue rank is always eight.

\section{The paired critical limit in the original quartet}
Retain $\gamma>0$, $\eta\in\{+1,-1\}$ and the original
factor $-1/2$. At $\delta=0$ put
\begin{equation}
 h_0(r)=-\tfrac12\bigl((r-\tfrac12)^2+\gamma^2\bigr)^2,
 \quad r_\eta=\tfrac12+i\eta\gamma,\quad
 g_\eta(\epsilon)=-\tfrac12(\epsilon+2i\eta\gamma)^2 .
 \tag{RD13}
\end{equation}
At each of the two roots the exact relations are
$g_0=2\gamma^2$, $g_1=-2i\eta\gamma$,
$t^2=4\gamma^2\epsilon$. Hence the two explicit mutually
inverse local homomorphisms identify the algebra with
$\C[t]/(t^4)$ by
$\epsilon\mapsto t^2/(4\gamma^2)$ and $t\mapsto t$.
This is the double-root case only; the triple and quadruple
root algebras remain the two-generator algebras RD9.
In the unchanged powers $(1,t,t^2,t^3)$ one obtains
\begin{equation}
 \Lambda_\eta(1)=\Lambda_\eta(t^2)=0,\quad
 \Lambda_\eta(t)=\frac{i\eta}{2\gamma^3},\quad
 \Lambda_\eta(t^3)=2,\qquad
 B_\eta=\begin{pmatrix}
 0&i\eta/(2\gamma^3)&0&2\\
 i\eta/(2\gamma^3)&0&2&0\\
 0&2&0&0\\2&0&0&0
 \end{pmatrix},\quad\det B_\eta=16.
 \tag{RD14}
\end{equation}
For example $\Lambda(t)=-g_1/g_0^2$ by RD9, and
$\Lambda(t^3)=4\gamma^2\Lambda(t\epsilon)=2$.
These evaluations prove every entry and the determinant.
Ordinary trace in this basis is $4$ times the product of
constant coefficients, so its rank is one in each cluster.

Set $X(\gamma)=\xi(1/2+i\gamma)$. The entire completion
and its reflection law are those of AR10--12, with the
human sources cited there. Reflection gives
$\xi(r_\eta)=X(\gamma)$ and differentiation gives
$\xi'(r_\eta)=-i\eta X'(\gamma)$. Thus its exact class is
\begin{equation}
 j_\eta\xi=X(\gamma)-\frac{i\eta X'(\gamma)}{4\gamma^2}t^2,
 \quad
 \Lambda_\eta((j_\eta\xi)t^3)=2X(\gamma),\quad
 \Lambda_\eta((j_\eta\xi)t)=
 \frac{i\eta X(\gamma)}{2\gamma^3}
       -\frac{i\eta X'(\gamma)}{2\gamma^2}.
 \tag{RD15}
\end{equation}
Every equality follows by reducing the Taylor series modulo
$\epsilon^2$ and using RD14. These two measurements recover
both $X$ and $X'$ without division by either:
\begin{equation}
 X=\tfrac12\Lambda_\eta((j_\eta\xi)t^3),\quad
 X'=2i\eta\gamma^2\left[
 \Lambda_\eta((j_\eta\xi)t)
 -\frac{i\eta}{4\gamma^3}\Lambda_\eta((j_\eta\xi)t^3)
 \right].
 \tag{RD16}
\end{equation}
The coefficient of $t^2$ is therefore not erased by passing
to the completed fibre. This statement makes no assumption
about whether a particular $X$ or $X'$ is zero. For any
entire $F$ of order $\nu$ at a root of multiplicity $m$,
its multiplication rank on RD9 is exactly
$2\max(m-\nu,0)$: multiplication by its unit Taylor factor
is invertible, and multiplication by $\epsilon^\nu$ has
rank $m-\nu$ on each of the two free parity summands when
$\nu<m$, and rank zero otherwise.

Both local measurements are literal measurements on the
whole eight-dimensional fibre. Put $d_\eta=2i\eta\gamma$,
$\epsilon_\eta=r-r_\eta$ and define its root projector by
\begin{equation}
 e_\eta(r)=\frac{(r-r_{-\eta})^2}{d_\eta^2}
                \left(1-\frac{2\epsilon_\eta}{d_\eta}\right)
                  \pmod{h_0}.
 \tag{RD17}
\end{equation}
It is $1$ modulo $\epsilon_\eta^2$ and $0$ modulo
$(r-r_{-\eta})^2$. The two ideals are coprime since
$d_\eta\ne0$, so $e_++e_-=1$, $e_\eta^2=e_\eta$ and
$e_+e_-=0$. These identities prove the direct product
decomposition, with the original signs retained.
For every $F$, $\Lambda_\eta(F)=\Lambda_u(e_\eta F)$.
Thus RD15 uses the global test vectors $e_\eta t^3$ and
$e_\eta t$, with all products reduced in the exact RD1 basis.

\section{Every dual vector in the two original conductor spaces}
Fix the original period and its four distinct labelled reference
roots, as in FC26--32. The reference quartet is held fixed while
$u$ varies in RD1; it is not replaced by the colliding roots.
Write $\mathcal U=\mathcal P_{v+3}/\mathcal P_{v-1}$,
$\mathcal W=\mathcal P_3$, and keep the exact maps
$X,Y,\Phi,\Psi,E,O$ of FC32 and SF2--3. In particular
$Y=T_AX$, $\Psi=T_A\Phi$, and $O$ is the original invertible
odd frame. In the actual target variable $S'$ the center is
$c_*=a_{\rm amp}/2-4$; retain $\kappa=c_*-1/2$.
The spectral coefficient map is
\begin{equation}
 \mathcal R_u:\mathcal W\oplus\mathcal W\longrightarrow S_u,
 \quad(p,q)\longmapsto p(r+\kappa)+tq(r+\kappa),\qquad
 \mathcal R_u^{-1}(R+tS)=(R(S'-\kappa),S(S'-\kappa)).
 \tag{RD18}
\end{equation}
The remainders $R,S$ have degree at most three, so the
two substitution maps compose to the identity in both
directions. Their coefficient matrix is
$D_\kappa=\operatorname{diag}(C_\kappa,C_\kappa)$ with
the full binomial matrix RW5; its determinant is one.

For the original eight signed coordinates $z=(\alpha,\beta)$ set
\begin{align}
 \mathcal J_W z&=(Y\alpha,\Psi(E\alpha+O\beta)),\nonumber\\
 \mathcal J_U z&=(X\alpha,\Phi(E\alpha+O\beta)),\qquad
 \mathcal J_W=(T_A\oplus T_A)\mathcal J_U,\nonumber\\
 Q_u&=D_\kappa
 \begin{pmatrix}Y&0\\\Psi E&\Psi O\end{pmatrix},\qquad
 \mathcal Q_u=\mathcal R_u\mathcal J_W
      =\mathcal R_u(T_A\oplus T_A)\mathcal J_U .
 \tag{RD19}
\end{align}
The entries of $Y,\Psi$ in this matrix are in the original
target monomial basis. Its complete inverse on a class
$R+tS$ is
\begin{equation}
 \alpha=Y^{-1}C_{-\kappa}[R],\qquad
 \beta=O^{-1}\left(\Psi^{-1}C_{-\kappa}[S]-E\alpha\right).
 \tag{RD20}
\end{equation}
Square brackets here mean the cubic coefficient column.
Substitution in RD19 proves both inverse identities. Its
actual source state is, without changing its quotient norm,
$(T_A^{-1}R(S'-\kappa),T_A^{-1}S(S'-\kappa))$.
Every coefficient of this inverse is the finite recurrence
FC24 with its original moments and weights. Thus RD17 and
RD20 explicitly return every critical derivative test to
the original source, including its retained label component.

The residue and trace matrices in the signed coordinates are
\begin{equation}
 \widetilde B_u=Q_u^{\mathsf T}B_uQ_u,\quad
 \widetilde B_u^{-1}=Q_u^{-1}B_u^{-1}Q_u^{-\mathsf T},\quad
 \widetilde T_u=\widetilde B_u\,\widetilde D_u,\quad
 \widetilde D_u=Q_u^{-1}M_{2t^3}Q_u,\quad
 \det\widetilde B_u=\frac{(\det Y\det\Psi\det O)^2}{A^8}.
 \tag{RD21}
\end{equation}
These are transpose congruences because residue is bilinear.
All inverses exist for every $A\ne0$, including $\delta=0$
in RD13. In contrast, the exact kernel of the trace matrix
is $Q_u^{-1}\ker M_{2t^3}$, with every local component
given by RD12. No kernel of the original $T_A$ has been
introduced by this composition.

The original positive norms are retained separately. Set
$L_0=\left(\begin{smallmatrix}I&0\\E&O\end{smallmatrix}\right)$.
With the original full Gram matrices of RW14 and RW19,
\begin{equation}
 \mathbb H_W=L_0^*\operatorname{diag}(G_Y,G)L_0,
 \quad
 \mathbb H_U=L_0^*\operatorname{diag}(G_X,G_U)L_0,
 \quad
 \|\ell\|_{W^*}^2=\ell\mathbb H_W^{-1}\ell^*,\quad
 \|\ell\|_{U^*}^2=\ell\mathbb H_U^{-1}\ell^* .
 \tag{RD22}
\end{equation}
Here $\ell$ is any one of the exact linear rows
$f^{\mathsf T}\widetilde B_u$ furnished by RD11 or RD17--21.
The dual norm follows by writing a positive Gram as $H^{1/2}H^{1/2}$
and applying Cauchy--Schwarz to $\ell H^{-1/2}$, with equality
on its adjoint column. This derivation changes no input norm.
It includes all Gamma masses, original centers and moments
already evaluated in FC28--30 and RW19--20.

The existing full Weil form also has an exact expression
through this nondegenerate residue pairing. For its actual
Hermitian matrix $\mathbb W$ from RW14, define the
conjugate-linear map
\begin{equation}
 \mathscr A_u(z)=\widetilde B_u^{-1}\mathbb W^{\mathsf T}\bar z,
 \qquad
 \mathcal B_u(\mathcal Q_u\mathscr A_u(z),\mathcal Q_u w)
       =z^*\mathbb Ww.
 \tag{RD23}
\end{equation}
Transpose the displayed formula and use the symmetry of
$\widetilde B_u$ to prove the equality. Since the latter
matrix is invertible and $\mathbb W$ is Hermitian,
$\ker\mathscr A_u=\ker\mathbb W$; its exact kernel is
RW17, and its rank equals that of $\mathbb W$.
Equation RD23 is a proved correspondence, not a claim that
bilinear nondegeneracy implies Hermitian positivity.

\section{The remaining root charts and the seven-point fibre}
No zero of $A$ is discarded from the finite completion.
On the open $h(p)\ne0$, $p\in\{0,1,-1,2\}$, use the
actual determinant-one chart FS7:
\begin{equation}
 R=(p-r)^{-1},\quad T=t/(p-r),\quad
 h_p(R)=R^4h(p-R^{-1})=aR^4+bR^3+cR^2+dR+e .
 \tag{RD24}
\end{equation}
Its leading coefficient is $a=h(p)\ne0$ and its signed
equation remains $T^2=h_p'(R)$. Define
$\Lambda_p(U+TV)=[R^3](V\bmod h_p)/a$.
The complete matrices in this chart are
\begin{equation}
 H_p=\begin{pmatrix}
 0&0&0&a^{-1}\\0&0&a^{-1}&-ba^{-2}\\
 0&a^{-1}&-ba^{-2}&b^2a^{-3}-ca^{-2}\\
 a^{-1}&-ba^{-2}&b^2a^{-3}-ca^{-2}&-b^3a^{-4}+2bca^{-3}-da^{-2}
 \end{pmatrix},\quad
 H_p^{-1}=\begin{pmatrix}d&c&b&a\\c&b&a&0\\b&a&0&0\\a&0&0&0\end{pmatrix}.
 \tag{RD25}
\end{equation}
The same three remainder recurrences prove every entry and
their product is the identity. The signed Gram is
$B_p=\left(\begin{smallmatrix}0&H_p\\H_p&0\end{smallmatrix}\right)$,
$\det B_p=a^{-8}$. The trace map is again exactly $M_{2T^3}$.
On the overlap with $A\ne0$, let $\iota$ substitute
$r=p-1/R$, $t=T/R$. The derivative identity in FS7 proves
that this is an algebra isomorphism, with the displayed
inverse, and the residue change of variable gives
\begin{equation}
 \Lambda_u(F)=\Lambda_p(R^3\iota(F)).
 \tag{RD26}
\end{equation}
Indeed $dr/h=R^2dR/h_p$, while the odd coefficient of
$\iota(F)$ is $S(p-1/R)/R$. Multiplication by $R^3$
therefore gives precisely the transformed residue. On this
overlap $R$ is a unit, so the transition is invertible.
This identifies the two dualities with the exact factor
$R^3$ retained.

There is also a full physical coefficient receiver on this
chart, including $R=0$. For $f\in\mathcal P_3$ define
\begin{equation}
 (Z_pf)(R)=R^3 f(p+\kappa-R^{-1}),\qquad
 f(S')=(p+\kappa-S')^3
       (Z_pf)((p+\kappa-S')^{-1}).
 \tag{RD27}
\end{equation}
Each expression is a polynomial of degree at most three;
the formulas compose to the identity by substitution.
For every monomial its forward image is exactly
$R^{3-j}((p+\kappa)R-1)^j$, $0\le j\le3$.
Reversing the column order gives diagonal $(-1)^j$;
both resulting signs are $(-1)^6=1$, so $\det Z_p=1$.
Thus $(f,g)\mapsto Z_pf+T Z_pg$ is an isomorphism
$\mathcal W^2\to S_{u,p}$ even at the infinity root.
Composing with $\mathcal J_W$ and its exact inverse RD20,
using $Z_p$ in place of $C_\kappa$, proves every analogue
of RD19--22 with $B_p$ and determinant denominator $a^8$.

This coefficient receiver is not silently identified with
the algebra change of variable. On their overlap its two
parity parts are $R^3\iota(F_{\mathrm{even}})$ and
$R^4\iota(F_{\mathrm{odd}})$. Therefore for two physical inputs
with finite-chart images $F,G$, the two pulled-back pairings
are related by the explicit invertible multiplier
\begin{equation}
 \Lambda_p(F_pG_p)
       =\Lambda_u((p-r)^{-4}FG).
 \tag{RD28}
\end{equation}
The odd product acquires $R^7$, and RD26 removes $R^3$,
leaving $R^4=(p-r)^{-4}$, which proves this identity.
The five opens cover the entire target by FS6. Hence these
formulas supply receivers and their exact comparisons on
every fibre of the original finite completion.

At the original target $(A,B,C,D)=(0,0,-1,0)$ take $p=2$.
Then the full polynomial, roots and residue determinant are
\begin{equation}
 h_2(R)=6R^4-11R^3+6R^2-R,\quad
 R\in\{0,\tfrac13,\tfrac12,1\},\quad
 \det B_2=6^{-8},\quad
 \det T_2=256\cdot4^3/6^{14}.
 \tag{RD29}
\end{equation}
Expansion of $(2R-1)^3R-(2R-1)R^3$ proves the polynomial.
Its factorization $R(R-1)(2R-1)(3R-1)$ proves the roots;
their difference product gives discriminant $4$.
At $R=0$, $T^2=h_2'(0)=-1$ gives both $T=\pm i$.
FS7 sends these to the two original infinity signs;
$T=-i$ is the original affine state and $T=i$ is the
omitted state. Thus all eight completed states, the seven
original affine states and the exact missing sign are
present in the same invertible residue calculation.

\section{Complete original-source inverse constants and marked return}


The finite signed-root algebra and its exact trace-to-residue
multiplier are those of RD1--29 and the received continuation
RE1--30 \cite{Fable037,ReceivedRE}. We derive the complete inverse singular
constants for every retained positive coefficient Gram, then
evaluate that Gram through the full original conductor inverse.
The original fixed conductor, its moments and both Gamma centers
are retained. The Gamma polynomial dictionary is the one credited
to the human authors of DLMF Chapter 18 in WCF1--6
\cite{WCF,DLMF18Source}; its application to the original conductor is
FC20--24 and NH26--34 \cite{Fable037}. The new calculation here is the
arbitrary-Gram singular decomposition and its exact original-source
specialization, rather than a new orthogonality result.

\section{The full signed multiplier with its complex phase}
Fix the original reference period and its conductor throughout.
Independently vary the real collision parameter $w=\delta^2>0$,
keeping $g=\gamma^2>0$ fixed. In the unchanged coordinate
$x=r-1/2$, set
\begin{equation}
 h_w(x)=-\tfrac12\bigl(x^4+2(g-w)x^2+(g+w)^2\bigr),\qquad
 S_w=\C[x,t]/(h_w,t^2-h_w'),\qquad
 \mathbf e=(1,x,x^2,x^3,t,tx,tx^2,tx^3).
 \tag{RS1}
\end{equation}
The factor $-1/2$ is unchanged: after replacing $x$ by $r-1/2$
the coefficient of $r^3$ is precisely one. The matrix of
multiplication by $x$ in each parity summand and the derivative
multiplier are
\begin{equation}
 X_w=\begin{pmatrix}
 0&0&0&-(g+w)^2\\1&0&0&0\\0&1&0&-2(g-w)\\0&0&1&0
 \end{pmatrix},\qquad
 H_w=-2X_w\bigl(X_w^2+(g-w)I\bigr).
 \tag{RS2}
\end{equation}
These follow by reducing $x^4$ in the exact algebra RS1 and
differentiating its unchanged $h_w$. In particular,
\begin{equation}
 H_w^2=-16gwX_w^2,\qquad
 \mathsf J_w=M_{2t^3}=2\begin{pmatrix}0&H_w^2\\H_w&0\end{pmatrix},
 \qquad c=\frac{1+i}{\sqrt2},\qquad \Theta_w=\mathsf J_w/c.
 \tag{RS3}
\end{equation}
To verify the first identity, the relation RS1 gives
$(X_w^2+(g-w)I)^2=-4gwI$; multiplying by $4X_w^2$
gives the asserted expression. The second follows by cubing
$M_t=\left(\begin{smallmatrix}0&H_w\\I&0\end{smallmatrix}\right)$.
The phase $c$ is the reflected-residue convention RE10--12:
$x^\sharp=-x$, $t^\sharp=it$, with conjugation of coefficients.
It respects $t^2=h_w'(x)$ and has square one. In that convention
the reflected trace is exactly the reflected residue composed
with $\Theta_w$, not with a multiplier in which $c$ is omitted.
We retain this $c$ in every operator identity below.

Since $\det X_w=(g+w)^2$, RS3 gives for $w>0$
\begin{equation}
 \mathsf J_w^{-1}
 =\frac12\begin{pmatrix}0&H_w^{-1}\\H_w^{-2}&0\end{pmatrix}
 =\begin{pmatrix}
 0&\dfrac{X_w^2+(g-w)I}{16gwX_w}\\[4pt]
 -\dfrac{X_w^{-2}}{32gw}&0
 \end{pmatrix},\qquad
 \Theta_w^{-1}=c\mathsf J_w^{-1}.
 \tag{RS4}
\end{equation}
For the upper-right equality multiply by $H_w$ and use
$(X_w^2+(g-w)I)^2=-4gwI$. The lower-left equality is the
inverse of the first identity RS3. These checks prove both
inverse products and preserve every complex phase.

At $w=0$ define $X_0=X_w|_{w=0}$ and
\begin{equation}
 N=X_0^2+gI=
 \begin{pmatrix}g&0&-g^2&0\\0&g&0&-g^2\\1&0&-g&0\\0&1&0&-g\end{pmatrix},
 \quad H_0=-2X_0N,\quad
 N^2=0,\quad\rank N=\rank H_0=2.
 \tag{RS5}
\end{equation}
The first and third columns of each parity block are proportional,
and the first two columns are independent, so $N$ has rank two.
Because $X_0$ is invertible, $H_0$ has rank two as well.
Since $X_0^2=-gI+N$ and $N^2=0$, its inverse square is
$X_0^{-2}=-(gI+N)/g^2$. Therefore RS4 gives the full matrix limit
\begin{equation}
 w\Theta_w^{-1}\longrightarrow
 c\begin{pmatrix}0&R_1\\R_2&0\end{pmatrix},\qquad
 R_1=\frac{H_0}{32g^2},\qquad R_2=\frac{gI+N}{32g^3}.
 \tag{RS6}
\end{equation}
For the first block, $H_0X_0^{-2}=-H_0/g$, since $H_0N=0$.
For the second block use the displayed inverse square. This also
proves that the residue of the inverse has rank six.

\section{Every singular constant for an arbitrary retained Gram}
Let $G$ be any fixed positive definite Hermitian four-by-four
matrix in the exact increasing powers $(1,x,x^2,x^3)$.
No parity condition is imposed. Endow $S_w$ with
$G_8=\diag(G,G)$ in the order RS1. For a four-dimensional
operator $A$, put
\begin{equation}
 K_A=G^{-1}A^*GA,\qquad
 a_A=\tr K_A,\qquad
 b_A=\frac{(\tr K_A)^2-\tr(K_A^2)}2.
 \tag{RS7}
\end{equation}
These expressions are computed from the original Gram.
$K_A$ is similar to the positive semidefinite Hermitian matrix
$(G^{1/2}AG^{-1/2})^*(G^{1/2}AG^{-1/2})$, so its eigenvalues
are the squared singular values in that Gram. For $A=N,H_0$
exactly two are positive. Thus the following are all of their
positive singular values, in decreasing order:
\begin{equation}
 \chi_{1,2}=
 \sqrt{\frac{a_N\pm\sqrt{a_N^2-4b_N}}2},\qquad
 \eta_{1,2}=
 \sqrt{\frac{a_{H_0}\pm\sqrt{a_{H_0}^2-4b_{H_0}}}2}.
 \tag{RS8}
\end{equation}
Both inner discriminants are nonnegative: for two positive
eigenvalues $\lambda_1,\lambda_2$ they equal
$(\lambda_1-\lambda_2)^2$. Their products $b_A$ are strictly
positive. Thus all radicals in RS8 are well defined, including
equality of the two singular values.

There is an exact two-by-two evaluation of every invariant in RS8.
It keeps every entry of $G$ and of $G^{-1}$:
\begin{equation}
 V=\begin{pmatrix}g&0\\0&g\\1&0\\0&1\end{pmatrix},\qquad
 R=\begin{pmatrix}1&0&-g&0\\0&1&0&-g\end{pmatrix},\qquad
 J_g=\begin{pmatrix}0&-g\\1&0\end{pmatrix},\qquad
 Q=V^*GV,\quad P=RG^{-1}R^*.
 \tag{RS9}
\end{equation}
Both $P,Q$ are positive definite. Matrix multiplication gives
$N=VR$, $RV=0$, $X_0V=VJ_g$ and $H_0=-2VJ_gR$.
The nonzero eigenvalues of products $AB$ and $BA$ agree, with
algebraic multiplicity: for a nonzero eigenvalue, $B$ maps its
generalized eigenspace injectively to that of $BA$, with inverse
obtained from the finite inverse of $AB$ on that eigenspace.
Applying this to the displayed factorizations, or taking the
coefficient identity $\det(I+zAB)=\det(I+zBA)$, yields
\begin{equation}
 \begin{aligned}
 a_N&=\tr(PQ),& b_N&=\det P\det Q,\\
 a_{H_0}&=4\tr(PJ_g^*QJ_g),&
 b_{H_0}&=16g^2\det P\det Q,\qquad
 \eta_1\eta_2=4g\chi_1\chi_2.
 \end{aligned}
 \tag{RS10}
\end{equation}
For example $\det J_g=g$, so the last determinant factor is
$4^2|\det J_g|^2=16g^2$. These are finite two-by-two formulas,
not unresolved singular-value problems. To make their entry use
explicit, indices $0,1,2,3$ refer to the original monomials and
\begin{equation}
 \begin{aligned}
 Q_{11}&=g^2G_{00}+g(G_{02}+G_{20})+G_{22},\\
 Q_{12}&=g^2G_{01}+gG_{03}+gG_{21}+G_{23},\\
 Q_{22}&=g^2G_{11}+g(G_{13}+G_{31})+G_{33},\\
 P_{11}&=(G^{-1})_{00}-g((G^{-1})_{02}+(G^{-1})_{20})
                                      +g^2(G^{-1})_{22},\\
 P_{12}&=(G^{-1})_{01}-g(G^{-1})_{03}-g(G^{-1})_{21}
                                      +g^2(G^{-1})_{23},\\
 P_{22}&=(G^{-1})_{11}-g((G^{-1})_{13}+(G^{-1})_{31})
                                      +g^2(G^{-1})_{33},
 \quad P_{21}=\overline{P_{12}},\quad Q_{21}=\overline{Q_{12}}.
 \end{aligned}
 \tag{RS11}
\end{equation}

We next evaluate all four singular values of $R_2$.
Let $\mathcal L=\im N$. Its dimension is two and $N^2=0$,
so $\mathcal L=\ker N$. The $G$-orthogonal decomposition
$\C^4=\mathcal L\oplus\mathcal L^{\perp_G}$ gives
\begin{equation}
 N=\begin{pmatrix}0&S\\0&0\end{pmatrix},\qquad
 gI+N=\begin{pmatrix}gI&S\\0&gI\end{pmatrix},
 \tag{RS12}
\end{equation}
where $S:\mathcal L^{\perp_G}\to\mathcal L$ is invertible.
Use orthonormal bases on these two actual metric subspaces.
Diagonalize the positive matrix $S^*S$, then take the normalized
images under $S$ as the basis on $\mathcal L$. This makes $S$
diagonal with positive entries $\chi_1,\chi_2$ by isometries.
It does not alter the original metric. After interleaving the
two pairs, $gI+N$ is the orthogonal direct sum of
$\left(\begin{smallmatrix}g&\chi_j\\0&g\end{smallmatrix}\right)$.
The squared singular-value sum of that two-by-two matrix is
$2g^2+\chi_j^2$, and its singular-value product is $g^2$.
The two positive numbers
$\bigl(\sqrt{4g^2+\chi_j^2}\pm\chi_j\bigr)/2$
have exactly that sum of squares and product, so they are
its two singular values. Consequently the four constants are
\begin{equation}
 \zeta_{j,\pm}
 =\frac{\sqrt{4g^2+\chi_j^2}\pm\chi_j}{64g^3}>0,
 \qquad j=1,2.
 \tag{RS13}
\end{equation}

Form the six positive numbers and the two finite constants
\begin{equation}
 \mathcal C=
 \left\{\frac{\eta_1}{32g^2},\frac{\eta_2}{32g^2},
 \zeta_{1,+},\zeta_{1,-},\zeta_{2,+},\zeta_{2,-}\right\},
 \qquad d_1=\frac1{2\eta_2},\quad d_2=\frac1{2\eta_1}.
 \tag{RS14}
\end{equation}
Write $c_1\ge\cdots\ge c_6$ for the decreasing ordering of
the displayed six numbers; $d_1\ge d_2$. Every ordered singular
value of $\Theta_w^{-1}$ in $G_8$ is now determined:
\begin{equation}
 \sigma_j(\Theta_w^{-1})=\frac{c_j}{w}(1+o(1)),\quad 1\le j\le6,
 \qquad
 \sigma_7(\Theta_w^{-1})=d_1(1+o(1)),\quad
 \sigma_8(\Theta_w^{-1})=d_2(1+o(1)).
 \tag{RS15}
\end{equation}
Here is a proof covering the finite constants as well as the
divergent ones. In $\diag(G,G)$, the squared singular spectrum
of a block matrix $\left(\begin{smallmatrix}0&A\\B&0\end{smallmatrix}\right)$
is the union of the squared singular spectra of $A$ and $B$,
because its metric adjoint times itself is block diagonal.
For the upper-right block of RS4 the singular values are the
reciprocals of twice the singular values of $H_w$. The two
largest singular values of $H_w$ tend to $\eta_1,\eta_2$,
so two inverse singular values tend to $d_2,d_1$.
Multiplying this block by $w$ gives $R_1$, which has precisely
two positive singular values $\eta_j/(32g^2)$. Thus its other
two singular values diverge with those constants. The lower-left
block times $w$ tends to the invertible $R_2$, whose four
singular values are RS13. Singular values are continuous under
matrix convergence: the minimum-maximum description of the
singular values gives their change bounded by the operator norm
of the difference. Finally $|c|=1$, so its unchanged phase in
RS4 changes none of these singular values. This proves RS15.

\section{Every exterior rank and the exact determinant}
The exterior norm in the metric induced by $G_8$ is the product
of the largest singular values. Indeed orthonormal singular
bases induce orthonormal wedge bases, and the matrix entries
on those bases are all corresponding products. Define
\begin{equation}
 C_0=1,\qquad C_r=\prod_{j=1}^r c_j
 =\max_{\substack{I\subset\{1,\ldots,6\}\\|I|=r}}
                         \prod_{j\in I}\mathcal C_j,
 \quad1\le r\le6.
 \tag{RS16}
\end{equation}
The maximum is a finite explicit expression in the radicals
RS8 and RS13 and makes no assumption about their ordering for
different original moments. For $r=6$ the product evaluates to
\begin{equation}
 \zeta_{j,+}\zeta_{j,-}=\frac1{2^{10}g^4},\qquad
 C_6=\frac{\eta_1\eta_2}{2^{30}g^{12}}
     =\frac{\sqrt{b_{H_0}}}{2^{30}g^{12}},\qquad
 C_6d_1=\frac{\eta_1}{2^{31}g^{12}}.
 \tag{RS17}
\end{equation}
The first equality follows by subtracting the two squares in
the numerator of RS13. Multiplying its two instances and
the first two entries of RS14 proves the others.

The determinant is also exact for every $w>0$:
\begin{equation}
 \det H_w=256g^2w^2(g+w)^2,\qquad
 \det\Theta_w=2^{32}g^6w^6(g+w)^6,\qquad
 \det\Theta_w^{-1}=\frac1{2^{32}g^6w^6(g+w)^6}.
 \tag{RS18}
\end{equation}
For the first formula, the roots of the monic quartic in RS1
are $\varepsilon\sqrt w+i\eta\sqrt g$, with independent signs.
At each root, the factor $x^2+g-w$ equals
$2i\varepsilon\eta\sqrt{gw}$. Their product is $16g^2w^2$;
$\det(-2X_w)=16(g+w)^2$. Multiplication gives the first
identity. The determinant of the two swapped four-by-four
blocks of $\mathsf J_w$ is $2^8(\det H_w)^3$, with positive
block-swap sign. Also $c^8=1$. This proves both remaining
identities, including their phases.

Combining RS15--18 gives every inverse exterior asymptote:
\begin{equation}
 \boxed{\begin{aligned}
 \|\wedge^0\Theta_w^{-1}\|_{G_8}&=1,\\
 \|\wedge^r\Theta_w^{-1}\|_{G_8}&\sim C_rw^{-r},&&1\le r\le6,\\
 \|\wedge^7\Theta_w^{-1}\|_{G_8}&\sim
                       \frac{\eta_1}{2^{31}g^{12}}w^{-6},\\
 \|\wedge^8\Theta_w^{-1}\|_{G_8}&=
                       \frac1{2^{32}g^6w^6(g+w)^6}.
 \end{aligned}}
 \tag{RS19}
\end{equation}
The complete leading product is
$C_6d_1d_2=1/(2^{32}g^{12})$, agreeing with RS18.
For the original variable $\delta$, every exponent is
$2\min(r,6)$, with exactly the constants displayed above.

\section{The exact Gram furnished by the original source}
We now evaluate $G$, rather than leave it as an auxiliary
positive matrix. Keep the original conductor $L=T_A$ on
\begin{equation}
 U=\mathcal P_{v+3}/\mathcal P_{v-1},\qquad
 W=\mathcal P_3,\qquad
 c_*=a_{\rm amp}/2-4,\quad c_0=c_*+4,\quad
 b=s/2>0,\quad M=(2\pi)^{s/2}.
 \tag{RS20}
\end{equation}
The source has its minimum quotient norm at center $c_0$;
the target has its original Gamma norm at center $c_*$.
Both use the same fixed original order $s$ in the allowed
programme domain. In target coordinates
$y=(S'-c_*)/i$, $x=S'-c_*=iy$.
Thus the spectral coordinate $x=r-1/2$ is sent by the actual
dictionary $r=S'-\kappa$, $\kappa=c_*-1/2$, to exactly
$S'-c_*$. No center is set to zero in the original spaces.

Write $\mu_n=E_A^{(n)}(0)$ and let $v$ be its actual
vanishing order, so $\mu_j=0$ for $j<v$ and $\mu_v\ne0$.
To distinguish the conductor coefficients from $g=\gamma^2$,
use $\mathfrak g_j$ for the coefficients of its exact FC20 unit:
\begin{equation}
 \begin{gathered}
 z(q)=-i\arctan q,\qquad
 \mathfrak g(q)=q^{-v}e^{-4z(q)}E_A(z(q)),\\
 \nu_{v+j}=\sum_{k=0}^{j}\binom{v+j}{v+k}(-4)^{j-k}\mu_{v+k},\\
 \mathfrak g_0=(-i)^v\mu_v/v!,\qquad
 \mathfrak g_1=(-i)^{v+1}\nu_{v+1}/(v+1)!,\\
 \mathfrak g_2=(-i)^{v+2}\nu_{v+2}/(v+2)!-(v/3)\mathfrak g_0,\\
 \mathfrak g_3=(-i)^{v+3}\nu_{v+3}/(v+3)!-((v+1)/3)\mathfrak g_1,\\
 \xi_0=\mathfrak g_0^{-1},\quad
 \xi_1=-\mathfrak g_1\mathfrak g_0^{-2},\quad
 \xi_2=\mathfrak g_1^2\mathfrak g_0^{-3}-\mathfrak g_2\mathfrak g_0^{-2},\\
 \xi_3=-\mathfrak g_1^3\mathfrak g_0^{-4}
       +2\mathfrak g_1\mathfrak g_2\mathfrak g_0^{-3}
       -\mathfrak g_3\mathfrak g_0^{-2}.
 \end{gathered}
 \tag{RS21}
\end{equation}
The finite moment formula for $\nu$ follows by multiplying
$e^{-4z}E_A(z)$ as power series. Since
$z(q)^n=(-i)^nq^n(1-nq^2/3+O(q^4))$, it gives every
$\mathfrak g_j$ displayed through degree three. Multiplying
the two cubic truncations $\sum\mathfrak g_jq^j$ and
$\sum\xi_jq^j$ gives $1$ modulo $q^4$, proving all inverse
coefficients with their original phases.

For completeness the original orthonormal source polynomials
and their minimum quotient representatives are
\begin{equation}
 \begin{gathered}
 p_0(y)=1,\quad p_1(y)=y,\quad
 p_{n+1}(y)=yp_n(y)-n(n-1+b)p_{n-1}(y),\\
 \phi_n(S)=\frac{p_n((S-c_0)/i)}{\sqrt{Mn!(b)_n}},\qquad
 [\phi_v],\ldots,[\phi_{v+3}]\text{ an orthonormal frame of }U,\\
 \rho_n=\sqrt{n!/(b)_n},\qquad
 k_j=\rho_{v+j}^{-2}=(b)_{v+j}/(v+j)!,\quad0\le j\le3.
 \end{gathered}
 \tag{RS22}
\end{equation}
The orthogonality and original measure are the fixed Gamma
dictionary of \cite{WCF,DLMF18Source}; its squared norms include $M$.
Each listed polynomial is orthogonal to $\mathcal P_{v-1}$.
Consequently Pythagoras proves that its linear combinations
are precisely the minimum-norm representatives of their classes.

In the original source and target orthonormal frames FC24 has
entry $\xi_{j-r}\rho_j/\rho_{v+r}$ for $0\le r\le j\le3$.
To evaluate its action on monomials, the target orthonormal
coordinate columns of $(1,y,y^2,y^3)$ are
\begin{equation}
 V_W=\begin{pmatrix}
 \sqrt M&0&b\sqrt M&0\\
 0&\sqrt{Mb}&0&(3b+2)\sqrt{Mb}\\
 0&0&\sqrt{2Mb(b+1)}&0\\
 0&0&0&\sqrt{6Mb(b+1)(b+2)}
 \end{pmatrix}.
 \tag{RS23}
\end{equation}
Indeed $p_2=y^2-b$ and $p_3=y^3-(3b+2)y$ follow
from RS22. Multiplying these columns by the complete FC24
inverse, and using
$\rho_j\sqrt{Mj!(b)_j}=\sqrt M\,j!$, gives
\begin{equation}
 \begin{gathered}
 F=\begin{pmatrix}
 \xi_0&\xi_1&b\xi_0+2\xi_2&(3b+2)\xi_1+6\xi_3\\
 0&\xi_0&2\xi_1&(3b+2)\xi_0+6\xi_2\\
 0&0&2\xi_0&6\xi_1\\0&0&0&6\xi_0
 \end{pmatrix},\quad D_i=\diag(1,i,-1,-i),\\
 U_x=\sqrt M\diag(\sqrt{k_0},\sqrt{k_1},\sqrt{k_2},\sqrt{k_3})FD_i,\\
 L^{-1}(S'-c_*)^j=\sum_{r=0}^3(U_x)_{rj}[\phi_{v+r}],\qquad
 G_{U,x}=U_x^*U_x=M D_i^*F^*\diag(k_0,k_1,k_2,k_3)FD_i.
 \end{gathered}
 \tag{RS24}
\end{equation}
The additional $D_i$ is required because $x^j=i^jy^j$.
Thus $G_{U,x}$ is the full original source Gram in precisely
the $x$ coefficients of RS1. The subscript records this receiver;
the programme Gram $\Phi^*\langle\, ,\,\rangle_U\Phi$ in
original ES coordinates keeps its existing notation.

Its entries and inverse can be evaluated without a numerical
matrix inversion. In addition to
$(G_{U,x})_{pq}=M(-i)^p i^q
\sum_{r=0}^{\min(p,q)}k_r\overline{F_{rp}}F_{rq}$,
the complete inverse of $F$ is
\begin{equation}
 B_F=\begin{pmatrix}
 \mathfrak g_0&\mathfrak g_1&\mathfrak g_2-b\mathfrak g_0/2
                       &\mathfrak g_3-b\mathfrak g_1/2\\
 0&\mathfrak g_0&\mathfrak g_1
                       &\mathfrak g_2-(3b+2)\mathfrak g_0/6\\
 0&0&\mathfrak g_0/2&\mathfrak g_1/2\\
 0&0&0&\mathfrak g_0/6
 \end{pmatrix},\qquad
 G_{U,x}^{-1}=M^{-1}D_i^*B_F\diag(k_j^{-1})B_F^*D_i.
 \tag{RS25}
\end{equation}
To prove the inverse, write
$F=T_3(\xi)V_0$, where $V_0$ has diagonal $(1,1,2,6)$,
entry $(0,2)$ equal to $b$ and entry $(1,3)$ equal to $3b+2$.
All other off-diagonal entries vanish. The reciprocal identity
in RS21 gives $T_3(\xi)^{-1}=T_3(\mathfrak g)$, and direct
upper-triangular multiplication gives $V_0^{-1}T_3(\mathfrak g)=B_F$.
Thus $FB_F=B_FF=I$, proving RS25. All original cross terms
remain in both products RS24--25. Positivity and invertibility
also follow from the exact determinant
\begin{equation}
 \det G_{U,x}=144M^4 k_0k_1k_2k_3|\xi_0|^8>0.
 \tag{RS26}
\end{equation}
In fact all four leading minors are respectively
$Mk_0|\xi_0|^2$, $M^2k_0k_1|\xi_0|^4$,
$4M^3k_0k_1k_2|\xi_0|^6$ and RS26, by the triangular
diagonal of $F$. The phases in $D_i$ have modulus one.

Substituting RS24--25 into RS9--11 and then RS8, RS13--19
therefore evaluates all eight source singular constants in
finite original moments, Gamma weights and real radicals.
Every matrix used in this substitution has all its entries
displayed. The common $M$ cancels from $G_{U,x}^{-1}A^*G_{U,x}A$
because it occurs on both sides of the same operator norm;
it remains unchanged in the source norms RS24 and RS26.

\section{The exact conductor receiver and the completed estimate}
Define the coefficient receivers, on each of the two parity
summands, by
\begin{equation}
 C_Wp=p(S'-c_*),\qquad C_Up=L^{-1}p(S'-c_*),\qquad
 \mathcal C_W=C_W\oplus C_W,\quad
 \mathcal C_U=C_U\oplus C_U.
 \tag{RS27}
\end{equation}
They are isomorphisms from the coefficient space of $S_w$
to $W^2$ and $U^2$ respectively; the inverse of $C_W$ is
$q(S')\mapsto q(x+c_*)$, and that of $C_U$ is
$u\mapsto (Lu)(x+c_*)$. The full inverse matrix in the original
source frame is supplied by RS24--25. The original relation is
$(L\oplus L)\mathcal C_U=\mathcal C_W$.
With $G=G_{U,x}$ in RS7--19, $\mathcal C_U$ is an isometry
from $\diag(G_{U,x},G_{U,x})$ onto the original quotient direct sum,
by its proved Gram RS24. The transported operation and its
inverse are consequently
\begin{equation}
 \Theta_{U,w}=\mathcal C_U\Theta_w\mathcal C_U^{-1},\quad
 \Theta_{U,w}^{-1}=\mathcal C_U\Theta_w^{-1}\mathcal C_U^{-1},\quad
 (L\oplus L)\Theta_{U,w}=\Theta_{W,w}(L\oplus L),
 \quad\Theta_{W,w}=\mathcal C_W\Theta_w\mathcal C_W^{-1}.
 \tag{RS28}
\end{equation}
Substitution and cancellation of the stated inverse maps prove
all identities. The singular constants of $\Theta_{U,w}^{-1}$
are exactly RS14--15 evaluated at the explicit RS24 Gram,
and every original-source exterior norm is exactly RS19 with
that substitution. No condition-number bound replaces a source
constant. The determinant RS18 is unchanged by this similarity.

As a direct equation-(11) consequence, for every rank $1\le r\le8$
the product
\begin{equation}
 w^{\min(r,6)}\|\wedge^r\Theta_{U,w}^{-1}\|_U
 \longrightarrow
 \begin{cases}
 C_r(G_{U,x}),&1\le r\le6,\\
 \eta_1(G_{U,x})/(2^{31}g^{12}),&r=7,\\
 1/(2^{32}g^{12}),&r=8.
 \end{cases}
 \tag{RS29}
\end{equation}
Every limit is strictly positive by RS8 and RS13. Hence this
specific completed trace-to-residue operation has no uniform
inverse exterior bound as the paired collision is approached.
The locus of this loss is the explicit multiplier RS3, whose
rank at $w=0$ is two. The original $L$ remains the fixed
isomorphism with determinant FC23; its nonzero leading moment
$\mu_v$ is retained in RS21--26. Equations RS27--28 are the exact
map between these two statements and do not change the original
conductor's order.


\section{The same singular constants in the complete marked receiver}
We now prove the exact change of coordinates from RS27 to RD19,
including the mixed metric. The original reference period and its
distinct quartet are fixed throughout. Its ordered labels remain
$(\mathrm M,\mathrm N,\mathrm T,\mathrm E)$, with both signs at
each label. The varying quartic is RS1; its parameter $w=\delta^2$
does not replace that reference period. The marked frames
$K,X,Y,\Phi,\Psi,E,O$ are exactly FC26--32 and SF2--3.
In particular $Y=\Psi K$, $X=\Phi K$, $Y=T_AX$ and
$\Psi=T_A\Phi$. All their original phases and entries are retained.
The human polynomial provenance remains Alp\"oge's announcement
and Tao's account \cite{Tao}; the four labels and their unchanged
marking come from the ES source \cite{ES488}.

For a scalar $a$, define the complete cubic substitution matrix
\begin{equation}
 C_a=\begin{pmatrix}1&a&a^2&a^3\\0&1&2a&3a^2\\
 0&0&1&3a\\0&0&0&1\end{pmatrix},\qquad
 [p(x+a)]=C_a[p(x)],\quad C_aC_b=C_{a+b},\quad C_a^{-1}=C_{-a}.
 \tag{MR1}
\end{equation}
The coefficient of $x^i$ in $(x+a)^j$ is
$\binom ji a^{j-i}$, proving the first identity. Substitution
twice proves the composition and inverse without discarding
any original center. Retain $c_*=a_{\rm amp}/2-4$ and
$\kappa=c_*-1/2$. In RS1 the coordinate is $x=r-1/2$, so
the full $r$-coefficient matrix $Q_r$ of RD19 becomes
\begin{equation}
 \begin{split}
 Z&=C_{c_*}\Psi,\\
 Q_x&=\operatorname{diag}(C_{1/2},C_{1/2})Q_r
   =\begin{pmatrix}C_{c_*}Y&0\\C_{c_*}\Psi E&C_{c_*}\Psi O\end{pmatrix}
   =\operatorname{diag}(Z,Z)\begin{pmatrix}K&0\\E&O\end{pmatrix},\\
 Q_x^{-1}&=\begin{pmatrix}
 K^{-1}Z^{-1}&0\\-O^{-1}EK^{-1}Z^{-1}&O^{-1}Z^{-1}
 \end{pmatrix},\qquad
 \det Q_x=(\det\Psi)^2\det K\det O\ne0.
 \end{split}\tag{MR2}
\end{equation}
Indeed $C_{1/2}C_\kappa=C_{c_*}$ and $Y=\Psi K$.
Multiplication of the two displayed block matrices proves
both inverse products. The original odd frame $O$ is
invertible by SE6; all remaining factors are the original
FC isomorphisms. These facts prove the last assertion with
the actual determinant, rather than selecting a new frame.

Use the two coefficient receivers $\mathcal C_U,\mathcal C_W$
of RS27. The exact identities with the original marked maps are
\begin{equation}
 \mathcal J_U=\mathcal C_UQ_x,\qquad
 \mathcal J_W=\mathcal C_WQ_x,\qquad
 (T_A\oplus T_A)\mathcal J_U=\mathcal J_W.
 \tag{MR3}
\end{equation}
For example, $C_WC_{c_*}Y\alpha=Y\alpha$ since
$(C_{c_*}p)(S'-c_*)=p(S')$. Applying the original $T_A^{-1}$
gives $X\alpha$. The second component gives
$\Phi(E\alpha+O\beta)$ by the same substitutions. This proves
each component of MR3 and both directions of the correspondence.

Write $G_{U,x}$ for the complete RS24 source Gram in the
$x$-coefficient basis and $G_{W,x}$ for the corresponding
original target Gram. Reserve $G_U=\Phi^*\langle\ ,\ \rangle_U\Phi$
for its earlier marked meaning; these two matrices are related by
\begin{equation}
 \begin{gathered}
 G_U=Z^*G_{U,x}Z,\quad G_X=K^*G_UK,\qquad
 G=Z^*G_{W,x}Z,\quad G_Y=K^*GK,\\
 \mathbb H_U=Q_x^*\operatorname{diag}(G_{U,x},G_{U,x})Q_x
 =\begin{pmatrix}
 G_X+E^*G_UE&E^*G_UO\\O^*G_UE&O^*G_UO
 \end{pmatrix},\\
 \mathbb H_W=Q_x^*\operatorname{diag}(G_{W,x},G_{W,x})Q_x
 =\begin{pmatrix}
 G_Y+E^*GE&E^*GO\\O^*GE&O^*GO
 \end{pmatrix}.
 \end{gathered}\tag{MR4}
\end{equation}
The first line follows from $C_UZ=\Phi$, $C_UZK=X$ and
their target counterparts, proved in MR3. Multiplication
of the complete block matrices then proves every remaining
entry. In particular the mixed terms $E^*G_UO$ and $O^*G_UE$
are present; neither parity orthogonality nor a diagonal
marked metric was assumed. The full moment and Gamma formula
for $G_{U,x}$, its inverse and all phases is RS21--26.

Let $\Theta_w=M_{2t^3}/c$ be RS3 with its exact phase
$c=(1+i)/\sqrt2$, and put
\begin{equation}
 \Theta_{{\rm mark},w}=Q_x^{-1}\Theta_wQ_x,\qquad
 \Theta_{{\rm mark},w}^{-1}=Q_x^{-1}\Theta_w^{-1}Q_x.
 \tag{MR5}
\end{equation}
All maps in this identity act on the same eight signed
coordinates as RD19. Their physical original-source and
target forms are precisely $\mathcal J_U\Theta_{{\rm mark},w}
\mathcal J_U^{-1}$ and $\mathcal J_W\Theta_{{\rm mark},w}
\mathcal J_W^{-1}$. MR3 proves their intertwining by
$T_A\oplus T_A$ in both directions.

Here is the full singular-value argument including the metric.
For any invertible $Q$, positive $G_8$, and matrix $A$, set
$H=Q^*G_8Q$ and $A_Q=Q^{-1}AQ$. Direct inversion and
multiplication give
\begin{equation}
 H^{-1}A_Q^*HA_Q
       =Q^{-1}G_8^{-1}A^*G_8AQ.
 \tag{MR6}
\end{equation}
Both sides therefore have the same eigenvalues. Those eigenvalues
are the squared singular values in their respective positive
metrics, because conjugation by the positive square root of
the metric gives the usual adjoint product. Thus every singular
value is retained. Exterior norms are retained as well: the
Gram of wedges is the determinant of the individual pairings,
and an isometry preserves each of these pairings. With MR4,
this proves, for every $w>0$ and $0\le j\le8$,
\begin{equation}
 \left\|\bigwedge^j\Theta_{{\rm mark},w}^{-1}\right\|_{\mathbb H_U}
 =\left\|\bigwedge^j\Theta_w^{-1}\right\|_{
             \operatorname{diag}(G_{U,x},G_{U,x})}
 =\left\|\bigwedge^j\Theta_{U,w}^{-1}\right\|_{U\oplus U}.
 \tag{MR7}
\end{equation}
Consequently RS14--19 give the exact marked constants, with
the explicit original-source substitution RS24--25. The
target statement follows by the same proved identity with
$G_{W,x}$. No condition-number estimate replaces either
source or target constants in MR7.

\section{The complete nilpotent flag of the inverse pole}
Keep $g=\gamma^2>0$, $X_0,N,H_0$ of RS5, and let
\begin{equation}
 \begin{gathered}
 v_1=(g,0,1,0)^{\mathsf T},\quad
 v_2=(0,g,0,1)^{\mathsf T},\quad
 V=\operatorname{span}\{v_1,v_2\}=\operatorname{im}N=\ker N,\\
 \mathscr L=\lim_{w\downarrow0}w\Theta_w^{-1}
 =c\begin{pmatrix}0&H_0/(32g^2)\\(gI+N)/(32g^3)&0\end{pmatrix},\\
 \Theta_0=\frac2c\begin{pmatrix}0&0\\H_0&0\end{pmatrix}.
 \end{gathered}\tag{MR8}
\end{equation}
The first two columns of $N$ are $v_1,v_2$; its remaining
columns are $-gv_1,-gv_2$. They are independent, and $N^2=0$,
proving both descriptions of $V$. Its polynomials are exactly
$x^2+g$ and $x(x^2+g)$, with their original coefficients.
Since $X_0$ commutes with $N$ and is invertible,
$\operatorname{im}H_0=\ker H_0=V$ and $H_0N=NH_0=0$.
Multiplication of the unchanged matrices gives
\begin{equation}
 \mathscr L^2=\frac{c^2}{2^{10}g^4}
            \begin{pmatrix}H_0&0\\0&H_0\end{pmatrix},\qquad
 \mathscr L^3=\frac{c^3}{2^{15}g^6}
            \begin{pmatrix}0&0\\H_0&0\end{pmatrix}
       =-\frac{\Theta_0}{2^{16}g^6},\qquad
 \mathscr L^4=0.
 \tag{MR9}
\end{equation}
For the first product use $H_0(gI+N)=gH_0$; for the second
also use $H_0^2=0$. In comparing with $\Theta_0$ the factor
is $c^4/2^{16}g^6=-1/2^{16}g^6$, so the phase is essential.

The complete image and kernel flags, in even then odd order, are
\begin{equation}
 \begin{array}{c|c|c|c}
 j&\operatorname{im}\mathscr L^j&\ker\mathscr L^j&
     \operatorname{rank}\mathscr L^j\\\hline
 0&\C^4\oplus\C^4&0&8\\
 1&V\oplus\C^4&0\oplus V&6\\
 2&V\oplus V&V\oplus V&4\\
 3&0\oplus V&V\oplus\C^4&2\\
 4&0&\C^4\oplus\C^4&0
 \end{array}
 \tag{MR10}
\end{equation}
Indeed $(gI+N)^{-1}=g^{-1}I-g^{-2}N$ since $N^2=0$,
whereas $H_0$ has precisely image and kernel $V$. This proves
the first row of nontrivial subspaces directly from MR8;
MR9 proves the others. In particular
$\ker\Theta_0=V\oplus\C^4$ and
$\operatorname{im}\Theta_0=0\oplus V$.

There are two explicit complete chains. Let $e_0,e_1$
denote the first two cubic coefficient vectors, and set
$u_j=(e_j,0)$, for $j=0,1$. Without rescaling any vector they are
\begin{equation}
 \begin{split}
 u_j&=(e_j,0),\\
 \mathscr L u_j&=\left(0,\frac{c(gI+N)e_j}{32g^3}\right),\\
 \mathscr L^2u_j&=\left(\frac{c^2H_0e_j}{1024g^4},0\right),\\
 \mathscr L^3u_j&=\left(0,\frac{c^3H_0e_j}{32768g^6}\right),\\
 H_0e_0&=-2v_2,\qquad H_0e_1=2gv_1.
 \end{split}\tag{MR11}
\end{equation}
Both last vectors are independent. In any linear relation
between these eight vectors, applying $\mathscr L^3$ first
forces the coefficients of $u_0,u_1$ to vanish. Applying
$\mathscr L^2$ then forces the coefficients of
$\mathscr L u_0,\mathscr L u_1$ to vanish; applying
$\mathscr L$ and then using the independence of the last
vectors eliminates the remaining coefficients. Thus the
ordered columns
$(\mathscr L^3u_0,\mathscr L^2u_0,\mathscr L u_0,u_0,
  \mathscr L^3u_1,\mathscr L^2u_1,\mathscr L u_1,u_1)$
are a basis and exhibit exactly two blocks of length four.

Every vector and subspace in this flag returns to the original
marked and physical spaces by the explicit isomorphisms
\begin{equation}
 \mathscr L_{\rm mark}=Q_x^{-1}\mathscr LQ_x,
 \quad \operatorname{im}\mathscr L_{\rm mark}^j
       =Q_x^{-1}\operatorname{im}\mathscr L^j,
 \quad \ker\mathscr L_{\rm mark}^j
       =Q_x^{-1}\ker\mathscr L^j,
 \quad \mathcal J_UQ_x^{-1}=\mathcal C_U.
 \tag{MR12}
\end{equation}
The corresponding target map is $\mathcal J_WQ_x^{-1}
=\mathcal C_W$. Applying $T_A\oplus T_A$ to the source
chain gives the target chain term by term. MR2 supplies the
entire inverse, so no marked state or phase is omitted.
For any chain column matrix $B_j$, its actual source Gram
is exactly $B_j^*\operatorname{diag}(G_{U,x},G_{U,x})B_j$,
which equals $(Q_x^{-1}B_j)^*\mathbb H_U(Q_x^{-1}B_j)$
by MR4. This also proves the metric statement for every flag.

\section{Exact finite-parameter inverse bounds with the original metric}
The singular asymptotes admit a fully evaluated error term.
For every $w>0$, the matrices RS4 satisfy
\begin{equation}
 \Theta_w^{-1}=\frac{\mathscr L}{w}
       +c\begin{pmatrix}0&A_w\\B_w&0\end{pmatrix},
 \tag{MR13}
\end{equation}
where every entry of the remainder is
\begin{equation}
 A_w=\begin{pmatrix}
 0&-\frac1{16g}&0&-\frac{2g+w}{16g}\\
 \frac1{2(g+w)^2}&0&-\frac1{16g}&0\\
 0&0&0&\frac1{16g}\\
 \frac{3g+w}{16g^2(g+w)^2}&0&0&0
 \end{pmatrix},\qquad
 B_w=\begin{pmatrix}
 -\frac{3g+w}{16g^2(g+w)^2}&0&0&0\\
 0&-\frac{3g+w}{16g^2(g+w)^2}&0&0\\
 -\frac{2g+w}{32g^3(g+w)^2}&0&0&0\\
 0&-\frac{2g+w}{32g^3(g+w)^2}&0&0
 \end{pmatrix}.
 \tag{MR14}
\end{equation}
For a direct verification put $a=g-w$, $p=(g+w)^2$.
The exact inverse of RS2 is
\begin{equation}
 X_w^{-1}=\begin{pmatrix}
 0&1&0&0\\-2a/p&0&1&0\\0&0&0&1\\-1/p&0&0&0
 \end{pmatrix}.
 \tag{MR15}
\end{equation}
Multiplying both products with $X_w$ proves this inverse.
Subtracting $H_0/(32g^2w)$ from
$(X_w+aX_w^{-1})/(16gw)$ gives exactly $A_w$.
Subtracting $(gI+N)/(32g^3w)$ from
$-X_w^{-2}/(32gw)$ gives exactly $B_w$.
These are entrywise calculations and prove MR13--14, including
the factors of $g$, $w$ and $c$. Every entry extends continuously
to $w=0$ for the unchanged $g>0$.

For the original source set $G_0=G_{U,x}$ with its complete
RS24--25 evaluation, and define the following explicit finite
quantities:
\begin{equation}
 \begin{split}
 e(w)^2&=\operatorname{tr}(G_0^{-1}A_w^*G_0A_w)
            +\operatorname{tr}(G_0^{-1}B_w^*G_0B_w),\\
 b(w)^2&=4\operatorname{tr}(G_0^{-1}(H_w^2)^*G_0H_w^2)
            +4\operatorname{tr}(G_0^{-1}H_w^*G_0H_w),\\
 p_j(w)&=\max\{0,c_j-we(w)\},\qquad
 q_j(w)=c_j+we(w),\quad 1\le j\le6 .
 \end{split}\tag{MR16}
\end{equation}
The $c_j$ are the six explicitly evaluated RS14 constants
for this $G_0$. All entries in MR16 are displayed in
RS2, RS24--25 and MR14. The two traces in each line are
sums of squares in the original positive metric. Hence $e(w)$
is the full Hilbert--Schmidt norm of the remainder and $b(w)$
is that of $\Theta_w$. Both are finite, and $b(w)>0$ for
every $w>0$. The phase has modulus one and remains present
in MR13 even though it cancels from these norms.

The variational characterization of singular values gives
$|\sigma_j(A)-\sigma_j(B)|\le\|A-B\|$: apply the triangle
inequality on each subspace in the minimum-maximum formula,
and then reverse $A,B$. Apply this to $w\Theta_w^{-1}$ and
$\mathscr L$. Its first six singular values are $c_j$, and
its last two are zero. The operator norm is bounded by the
Hilbert--Schmidt norm, so MR13 proves
\begin{equation}
 p_j(w)\le w\sigma_j(\Theta_{{\rm mark},w}^{-1})\le q_j(w)
 \quad(1\le j\le6),\qquad
 b(w)^{-1}\le\sigma_j(\Theta_{{\rm mark},w}^{-1})\le e(w)
 \quad(j=7,8).
 \tag{MR17}
\end{equation}
For the lower bound on the last two, the smallest singular
value of an inverse is the reciprocal of the largest singular
value of the original. This follows by inverting orthonormal
singular bases. The latter largest value is at most $b(w)$.
MR7 transfers all these inequalities to the actual marked
metric without any comparison loss.

Multiplying the decreasing singular values yields, for every
$w>0$, the complete inverse exterior bounds
\begin{equation}
 \begin{aligned}
 \frac{\prod_{j=1}^r p_j(w)}{w^r}
 &\le\|\bigwedge^r\Theta_{{\rm mark},w}^{-1}\|_{\mathbb H_U}
 \le\frac{\prod_{j=1}^r q_j(w)}{w^r},&&1\le r\le6,\\
 \frac{\prod_{j=1}^6p_j(w)}{w^6b(w)^{r-6}}
 &\le\|\bigwedge^r\Theta_{{\rm mark},w}^{-1}\|_{\mathbb H_U}
 \le\frac{\prod_{j=1}^6q_j(w)e(w)^{r-6}}{w^6},&&r=7,8,\\
 \|\bigwedge^8\Theta_{{\rm mark},w}^{-1}\|_{\mathbb H_U}
 &=\frac1{2^{32}g^6w^6(g+w)^6}.
 \end{aligned}\tag{MR18}
\end{equation}
Rank zero has norm one. Continuity of MR14 at zero implies
$we(w)\to0$; all $c_j$ are positive by RS8 and RS13.
Thus the lower bounds on the first six factors become positive
near zero and approach their exact constants. RS19 supplies
the two remaining exact leading constants. Equations MR7 and
MR18 are the full original-metric conclusion for the completed
trace-to-residue operation, at finite parameters as well as in
the limit.

The loss occurs in the explicit multiplication $2t^3/c$, whose
critical rank is two and whose inverse has the chains MR11.
The original conductor itself is the fixed isomorphism FC23--24.
MR3, MR5 and MR12 prove the exact relation to that isomorphism;
they do not replace its moment $\mu_v$ or change its order.

\section{Uniform original-metric error and the two finite singular values}
The exact remainder MR14 also gives bounds with constants
independent of $w$ on a specified interval. Keep $g>0$ unchanged,
put $G_0=G_{U,x}$ with all its original entries from RS24--25,
and define
\begin{equation}
 \begin{gathered}
 \kappa_{G_0}=
   \sqrt{\operatorname{tr}G_0\operatorname{tr}G_0^{-1}},\\
 B(g)=\frac9{256}+\frac3{256g^2}+\frac1{4g^4}
             +\frac3{16g^6}+\frac9{512g^8},\\
 K_{G_0}=\kappa_{G_0}\sqrt{B(g)},\qquad
 D_{G_0}=\kappa_{G_0}\sqrt{4g^4+216g^2+12}.
 \end{gathered}
 \tag{MR19}
\end{equation}
All quantities are finite positive expressions in the retained
original moments, Gamma weights and $g$. For an endomorphism
$A$ of the cubic coefficient space,
\[
 \|A\|_{G_0}
 =\|G_0^{1/2}AG_0^{-1/2}\|_2
 \le\sqrt{\lambda_{\max}(G_0)\lambda_{\max}(G_0^{-1})}
                 \,\|A\|_{\mathrm F}
 \le\kappa_{G_0}\|A\|_{\mathrm F}.
\]
The first equality follows from the isometric map
$z\mapsto G_0^{1/2}z$. The first inequality follows by
submultiplication and the bound of the operator norm by the
sum of squared matrix entries. The second follows because
each positive eigenvalue is at most the sum of all eigenvalues.
For $\operatorname{diag}(G_0,G_0)$ the largest eigenvalue and
the largest inverse eigenvalue are unchanged, so the same
$\kappa_{G_0}$ applies to the full eight-dimensional operator.
The matrices $G_0$ and $G_0^{-1}$ remain those in RS24--25.

For $0<w\le g$, the absolute value of each entry of $A_w,B_w$
in MR14 is bounded by the corresponding entry below:
\begin{equation}
 A^\flat=
 \begin{pmatrix}
 0&\frac1{16g}&0&\frac3{16}\\
 \frac1{2g^2}&0&\frac1{16g}&0\\
 0&0&0&\frac1{16g}\\
 \frac1{4g^3}&0&0&0
 \end{pmatrix},\qquad
 B^\flat=
 \begin{pmatrix}
 \frac1{4g^3}&0&0&0\\
 0&\frac1{4g^3}&0&0\\
 \frac3{32g^4}&0&0&0\\
 0&\frac3{32g^4}&0&0
 \end{pmatrix}.
 \tag{MR20}
\end{equation}
To verify every bound, use $g+w\ge g$, $2g+w\le3g$ and
$3g+w\le4g$. The entries $\pm1/(16g)$ are unchanged,
the $(0,3)$ entry of $A_w$ is at most $3/16$, its $(1,0)$
entry is at most $1/(2g^2)$, and its $(3,0)$ entry is at most
$1/(4g^3)$. The first two nonzero entries of $B_w$ have
the same last bound, and the remaining two are at most
$3/(32g^4)$. This accounts for every nonzero entry.
Their squared sums are
\[
 \|A^\flat\|_{\mathrm F}^2
 =\frac9{256}+\frac3{256g^2}+\frac1{4g^4}+\frac1{16g^6},
 \qquad
 \|B^\flat\|_{\mathrm F}^2
 =\frac1{8g^6}+\frac9{512g^8}.
\]
Together they equal $B(g)$ in MR19. Since $|c|=1$, MR13
and the proved metric inequality yield the uniform estimate
\begin{equation}
 \left\|\Theta_w^{-1}-\frac{\mathscr L}{w}
       \right\|_{\operatorname{diag}(G_0,G_0)}
 \le K_{G_0},\qquad 0<w\le g.
 \tag{MR21}
\end{equation}
The phase $c$ remains the exact one in MR13; only its
modulus enters this inequality. MR6--7 give the identical
inequality after conjugation by the actual $Q_x$, in
$\mathbb H_U$, with
$\mathscr L_{\rm mark}=Q_x^{-1}\mathscr LQ_x$.

We also need the finite nonzero singular values of the
forward derivative multiplier, rather than only the six
values of the inverse pole. Subtracting the two full
matrices RS2 gives
\begin{equation}
 \frac{H_w-H_0}{w}=
 \begin{pmatrix}
 0&2(2g+w)&0&2(-g^2+gw+w^2)\\
 2&0&2(2g+w)&0\\
 0&-2&0&2(6g-w)\\
 0&0&-2&0
 \end{pmatrix}.
 \tag{MR22}
\end{equation}
This is an identity for every $w>0$, obtained by multiplying
the companion polynomial in RS2 before subtracting its value
at zero. On $0\le w\le g$,
\[
 -g^2\le -g^2+gw+w^2\le g^2.
\]
The lower inequality is $gw+w^2\ge0$; the upper is
$(g-w)(2g+w)\ge0$. Hence the absolute entries of MR22 are
bounded by
\[
 \begin{pmatrix}
 0&6g&0&2g^2\\2&0&6g&0\\0&2&0&12g\\0&0&2&0
 \end{pmatrix}.
\]
Its squared entry sum is exactly $4g^4+216g^2+12$.
Applying the metric estimate from MR19 proves
\begin{equation}
 \|H_w-H_0\|_{G_0}\le wD_{G_0},\qquad
 |\sigma_j^{G_0}(H_w)-\eta_j|\le wD_{G_0},
 \quad j=1,2,\quad 0<w\le g.
 \tag{MR23}
\end{equation}
The second inequality is the same singular-value variational
argument used in MR17; $\eta_1\ge\eta_2>0$ are the two
complete RS8 values in the same $G_0$.

Let $c_6>0$ be the smallest of the six explicitly specified
RS14 constants. Define the following positive threshold:
\begin{equation}
 w_*=\min\left\{
 g,\ \frac{\eta_2}{2D_{G_0}},\
 \frac{c_6}{2(K_{G_0}+1/\eta_2)}
 \right\},\qquad \epsilon_w=wD_{G_0}.
 \tag{MR24}
\end{equation}
This retains the original $w=\delta^2$ and $g=\gamma^2$;
every entry of the threshold is already a finite expression
in the original source data. For $0<w\le w_*$,
$\epsilon_w\le\eta_2/2$. By MR23 the first two singular
values of $H_w$ therefore lie in the positive intervals
$[\eta_j-\epsilon_w,\eta_j+\epsilon_w]$.
Their inverse values $1/(2\sigma_j^{G_0}(H_w))$ are both
at most $1/\eta_2$.

The upper-right block of $\Theta_w^{-1}$ is
$cH_w^{-1}/2$, while its lower-left block is
$cH_w^{-2}/2$. Because the two summands have the same
original Gram $G_0$, their singular spectra combine to
the singular spectrum of the whole inverse, as proved in
RS15. Thus these two bounded reciprocal values occur among
the eight actual singular values. On the other hand MR21
and the variational inequality give
\[
 \sigma_6(\Theta_w^{-1})
 \ge c_6/w-K_{G_0}
 \ge K_{G_0}+2/\eta_2>1/\eta_2.
\]
There are therefore exactly six larger values before those
two reciprocal values. This proves their indices, including
possible equality between $\eta_1$ and $\eta_2$:
\begin{equation}
 \begin{aligned}
 \sigma_7(\Theta_{{\rm mark},w}^{-1})
   &=\frac1{2\sigma_2^{G_0}(H_w)},&
 \left|\sigma_7-\frac1{2\eta_2}\right|
   &\le\frac{wD_{G_0}}{\eta_2^2},\\
 \sigma_8(\Theta_{{\rm mark},w}^{-1})
   &=\frac1{2\sigma_1^{G_0}(H_w)},&
 \left|\sigma_8-\frac1{2\eta_1}\right|
   &\le\frac{wD_{G_0}}{\eta_1^2},\\
 \left|\sigma_j(\Theta_{{\rm mark},w}^{-1})-\frac{c_j}{w}\right|
   &\le K_{G_0},&&1\le j\le6,\qquad 0<w\le w_*.
 \end{aligned}
 \tag{MR25}
\end{equation}
For the finite errors, if $|s-\eta|\le\epsilon_w$ then
\[
 \left|\frac1{2s}-\frac1{2\eta}\right|
 =\frac{|s-\eta|}{2\eta s}
 \le\frac{\epsilon_w}{2\eta(\eta-\epsilon_w)}
 \le\frac{\epsilon_w}{\eta^2}.
\]
The last step uses $\epsilon_w\le\eta_2/2\le\eta/2$.
MR7 transfers both the indexing and inequalities to the
full marked metric, retaining its $E,O$ cross terms.

These estimates also give uniformly positive finite-parameter
exterior bounds. For $0<w\le w_*$ put
\[
 P_r^\pm(w)=\prod_{j=1}^r(c_j\pm wK_{G_0}),
 \qquad 1\le r\le6.
\]
Every minus factor is positive: MR24 gives
$c_j/w-K_{G_0}\ge c_6/w-K_{G_0}>1/\eta_2>0$.
For $1\le r\le6$ and $s=1,2$ the complete bounds are
\begin{equation}
 \begin{aligned}
 \frac{P_r^-(w)}{w^r}
 &\le\|\wedge^r\Theta_{{\rm mark},w}^{-1}\|_{\mathbb H_U}
 \le\frac{P_r^+(w)}{w^r},\\
 \frac{P_6^-(w)}
 {w^6\,2^s\prod_{j=1}^s(\eta_{3-j}+\epsilon_w)}
 &\le\|\wedge^{6+s}\Theta_{{\rm mark},w}^{-1}\|_{\mathbb H_U}
 \le
 \frac{P_6^+(w)}
 {w^6\,2^s\prod_{j=1}^s(\eta_{3-j}-\epsilon_w)} .
 \end{aligned}
 \tag{MR26}
\end{equation}
Every denominator on both sides is positive by MR24.
The first line multiplies the first six inequalities
MR25. The second also uses the exact reciprocal intervals
$[1/(2(\eta_j+\epsilon_w)),1/(2(\eta_j-\epsilon_w))]$
proved above, in their established index order.
Rank zero still has norm one. At rank eight the exact
determinant in MR18 remains the sharper identity
$1/(2^{32}g^6w^6(g+w)^6)$. Thus MR25--26 control both
finite singular directions and all six divergent directions,
with the original source Gram evaluated through every
moment and with its complete marked return.

\section{Global signed duality and the complete ES family}


The proposed eight sections were supplied in the received
ramification continuation, Section 2 \cite{RamificationReceived}. Their
even basis and odd coefficient module have established human
provenance: specialize Melanie Matchett Wood's equations
\texttt{E:basis} and \texttt{E:Isbasis} to
$(f_0,f_1,f_2,f_3,f_4)=(A,1,B,C,D)$ \cite{WoodBinary}.
Wood credits the even basis to Birch--Merriman and its
multiplication to Nakagawa; her equation \texttt{E:mult}
is their explicit ring law. The form here is primitive
because $f_1=1$, and Wood's proposition \texttt{P:Iinvert}
identifies its odd coefficient module with $I_f^{-1}$.
We now prove directly that the proposed sections are a
global basis of the actual signed extension from FS2--10,
including its original sign convention, and calculate its
entire multiplication tensor and trace. No theorem about
a different choice of sign coordinate is substituted.

\section{Regular sections in both original root charts}
Let $\mathfrak B=\C[A,B,C,D]$ and retain
\begin{equation}
 \begin{gathered}
 h(r)=Ar^4+r^3+Br^2+Cr+D,\qquad
 k(v)=A+v+Bv^2+Cv^3+Dv^4,\\
 h(r)=0,\quad t^2=h'(r),\qquad
 k(v)=0,\quad\theta^2=-k'(v),\\
 v=1/r,\quad\theta=-t/r,\qquad
 e=(1,Ar,Ar^2+r,Ar^3+r^2+Br),\\
 o=t(A,Ar+1,Ar^2+r+B,Ar^3+r^2+Br+C).
 \end{gathered}
 \tag{GB1}
\end{equation}
The completed cover $X$ is obtained by these two charts
and exact inverses $r=1/v$, $t=-\theta/v$.
This is the finite cover constructed in FS2--10
\cite{Fable037}; in particular the reciprocal original factor
is exactly $t=1/a$. The gluing identity follows by
differentiating $k(v)=v^4h(1/v)$ at a root.

On the infinity chart the same sections are the polynomials
\begin{equation}
 \begin{gathered}
 e_0=1,\qquad e_1=-1-Bv-Cv^2-Dv^3,\qquad
 e_2=-B-Cv-Dv^2,\qquad e_3=-C-Dv,\\
 (o_0,o_1,o_2,o_3)=(-e_1,-e_2,-e_3,D)\theta.
 \end{gathered}
 \tag{GB2}
\end{equation}
For instance $A/v=-1-Bv-Cv^2-Dv^3$ follows by dividing
the root equation on the overlap; its right side is regular
on the whole infinity chart. The same relation gives
$(A+v)/v^2=-B-Cv-Dv^2$ and
$(A+v+Bv^2)/v^3=-C-Dv$.
Multiplying these expressions by $t=-\theta/v$ proves
the odd formulas, including the last value $D\theta$.
They agree on the overlap and so define all eight global
regular sections. The minus signs in GB2 are necessary
for the retained $\theta=-t/r$ convention; the incoming
continuation used the opposite infinity coordinate.

On the open $A\ne0$, coefficient columns in the basis
$(1,r,r^2,r^3,t,tr,tr^2,tr^3)$ are given by
\begin{equation}
 E=\begin{pmatrix}
 1&0&0&0\\0&A&1&B\\0&0&A&1\\0&0&0&A
 \end{pmatrix},\qquad
 Q=\begin{pmatrix}
 A&1&B&C\\0&A&1&B\\0&0&A&1\\0&0&0&A
 \end{pmatrix},\qquad
 M_{\mathrm{fin}}=\diag(E,Q),\quad
 \det M_{\mathrm{fin}}=A^7.
 \tag{GB3}
\end{equation}
The triangular diagonals prove $\det E=A^3$ and
$\det Q=A^4$, so they give a basis there. Complete
inverses are the following coefficient recurrences:
\[
 1=e_0,\quad r=e_1/A,\quad
 r^2=(e_2-r)/A,\quad r^3=(e_3-r^2-Br)/A,
\]
\[
 t=o_0/A,\quad tr=(o_1-t)/A,\quad
 tr^2=(o_2-tr-Bt)/A,\quad
 tr^3=(o_3-tr^2-Btr-Ct)/A.
\]
Every term on each right side was specified earlier in its
line. Thus these are finite explicit inverse columns, with
all their denominators admitted on the stated open.

\section{Every other chart and its complete inverse}
For $p\in\{0,1,-1,2\}$ work on $a_p=h(p)\ne0$, and set
\begin{equation}
 \begin{gathered}
 R=(p-r)^{-1},\quad T=t/(p-r),\qquad
 R=-v/(1-pv),\quad T=\theta/(1-pv),\\
 r=p-R^{-1},\quad t=T/R,\qquad
 v=R/(pR-1),\quad\theta=T/(1-pR),\\
 h_p(R)=R^4h(p-R^{-1})
        =a_pR^4+b_pR^3+c_pR^2+d_pR+A,\\
 a_p=Ap^4+p^3+Bp^2+Cp+D,\quad
 b_p=-4Ap^3-3p^2-2Bp-C,\\
 c_p=6Ap^2+3p+B,\qquad d_p=-4Ap-1,\qquad
 T^2=h_p'(R).
 \end{gathered}
 \tag{GB4}
\end{equation}
These are the original FS7 maps. At a root,
$h_p'(R)=R^2h'(p-R^{-1})$, proving the sign relation.
Every displayed inverse is used on its explicit nonzero
denominator; the inverse opens $R\ne0$ and $pR-1\ne0$
cover the new chart. On it the original finite cover has
the free monomial basis $(1,R,R^2,R^3,T,TR,TR^2,TR^3)$.
The derivative remains $h_p'$, including its leading
coefficient; it is not replaced by the derivative of a
polynomial divided by $a_p$.

Put
$k_1=Ap+d_p$,
$k_2=Ap^2+p-c_p-2pd_p$,
$k_3=Ap^3+p^2+Bp+b_p+pc_p+p^2d_p$.
The complete even coefficient matrix of the global sections
in this chart is
\begin{equation}
 E_p=\begin{pmatrix}
 1&k_1&k_2&k_3\\
 0&c_p&-b_p-2pc_p&a_p+pb_p+p^2c_p\\
 0&b_p&-a_p-2pb_p&pa_p+p^2b_p\\
 0&a_p&-2pa_p&p^2a_p
 \end{pmatrix}.
 \tag{GB5}
\end{equation}
The complete odd coefficient matrix, whose rows refer to
$(T,TR,TR^2,TR^3)$, is
\begin{equation}
 Q_p=\begin{pmatrix}
 -d_p&c_p+3pd_p&-b_p-2pc_p-3p^2d_p&
                        a_p+pb_p+p^2c_p+p^3d_p\\
 -c_p&b_p+3pc_p&-a_p-2pb_p-3p^2c_p&
                        pa_p+p^2b_p+p^3c_p\\
 -b_p&a_p+3pb_p&-2pa_p-3p^2b_p&
                        p^2a_p+p^3b_p\\
 -a_p&3pa_p&-3p^2a_p&p^3a_p
 \end{pmatrix}.
 \tag{GB6}
\end{equation}
These identities can be derived without a denominator $A$.
Substitute $r=p-1/R$, $t=T/R$ into GB1, and successively use
\[
 A/R=-a_pR^3-b_pR^2-c_pR-d_p
\]
and its products by $R^{-1}$ to remove the apparent negative
powers. The coefficients of $1/R$ that occur next are
integer multiples of $Ap,Ap^2,Ap^3$ and are eliminated by
the same equation. For example this gives
$e_1=k_1+c_pR+b_pR^2+a_pR^3$ and
$o_0/T=-d_p-c_pR-b_pR^2-a_pR^3$.
Repeating for the remaining columns gives precisely GB5--6.
All resulting entries are polynomials in the original
coefficients and $p$. Equality on $R\ne0$ proves their equality
on the whole completed chart: FS4--5 prove that $X$ is
integral and its finite-root open is dense, so multiplication
by this nonzero coordinate has no torsion. Equivalently the
identities follow by the explicit finite and infinity
substitutions in GB2 and GB4 on their covering opens.

Column additions followed by expansion give
\begin{equation}
 \det E_p=a_p^3,\qquad \det Q_p=a_p^4,\qquad
 M_p=\diag(E_p,Q_p),\qquad \det M_p=a_p^7.
 \tag{GB7}
\end{equation}
The following inverse expressions also verify these matrices
fully. Write $q_1=3Ap+1$, $q_2=3Ap^2+2p+B$ and
$q_3=Ap^3+p^2+Bp+C$. Then, as elements of the global
even basis,
\begin{equation}
 \begin{aligned}
 R&=(q_3e_0+p^2e_1+pe_2+e_3)/a_p,\\
 R^2&=-(q_2e_0+2pe_1+e_2+b_pR)/a_p,\\
 R^3&=(q_1e_0+e_1-c_pR-b_pR^2)/a_p .
 \end{aligned}
 \tag{GB8}
\end{equation}
The even inverse also has its first column $1=e_0$.
For the odd inverse every column is explicitly
\begin{equation}
 \begin{aligned}
 T&=(p^3o_0+p^2o_1+po_2+o_3)/a_p,\\
 TR&=-(3p^2o_0+2po_1+o_2+b_pT)/a_p,\\
 TR^2&=(3po_0+o_1-c_pT-b_pTR)/a_p,\\
 TR^3&=-(o_0+d_pT+c_pTR+b_pTR^2)/a_p .
 \end{aligned}
 \tag{GB9}
\end{equation}
Each is a four-step finite formula: substituting the
previously defined columns removes every occurrence of
$R,R^2,T,TR,TR^2$ from the coefficient expressions.
Multiplication by GB5--6 then gives respectively the
four monomial columns, proving both inverse products.
Thus these inverses use only the admitted denominator
$a_p$ and remain valid at $R=0$.

The five base opens really cover all coefficients, because
\[
 1=\tfrac12h(0)-\tfrac12h(1)-\tfrac16h(-1)
                         +\tfrac16h(2)-2A.
\]
The eight global sections are a basis on each of them by
GB3 and GB7. The map from the free module $\mathfrak B^8$
to $\pi_*\mathcal O_X$ that sends its columns to those
sections is therefore an isomorphism on this covering and
hence an isomorphism globally. Explicitly, both its kernel
and cokernel vanish after localization at each of the five
elements; the unit-ideal identity makes any element of either
module vanish, by taking a sufficiently large power of that
identity. Consequently
\begin{equation}
 \Gamma(X,\mathcal O_X)
 =\mathfrak B e_0\oplus\mathfrak B e_1\oplus
  \mathfrak B e_2\oplus\mathfrak B e_3
 \oplus\mathfrak B o_0\oplus\mathfrak B o_1\oplus
  \mathfrak B o_2\oplus\mathfrak B o_3
 \quad\text{is free of rank eight.}
 \tag{GB10}
\end{equation}
This proves the assertion at $A=0$ and at every collision
without dividing by the discriminant.

\section{The complete multiplication tensor}
For reproducibility, the coefficient rule underlying every
entry below is explicit. On $A\ne0$ let
\begin{equation}
 Z=\begin{pmatrix}
 0&0&0&-D/A\\1&0&0&-C/A\\0&1&0&-B/A\\0&0&1&-1/A
 \end{pmatrix},\qquad
 L=4AZ^3+3Z^2+2BZ+CI.
 \tag{GB11}
\end{equation}
Let $f_i(r)=e_i$ and $q_j(r)=o_j/t$ be precisely GB1.
In the global ordered eight-section basis, the eight
multiplication matrices are given by the full block formulas
\begin{equation}
 M_{e_i}=\diag(E^{-1}f_i(Z)E,Q^{-1}f_i(Z)Q),\qquad
 M_{o_j}=
 \begin{pmatrix}
 0&E^{-1}Lq_j(Z)Q\\Q^{-1}q_j(Z)E&0
 \end{pmatrix}.
 \tag{GB12}
\end{equation}
GB3 gives every inverse coefficient, so these are complete
finite expressions. The following tables evaluate them to
polynomial coefficients with no denominator. The column of
$M_a$ indexed by a basis vector $b$ is exactly the coefficient
column on the right of the displayed product $ab$.
Thus the tables specify all entries of all eight matrices;
$M_{e_0}=I_8$ supplies the remaining unit products.

To prove the entries, multiply the displayed polynomials in
GB1. Reduce $r^4$ by $Ar^4=-r^3-Br^2-Cr-D$ and use
$t^2=h'(r)$ for an odd--odd product. Applying GB3 produces
exactly GB12 and the tables. The exact checker performs every
such product and verifies all identities
$M_aM_b=\sum_k[ab]_kM_k$.
Because every coefficient in the evaluated tables is a
polynomial, the identities extend from $A\ne0$ to the full
free module GB10: that module has no $A$-torsion.
This proves the entire algebra law on the actual completion,
including all multiple-root fibres.

% BEGIN GENERATED GLOBAL TABLES
\subsection{Every multiplication coefficient}
The following table gives every product other than multiplication by $e_0=1$.
Commutativity gives the entries below the diagonal; each right-hand side is in the unchanged eight-section basis.
\paragraph{Even products}
\begin{align}
e_1 e_1 &=A e_{2} - e_{1}\tag{GB13a}\\
e_1 e_2 &=A e_{3} - B e_{1}\tag{GB13b}\\
e_1 e_3 &=- A D e_{0} - C e_{1}\tag{GB13c}\\
e_2 e_2 &=- A D e_{0} - B e_{2} - C e_{1} + e_{3}\tag{GB13d}\\
e_2 e_3 &=- C e_{2} - D e_{0} - D e_{1}\tag{GB13e}\\
e_3 e_3 &=- B D e_{0} - C e_{3} - D e_{2}\tag{GB13f}
\end{align}
\paragraph{Even--odd products}
\begin{align}
e_1 o_0 &=A o_{1} - o_{0}\tag{GB14a}\\
e_1 o_1 &=A o_{2} - B o_{0}\tag{GB14b}\\
e_1 o_2 &=A o_{3} - C o_{0}\tag{GB14c}\\
e_1 o_3 &=- D o_{0}\tag{GB14d}\\
e_2 o_0 &=A o_{2} - B o_{0}\tag{GB14e}\\
e_2 o_1 &=A o_{3} - B o_{1} - C o_{0} + o_{2}\tag{GB14f}\\
e_2 o_2 &=- C o_{1} - D o_{0} + o_{3}\tag{GB14g}\\
e_2 o_3 &=- D o_{1}\tag{GB14h}\\
e_3 o_0 &=A o_{3} - C o_{0}\tag{GB14i}\\
e_3 o_1 &=- C o_{1} - D o_{0} + o_{3}\tag{GB14j}\\
e_3 o_2 &=B o_{3} - C o_{2} - D o_{1}\tag{GB14k}\\
e_3 o_3 &=- D o_{2}\tag{GB14l}
\end{align}
\paragraph{Odd products}
\begin{align}
o_0 o_0 &=A^{2} C e_{0} + 4 A^{2} e_{3} - A e_{2} + e_{1} \left(- 2 A B + 1\right)\tag{GB15a}\\
o_0 o_1 &=- 2 A B e_{2} + 3 A e_{3} + e_{0} \left(- 4 A^{2} D + A C\right) + e_{1} \left(- 3 A C + B\right)\tag{GB15b}\\
o_0 o_2 &=2 A B e_{3} - 3 A C e_{2} + e_{0} \left(A B C - 3 A D\right) + e_{1} \left(- 4 A D + C\right)\tag{GB15c}\\
o_0 o_3 &=A C e_{3} - 4 A D e_{2} + D e_{1} + e_{0} \left(- 2 A B D + A C^{2}\right)\tag{GB15d}\\
o_1 o_1 &=e_{0} \left(- 7 A D + C\right) + e_{1} \left(- 4 A D + 2 B^{2} - 2 C\right) + e_{2} \left(- 3 A C - B\right) + e_{3} \left(- 2 A B + 3\right)\tag{GB15e}\\
o_1 o_2 &=e_{0} \left(- 2 A B D + B C - 3 D\right) + e_{1} \left(2 B C - 3 D\right) + e_{2} \left(- 4 A D - 2 C\right) + e_{3} \left(- 3 A C + 2 B\right)\tag{GB15f}\\
o_1 o_3 &=2 B D e_{1} - 3 D e_{2} + e_{0} \left(- A C D - 2 B D + C^{2}\right) + e_{3} \left(- 4 A D + C\right)\tag{GB15g}\\
o_2 o_2 &=e_{0} \left(3 A C D + B^{2} C - 5 B D\right) + e_{1} \left(- 2 B D + 3 C^{2}\right) + e_{2} \left(- B C - 3 D\right) + e_{3} \left(- 4 A D + 2 B^{2} - 2 C\right)\tag{GB15h}\\
o_2 o_3 &=- 2 B D e_{2} + 3 C D e_{1} + e_{0} \left(4 A D^{2} - 2 B^{2} D + B C^{2} - C D\right) + e_{3} \left(B C - 3 D\right)\tag{GB15i}\\
o_3 o_3 &=- C D e_{2} + 4 D^{2} e_{1} + e_{0} \left(- 3 B C D + C^{3} + 3 D^{2}\right) + e_{3} \left(- 2 B D + C^{2}\right)\tag{GB15j}
\end{align}
\subsection{The complete trace blocks}
\begin{equation}T_{ee}=\left[\begin{matrix}8 & -2 & - 4 B & - 6 C\\-2 & - 4 A B + 2 & - 6 A C + 2 B & - 8 A D + 2 C\\- 4 B & - 6 A C + 2 B & - 8 A D + 4 B^{2} - 4 C & 4 B C - 6 D\\- 6 C & - 8 A D + 2 C & 4 B C - 6 D & - 4 B D + 6 C^{2}\end{matrix}\right],\qquad T_{eo}=0.\tag{GB16}\end{equation}
The symmetric odd block has the following ten entries; $T_{ji}=T_{ij}$ supplies the other six.
\begin{align}
(T_{oo})_{00}&=- 16 A^{2} C + 8 A B - 2\tag{GB17a}\\
(T_{oo})_{01}&=- 32 A^{2} D + 8 A B^{2} - 4 A C - 2 B\tag{GB17b}\\
(T_{oo})_{02}&=8 A B C - 16 A D - 2 C\tag{GB17c}\\
(T_{oo})_{03}&=2 A C^{2} - 2 D\tag{GB17d}\\
(T_{oo})_{11}&=24 A B C - 48 A D - 6 C\tag{GB17e}\\
(T_{oo})_{12}&=18 A C^{2} - 18 D\tag{GB17f}\\
(T_{oo})_{13}&=16 A C D - 8 B D + 2 C^{2}\tag{GB17g}\\
(T_{oo})_{22}&=48 A C D - 24 B D + 6 C^{2}\tag{GB17h}\\
(T_{oo})_{23}&=32 A D^{2} - 8 B^{2} D + 2 B C^{2} + 4 C D\tag{GB17i}\\
(T_{oo})_{33}&=- 8 B C D + 2 C^{3} + 16 D^{2}\tag{GB17j}
\end{align}

% END GENERATED GLOBAL TABLES

\section{Trace determinants without a missing leading coefficient}
Let $\Delta=\Disc(h)$ mean the binary quartic discriminant:
\begin{equation}
 \begin{aligned}
 \Delta={}&256A^3D^3-192A^2CD^2-128A^2B^2D^2
 +144A^2BC^2D-27A^2C^4\\
 &+144ABD^2-6AC^2D-80AB^2CD+18ABC^3
 +16AB^4D-4AB^3C^2\\
 &-27D^2+18BCD-4C^3-4B^3D+B^2C^2 .
 \end{aligned}
 \tag{GB18}
\end{equation}
The trace blocks above are computed by the ordinary
multiplication trace in this rank-eight algebra. The single
element trace vector is
$(8,-2,-4B,-6C,0,0,0,0)$, obtained by taking the diagonals
of GB12. In particular every displayed block entry can also
be obtained by applying this vector to its product table.
Parity makes each mixed multiplication matrix off-diagonal,
so its trace is zero.

For $A\ne0$ and distinct roots let $V_r=(r_j^i)_{j,i=0}^3$.
The root-algebra trace Gram in monomials is $V_r^{\mathsf T}V_r$,
whose determinant is $\Delta/A^6$. Its signed extension
doubles that trace on the even summand. On the odd summand
it gives twice the root trace after multiplication by $h'$.
The determinant of multiplication by $h'$ is
$\prod_jh'(r_j)=\Delta/A^2$; the six unordered root-pair
signs have product $(-1)^6=1$. Using GB3 therefore gives
\begin{equation}
 \det T_{ee}=2^4(A^3)^2\frac{\Delta}{A^6}=16\Delta,\qquad
 \det T_{oo}=2^4(A^4)^2
           \frac{\Delta}{A^6}\frac{\Delta}{A^2}=16\Delta^2,
 \qquad \det T=256\Delta^3.
 \tag{GB19}
\end{equation}
Both sides are polynomial coefficient identities by the
full multiplication calculation, so they hold at $A=0$
and on $\Delta=0$ as well. Degeneration of this trace
form does not remove any vector from the free basis GB10.

\section{The original seven-point target with its eighth completed state}
At $(A,B,C,D)=(0,0,-1,0)$ choose $p=2$.
Then $(a_p,b_p,c_p,d_p)=(6,-11,6,-1)$ and
$h_p(R)=R(R-1)(2R-1)(3R-1)$. The transition matrices are
\begin{equation}
 E_2=\begin{pmatrix}
 1&-1&0&1\\0&6&-13&8\\0&-11&38&-32\\0&6&-24&24
 \end{pmatrix},\qquad
 Q_2=\begin{pmatrix}
 1&0&-1&0\\-6&25&-34&16\\11&-60&108&-64\\-6&36&-72&48
 \end{pmatrix},
 \qquad \det E_2=6^3,\quad\det Q_2=6^4.
 \tag{GB20}
\end{equation}
Their complete inverses are GB8--9 with these exact coefficients.
All eight vectors remain independent. The discriminant is
$\Delta=4$, so the trace determinants are $64$, $256$ and
$16384$ for the even, odd and full blocks respectively.
Every entry at this target is
\begin{equation}
 T_{ee}=\begin{pmatrix}
 8&-2&0&6\\-2&2&0&-2\\0&0&4&0\\6&-2&0&6
 \end{pmatrix},\qquad
 T_{oo}=\begin{pmatrix}
 -2&0&2&0\\0&6&0&2\\2&0&6&0\\0&2&0&-2
 \end{pmatrix},\qquad T_{eo}=0.
 \tag{GB21}
\end{equation}
The six finite states and the two infinity states have the
following complete evaluation columns:
\begin{equation}
 \begin{array}{c|c|c|c}
 \text{selected root}&\text{sign coordinate}&(e_0,e_1,e_2,e_3)&
                    (o_0,o_1,o_2,o_3)\\\hline
 r=-1&t=\pm\sqrt2&(1,0,-1,1)&(0,t,-t,0)\\
 r=0&t=\pm i&(1,0,0,0)&(0,t,0,-t)\\
 r=1&t=\pm\sqrt2&(1,0,1,1)&(0,t,t,0)\\
 v=0&\theta=\pm i&(1,-1,0,1)&(\theta,0,-\theta,0)
 \end{array}
 \tag{GB22}
\end{equation}
Each row is obtained from GB1 or GB2; the two values of its
sign coordinate are distinct states. At infinity $T=\theta$
by GB4. The original affine source has $b=i=1/\theta$, hence
$\theta=-i$ is its included state and $\theta=i$ is the
omitted state restored by the completion. The original seven
affine states and all eight completed states are therefore
represented with their unchanged sign orientation.


\section{The global residue section and its two infinity zeros}
Let $\mathfrak B=\C[A,B,C,D]$, let $X$ be the actual smooth
finite completion FS2--15, and put $\mathscr S=\Gamma(X,\mathcal O_X)$.
The global ordered basis $b=(e_0,e_1,e_2,e_3,o_0,o_1,o_2,o_3)$
is GB1. Its even part is the classical binary-form basis, and
its odd coefficient module is the inverse ideal in Wood's
equations \texttt{E:basis}, \texttt{E:Isbasis} and proposition
\texttt{P:Iinvert} \cite{WoodBinary}. Wood's dual-module theorem
and equation \texttt{E:dualmap} identify the classical source of
this residue construction. We prove the signed global formulas
and their original-conductor return below with all coefficients
retained. No claim that the classical basis is new is made.

To avoid confusing the odd coefficient matrix with the original
marked odd frame $O$ of SF2, write the GB3 matrices as
\begin{equation}
 E_g=\begin{pmatrix}1&0&0&0\\0&A&1&B\\0&0&A&1\\0&0&0&A\end{pmatrix},\quad
 O_g=\begin{pmatrix}A&1&B&C\\0&A&1&B\\0&0&A&1\\0&0&0&A\end{pmatrix},\quad
 T_g=\operatorname{diag}(E_g,O_g),\qquad\det T_g=A^7.
 \tag{GD1}
\end{equation}
Thus the finite $r,t$ coefficient column of an element with
global column $z$ is $T_gz$ on $A\ne0$. The full inverses are
\begin{equation}
\begin{aligned}
 E_g^{-1}&=\begin{pmatrix}1&0&0&0\\
 0&A^{-1}&-A^{-2}&A^{-3}-BA^{-2}\\
 0&0&A^{-1}&-A^{-2}\\0&0&0&A^{-1}\end{pmatrix},\\
 O_g^{-1}&=\begin{pmatrix}
 A^{-1}&-A^{-2}&A^{-3}-BA^{-2}&-A^{-4}+2BA^{-3}-CA^{-2}\\
 0&A^{-1}&-A^{-2}&A^{-3}-BA^{-2}\\
 0&0&A^{-1}&-A^{-2}\\0&0&0&A^{-1}
 \end{pmatrix}.
\end{aligned}\tag{GD2}
\end{equation}
Multiplication against GD1 cancels every off-diagonal entry
and gives the identity on both sides. GB5--10 provide the
complete regular replacement on every open $h(p)\ne0$;
the global basis itself does not cease to exist at $A=0$.

Define the $\mathfrak B$-linear functional and its associative
bilinear pairing by
\begin{equation}
 \Lambda_0\left(\sum_{j=0}^3z_je_j+\sum_{j=0}^3z_{4+j}o_j\right)=z_7,
 \qquad\mathcal B_0(f,g)=\Lambda_0(fg).
 \tag{GD3}
\end{equation}
On $A\ne0$ this equals RD2's functional
$[r^3]S/A$ on $R(r)+tS(r)$: only the column $o_3$ of
$O_g$ has a cubic coefficient, and that coefficient is $A$.
Multiplying the complete RD3 matrix on the left by $E_g^{\mathsf T}$
and on the right by $O_g$ gives
\begin{equation}
 B_0=\begin{pmatrix}0&C_0\\C_0^{\mathsf T}&0\end{pmatrix},\qquad
 C_0=\begin{pmatrix}0&0&0&1\\0&0&A&0\\0&A&1&0\\A&1&B&0\end{pmatrix},
 \qquad\det C_0=A^3,\quad\det B_0=A^6.
 \tag{GD4}
\end{equation}
All entries are polynomial. Since GB proves that $\mathscr S$
is free over the integral ring $\mathfrak B$, the equality
proved after inverting $A$ is the equality over the whole base.
Its full inverse on that open has off-diagonal blocks
$C_0^{-\mathsf T},C_0^{-1}$, where
\begin{equation}
 C_0^{-1}=\begin{pmatrix}
 0&A^{-3}-BA^{-2}&-A^{-2}&A^{-1}\\
 0&-A^{-2}&A^{-1}&0\\0&A^{-1}&0&0\\1&0&0&0
 \end{pmatrix}.
 \tag{GD5}
\end{equation}
Direct matrix multiplication proves both inverses; no root
separation or nonzero discriminant is needed.

At $A=0$, $C_0$ has rank three, so $B_0$ has rank six and
its radical is precisely $\langle e_1,o_0\rangle$. The exact
two infinity projectors are
\begin{equation}
 e_\infty=-e_1,\qquad
 \epsilon_+=\frac{-e_1+i o_0}{2},\quad
 \epsilon_- =\frac{-e_1-i o_0}{2},\qquad
 \operatorname{rad}\mathcal B_0=\C\epsilon_+\oplus\C\epsilon_-.
 \tag{GD6}
\end{equation}
Here all products are in the specialized global algebra.
For every finite root $e_1=Ar=0$ and $o_0=At=0$.
At infinity GB2 gives $e_1=-1$, $o_0=\theta$ and
$\theta^2=-1$. These identities also hold for the complete
possibly nonreduced fibre: the simple infinity root is
disjoint from the finite cubic, hence is an open and closed
factor; the finite cubic factor retains all its nilpotents.
Equivalently the full GB multiplication table gives
$e_\infty^2=e_\infty$, $o_0^2=-e_\infty$ and
$e_\infty o_0=o_0$. These equations prove that the two
elements GD6 are orthogonal idempotents with sum $e_\infty$.
Their $o_0$ eigenvalues are respectively $-i,+i$.
Thus $\epsilon_+$ is the included original infinity state,
and $\epsilon_-$ is the omitted state. The loss of the
finite-coordinate residue includes both of them.

This is a different exact divisor from trace ramification.
At the original target $u_0=(0,0,-1,0)$ the trace Gram in
this very global basis is
\begin{equation}
 \operatorname{Tr}_{u_0}=\operatorname{diag}\left(
 \begin{pmatrix}8&-2&0&6\\-2&2&0&-2\\0&0&4&0\\6&-2&0&6\end{pmatrix},
 \begin{pmatrix}-2&0&2&0\\0&6&0&2\\2&0&6&0\\0&2&0&-2\end{pmatrix}\right),
 \quad\det\operatorname{Tr}_{u_0}=2^{14}.
 \tag{GD7}
\end{equation}
Each entry is the trace of multiplication by the corresponding
product in GB's complete table. Its determinant is also
$256\mathfrak D(u_0)^3=256\cdot4^3$ by GB's trace formula.
There are the six finite states $r=0,1,-1$ with
$t^2=3r^2-1$, the included original point
$(a,y,z,w)=(0,0,i/2,0)$, and the omitted infinity sign.
The $p=2$ chart of RD29 has perfect residue determinant
$6^{-8}$ on all eight states. GD6--7 identify exactly why
using only $\Lambda_0$ would falsely lose an original state
even though this completed fibre is etale.

\section{The precise global dualizing obstruction}
Define
\begin{equation}
 \omega=\operatorname{Hom}_{\mathfrak B}(\mathscr S,\mathfrak B),
 \qquad (f\ell)(g)=\ell(fg).
 \tag{GD8}
\end{equation}
This specifies the $\mathscr S$-module structure, including
its action on every functional. On $h(p)\ne0$ the RD25
functional $\Lambda_p$ generates this module: its Gram is
invertible with determinant $h(p)^{-8}$, so
$f\mapsto f\Lambda_p$ is an isomorphism onto $\omega$.
The same statement holds on $A\ne0$ for $\Lambda_0$.
FS6 proves that these opens cover the base. Therefore
$\omega$ is an invertible $\mathscr S$-module, with these
fully specified local generators and inverse Gram matrices.
RD26 and the complete trace calculation RD6 give
\begin{equation}
 \Lambda_0=R^3\Lambda_p,\qquad
 \operatorname{Tr}=2T^3\Lambda_p,\qquad
 t=T/R,\qquad
 \Lambda_0=\frac{1}{2t^3}\operatorname{Tr}.
 \tag{GD9}
\end{equation}
The first two equations are identities of regular sections
on the whole $h(p)\ne0$ chart. For the first, equality
holds on the dense $A\ne0$ open by RD26; the difference
is a section of an invertible module over the integral
scheme $X$ and is zero there, hence is zero everywhere.
The trace equation follows directly from the local
monomial calculation and holds even at $T=0$.
The final equation is an identity of rational sections;
it does not assert that $1/t$ is regular at ramification.

Let $R_0$ and $D_\pm$ have their exact FS5,FS15 meanings.
On $D_\pm$ the parameter $v$ is a simple transverse
coordinate, since $2\theta\ne0$ in FS5. In a $p$ chart
$R=-v/(1-pv)$ has a simple zero on each of them and no
other zero. At $R=0$, $T=\theta=\mp i$ is a unit.
On the finite chart $t$ is one of the polynomial coordinates
FS4, so it has a simple zero exactly on $R_0$.
The equation $t=-\theta/v$ shows a simple pole on each
infinity divisor and no further zero or pole. It follows,
with multiplicities, that
\begin{equation}
\begin{aligned}
 \operatorname{div}(\Lambda_0)&=3D_++3D_-,&
 \operatorname{div}(\operatorname{Tr})&=3R_0,\\
 \operatorname{div}(t)&=R_0-D_+-D_-,&
 \operatorname{div}(a)&=D_++D_--R_0,\quad a=1/t.
\end{aligned}\tag{GD10}
\end{equation}
The first two are divisors of sections of $\omega$, not
divisors of scalar functions. The local formulas GD9
prove them because every $\Lambda_p$ is a generator.

There is no global $\ell\in\omega$ whose associative
pairing $(f,g)\mapsto\ell(fg)$ is perfect over every
base point. Here is a complete obstruction proof. Such
a pairing would make the $\mathscr S$-linear map
$f\mapsto f\ell$ an isomorphism: in free base-module
bases its determinant is nonzero at every complex base
point, hence is a unit polynomial, so its matrix has a
global inverse. Consequently $\Lambda_0=f\ell$ for a
global regular $f\in\mathscr S$, and GD10 would imply
\begin{equation}
 \operatorname{div}(f)=3D_++3D_-.
 \tag{GD11}
\end{equation}
Restrict to the exact original open
$U=X\setminus(R_0\cup D_-)=\mathbb A^4_{a,y,z,w}$
of FS15. On it $f/a^3$ has neither zeros nor poles
along any prime divisor. Write this rational function
as a quotient of coprime polynomials in $a,y,z,w$.
Any nonconstant irreducible factor of the numerator
would give a zero divisor; any such denominator factor
would give a pole divisor. Both are excluded. Thus the
quotient is a nonzero constant $c$, and equality on this
dense open gives $f=c a^3$ in the function field of $X$.
But GD10 gives this function divisor
$3D_++3D_--3R_0$, contradicting GD11 because $R_0$
is a nonempty divisor. This proves the assertion.
It neither denies local perfect residue pairings nor
denies other nonassociative base-module pairings.

The exact global replacement is the evaluation pairing
\begin{equation}
 \mathscr S\times\omega\longrightarrow\mathfrak B,
 \quad(f,\ell)\longmapsto\ell(f),\qquad
 \langle af,\ell\rangle=\langle f,a\ell\rangle.
 \tag{GD12}
\end{equation}
In $b$ and its complete base-dual basis $b^\vee$ its
matrix is $I_8$, including every boundary fibre.
If multiplication by $a\in\mathscr S$ has matrix $M_a$
on $b$ columns, its action on dual columns is exactly
$M_a^{\mathsf T}$: evaluate both sides on an arbitrary
column. This proves the stated balancing identity and
all eight dual coordinate maps. It does not replace
$\omega$ by $\mathscr S$ as an algebra module.

\section{The global basis in the original marked metric}
Fix the original period, its distinct labelled reference
quartet, all its moments and Gamma masses as in MR1--7.
Write $u_\circ$ for its quartic coefficients. In the
RS1 convention with $g_\circ=\gamma_\circ^2$ and
$w_\circ=\delta_\circ^2$ these coefficients are
\begin{equation}
 \begin{aligned}
 A_\circ&=-\tfrac12,&B_\circ&=-\tfrac34-g_\circ+w_\circ,\\
 C_\circ&=\tfrac14+g_\circ-w_\circ,&
 D_\circ&=-\tfrac1{32}-\tfrac14(g_\circ-w_\circ)
                         -\tfrac12(g_\circ+w_\circ)^2.
 \end{aligned}\tag{GD13}
\end{equation}
They are obtained by expanding the unchanged polynomial
$-\tfrac12[(r-\tfrac12)^4+2(g_\circ-w_\circ)(r-\tfrac12)^2
 +(g_\circ+w_\circ)^2]$. Set $T_\circ=T_g(u_\circ)$;
it is invertible by GD1. Its complete inverse is GD2
evaluated at precisely these coefficients.

For coefficient columns of $R(r)+tS(r)$, define the
original maps, with $\kappa=a_{\rm amp}/2-4-1/2$ unchanged,
\begin{equation}
\begin{aligned}
 F_U([R],[S])&=(T_A^{-1}R(S'-\kappa),T_A^{-1}S(S'-\kappa)),\\
 F_W([R],[S])&=(R(S'-\kappa),S(S'-\kappa)),\\
 J_{\circ,U}&=F_UT_\circ,&J_{\circ,W}&=F_WT_\circ,
 \qquad J_{\circ,W}=(T_A\oplus T_A)J_{\circ,U}.
\end{aligned}\tag{GD14}
\end{equation}
Every coefficient of the original $T_A^{-1}$ is FC24.
The inverse of $F_W$ substitutes $S'=r+\kappa$;
the inverse of $F_U$ first applies $T_A$ in each
component and then the same substitution. Composing
these inverses with GD2 at $u_\circ$ proves both exact
inverses of $J_{\circ,U},J_{\circ,W}$ on every fibre.
These are fixed base-module receivers, not varying
algebra identifications with the reference fibre.

The earlier marked column $z=(\alpha,\beta)$ corresponds
at the reference fibre to the global column
\begin{equation}
 z_g=T_\circ^{-1}Q_rz,\qquad
 J_{\circ,U}z_g=\mathcal J_Uz,\quad
 J_{\circ,W}z_g=\mathcal J_Wz.
 \tag{GD15}
\end{equation}
These equalities follow from $F_UQ_r=\mathcal J_U$
and $F_WQ_r=\mathcal J_W$, proved in RD18--20 and
MR1--3. Thus all four labels $(\mathrm M,\mathrm N,
\mathrm T,\mathrm E)$ and both signs are retained.
For every global multiplication matrix $M_f(u)$ the
complete original action is
\begin{equation}
 M_{f,U}(u)=J_{\circ,U}M_f(u)J_{\circ,U}^{-1},\qquad
 M_{f,W}(u)=J_{\circ,W}M_f(u)J_{\circ,W}^{-1},\qquad
 (T_A\oplus T_A)M_{f,U}=M_{f,W}(T_A\oplus T_A).
 \tag{GD16}
\end{equation}
Insert GD14 in the two sides to prove the last identity.
The entire multiplication table is polynomial in $u$,
so this action remains defined at $A=0$. A dual column
$\eta$ returns to the functional
$\eta^{\mathsf T}J_{\circ,U}^{-1}$ on the actual source,
or $\eta^{\mathsf T}J_{\circ,W}^{-1}$ on the actual
target. Substitution proves GD12 after both returns.

Let $D_{1/2}=\operatorname{diag}(C_{1/2},C_{1/2})$ with
the full binomial matrix MR1. The complete original
source Gram and the Gram of its complex-linear dual are
\begin{equation}
\begin{aligned}
 \Gamma_r&=D_{1/2}^*\operatorname{diag}(G_{U,x},G_{U,x})D_{1/2},\\
 \Gamma_b&=T_\circ^*\Gamma_rT_\circ
                 =\operatorname{diag}(G_e,G_o),\qquad
 \Gamma_{b^\vee}=\Gamma_b^{-\mathsf T},\\
 (T_\circ^{-1}Q_r)^*\Gamma_b(T_\circ^{-1}Q_r)&=\mathbb H_U.
\end{aligned}\tag{GD17}
\end{equation}
Here $G_{U,x}$ is the complete original RS24 matrix,
with its original mass, every inverse moment, phase
and center. Block diagonality follows from the displayed
parity blocks in $D_{1/2}$ and $T_\circ$, not from a
claim that the marked mixed Gram is diagonal. The last
identity is MR4 with $Q_x=D_{1/2}Q_r$. For a dual column
$\eta$, its squared norm is
$\eta^{\mathsf T}\Gamma_b^{-1}\bar\eta
 =\eta^*\Gamma_b^{-\mathsf T}\eta$, which proves
the dual Gram. The target formulas follow by replacing
$G_{U,x}$ with the actual $G_{W,x}$ throughout.
The residue and trace in physical coordinates are
$J_{\circ,U}^{-\mathsf T}B_0(u)J_{\circ,U}^{-1}$ and
$J_{\circ,U}^{-\mathsf T}\operatorname{Tr}_uJ_{\circ,U}^{-1}$;
they are bilinear forms, not the positive Gram GD17.

The relation to the earlier varying finite-coordinate
receiver is fully specified. On $A\ne0$ that receiver
is $F_UT_g(u)$, so the map from its physical state to
GD14's global fixed receiver is
\begin{equation}
 K_U(u)=F_UT_\circ T_g(u)^{-1}F_U^{-1},\qquad
 K_U(u)^{-1}=F_UT_g(u)T_\circ^{-1}F_U^{-1},\qquad
 \det K_U(u)=(A_\circ/A)^7.
 \tag{GD18}
\end{equation}
All factors and inverses have been displayed. Their
products are the identity, and applying the first to
$F_UT_g(u)z_g$ gives $J_{\circ,U}z_g$ exactly.
Thus the change of receiver is not silently omitted.

\section{The full inverse cost of the global scalar residue}
Regard $B_0(u)$ as a map from the global states with
Gram $\Gamma_b$ to their actual dual with Gram
$\Gamma_b^{-\mathsf T}$. Hold $B,C,D$ fixed and let
the complex parameter $A$ tend to zero. Write $C_{00}=C_0|_{A=0}$.
The four singular values of the cross map are those of
$Z_A=G_e^{-\mathsf T/2}C_0G_o^{-1/2}$. Its three
positive boundary singular values are $q_1\ge q_2\ge q_3>0$.
They are determined exactly, with all original metric
entries retained, by
\begin{equation}
\begin{aligned}
 L&=G_o^{-1}C_{00}^*G_e^{-\mathsf T}C_{00},\qquad
 a_1=\operatorname{tr}L,\quad
 a_2=\tfrac12[(\operatorname{tr}L)^2-\operatorname{tr}(L^2)],\\
 k^2&=(G_o)_{00}(G_e)_{11},\qquad
 a_3=\frac{k^2}{\det G_e\det G_o},\qquad
 \{q_1^2,q_2^2,q_3^2\}
   =\{\hbox{roots of }z^3-a_1z^2+a_2z-a_3\}.
\end{aligned}\tag{GD19}
\end{equation}
All three roots are positive because $L$ is similar to
$Z_0^*Z_0$, which is positive semidefinite of rank three.
To verify the stated constant, GD5 gives
$\operatorname{adj}(C_{00})=e_0e_1^{\mathsf T}$,
where here $e_j$ denotes the numerical coordinate unit
column of length four. For a rank-three four-by-four
matrix, the squared Frobenius norm of its adjugate is
the product of its three nonzero singular squares;
this follows by taking its singular value decomposition
and the induced third exterior map. Applying
$\operatorname{adj}(UXV)=\det(UV)V^{-1}\operatorname{adj}(X)U^{-1}$
to $Z_0$ gives that norm as precisely
$(G_o)_{00}(G_e)_{11}/(\det G_e\det G_o)$.
Trace and trace-square give the other two elementary
symmetric functions, proving GD19 completely.

GD5 gives the complete inverse pole
\begin{equation}
 C_0^{-1}=A^{-3}e_0e_1^{\mathsf T}
 +A^{-2}(-B e_0e_1^{\mathsf T}-e_0e_2^{\mathsf T}-e_1e_1^{\mathsf T})
 +A^{-1}(e_0e_3^{\mathsf T}+e_1e_2^{\mathsf T}+e_2e_1^{\mathsf T})
 +e_3e_0^{\mathsf T}.
 \tag{GD20}
\end{equation}
The inverse of $B_0$ after passage to Euclidean coordinates
has the off-diagonal blocks
$G_e^{1/2}C_0^{-\mathsf T}G_o^{\mathsf T/2}$ and
$G_o^{1/2}C_0^{-1}G_e^{\mathsf T/2}$. They are mutual
transposes and have identical singular values. The
leading rank-one term of the second has norm $k$.
Hence exactly two of the eight inverse singular values
diverge, each asymptotic to $k|A|^{-3}$; the other six
converge, in descending order, to
$q_3^{-1},q_3^{-1},q_2^{-1},q_2^{-1},q_1^{-1},q_1^{-1}$.
Indeed the three positive singular values of $Z_A$
converge to those of $Z_0$; their inverse limits follow
by continuity, while the remaining inverse singular
value has the rank-one limit just computed.

Define the ordered list
$d=(k,k,q_3^{-1},q_3^{-1},q_2^{-1},q_2^{-1},q_1^{-1},q_1^{-1})$.
The first two entries represent the divergent pair and
need not exceed the finite constants before multiplying
by $|A|^{-3}$. Every inverse exterior norm and the
full determinant are now exactly controlled by
\begin{equation}
 \lim_{A\to0}|A|^{3\min(j,2)}\|\wedge^j B_0(A)^{-1}\|
       =\prod_{\nu=1}^j d_\nu\quad(1\le j\le8),\qquad
 \|\wedge^8 B_0(A)^{-1}\|=\frac{\det\Gamma_b}{|A|^6}.
 \tag{GD21}
\end{equation}
The norm of an exterior map is the product of its
largest singular values, proved by applying the singular
value decomposition to the induced exterior bases.
This proves every limit. For the last equality, the
Euclidean inverse determinant is
$\det(\Gamma_b^{1/2}B_0^{-1}\Gamma_b^{\mathsf T/2})
 =\det\Gamma_b/\det B_0$; its absolute value is the
stated formula. GD19 also gives agreement of the
full-rank limit constant. These are costs of the
scalar residue map into its dual; GD12's evaluation
pairing and the fixed original $T_A$ remain invertible.

\section{The same cubic factor in the original escaping sign state}
The original sign involution has already been derived
from the polynomial FS1 and its factor signs. Retaining
all its original source coordinates, it is
\begin{equation}
 \mathcal D(a,y,z,w)=\left(-a,\ y+\frac{2i}{a},\
 -z+\frac{6i}{a^2},\
 w-\frac{14iy^2}{a}+\frac{28y}{a^2}+\frac{40i}{a^3}\right)
 \quad(a\ne0).
 \tag{GD22}
\end{equation}
This is the exact rational sign change from FS31 in
the original polynomial source, not a newly introduced
linear map. On $U$ the trace section is a local generator:
FS16 proves that the original polynomial has determinant
$-2$ everywhere. GD9 therefore gives the exact regular
ratio $\Lambda_0=(a^3/2)\operatorname{Tr}$ on that open.
Multiplying every coordinate of GD22 by this very ratio
gives the polynomial map
\begin{equation}
 W(a,y,z,w)=\frac{a^3}{2}\mathcal D(a,y,z,w)
  =\left(-\frac{a^4}{2},\ \frac{a^3y}{2}+ia^2,\
  -\frac{a^3z}{2}+3ia,\
  \frac{a^3w}{2}-7ia^2y^2+14ay+20i\right),\qquad
 W(0,y,z,w)=(0,0,0,20i).
 \tag{GD23}
\end{equation}
Each equality follows by multiplication with no discarded
lower term. Since the right side is polynomial, it
extends the weighted sign map to the whole original
source and proves its stated boundary value.

Let $\mathbf e_4=(0,0,0,1)^{\mathsf T}$ in the original
ordered four-dimensional source coordinates of $P$.
Use the unchanged original maps $\Phi=XK^{-1}$ and
$\Psi=YK^{-1}=T_A\Phi$. Their full weighted images obey
\begin{equation}
 W_U(q)=\Phi W(q),\qquad W_W(q)=\Psi W(q),\qquad
 T_AW_U(q)=W_W(q),\qquad
 W_U(0,y,z,w)=20i\Phi\mathbf e_4,\quad
 W_W(0,y,z,w)=20i\Psi\mathbf e_4.
 \tag{GD24}
\end{equation}
These are polynomial identities and include every
original moment, phase and Gamma mass through FC23--32.
They provide a global extension of the weighted map in
both original receiver spaces, with the exact conductor
between them. For $a\to0$ and bounded $y,z,w$, GD23
gives, in the retained positive norms,
\begin{equation}
 \lim |a|^3\|\Phi\mathcal D(q)\|_U=40\|\Phi\mathbf e_4\|_U,
 \qquad
 \lim |a|^3\|\Psi\mathcal D(q)\|_W=40\|\Psi\mathbf e_4\|_W.
 \tag{GD25}
\end{equation}
To prove the limits, all nonconstant terms in GD23 tend
to zero under the stated boundedness assumption.
Continuity of each fixed linear map and norm gives
$\|(a^3/2)\Phi\mathcal D(q)\|_U\to20\|\Phi\mathbf e_4\|_U$
and the analogous target statement. Absolute homogeneity
of the norm yields GD25, even for complex approaches.
The constants are positive since $\Phi,\Psi$ are
isomorphisms. Thus the unweighted escape persists
with its exact cubic rate. Its residue-weighted
extension is regular because the residue section
vanishes to exactly the matching order on $D_+$.
The perfect global correspondence is GD12 and its
original return GD14--17; no singularity of the fixed
original conductor or arithmetic RH conclusion is
deduced by equating these different maps.

\section{Explicit finite-parameter bounds for every global inverse exterior power}
The preceding asymptotes have the following complete
finite-parameter version in exactly the same original metric.
Retain fixed $B,C,D,G_e,G_o$ and write $s=|A|>0$.
Let $K_3=e_0e_1^{\mathsf T}$ and let $K_2,K_1,K_0$
be respectively the coefficients of $A^{-2},A^{-1},1$
in GD20, without altering them. Put
\begin{equation}
\begin{aligned}
 C_1&=\begin{pmatrix}0&0&0&0\\0&0&1&0\\0&1&0&0\\1&0&0&0\end{pmatrix},
 &C_0(A)&=C_{00}+AC_1,\\
 d^2&=\operatorname{tr}(G_o^{-1}C_1^*G_e^{-\mathsf T}C_1),&
 b_j^2&=\operatorname{tr}(G_e^{\mathsf T}K_j^*G_oK_j)\quad(j=0,1,2),\\
 E(s)&=b_2s+b_1s^2+b_0s^3,&
 s_*&=\min\left\{1,\frac{q_3}{2d},
 \frac{k}{2(b_2+b_1+b_0)},\left(\frac{kq_3}{8}\right)^{1/3}\right\}>0.
\end{aligned}\tag{GD26}
\end{equation}
All square roots mean the nonnegative real roots of the
displayed exact positive quantities. In particular $d>0$
because $C_1\ne0$ and the two metrics are positive;
$b_0>0$ because $K_0=e_3e_0^{\mathsf T}\ne0$.
Thus no denominator in $s_*$ vanishes.

These trace expressions are the complete squared
Hilbert--Schmidt norms after passing to the actual
Euclidean source and target coordinates; cyclically
moving the positive square-root factors proves the
formulas. The operator norm is at most that norm.
Consequently the triangle inequality in GD20 and
the rank-one leading term imply that the largest
singular value $v(A)$ of
$G_o^{1/2}C_0(A)^{-1}G_e^{\mathsf T/2}$ satisfies
\begin{equation}
 \frac{k-E(s)}{s^3}\le v(A)\le\frac{k+E(s)}{s^3},\qquad
 \big|\sigma_j(Z_A)-q_j\big|\le ds\quad(1\le j\le3).
 \tag{GD27}
\end{equation}
For completeness the second inequality is the singular-value
perturbation bound: for every unit vector the norm changes
by at most $\|Z_A-Z_0\|\le ds$; taking the minimum over
vectors in each fixed subspace and then the maximum over
subspaces in the variational characterization of the
ordered singular value gives the same bound. Interchanging
$Z_A,Z_0$ gives its absolute-value version.

If $0<s\le s_*$, then $E(s)\le k/2$ and $ds\le q_3/2$.
Moreover $v(A)\ge k/(2s^3)\ge4/q_3$, while each of the
other three inverse values is at most $2/q_3$.
Thus the inverse ordering used below is proved on this
whole explicit interval. For $\nu=3,\ldots,8$ put
$(p_3,\ldots,p_8)=(q_3,q_3,q_2,q_2,q_1,q_1)$ and set
$m=\min(j,2)$. With an empty product interpreted as one,
every exterior rank satisfies the positive bounds
\begin{equation}
 \left(\frac{k-E(s)}{s^3}\right)^m
       \prod_{\nu=3}^{j}\frac1{p_\nu+ds}
 \ \le\ \|\wedge^jB_0(A)^{-1}\|
 \ \le\
 \left(\frac{k+E(s)}{s^3}\right)^m
       \prod_{\nu=3}^{j}\frac1{p_\nu-ds},
 \qquad 1\le j\le8,\quad0<s\le s_*.
 \tag{GD28}
\end{equation}
GD27 bounds each of the two copies of the largest inverse
value. Inverting the positive bounds on the three largest
forward values bounds the remaining six inverse values
in their established order. Multiplying those inequalities
and using the exterior singular-value identity proves
GD28. This supplies explicit inverse-exterior estimates
through the actual original receiver and its dual,
with no discarded metric entries or unspecified error.



This review proves the proposed obstruction on the actual signed
completion, with its unchanged infinity orientation. It also identifies
exactly what multiplication by the cubic residue weight does to the
original rational deck map. The basis of the underlying binary-form
ring and the corresponding coefficient ideal have human provenance
in Wood's author-source equations \texttt{E:basis} and
\texttt{E:Isbasis} \cite{WoodBinary}; Wood attributes the ring basis to
Birch--Merriman and its multiplication to Nakagawa. The present
proof uses the explicit signed-extension basis and charts proved in
GB1--10 \cite{Fable042}. The original polynomial, factor states and maps
are those of FS1--17 and FS30--31 \cite{Fable037}. The perfect local
residue functional and its precise transition are RD2--4 and
RD24--26 \cite{Fable037}. No global scalar residue generator is assumed.

\section{The scheme, its original open and its three divisors}
Set $\mathfrak B=\C[A,B,C,D]$ and retain
\begin{equation}
 h(r)=Ar^4+r^3+Br^2+Cr+D,\qquad
 k(v)=A+v+Bv^2+Cv^3+Dv^4 .
 \tag{GDR1}
\end{equation}
The finite completion $X$ is covered by
\[
 X_f=\Spec\mathfrak B[r,t]/(h,t^2-h'),\qquad
 X_\infty=\Spec\mathfrak B[v,\theta]/(k,\theta^2+k'),
\]
glued by $v=1/r$, $\theta=-t/r$ with inverse $r=1/v$,
$t=-\theta/v$. On $X_f$ the equations solve exactly as
\[
 C=t^2-4Ar^3-3r^2-2Br,\qquad
 D=3Ar^4+2r^3+Br^2-rt^2,
\]
so $X_f=\mathbb A^4_{t,r,A,B}$. At $v=0$, the other chart
has $A=0$ and $\theta=\pm i$; its equation after elimination
of $A$ is $1+\theta^2+2Bv+3Cv^2+4Dv^3=0$.
The derivative in $\theta$ is $2\theta\ne0$ there. Hence $v$
vanishes to order exactly one on each infinity component:
\begin{equation}
 R_0=\{t=0\}\simeq\mathbb A^3_{r,A,B},\qquad
 D_+=\{v=0,\theta=-i\},\qquad D_-=\{v=0,\theta=i\}.
 \tag{GDR2}
\end{equation}
The infinity components are each $\mathbb A^3_{B,C,D}$;
they are disjoint from $R_0$. The multiplicity of $R_0$ in
the zero divisor of $t$ is one, because $t$ is one of the
four coordinates on $X_f$.

The resulting $X$ is smooth and integral, and finite locally free
of degree eight over $\Spec\mathfrak B$. For clarity, its finite
covering proof uses $D(A)$ and the four opens $D(h(p))$,
$p=0,1,-1,2$, which cover by
\[
 1=\tfrac12h(0)-\tfrac12h(1)-\tfrac16h(-1)+\tfrac16h(2)-2A.
\]
On $D(h(p))$, the exact coordinates are
\begin{equation}
 R=(p-r)^{-1},\quad T=t/(p-r),\qquad
 R=-v/(1-pv),\quad T=\theta/(1-pv),\qquad
 h_p(R)=R^4h(p-R^{-1}),\quad T^2=h_p'(R).
 \tag{GDR3}
\end{equation}
The leading coefficient $h(p)$ is a unit, so the monomials
$1,R,R^2,R^3,T,TR,TR^2,TR^3$ are a free basis. These
charts prove finiteness, including every multiple root.
Integrality follows from the dense affine chart $X_f$;
there is no component supported at $v=0$, since the infinity
equation is monic in $\theta$ over $\C[v,B,C,D]$.

The explicit original open immersion in FS11--15 gives
\begin{equation}
 U=X\setminus(R_0\cup D_-)\simeq\mathbb A^4_{a,y,z,w},
 \qquad a=1/t=-v/\theta.
 \tag{GDR4}
\end{equation}
In particular $D_+\cap U=\{a=0\}$, with its actual reduced
scheme structure. In the infinity chart $\theta$ is a unit
near either component, so the rational functions $t$ and $a$
have exact divisors
\begin{equation}
 \Div_X(t)=R_0-D_+-D_-,\qquad
 \Div_X(a)=D_++D_--R_0.
 \tag{GDR5}
\end{equation}
There are no other terms: $t$ is a coordinate on all of $X_f$,
and the complement of that chart consists of the two infinity
components already inspected.

\section{The coefficient functional and both lost infinity states}
Write $S=\Gamma(X,\mathcal O_X)$. Its retained global
$\mathfrak B$-basis is
\begin{equation}
 \begin{aligned}
 (e_0,e_1,e_2,e_3)&=(1,Ar,Ar^2+r,Ar^3+r^2+Br),\\
 (o_0,o_1,o_2,o_3)&=t(A,Ar+1,Ar^2+r+B,Ar^3+r^2+Br+C).
 \end{aligned}
 \tag{GDR6}
\end{equation}
The proof of freeness uses the chart matrices GB3 and GB5--9,
whose even and odd determinants are respectively $A^3,A^4$
or $h(p)^3,h(p)^4$. In particular this basis is valid at $A=0$.
Define $\Lambda_0:S\to\mathfrak B$ to be the coefficient of
$o_3$ in this global basis.

On $A\ne0$, a class $F=R+tV$ has
$\Lambda_0(F)=[r^3]V/A$, since only $o_3$ contributes to
the coefficient of $tr^3$, and that coefficient is $A$.
Consequently $\Lambda_0$ is exactly the RD finite-coordinate
residue on this open. Its full global product Gram is
\begin{equation}
 B_0=\begin{pmatrix}0&C_0\\C_0^{\mathsf T}&0\end{pmatrix},
 \qquad
 C_0=\begin{pmatrix}
 0&0&0&1\\0&0&A&0\\0&A&1&0\\A&1&B&0
 \end{pmatrix},\qquad
 \det C_0=A^3,\quad\det B_0=A^6.
 \tag{GDR7}
\end{equation}
To verify all entries directly, put
$\ell_j=[r^3](r^j\bmod h)/A$. Then
$\ell_0=\ell_1=\ell_2=0$, $\ell_3=A^{-1}$,
$\ell_4=-A^{-2}$, $\ell_5=A^{-3}-BA^{-2}$ and
$\ell_6=-A^{-4}+2BA^{-3}-CA^{-2}$ by multiplying
$h=0$ successively by $1,r,r^2$. Expanding each
$e_i(o_j/t)$ in these seven coefficients gives $C_0$.
Even--even and odd--odd products are even, and hence have
zero $\Lambda_0$. This proves GDR7 on $A\ne0$; both
sides are global regular functions, so integrality extends
the identity everywhere. Expansion along the nonzero reversed
diagonal gives its two determinant formulas. The complete
inverse cross matrix, on its admitted open, is
\begin{equation}
 C_0^{-1}=
 \begin{pmatrix}
 0&A^{-3}-BA^{-2}&-A^{-2}&A^{-1}\\
 0&-A^{-2}&A^{-1}&0\\0&A^{-1}&0&0\\1&0&0&0
 \end{pmatrix}.
 \tag{GDR8}
\end{equation}
Multiplication in both orders verifies this inverse.

At every point of $A=0$, $C_0$ has rank three and
\begin{equation}
 \operatorname{rad}B_0=\C e_1\oplus\C o_0.
 \tag{GDR9}
\end{equation}
Indeed the kernel of $C_0$ is its first odd column and the
kernel of $C_0^{\mathsf T}$ is its second even column.
The same argument works over the coefficient ring of $A=0$;
three unit pivots remain. On this whole fibre family the
relevant products are
\[
 e_1^2=-e_1,\qquad e_1o_0=-o_0,\qquad o_0^2=e_1.
\]
They follow by reducing the GDR6 polynomials modulo $h$
before setting $A=0$, or by the full multiplication tensor
GB12--17. Put $e_\infty=-e_1$. The two projectors are
\begin{equation}
 e_{\rm inc}=\frac{-e_1+i o_0}{2},\qquad
 e_{\rm omit}=\frac{-e_1-i o_0}{2}.
 \tag{GDR10}
\end{equation}
The three product identities prove that these are idempotent,
orthogonal and sum to $e_\infty$. On a finite root $e_1=o_0=0$.
On either infinity root $e_1=-1$ and $o_0=\theta$, so
$e_{\rm inc}$ equals one exactly at $\theta=-i$, while
$e_{\rm omit}$ equals one exactly at $\theta=i$.
Thus the radical consists of both entire infinity factors.
It removes the included original state as well as the missing
one; it cannot be interpreted as merely forgetting the omitted
original sign.

At $(A,B,C,D)=(0,0,-1,0)$, the six finite states have
$r=-1,0,1$ and $t^2=3r^2-1$, and the other two have
$\theta=\pm i$. The binary discriminant is $4$.
In the GDR6 basis the trace determinant is
$256\cdot4^3=2^{14}$ by GB19, although $B_0$ has rank six.
This is a direct distinction between the trace degeneration
locus and the zero locus of this particular coefficient
residue identification. The local residue in the $p=2$ chart
has determinant $6^{-8}$ by RD29 and sees all eight states.

\section{The dualizing line and the exact global obstruction}
Define the $S$-module
\begin{equation}
 \omega=\Hom_{\mathfrak B}(S,\mathfrak B),\qquad
 (f\lambda)(g)=\lambda(fg).
 \tag{GDR11}
\end{equation}
On the $p$ chart let
$\lambda_p(U+TV)=[R^3](V\bmod h_p)/h(p)$.
Its Gram matrix is the RD25 block matrix, whose determinant
is $h(p)^{-8}$, a unit on this base open. Therefore the map
$S\to\omega$, $f\mapsto f\lambda_p$, is an isomorphism
there. The corresponding assertion on $A\ne0$ follows from
the determinant $A^{-8}$ in RD4. The five opens cover the
base, so $\omega$ is an invertible $S$-module. This conclusion
uses explicit perfect local pairings even on the discriminant.

On the overlap with $A\ne0$, change of variable gives
$dr/h=R^2\,dR/h_p$; the odd coefficient acquires the further
factor $R^{-1}$ from $t=T/R$. Hence
\begin{equation}
 \Lambda_0=R^3\lambda_p,\qquad
 \Tr=2T^3\lambda_p.
 \tag{GDR12}
\end{equation}
The first equality means equality as sections of $\omega$,
with the $S$-action specified in GDR11. Both sections are
regular on the entire $p$ chart, and equality on the dense
$A\ne0$ open proves it on that chart, including $R=0$.
The trace equality follows either from RD6 applied to $h_p$
or directly as follows. The trace of $U+TV$ is twice the
root-algebra trace of $U$. At a root of multiplicity $m$,
that root trace is $mU(r_0)$, which is the residue of
$h_p'U/h_p$. The odd coefficient of $2T^3(U+TV)$ is
$2h_p'U$. This proves the second equality on nonreduced
fibres as well.

Because $R=-v/(1-pv)$, its zero is simple on each infinity
component. It has no zero at a finite selected root in this
chart, where $R=(p-r)^{-1}$ is a unit. Since $T=tR$, $T$
vanishes simply on $R_0$ and is the unit $\theta=\pm i$ at
infinity. The local generators in GDR12 therefore prove
the exact divisors of these sections of the line $\omega$:
\begin{equation}
 \Div_X(\Lambda_0)=3D_++3D_-,\qquad
 \Div_X(\Tr)=3R_0,\qquad
 \Lambda_0=\frac{a^3}{2}\Tr
 \quad\hbox{as rational sections of }\omega.
 \tag{GDR13}
\end{equation}
The last equality uses $T=tR$ and $a=1/t$; it is a regular
equality on $X\setminus R_0$. In GDR13 a divisor of a section
is taken on $X$, where $\omega$ has rank one. It is not the
zero divisor of a rank-eight vector over the base.

There is no global generator of $\omega$. Here is the complete
argument. If $\ell$ were such a generator, there would be a
unique regular $f\in S$ with $\Lambda_0=f\ell$. Since $\ell$
vanishes nowhere, GDR13 would force
$\Div_X(f)=3D_++3D_-$. Restrict to the actual original open
$U=\Spec\C[a,y,z,w]$. Its coordinate ring is a unique
factorization domain. The nonzero polynomial $f|_U$ has the
single zero divisor $3\{a=0\}$, so its irreducible-factor
decomposition is exactly $f|_U=c a^3$ with $c\in\C^*$.
Equality on the dense open $U$ proves $f=c a^3$ as a rational
function on $X$. But GDR5 gives order $-3$ along $R_0$ for
that function. A regular function cannot have negative order
along a divisor of the smooth $X$. This is a contradiction.

The same proof applies to every positive tensor power. A global
generator of $\omega^{\otimes m}$ would express
$\Lambda_0^{\otimes m}=f\ell$ with
$\Div f=3m(D_++D_-)$. Restriction to $U$ forces
$f=c a^{3m}$, of order $-3m$ along $R_0$. Consequently
\begin{equation}
 \omega^{\otimes m}\not\simeq\mathcal O_X
 \quad\hbox{for every integer }m\ge1.
 \tag{GDR14}
\end{equation}
In particular the class of $\omega$ has infinite order in
$\operatorname{Pic}(X)$. This statement neither asserts nor
requires a presentation of the entire Picard group.

A $\mathfrak B$-bilinear associative pairing $b:S\times S\to
\mathfrak B$ is necessarily $b(f,g)=\ell(fg)$ with
$\ell(h)=b(1,h)$: associativity gives
$b(f,g)=b(1,fg)$. Its identification $S\to\omega$ is
exactly $f\mapsto f\ell$. It is perfect if and only if this
map is an isomorphism, equivalently if $\ell$ generates
$\omega$. The obstruction just proved therefore rules out
every globally perfect associative scalar pairing, not merely
the candidate $\Lambda_0$. This does not rule out arbitrary
nonassociative bilinear forms on the free base module.

There is a useful exact product consequence. For $m$ copies of
the cover over the product base, finite freeness gives
\[
 \Hom_{\mathfrak B^{\otimes_\C m}}
 (S^{\otimes_\C m},\mathfrak B^{\otimes_\C m})
 \simeq \omega\boxtimes\cdots\boxtimes\omega.
\]
The map sends a tensor of functionals to its product on a tensor
of arguments. The tensor products of the GDR6 basis and its
dual prove that map is an isomorphism, including its module
action. If this line on $X^m$ were trivial, pullback along the
diagonal $X\to X^m$ would trivialize $\omega^{\otimes m}$,
contradicting GDR14. Thus taking products cannot by itself
produce a global generator of this dualizing line. This is
an obstruction to that specified scalar construction, with
the exact pullback map proved; it is not a conclusion about
other amplification arguments.

The domain of this obstruction matters. On the actual open
$A\ne0$, $\Lambda_0$ itself is a generator, by GDR7.
Restricting further to the original quartic slice $A=-1/2$
leaves its pairing determinant exactly $1/64$, a unit.
Thus GDR14 is a global assertion on the whole completion,
not an obstruction to the already constructed perfect
scalar residue on that fixed slice. Its map to the fixed
slice is the explicit base restriction with generator
$\Lambda_0|_{A=-1/2}$ and inverse GDR8 at $A=-1/2$.

\section{The global replacement and the original receivers}
The canonical evaluation
\begin{equation}
 S\times\omega\longrightarrow\mathfrak B,
 \qquad (f,\ell)\longmapsto\ell(f)
 \tag{GDR15}
\end{equation}
is globally perfect as a pairing of base modules. In the basis
$b=(e_0,\ldots,e_3,o_0,\ldots,o_3)$ and its coefficient-dual
basis $b^*$, its matrix is $I_8$. For multiplication by any
$s\in S$, if $M_s$ is the coefficient matrix on $b$, the
dual action is $M_s^{\mathsf T}$, since
$(s\ell)(f)=\ell(sf)$. This proves both the balance identity
and the exact correspondence between the two representations.
It retains the dual module as the second argument.

More explicitly, let $J:\C^8\to\mathcal H$ be any one of
the fixed invertible original receivers, with all its moments
and masses retained. A primal coefficient column $x$ is sent
to $Jx$ and a dual coefficient column $\eta$ is sent to the
linear functional $\eta^{\mathsf T}J^{-1}$ on $\mathcal H$.
Then
\begin{equation}
 (\eta^{\mathsf T}J^{-1})(Jx)=\eta^{\mathsf T}x,
 \qquad M_s\longmapsto JM_sJ^{-1}.
 \tag{GDR16}
\end{equation}
These two displayed inverse maps prove evaluation and
multiplication transport exactly. If the retained positive
Gram in the $b$ coordinates is $G=J^*J$ (with the original
inner product on $\mathcal H$ understood), the norm of the
functional is
\begin{equation}
 \|\eta^{\mathsf T}J^{-1}\|^2
   =\eta^{\mathsf T}G^{-1}\overline\eta
   =\eta^*G^{-\mathsf T}\eta.
 \tag{GDR17}
\end{equation}
To see this, write an input as $G^{-1/2}v$ with $\|v\|=1$;
Cauchy--Schwarz gives the first value, attained by the
adjoint of the row $\eta^{\mathsf T}G^{-1/2}$. Hermitian
symmetry gives the second expression. This uses the original
Gram, not the residue product matrix $B_0$, as the positive
metric.

\section{What the cubic weight does to the actual deck map}
Use precisely the original source coordinates $q=(a,y,z,w)$.
The rational factor-sign involution is
\begin{equation}
 \Sigma(q)=\left(-a,\ y+\frac{2i}{a},\ -z+\frac{6i}{a^2},\
 w-\frac{14iy^2}{a}+\frac{28y}{a^2}+\frac{40i}{a^3}\right)
 \quad(a\ne0).
 \tag{GDR18}
\end{equation}
FS30--31 prove that it extends to the regular sign involution
on $X$, interchanging $D_+$ and $D_-$. In the original open
it is only defined where $a\ne0$. Substitution verifies
$\Sigma^2=1$ and $P\Sigma=P$ with every original target
coordinate unchanged. The exact checker accompanying this
proof independently reconstructs FS1, verifies its
resultant-one relation and $\det DP=-2$, and checks both
identities as rational identities.

On a split etale fibre, the trace functional has weight one
at each sheet. GDR12 gives the residue weight there as
$R^3/(2T^3)=a^3/2$. At an infinity sheet it is zero, since
$R=0$ and $T=\pm i$. Thus GDR13 identifies the exact cubic
weight of $\Lambda_0$ on the selected unramified sheet,
including its vanishing on $D_+$.
Multiplying the original coordinate vector GDR18 by this
weight gives the polynomial map
\begin{equation}
 F(q)=\frac{a^3}{2}\Sigma(q)=
 \begin{pmatrix}
 -a^4/2\\a^3y/2+ia^2\\-a^3z/2+3ia\\
 a^3w/2-7ia^2y^2+14ay+20i
 \end{pmatrix},\qquad
 F(0,y,z,w)=(0,0,0,20i).
 \tag{GDR19}
\end{equation}
This is equality of the four coordinate functions, including
all constants. It proves the complete cancellation of their
displayed pole orders by the exact residue weight. It does not
replace the rational involution by an invertible extension in
the original affine coordinates.

In fact the complete differential records the resulting loss:
\begin{equation}
 DF=\begin{pmatrix}
 -2a^3&0&0&0\\
 3a^2y/2+2ia&a^3/2&0&0\\
 -3a^2z/2+3i&0&-a^3/2&0\\
 3a^2w/2-14iay^2+14y&-14ia^2y+14a&0&a^3/2
 \end{pmatrix},
 \quad\det DF=\frac{a^{12}}4.
 \tag{GDR20}
\end{equation}
The determinant is the product of the displayed diagonal
entries. At $a=0$ the sole nonzero column is
$(0,0,3i,14y)^{\mathsf T}$, so the rank is exactly one
at every such point and the kernel is
$\operatorname{span}(\partial_y,\partial_z,\partial_w)$.
The map collapses the entire $D_+\simeq\mathbb A^3$ in the
original open to the single vector in GDR19. This singularity
belongs to the explicitly constructed weighted map $F$.
The independently verified $\det DP=-2$ is unchanged.

Retain the actual four-dimensional original linear receivers
$\Phi$ and $\Psi=T_A\Phi$, with their positive source and
target norms. They are held fixed while $q$ tends to its
boundary. Put
$G_U=(\langle\Phi e_j,\Phi e_k\rangle)_{j,k=1}^4$ and
$G_W=(\langle\Psi e_j,\Psi e_k\rangle)_{j,k=1}^4$,
using all the original moments and masses in those inner
products. The exact transported weighted maps are
\begin{equation}
 F_U=\Phi F\Phi^{-1},\qquad F_W=\Psi F\Psi^{-1},\qquad
 T_AF_U=F_WT_A.
 \tag{GDR21}
\end{equation}
Their derivatives are the corresponding conjugates of GDR20;
they therefore have determinant $a^{12}/4$ and rank one at
the corresponding boundary. Their kernels there are exactly
$\Phi\operatorname{span}(e_2,e_3,e_4)$ and
$\Psi\operatorname{span}(e_2,e_3,e_4)$.

For any approach with $a\to0$, $a\ne0$, and bounded
$y,z,w$, the four explicit coordinates in GDR19 converge to
$20ie_4$. Continuity and complex homogeneity of the fixed
norms prove the following limits
\begin{equation}
 \begin{aligned}
 \frac{a^3}{2}\Phi\Sigma(q)&\longrightarrow20i\Phi e_4,
 &\left\|\frac{a^3}{2}\Phi\Sigma(q)\right\|_U
     &\longrightarrow20\sqrt{(G_U)_{44}},\\
 \frac{a^3}{2}\Psi\Sigma(q)&\longrightarrow20i\Psi e_4,
 &\left\|\frac{a^3}{2}\Psi\Sigma(q)\right\|_W
     &\longrightarrow20\sqrt{(G_W)_{44}},\\
 |a|^3\|\Phi\Sigma(q)\|_U&\longrightarrow40\sqrt{(G_U)_{44}},
 &|a|^3\|\Psi\Sigma(q)\|_W
     &\longrightarrow40\sqrt{(G_W)_{44}}.
 \end{aligned}
 \tag{GDR22}
\end{equation}
No direction or phase of the complex approach is restricted:
the vector limits use $a^3$, and the norm limits use
$|a|^3$. The nonzero leading constants are strictly positive
because both receivers are invertible and both original
norms are positive. Thus the unweighted original source and
target vectors still escape cubically. The residue multiplier
removes that pole in the weighted coordinate vector by its
own cubic zero; it does not cancel the divergence of the
unweighted norms and does not prove a kernel or order jump
of the original linear $T_A$.

\section{Independent verification of the complete original-metric inverse cost}
This section reviews GD19--21 in the root calculation
\texttt{GLOBAL\_RESIDUE\_DUALITY\_BODY.tex}. Let the retained
global state Gram be $\Gamma=\diag(G_e,G_o)$, with $G_e,G_o$
positive Hermitian four-by-four matrices. They are held fixed
while the complex parameter $A$ tends to zero and $B,C,D$
are held fixed. By GDR17, the dual Gram is $\Gamma^{-\mathsf T}$.
For a positive Hermitian matrix, every half power below denotes
the unique positive Hermitian square root, obtained by unitary
diagonalization. The notation $G^{\mathsf T/2}$ means
$(G^{1/2})^{\mathsf T}$, the positive square root of
$G^{\mathsf T}$.

The map $B_0$ sends primal coefficient columns to dual columns.
Its matrix after isometric passage to Euclidean coordinates is
\begin{equation}
 \widehat B_A=\Gamma^{-\mathsf T/2}B_0\Gamma^{-1/2}
  =\begin{pmatrix}0&Z_A\\Z_A^{\mathsf T}&0\end{pmatrix},
 \qquad Z_A=G_e^{-\mathsf T/2}C_0G_o^{-1/2}.
 \tag{GDR23}
\end{equation}
Indeed multiplication by $\Gamma^{1/2}$ turns the primal norm
into the Euclidean norm, and multiplication by
$\Gamma^{-\mathsf T/2}$ does the same for dual columns.
The bottom-left block is the transpose of the top-right block
by direct transposition of its three factors. A complex matrix
and its transpose have the same singular values: transpose a
singular value decomposition and conjugate its unitary factors.
The two orthogonal block subspaces in GDR23 therefore show
that each of the four singular values of $Z_A$ occurs twice
in $\widehat B_A$.

At $A=0$, $C_{00}=C_0|_{A=0}$ has rank three. The three
positive singular values of $Z_0$ will be denoted
$q_1\ge q_2\ge q_3>0$. They are finite and unchanged by
the vanishing fourth singular value. Put
\begin{equation}
 \begin{aligned}
 L&=G_o^{-1}C_{00}^*G_e^{-\mathsf T}C_{00},\qquad
 k^2=(G_o)_{00}(G_e)_{11},\\
 a_1&=\Tr L,\quad
 a_2=\tfrac12\big((\Tr L)^2-\Tr(L^2)\big),\quad
 a_3=\frac{k^2}{\det G_e\det G_o}.
 \end{aligned}
 \tag{GDR24}
\end{equation}
The matrices $L$ and $Z_0^*Z_0$ are similar: conjugate $L$
by $G_o^{1/2}$. Thus the eigenvalues of $L$ are
$q_1^2,q_2^2,q_3^2,0$. The first two elementary symmetric
functions are $a_1,a_2$. For the third, multiply GDR8 by
$\det C_0=A^3$ and set $A=0$ to get
$\operatorname{adj}C_{00}=u_0u_1^{\mathsf T}$, where the
$u_j$ are the numerical coordinate unit columns in $\C^4$.
For invertible $U,V$ the cofactor identity is
\[
 \operatorname{adj}(UXV)=\det(UV)V^{-1}
          \operatorname{adj}(X)U^{-1}.
\]
It follows first for invertible $X$ from
$\operatorname{adj}X=(\det X)X^{-1}$, and then for every
$X$ because both sides are polynomials in its entries.
Apply it to $U=G_e^{-\mathsf T/2}$, $V=G_o^{-1/2}$
and $X=C_{00}$. The adjugate has rank one; its squared
Frobenius norm is
\[
 \frac{\|G_o^{1/2}u_0\|^2\,
        \|u_1^{\mathsf T}G_e^{\mathsf T/2}\|^2}
      {\det G_e\det G_o}
 =\frac{(G_o)_{00}(G_e)_{11}}{\det G_e\det G_o}.
\]
By the singular value decomposition its unique nonzero
singular value is $q_1q_2q_3$: the adjugate is, up to the
unitary exterior identifications, the third exterior map.
This proves $q_1^2q_2^2q_3^2=a_3$. Consequently all three
values are specified exactly as the positive roots of
$z^3-a_1z^2+a_2z-a_3$, with their order retained.

The inverse Euclidean matrix has blocks
\begin{equation}
 \widehat B_A^{-1}=
 \begin{pmatrix}
 0&G_e^{1/2}C_0^{-\mathsf T}G_o^{\mathsf T/2}\\
 G_o^{1/2}C_0^{-1}G_e^{\mathsf T/2}&0
 \end{pmatrix}.
 \tag{GDR25}
\end{equation}
These blocks are mutual transposes. GDR8 gives the exact
limit, including the complex phase of $A$,
\[
 A^3G_o^{1/2}C_0^{-1}G_e^{\mathsf T/2}
 \longrightarrow G_o^{1/2}u_0u_1^{\mathsf T}G_e^{\mathsf T/2}.
\]
The right side has rank one and norm $k$. The other three
singular values of this inverse block converge to
$q_3^{-1},q_2^{-1},q_1^{-1}$, because the three positive
singular values of its inverse converge to $q_1,q_2,q_3$.
The fourth inverse singular value is therefore asymptotic
to $k|A|^{-3}$. Doubling gives exactly two divergent singular
values for the full eight-dimensional map, and exactly six
finite values in the descending order
\begin{equation}
 k|A|^{-3},\ k|A|^{-3},\
 q_3^{-1},\ q_3^{-1},\ q_2^{-1},\ q_2^{-1},\ q_1^{-1},\ q_1^{-1},
 \tag{GDR26}
\end{equation}
where the first two entries denote asymptotic equivalents
and the other six denote limits. The use of a complex
transpose in GDR23--25 is essential; replacing it by an
adjoint changes the map to the dual.

For $d=(k,k,q_3^{-1},q_3^{-1},q_2^{-1},q_2^{-1},q_1^{-1},q_1^{-1})$,
the induced exterior maps consequently satisfy
\begin{equation}
 \lim_{A\to0}|A|^{3\min(j,2)}
       \|\wedge^j B_0(A)^{-1}\|
   =\prod_{\nu=1}^j d_\nu\quad(1\le j\le8),\qquad
 \|\wedge^8 B_0(A)^{-1}\|
       =\frac{\det\Gamma}{|A|^6}.
 \tag{GDR27}
\end{equation}
The exterior norm is the product of the largest singular
values, since the exterior basis of a singular value
decomposition diagonalizes that map. This proves every
limit from GDR26. For the exact full-rank equality, the
determinant of GDR25 is
$\det\Gamma/\det B_0=\det\Gamma/A^6$; its absolute value
is the desired norm. The limiting product is the same
constant because $k^2/(q_1^2q_2^2q_3^2)=\det\Gamma$.

The finite-parameter bounds GD26--28 also hold in the same
metric. To provide their independent proof, write
\begin{equation}
 \begin{aligned}
 C_0(A)&=C_{00}+AC_1,\qquad
 C_1=\begin{pmatrix}0&0&0&0\\0&0&1&0\\0&1&0&0\\1&0&0&0\end{pmatrix},\\
 C_0(A)^{-1}&=A^{-3}K_3+A^{-2}K_2+A^{-1}K_1+K_0,\\
 K_3&=u_0u_1^{\mathsf T},\qquad
 K_2=-Bu_0u_1^{\mathsf T}-u_0u_2^{\mathsf T}-u_1u_1^{\mathsf T},\\
 K_1&=u_0u_3^{\mathsf T}+u_1u_2^{\mathsf T}+u_2u_1^{\mathsf T},
 \qquad K_0=u_3u_0^{\mathsf T}.
 \end{aligned}
 \tag{GDR28}
\end{equation}
For $s=|A|>0$ put
\begin{equation}
 \begin{aligned}
 d^2&=\Tr(G_o^{-1}C_1^*G_e^{-\mathsf T}C_1),\qquad
 b_j^2=\Tr(G_e^{\mathsf T}K_j^*G_oK_j)\quad(0\le j\le2),\\
 E(s)&=b_2s+b_1s^2+b_0s^3,\\
 s_*&=\min\left\{1,\frac{q_3}{2d},
 \frac{k}{2(b_2+b_1+b_0)},
 \left(\frac{kq_3}{8}\right)^{1/3}\right\}.
 \end{aligned}
 \tag{GDR29}
\end{equation}
Each trace is the squared Hilbert--Schmidt norm of its
actual Euclidean coefficient matrix: for $d$ that matrix
is $G_e^{-\mathsf T/2}C_1G_o^{-1/2}$, and for $b_j$ it
is $G_o^{1/2}K_jG_e^{\mathsf T/2}$. Cyclic invariance of
trace proves the displayed expressions. These norms are
positive when the middle matrix is nonzero; $C_1,K_0$ are
nonzero. Therefore $d>0$, $b_0>0$ and $s_*>0$.

Every matrix operator norm is at most its Hilbert--Schmidt
norm, since for a unit vector the squared norm of its image
is bounded by the sum of the squared singular values.
The triangle inequality applied to $A^3$ times the inverse
block in GDR25 consequently gives
\[
 \frac{k-E(s)}{s^3}\le v(A)\le\frac{k+E(s)}{s^3},
 \qquad |\sigma_j(Z_A)-q_j|\le ds\quad(1\le j\le3),
\]
where $v(A)$ is the largest inverse singular value.
For the second assertion, on every unit vector the two
norms differ by at most $\|Z_A-Z_0\|\le ds$. Taking
the minimum within each fixed-dimensional subspace and
then the maximum over such subspaces in the singular
value variational characterization preserves this bound.
Reversing the two matrices proves the absolute-value form.

For $0<s\le s_*$, the first three terms defining the threshold
give $E(s)\le k/2$ and $ds\le q_3/2$. Its fourth term gives
$v(A)\ge k/(2s^3)\ge4/q_3$. The other three inverse values
are at most $2/q_3$ because each relevant forward value is
at least $q_3/2$. This proves the ordering throughout that
whole interval. For $p_3,p_4,p_5,p_6,p_7,p_8=
q_3,q_3,q_2,q_2,q_1,q_1$ and $m_j=\min(j,2)$, multiplying
the individual positive bounds gives
\begin{equation}
 \left(\frac{k-E(s)}{s^3}\right)^{m_j}
 \prod_{\nu=3}^{j}\frac1{p_\nu+ds}
 \le\|\wedge^j B_0(A)^{-1}\|
 \le\left(\frac{k+E(s)}{s^3}\right)^{m_j}
 \prod_{\nu=3}^{j}\frac1{p_\nu-ds},
 \quad 1\le j\le8,\quad 0<s\le s_*.
 \tag{GDR30}
\end{equation}
An empty product is one. All denominators are strictly
positive by the preceding threshold. This proves every
finite-parameter bound and its domain, including complex
$A$ with arbitrary phase.

\section{Verification scope}
\texttt{GLOBAL\_DUALIZING\_REVIEW.py} computes the sixteen
residue entries independently from polynomial division,
checks both inverse products and all radical/projector
relations, reconstructs the eight-state evaluation at the
original target, verifies every original polynomial and deck
identity used above, and computes the weighted differential
and its rank exactly. Its certificate records the checked
source hashes. The geometric arguments proving divisor
multiplicities, the global obstruction and the tensor-power
obstruction are given in full in GDR1--14. They are not
represented as consequences of a finite numerical test.


\section{A complete eight-state ES-to-paired-quartet family}
The received local comparison \cite[Section 5]{ESRHReceived}
maps one length-four ES primary algebra to one of the two
paired-quartet primary algebras. Here we construct the whole
finite flat family, including the other four ES states and
their final collision. The endpoint is the formal quartet
polynomial of RD13; no assertion that its roots are zeros
of the arithmetic zeta function is made. All original
labels, factors, cofactor corrections and receiver metrics
are specified below. The original ES marking and polynomial
provenance remain \cite{ES488,Tao}.

Fix $p>0$ and $\gamma>0$. The exact reciprocal identity is
$4/p=1/p+1/(2p/5)+1/(2p)$. The ES tuple is therefore
$(p;x,y,z_{\rm ES})=(p;p,2p/5,2p)$; $p=5$ gives the
integral example $(5;5,2,10)$. This is not claimed to be
in a hard-prime subdomain. Explicitly place its four
root slots in the order
$(\mathrm M,\mathrm N,\mathrm T,\mathrm E)=(p,2p/5,p,2p)$.
This declares the permutation from $(p,x,y,z_{\rm ES})$;
no historical label is silently replaced. For $0\le z\le1$
and $\tau=1-z$, set
\begin{equation}
\begin{aligned}
 \alpha&=p\tau+z(\tfrac12+i\gamma),&
 \beta&=\tfrac{2p}{5}\tau+z(\tfrac12-i\gamma),\\
 \chi&=2p\tau+z(\tfrac12-i\gamma),&
 S&=2\alpha+\beta+\chi=\tfrac{22p}{5}\tau+2z>0,\\
 A_z&=-S^{-1},&
 h_z(r)&=A_z(r-\alpha)^2(r-\beta)(r-\chi).
\end{aligned}\tag{EF1}
\end{equation}
The marked slots are now $(\alpha,\beta,\alpha,\chi)$.
Expanding gives precisely
\begin{equation}
\begin{aligned}
 h_z(r)&=A_zr^4+r^3+B_zr^2+C_zr+D_z,\\
 B_z&=A_z[\alpha^2+2\alpha(\beta+\chi)+\beta\chi],\\
 C_z&=-A_z[\alpha^2(\beta+\chi)+2\alpha\beta\chi],\qquad
 D_z=A_z\alpha^2\beta\chi,\\
 h_0(r)&=-\frac5{22p}(r-p)^2(r-2p/5)(r-2p),\qquad
 h_1(r)=-\tfrac12[(r-\tfrac12)^2+\gamma^2]^2.
\end{aligned}\tag{EF2}
\end{equation}
In particular the coefficient of $r^3$ is one at every
parameter, and the leading coefficient is retained.
The entire family lies on the discriminant because
the M/T root is double throughout.

For all computations keep the explicit change and inverse
\begin{equation}
 m=(\beta+\chi)/2=\tfrac{6p}{5}\tau+z(\tfrac12-i\gamma),\quad
 x=r-m,\quad r=x+m,\quad
 d=\tfrac{4p}{5}\tau,\quad k=\alpha-m=-\tfrac p5\tau+2i\gamma z,
 \quad\kappa=k^2-d^2.
 \tag{EF3}
\end{equation}
Then $h_z=A_z(x-k)^2(x^2-d^2)$. Neither $k-d$
nor $k+d$ vanishes on the parameter interval: their
real parts are respectively $-p\tau$ and $3p\tau/5$,
and their imaginary part is $2\gamma z$. Vanishing
would require $\tau=z=0$, which is impossible.
Thus $\kappa$ is a unit on an open complex neighborhood
of the entire interval. Over the ring of polynomial
functions of $z$ with $S\kappa$ inverted, the algebra
\begin{equation}
 \mathscr S_z=\mathbb C[z,(S\kappa)^{-1}][r,t]/
                  (h_z(r),t^2-h_z'(r))
 \tag{EF4}
\end{equation}
is free of rank eight with the original basis RD1.
Successive division by the two relations proves freeness;
all eight coefficient states persist at both endpoints.

\section{Complete primary projectors, their inverse and original return}
Put $a_0=k^2+d^2$ in this section only; this symbol is
not the original source coordinate $a$. The exact two
primary projectors are
\begin{equation}
 e_U=\frac{(x^2-d^2)(3k^2-d^2-2kx)}{\kappa^2},\qquad
 e_L=\frac{(x-k)^2(a_0+2kx)}{\kappa^2},\qquad
 e_U+e_L=1.
 \tag{EF5}
\end{equation}
The last equation is an identity of polynomials.
Modulo $(x-k)^2$, the first is one, because it is
the first-order inverse of $x^2-d^2$; modulo
$x^2-d^2$ it is zero. The second has the reverse
congruences, since
$(a_0-2kx)(a_0+2kx)\equiv\kappa^2$ modulo $x^2-d^2$.
The two factors are coprime because $\kappa$ is a unit.
Consequently $e_U^2=e_U$, $e_L^2=e_L$, $e_Ue_L=0$
in the exact quartic algebra, including $d=0$.
Their central extension gives
$\mathscr S_z=e_U\mathscr S_z\oplus e_L\mathscr S_z$.

Here are every coefficient and the inverse of that
decomposition. In increasing powers of $x$, the
columns of $e_U,e_U(x-k)$ and $e_L,e_Lx$ are
\begin{equation}
 P_U=\frac1{\kappa^2}\begin{pmatrix}
 d^2(d^2-3k^2)&kd^2\kappa\\2kd^2&-d^2\kappa\\
 3k^2-d^2&-k\kappa\\-2k&\kappa
 \end{pmatrix},\qquad
 P_L=\frac1{\kappa^2}\begin{pmatrix}
 k^2a_0&2k^3d^2\\-2kd^2&k^2(k^2-3d^2)\\
 d^2-3k^2&-2k^3\\2k&a_0
 \end{pmatrix}.
 \tag{EF6}
\end{equation}
Multiplication by $x$ is reduced only by the original
quartic relation
$x^4=2kx^3-\kappa x^2-2kd^2x+k^2d^2$;
expanding EF5 and using that relation gives EF6.
For a cubic coefficient column $c$ the inverse is
\begin{equation}
 [P_U\ P_L]^{-1}c=
 \begin{pmatrix}1&k&k^2&k^3\\0&1&2k&3k^2\\
 1&0&d^2&0\\0&1&0&d^2\end{pmatrix}c,
 \qquad\det[P_U\ P_L]=\kappa^{-2}.
 \tag{EF7}
\end{equation}
The first two rows take the value and derivative at $x=k$.
The last two take the remainder modulo $x^2-d^2$.
Applying these four functionals to EF6 gives $I_4$;
the displayed square matrices therefore are mutual
inverses. Their determinant follows either by expansion
or the same four rows, whose determinant is $\kappa^2$.

In original $r$ coefficients use the full binomial
matrix $C_{-m}$ of MR1, since $x=r-m$. The complete
eight-dimensional primary inclusion and inverse are
\begin{equation}
Q_{\mathrm{prim}}=
 \operatorname{diag}(C_{-m}[P_U\ P_L],C_{-m}[P_U\ P_L]),\qquad
Q_{\mathrm{prim}}^{-1}=
 \operatorname{diag}([P_U\ P_L]^{-1}C_m,[P_U\ P_L]^{-1}C_m).
 \tag{EF8}
\end{equation}
The column order is upper even, lower even, upper odd,
lower odd, with two columns in each group in the order
shown in EF6. Its source and target returns are the exact
$F_UQ_{\mathrm{prim}}$ and $F_WQ_{\mathrm{prim}}$ of
GD14, with $T_A\oplus T_A$ between them. Its marked
column is $Q_r^{-1}Q_{\mathrm{prim}}z_{\mathrm{prim}}$.
All inverse coefficients are EF7--8 and FC24, and
their products are the identity. Thus the changing
root center $m$ is not discarded in the physical return.

The upper and lower primary relations are respectively
\begin{equation}
\begin{aligned}
 U_z&=\mathbb C[\varepsilon,t]/(\varepsilon^2,t^2-2g_0\varepsilon),
 &\varepsilon&=x-k,&g_0&=A_z\kappa,\quad g_1=2A_zk,\\
 L_z&=\mathbb C[x,t]/(x^2-d^2,t^2-u x-v),
 &u&=2A_z(k^2+d^2),&v&=-4A_zkd^2.
\end{aligned}\tag{EF9}
\end{equation}
The formulas follow by reducing the full derivative
$h_z'$ modulo the indicated root factors. Both algebras
have dimension four, with bases $(1,\varepsilon,t,t\varepsilon)$
and $(1,x,t,tx)$. Since $g_0\ne0$, the first is the
length-four algebra with $\varepsilon=t^2/(2g_0)$,
$t^4=0$ and basis $(1,t,t^2,t^3)$.
At $z=1$, $d=0$, $k=2i\gamma$, $A_z=-1/2$,
and the lower relation becomes $t^2=4\gamma^2x$,
$x^2=0$. Thus the whole endpoint has two length-four
primary algebras. At $z=0$ the lower roots are
distinct and $h_0'$ is nonzero there, so that summand
has four distinct signed states. The upper summand is
already nonreduced at every parameter.

\section{The complete residue corrections along the family}
The upper root cofactor at $\varepsilon=0$ is
$A_z[(\varepsilon+k)^2-d^2]=g_0+g_1\varepsilon+A_z\varepsilon^2$.
Expansion of its reciprocal through the exact quotient
$\varepsilon^2=0$ gives
\begin{equation}
 \Lambda_U(t(b_0+b_1\varepsilon))
   =\frac{b_1}{A_z\kappa}-\frac{2k b_0}{A_z\kappa^2},\qquad
 \Lambda_U(1)=\Lambda_U(\varepsilon)=0.
 \tag{EF10}
\end{equation}
For the lower component the inverse of its full cofactor
$A_z(x-k)^2$ modulo $x^2-d^2$ is
$(a_0+2kx)/(A_z\kappa^2)$. Taking the linear coefficient
of its product with $b_0+b_1x$ proves
\begin{equation}
 \Lambda_L(t(b_0+b_1x))
   =\frac{a_0b_1+2kb_0}{A_z\kappa^2},\qquad
 B_L=\begin{pmatrix}0&J_L\\J_L&0\end{pmatrix},\quad
 J_L=\frac1{A_z\kappa^2}
       \begin{pmatrix}2k&a_0\\a_0&2kd^2\end{pmatrix},\quad
 \det B_L=\frac1{A_z^4\kappa^4}.
 \tag{EF11}
\end{equation}
Both formulas are the restrictions of the full RD2
residue under the exact projectors EF5. They detect
every nilpotent at $d=0$. All cofactor terms are present.

At the original ES upper root, $g_{0,E}=3p/22$ and
$g_{1,E}=1/11$. Along the interval choose the continuous
square root $c(z)^2=g_{0,E}/g_0(z)$ with $c(0)=1$;
it exists since the nonzero continuous function has a
continuous argument on an interval. The complete local
algebra map and residue correction are
\begin{equation}
\begin{aligned}
 \phi_z(\varepsilon_E)&=\varepsilon,&\phi_z(t_E)&=c(z)t,\\
 \Lambda_U(\phi_z f)&=\Lambda_E(w_E(z)f),&
 w_E(z)&=c(z)^3\left[1+
       \left(\frac2{3p}-\frac{2k}{\kappa}\right)\varepsilon_E\right],\\
 w_E(z)^{-1}&=c(z)^{-3}\left[1-
       \left(\frac2{3p}-\frac{2k}{\kappa}\right)\varepsilon_E\right].
\end{aligned}\tag{EF12}
\end{equation}
The two relations in EF9 are preserved because
$c^2g_0=g_{0,E}$; the inverse sends $t$ to $c^{-1}t_E$
and $\varepsilon$ to $\varepsilon_E$. Evaluate both
residue functionals on the full odd basis
$t_E,t_E\varepsilon_E$ using EF10. Their two values
agree exactly by the displayed coefficient; the even
basis has zero residue. This proves the identity on
every element. Its stated inverse uses
$\varepsilon_E^2=0$. At $z=1$ the bracket is
$1+(2/(3p)+i/\gamma)\varepsilon_E$, and
$c(1)^2=3p/(44\gamma^2)$, reproducing the full
received upper local correction with its original sign.
EF5--11 retain the four lower states absent from a
comparison of just those two local blocks.

\section{The four remaining signed branches and their exact inverse cost}
For $z<1$, $d>0$ and the lower root/sign states are
\begin{equation}
 x_\epsilon=\epsilon d,\quad r_\epsilon=m+\epsilon d,\quad
 t_{\epsilon,\sigma}=\sigma(\epsilon d-k)\sqrt{2A_z\epsilon d},
 \qquad\epsilon,\sigma\in\{1,-1\}.
 \tag{EF13}
\end{equation}
Choose $\sqrt{2A_zd}=i\sqrt{-2A_zd}$ and
$\sqrt{-2A_zd}>0$. Both are continuous on the interval.
Then $t^2=2A_z\epsilon d(\epsilon d-k)^2=h_z'(r_\epsilon)$,
and every one of the four values is nonzero.
The exact original source state is FS13,
\begin{equation}
 q_{\epsilon,\sigma}=
 \left(t^{-1},-r-it,A_zt^3+2rt+3it^2,
 7ir^2t+(B_z-17r+A_zr^2)t^2-13it^3-2A_zt^4\right),
 \quad(r,t)=(r_\epsilon,t_{\epsilon,\sigma}).
 \tag{EF14}
\end{equation}
In particular these are actual finite inverse states of
the original polynomial for every $z<1$, and not just
points of a coefficient diagram. The N slot is
$\epsilon=-1$ and the E slot is $\epsilon=1$, with
both factor signs $\sigma$ retained. The M/T primary
state has $t=0$ throughout and belongs to the completed
boundary, not to an additional finite inverse branch.

As $z\to1$, $r_\epsilon\to1/2-i\gamma$ and
$t_{\epsilon,\sigma}/\sqrt d\to
\sigma(-2i\gamma)\sqrt{-\epsilon}$, with the square
roots just specified. Multiplication of every coordinate
in EF14 by $\sqrt d$ gives
\begin{equation}
 \sqrt d\,q_{\epsilon,\sigma}
 \longrightarrow\frac{\mathbf e_1}
       {\sigma(-2i\gamma)\sqrt{-\epsilon}},\qquad
 \sqrt d\,\|\Phi q_{\epsilon,\sigma}\|_U
 \longrightarrow\frac{\|\Phi\mathbf e_1\|_U}{2\gamma},\quad
 \sqrt d\,\|\Psi q_{\epsilon,\sigma}\|_W
 \longrightarrow\frac{\|\Psi\mathbf e_1\|_W}{2\gamma}.
 \tag{EF15}
\end{equation}
All remaining terms tend to zero by their displayed
powers of $t$ and bounded $r,A_z,B_z$. Fixed linearity
and continuity of the original norms prove both limits.
Equal norm limits therefore do not identify the four
labelled and signed branches or their distinct vector phases.

In the exact lower basis $(1,x,t,tx)$ let
\begin{equation}
 X_d=\begin{pmatrix}0&d^2\\1&0\end{pmatrix},\quad
 H_d=uX_d+vI,\quad D_d=v^2-u^2d^2=-4A_z^2d^2\kappa^2,\quad
 \Theta_L=M_{2t^3}=2\begin{pmatrix}0&H_d^2\\H_d&0\end{pmatrix}.
 \tag{EF16}
\end{equation}
Its exact inverse for $z<1$ is
\begin{equation}
 \Theta_L^{-1}=\tfrac12\begin{pmatrix}0&H_d^{-1}\\H_d^{-2}&0\end{pmatrix},\quad
 H_d^{-1}=\frac{vI-uX_d}{D_d},\quad
 H_d^{-2}=\frac{(v^2+u^2d^2)I-2uvX_d}{D_d^2},\quad
 \det\Theta_L=-1024A_z^6d^6\kappa^6.
 \tag{EF17}
\end{equation}
The identity $X_d^2=d^2I$ proves every product and the
inverse formulas. The determinant follows from
$\det M_t=\det H_d=D_d$ in this dimension, so
$\det M_{2t^3}=2^4D_d^3$. The lower trace Gram is
\begin{equation}
 T_L=B_L\Theta_L=
 \operatorname{diag}\left(
 \begin{pmatrix}4&0\\0&4d^2\end{pmatrix},
 \begin{pmatrix}4v&4ud^2\\4ud^2&4vd^2\end{pmatrix}\right),
 \qquad\det T_L=-1024A_z^2d^6\kappa^2.
 \tag{EF18}
\end{equation}
Multiplication in EF9 gives traces $4,0,0,0$ on
$(1,x,t,tx)$; reducing their ten products proves the
matrix directly and also verifies the residue factorization.
The full eight-dimensional trace is singular at every
parameter because of the upper nilpotents. Only the
specified lower four-dimensional inverse is used here.

The exact limiting inverse coefficient, with no phase
removed, is
\begin{equation}
 K_L=\lim_{z\to1}d^2\Theta_L^{-1}
  =\tfrac12\begin{pmatrix}
       0&(4\gamma^2)^{-1}N_2\\
       (16\gamma^4)^{-1}(I-2iN_2/\gamma)&0
     \end{pmatrix},\qquad
 N_2=\begin{pmatrix}0&0\\1&0\end{pmatrix},\qquad\operatorname{rank}K_L=3.
 \tag{EF19}
\end{equation}
Insert $u\to4\gamma^2$, $k\to2i\gamma$,
$A_z\to-1/2$ into EF17 after multiplying by $d^2$.
The upper block has rank one and the lower block is
invertible; their different parity images prove rank three.
This calculation retains the derivative of the cofactor
in the nonzero $N_2$ term of the lower block.

Here are the actual original norms for these constants.
Let $\Gamma_{r,1}=C_{1/2}^*G_{U,x}C_{1/2}$,
$V_L=C_{-m}P_L$, and
\begin{equation}
 G_L(z)=V_L^*\Gamma_{r,1}V_L,\qquad
 \Gamma_L(z)=\operatorname{diag}(G_L(z),G_L(z)),\qquad
 \Gamma_* =\Gamma_L(1),\quad G_*=G_L(1).
 \tag{EF20}
\end{equation}
The inclusion into the original source is
$J_L(z)=F_U\operatorname{diag}(V_L,V_L)$.
Equations EF6--8 prove it is injective through the
endpoint, so $G_L(z)$ is positive and converges to
the positive $G_*$. This is its induced original norm,
not a replacement metric. The target return uses
$F_W$ and $G_{W,x}$ with the same exact formula.
The equation $(T_A\oplus T_A)J_L(z)=J_{L,W}(z)$
proves the conductor comparison on this moving summand.

Set $L_*=\Gamma_*^{-1}K_L^*\Gamma_*K_L$.
Its three positive eigenvalues are $c_1^2\ge c_2^2\ge c_3^2>0$,
specified without missing metric entries by
\begin{equation}
\begin{aligned}
 a_1&=\operatorname{tr}L_*,\quad
 a_2=\tfrac12(a_1^2-\operatorname{tr}L_*^2),\quad
 a_3=\tfrac16(a_1^3-3a_1\operatorname{tr}L_*^2+2\operatorname{tr}L_*^3),\\
 \{c_1^2,c_2^2,c_3^2\}&=
    \{\text{roots of }s^3-a_1s^2+a_2s-a_3\},\qquad
 \eta=8\gamma^2\sqrt{(G_*)_{11}(G_*^{-1})_{00}}.
\end{aligned}\tag{EF21}
\end{equation}
Similarity with the positive semidefinite Euclidean
square of $K_L$ proves positivity and exactly one zero
eigenvalue. The trace identities are the elementary
symmetric functions, proving the complete cubic.
The limiting forward map is
$8\gamma^2\left(\begin{smallmatrix}0&0\\N_2&0\end{smallmatrix}\right)$.
Its unique positive singular value in $\Gamma_*$ is
$\eta$, by the norm of this rank-one map: its image
vector has squared norm $(G_*)_{11}$ and its input
coordinate functional has squared norm $(G_*^{-1})_{00}$.
Continuity of the positive square roots of $\Gamma_L(z)$,
EF19, and this rank-one forward limit prove the entire
inverse spectrum and every exterior asymptote:
\begin{equation}
\begin{aligned}
 \sigma(\Theta_L^{-1})&\sim
       (c_1d^{-2},c_2d^{-2},c_3d^{-2},\eta^{-1}),\\
 \lim_{z\to1}d^{2\min(j,3)}\|\wedge^j\Theta_L^{-1}\|
       &=\prod_{\nu=1}^{\min(j,3)}c_\nu\,
           \eta^{-\max(j-3,0)}\quad(1\le j\le4),\\
 \|\wedge^4\Theta_L^{-1}\|
       &=\frac1{1024|A_z|^6d^6|\kappa|^6},\qquad
 \frac{c_1c_2c_3}{\eta}=\frac1{65536\gamma^{12}}.
\end{aligned}\tag{EF22}
\end{equation}
The three growing values follow from the rank-three
inverse limit, and the last value is the reciprocal
of the unique positive forward value. Exterior norms
are products of the largest singular values. At full
rank the source and target are the same retained
metric, so the norm is the absolute reciprocal
determinant EF17. Substituting the exact endpoint
$|A_z|=1/2$, $|\kappa|=4\gamma^2$ gives the final
constant and proves agreement with the spectral product.
These are all four remaining ES branches and their
full primary inverse cost in the original receiver.
They do not assert invertibility of the already
singular upper block or identify this auxiliary
family with an arithmetic family of zeta zeros.

The exact constants in EF21 admit a full quadratic
evaluation because the two retained parity metrics
in EF20 are equal. Define
$\rho=(G_*)_{11}(G_*^{-1})_{00}>0$, retaining $G_*$
as the full moment matrix in EF20. Then
\begin{equation}
 \{c_1,c_2,c_3\}=\left\{
 \frac{\sqrt\rho}{8\gamma^2},\quad
 \frac{\sqrt{1+\rho/\gamma^2}+\sqrt\rho/\gamma}{32\gamma^4},\quad
 \frac{\sqrt{1+\rho/\gamma^2}-\sqrt\rho/\gamma}{32\gamma^4}
 \right\},\qquad \eta=8\gamma^2\sqrt\rho.
 \tag{EF23}
\end{equation}
The braces specify a multiset, to be sorted as in EF21.
To prove it, the upper block in EF19 has the single
positive singular value $\sqrt\rho/(8\gamma^2)$.
For $V=I-2iN_2/\gamma$, its determinant is one and
$\operatorname{tr}(G_*^{-1}V^*G_*V)=2+4\rho/\gamma^2$:
the two terms linear in $N_2$ vanish because its trace
and the trace of its adjoint are zero. Thus the two
singular values of $V$ have product one and squared
sum $2+4\rho/\gamma^2$. Solving these two positive
equations gives $\sqrt{1+\rho/\gamma^2}\pm\sqrt\rho/\gamma$.
The unchanged factor $1/(32\gamma^4)$ in EF19 proves
the last two constants. Their product and the first
constant give $c_1c_2c_3=\sqrt\rho/(8192\gamma^{10})$,
which proves the final identity in EF22 again with
every metric factor present.

There is also an exact finite-parameter evaluation,
not only a limiting estimate. For $0\le z<1$, put
\begin{equation}
\begin{aligned}
 a_H&=\operatorname{tr}(G_L(z)^{-1}H_d^*G_L(z)H_d),&
 b_H&=|D_d|^2,\\
 a_{H^2}&=\operatorname{tr}(G_L(z)^{-1}(H_d^2)^*G_L(z)H_d^2),&
 b_{H^2}&=|D_d|^4,\\
 \lambda_{F,\pm}&=\tfrac12(a_F\pm\sqrt{a_F^2-4b_F})
       &&(F=H,H^2),\\
 \sigma(\Theta_L^{-1})&=
 \left\{\frac1{2\sqrt{\lambda_{H,+}}},
 \frac1{2\sqrt{\lambda_{H,-}}},
 \frac1{2\sqrt{\lambda_{H^2,+}}},
 \frac1{2\sqrt{\lambda_{H^2,-}}}\right\}.
\end{aligned}\tag{EF24}
\end{equation}
The positive singular squares of a two-dimensional
map $F$ have sum $a_F$ and product $|\det F|^2$,
as follows by conjugating its metric square to a
positive Hermitian matrix. Its characteristic
polynomial therefore proves the first three lines.
Both $H_d,H_d^2$ are invertible by EF16, so all
four values are positive. The last line follows
from the exact two orthogonal parity blocks of
EF17 with the identical retained $G_L(z)$.
Ordering these four explicitly evaluated values
$v_1\ge v_2\ge v_3\ge v_4$ gives the exact
finite-parameter exterior norm
$\|\wedge^j\Theta_L^{-1}\|=\prod_{\nu=1}^jv_\nu$
for every $1\le j\le4$. Every original center,
cofactor, Gamma mass and inverse moment enters
through the displayed $A_z,k,d,V_L,G_L(z)$.

Finally the complete eight-state residue and trace
are evaluated in the full primary order specified
after EF8. Define
$J_U=\left(\begin{smallmatrix}-2k/(A_z\kappa^2)&1/(A_z\kappa)\\
1/(A_z\kappa)&0\end{smallmatrix}\right)$.
The complete matrices and ranks are
\begin{equation}
\begin{aligned}
 B_{\mathrm{prim}}&=\begin{pmatrix}
 0&\operatorname{diag}(J_U,J_L)\\
 \operatorname{diag}(J_U,J_L)&0\end{pmatrix},
 &\det B_{\mathrm{prim}}&=\frac1{A_z^8\kappa^8},\\
 T_{\mathrm{prim}}&=\operatorname{diag}\left(
 4,0,4,4d^2,0,0,
 \begin{pmatrix}4v&4ud^2\\4ud^2&4vd^2\end{pmatrix}\right),
 &\operatorname{rank}T_{\mathrm{prim}}&=
 \begin{cases}5,&0\le z<1,\\2,&z=1.\end{cases}
\end{aligned}\tag{EF25}
\end{equation}
The two residue cross blocks are exactly EF10--11;
their off-primary products vanish because $e_Ue_L=0$.
Each cross determinant is $-1/(A_z^2\kappa^2)$,
which proves the full determinant. Multiplication
by a scalar plus a nilpotent on the upper length-four
algebra has trace four times its scalar part, since
its matrix in powers of $t$ is triangular with that
scalar on every diagonal entry. Hence the upper
trace is $\operatorname{diag}(4,0,0,0)$. Combining
it with the full lower trace EF18 in the specified
primary order gives the displayed matrix. For $d>0$
the lower trace determinant is nonzero by EF18,
whereas at $d=0$ its rank is exactly one. This proves
both full ranks while the residue remains perfect.
Conjugation and bilinear congruence through EF8 and
the original receivers preserve all the stated ranks.

\section{The ES coefficient boundary with every signed state retained}
\label{sec:es-total-extension}
The starting coefficient equation and rank-three inclusion were proposed
in the received defining-prime account \cite[eqs.\ (3)--(8)]{DefiningPrimeReceived}.
We prove them here and extend them to the entire signed quartic algebra,
its two original conductor receivers, and every inverse exterior norm.
The original four labels and polynomial provenance remain
\cite{ES488,Tao}; FS1--25 and GD14 specify the complete signed algebra
and the unchanged physical return. Integral closure below is a specified
finite morphism of rings, not a rescaling or omission of their coordinates.

\subsection{The exact coefficient inclusion and its conductor ideal}
Put
\[
 R=\mathbb C[A,A^{-1},B,C],\qquad
 \mathscr S=R[D]/(625AD^3-125CD^2+25BC^2D-4C^4).
 \tag{ET1}
\]
Here $A$ is the original quartic coefficient, not the amplification
$a_{\rm amp}$. The retained distinguished prime parameter is $P$.
Define
\[
 \mathscr N=R[P]/\left(P^3+A^{-1}P^2+(B/A)P+\frac{4C}{5A}\right),
 \quad D=-CP/5,\quad W=A^{-1}(B+P+AP^2).
 \tag{ET2}
\]
Both displayed polynomials become monic by multiplying by units.
Consequently $\mathscr S$ and $\mathscr N$ are free over $R$ in the
ordered bases $(1,D,D^2)$ and $(1,P,P^2)$, respectively. Substitution
in ET1 gives zero. The inclusion matrix and its determinant are
\[
 K(C)=\diag(1,-C/5,C^2/25),\qquad \det K=-C^3/125.
 \tag{ET3}
\]
This diagonal map is injective over the domain $R$; after inverting
$C$ it is an isomorphism. Solving ET2 for $B,C,D$ gives
\[
 B=AW-P-AP^2,\quad C=-5APW/4,\quad D=AP^2W/4,
 \qquad \mathscr N=\mathbb C[A,A^{-1},P,W].
 \tag{ET4}
\]
The inverse ring maps follow by substituting these three expressions
in ET2. Thus $\mathscr N$ is a domain and is integrally closed:
a polynomial ring over a field is a unique factorization domain,
localization preserves that property, and a unique factorization
domain is integrally closed by applying prime factorization to the
denominator of an element satisfying a monic equation.
Since the rings have the same fraction field and $\mathscr N$ is
finite integral over $\mathscr S$, it is the full integral closure.
Indeed every fraction integral over $\mathscr S$ is integral over
$\mathscr N$ by the same monic equation, hence belongs to $\mathscr N$.

The quotient and the two expressions of the conductor ideal are
\[
 \mathscr N/\mathscr S=R/(C)\oplus R/(C^2),\qquad
 (\mathscr S:\mathscr N)=C^2\mathscr N=(C,D)^2\mathscr S .
 \tag{ET5}
\]
Here $(\mathscr S:\mathscr N)=\{x\in\mathscr N:x\mathscr N\subset\mathscr S\}$.
For a proof retaining every coefficient, write an element of $\mathscr S$
as $a+CbP+C^2cP^2$; the harmless constant units $-1/5,1/25$
are retained in ET3, and do not change divisibility by $C$.
Set $a_0=A^{-1},b_0=B/A,c_0=4C/(5A)$.
Its product with $P$ has coefficients
$(-C^2cc_0,\ a-C^2cb_0,\ Cb-C^2ca_0)$.
Membership in $\mathscr S$ forces $C\mid a$ and $C\mid b$.
In its product with $P^2$, the $P^2$ coefficient is
$a-Cba_0+C^2c(a_0^2-b_0)$, so $C^2\mid a$.
These conditions say exactly that the original element is in $C^2\mathscr N$.
Conversely $C^2\mathscr N$ is an ideal of $\mathscr N$ contained in
$\mathscr S$. Its $R$ generators $C^2,C^2P,C^2P^2$ are
$C^2,-5CD,25D^2$, proving the last equality.

\subsection{The complete differential and its signed fibres}
The coefficient map $\nu:(A,P,W)\mapsto(A,B,C,D)$ of ET4 has Jacobian
\[
 J_\nu=\begin{pmatrix}
 1&0&0\\
 W-P^2&-1-2AP&A\\
 -5PW/4&-5AW/4&-5AP/4\\
 P^2W/4&APW/2&AP^2/4
 \end{pmatrix}.
 \tag{ET6}
\]
In row orders $ABC,ABD,ACD,BCD$, its four maximal minors are
\[
 \left(
 \frac{5A}{4}(2AP^2+AW+P),\
 -\frac{AP}{4}(2AP^2+2AW+P),\
 \frac{5A^2P^2W}{16},\
 -\frac{5A^2P^4W}{16}
 \right).
 \tag{ET7}
\]
Expansion of each two-by-two minor after the first row gives the first
three; expanding the remaining determinant gives the fourth.
Their common zero locus is exactly
\[
 L_0=\{P=W=0\},\qquad
 L_1=\{W=0,\ P=-1/(2A)\}.
 \tag{ET8}
\]
Indeed the third minor forces $P=0$ or $W=0$.
If $P=0$, the first forces $W=0$; if $W=0$, the first two force
$P(1+2AP)=0$. The first row and the nonzero $W$ derivative of $B$
show that the rank is exactly two at both loci.
In coordinate order $(dA,dP,dW)$ their kernels are, respectively,
$\mathbb C(0,A,1)$ and $\mathbb C(0,1,0)$.

These loci have explicit fibres in the original signed cover
$\mathbb C[r,\sigma]/(h,\sigma^2-h')$, where
$h=Ar^4+r^3+Br^2+Cr+D$. They are
\[
 \begin{array}{ll}
 L_0:&h=r^3(Ar+1),\\[2pt]
 L_1:&h=Ar^2(r+1/(2A))^2 .
 \end{array}
 \tag{ET9}
\]
At $L_0$, reduction in the root-zero primary gives
$\mathbb C[r,\sigma]/(r^3,\sigma^2-3r^2)$, of length six.
The other root is $-1/A$, with $h'=-1/A^2$, hence its two
distinct signs are $\sigma=\pm i/A$. At $L_1$, each primary
has unchanged local coordinate $\epsilon=r-r_0$ and algebra
$\mathbb C[\epsilon,\sigma]/(\epsilon^2,\sigma^2-\epsilon/(2A))$,
of length four, for $r_0=0,-1/(2A)$. These follow by expanding $h'$
and reducing modulo the displayed primary powers, not by discarding
the cofactor. Distinct primary factors are separated by Bezout
polynomials since their root differences are units. Dimensions add
to eight in both cases.

The original coefficient receivers give a further exact differential
square:
\[
 \nu_U=\Phi\nu,\quad \nu_W=\Psi\nu,\quad
 T_A\nu_U=\nu_W,\quad T_A\,d\nu_U=d\nu_W .
 \tag{ET10}
\]
Both $\Phi,\Psi$ are the fixed isomorphisms of FC, so the ranks and
kernels in ET8 remain exactly those of the transported differentials.
This is a calculation of $d\nu$, not a replacement of the nonzero
conductor determinant FC23.

\subsection{A rank-twenty-four comparison retaining the whole cover}
Define over the actual coefficient ring $\mathscr S$
\[
 \mathscr E=\mathscr S[r,\sigma]/
 \left(r^4+A^{-1}r^3+(B/A)r^2+(C/A)r+D/A,\
       \sigma^2-(4Ar^3+3r^2+2Br+C)\right).
 \tag{ET11}
\]
Successive division by the two monic polynomials proves unique
remainders in the full ordered basis
\[
 e=(1,r,r^2,r^3,\sigma,\sigma r,\sigma r^2,\sigma r^3).
 \tag{ET12}
\]
It is therefore free of rank eight over $\mathscr S$.
Set $\mathscr E_{\mathscr N}=\mathscr N\otimes_{\mathscr S}\mathscr E$.
Both algebras are free of rank twenty-four over $R$, with ordered
bases $(e,De,D^2e)$ and $(e,Pe,P^2e)$. Their exact inclusion is
\[
 \mathcal K=K(C)\otimes I_8,\quad
 \det\mathcal K=C^{24}/5^{24},\quad
 \mathscr E_{\mathscr N}/\mathscr E
 =(R/(C))^8\oplus(R/(C^2))^8 .
 \tag{ET13}
\]
This follows coefficient by coefficient in ET12.
No choice of a root or a sign occurs in this map. It base-changes the
entire signed cover along the distinguished-parameter inclusion.

The exact algebra conductor of this extension is
\[
 (\mathscr E:\mathscr E_{\mathscr N})=C^2\mathscr E_{\mathscr N}.
 \tag{ET14}
\]
To prove necessity, write $x=\sum_{j=0}^7 n_j e_j$ with $n_j\in\mathscr N$.
If $x\mathscr E_{\mathscr N}\subset\mathscr E$, multiplication by every
scalar $n\in\mathscr N$ and uniqueness in ET12 give
$n_j\mathscr N\subset\mathscr S$ for each $j$. ET5 gives
$n_j\in C^2\mathscr N$. Conversely every coefficient of
$C^2\mathscr E_{\mathscr N}$ lies in $\mathscr S$, and this is an
ideal of $\mathscr E_{\mathscr N}$, so it is contained in the conductor.
ET14 describes this particular extension; it does not assert that
$\mathscr E_{\mathscr N}$ is the total integral closure of the signed
cover at its own ramification.

At $C=0$, $\mathcal K$ has rank eight and kernel the sixteen-dimensional
span of $(De,D^2e)$. Tensoring the quotient exact sequence with $R/(C)$
gives the full exact sequence
\[
 0\longrightarrow (R/(C))^{16}\longrightarrow
 \mathscr E/C\mathscr E\longrightarrow
 \mathscr E_{\mathscr N}/C\mathscr E_{\mathscr N}
 \longrightarrow(R/(C))^{16}\longrightarrow0 .
 \tag{ET15}
\]
Here the left term is $\operatorname{Tor}_1^R(
 \mathscr E_{\mathscr N}/\mathscr E,R/(C))$.
For each summand $R/(C^j)$, the free resolution
$0\to R\xrightarrow{C^j}R\to R/(C^j)\to0$
computes both its specialization and Tor term as $R/(C)$.
Thus the left injection retains all sixteen killed directions;
they have not been identified with absent root labels.

\subsection{All root and sign actions, and the distinguished-parameter pole}
In the increasing-power root basis define
\[
 X(D)=\begin{pmatrix}
 0&0&0&-D/A\\1&0&0&-C/A\\0&1&0&-B/A\\0&0&1&-1/A
 \end{pmatrix},\quad
 Z(D)=4AX(D)^3+3X(D)^2+2BX(D)+CI_4.
 \tag{ET16}
\]
Every entry of the latter matrix is
\[
 Z(D)=\begin{pmatrix}
 C&-4D&D/A&D(2AB-1)/A^2\\
 2B&-3C&-(4AD-C)/A&(2ABC+AD-C)/A^2\\
 3&-2B&-(3AC-B)/A&-(4A^2D-2AB^2-AC+B)/A^2\\
 4A&-1&-(2AB-1)/A&-(3A^2C-3AB+1)/A^2
 \end{pmatrix}.
 \tag{ET17}
\]
The complete actions on ET12 are
\[
 X_8(D)=\diag(X(D),X(D)),\qquad
 Y_8(D)=\begin{pmatrix}0&Z(D)\\I_4&0\end{pmatrix}.
 \tag{ET18}
\]
They commute, satisfy $h(X_8)=0$ and
$Y_8^2=4AX_8^3+3X_8^2+2BX_8+CI_8$.
These identities follow from monic reduction for $X$ and direct block
multiplication for $Y_8$.

The two scalar matrices for $D$ are
\[
 M_D=\begin{pmatrix}
 0&0&4C^4/(625A)\\1&0&-BC^2/(25A)\\0&1&C/(5A)
 \end{pmatrix},\quad
 M_P^{\mathscr N}=\begin{pmatrix}
 0&0&-4C/(5A)\\1&0&-B/A\\0&1&-1/A
 \end{pmatrix},\quad M_D^{\mathscr N}=-CM_P^{\mathscr N}/5.
 \tag{ET19}
\]
For $Q=X_8,Y_8$, let $Q_0=Q(0),Q_1=\partial_DQ$.
ET16--18 show that $Q=Q_0+DQ_1$ exactly.
Thus all twenty-four dimensional root and sign actions are
\[
 \mathcal Q_{\mathscr S}=I_3\otimes Q_0+M_D\otimes Q_1,\quad
 \mathcal Q_{\mathscr N}=I_3\otimes Q_0+M_D^{\mathscr N}\otimes Q_1,
 \quad \mathcal K\mathcal Q_{\mathscr S}
 =\mathcal Q_{\mathscr N}\mathcal K .
 \tag{ET20}
\]
The last identity follows from $KM_D=M_D^{\mathscr N}K$,
which is checked entrywise using ET3 and ET19.
It is also the scalar multiplication identity in the proved inclusion.

On $C\ne0$, multiplication by the retained $P=-5D/C$ in the
original coefficient basis has the exact matrix and inverse
\[
 M=K^{-1}M_P^{\mathscr N}K
 =\begin{pmatrix}
 0&0&-4C^3/(125A)\\-5/C&0&BC/(5A)\\0&-5/C&-1/A
 \end{pmatrix},
 \qquad
 M^{-1}=\begin{pmatrix}
 -5B/(4C)&-C/5&0\\25/(4C^2)&0&-C/5\\-125A/(4C^3)&0&0
 \end{pmatrix}.
 \tag{ET21}
\]
Multiplication in both orders gives $I_3$.
The characteristic polynomial is the unchanged cubic ET2;
$\det M=-4C/(5A)$ and $\det M^{-1}=-5A/(4C)$.
The full signed-cover action is $\mathcal M=M\otimes I_8$.
It commutes with both $\mathcal X_{\mathscr S}$ and
$\mathcal Y_{\mathscr S}$ in ET20: $M=-5M_D/C$ commutes with
$M_D$, and scalar and root/sign tensor factors commute.
Thus the original signed action, its scalar extension and its inverse
are specified on the same twenty-four dimensional object.

Fix $A\ne0$ and $B$ and let complex $C$ tend to zero.
Retaining its phase, the exact residue and inverse leading term are
\[
 C M\longrightarrow N=-5(E_{10}+E_{21}),\quad
 \rank(N,N^2,N^3)=(2,1,0),\qquad
 C^3M^{-1}\longrightarrow-\frac{125A}{4}E_{20}.
 \tag{ET22}
\]
Indices are zero-based. For $\mathcal M$ these ranks are
$(16,8,0)$ and the leading inverse rank is eight.
The eight chains are $e_j\mapsto-5De_j\mapsto25D^2e_j\mapsto0$.
This is the coefficient lattice residue; it does not change the
characteristic cubic or replace the actual root and sign operators.

\subsection{The full original conductor and metric return}
Let $F_U,F_W$ be exactly GD14, with the original shift $\kappa$,
all FC24 inverse moments and the original Gamma masses.
For a column in three consecutive copies of ET12, define
\[
 \mathfrak F_U=I_3\otimes F_U,\qquad
 \mathfrak F_W=I_3\otimes F_W,\qquad
 \mathfrak T=T_A^{\oplus6},\qquad
 \mathfrak T\mathfrak F_U=\mathfrak F_W .
 \tag{ET23}
\]
Each pair in one copy returns $R(r)+\sigma S(r)$ as
$(T_A^{-1}R(S'-\kappa),T_A^{-1}S(S'-\kappa))$ at the source
and $(R(S'-\kappa),S(S'-\kappa))$ at the target.
These substitutions and FC24 give the inverses of both receivers.
No moments vary with $A,B,C,D$.

For every one of the exact maps $\mathcal Z$ in ET13, ET20 or ET21,
its two physical versions are
\[
 \mathcal Z_U=\mathfrak F_U\mathcal Z\mathfrak F_U^{-1},\quad
 \mathcal Z_W=\mathfrak F_W\mathcal Z\mathfrak F_W^{-1},\quad
 \mathfrak T\mathcal Z_U=\mathcal Z_W\mathfrak T .
 \tag{ET24}
\]
For an inclusion, source and target copies in this formula carry their
specified $\mathscr S$ and $\mathscr N$ bases. Multiplying ET23 proves
the equality in either interpretation. This supplies the exact
morphism to the fixed conductor for the entire cover.

Write $G_{U,r},G_{W,r}$ for the full degree-three coefficient Grams
of the two single-component receivers in GD14, with all original
moments, centers and Gamma masses as evaluated in RS21--26 and MR1--7.
The metrics on this actual direct sum are, respectively,
\[
 \Gamma_U^{(24)}=I_3\otimes\diag(G_{U,r},G_{U,r}),\qquad
 \Gamma_W^{(24)}=I_3\otimes\diag(G_{W,r},G_{W,r}).
 \tag{ET25}
\]
This follows by adding the squared norms of the six actual components.
In particular $G_{U,r}$ and $G_{W,r}$ are not equated.
For either positive inner block $\Gamma$, the metric square of
$M^{-1}\otimes I_8$ is
\[
 (I_3\otimes\Gamma)^{-1}
 (M^{-1}\otimes I_8)^*(I_3\otimes\Gamma)
 (M^{-1}\otimes I_8)
 =((M^{-1})^*M^{-1})\otimes I_8 .
 \tag{ET26}
\]
Every original metric entry cancels by this exact tensor identity;
no metric was replaced by a convenient one. Thus its three Euclidean
singular values occur with multiplicity eight in both physical spaces.

Here is a fully evaluated finite-parameter formula. For $C\ne0$ put
\[
 x=-5B/(4C),\quad y=-C/5,\quad z=25/(4C^2),\quad w=-125A/(4C^3).
\]
Then $(M^{-1})^*M^{-1}$ is
\[
 Q=\begin{pmatrix}
 |x|^2+|z|^2+|w|^2&\bar x y&\bar z y\\
 \bar y x&|y|^2&0\\
 \bar y z&0&|y|^2
 \end{pmatrix}.
 \tag{ET27}
\]
Its positive eigenvalues $\lambda_1\ge\lambda_2\ge\lambda_3>0$
are exactly the three roots, with their multiplicities, of
\[
 \lambda^3-a_1\lambda^2+a_2\lambda-a_3=0,\qquad
 \begin{aligned}
 a_1&=|x|^2+|z|^2+|w|^2+2|y|^2,\\
 a_2&=|y|^2(|x|^2+|z|^2+2|w|^2)+|y|^4,\\
 a_3&=|y|^4|w|^2=25|A|^2/(16|C|^2).
 \end{aligned}
 \tag{ET28}
\]
Indeed expansion of the determinant of $Q-\lambda I$ gives precisely
these coefficients; positivity follows since $M^{-1}$ is invertible.
The roots can all be evaluated by one quadratic. Set
$t=|y|^2$, $r=|x|^2+|z|^2+|w|^2$, $q=|w|^2$.
Expansion, or substitution in ET28, factors the cubic as
$(\lambda-t)(\lambda^2-(r+t)\lambda+tq)$. Hence
\[
 \lambda_1=\frac{r+t+\sqrt{(r+t)^2-4tq}}2,\quad
 \lambda_2=t,\quad
 \lambda_3=\frac{2tq}{r+t+\sqrt{(r+t)^2-4tq}} .
 \tag{ET28a}
\]
The quadratic is positive at zero and negative at $t$, since
$q-r=-|x|^2-|z|^2<0$. Thus its roots satisfy
$\lambda_1>t>\lambda_3>0$; this also proves that its discriminant
is strictly positive. The rationalized last expression retains every
coefficient and avoids cancellation. In particular eight inverse
singular values are exactly $|C|/5$, and eight more merely approach
that scale. No root or metric constant remains unnamed.
Let $s_i=\sqrt{\lambda_i}$. For $1\le j\le24$, set
\[
 n_1=\min(j,8),\quad n_2=\min(\max(j-8,0),8),\quad
 n_3=\max(j-16,0).
\]
The exact original-metric exterior inverse norm is
\[
 \left\|\bigwedge^j\mathcal M_U^{-1}\right\|
 =\left\|\bigwedge^j\mathcal M_W^{-1}\right\|
 =s_1^{n_1}s_2^{n_2}s_3^{n_3}.
 \tag{ET29}
\]
A singular-value decomposition in orthonormal bases proves that
exterior singular values are the corresponding products, and ET26
supplies the multiplicity eight. Thus ET29 covers every exterior degree,
including the full determinant
$\|\bigwedge^{24}\mathcal M^{-1}\|=(5|A|/(4|C|))^8$.

There are two equal positive singular values $5$ in the residue $N$.
Continuity of the metric square applied to $CM$ shows that the two
largest singular values of $M$ are $5/|C|(1+o(1))$.
Using the exact determinant in ET21 gives the third one as
$4|C|^3/(125|A|)(1+o(1))$. Inverting these three values proves
\[
 s_1=\frac{125|A|}{4|C|^3}(1+o(1)),\qquad
 s_2=\frac{|C|}{5}(1+o(1)),\qquad
 s_3=\frac{|C|}{5}(1+o(1)).
 \tag{ET30}
\]
This also follows for the largest inverse value directly from ET22.
For $d_A=125|A|/4$ and all twenty-four exterior degrees, ET29 gives
\[
 \left\|\bigwedge^j\mathcal M^{-1}\right\|\sim
 \begin{cases}
 d_A^j|C|^{-3j},&1\le j\le8,\\
 d_A^8\,5^{-(j-8)}|C|^{j-32},&8\le j\le24.
 \end{cases}
 \tag{ET31}
\]
The formulas agree at $j=8$; at $j=24$ the constant equals
$(5|A|/4)^8$, checking the complete determinant.
They hold for every complex approach with fixed $A\ne0,B$.
Equations ET21--31 establish an actual cubic inverse pole and all its
complementary directions in the stated original receiving norms.
The unchanged linear conductor is still the explicit FC24 isomorphism;
the pole belongs to the displayed distinguished-parameter action.

For comparison, the full coefficient inclusion itself has inverse
singular values $1,5/|C|,25/|C|^2$, each with multiplicity eight in
ET25, sorted according to their actual sizes. This follows from ET3
and the same tensor metric identity. Its determinant inverse norm is
$5^{24}/|C|^{24}$, consistent with ET13. The inclusion pole, the
distinguished-parameter action pole and the fixed conductor have thus
been related by exact maps while retaining their different matrices.

\subsection{The inverse-exterior estimate with the actual conductor included}
Let $\gamma_1\ge\cdots\ge\gamma_8>0$ be the singular values of
$(T_A\oplus T_A)^{-1}$ in its original target-to-source metrics.
They are the four singular values of the exact FC24 inverse
$A_{3,s}^{-1}$, each repeated twice, using precisely its original
$\rho_j^{(s)},h_j,\mu_v$ and centers. In particular
$\prod_{j=1}^8\gamma_j=|\det A_{3,s}|^{-2}$.
Define the actual source-to-target map and its inverse by
\[
 \mathcal B_C=\mathfrak F_W(M\otimes I_8)\mathfrak F_U^{-1}
             =\mathfrak T\,\mathcal M_U,\qquad
 \mathcal B_C^{-1}=\mathfrak F_U(M^{-1}\otimes I_8)
                  \mathfrak F_W^{-1}
             =\mathcal M_U^{-1}\mathfrak T^{-1}.
 \tag{ET32}
\]
Their products are the identity on their full original domains.
The inverse metric square in coefficient coordinates is
\[
 ((M^{-1})^*M^{-1})\otimes
       (\Gamma_W^{-1}\Gamma_U),\qquad
 \Gamma_b=\diag(G_{b,r},G_{b,r}) .
 \tag{ET33}
\]
For a map from metric $I_3\otimes\Gamma_W$ to metric
$I_3\otimes\Gamma_U$, this follows by the definition of the metric
adjoint. The second factor is positive in its $\Gamma_W$ inner product;
its eigenvalues are exactly $\gamma_j^2$ by GD14.
Diagonalizing both positive metric squares proves that the complete
finite-$C$ inverse singular list is
\[
 \{s_i(C)\gamma_j:1\le i\le3,\ 1\le j\le8\}.
 \tag{ET34}
\]
Here the $s_i$ are the explicitly evaluated ET28a square roots.
Sorting these twenty-four positive numbers and multiplying the largest
$k$ gives every exact inverse exterior norm, with all original conductor
weights present, even where the order within the list changes.

Put $\widehat\gamma_{2j-1}=\widehat\gamma_{2j}=\gamma_j$ for
$1\le j\le8$. For $C\to0$, fixed nonzero $A$, fixed $B$ and the
unchanged conductor, all eight large factors eventually exceed
all sixteen small factors, because their scale ratio is $|C|^{-4}$
times a positive fixed constant. Products in ET34 therefore give
\[
 \left\|\bigwedge^k\mathcal B_C^{-1}\right\|\sim
 \begin{cases}
 d_A^k\left(\prod_{j=1}^k\gamma_j\right)|C|^{-3k},
                    &1\le k\le8,\\[3pt]
 d_A^8\,5^{-(k-8)}
 \left(\prod_{j=1}^8\gamma_j\right)
 \left(\prod_{j=1}^{k-8}\widehat\gamma_j\right)|C|^{k-32},
                    &8\le k\le24 .
 \end{cases}
 \tag{ET35}
\]
Repeated $\gamma_j$ do not affect this conclusion: ordered positive
finite lists are continuous and tied limiting products coincide.
In full degree there is the exact, not merely limiting, equality
\[
 \left\|\bigwedge^{24}\mathcal B_C^{-1}\right\|
 =\left(\frac{5|A|}{4|C|}\right)^8|\det A_{3,s}|^{-6},
 \quad
 |\det A_{3,s}|=
 \left|\frac{(-i)^v\mu_v}{v!}\right|^4
 \prod_{j=0}^3\frac{\rho_{v+j}^{(s)}}{\rho_j^{(s)}} .
 \tag{ET36}
\]
This is the actual inverse-exterior comparison with audit equation (11).
The cubic coefficient-lattice loss multiplies the retained conductor
cost; it does not set the nonzero original moment $\mu_v$ to zero.

There is a uniform finite-$C$ estimate with these exact leading scales.
Put
\[
 a_C=\frac{125|A|}{4|C|^3},\quad b_C=\frac{|C|}{5},\quad
 \delta_C=\frac{|C|^2}{25|A|^2}
       +\frac{|B|^2|C|^4}{625|A|^2}
       +\frac{16|C|^8}{390625|A|^2},\quad
 e_k=\min(k,8,24-k).
 \tag{ET37}
\]
The independent derivation ETR21--22 gives the same expression;
here is its full cross-receiver consequence.
Let $L_k$ be the product of the largest $k$ entries of the exact
baseline multiset
$\{a_C\gamma_j,b_C\gamma_j,b_C\gamma_j:1\le j\le8\}$.
Then, for every $A,C\ne0$ and every $B$,
\[
 L_k\le\left\|\bigwedge^k\mathcal B_C^{-1}\right\|
          \le L_k(1+\delta_C)^{e_k/2},
 \qquad 0\le k\le24 .
 \tag{ET38}
\]
To prove this, write $\theta=s_1/a_C$.
The largest eigenvalue of $uu^*+|y|^2\diag(1,1,0)$, where
$u=(x,z,w)^{\mathsf T}$, lies between $|w|^2$ and
$\|u\|^2+|y|^2=|w|^2(1+\delta_C)$.
Thus $1\le\theta\le\sqrt{1+\delta_C}$, and the exact finite
singular triple is
\[
 (s_1,s_2,s_3)=(a_C\theta,b_C,b_C/\theta).
 \tag{ET39}
\]
Its first entry exceeds both $a_C$ and $b_C$ by ET28a; its third
is smaller than both. Its total product is $a_Cb_C^2$.
Consequently, in each fixed $j$ block, the product of the largest
zero, one, two or three actual entries is at least the corresponding
baseline product. For one entry this uses $s_1\ge\max(a_C,b_C)$.
For two entries, $s_1b_C$ is at least $a_Cb_C$ and $b_C^2$.
For zero or three entries there is equality.
Take a subset attaining the largest baseline $k$-product, count its
entries in each block, and replace them by that many largest actual
entries in the same block. This proves the lower bound.
For the upper bound, each actual entry is its specified baseline
entry times $\theta$, $1$, or $\theta^{-1}$.
Across the eight blocks the exponent list has eight $+1$'s,
eight zeros and eight $-1$'s. A $k$-element subset has exponent
at most $\min(k,8,24-k)$, the last bound following by taking its
complement. Its baseline product is at most $L_k$.
Maximizing over all subsets proves the upper bound in ET38.
This proof holds even when baseline values change order and when
original conductor singular values coincide. At $k=24$ the two
bounds agree with the exact determinant ET36.

% Independent derivation of the complete coefficient-normalization extension.
% This file is an includable mathematical body; it performs no publication.
\section{Independent review: the complete ES coefficient extension}

The received derivation, \emph{Defining-prime continuation}, equations
(3)--(8), is the immediate source of the coefficient construction below.
The exact original receiving maps are those of the cumulative manuscript,
RD18--22 and GD14--17. All assertions made in this review are proved here.
In particular, a pullback to the normalization of the coefficient base is
not asserted to be the integral closure of the entire signed cover.

Put
\[
 R=\mathbb C[A,A^{-1},B,C],\qquad
 S=R[D]/(D^3-\tfrac{C}{5A}D^2+\tfrac{BC^2}{25A}D
                         -\tfrac{4C^4}{625A}),
 \quad
 N=R[P]/(P^3+A^{-1}P^2+(B/A)P+4C/(5A)).                 \tag{ETR1}
\]
Every occurrence of $A^{-1}$ has the literal domain $A\ne0$.
No restriction on $B$ or on $C$ is made until an inverse using $C^{-1}$
is explicitly introduced. Both displayed polynomials are monic; division
by them proves that the respective ordered bases over $R$ are
$(1,D,D^2)$ and $(1,P,P^2)$. Solving the relation for $C$ gives
\[
 C=-\tfrac54(AP^3+P^2+BP),\quad
 W=A^{-1}(B+P+AP^2),\quad
 B=AW-P-AP^2,
 \quad C=-\tfrac54APW,\quad D=\tfrac14AP^2W=-CP/5.       \tag{ETR2}
\]
Thus $N=\mathbb C[A,A^{-1},P,W]$. These formulas and their displayed
inverses prove the isomorphism, rather than only a parameter count.
The relation defining $S$ vanishes after substituting $D=-CP/5$.
The resulting inclusion has the complete matrix
\[
 K=\operatorname{diag}(1,-C/5,C^2/25),\qquad
 S\lhook\joinrel\longrightarrow N.                    \tag{ETR3}
\]
Since $R$ is a domain, this matrix is injective. It becomes invertible
over $R[C^{-1}]$, so the two domains have the same fraction field.
$N$ is finite over $S$: its three $R$-generators also generate it over
$S$, because $R\subset S$. A Laurent polynomial ring is integrally
closed: polynomial factorization over a field is unique, localization
preserves unique factorization, and an element $a/b$ with relatively
prime numerator and denominator satisfying a monic polynomial would
force every prime factor of $b$ to divide $a$. Hence $N$ is exactly the
normalization of $S$. Indeed, it is integral over $S$, and every element
of their common fraction field integral over $S$ is integral over $N$
by the same monic equation and consequently belongs to $N$.

The cokernel of ETR3, with its entire $R$-module structure, is
\[
 N/S=R/(C)\oplus R/(C^2).                                \tag{ETR4}
\]
Here and below the nonzero constants $5$ and $25$ are retained in the
matrices even though they are units for this cokernel calculation.
To calculate the algebra conductor, write an arbitrary element of $S$
uniquely as $a+CbP+C^2cP^2$, with $a,b,c\in R$. Write
$\alpha=A^{-1}$, $\beta=B/A$, $\gamma=4/(5A)$, so
$P^3=-\alpha P^2-\beta P-\gamma C$. The $P$ and $P^2$
coefficients of its product with $P$ are
$a-\beta C^2c$ and $Cb-\alpha C^2c$. Membership in $S$ forces
$a\in CR$ and $b\in CR$. The $P^2$ coefficient of its product with
$P^2$ is $a-\alpha Cb+(\alpha^2-\beta)C^2c$, which then forces
$a\in C^2R$. Conversely, $C^2N\subset S$ and is an ideal of $N$.
Since $1,P,P^2$ generate $N$, these tests prove the exact conductor
identities
\[
 (S:N)=C^2N=(C,D)^2\subset S.                            \tag{ETR5}
\]
The last equality follows from the three generators
$C^2,CD=-C^2P/5,D^2=C^2P^2/25$. It holds without assuming that
$B$ is invertible. At $C=D=0$ the distinguished parameter is retained
by $P(AP^2+P+B)=0$; assigning a value to $-5D/C$ would not define this
fiber.

\subsection{All twenty-four coordinates and the algebra conductor}
Keep the literal polynomial
\[
 h(r)=Ar^4+r^3+Br^2+Cr+D,\quad
 E=S[r,\sigma]/(h(r)/A,\ \sigma^2-h'(r)),\quad
 E_N=N\otimes_S E.                                     \tag{ETR6}
\]
Division first by the monic polynomial $h/A$ and then by the monic
polynomial in $\sigma$ proves that
$b=(1,r,r^2,r^3,\sigma,\sigma r,\sigma r^2,\sigma r^3)$ is an
$S$-basis of $E$ and an $N$-basis of $E_N$. Use the $R$-bases
$(b,Db,D^2b)$ and $(b,Pb,P^2b)$, retaining this order. Then
\[
 \iota=K\otimes I_8,
 \quad E_N/E=(R/(C))^8\oplus(R/(C^2))^8,
 \quad (E:E_N)=C^2E_N=(C,D)^2E.                         \tag{ETR7}
\]
For the conductor assertion, write $z=\sum_{i=1}^8z_i b_i\in E$.
The conditions $zP,zP^2\in E$ are exactly the eight separate conditions
$z_iP,z_iP^2\in S$, because $b$ is an $N$-basis of $E_N$. ETR5 makes
each $z_i$ belong to $C^2N$. Conversely, $E_N=E+EP+EP^2$, so those
conditions suffice for $zE_N\subset E$. This proves the conductor,
including its necessity, and does not invoke a claim that normalization
commutes with an arbitrary base change. The two algebras have rank
twenty-four over $R$, and the signed fiber still has all eight states
over each coefficient point. No pair of root labels or factor signs
has been identified.

\subsection{The distinguished action and its exact signed-root commutants}
In the indicated three-dimensional coefficient bases let
\[
 D_S=\begin{pmatrix}0&0&4C^4/(625A)\\1&0&-BC^2/(25A)\\0&1&C/(5A)\end{pmatrix},
 \quad
 P_N=\begin{pmatrix}0&0&-4C/(5A)\\1&0&-B/A\\0&1&-1/A\end{pmatrix}.
                                                               \tag{ETR8}
\]
On $C\ne0$ multiplication by $P=-5D/C$ on $E$ is
\[
 \mathcal M=M\otimes I_8,\quad
 M=\begin{pmatrix}0&0&-4C^3/(125A)\\-5/C&0&BC/(5A)\\0&-5/C&-1/A\end{pmatrix},
 \quad
 L=M^{-1}=\begin{pmatrix}-5B/(4C)&-C/5&0\\25/(4C^2)&0&-C/5\\-125A/(4C^3)&0&0\end{pmatrix}.
                                                               \tag{ETR9}
\]
Direct multiplication gives $ML=LM=I_3$ and
$\det M=-4C/(5A)$, $\det L=-5A/(4C)$. Also
$P_NK=KM$ and $D_S=-CM/5$. Thus $(P_N\otimes I_8)\iota
=\iota\mathcal M$ on the entire stated domain; the left-hand side
extends regularly to $C=0$.

Here are all operators needed to check that this distinguished action
retains the signed-root algebra. In zero-based indices put
\[
 J=\begin{pmatrix}0&0&0&0\\1&0&0&-C/A\\0&1&0&-B/A\\0&0&1&-1/A\end{pmatrix},
 \quad
 \mathcal R_S=I_3\otimes I_2\otimes J
       -A^{-1}D_S\otimes I_2\otimes E_{03},
 \quad
 \mathcal H_S=4A\mathcal R_S^3+3\mathcal R_S^2+2B\mathcal R_S+CI_{24},
 \quad
 \mathcal S_S=I_3\otimes E_{10}\otimes I_4
          +(I_3\otimes E_{01}\otimes I_4)\mathcal H_S .    \tag{ETR10}
\]
The first matrix is multiplication by $r$: its last column is exactly
the relation $r^4=-(r^3+Br^2+Cr+D)/A$. The second is multiplication
by $h'(r)$, and the last is multiplication by $\sigma$ in the ordered
even and odd parts. In particular,
\[
 A\mathcal R_S^4+\mathcal R_S^3+B\mathcal R_S^2+C\mathcal R_S
                      +D_S\otimes I_8=0,\quad
 \mathcal S_S^2=\mathcal H_S,\quad
 [\mathcal R_S,\mathcal S_S]=[\mathcal M,\mathcal R_S]
                         =[\mathcal M,\mathcal S_S]=0.   \tag{ETR11}
\]
The first identity follows column by column from the monic relation.
For the second, $\mathcal H_S$ is the same operator on the even and
odd copies, so multiplying the two by two matrix
$\left(\begin{smallmatrix}0&\mathcal H_S\\I&0\end{smallmatrix}\right)$
gives the claim in either copy. The commutators follow because every
coefficient in $\mathcal R_S$ is a polynomial in $D_S$, whereas
$M=-5D_S/C$; the even-odd matrices commute with both. Replace $D_S$
in ETR10 by $D_N=-CP_N/5$ to obtain $\mathcal R_N,\mathcal S_N$.
The equality $KD_S=D_NK$ proves, with no omitted correction terms,
\[
 \iota\mathcal R_S=\mathcal R_N\iota,\qquad
 \iota\mathcal S_S=\mathcal S_N\iota.                    \tag{ETR12}
\]
These equations are the exact relation with the root and signed-factor
operators; no assertion identifies these new operators with the
original conductor's multiplication by its spectral variable.

The pole residue is particularly explicit:
\[
 R_P=\left.(C\mathcal M)\right|_{C=0}
       =-5(E_{10}+E_{21})\otimes I_8,
 \quad R_P^2=25E_{20}\otimes I_8,
 \quad R_P^3=0,
 \quad (\operatorname{rank}R_P,\operatorname{rank}R_P^2,
             \operatorname{rank}R_P^3)=(16,8,0).         \tag{ETR13}
\]
These ranks follow from the independent images of the three coefficient
blocks: block zero maps onto block one, block one onto block two,
and block two to zero. There are eight Jordan blocks of size three.
This residue is the residue of the stated meromorphic action in the
coefficient lattice; the original conductor has not acquired this
nilpotent residue.

\subsection{Return to both fixed original metrics}
Let $\mathcal U,\mathcal W,T_A,\kappa$ have precisely their meanings
and original moments in RD18--22 and GD14. For a coefficient column
of $R(r)+\sigma S(r)$ define
\[
 F_U(R,S)=(T_A^{-1}R(S'-\kappa),T_A^{-1}S(S'-\kappa)),
 \quad F_W(R,S)=(R(S'-\kappa),S(S'-\kappa)),
 \quad G_8=T_A\oplus T_A,\quad F_W=G_8F_U.              \tag{ETR14}
\]
Their inverses first apply $T_A$ componentwise when necessary and
then substitute $S'=r+\kappa$. Thus the return is an actual linear
isomorphism on all eight coordinates. Let $\Gamma_U=F_U^*H_UF_U$
and $\Gamma_W=F_W^*H_WF_W$ denote the exact positive Gram matrices
from these original norms; neither is set equal to an identity.
On the full extension use
\[
 \mathcal F_U=I_3\otimes F_U,\quad
 \mathcal F_W=I_3\otimes F_W,\quad
 \Gamma_{24,U}=I_3\otimes\Gamma_U,\quad
 \Gamma_{24,W}=I_3\otimes\Gamma_W,\quad
 G_{24}=I_3\otimes G_8.                                 \tag{ETR15}
\]
The original center, Gamma masses and all inverse moments are retained
by these formulas; there are three copies of the exact eight-state
receiver. Define $\mathcal M_U=\mathcal F_U\mathcal M\mathcal F_U^{-1}$
and $\mathcal M_W=\mathcal F_W\mathcal M\mathcal F_W^{-1}$, with the same
definition for the root, sign and inclusion matrices. Substitution gives
\[
 G_{24}\mathcal M_U=\mathcal M_WG_{24},\quad
 G_{24}\mathcal R_U=\mathcal R_WG_{24},\quad
 G_{24}\mathcal S_U=\mathcal S_WG_{24}.                   \tag{ETR16}
\]
This gives the exact original conductor return and every commutant,
including the nilpotent residue. The norm on the coefficient block of
length three is the specified orthogonal direct sum, not a new metric
on any of the original eight-state spaces.

\subsection{Finite singular values and every exterior inverse norm}
For $A,C\ne0$ and arbitrary complex $B$, put
\[
 u=\begin{pmatrix}-5B/(4C)\\25/(4C^2)\\-125A/(4C^3)\end{pmatrix},\quad
 t=|C|^2/25,\quad r_u=u^*u,\quad q=|u_3|^2
      ={15625|A|^2\over16|C|^6}.
                                                               \tag{ETR17}
\]
Here $u_3$ denotes the third entry of the displayed column. Since
$L=(u,-Ce_1/5,-Ce_2/5)$, one has
$LL^*=uu^*+t\operatorname{diag}(1,1,0)$. The nonzero vector in
$\operatorname{span}(e_1,e_2)$ orthogonal to $(u_1,u_2,0)$ is an
eigenvector of eigenvalue $t$. On its orthogonal complement the
trace is $r_u+t$ and the determinant is $tq$: in the orthonormal basis
consisting of $(u_1,u_2,0)/\sqrt{|u_1|^2+|u_2|^2}$ and $e_3$, this
follows by multiplying the full two by two Hermitian matrix. Thus
\[
 \lambda_\pm={r_u+t\pm\sqrt{(r_u+t)^2-4tq}\over2},\quad
 \lambda_-={2tq\over r_u+t+\sqrt{(r_u+t)^2-4tq}},\quad
 (\alpha,\beta,\chi)=(\sqrt{\lambda_+},\sqrt t,\sqrt{\lambda_-}).
                                                               \tag{ETR18}
\]
The three symbols on the right are singular values, not ES root labels.
The characteristic polynomial evaluated at $t$ is
$-t(|u_1|^2+|u_2|^2)<0$, because $u_2\ne0$. Its positive roots therefore
satisfy $\lambda_+>t>\lambda_->0$, proving their order and the strict
positivity of the displayed square root. These are exact finite formulas.

For any positive original Gram $\Gamma$, conjugation by
$I_3\otimes\Gamma^{1/2}$ sends $L\otimes I_8$ to itself. This is an
isometry from the literal Gram $I_3\otimes\Gamma$ to the Euclidean
calculation and proves that its twenty-four singular values in either
original receiver are
\[
 \underbrace{\alpha,\ldots,\alpha}_{8},\quad
 \underbrace{\beta,\ldots,\beta}_{8},\quad
 \underbrace{\chi,\ldots,\chi}_{8}.                     \tag{ETR19}
\]
The cancellation here is a proved tensor identity, not an assumption
that the original Gram is the identity. In particular only the middle
eight values are exactly $|C|/5$. The last eight are strictly smaller
for every $A,C\ne0$.

For $0\le j\le24$, with an exterior power of degree zero assigned
norm one, define $a_j=\min(j,8)$, $b_j=\min(\max(j-8,0),8)$ and
$c_j=\max(j-16,0)$. The complete exterior inverse answer is
\[
 \|\mathop{\bigwedge}\nolimits^j\mathcal M^{-1}\|
       =\alpha^{a_j}\beta^{b_j}\chi^{c_j}.               \tag{ETR20}
\]
To prove this, apply an orthonormal singular basis in both the domain
and codomain; its exterior basis is again orthonormal, and the exterior
matrix is diagonal with every product of $j$ distinct singular values.
Its largest entry is exactly the ordered product in ETR20.

There is also a finite explicit two-sided bound with the exact leading
constants. Set
\[
 k={125|A|\over4},\quad
 \delta={|C|^2\over25|A|^2}
       +{|B|^2|C|^4\over625|A|^2}
       +{16|C|^8\over390625|A|^2},\quad
 Q_j=(k|C|^{-3})^{\min(j,8)}(|C|/5)^{\max(j-8,0)},\quad
 e_j=\min(j,8,24-j).
                                                               \tag{ETR21}
\]
Because the largest eigenvalue of $uu^*+t\operatorname{diag}(1,1,0)$
is at least $r_u\ge q$ and at most $r_u+t=q(1+\delta)$, and
$\lambda_+\lambda_-=tq$, one has, on the entire finite domain,
\[
 Q_j\ \le\ \|\mathop{\bigwedge}\nolimits^j\mathcal M^{-1}\|
       \ \le\ Q_j(1+\delta)^{e_j/2}.                    \tag{ETR22}
\]
Indeed, writing $\theta=\alpha/\sqrt q\in[1,\sqrt{1+\delta}]$ gives
$\chi=\beta/\theta$ and the exact identity
$\|\bigwedge^j\mathcal M^{-1}\|=Q_j\theta^{a_j-c_j}$;
$a_j-c_j=e_j$ in all three ranges. This proves ETR22 without discarding
any direction. At fixed $A\ne0$ and fixed $B$, $\delta\to0$ as $C\to0$,
so the complete leading terms are
\[
 \|\mathop{\bigwedge}\nolimits^j\mathcal M^{-1}\|
 \sim
 \begin{cases}
 k^j|C|^{-3j},&0\le j\le8,\\
 k^8 5^{-(j-8)}|C|^{j-32},&8\le j\le24.
 \end{cases}\qquad
 \|\mathop{\bigwedge}\nolimits^{24}\mathcal M^{-1}\|
       =\left|{5A\over4C}\right|^8.                     \tag{ETR23}
\]
If $A$ varies, the limit requires $A$ to have a nonzero limiting value;
the exact ETR18--22 remain valid without this limiting hypothesis.

Finally retain the original conductor in a cross-receiver calculation.
Let $g_1,\ldots,g_8$ be the singular values of the literal
$G_8=T_A\oplus T_A$ with its given original source and target norms.
Each of the four singular values of $T_A$ occurs twice. In coefficient
coordinates these are the singular values of
$\Gamma_W^{1/2}\Gamma_U^{-1/2}$, because both coefficient columns
describe the same transported vector under ETR14. The full operator
$G_{24}\mathcal M_U^{-1}$, from the source receiver to the target
receiver, consequently has exactly the multiset
\[
 \{\alpha g_i,\ \beta g_i,\ \chi g_i:1\le i\le8\}.    \tag{ETR24}
\]
An orthonormal singular basis for each tensor factor proves this by
direct multiplication. Every cross-receiver exterior norm is the product
of the largest $j$ members of this displayed, entirely retained
multiset. Thus ETR24 retains the original moments and Gamma factors in
the cross-receiver estimate. It does not replace them by one. The fixed
conductor $G_{24}$ itself remains invertible; the pole and the conductor
ideal in ETR7 and ETR13 belong to the explicitly introduced coefficient
extension and action. No order jump of the original $T_A$, and no
statement settling RH, follows from these calculations alone.
\subsection{The inverse of the full source-to-target comparison}
For the actual source-to-target operator, retain its order:
\[
 \mathcal B_C=\mathcal F_W(M\otimes I_8)\mathcal F_U^{-1}
              =G_{24}\mathcal M_U,\qquad
 \mathcal B_C^{-1}=\mathcal F_U(L\otimes I_8)\mathcal F_W^{-1}
              =\mathcal M_U^{-1}G_{24}^{-1}.             \tag{ETR25}
\]
Both products are the identity by ETR9 and the exact receiver inverses.
Let $\gamma_1\ge\cdots\ge\gamma_8>0$ be the singular values of the
original $G_8^{-1}$, from its target metric to its source metric. They
are the original four singular values of $T_A^{-1}$ each repeated
twice. Applying the definition of the metric adjoint to ETR25 gives
\[
 (L^*L)\otimes(\Gamma_W^{-1}\Gamma_U),\qquad
 \operatorname{Sing}(\mathcal B_C^{-1})
     =\{\alpha\gamma_j,\beta\gamma_j,\chi\gamma_j:1\le j\le8\}.
                                                               \tag{ETR26}
\]
The second factor is self-adjoint and positive for the $\Gamma_W$
inner product. It is similar to the positive Hermitian matrix
$\Gamma_W^{-1/2}\Gamma_U\Gamma_W^{-1/2}$ and hence has a full
orthonormal singular basis for that metric. Its eigenvalues are exactly
$\gamma_j^2$, as follows by pulling back the unchanged map $G_8^{-1}$
through ETR14. Tensoring the orthonormal eigenbases proves ETR26,
including all multiplicities. Every finite exterior inverse norm is
therefore the product of the largest $j$ members of this explicit list.

Define the decreasing sixteen-element list
$\widehat\gamma_{2i-1}=\widehat\gamma_{2i}=\gamma_i$.
At fixed $A\ne0$, fixed $B$ and fixed original conductor, the smallest
of the eight large values divided by the largest of the sixteen small
ones is at least $\alpha\gamma_8/(\beta\gamma_1)$, which tends to
infinity. The large and small groups thus separate for all sufficiently
small $|C|>0$. Within the large group their order is exactly that of
the $\gamma_i$. Dividing the small group by $|C|/5$ gives a list
converging entry by entry to the duplicated list $\widehat\gamma$,
because $\chi/\beta\to1$. Sorting is continuous for finite real lists:
if every entry changes by at most $\varepsilon$, every order statistic
changes by at most $\varepsilon$, since the counts above any threshold
can change only in its $\varepsilon$ neighborhood. This argument also
covers tied values. Consequently, with $k=125|A|/4$,
\[
 \|\mathop{\bigwedge}\nolimits^j\mathcal B_C^{-1}\|
 \sim\begin{cases}
 k^j(\prod_{i=1}^{j}\gamma_i)|C|^{-3j},&0\le j\le8,\\
 k^8 5^{-(j-8)}(\prod_{i=1}^{8}\gamma_i)
      (\prod_{i=1}^{j-8}\widehat\gamma_i)|C|^{j-32},&8\le j\le24.
 \end{cases}                                                \tag{ETR27}
\]
Empty products are one. In degree twenty-four, determinant
multiplicativity gives the exact finite equality
\[
 \|\mathop{\bigwedge}\nolimits^{24}\mathcal B_C^{-1}\|
   =\left|\frac{5A}{4C}\right|^8 |\det A_{3,s}|^{-6},\qquad
 |\det A_{3,s}|=
 \left|\frac{(-i)^v\mu_v}{v!}\right|^4
       \prod_{i=0}^{3}\frac{\rho_{v+i}^{(s)}}{\rho_i^{(s)}}.
                                                               \tag{ETR28}
\]
The last formula is the full original FC23 determinant, which is
nonzero and has all its original weights. The exponent six is forced
by the six original four-dimensional conductor components, equivalently
by $(\prod_i\gamma_i)^3=|\det A_{3,s}|^{-6}$. Thus the total exterior
estimate preserves the entire original conductor and every one of the
twenty-four extension directions.

\subsection{Independent check of the coefficient differential and its fibers}
For completeness the full differential of ETR2, with target order
$(A,B,C,D)$ and source order $(A,P,W)$, is
\[
 J_\nu=\begin{pmatrix}
 1&0&0\\ W-P^2&-1-2AP&A\\-5PW/4&-5AW/4&-5AP/4\\
 P^2W/4&APW/2&AP^2/4
 \end{pmatrix}.
                                                               \tag{ETR29}
\]
Its maximal minors in row orders $ABC,ABD,ACD,BCD$ are respectively
$5A(2AP^2+AW+P)/4$, $-AP(2AP^2+2AW+P)/4$,
$5A^2P^2W/16$, and $-5A^2P^4W/16$, by direct expansion.
Their common zero locus is the union $L_0=\{P=W=0\}$ and
$L_1=\{W=0,P=-1/(2A)\}$: the third minor first forces $P=0$ or
$W=0$, and the first two then give exactly these alternatives. The
first row and the nonzero $A$ entry in the second row show rank two,
and direct substitution gives kernel vectors $(0,A,1)$ and $(0,1,0)$.
The full signed fibers are
\[
 L_0:\quad h=r^3(Ar+1),\quad
 E_0=\mathbb C[r,\sigma]/(r^3,\sigma^2-3r^2)
       \ \times\ \mathbb C_{\sigma=i/A}\ \times\ \mathbb C_{\sigma=-i/A};
 \qquad
 L_1:\quad h=Ar^2(r+1/(2A))^2,\quad
 E_1=\prod_{r_0=0,-1/(2A)}
 \mathbb C[\epsilon,\sigma]/(\epsilon^2,\sigma^2-\epsilon/(2A)),
 \quad\epsilon=r-r_0.                                      \tag{ETR30}
\]
For $L_0$, reducing $h'$ modulo $r^3$ gives $3r^2$; evaluating at
$-1/A$ gives $-1/A^2$, with exactly the two stated signs. For $L_1$,
expansion at either original root and reduction modulo $\epsilon^2$
gives $\epsilon/(2A)$. The factors are pairwise coprime because
$A\ne0$, so the ordinary Bezout identity gives the indicated direct
products. Their dimensions are respectively $6+1+1$ and $4+4$.
No signed state is deleted by either differential rank loss.

\subsection{A uniform finite bound for the full cross-receiver inverse}
Keep the complete original positive list $\gamma_1,\ldots,\gamma_8$
from ETR26. Write $a=125|A|/(4|C|^3)$, $b=|C|/5$, and use the
unchanged $\delta$ of ETR21. Let $L_j$ be the product of the largest
$j$ entries of the twenty-four-element multiset
$\{a\gamma_i,b\gamma_i,b\gamma_i:1\le i\le8\}$, with $L_0=1$.
Then the following bounds hold for every $A,C\ne0$, every complex
$B$, and every exterior degree, without assuming any ordering of
$a$ and $b$:
\[
 L_j\ \le\ \|\mathop{\bigwedge}\nolimits^j\mathcal B_C^{-1}\|
       \ \le\ L_j(1+\delta)^{\min(j,8,24-j)/2},
       \qquad 0\le j\le24.                              \tag{ETR31}
\]
For a full proof put $\theta=\alpha/a$. ETR18--22 give
$1\le\theta\le\sqrt{1+\delta}$ and the exact actual triple
$(\alpha,\beta,\chi)=(a\theta,b,b/\theta)$.
The finite order proved in ETR18 gives $\alpha>b$; the Rayleigh
quotient used in ETR21 gives $\alpha\ge a$. Thus $\alpha$ is at
least both $a$ and $b$. In one fixed $i$ block, the largest one-entry
actual product is $\alpha\gamma_i$, at least
$\max(a,b)\gamma_i$, which is the largest baseline one-entry
product. The largest two-entry actual product is $\alpha b\gamma_i^2$;
its difference from each candidate baseline two-entry product is
$b(\alpha-a)\gamma_i^2$ or $b(\alpha-b)\gamma_i^2$, both nonnegative.
For zero entries both products are one; for three entries both are
$ab^2\gamma_i^3$, since the $\theta$ and $\theta^{-1}$ factors cancel.
Choose a subset of baseline entries attaining $L_j$ and retain its
number of entries in each of the eight blocks. Replacing those entries
in each block by that many largest actual entries can only increase
the product, by the four cases just proved. The resulting actual
$j$-element product is bounded above by its largest such product,
$\|\bigwedge^j\mathcal B_C^{-1}\|$. This proves the lower bound.

For the upper bound, assign to each actual entry the identically
indexed baseline entry. Their ratios are respectively $\theta$, $1$
and $\theta^{-1}$ in every block. For a selected $j$-element subset
let $n_+,n_0,n_-$ count these three types. Its exponent is
$e=n_+-n_-$. The inequalities $e\le j$ and $e\le8$ follow from
$n_++n_0+n_-=j$ and $n_+\le8$. The full list has total exponent zero;
its complement contains $24-j$ entries, each with exponent at least
$-1$, so $e$ is at most $24-j$. Hence
$e\le\min(j,8,24-j)$ for every subset. Since $\theta\ge1$, its
actual product is at most its baseline product times
$\theta^{\min(j,8,24-j)}$, and the baseline product is at most
$L_j$. Maximizing over the subsets proves the upper bound using
$\theta\le\sqrt{1+\delta}$. This proof retains coincident conductor
singular values and any changes in their order. At $j=24$ the exponent
is zero and both bounds equal the exact determinant ETR28.


\section{Defining-prime and collision continuation: domains and strengthened conclusions}
The following complete proofs extend the coefficient result ET1--39.
PS and ISR calculate the exact arithmetic comparison, including the
fractional chart when $p\mid S$. PB and PES then use positivity:
the one-divisible gap satisfies $n\le1$, and every actual two-divisible
witness has $n=1$. The higher gap formulas in PS16 and ISR16--17 are
retained as general algebraic calculations; PB1--8 proves that they
create no exception to the coefficient-ideal inclusion for an actual
positive witness on the original $p\nmid S$ chart.
ECC retains the complete original signed collision and both original
metrics. OCQ computes the actual original conductor on their quotient:
its four inverse singular values have finite limits, and each inverse
exterior norm is bounded by that of the unchanged original conductor.
These domains and conclusions strengthen the earlier comparison without
asserting that a local algebra conductor is the original linear operator.

\section{The defining-prime order and the exact coefficient specialization}
The received defining-prime account \cite[eqs.\ (9)--(31)]{DefiningPrimeReceived}
provides two integral ES witnesses and proposes the signed-order comparison.
The independent derivation below proves its arithmetic assertions with all
units retained. Here we give the full residue/conductor matrix and the exact
map from the coefficient extension ET1--39 to the actual witness. The
original finite inverse FS13 and residue RD1--4 retain their human
polynomial and residue provenance \cite{ES488,Tao,JouveRV}.

\subsection{All eight integral conductor generators}
Let $p\equiv1\pmod {12}$ be prime, and let positive integers $x,y,z$
satisfy $4/p=1/x+1/y+1/z$. Set
\[
 (r_0,r_1,r_2,r_3)=(p,x,y,z),\quad
 S=p+x+y+z,\quad A=-S^{-1},\quad
 f(r)=\prod_{i=0}^3(r-r_i),\quad h=Af .
 \tag{PS1}
\]
The four roots are distinct by the elementary ES argument in the
independent integral derivation. For now assume $p\nmid S$, so $A$
is a unit in $R_p=\mathbb Z_p$, with fraction field $K_p=\mathbb Q_p$.
Put
\[
 \begin{aligned}
 d_i&=Af'(r_i),&
 b_i&=v_p(d_i)=\sum_{j\ne i}v_p(r_i-r_j),\\
 m_i&=\lfloor b_i/2\rfloor,&
 \epsilon_i&=b_i-2m_i,&u_i&=p^{-b_i}d_i .
 \end{aligned}
 \tag{PS2}
\]
The symbols $d_i$ retain their full units; $d_i$ is not replaced by $p^{b_i}$.
The original integral algebra and its integral closure in its generic
fibre are
\[
 E=R_p[r,\sigma]/(f,\sigma^2-h'),\qquad
 O=\prod_{i=0}^3R_p[\tau_i],\qquad
 \tau_i^2=p^{\epsilon_i}u_i,\quad \sigma_i=p^{m_i}\tau_i .
 \tag{PS3}
\]
Each factor can be a quadratic field order or a split quadratic algebra.
The independent proof verifies its full integral closure directly, including
the split case. With $V=(r_i^j)_{0\le i,j\le3}$, the inclusion is
\[
 C_p=\diag(V,\diag(p^{m_i})V)
 \tag{PS4}
\]
in the unchanged basis $(1,r,r^2,r^3,\sigma,\sigma r,\sigma r^2,\sigma r^3)$
and the normalization basis $(1_0,\ldots,1_3,\tau_0,\ldots,\tau_3)$.
Thus every label and both geometric signs are retained.

Define $f_i=f/(r-r_i)$ and let $Q$ be the four-by-four matrix with
$i$th column the increasing-power coefficients of $f_i$:
\[
 Q_i=\begin{pmatrix}
 -\prod_{j\ne i}r_j\\
 \sum_{\substack{j<k\\j,k\ne i}}r_jr_k\\
 -\sum_{j\ne i}r_j\\1
 \end{pmatrix},\quad
 VQ=\diag(f'(r_i)),\quad\det Q=\det V .
 \tag{PS5}
\]
The evaluation identity proves the middle formula. Since
$\prod_i f'(r_i)=(-1)^6\prod_{i<j}(r_i-r_j)^2=(\det V)^2$,
taking determinants proves the last equality, with this precise quartic sign.

The exact algebra conductor $I=(E:O)=\{z\in O:zO\subset E\}$ is the
$R_p$-lattice with these eight original-coordinate generators:
\[
 \boxed{A p^{m_i}f_i(r),\qquad A\sigma f_i(r),\quad0\le i\le3.}
 \qquad
 J_I=\diag(AQ\diag(p^{m_i}),AQ).
 \tag{PS6}
\]
To prove equality with the conductor, retain the unchanged residue
\[
 \Lambda(R+\sigma T)=A^{-1}[r^3](T\bmod f),\quad
 \Lambda(R+\sigma T)=\sum_i\frac{T(r_i)}{d_i}.
 \tag{PS7}
\]
The second equality is Lagrange interpolation: the coefficient of $r^3$
in $f_i/f'(r_i)$ is $1/f'(r_i)$. The full Gram $B_E$ is exactly RD3,
whose determinant $A^{-8}$ is a unit. All of its entries and its inverse
are integral here. Hence $E$ is its own residue-dual lattice:
$E^\#=\{z:\Lambda(zE)\subset R_p\}=E$.

On the $i$th factor of $O$, the residue Gram in $(1_i,\tau_i)$ is
\[
 B_{O,i}=\begin{pmatrix}
 0&(p^{m_i}d_i)^{-1}\\(p^{m_i}d_i)^{-1}&0
 \end{pmatrix}.
 \tag{PS8}
\]
Indeed only the odd component contributes, and
$\tau_i=\sigma_i/p^{m_i}$. Thus
$O^\#=\prod_i(p^{m_i}d_i)O_i$.
For $z\in O^\#$ and $o\in O$, every $e\in E$ gives
$\Lambda(zoe)\in R_p$ because $oe\in O$; hence $zo\in E^\#=E$.
Conversely $zO\subset E$ implies $\Lambda(zO)\subset R_p$.
Therefore
\[
 I=O^\#=\prod_i(p^{m_i}d_i)O_i
       =\prod_i p^{\,3m_i+\epsilon_i}u_iO_i .
 \tag{PS9}
\]
The two generators on branch $i$ pull back by the labelled Lagrange
polynomial $L_i=f_i/f'(r_i)$ to
$p^{m_i}d_iL_i=Ap^{m_i}f_i$ and
$d_i\sigma L_i=A\sigma f_i$. This proves PS6 with every scalar unit.

More than the index is determined: the full Smith factors of $E/I$
are the same as those of $O/E$. Let $B_O$ be PS8 in the order PS4.
Then
\[
 B_E=C_p^{\mathsf T}B_OC_p,\qquad
 C_p^{-1}B_O^{-1}=B_E^{-1}C_p^{\mathsf T}.
 \tag{PS10}
\]
The left matrix generates $I$ in the original basis. Since $B_E^{-1}$
is an integral invertible matrix, its Smith factors are those of
$C_p^{\mathsf T}$, hence those of $C_p$.
Equivalently the exact pairing
\[
 E/I\times O/E\longrightarrow K_p/R_p,\qquad
 (\bar e,\bar o)\longmapsto\Lambda(eo)\bmod R_p
 \tag{PS11}
\]
is perfect. It is well defined by $I=O^\#$ and $E=E^\#$.
Its left kernel is $I$ and its right kernel is $E$ by these same
equalities; the equal elementary divisors from PS10 prove perfectness
on each finite cyclic factor.

In particular, for $g=\sum_{i<j}v_p(r_i-r_j)$,
\[
 \begin{aligned}
 \ell(O/E)=\ell(E/I)&=2g+\sum_i m_i
                   =3g-\tfrac12\sum_i\epsilon_i,\\
 \ell(O/I)&=2\ell(O/E),\\
 \det J_I&=A^8(\det V)^2p^{\sum_i m_i}.
 \end{aligned}\tag{PS12}
\]
The first equality follows also by taking determinants in PS4 and PS6.
The branch length in PS9 is $2(3m_i+\epsilon_i)$, confirming the second.

\subsection{The exact rank-twenty-four to rank-eight specialization}
Keep the coefficient names $C_0,D_0$ for the actual witness, so
\[
 C_0=-5Axyz,\qquad D_0=Apxyz,\qquad D_0=-pC_0/5 .
 \tag{PS13}
\]
The ring constructions ET1--20 work over
$\mathbb Z[1/10,A,A^{-1},B,C]$, because their displayed matrices
use only these inverted units and their proofs use monic division.
Specialize $(A,B,C)$ to the actual values $(A,B_0,C_0)\in R_p^3$.
The resulting two rank-twenty-four orders have exactly the original
ET bases and the inclusion $K(C_0)\otimes I_8$. Selecting the given
witness is the following pair of surjective algebra maps:
\[
 \begin{aligned}
 \Pi_D&=(I_8,\ D_0I_8,\ D_0^2I_8),&
 \Pi_P&=(I_8,\ pI_8,\ p^2I_8),\\
 \Pi_P(K(C_0)\otimes I_8)&=\Pi_D,&
 C_p\Pi_P(K(C_0)\otimes I_8)&=C_p\Pi_D .
 \end{aligned}\tag{PS14}
\]
The $P$ substitution satisfies its original monic cubic because the
witness satisfies the ES relation. The $D$ substitution satisfies ET1
by the same identity. They are algebra maps on the whole signed algebra:
the $r,\sigma$ images retain exactly $f$ and $h'$ at this coefficient
point. Entrywise substitution of $D_0=-pC_0/5$ proves both matrix
identities; the leading $I_8$ proves surjectivity.

The full kernels are free of rank sixteen, with no lost components:
\[
 \begin{aligned}
 \ker\Pi_D&=\{(-D_0u-D_0^2v,u,v):u,v\in R_p^8\},\\
 \ker\Pi_P&=\{(-pu-p^2v,u,v):u,v\in R_p^8\}.
 \end{aligned}\tag{PS15}
\]
Indeed these formulas solve the displayed equations without division.
Thus the rank-twenty-four coefficient comparison and the rank-eight
integral comparison have an exact commuting specialization, rather
than an identification of their different cokernels.

Its integral-ideal effect is fully calculable. Put $c=v_p(C_0)$.
The coefficient conductor $C_0^2 E_N$ specializes under $\Pi_P$ to
$C_0^2E$. Its inclusion in the arithmetic conductor $I$ holds exactly when
\[
 \boxed{2c\ge\max_i(3m_i+\epsilon_i).}
 \tag{PS16}
\]
Since $I$ is an $E$-ideal, the inclusion is equivalent to $C_0^2\in I$.
Branchwise PS9 turns that assertion into precisely the displayed
inequalities; the unit $u_i$ does not change a valuation.
In the one-divisible case $c=1$ and the other unit-root gap has
valuation $n=v_p(x-y)$; the branch exponents are
$1,1,3\lfloor n/2\rfloor+(n\bmod2)$ twice, up to the labelled order.
Hence PS16 holds exactly for $n\le1$.
In the two-divisible case $c=2$ and the three-root cluster has two gaps
of valuation one and one gap of valuation $n\ge1$.
Its exponents are zero at the isolated unit root, three at the separated
cluster root, and
$3\lfloor(n+1)/2\rfloor+((n+1)\bmod2)$ at both closer roots.
Hence PS16 holds exactly for $n\le2$.
These are complete alternatives on the classified ES strata.
Failure of this inclusion for larger gaps is an exact obstruction
between these two particular conductor ideals, while PS14 continues
to give their full algebraic relation.

\subsection{The two actual witnesses, all labels and signs}
At $p=1201$ the two input witnesses are
\[
 (x,y,z)=(306,16218,1082101),\qquad
 (x,y,z)=(306,21618,61251).
 \tag{PS17}
\]
For the first,
$1/16218+1/1082101=23/(306\cdot1201)$; adding $1/306$
gives $4/1201$. For the second, the last two denominators are
$1201\cdot18,1201\cdot51$ and $1/18+1/51=23/306$, giving
the same exact equality. The literal sums are $1099826$ and $84376$.
In the ordered roots $(p,x,y,z)$ their complete coefficient vectors are
\[
 \begin{aligned}
 u^{(1)}={}&\left(
 -\frac1{1099826},-\frac{19205048257}{1099826},
 \frac{13425378223770}{549913},
 -\frac{3224775849349554}{549913}\right),\\
 u^{(2)}={}&\left(
 -\frac1{84376},-\frac{1449375207}{84376},
 \frac{506477475135}{21094},
 -\frac{121655889527427}{21094}\right).
 \end{aligned}\tag{PS18}
\]
Multiplying the four factors in PS1 gives every entry.
Modulo $1201$ these are $(410,100,0,0)$ and $(522,0,0,0)$.
The complete valuation and branch-unit lists are
\[
 \begin{array}{c|c|c|c|c}
 & (b_i)&(m_i)&(\epsilon_i)&(u_i\bmod1201)\\ \hline
 1&(1,0,0,1)&(0,0,0,0)&(1,0,0,1)&(75,93,581,1126)\\
 2&(2,0,2,2)&(1,0,1,1)&(0,0,0,0)&(850,42,640,449).
 \end{array}\tag{PS19}
\]
Here $u_i$ in the last column is precisely $p^{-b_i}Af'(r_i)$,
not the coefficient-vector notation in PS18.
These lists can be checked by the exact product
$Af'(r_i)=A\prod_{j\ne i}(r_i-r_j)$; every difference is an explicit
integer from PS17. The Smith exponent lists of both $O/E$ and $E/I$ are
\[
 (0,0,0,0,0,0,1,1),\qquad
 (0,0,0,1,1,2,2,3),
 \tag{PS20}
\]
respectively. The branch conductor exponents are $(1,0,0,1)$ and
$(3,0,3,3)$. The entire original generator matrix in either example
is PS5--6 evaluated at the four stated integers; these formulas
specify all sixty-four rational entries without deleting any direction.
The accompanying exact certificate additionally records every entry.

At the common reduced root zero, the two full original primary algebras are
\[
 \mathbb F_{1201}[r,\sigma]/(r^2,\sigma^2-200r),\qquad
 \mathbb F_{1201}[r,\sigma]/(r^3,\sigma^2-3r^2).
 \tag{PS21}
\]
They follow from $h'$ reduced modulo $r^2$ and $r^3$,
respectively, retaining the coefficients $410,100$ and $522$.
The remaining simple roots and signs are still the other PS3 factors;
PS21 does not stand for the entire eight-dimensional algebra.

Choose $i\in K_p$ with $i^2=-1$, possible since $p\equiv1\pmod4$,
and for each labelled root choose both signs
$\sigma_{i,\pm}=\pm\sqrt{d_i}$ in its quadratic algebraic closure.
Every original affine inverse coordinate is retained by
\[
 \begin{aligned}
 q_{i,\pm}=\big(&\sigma_{i,\pm}^{-1},\
 -r_i-i\sigma_{i,\pm},\
 A\sigma_{i,\pm}^3+2r_i\sigma_{i,\pm}+3i\sigma_{i,\pm}^2,\\
 &7ir_i^2\sigma_{i,\pm}
 +(B_0-17r_i+Ar_i^2)\sigma_{i,\pm}^2
 -13i\sigma_{i,\pm}^3-2A\sigma_{i,\pm}^4\big).
 \end{aligned}\tag{PS22}
\]
FS13 verifies $P(q_{i,\pm})=(A,B_0,C_0,D_0)$ exactly.
All $d_i$ are nonzero in characteristic zero, so all eight points exist.
Their first-coordinate valuations are $-b_i/2$:
in the first example $(-1/2,0,0,-1/2)$ and in the second
$(-1,0,-1,-1)$, for each of the two signs.
The integral boundary therefore records escaping affine coordinates
without changing the original polynomial Jacobian $-2$.

\subsection{The same rational algebra in the two original complex receivers}
For these rational witnesses define $E_{\mathbb Q}$ by PS3 over
$\mathbb Q$ with its literal rational $A,B_0,C_0,D_0$.
Its two base changes to $K_p$ and $\mathbb C$ retain the very same
eight coefficient basis vectors. The first has the maps PS4--22.
For the second, use exactly $F_U,F_W$ of GD14.
Thus, for example, the full arithmetic conductor generator matrix has
the two original complex realizations
\[
 \mathcal I_U=F_UJ_I,\qquad
 \mathcal I_W=F_WJ_I,\qquad
 (T_A\oplus T_A)\mathcal I_U=\mathcal I_W .
 \tag{PS23}
\]
The columns are literally
$(T_A^{-1}(Ap^{m_i}f_i)(S'-\kappa),0)$ and
$(0,T_A^{-1}(Af_i)(S'-\kappa))$ at the source, with the
$T_A^{-1}$ removed at the target. Their exact squared Gram matrices are
\[
 J_I^*\Gamma_UJ_I,\qquad J_I^*\Gamma_WJ_I,\quad
 \det(J_I^*\Gamma_bJ_I)
   =|A|^{16}|\det V|^4p^{2\sum m_i}\det\Gamma_b .
 \tag{PS24}
\]
Here $\Gamma_b=\diag(G_{b,r},G_{b,r})$ contains the entire original
moments and Gamma masses, and the displayed determinant is PS12
times its complex conjugate. It is nonzero for both actual witnesses.
These are exact complex embeddings of an integral lattice, not an
assignment of a $p$-adic valuation to a transcendental Gamma mass.

The coefficient specialization itself also has a complete complex
minimum in those same metrics. For either receiver set
$\Gamma=\Gamma_b$ and
\[
 \Pi_D=(I,D_0I,D_0^2I),\qquad
 d_D=1+|D_0|^2+|D_0|^4 .
\]
For every prescribed column $v\in\mathbb C^8$, the exact original
quotient minimum and its unique orthogonal representative are
\[
 \min_{\Pi_D z=v}\sum_{a=0}^2 z_a^*\Gamma z_a
       =\frac{v^*\Gamma v}{d_D},\qquad
 z_{\min}=\frac1{d_D}(v,\bar D_0v,\bar D_0^{\,2}v).
 \tag{PS25}
\]
Multiplication proves the constraint. For a kernel vector
$w$ in PS15, the inner product with this $z_{\min}$ is
$d_D^{-1}v^*\Gamma(w_0+D_0w_1+D_0^2w_2)=0$.
Pythagoras proves the minimum, uniqueness and the exact excess
$\|z\|^2=\|z_{\min}\|^2+\|z-z_{\min}\|^2$.
The identical proof for $\Pi_P$ replaces $D_0$ by the actual prime $p$.
Thus PS14 and PS25 give both the algebraic specialization and the
attained full-kernel minimum, with the original metric retained.

For the source-to-target specialization
$F_W\Pi_D(I_3\otimes F_U^{-1})$, the eight nonzero singular values
are exactly $\sqrt{d_D}$ times the eight original singular values
of $T_A\oplus T_A$, and its kernel is sixteen-dimensional.
Indeed its orthonormal representation is the tensor product of
the row $(1,D_0,D_0^2)$ with the original conductor representation;
the row has the sole positive singular value $\sqrt{d_D}$.
This provides the complete map relating the rank-twenty-four
coefficient cost to the rank-eight actual arithmetic witness.
It does not transfer a coefficient-lattice pole into a vanishing
original conductor moment.


% Independent finite arithmetic derivation. Integration body, no document wrapper.
\section{The complete arithmetic signed order and its algebra conductor}
This section independently verifies equations (9)--(31) of the preserved
user-supplied defining-prime continuation
\cite[eqs.\ (9)--(31)]{DefiningPrimeReceived}. Its original signed algebra,
residue convention, and inverse factor coordinates are RD1--8 and FS13
of the cumulative manuscript, retaining the human polynomial, ES-reader,
and residue provenance \cite{ES488,Tao,JouveRV}. The original ES marking is retained as
$(\mathrm M,\mathrm N,\mathrm T,\mathrm E)=(p,x,y,z)$ in this section;
this is a declared marking and does not identify a historical operator
named $A$ with the formal label $\mathrm N$. The four labels and both
signs at each label are always retained. The algebra conductor calculated
below is an ideal of a finite arithmetic algebra. Its exact transport to
the original weighted linear conductor is supplied at the end; those two
uses of the word conductor are not identified by an assertion.

Let $p$ be a rational prime with $p\equiv1\pmod {12}$ and let $x,y,z$
be positive integers satisfying $4/p=1/x+1/y+1/z$. Put
\begin{equation}
 r_0=p,\quad r_1=x,\quad r_2=y,\quad r_3=z,\qquad
 S=p+x+y+z,\quad f(r)=\prod_{i=0}^3(r-r_i),\quad
 A=-S^{-1},\quad h=Af=Ar^4+r^3+Br^2+Cr+D.
 \tag{ISR1}
\end{equation}
All coefficients in ISR1 are literal coefficients of the supplied
quartic. In particular its cubic coefficient equals one.

\subsection{All actual witness valuation patterns}
The four integers $r_i$ are pairwise distinct. To prove distinctness of
two denominators, suppose $y=z$ and set $d=4x-p>0$. The reciprocal
equation gives $dy=2px$, hence $2dy=p(d+p)$ and $d\mid p^2$.
But $d\equiv-p\equiv3\pmod4$, whereas all positive divisors of $p^2$
are $1\pmod4$. If $x=p$, the equation gives
$(3y-p)(3z-p)=p^2$. Its two factors are positive and $2\pmod3$;
every positive divisor of $p^2$ is $1\pmod3$. These contradictions,
and permutation of the denominator labels, prove the assertion.

The exact equation $4xyz=p(xy+xz+yz)$ implies that at least one
denominator is divisible by $p$. All three cannot be, since writing
them $pX,pY,pZ$ would give $4=1/X+1/Y+1/Z\le3$. If only $z$ is
divisible, valuation comparison gives $z=pZ$ with $Z$ a unit, and
division of the equation by $p$ gives
\begin{equation}
 4xyZ=xy+pZ(x+y),\qquad Z\equiv1/4\pmod p,
 \qquad v_p(p-z)=1.                                      \tag{ISR2}
\end{equation}
For two divisible denominators write their valuations $a,b\ge1$.
If $a<b$, then the right side of the exact equation has valuation
$a+1$, so $a+b=a+1$, a contradiction to $b>a\ge1$.
Thus $a=b$. If $a=b\ge2$, positivity gives
$p/4<x\le p^2/(4p-2)<(p+1)/4$, an interval containing no integer
because $p\equiv1\pmod4$. Therefore, after declaring the denominator
permutation, the only two cases are
\begin{equation}
 \begin{array}{ll}
 \text{one divisible:}&z=pZ,\ p\nmid xyZ;\\
 \text{two divisible:}&y=pY,\ z=pZ,\ p\nmid xYZ,
       \quad Y+Z\equiv4YZ\pmod p.
 \end{array}                                             \tag{ISR3}
\end{equation}

\subsection{Integral closure, every inclusion entry, and residue duality}
First suppose $p\nmid S$. Set $R=\mathbb Z_p$, $K=\mathbb Q_p$ and
\begin{equation}
 E=R[r,\sigma]/(f(r),\sigma^2-Af'(r)),\qquad
 D_i=Af'(r_i),\quad n_{ij}=v_p(r_i-r_j),\quad
 b_i=\sum_{j\ne i}n_{ij},\quad
 m_i=\lfloor b_i/2\rfloor,\quad \epsilon_i=b_i-2m_i,\quad
 u_i=p^{-b_i}D_i\in R^\times.                            \tag{ISR4}
\end{equation}
The ordered basis of $E$ is $(1,r,r^2,r^3,\sigma,\sigma r,
\sigma r^2,\sigma r^3)$. Both relations are monic in the eliminated
variable, so successive polynomial division proves this is a free
basis, without a reduced-fibre assumption. Since the $r_i$ are distinct,
evaluation and labelled Lagrange interpolation give the generic
product and its full integral closure
\begin{equation}
 E_K=\prod_{i=0}^3 K[\sigma_i]/(\sigma_i^2-D_i),\qquad
 O=\prod_{i=0}^3 O_i,\qquad
 O_i=R[\tau_i]/(\tau_i^2-p^{\epsilon_i}u_i),\qquad
 \sigma_i=p^{m_i}\tau_i.                                \tag{ISR5}
\end{equation}
Here a quadratic algebra is allowed to split. To prove that $O_i$ is
the entire integral closure, every element of its generic algebra has
a unique expression $a+b\tau_i$, including in the split case. If it
is integral, its conjugate is integral; its trace $2a$ is in $R$.
Since $p\ne2$, $a\in R$. Then $(b\tau_i)^2=b^2p^{\epsilon_i}u_i$
is integral in $K$, so $2v_p(b)+\epsilon_i\ge0$ and $b\in R$.
Conversely every $a,b\in R$ give an integral element, using the monic
quadratic relation. This proves maximality in both cases. The case
$\epsilon_i=1$ is ramified: a square in $K$ cannot have odd valuation,
and $\tau_i$ has valuation $1/2$. The case $\epsilon_i=0$ is split
when $u_i$ is a residue square and is an unramified quadratic field
when it is a residue nonsquare. Indeed a residue root lifts uniquely
by the recurrence $a_{n+1}=a_n+p^n d_n$ solving
$2a_n d_n=(u_i-a_n^2)/p^n\pmod p$; a residue nonsquare gives a
degree-two residue field in the displayed maximal ring.

Use the normalization basis $(1_0,1_1,1_2,1_3,\tau_0,\tau_1,
\tau_2,\tau_3)$, so that its unit is $(1,1,1,1,0,0,0,0)$.
With $V=(r_i^j)_{0\le i,j\le3}$ the complete inclusion and inverse are
\begin{equation}
 C_p=\operatorname{diag}(V,\operatorname{diag}(p^{m_i})V),\qquad
 C_p^{-1}=\operatorname{diag}(V^{-1},
                      V^{-1}\operatorname{diag}(p^{-m_i})),
 \quad (V^{-1})_{ji}=[r^j]\frac{f_i(r)}{f'(r_i)},\quad
 f_i(r)=\prod_{j\ne i}(r-r_j).                           \tag{ISR6}
\end{equation}
Direct evaluation proves both inverse compositions over $K$.
For $N=\sum_{i<j}n_{ij}$, its determinant and index are
\begin{equation}
 \det C_p=(\det V)^2p^{\sum m_i},\qquad
 L:=\operatorname{length}_R(O/E)
       =2N+\sum_i m_i=3N-\tfrac12\sum_i\epsilon_i.
                                                               \tag{ISR7}
\end{equation}
The determinant formula uses precisely the chosen root order, with
$\det V=\prod_{i<j}(r_j-r_i)$, and $\sum b_i=2N$.

Write a generic element as $R(r)+\sigma Q(r)$. The original RD
residue is
\begin{equation}
 \Lambda(R+\sigma Q)=\frac{[r^3](Q\bmod f)}{A}
    =\sum_i\frac{Q(r_i)}{D_i};\qquad
 \Lambda_i(a+b\tau_i)=\frac{b}{p^{m_i}D_i}.               \tag{ISR8}
\end{equation}
The coefficient equality follows by inserting the complete Lagrange
formula in ISR6: the cubic coefficient of $f_i$ is one. In the
normalization basis its bilinear pairing is
\begin{equation}
 B_O=\begin{pmatrix}0&H_O\\H_O&0\end{pmatrix},\qquad
 H_O=\operatorname{diag}((p^{m_i}D_i)^{-1}),\qquad
 B_E=C_p^{\mathsf T}B_OC_p
    =\begin{pmatrix}0&H_E\\H_E&0\end{pmatrix},\quad
 H_E=V^{\mathsf T}\operatorname{diag}(D_i^{-1})V.         \tag{ISR9}
\end{equation}
For a proof of integral self-duality that does not appeal to a
Gorenstein theorem, write $f=r^4+a_3r^3+a_2r^2+a_1r+a_0$.
Let $H_f=A H_E$. Its entries are cubic coefficients of $r^{i+j}$
modulo the monic $f$. Polynomial division gives
\begin{equation}
 H_f^{-1}=\begin{pmatrix}
 a_1&a_2&a_3&1\\a_2&a_3&1&0\\a_3&1&0&0\\1&0&0&0
 \end{pmatrix},\qquad \det H_f=1,
 \qquad H_E^{-1}=A H_f^{-1},\quad \det B_E=A^{-8}.        \tag{ISR10}
\end{equation}
To check the inverse, the cubic coefficients in degrees zero through
six are $0,0,0,1,-a_3,a_3^2-a_2,-a_3^3+2a_3a_2-a_1$.
Inserting these seven values proves both matrix products. Since $A$
is a unit, $B_E$ and its inverse are over $R$. Thus
$E^\#=\{z:\Lambda(zE)\subset R\}=E$.

Define the algebra conductor $I=\{z\in O:zO\subset E\}$.
Residue duality proves its exact value, including both quadratic
directions in every branch:
\begin{equation}
 I=O^\#=\prod_i p^{m_i}D_iO_i
   =\prod_i p^{3m_i+\epsilon_i}u_iO_i.                  \tag{ISR11}
\end{equation}
Indeed $O^\#\subset E^\#=E$. For $z\in O^\#$ and $w\in O$,
$\Lambda(zwE)\subset R$ because $wE\subset O$, hence $zw\in E$.
Conversely $zO\subset E$ implies $\Lambda(zO)\subset R$ by testing
against $1\in E$. Finally, the two-by-two branch matrix in ISR9
shows directly that $O_i^\#=p^{m_i}D_iO_i$. This proves equality
of ideals, not just an annihilator containment.

Every generator is also explicit in the original eight-element
basis. Let $Q$ be the four-column matrix of the entire cubics $f_i$:
\begin{equation}
 Q_{\cdot i}=\begin{pmatrix}
 -\prod_{j\ne i}r_j\\\sum_{j<k;\ j,k\ne i}r_jr_k\\
 -\sum_{j\ne i}r_j\\1
 \end{pmatrix},\qquad
 G_I=\operatorname{diag}(A Q\operatorname{diag}(p^{m_i}),A Q).
                                                               \tag{ISR12}
\end{equation}
Its columns are $Ap^{m_i}f_i$ and $A\sigma f_i$. Their images in
$O$ are exactly the eight columns of
$\operatorname{diag}(p^{m_i}D_i,p^{m_i}D_i)$, because
$VQ=\operatorname{diag}(f'(r_i))$. Their independence and equality
with ISR11 follow. Since $\prod_i f'(r_i)=(\det V)^2$ in degree
four, $\det Q=\det V$ and
\begin{equation}
 \det G_I=A^8(\det V)^2p^{\sum m_i},\qquad
 \operatorname{length}_R(E/I)=L,\qquad
 \operatorname{length}_R(O/I)=2\sum_i(3m_i+\epsilon_i)=2L.
                                                               \tag{ISR13}
\end{equation}
Each scalar ideal in the free rank-two $O_i$ contributes twice its
$p$-exponent. This verifies the stated length equality directly.
There is also an equality of every invariant factor, with a perfect
finite-module pairing rather than just a length comparison:
\begin{equation}
 E/I\ \times\ O/E\longrightarrow K/R,\qquad
 ([e],[o])\longmapsto\Lambda(eo)\bmod R,
 \quad E/I\simeq\operatorname{Hom}_R(O/E,K/R).
                                                               \tag{ISR13a}
\end{equation}
Adding $I$ to the first input or $E$ to the second changes the value
by an element of $R$, using ISR11 and self-duality of $E$.
Its first kernel is $E\cap O^\#=I$ and its second is
$O\cap E^\#=E$. A finite module with diagonal factors
$R/(p^{e_j})$ has $K/R$-dual with the same diagonal factors,
by sending a generator to any class in $p^{-e_j}R/R$.
Together with the equal lengths already proved, these trivial
kernels prove perfection. Directly at the level of matrices,
the conductor generator in the original basis can be chosen as
$C_p^{-1}B_O^{-1}=B_E^{-1}C_p^{\mathsf T}$. Since $B_E^{-1}$
is invertible over $R$, its Smith factors are those of $C_p$.
This generator differs from ISR12 only by interchanging the two
normalization parity blocks, so ISR16--17 describe $E/I$ as well.

\subsection{The full multiplication trace, with every residue factor}
Let $Z$ be the companion matrix of $f$, $L_h=Af'(Z)$, and retain
the original parity basis. Multiplication and its trace satisfy
\begin{equation}
 M_\sigma=\begin{pmatrix}0&L_h\\I_4&0\end{pmatrix},\quad
 M_{2\sigma^3}=2\begin{pmatrix}0&L_h^2\\L_h&0\end{pmatrix},\quad
 \operatorname{Tr}(M_F)=\Lambda(2\sigma^3 F),\quad
 T_E=B_E M_{2\sigma^3}
       =\operatorname{diag}(2V^{\mathsf T}V,
                            2V^{\mathsf T}\operatorname{diag}(D_i)V).
                                                               \tag{ISR14}
\end{equation}
The first two matrices are the original quotient multiplication.
For $F=R+\sigma Q$, its branch trace is $2R(r_i)$, while its
odd coefficient after multiplication by $2\sigma^3$ is
$2D_iR(r_i)$; ISR8 proves the functional equality. The even-even
trace is therefore $2\sum r_i^{j+k}$, the odd-odd trace is
$2\sum D_i r_i^{j+k}$, and mixed entries vanish. No field-trace
weight has been omitted. In the normalized basis,
\begin{equation}
 T_O=\operatorname{diag}(2I_4,2\operatorname{diag}(p^{\epsilon_i}u_i)),
 \quad T_E=C_p^{\mathsf T}T_OC_p,\quad
 \det T_E=256A^4(\det V)^6,\quad
 v_p(\det T_E)=6N,\quad v_p(\det T_O)=\sum\epsilon_i.
                                                               \tag{ISR15}
\end{equation}
Taking determinants in the congruence gives their difference $2L$.
Equations ISR9 and ISR14 prove the exact residue-to-trace morphism;
reduction of trace can lose rank even when the residue pairing is
perfect. This does not create a kernel in the original complex
weighted conductor.

\subsection{All invariant factors, including closer pairs}
For a matrix over $R$, the sum of its first $j$ Smith exponents is
the least valuation of a nonzero $j$-minor. This follows from invertible
row and column operations: the ideals generated by all $j$-minors
are unchanged, and a diagonal matrix with divisibility-ordered entries
has least product its first $j$ entries. The diagonal reduction exists
by repeatedly choosing an entry of minimum valuation and clearing its
row and column; all other entries are its multiples in a DVR. This
also proves the criterion used here without leaving a classification
step implicit.

Residue clusters whose differences are units separate over $R$: the
resultant of their monic polynomials is a product of unit differences.
The Sylvester determinant gives polynomials whose combination is a
unit, so their ideals are comaximal and the root algebra splits.
This split acts on both parity copies. A two-root block with gap
valuation $n$ has even Smith exponents $0,n$, by subtracting its
first row and then its constant multiple. Both derivative valuations
are $n$ when all external differences are units. Its odd exponents
are therefore $\lfloor n/2\rfloor,n+\lfloor n/2\rfloor$.
In the one-divisible case ISR2 separates $(p,z)$ from $(x,y)$, so
for $n=v_p(x-y)\ge0$ the eight exponents, as a multiset, are
\begin{equation}
 \{0,0,1,1\}\ \cup\
 \{0,\lfloor n/2\rfloor,n,n+\lfloor n/2\rfloor\},\qquad
 L=2+3n-(n\bmod2).                                      \tag{ISR16}
\end{equation}
The second block is allowed to have $n=0$, in which case it is the
sum of two separated singleton blocks and the formula remains literal.

In the two-divisible case the three clustered roots are $p,pY,pZ$.
The relation in ISR3 prevents $1,Y,Z$ from all coinciding: that would
give $2=4\pmod p$, impossible. If $Y=Z\pmod p$ they both equal
$1/2$; if $Y=1$, then $Z=1/3$, and conversely. Thus one has either
three differences of valuation one, or exactly one closer pair. Put
$n$ equal to the largest valuation and $m=\lfloor(n+1)/2\rfloor$.
After a row permutation the valuations are $(1,1,n)$ and derivative
valuations $(2,n+1,n+1)$, with row multipliers $(1,m,m)$.
The ordinary three-root Vandermonde has successive minimum minor
valuations $0,1,n+2$, hence exponents $0,1,n+1$.
For the odd weighted block its one-minor minimum is $1$.
The two-minors in columns $1,r$ have valuations $m+2,m+2,2m+n$;
all other two-minors have at least these valuations because the
corresponding determinant is a root difference times an integral
symmetric polynomial. Since $m+n\ge2$, their minimum is $m+2$.
The determinant has valuation $n+2+1+2m=n+2m+3$. Thus the odd
exponents are $1,m+1,n+m+1$. Adding the isolated root's two unit
factors proves the entire list
\begin{equation}
 0,0,0,1,1,m+1,n+1,n+m+1,\qquad
 L=3n+6-((n+1)\bmod2).                                 \tag{ISR17}
\end{equation}
For $n=1$ the list is $0,0,0,1,1,2,2,3$ and all quadratic
branches have even derivative valuation: the normalized fields are
unramified or split although the integral fibre has a nonreduced
six-dimensional cluster.

\subsection{The nonunit leading coefficient and the exact quadratic twist}
Suppose now $k=v_p(S)>0$. The two-divisible case is impossible because
$S\equiv x\not\equiv0\pmod p$. In the one-divisible case,
$x+y\equiv0\pmod p$; therefore $x-y$ is a unit, since $p\ne2$.
The original derivative valuations, without changing any coefficient,
are
\begin{equation}
 b_i^0=\sum_{j\ne i}n_{ij}-k,\quad
 m_i^0=\lfloor b_i^0/2\rfloor,\quad
 \epsilon_i^0=b_i^0-2m_i^0,\quad u_i^0=p^{-b_i^0}Af'(r_i),
 \quad O_i^0=R[\tau_i^0],\quad
 \sigma_i^0=p^{m_i^0}\tau_i^0,\quad
 (\tau_i^0)^2=p^{\epsilon_i^0}u_i^0.                     \tag{ISR18}
\end{equation}
The same trace-and-square proof used in ISR5 proves that this is the
full integral closure of the original generic algebra, even when
$m_i^0<0$. In label order $(p,x,y,z)$ the valuations are
$(1-k,-k,-k,1-k)$. In particular the original $\sigma$ is nonintegral
on the two unit-denominator branches. Thus $R[r,\sigma]$ is not a
finite integral order; calling its eight coefficient columns an
integral order would be false.

Those eight columns still span a full fractional $R$-lattice $L_0$,
with exact comparison matrix
\begin{equation}
 C_0=\operatorname{diag}(V,\operatorname{diag}(p^{m_i^0})V).
                                                               \tag{ISR19}
\end{equation}
The same separated clusters give its fractional Smith exponent
multiset relative to $O^0$:
\begin{equation}
 \{0,0,0,1,\mu,\mu,\nu,\nu+1\},\qquad
 \mu=\lfloor-k/2\rfloor,\quad \nu=\lfloor(1-k)/2\rfloor.
                                                               \tag{ISR20}
\end{equation}
Negative exponents are retained. If these exponents are $e_j$, then
$\operatorname{length}(O^0/(O^0\cap L_0))=\sum\max(e_j,0)$ and
$\operatorname{length}(L_0/(O^0\cap L_0))=\sum\max(-e_j,0)$.
These follow by computing the intersection in each diagonal coordinate.
There is no assertion of an undefined quotient $O^0/L_0$.
The original residue matrix still equals $A^{-1}$ times its monic
matrix from ISR10, hence $L_0^\#=A L_0$, not $L_0$.

The shifted integral chart uses exactly
\begin{equation}
 r_i^1=r_i-1,\quad A_1=-1/(S-4),\quad
 f_1(r)=f(r+1),\quad h_1(r)=A_1f_1(r),\quad
 h_1(r)=q h_0(r+1),\qquad q=S/(S-4).                   \tag{ISR21}
\end{equation}
Since $p>3$ and $p\mid S$, $A_1$ is a unit. The differences are
unchanged, so ISR4--17 apply to this shifted integral order. They
do not retroactively turn $L_0$ into an order. Choose an element
$\eta$ satisfying $\eta^2=q$ and put $K'=K(\eta)$. Then the exact
generic algebra map, including both signs and its inverse, is
\begin{equation}
 E^1_K\otimes_KK'\longrightarrow E^0_K\otimes_KK',\quad
 r^1\longmapsto r^0-1,\quad \sigma^1\longmapsto\eta\sigma^0;
 \qquad r^0\longmapsto r^1+1,\quad
 \sigma^0\longmapsto\eta^{-1}\sigma^1.                 \tag{ISR22}
\end{equation}
Substituting in both relations proves the maps and both compositions.
The other root $-\eta$ interchanges the two sign labels uniformly.
There is no assumption that $q$ is a square in $K$ or that the two
generic covers are $K$-isomorphic with their stated labels. At each
branch the maximal-basis relation after extension is the literal
identity
\begin{equation}
 \tau_i^1=\eta p^{m_i^0-m_i^1}\tau_i^0,
 \qquad v_p(\eta p^{m_i^0-m_i^1})
       =(\epsilon_i^1-\epsilon_i^0)/2.                 \tag{ISR23}
\end{equation}
No equality of unsaturated integral tensor products after a ramified
extension is needed or claimed. The relations in ISR22 concern the
generic algebras, and ISR18 and ISR5 specify the full integral closure
of each before extending scalars.

\subsection{Two fully evaluated $p=1201$ witnesses, with both factor signs}
The two input witnesses are
\begin{equation}
 (p;x,y,z)=(1201;306,16218,1082101),\quad
 (1201;306,21618,61251).                                 \tag{ISR24}
\end{equation}
Exact multiplication gives $4xyz=p(xy+xz+yz)$ for each. Their full
quartic coefficient rows, in the order $(A,B,C,D)$, are
\begin{equation}
 \begin{array}{c|rrrr}
 1&-1/1099826&-19205048257/1099826&13425378223770/549913&
     -3224775849349554/549913\\
 2&-1/84376&-1449375207/84376&506477475135/21094&
     -121655889527427/21094
 \end{array}                                             \tag{ISR25}
\end{equation}
The derivative values and their full normalization data are
\begin{equation}
\begin{array}{c|r|r|c|c|c|r|r}
 &r_i&D_i&b_i&m_i&\epsilon_i&u_i\bmod1201&
          (u_i/1201)\\\hline
1&1201&-7263764196750/549913&1&0&1&75&1\\
1&306&592542393300/42301&0&0&0&93&-1\\
1&16218&1399413626676/6043&0&0&0&581&-1\\
1&1082101&-6848076990075750/6043&1&0&1&1126&1\\\hline
2&1201&-548653280375/42188&2&1&0&850&-1\\
2&306&145309944600/10547&0&0&0&42&1\\
2&21618&2155674064104/10547&2&1&0&640&1\\
2&61251&-72523381379625/42188&2&1&0&449&-1
\end{array}                                             \tag{ISR26}
\end{equation}
Here $(u_i/1201)$ is the quadratic residue sign, verified by raising
the displayed residue to power $600$ modulo $1201$. The full unit
is $u_i=D_i/1201^{b_i}$, not its residue. Thus witness one has two
ramified and two unramified branches; witness two has two unramified
and two split branches. For the first witness only the $(p,z)$ gap
has positive valuation, equal to one. For the second the three
clustered gaps all have valuation one. Their Smith lists are,
respectively, $(0,0,0,0,0,0,1,1)$ and $(0,0,0,1,1,2,2,3)$.
Their exact algebra conductor branch exponents are $(1,0,0,1)$
and $(3,0,3,3)$, giving indices $L=2,9$ and lengths $O/I=4,18$.

For each individual row of ISR26 select either root
$s_i$ of $s_i^2=D_i$. The two original FS13 factor states are the
following complete four-coordinate columns, with $\varepsilon=+1,-1$:
\begin{equation}
 q_{i,\varepsilon}=\begin{pmatrix}
 \varepsilon/s_i\\
 -r_i-i\varepsilon s_i\\
 \varepsilon(As_i^3+2r_i s_i)+3iD_i\\
 i\varepsilon(7r_i^2s_i-13s_i^3)
          +(B-17r_i+Ar_i^2)D_i-2AD_i^2
 \end{pmatrix}.                                        \tag{ISR27}
\end{equation}
Every $r_i,D_i,A,B$ in this formula is evaluated in ISR25--26, so
it specifies all sixteen original factor states across the two
witnesses without identifying either sign. Changing the choice of
$s_i$ interchanges the two stated columns and changes no label.
Over $\overline{\mathbb Q}_p$ take either root of $i^2=-1$; for
the present primes $i$ already lies in $\mathbb Q_p$. Over $\mathbb C$
use the original $i$. The first coordinate has valuation $-b_i/2$.
In particular the first witness has two prime/divisible-label pairs
of signs with $v_p(a)=-1/2$, and the second has three with $v_p(a)=-1$.

For completeness, the original polynomial factor construction is
\begin{align*}
 b&=i+ay,&c&=-i+2ay+a^2z,\\
 d&=-iy-a(iz+2y^2)-a^2yz,&e&=2z-7iy^2+aw,\\
 f_c&=iw+3iy^3-4yz+a(6iy^2z+wy+4y^4)+2a^2y^3z,\\
 P(a,y,z,w)&=(ac,ae+bd,af_c+be,bf_c).
\end{align*}
Insert $a=1/t$, $y=-r-it$, $z=At^3+2rt+3it^2$ and the last
FS13 coordinate. Polynomial expansion gives
$P=(A,B,t^2-4Ar^3-3r^2-2Br,
3Ar^4+2r^3+Br^2-rt^2)$. Thus $t^2=h'(r)$ and $h(r)=0$
give $Pq_{i,\varepsilon}=(A,B,C,D)$ for both signs exactly.

Reducing ISR25 gives $(410,100,0,0)$ and $(522,0,0,0)$ modulo
$1201$. The corresponding clustered local algebras are
\begin{equation}
 \mathbb F_{1201}[r,\sigma]/(r^2,\sigma^2-200r),\qquad
 \mathbb F_{1201}[r,\sigma]/(r^3,\sigma^2-3r^2).          \tag{ISR28}
\end{equation}
The first is a length-four local algebra with a single length-four
$\sigma$ chain; the second has $\sigma$ chain lengths four and two.
The equality $a=1/\sigma$ explains why these integral boundary points
are not integral lifts of the original affine factor coordinates.
It does not alter the original affine Jacobian $\det DP=-2$.

\subsection{The exact nilpotent statement and exterior cofactor field}
Let $m\ge2$ and let $k$ be a field. In the algebra
$k[r,\sigma]/(r^m,\sigma^2-r^{m-1}u(r))$ with $u(0)\ne0$,
the reduced relation is $\sigma^2=u(0)r^{m-1}$. Thus
\begin{equation}
 \operatorname{rank}M_\sigma=m+1,\quad
 \operatorname{rank}M_{\sigma^2}=2,\quad
 \operatorname{rank}M_{\sigma^3}=1,\quad M_{\sigma^4}=0;
 \quad\text{Jordan lengths }4,\underbrace{2,\ldots,2}_{m-2}.
                                                               \tag{ISR29}
\end{equation}
Indeed $M_\sigma(R+\sigma Q)=u(0)r^{m-1}Q+\sigma R$ has
independent images of dimensions one and $m$. Multiplication by
$\sigma^2$ has the two independent images $r^{m-1}$ and
$\sigma r^{m-1}$; multiplication by $\sigma^3$ has just the second.
These ranks determine the Jordan counts by successive rank
differences, and also give two length-$m$ blocks for $M_r$.
The restriction $m\ge2$ is essential. For a derivative relation
$\sigma^2=h'$ at a multiplicity-$m$ root with $h=r^m g$,
$u(0)=mg(0)$, so one also needs $\operatorname{char}k\nmid m$.
Our examples have $p=1201$ and $m=2,3$; both requirements hold.

Write the original exterior parameters as $h_e,r_e,s_e,\kappa_e$
to distinguish them from the signed-algebra root variable. The
preserved formulas are $x=h_er_es_e$, $y=h_es_e\kappa_e$ and
$z=ph_er_e\kappa_e$. Thus $Z=h_er_e\kappa_e$ and their exact
exterior cofactor is
$D_E=(r_e+\kappa_e)/s_e=Z(x+y)/(xy)$. In particular
$xy/(x+y)=Z/D_E$ follows by multiplying the displayed parameters,
without an additional reciprocal parametrization. In the unit-$S$
case $D_E$ is a $p$-unit, by this equality and ISR2. Moreover
$D_0/p=A(1-Z)(p-x)(p-y)$; reducing $A=-1/S$ and $Z=1/4$
modulo $p$ gives the first formula below, and dividing by $D_E$
gives the second:
\begin{equation}
 \frac{D_0}{p}\equiv-\frac3{16D_E}\pmod p,
 \qquad \frac{D_0}{pD_E}\equiv-\frac3{16D_E^2}\pmod p.
                                                               \tag{ISR30}
\end{equation}
Since $p\equiv1\pmod3$, the cyclic subgroup of cube roots in
$\mathbb F_p^\times$ contains a nontrivial $\zeta$ and
$(2\zeta+1)^2=-3$. The elementary lifting recurrence in the proof
of ISR5 therefore supplies an actual unit $c\in R^\times$ with
$c^2=D_0/(pD_E)$ and either chosen residue sign. Consequently
\begin{equation}
 K[\sigma]/(\sigma^2-D_0)\longrightarrow
 K[\omega]/(\omega^2-pD_E),\qquad
 \sigma\longmapsto c\omega,\qquad \omega\longmapsto c^{-1}\sigma
                                                               \tag{ISR31}
\end{equation}
are mutually inverse. For witness one $D_E=3$, $60^2=-3\pmod{1201}$,
so one can choose $c\equiv5\pmod{1201}$; the other choice is $-c$.
This retains the actual lifted unit, not just the quadratic class.

\subsection{Exact base change into the two original receiving metrics}
All the rational witness coefficients, matrices $C_p,G_I$, and
polynomial basis elements above belong to the rational algebra
$E_{\mathbb Q}=\mathbb Q[r,\sigma]/(f,\sigma^2-Af')$ when their
displayed powers of the fixed integer $p$ are included. Its two
base changes to $\mathbb Q_p$ and $\mathbb C$ are taken separately;
there is no embedding of a $p$-adic norm into a complex norm.
Use the actual original fixed maps from RD18--20,
\begin{equation}
 F_W(R+\sigma Q)=(R(S'-\kappa),Q(S'-\kappa)),\qquad
 F_U=(T_{\mathcal A}^{-1}\oplus T_{\mathcal A}^{-1})F_W,
 \quad \kappa=c_*-\tfrac12.                            \tag{ISR32}
\end{equation}
Both substitutions preserve the complete coefficients. In coefficient
matrices $F_W=\operatorname{diag}(C_{-\kappa},C_{-\kappa})$;
$F_U$ uses the full original inverse-conductor recurrence and its
unchanged moments. Every exact inclusion, generator and multiplication
has the concrete return
\begin{equation}
 J_W^O=F_W C_p^{-1},\quad J_U^O=F_U C_p^{-1},\quad
 J_W^I=F_WG_I,\quad J_U^I=F_UG_I,\quad
 J_W^O C_p=F_W=(T_{\mathcal A}\oplus T_{\mathcal A})F_U,
 \quad J_W^I=(T_{\mathcal A}\oplus T_{\mathcal A})J_U^I.
                                                               \tag{ISR33}
\end{equation}
These are identities on the complete eight-dimensional spaces.
For any multiplication matrix $M$, its actual receiving actions are
$F_WMF_W^{-1}$ and $F_UMF_U^{-1}$, intertwined by the same
original doubled conductor. In particular ISR14 gives the exact
residue-to-trace multiplier in those spaces; it is not declared to
be the original conductor itself.

If $G_W,G_U$ denote the unchanged original Gram matrices on these
two receiving spaces, the full returned Gram matrices are
\begin{equation}
 \Gamma_W=F_W^*G_WF_W,\quad\Gamma_U=F_U^*G_UF_U,\quad
 \Gamma_a^O=C_p^{-*}\Gamma_a C_p^{-1},\quad
 \Gamma_a^I=G_I^*\Gamma_a G_I\quad(a=U,W).               \tag{ISR34}
\end{equation}
For $k=0,\ldots,8$, the exact inverse exterior norm of any invertible
map $B$ from the first original receiving metric to the second is
the product of the largest $k$ singular values of
$G_U^{1/2}B^{-1}G_W^{-1/2}$. This follows by isometrically sending
the two metrics to Euclidean coordinates, applying singular value
decomposition, and taking exterior powers. Equations ISR33--34
therefore give the complete matrices for applying the original
equation (11), while retaining its two different metrics. Finite
$p$-adic invariant factors do not by themselves evaluate these
complex singular values or force an original conductor order jump.

\paragraph{Provenance and scope.}
The arithmetic witness data and proposed equations are from the
preserved user-supplied defining-prime continuation (9)--(31), not
attributed to an unnamed classical theorem. The original FS13 and
RD residue/trace construction are the cited prior programme proof.
The four formal ES labels retain the canonical version-69 488-page
reader's human provenance at its original proof locators. All
valuation classifications, matrix identities, conductor equalities,
negative-valuation correction, and quadratic-twist statements needed
here have been proved in this section; no claim is made to have read
new external literature in preparing this independent verification.

\section{Positivity closes the coefficient-to-arithmetic conductor inclusion}
This calculation applies the positivity of the original ES witness to the
entire arithmetic comparison PS1--25 and ISR1--34. The incoming account
\cite[eqs.\ (9)--(23)]{DefiningPrimeReceived} allowed additional residue
collisions. The following bounds determine which of those strata can
occur for actual positive witnesses. The four formal labels and their
eight signs remain the original ones \cite{ES488}; sorting denominator
values in an inequality means an explicitly chosen permutation of the
three denominator labels, never identifying two labels or changing the
quartic. All quantities in the final matrices remain in original order.

\subsection{The complete positive one-divisible bound}
Let $p\equiv1\pmod{12}$, $4/p=1/x+1/y+1/z$, and suppose just $z=pZ$
is divisible by $p$. ISR1--3 proves that $p\nmid xyZ$, that the four
roots are distinct, and that $4Z\equiv1\pmod p$. Choose the permutation
of the two unit denominator labels with $x<y$. Since $p\equiv1\pmod4$,
the least positive residue of $1/4$ is $(3p+1)/4$. Also each denominator
is strictly greater than $p/4$. Consequently
\[
 Z\ge\frac{3p+1}{4},\qquad x\ge\frac{p+3}{4},\qquad
 \frac1y=\frac4p-\frac1x-\frac1{pZ}
 \ge\frac{32}{(p+3)(3p+1)}.
 \tag{PB1}
\]
The last inequality follows by substituting the first two lower bounds
in the two subtracted reciprocals; its numerator is exactly $32p$ before
cancelling the common $p$. Therefore
\[
 0<y-x\le\frac{(p+3)(3p-7)}{32}<p^2,
 \qquad n=v_p(y-x)\in\{0,1\}.
 \tag{PB2}
\]
The strict inequality follows from
$32p^2-(p+3)(3p-7)=29p^2-2p+21>0$.
This argument retains the integer difference itself; a higher valuation
would require the positive difference to be at least $p^2$.
There is no assertion here that every remaining stratum is attained.

\subsection{Every positive two-divisible witness has separated residues}
Now let the only unit denominator be $x$, and write the other two as
$pY,pZ$, choosing the permutation with $Y<Z$. The exact equation is
\[
 \frac{4x-p}{x}=\frac1Y+\frac1Z,\qquad
 Y+Z\equiv4YZ\pmod p,\qquad
 1<Y\le\frac{2x}{4x-p}\le\frac{p+3}{6}<p.
 \tag{PB3}
\]
Here $Y=1$ would make a denominator equal to $p$, excluded in ISR1.
The reciprocal inequality gives the middle bound. The function
$2u/(4u-p)$ has derivative $-2p/(4u-p)^2<0$ for $u>p/4$;
at the least admissible integer $u=(p+3)/4$ it equals $(p+3)/6$.
This proves every displayed inequality.

The residues of $1,Y,Z$ are pairwise distinct. First $Y\not\equiv1$
because $1<Y<p$. If $Z\equiv1$, the original congruence would give
$3Y\equiv1$. Its least positive solution, since $p\equiv1\pmod3$,
is $(2p+1)/3$, strictly above $(p+3)/6$. If $Z\equiv Y$, it would
give $2Y\equiv1$ since $Y$ is a unit. Its least positive solution
is $(p+1)/2$, also strictly above $(p+3)/6$. These exhaust all possible
pairwise residue coincidences. Thus, in the original labelled cluster,
\[
 v_p(p-pY)=v_p(p-pZ)=v_p(pY-pZ)=1.
 \tag{PB4}
\]
In particular the closer-pair cases $n\ge2$ in the general ISR17
calculation never occur for a positive witness. Positivity removes them;
the validity of the more general matrix calculation is unchanged.

\subsection{The full ideal inclusion and its least coefficient power}
Continue with the original integral chart $p\nmid S$, so $A=-1/S$
is a unit. No coordinate or coefficient is changed to impose this
condition; ISR18--23 treats the different fractional chart explicitly.
Let $I=(E:O)$ be the complete algebra conductor PS6--12 and let $C_0$
be the original quartic coefficient. Up to the stated permutation,
the only possible lists are
\[
\begin{array}{c|c|c|c|c}
 \text{divisible denominators}&n&(b_i)&(3m_i+\epsilon_i)&
              \text{Smith exponents of }O/E\text{ and }E/I\\\hline
 1&0&(1,0,0,1)&(1,0,0,1)&(0,0,0,0,0,0,1,1)\\
 1&1&(1,1,1,1)&(1,1,1,1)&(0,0,0,0,1,1,1,1)\\
 2&1&(2,0,2,2)&(3,0,3,3)&(0,0,0,1,1,2,2,3).
\end{array}\tag{PB5}
\]
For the first two rows the label order is $(p,x,y,z)$; for the last
row it is $(p,x,pY,pZ)$. The complete full units remain
$u_i=p^{-b_i}Af'(r_i)$ and the ideal is still
$I=\prod_i p^{3m_i+\epsilon_i}u_iO_i$, not just its list of valuations.
PB5 follows by substituting PB2 and PB4 into the fully proved ISR16--17.
The two-divisible quadratic factors are all unramified or split,
because their derivative valuations are even.

The original coefficient is $C_0=-5Axyz$, so its valuation is respectively
one or two. For every integer $j\ge0$, PS9 gives the exact equivalence
\[
 C_0^jE\subset I\quad\Longleftrightarrow\quad
 jv_p(C_0)\ge\max_i(3m_i+\epsilon_i).
 \tag{PB6}
\]
Indeed $I$ is an $E$-ideal containing $C_0^jE$ exactly when it contains
the element $C_0^j$, and its branch ideals give the displayed criterion.
Combining PB5--6 proves, for every positive witness in this integral chart,
\[
 \boxed{\Pi_P(C_0^2E_N)=C_0^2E\subset I.}\qquad
 \min\{j\ge0:C_0^jE\subset I\}=
 \begin{cases}1,&\text{one divisible denominator},\\
 2,&\text{two divisible denominators}.
 \end{cases}\tag{PB7}
\]
The equality is the full surjective specialization PS14; it is not an
identification of the two different conductor ideals. The minima hold
because the largest branch exponent in PB5 is one or three, respectively.
Thus the coefficient conductor inclusion, formerly restricted by an
unevaluated gap condition, now holds for the whole actual positive
witness class in its integral chart.

All matrices realizing this inclusion in the original conductor remain
explicit: $F_UJ_I,F_WJ_I$ and $F_U\Pi_P(C_0^2I),F_W\Pi_P(C_0^2I)$,
with the unchanged original receivers PS23--25. For a direct eight-column
factorization put $J=J_I$ and define
\[
 Z_j=J^{-1}C_0^jI_8
 =\diag\!\left(C_0^j A^{-1}\diag(p^{-m_i})Q^{-1},
                   C_0^j A^{-1}Q^{-1}\right),\qquad
 JZ_j=C_0^jI_8.
 \tag{PB8}
\]
Every entry is integral for the values of $j$ in PB7 by the proved
lattice containment, and the displayed inverse specifies the full
matrix without losing any direction. Multiplication by the original
receivers gives $(F_bJ)Z_j=C_0^jF_b$ for $b=U,W$, and the two
identities intertwine by $T_A\oplus T_A$ exactly as in PS23.
The norms and quotient minima remain those in PS24--25. In particular
PB7 closes an arithmetic ideal comparison; it does not replace the
original linear conductor by a multiplication operator or force its
first nonzero moment to vanish.

% Independent extension of ISR. All original labels are retained.
\section{Positivity eliminates the deep defining-prime strata}
Retain the actual positive ES hypotheses of ISR1 and the classification
proved in ISR2--3. Permuting two denominator labels below is explicitly
only for the inequality proof; all original ordered matrices are returned
by the same permutation of their rows and branches. No two root or factor
sign labels are identified.

\subsection{The two unit denominators differ by less than $p^2$}
In the one-divisible case write $z=pZ$ and order the two unit denominators
as $x<y$. ISR2 gives $Z\equiv1/4\pmod p$. Since $p\equiv1\pmod4$,
the least positive residue of $1/4$ is $(3p+1)/4$. Positivity and the
strict inequality $1/x<4/p$ therefore imply
\begin{equation}
 Z\ge Z_0=(3p+1)/4,\qquad x\ge x_0=(p+3)/4,\qquad
 y=\frac{pZx}{Z(4x-p)-x}.                                \tag{PES1}
\end{equation}
The denominator at $(x_0,Z_0)$ is $2p>0$. On the entire rectangle
$x\ge x_0,Z\ge Z_0$, its derivative in $x$ is $4Z-1>0$ and its
derivative in $Z$ is $4x-p>0$, so the denominator remains positive.
For its exact gap function $g(x,Z)=y-x$, differentiation gives
\begin{equation}
 \partial_x g=-\frac{p^2Z^2}{(Z(4x-p)-x)^2}-1<0,\qquad
 \partial_Z g=-\frac{p x^2}{(Z(4x-p)-x)^2}<0.             \tag{PES2}
\end{equation}
Consequently every actual one-divisible witness satisfies
\begin{equation}
 y\le\frac{(p+3)(3p+1)}{32},\qquad
 0<y-x\le\frac{(p+3)(3p-7)}{32}<p^2,
 \quad v_p(y-x)\le1.                                  \tag{PES3}
\end{equation}
The last strict inequality is equivalent to
$29p^2-2p+21>0$. These bounds do not assume $p\nmid S$.
The real upper gap bound is attained by the actual witness
$(p;x,y,z)=(13;4,20,130)$, for which $y-x=16$.
This is sharpness of the displayed real inequality; it does not
assert that valuation one occurs among actual witnesses.

For later finite searches, ordering $x\le y$ also gives
$2/x\ge4/p-1/(pZ)$, hence
\begin{equation}
 x\le\frac{2pZ}{4Z-1}
   =\frac p2+\frac{p}{2(4Z-1)}
   \le\frac p2+\frac16,
 \quad x\le(p-1)/2.                                    \tag{PES4}
\end{equation}
All these inequalities are consequences of the original reciprocal
equation and positivity. They exclude every formal one-divisible
Smith stratum with $n\ge2$ from actual positive witnesses.

\subsection{Every two-divisible witness has separated scaled residues}
In the two-divisible case order $1\le Y<Z$ after retaining the
corresponding permutation of $(y,z)=(pY,pZ)$. Put $r=4x-p>0$.
The exact reciprocal equation and ISR3 give
\begin{equation}
 \frac r x=\frac1Y+\frac1Z,\quad
 Y+Z\equiv4YZ\pmod p,
 \quad Y\le\frac{2x}{4x-p}\le\frac{p+3}{6}<p.           \tag{PES5}
\end{equation}
The second inequality follows because $2x/(4x-p)$ has negative
derivative $-2p/(4x-p)^2$ and $x\ge(p+3)/4$.
The root distinctness proof in ISR1 implies $Y,Z\ne1$ as integers.
There can be no repeated residue among $1,Y,Z$:

If $Y\equiv1\pmod p$, its positive bound $Y<p$ forces $Y=1$,
contradicting the preceding fact. If $Z\equiv1\pmod p$, the reduced
equation forces $3Y\equiv1$. Since $p\equiv1\pmod3$, its least
positive solution is $(2p+1)/3$, larger than $(p+3)/6$.
Finally, if $Y\equiv Z\pmod p$, the same equation and their unit
property imply $2Y\equiv1$. Its least positive solution is
$(p+1)/2$, again larger than $(p+3)/6$. Therefore
\begin{equation}
 v_p(p-pY)=v_p(p-pZ)=v_p(pY-pZ)=1.                       \tag{PES6}
\end{equation}
This proves that the closer-pair parameter in ISR17 is always $n=1$
for an actual positive hard-prime witness. In particular $p\nmid S$
automatically, and the entire invariant-factor list and conductor data
are universal on this actual class:
\begin{equation}
 \begin{aligned}
 &\operatorname{Smith}(O/E)=\operatorname{Smith}(E/I)
       =(0,0,0,1,1,2,2,3),\qquad L=9,\\
 &(b_p,b_x,b_y,b_z)=(2,0,2,2),\qquad
 (m_p,m_x,m_y,m_z)=(1,0,1,1),\\
 &(\epsilon_p,\epsilon_x,\epsilon_y,\epsilon_z)=(0,0,0,0),
 \qquad (v_p I_i)=(3,0,3,3).
 \end{aligned}                                         \tag{PES7}
\end{equation}
Thus every normalized quadratic branch is unramified or split; the
nonreduced length-six clustered integral fibre remains present.

\subsection{The earlier coefficient ideal always enters the actual ideal}
For any actual positive witness with $p\nmid S$, the criterion PS16
now holds without an additional gap hypothesis. In the one-divisible
case $v_p(C_0)=1$ and PES3 gives $n\in\{0,1\}$, so the largest
branch exponent $3m_i+\epsilon_i$ is one. In the two-divisible case
$v_p(C_0)=2$ and PES7 gives largest branch exponent three. Hence
\begin{equation}
 C_0^2 E\subset I
 \quad\text{for every actual positive witness with }p\nmid S.
                                                               \tag{PES8}
\end{equation}
This is the exact specialization of the coefficient conductor
from PS14--16, not an equality of the two ideals. Both original
complex receiving maps of ISR32--34 apply unchanged to this proven
inclusion. For $p\mid S$ the original coefficient columns are the
fractional lattice of ISR18--23, so PES8 is not incorrectly stated
as an inclusion of an undefined original integral order.

\subsection{Exact parameters for the remaining two bounded questions}
The proofs above do not establish that the remaining one-divisible
case $n=1$ is realizable or impossible, and do not prove that $p\mid S$
is impossible. The following exact derivation constrains both questions.
Write $x=g a$, $y=g b$ with $a<b$, $\gcd(a,b)=1$ and $p\nmid gab$.
Set $d=4gab-p(a+b)>0$. From the ES equation,
$Z=gab/d$. The integer $d$ is coprime to both $a$ and $b$ because
its reductions modulo them are $-pb$ and $-pa$. Thus $d\mid g$.
Write $g=hd$. Rearranging the definition of $d$ gives
$d(4hab-1)=p(a+b)$. Also $p\nmid d$, since
$d\equiv4gab\not\equiv0\pmod p$. Consequently there are positive
integers $k,r$ with the complete identities
\begin{equation}
 \begin{aligned}
 &x=hda,\quad y=hdb,\quad Z=hab,\quad
 k=(4hab-1)/p,\quad a+b=dk,\\
 &r=4x-p=4had-p,\quad rk=4ha^2+1,\quad
 b=dk-a,\quad p=4had-r,\quad r\equiv k\equiv3\pmod4.
 \end{aligned}                                         \tag{PES9}
\end{equation}
For $rk$, substitute $b=dk-a$ into $pk=4hab-1$ and cancel to obtain
$(4had-p)k=4ha^2+1$. The two residue classes modulo four follow
from $p\equiv1\pmod4$ and the product $rk\equiv1\pmod4$.
No parameter has been replaced by a unit representative.
The two still possible residue events have exact tests
\begin{equation}
 \begin{aligned}
 x\equiv y\pmod p
   &\Longleftrightarrow p\mid(dk-2a),
       &&\Longrightarrow p\mid(4ha^2-1),\\
 p\mid S
   &\Longleftrightarrow p\mid k,
       &&\Longrightarrow p\mid(4ha^2+1).
 \end{aligned}                                         \tag{PES10}
\end{equation}
The equivalences use $p\nmid hd$ and $S\equiv x+y\equiv hd^2k$.
For the first implication, $a\equiv b$ and $4hab-1=pk$ give
$4ha^2-1=0\pmod p$. The second follows from $rk=4ha^2+1$.
These are exact necessary conditions and a finite-search parametrization;
neither event is excluded merely because a bounded search finds none.

\subsection{Every inclusion entry and the least coefficient power}
Let $c=v_p(C_0)$ and keep the complete original ideal basis $G_I$ of
ISR12. From $VQ=\operatorname{diag}(f'(r_i))$, the full eight-column
factorization of the scalar ideal inclusion is
\begin{equation}
 Z_j=G_I^{-1}C_0^j I_8
   =\operatorname{diag}\left(
       \operatorname{diag}\left(\frac{C_0^j}{p^{m_i}D_i}\right)V,
       \operatorname{diag}\left(\frac{C_0^j}{D_i}\right)V\right),
 \quad G_IZ_j=C_0^jI_8.                                \tag{PES11}
\end{equation}
Thus its complete even block entry $(i,n)$ is
$C_0^j r_i^n/(p^{m_i}D_i)$ and its complete odd block entry is
$C_0^j r_i^n/D_i$, for $0\le i,n\le3$. Because each original root
is integral and $m_i\ge0$, all entries are integral if
$jc\ge\max_i(m_i+b_i)$. Conversely the even entries with $n=0$
force exactly these four inequalities. Therefore the least power
and both unchanged receiving identities are
\begin{equation}
 \min\{j\ge0:C_0^jE\subset I\}
  =\begin{cases}1&\text{one divisible denominator},\\
                 2&\text{two divisible denominators},\end{cases}
 \qquad (F_U G_I)Z_j=C_0^jF_U,\quad
        (F_W G_I)Z_j=C_0^jF_W.                          \tag{PES12}
\end{equation}
Here $c=1$ and the largest $m_i+b_i$ equals one in the first case;
$c=2$ and the largest equals three in the second. The matrices
$F_U,F_W$ retain all original moments, centres and Gamma masses
from ISR32--34. This proves every assertion of PB6--8 through the
entire original coefficient matrices, including necessity of the
least exponents, not just an index comparison.

\paragraph{Provenance.}
The ES relation, actual-witness marking, and finite signed algebra
retain the human-source citations and original programme definitions
of ISR1 and PS1. PES1--12 are the independent positivity and parameter
derivations written above. No new external literature is claimed read,
and no Frobenius or RH conclusion is assigned to this finite calculation.

\section{Scope, source, and original objects}
This proof derives the gluing and collision assertions (32)--(46) of the
received account \cite{DefiningPrimeReceived}, retaining the original odd generator
$\sigma$. The coefficient change to $\zeta$ below is accompanied by its full
inverse, residue cofactor, connection residue, and both original physical
receivers. The original ES labels are not merged when their coordinates
coincide. The marked prime is denoted $P$; the polynomial variable local to
it is $x=T-P$; the collision parameter is $t$. None is the fixed conductor
amplification $a_{\rm amp}$. The inherited polynomial provenance is the
Alp\"oge--Fable example as recorded by Tao \cite{Tao}; the four formal
$\mathrm M,\mathrm N,\mathrm T,\mathrm E$ labels and unchanged coordinate
marking are from the specified ES source \cite{ES488}. The Gamma orthogonality
is the human source \cite{DLMF18Source}; the original programme's FC20--24,
RS21--25 and GD14 specify its exact use \cite{Fable050}. Classical
residue/dual-form provenance is \cite[Section 2]{JouveRV}. All finite
algebra and metric formulas required here are proved below.

Let the coefficient base be a characteristic-zero domain with $A,P$ units,
and let
\begin{equation}
 g(T)=T^3-s_1T^2+wT-s_3,\quad w=4s_3/P,\quad
 F(T)=A(T-P)g(T),\quad R=g(P).
 \tag{ECC1}
\end{equation}
Work where $R$ is not identically zero. The signed algebra is
$\mathscr E=B[T,\sigma]/((T-P)g,\sigma^2-F')$, in the unchanged even-then-odd
basis $(1,T,T^2,T^3,\sigma,\sigma T,\sigma T^2,\sigma T^3)$.
Its quotients are
\begin{equation}
 \mathscr E_P=B[\sigma]/(\sigma^2-AR),\qquad
 \mathscr E_g=B[T,\sigma]/(g,\sigma^2-A(T-P)g').
 \tag{ECC2}
\end{equation}
Division by the monic equations gives ranks $8,2,6$ respectively.

\section{Full gluing and the inverse CRT map}
The exact maps, including their target at the divisor, are
\begin{equation}
 0\longrightarrow\mathscr E\xrightarrow{j}\mathscr E_P\oplus\mathscr E_g
 \xrightarrow{d} B/(R)[\sigma]/(\sigma^2)\longrightarrow0.
 \tag{ECC3}
\end{equation}
Here $j$ restricts both coefficient polynomials and $d$ subtracts their
restrictions at $T=P$ modulo $R$. To prove injectivity, a polynomial
divisible by both $g$ and $T-P$ has form $gq$ with $R q(P)=0$.
The nonzerodivisor assumption on $R$ gives $q(P)=0$, hence divisibility by
$(T-P)g$. This applies to each even and odd coefficient. Their common
quotient has $(T-P,g)=(T-P,R)$ and $F'=AR=0$ there, giving the displayed
overlap including its odd nilpotent. For a compatible pair $(u,b(T))$,
divide $u-b(P)$ by $R$ and add that multiple of $g$ to $b$; this gives its
preimage. Every overlap element lifts from the prime factor, proving
surjectivity. This proves exactness and the fibre-product multiplication.

In the coefficient order $(u,b_0,b_1,b_2)$ the even comparison is
\begin{equation}
 \begin{gathered}
 C_R=\begin{pmatrix}1&P&P^2&P^3\\1&0&0&s_3\\
 0&1&0&-w\\0&0&1&s_1\end{pmatrix},\quad
 \det C_R=-R,\\
 C_R^{-1}=\begin{pmatrix}0&1&0&0\\0&0&1&0\\0&0&0&1\\0&0&0&0\end{pmatrix}
 +\frac1R\begin{pmatrix}-s_3\\w\\-s_1\\1\end{pmatrix}
             \begin{pmatrix}1&-1&-P&-P^2\end{pmatrix}.
 \end{gathered}
 \tag{ECC4}
\end{equation}
Indeed the inverse polynomial is
$b(T)+[u-b(P)]g(T)/R$. Evaluation and reduction prove both matrix products
equal $I_4$. The full signed matrix is $\diag(C_R,C_R)$, with determinant
$R^2$, full inverse $\diag(C_R^{-1},C_R^{-1})$, and two copies of the
displayed overlap. No sign state was deleted.

Put $\Delta_g=\operatorname{disc}(g)$. Then
\begin{equation}
 \det\Tr_{\mathscr E}=256A^4R^6\Delta_g^3,\qquad
 \det\Tr_{\mathscr E_P\oplus\mathscr E_g}=256A^4R^2\Delta_g^3.
 \tag{ECC5}
\end{equation}
For a monic quartic $f$, the ordinary even trace matrix has determinant
$\operatorname{disc}(f)$. In the signed algebra the trace matrix is block
diagonal with blocks $2T_f$ and $2T_fM_{Af'}$. Its determinant is
$2^8\operatorname{disc}(f)^2\prod_i Af'(r_i)
=256 A^4\operatorname{disc}(f)^3$, first on distinct roots and then
identically by polynomial continuation. Since
$\operatorname{disc}((T-P)g)=R^2\Delta_g$, the first formula follows.
The trace matrices under the inclusion differ by $C_R^{\oplus2}$ on
both sides, proving the second. Here ``trace'' is multiplication trace,
including on nonreduced fibres. These determinant identities by themselves
do not prove normality for an arbitrary base; the local normality statement
below includes a regular one-parameter base and a transverse collision.

\section{Original signed collision, exact forward and inverse maps}
Take a sufficiently small complex disc with coordinate $t$, and retain
analytic $P,A$ and the two remaining roots, such that $A$ never vanishes
and the latter roots remain distinct from one another, $P$, and $P+t$.
All formulas allow these units to depend analytically on $t$.
Write $h(P+x)=h_0+h_1x+x^2$, $h_t=h(P+t)$ and
\begin{equation}
 F=A x(x-t)h(P+x),\quad x^2=tx,\quad
 \sigma^2=A\{x(x-t)h'(P+x)+(2x-t)h(P+x)\}
          =A(2x-t)h(P+x).
 \tag{ECC6}
\end{equation}
The equality on the right is only in the displayed root quotient;
the full derivative on its left remains the original derivative.
Choose a single analytic square root of the nonzero function $Ah$ near
this disc and collision. Denote its values by
$q_0^2=Ah_0$, $q_t^2=Ah_t$, with $q_t=q_0$ at $t=0$.
The two values cannot be assigned unrelated signs if an analytic return
through $t=0$ is intended. The common sign remains a retained choice.
The unit in the root quotient and its inverse are exactly
\begin{equation}
 q(x)=q_0+\frac{q_t-q_0}{t}x,\qquad
 q(x)^{-1}=q_0^{-1}+\frac{q_t^{-1}-q_0^{-1}}t x.
 \tag{ECC7}
\end{equation}
The divided differences extend analytically. Multiplication followed by
$x^2=tx$ proves the inverse. Squaring and evaluating at $0,t$ proves
$q(x)^2=Ah(P+x)$; the same identity at zero follows by analytic continuity.

Define $\zeta=q(x)^{-1}\sigma$. Both directions of the isomorphism are
\begin{equation}
 \zeta^4=t^2,\quad x=(\zeta^2+t)/2,\quad
 \sigma=a\zeta+b\zeta^3,\quad
 a=(q_0+q_t)/2,\quad b=(q_t-q_0)/(2t),\quad
 \zeta=q(x)^{-1}\sigma.
 \tag{ECC8}
\end{equation}
Indeed $\zeta^2=2x-t$ and $x(x-t)=0$ are equivalent to the first two
relations. Substituting $q(x)$ proves the formula for $\sigma$, and
the inverse returns $x$ and $\sigma$ exactly. In the original basis
$e=(1,x,\sigma,x\sigma)$ and $z=(1,\zeta,\zeta^2,\zeta^3)$, let columns
transform by $[f]_z=S[f]_e$. Then
\begin{equation}
 S=\begin{pmatrix}1&t/2&0&0\\0&0&a&tq_t/2\\
 0&1/2&0&0\\0&0&b&q_t/2\end{pmatrix},\quad
 S^{-1}=\begin{pmatrix}
 1&0&-t&0\\0&0&2&0\\
 0&q_0^{-1}&0&-tq_0^{-1}\\
 0&(q_t^{-1}-q_0^{-1})/t&0&q_t^{-1}+q_0^{-1}
 \end{pmatrix},\quad \det S=-q_0q_t/4.
 \tag{ECC9}
\end{equation}
The last column of $S$ uses $x\sigma=(tq_t\zeta+q_t\zeta^3)/2$.
Each column of the inverse is ECC7 and
$\zeta^3=(2x-t)q(x)^{-1}\sigma$; this proves both inverse products and
invertibility at $t=0$. In particular, with $q=q_0|_{t=0}$ and
$\eta=h_1/h_0|_{t=0}$,
\begin{equation}
 \sigma=q\zeta+(q\eta/4)\zeta^3\quad(t=0).
 \tag{ECC10}
\end{equation}
Deleting the cubic term does not preserve the original odd coordinate.

The exact embedding of this local block in the full quartic algebra is
also retained. Put
$k_1=(h_t^{-1}-h_0^{-1})/t=-(h_1+t)/(h_0h_t)$ and
\begin{equation}
 e_c=h(P+x)(h_0^{-1}+k_1x),\quad e_h=1-e_c,\quad
 I_t=\begin{pmatrix}
 1&0\\h_1/h_0+k_1h_0&h_0/h_t\\
 1/h_0+k_1h_1&h_1/h_t\\k_1&1/h_t
 \end{pmatrix}.
 \tag{ECC11}
\end{equation}
Modulo $x(x-t)$, $e_c=1$; modulo $h$, it is zero. Hence $e_c^2=e_c$,
$e_h^2=e_h$, and $e_ce_h=0$ modulo the full quartic. The first column of
$I_t$ is $e_c$ and the second is $xe_c=xh/h_t$ modulo that quartic.
These claims follow either by division or by multiplication with
$x^4+(h_1-t)x^3+(h_0-th_1)x^2-th_0x=0$.
Consequently $\diag(I_t,I_t)S^{-1}$ is the entire four-to-eight original
coefficient embedding in $x$ coordinates, and translation $x=T-P$ uses
the full cubic binomial matrix, not a change of the original centres.
The complementary block $e_h\mathscr E$ has rank four and is retained.

\section{Normalization, specialization, and original connection residue}
Over the analytic one-dimensional base $B=\C\{t\}$, the normalization of
$B[\zeta]/(\zeta^4-t^2)$ is the sum of
$B[\zeta]/(\zeta^2+t)$ and $B[\zeta]/(\zeta^2-t)$.
Each summand is a regular local ring, isomorphic to $\C\{\zeta\}$ with
$t=\mp\zeta^2$; it is therefore integrally closed. The product is finite,
contains the original algebra and equals its total quotient algebra
after $t$ is inverted. Every integral element restricts to an integral
element of each normal factor, which proves the assertion.
In the order $(1,\zeta)$ at $x=0$ followed by $(1,\zeta)$ at $x=t$,
\begin{equation}
 C_\zeta=\begin{pmatrix}1&0&-t&0\\0&1&0&-t\\
 1&0&t&0\\0&1&0&t\end{pmatrix},\quad\det C_\zeta=4t^2,
 \qquad
 C_\sigma=\begin{pmatrix}1&0&0&0\\0&0&1&0\\
 1&t&0&0\\0&0&1&t\end{pmatrix},\quad\det C_\sigma=-t^2,
 \tag{ECC12}
\end{equation}
where $C_\sigma$ uses the original $(1,\sigma)$ in each factor.
They obey $C_\zeta S=\diag(1,q_0,1,q_t)C_\sigma$.
Their inverses send $(u_0,v_0,u_t,v_t)$ respectively to
\begin{equation}
 \left((u_0+u_t)/2,(v_0+v_t)/2,(u_t-u_0)/(2t),(v_t-v_0)/(2t)\right),
 \quad
 \left(u_0,(u_t-u_0)/t,v_0,(v_t-v_0)/t\right).
 \tag{ECC13}
\end{equation}
Elementary row subtraction and division by the units $2,q_0,q_t$
give Smith factors $(1,1,t,t)$, with cokernel $(B/(t))^2$.
Tensoring the short exact sequence with $B/(t)$ gives
\begin{equation}
 0\to\C^2\to\C[\zeta]/(\zeta^4)
 \to\C[\zeta]/(\zeta^2)\oplus\C[\zeta]/(\zeta^2)
 \to\C^2\to0.
 \tag{ECC14}
\end{equation}
The middle arrow sends a polynomial to its two remainders modulo
$\zeta^2$. Its kernel is precisely $(\zeta^2,\zeta^3)$; hence the initial
$\operatorname{Tor}_1$ injection has exactly this image. Its cokernel is
the difference of the two constant-linear coefficient pairs. Thus all
four original directions and the two lost specialization directions
are explicitly accounted for.

On the punctured disc differentiate the four labelled evaluation
branches. Since $t\partial_t\zeta^j=(j/2)\zeta^j$, the connection matrix
in $z$ is $D/t$, $D=\diag(0,1/2,1,3/2)$. In original coordinates it is
\begin{equation}
 \Omega_e=S^{-1}(DS/t+S'),\qquad
 \Res_{t=0}\Omega_e=
 \begin{pmatrix}0&0&0&0\\0&1&0&0\\
 0&0&1/2&0\\0&0&\eta/2&3/2\end{pmatrix}.
 \tag{ECC15}
\end{equation}
This is the derivative of $zS$ column by column. The residue is
$S(0)^{-1}DS(0)$ because $S'$ is analytic. Substitution of ECC9 proves
every entry, including the off-diagonal $\eta/2$. The normalized extension
has residue $\diag(0,1/2,0,1/2)$, so the traces differ by $2$, equal to
$t\partial_t\log\det C_\zeta$. Original $\sigma$ in its two factors
only adds analytic logarithmic derivatives of $q_0,q_t$, leaving these
residue eigenvalues unchanged. A loop in $t$ changes the signs over both
$P$ and $P+t$ without exchanging those labelled roots. This proves the
local monodromy assertion. ECC46--48 below supplies the further global
argument for the full normalized ES parameter space.

\section{Full original residue and trace inverse}
Use the original programme residue
$\Lambda(R+\sigma B)=\sum\Res B(T)dT/F(T)$, restricted to the two-root
block. For $B=b_0+b_1x$, at $t\ne0$ its exact values and continuation are
\begin{equation}
 \Lambda(\sigma B)=\frac{B(t)}{tq_t^2}-\frac{B(0)}{tq_0^2},\qquad
 \Lambda(\sigma)=\frac{q_t^{-2}-q_0^{-2}}t,\quad
 \Lambda(x\sigma)=q_t^{-2},\quad
 \Lambda(x^2\sigma)=tq_t^{-2}.
 \tag{ECC16}
\end{equation}
The formula follows from the two simple-root residues and extends by
the analytic divided difference; at zero it equals the double-root
residue $[x]B/(Ah(P+x))$. The full pairing in $e$ is
\begin{equation}
 \begin{gathered}
 B_e=\begin{pmatrix}0&J_e\\J_e&0\end{pmatrix},\quad
 J_e=\begin{pmatrix}(q_t^{-2}-q_0^{-2})/t&q_t^{-2}\\
 q_t^{-2}&tq_t^{-2}\end{pmatrix},\\
 \det J_e=-1/(q_0^2q_t^2),\quad
 J_e^{-1}=\begin{pmatrix}-tq_0^2&q_0^2\\
 q_0^2&(q_t^2-q_0^2)/t\end{pmatrix}.
 \end{gathered}
 \tag{ECC17}
\end{equation}
Odd-odd products are even and have zero residue; even-even products
also have zero residue. ECC16 proves the cross block, and direct
multiplication proves the displayed inverse. The residue pairing remains
perfect at the collision.

In $z$ the only nonzero values on its basis are
\begin{equation}
 \ell_1=\Lambda(\zeta)=\frac{q_t^{-3}-q_0^{-3}}t,\qquad
 \ell_3=\Lambda(\zeta^3)=q_t^{-3}+q_0^{-3}.
 \tag{ECC18}
\end{equation}
Indeed $\zeta=q^{-1}\sigma$, while $\zeta^3=(2x-t)q^{-1}\sigma$;
insert these two expressions into ECC16. With the monic transfer
functional $\lambda_z(f)=[\zeta^3](f\bmod(\zeta^4-t^2))$, the precise
residue cofactor and its inverse are
\begin{equation}
 \Lambda(f)=\lambda_z(cf),\quad
 c=\ell_3+\ell_1\zeta^2,\quad
 c^{-1}=\frac{q_0^3q_t^3}{4}(\ell_3-\ell_1\zeta^2),\quad
 c|_{t=0}=\frac2{q^3}\left(1-\frac{3\eta}{4}\zeta^2\right).
 \tag{ECC19}
\end{equation}
Here $\ell_3^2-t^2\ell_1^2=4/(q_0^3q_t^3)$ proves the inverse. Both
functionals vanish on even powers, so testing the two odd basis powers
proves the functional identity completely. The limit uses
$\partial_x q/q=h_1/(2h_0)$. The full Gram in $z$ is
\begin{equation}
 B_z=\begin{pmatrix}0&\ell_1&0&\ell_3\\
 \ell_1&0&\ell_3&0\\0&\ell_3&0&t^2\ell_1\\
 \ell_3&0&t^2\ell_1&0\end{pmatrix},\quad
 \det B_z=16/(q_0^6q_t^6),\quad B_e=S^{\mathsf T}B_zS.
 \tag{ECC20}
\end{equation}
This also checks the constant and sign in $\det S$ against ECC17.

Put $a_0=q_0^2$, $a_t=q_t^2$, and use the original even order $(1,x)$.
The full odd multiplication, the trace multiplier, and their inverses are
\begin{equation}
 \begin{gathered}
 H=\begin{pmatrix}-ta_0&0\\a_0+a_t&ta_t\end{pmatrix},\qquad
 M_\sigma=\begin{pmatrix}0&H\\I_2&0\end{pmatrix},\quad
 M_\sigma^{-1}=\begin{pmatrix}0&I_2\\H^{-1}&0\end{pmatrix},\\
 H^{-1}=\begin{pmatrix}-1/(ta_0)&0\\
 (a_0+a_t)/(t^2a_0a_t)&1/(ta_t)\end{pmatrix},\\
 H^{-2}=\begin{pmatrix}1/(t^2a_0^2)&0\\
 (a_0+a_t)(a_0-a_t)/(t^3a_0^2a_t^2)&1/(t^2a_t^2)\end{pmatrix},\\
 \Theta=M_{2\sigma^3}=2\begin{pmatrix}0&H^2\\H&0\end{pmatrix},\qquad
 \Theta^{-1}=\frac12\begin{pmatrix}0&H^{-1}\\H^{-2}&0\end{pmatrix},\quad
 \det\Theta=-16t^6a_0^3a_t^3.
 \end{gathered}\tag{ECC21}
\end{equation}
The original relation ECC6 gives $\sigma^2=-ta_0+(a_0+a_t)x$,
which proves $H$. The inverses follow by triangular multiplication;
the determinant follows from $\det H=-t^2a_0a_t$ and the even block swap.
Multiplication trace obeys $\Tr M_f=\Lambda(2\sigma^3f)$: it is the sum
of all four values when $t\ne0$, and both sides extend analytically.
Thus $B_e\Theta$ is the full trace pairing, not the perfect residue
pairing. At $t=0$, with $a=q^2$, it has rank one, while $B_e$ has rank four.

\section{Both original Gamma receivers, with their full finite matrices}
For any scalar $c$ define
\begin{equation}
 C_c=\begin{pmatrix}1&c&c^2&c^3\\0&1&2c&3c^2\\
 0&0&1&3c\\0&0&0&1\end{pmatrix},\quad
 [f(r+c)]=C_c[f(r)],\quad C_c^{-1}=C_{-c}.
 \tag{ECC22}
\end{equation}
Binomial expansion proves all entries and both inverse products.
Hold the original conductor, its period, actual order $v$, and moments
$\mu_v,\ldots,\mu_{v+3}$ fixed, with $\mu_v\ne0$. Put
$c_*=a_{\rm amp}/2-4$, $\kappa=c_*-1/2$, $b=s/2>0$ and
$M=(2\pi)^{s/2}$ as in the original Gamma measure. Its actual source and
target are $U=\mathcal P_{v+3}/\mathcal P_{v-1}$ and $W=\mathcal P_3$.
Their conductor is the given $T_A$ (the subscript here is the original
conductor notation, not multiplication by the quartic coefficient).
The scalar halves of GD14, on coefficient columns in $r$, are
\begin{equation}
 F_W f=f(S'-\kappa),\quad F_U f=T_A^{-1}f(S'-\kappa),\quad
 F_W=T_AF_U,\quad F_W^{-1}w=w(r+\kappa),\quad
 F_U^{-1}u=(T_Au)(r+\kappa).
 \tag{ECC23}
\end{equation}
Their signed receivers are $F_W^{\oplus2},F_U^{\oplus2}$, not an
identification of multiplication by $r$ with the original conductor.

Here are finite explicit original Gram matrices and their inverses.
Use the unchanged shifted moments and phases
\begin{equation}
 \begin{gathered}
 \nu_{v+j}=\sum_{k=0}^{j}\binom{v+j}{v+k}(-4)^{j-k}\mu_{v+k},\\
 g_0=(-i)^v\mu_v/v!,\quad g_1=(-i)^{v+1}\nu_{v+1}/(v+1)!,\\
 g_2=(-i)^{v+2}\nu_{v+2}/(v+2)!-(v/3)g_0,\quad
 g_3=(-i)^{v+3}\nu_{v+3}/(v+3)!-((v+1)/3)g_1,\\
 \xi_0=g_0^{-1},\quad\xi_1=-g_1/g_0^2,\quad
 \xi_2=g_1^2/g_0^3-g_2/g_0^2,\\
 \xi_3=-g_1^3/g_0^4+2g_1g_2/g_0^3-g_3/g_0^2,\\
 k_j=(b)_{v+j}/(v+j)!,\quad D_i=\diag(1,i,-1,-i).
 \end{gathered}\tag{ECC24}
\end{equation}
The cubic products $(\sum g_jz^j)(\sum\xi_jz^j)=1\bmod z^4$
prove every inverse coefficient; the shifted moment formula is the
coefficient expansion of $e^{-4z}E_A(z)$. Define
\begin{equation}
 \begin{gathered}
 V=\begin{pmatrix}\sqrt M&0&b\sqrt M&0\\
 0&\sqrt{Mb}&0&(3b+2)\sqrt{Mb}\\
 0&0&\sqrt{2Mb(b+1)}&0\\0&0&0&\sqrt{6Mb(b+1)(b+2)}\end{pmatrix},\\
 F=\begin{pmatrix}\xi_0&\xi_1&b\xi_0+2\xi_2&(3b+2)\xi_1+6\xi_3\\
 0&\xi_0&2\xi_1&(3b+2)\xi_0+6\xi_2\\
 0&0&2\xi_0&6\xi_1\\0&0&0&6\xi_0\end{pmatrix},\\
 B_F=\begin{pmatrix}g_0&g_1&g_2-bg_0/2&g_3-bg_1/2\\
 0&g_0&g_1&g_2-(3b+2)g_0/6\\0&0&g_0/2&g_1/2\\0&0&0&g_0/6\end{pmatrix},\\
 W_x=VD_i,\qquad U_x=\sqrt M\diag(\sqrt{k_j})FD_i,\qquad
 U_x^{-1}=M^{-1/2}D_i^*B_F\diag(k_j^{-1/2}).
 \end{gathered}\tag{ECC25}
\end{equation}
The target orthonormal polynomials are $p_j((S'-c_*)/i)/\sqrt{Mj!(b)_j}$,
where $p_0=1,p_1=y,p_2=y^2-b,p_3=y^3-(3b+2)y$.
Expanding their monomials gives $V$, and their source inverse images
under FC24 give $U_x$. More explicitly that inverse in orthonormal
coordinates has entries $\xi_{j-r}\rho_j/\rho_{v+r}$ for $r\le j$, where
$\rho_j=\sqrt{j!/(b)_j}$; multiplying $VD_i$ gives the displayed $U_x$.
The convolution identity in ECC24 proves $FB_F=B_FF=I_4$.
This supplies a direct finite check, with all original masses present.

In $r$ coefficients the actual two metrics and complete covariance
factors are
\begin{equation}
 \begin{gathered}
 H_U=C_{1/2}^*U_x^*U_xC_{1/2},\quad
 H_W=C_{1/2}^*W_x^*W_xC_{1/2},\\
 Q_U=C_{-1/2}U_x^{-1},\quad Q_W=C_{-1/2}D_i^*V^{-1},\quad
 H_X^{-1}=Q_XQ_X^*\quad(X=U,W).
 \end{gathered}\tag{ECC26}
\end{equation}
For specificity $V^{-1}$ has diagonal
$M^{-1/2},(Mb)^{-1/2},[2Mb(b+1)]^{-1/2},[6Mb(b+1)(b+2)]^{-1/2}$,
entry $(0,2)$ equal $-b[2Mb(b+1)]^{-1/2}$, entry $(1,3)$ equal
$-(3b+2)[6Mb(b+1)(b+2)]^{-1/2}$, and all other off-diagonal entries zero.
Direct multiplication proves this inverse. ECC26 follows because
$r=S'-\kappa=(S'-c_*)+1/2$. Every displayed matrix is invertible, so
both metrics are strictly positive. No original moment has been assigned
a substitute value.

The received coefficient receiver $\Psi_*$ is different from literal
substitution. Retain its exact definition
$\Phi_*=(V_*^{\mathsf T})^{-1}K^{-1}$ and
$\Psi_*=B_{v,*}(V_*^{\mathsf T})^{-1}K^{-1}=T_A\Phi_*$ in the inherited
coordinate bases. Its exact correction on $r$ columns is
\begin{equation}
 M_* = C_\kappa\Psi_*,\quad
 M_*^{-1}=K V_*^{\mathsf T} B_{v,*}^{-1}C_{-\kappa},\quad
 F_WM_* = \Psi_*,\quad F_UM_* = \Phi_*,\quad
 H_X^{\rm rec}=M_*^*H_XM_*.
 \tag{ECC27}
\end{equation}
The fixed matrices $K,V_*,B_{v,*}$ here are precisely the original
ES marking, reference evaluation and conductor matrices of FC26--32,
not matrices recomputed from the colliding family. Multiplying ECC23
by $M_*$ proves each identity; their inverses prove equivalence of the
two receivers. Consequently the received covariance is obtained by
replacing $Q_X$ with $M_*^{-1}Q_X$. On physical spaces the exact maps
$L_W=\Psi_*F_W^{-1}$ and $L_U=\Phi_*F_U^{-1}$ obey
$T_AL_U=L_WT_A$ by multiplication. All results below apply separately
to $H_U,H_W$ and, with this correction, their received versions.

\section{The complete minimum and all raw signed eigenvalues}
Fix one of the positive matrices $H=H_X$ from ECC26, or the explicitly
corrected received matrix. Let $v(r)=(1,r,r^2,r^3)^{\mathsf T}$ and
\begin{equation}
 q_j(r)=v(r)^{\mathsf T}Qe_j,\qquad
 K(r,s)=\sum_{j=0}^3q_j(r)\overline{q_j(s)}
       =v(r)^{\mathsf T}H^{-1}\overline{v(s)}.
 \tag{ECC28}
\end{equation}
Here $Q$ is explicitly ECC26, or $M_*^{-1}Q_X$. Thus this is a finite
formula in the original moments and masses, not an unspecified inverse.
The signed norm of $(f_0,f_1)$ is $[f_0]^*H[f_0]+[f_1]^*H[f_1]$.
For the observation $(f_0(P),f_1(P))=(u,v)$ the complete minimum is
\begin{equation}
 \frac{|u|^2+|v|^2}{K(P,P)},\quad
 f_0(r)=uK(r,P)/K(P,P),\quad f_1(r)=vK(r,P)/K(P,P).
 \tag{ECC29}
\end{equation}
Indeed $H^{-1}\overline{v(P)}$ is the representing vector of evaluation;
its squared norm is $K(P,P)>0$. Every other representative is this
one plus a vector in the entire evaluation kernel, orthogonal to it.
Pythagoras proves attainment, uniqueness, and the formula separately
for each coefficient. It includes the complementary quartic directions.
In particular the CRT section $g/R$ has the exact excess
\begin{equation}
 \|g/R\|_H^2=\frac1{K(P,P)}+
 \left\|g/R-K(\,\cdot,P)/K(P,P)\right\|_H^2,\qquad
 \frac{\det G_{\rm CRT}}{\det G_{\min}}=
 \left(\frac{K(P,P)\|g\|_H^2}{|R|^2}\right)^2.
 \tag{ECC30}
\end{equation}
The two signed section columns occupy the separate coefficient copies,
so their Gram is $\|g/R\|_H^2I_2$ and the minimum Gram is
$K(P,P)^{-1}I_2$. This proves the determinant ratio with no removal of
kernel directions. The minimum stays positive at the collision although
the CRT section has the displayed pole.

For $t\ne0$ let $s_0^2=-tq_0^2$, $s_t^2=tq_t^2$, retaining both choices
of $s_0,s_t$. Order the four raw values as $(P,+),(P,-),(P+t,+),(P+t,-)$.
Their observation matrix on the eight original coefficients is
\begin{equation}
 O_t=\begin{pmatrix}
 v(P)^{\mathsf T}&s_0v(P)^{\mathsf T}\\
 v(P)^{\mathsf T}&-s_0v(P)^{\mathsf T}\\
 v(P+t)^{\mathsf T}&s_tv(P+t)^{\mathsf T}\\
 v(P+t)^{\mathsf T}&-s_tv(P+t)^{\mathsf T}
 \end{pmatrix},\quad
 C_t=O_t\diag(H^{-1},H^{-1})O_t^*.
 \tag{ECC31}
\end{equation}
It has rank four: invert each signed sum/difference and then interpolate
the two distinct root values separately. A unitary signed sum/difference,
followed by ordering the two even observations first, gives the full
covariance blocks $2K_2$ and $2D_sK_2D_s^*$, where
\begin{equation}
 K_2=\begin{pmatrix}K_0&C\\\overline C&K_t\end{pmatrix},\quad
 K_0=K(P,P),\ K_t=K(P+t,P+t),\ C=K(P,P+t),\quad
 D_s=\diag(s_0,s_t).
 \tag{ECC32}
\end{equation}
Consequently all four finite covariance eigenvalues, with no asymptotic
replacement, are the sorted union
\begin{equation}
 \begin{aligned}
 \lambda_{E,\pm}&=K_0+K_t\ \pm\sqrt{(K_0-K_t)^2+4|C|^2},\\
 \lambda_{O,\pm}&=a+b\ \pm\sqrt{(a-b)^2+4|s_0s_t|^2|C|^2},\qquad
 a=|s_0|^2K_0,\quad b=|s_t|^2K_t.
 \end{aligned}\tag{ECC33}
\end{equation}
These follow from the quadratic characteristic polynomials of the two
Hermitian blocks. Positivity follows from rank two of the two evaluation
rows and positivity of $H$. In particular
$\lambda_{E,+}\lambda_{E,-}=4\det K_2$ and
$\lambda_{O,+}\lambda_{O,-}=4|s_0s_t|^2\det K_2$.

Take $P$ fixed in the limit, or use its analytic limiting value. Write
\begin{equation}
 K=\sum_j|q_j(P)|^2,\quad
 J=\begin{pmatrix}
 K&\sum_jq_j(P)\overline{q_j'(P)}\\
 \sum_jq_j'(P)\overline{q_j(P)}&\sum_j|q_j'(P)|^2
 \end{pmatrix},\quad L=\det J/K>0,\quad H_0=|Ah_0|_{t=0}.
 \tag{ECC34}
\end{equation}
The two rows $v(P)^{\mathsf T}Q$ and $v'(P)^{\mathsf T}Q$ are independent,
because $Q$ is invertible and evaluation/derivative are independent on
cubics. Thus $J$ is positive definite. Expanding its determinant gives
$KL=\sum_{i<j}|q_i(P)q_j'(P)-q_j(P)q_i'(P)|^2$, proving the exact
positive constant. Divided differences give
$\det K_2=|t|^2(KL+o(1))$, while $K_t,C\to K$ and
$|s_0|^2/|t|,|s_t|^2/|t|\to H_0$. Therefore, in decreasing order,
\begin{equation}
 \lambda_1\to4K,\quad
 \lambda_2/|t|\to4H_0K,\quad
 \lambda_3/|t|^2\to L,\quad
 \lambda_4/|t|^3\to H_0L.
 \tag{ECC35}
\end{equation}
For the smaller eigenvalue of each block use its exact determinant
divided by its larger eigenvalue in ECC33. This proves all four limits
also for arbitrary complex approach to zero; no phase is suppressed in
the exact matrices. The four different powers force the stated order
for all sufficiently small nonzero $t$.

The attained right inverse, its metric, and every exterior cost are
\begin{equation}
 \begin{gathered}
 R_t=\diag(H^{-1},H^{-1})O_t^*C_t^{-1},\quad
 O_tR_t=I_4,\\ R_t^*\diag(H,H)R_t=C_t^{-1},\quad
 \|\wedge^jR_t\|=\prod_{k=5-j}^4\lambda_k^{-1/2}\quad(1\le j\le4).
 \end{gathered}
 \tag{ECC36}
\end{equation}
Every other solution of the raw values is $R_tu+\ker O_t$, orthogonal
in the original metric, as multiplication of the displayed matrices
shows. Thus this is the full original quotient inverse, not the inverse
of a selected algebraic section. Its singular values are the square
roots of the eigenvalues of its Gram $C_t^{-1}$; singular-value
decomposition proves the exterior product formula. In particular
\begin{equation}
 \begin{aligned}
 |t|^{3/2}\|R_t\|&\longrightarrow(H_0L)^{-1/2},\\
 |t|^{5/2}\|\wedge^2R_t\|&\longrightarrow(H_0^{1/2}L)^{-1},\\
 |t|^3\|\wedge^3R_t\|&\longrightarrow(2H_0L\sqrt K)^{-1},\\
 |t|^3\|\wedge^4R_t\|&\longrightarrow(4H_0KL)^{-1},\qquad
 \|\wedge^4R_t\|=\frac1{4|s_0s_t|\det K_2}.
 \end{aligned}\tag{ECC37}
\end{equation}
The determinant expression follows directly from the two covariance
blocks. ECC33 and ECC36 are exact at every admissible nonzero parameter,
so the bounds are not limited to these leading terms.

\section{Retained jet map and original quotient trace cost}
The exact raw-to-jet matrix, preserving original $\sigma$ values, is
\begin{equation}
 Z_t=\begin{pmatrix}
 1/2&1/2&0&0\\-1/(2t)&-1/(2t)&1/(2t)&1/(2t)\\
 1/(2s_0)&-1/(2s_0)&0&0\\
 -1/(2ts_0)&1/(2ts_0)&1/(2ts_t)&-1/(2ts_t)
 \end{pmatrix},\quad
 |\det Z_t|=\frac1{4|s_0s_tt^2|}.
 \tag{ECC38}
\end{equation}
It sends the four values to
$(f_0(P),(f_0(P+t)-f_0(P))/t,f_1(P),(f_1(P+t)-f_1(P))/t)$.
Its inverse sends $(a,b,c,d)$ to
$(a+s_0c,a-s_0c,a+tb+s_t(c+td),a+tb-s_t(c+td))$.
These four identities prove both inverse products and the determinant.
Let
\begin{equation}
 E_t=\begin{pmatrix}v(P)^{\mathsf T}\\
 [v(P+t)^{\mathsf T}-v(P)^{\mathsf T}]/t\end{pmatrix},\quad
 J_t=E_tH^{-1}E_t^*,\quad G_t=J_t^{-1}.
 \tag{ECC39}
\end{equation}
Then $Z_tO_t=\diag(E_t,E_t)$ and
$Z_tC_tZ_t^*=\diag(J_t,J_t)$. At zero, $E_0$ has rows
$v(P)^{\mathsf T},v'(P)^{\mathsf T}$, so $J_t\to J>0$ and
$G_t\to G=J^{-1}$. The full minimum metric on the original local
coefficient algebra $(1,x,\sigma,x\sigma)$ is
$\diag(G_t,G_t)$. This follows from the same representing-vector
calculation as ECC36 with $E_t$ in place of $O_t$ and includes the entire
kernel of the quartic-to-local quotient. It is different from assigning
a Euclidean metric to local coordinates or restricting to the idempotent
section ECC11. Explicitly $G_t$ is the two-by-two adjugate of $J_t$
divided by $\det J_t>0$.

ECC21 now has a complete original quotient-metric inverse cost. For
$F=H^{-1},H^{-2}$ put
\begin{equation}
 a_F=\Tr(G_t^{-1}F^*G_tF),\quad b_F=|\det F|^2,\quad
 \lambda_{F,\pm}=(a_F\pm\sqrt{a_F^2-4b_F})/2.
 \tag{ECC40}
\end{equation}
The four singular values of $\Theta^{-1}$ in
$\diag(G_t,G_t)$ are the sorted union
$\tfrac12\sqrt{\lambda_{H^{-1},\pm}}$ and
$\tfrac12\sqrt{\lambda_{H^{-2},\pm}}$. Indeed the two off-diagonal
blocks act between orthogonal identical metric copies, and the
characteristic polynomial of each positive squared singular matrix is
$z^2-a_Fz+b_F$. This proves exact finite values and all exterior norms
by products, with no discarded direction.

Here are also all leading constants. Put $a=(Ah_0)|_{t=0}$,
$\eta=(h_1/h_0)|_{t=0}$, $N=\left(\begin{smallmatrix}0&0\\1&0\end{smallmatrix}\right)$,
and $n^2=G_{11}(G^{-1})_{00}=K/L$ (indices are zero-based). ECC21 gives
\begin{equation}
 \lim_{t\to0}t^2\Theta^{-1}=
 \frac12\begin{pmatrix}0&2N/a\\a^{-2}(I_2-2\eta N)&0\end{pmatrix},\qquad
 \Theta(0)=\begin{pmatrix}0&0\\4aN&0\end{pmatrix}.
 \tag{ECC41}
\end{equation}
The difference $a_t-a_0=A t(h_1+t)$ proves the limit even when $A,h_0$
depend on $t$. The first limit has rank three and its positive singular
values are the sorted list $(c_1,c_2,c_3)$ of
\begin{equation}
 \left\{\frac n{|a|},\quad
 \frac{\sqrt{1+|\eta|^2n^2}+|\eta|n}{2|a|^2},\quad
 \frac{\sqrt{1+|\eta|^2n^2}-|\eta|n}{2|a|^2}\right\}.
 \tag{ECC42}
\end{equation}
To verify it, the rank-one $N$ has squared metric norm
$\Tr(G^{-1}N^*GN)=G_{11}(G^{-1})_{00}=n^2$ and trace zero.
Consequently $I-2\eta N$ has determinant one and squared Frobenius norm
$2+4|\eta|^2n^2$, yielding the two displayed singular values by the
quadratic formula. The limiting forward matrix in ECC41 has its only
nonzero singular value $4|a|n$. Inverting singular values and using
continuity gives three inverse singular values $c_j|t|^{-2}(1+o(1))$
and the fourth tending to $(4|a|n)^{-1}$. Thus, with
$d=(c_1,c_2,c_3,(4|a|n)^{-1})$,
\begin{equation}
 \lim_{t\to0}|t|^{2\min(j,3)}\|\wedge^j\Theta^{-1}\|
 =\prod_{k=1}^j d_k\quad(1\le j\le4),\qquad
 \|\wedge^4\Theta^{-1}\|=\frac1{16|t|^6|a_0a_t|^3}.
 \tag{ECC43}
\end{equation}
The determinant norm equals the absolute determinant because the same
full metric is used in domain and codomain. The product of the constants
is $1/(16|a|^6)$, agreeing with ECC21. This is the trace-to-residue
multiplier in the actual local quotient; the original fixed $T_A$
has not been replaced by $\Theta$.

\section{An exact ES path and its original constants}
The explicit labelled algebraic ES family in the received source is
\begin{equation}
 (P;P+t,2P,z_t),\quad z_t=\frac{2P(P+t)}{5P+7t},\quad
 A=-\frac{5P+7t}{22P^2+35Pt+7t^2},\quad
 h_0=-\frac{P^2(3P+5t)}{5P+7t},\quad h_1=-z_t.
 \tag{ECC44}
\end{equation}
Direct addition gives
$1/(P+t)+1/(2P)+1/z_t=4/P$; adding all four roots gives the displayed
unchanged leading coefficient $A$. Also
$h(P+x)=(x-P)(x+P-z_t)$, proving $h_0,h_1$. Near zero all required
units are nonzero when $P\ne0$. At $t=0$ the exact constants are
\begin{equation}
 a=3P/22,\quad H_0=|3P/22|,\quad\eta=2/(3P),\quad
 \sigma=q\zeta+\frac q{6P}\zeta^3,\quad
 c=\frac2{q^3}\left(1-\frac{\zeta^2}{2P}\right),\quad q^2=3P/22.
 \tag{ECC45}
\end{equation}
These specialize ECC10, ECC19 and every original-metric constant in
ECC35--43, retaining $K_U,L_U$ and $K_W,L_W$ separately. The local residue
shear in ECC15 is $1/(3P)$. The four signed branches have values
$\sigma=\pm q_0\sqrt{-t}$ and $\sigma=\pm q_t\sqrt t$, with their two
root labels unchanged; thus the minimum, covariance, jet and trace costs
all concern the same four states with the original odd coordinate.
This is a complex algebraic family, not an assertion of integer ES
witnesses at arbitrary $t$.

Every map into either original physical space is now explicit:
translate ECC11 from $x$ to $r$ with $C_{-P}$, apply the two ECC23
receivers, and use ECC9 for $\zeta$ columns when required. The source and
target maps intertwine by $T_A\oplus T_A$ identically, since ECC23 does.
For local quotients their respective minimum metrics and all costs are
ECC26--43. These observations prove the exact return and all its finite
costs. The nonproper signed chart, the degenerating trace, and the
raw-point observation have explicit losses; the fixed original conductor
itself remains invertible with its unchanged nonzero moment $\mu_v$.
No arithmetic order jump or Riemann-hypothesis conclusion is inferred
from a different map's pole.

\section{The complete ES signed monodromy}
Here use the actual normalized ES parameter space
$\operatorname{Spec}\C[A,A^{-1},P,P^{-1},W]$, with
$s_1=-A^{-1}-P$ and $s_3=PW/4$, removing the divisors where roots
coincide. The four roots keep their prime/denominator marking. Put
$c=A^{-1}+P$. Over the function field $\C(P,c,W)$ the denominator
polynomial is
\begin{equation}
 g(T)=T^2(T+c)+W(T-P/4).
 \tag{ECC46}
\end{equation}
It is irreducible. In fact the two coefficients of this polynomial as
a degree-one polynomial in $W$ are coprime in $\C(P,c)[T]$:
$T=P/4$ is not a root of $T^2(T+c)$. A factorization in the two-variable
polynomial ring would therefore have a factor of $W$-degree zero
dividing both coefficients, which is impossible unless it is a unit.
Gauss's lemma then proves irreducibility in $\C(P,c,W)[T]$. Thus the
permutation group on the three denominators is transitive.

There is an explicit denominator-transposition loop within this ES
base. Fix $P\ne0$ and $a=2P$ and use the parameter $u$ with roots
\begin{equation}
 P,\quad a+\sqrt u,\quad a-\sqrt u,\quad
 z(u)=\frac{P(a^2-u)}{4(a^2-u)-2Pa},\qquad
 A(u)=-\frac1{P+2a+z(u)}.
 \tag{ECC47}
\end{equation}
Their denominator reciprocal sum is $4/P$ by direct addition. All
quartic coefficients are analytic in $u$, because the two moving roots
enter symmetrically. At zero the latter denominator is $P/3$, distinct
from both $P$ and $a$. The denominator discriminant is
$4u[(a-z(u))^2-u]^2$, with nonzero coefficient of $u$ at zero. Therefore
a small loop in $u$ exchanges exactly the two coalescing denominators.
A transitive subgroup of $S_3$ containing a transposition is $S_3$:
its order is divisible by both $3$ and $2$, hence equals $6$.

At a base point label the four roots by $r_0=P,r_1,r_2,r_3$ and their
odd square roots by $\sigma_i^2=A\prod_{j\ne i}(r_i-r_j)$. The product
of these four equations is
$(\prod\sigma_i)^2=A^4\prod_{i<j}(r_j-r_i)^2$. The quotient
$\prod\sigma_i/[A^2\prod_{i<j}(r_j-r_i)]$ has values in $\{1,-1\}$
and is constant along each lifted path. Thus every monodromy element
$(\epsilon,\pi)$ satisfies
\begin{equation}
 \pi\in S_3\text{ fixes }0,\qquad
 \epsilon_i\in\{1,-1\},\qquad
 \prod_{i=0}^3\epsilon_i=\operatorname{sgn}\pi.
 \tag{ECC48}
\end{equation}
This is an upper bound of $8\cdot6=48$ elements. The prime--denominator
loop ECC44, with ECC8's analytic unit, changes precisely the two signs
at $P$ and $P+t$ and fixes the root permutation; the other two signed
roots stay in nonvanishing analytic discs and do not change sign.
All loops may be based in the same connected complement of the
discriminant by joining their initial points with paths; this complement
is connected because it is a nonempty Zariski-open subset of the
irreducible smooth complex parameter space. Conjugation by the already
proved denominator $S_3$ gives the three flips $(0,1),(0,2),(0,3)$.
They are independent vectors in the four-bit sign space and generate
its full eight-element even subgroup. Together with a surjection onto
$S_3$, this gives at least $48$ elements. The upper bound proves exact
equality with ECC48. Thus the global group assertion is proved on this
precisely specified ES base, in addition to the full local original
coordinate and metric maps above.


\section{The actual original conductor on the complete collision quotient}
ECC23--43 gives both original receiving metrics and every raw observation,
jet and trace cost. Here we evaluate the induced original conductor itself
between those two quotient metrics. This retains its original moments,
Gamma masses, order, phase and centres \cite{WCF}; the four marked signed
states are the ones in ECC31, with the original ES marking \cite{ES488}.
The resulting estimate is a direct application to the original
inverse-exterior quantity, rather than to a different multiplication map.

Let $\mathcal U=U\oplus U$, $\mathcal W=W\oplus W$ have their unchanged
original Hilbert metrics and let $\mathcal T=T_A\oplus T_A$. Write
$\mathcal F_b=F_b\oplus F_b$ for $b=U,W$, so
$\mathcal T\mathcal F_U=\mathcal F_W$. Retain the literal local
evaluation matrix $E_t$ and the two positive covariance matrices from ECC39:
\[
 L_t=\diag(E_t,E_t),\qquad
 J_{b,t}=E_tH_b^{-1}E_t^*,\qquad G_{b,t}=J_{b,t}^{-1},\qquad
 q_{b,t}=L_t\mathcal F_b^{-1}.
 \tag{OCQ1}
\]
The matrix $E_t$ has rank two including at $t=0$: its rows are evaluation
at $P$ and the exact divided difference, with the latter equal to the
derivative at zero. On the polynomials $1,r$ those rows form
$\left(\begin{smallmatrix}1&P\\0&1\end{smallmatrix}\right)$.
Thus each $q_{b,t}$ is onto $\mathbb C^4$ and has its complete
four-dimensional kernel. Exact multiplication proves
\[
 q_{W,t}\mathcal T=q_{U,t},\qquad
 \mathcal T\ker q_{U,t}=\ker q_{W,t}.
 \tag{OCQ2}
\]
The reverse kernel inclusion follows by applying $\mathcal T^{-1}$;
there is no assertion based only on matching dimensions.
Equip the quotient by each kernel with its attained minimum norm.
ECC39 proves that its coordinate Gram is
$\mathcal G_{b,t}=\diag(G_{b,t},G_{b,t})$. Therefore the actual induced map is
\[
 \overline{\mathcal T}_t:
 (\mathbb C^4,\mathcal G_{U,t})\longrightarrow
 (\mathbb C^4,\mathcal G_{W,t}),\qquad z\longmapsto z.
 \tag{OCQ3}
\]
Identity coefficients do not equate the two metrics. OCQ2 proves
that OCQ3 is the quotient of the original conductor, with both kernels
and both norms retained, rather than an introduced operator.

\subsection{Every finite singular value and inverse exterior norm}
Define the exact scalar functions
\[
 a_t=\operatorname{tr}(J_{W,t}J_{U,t}^{-1}),\qquad
 b_t=\frac{\det J_{W,t}}{\det J_{U,t}},\qquad
 \kappa_\pm(t)=\frac{a_t\pm\sqrt{a_t^2-4b_t}}2.
 \tag{OCQ4}
\]
All matrices in this formula are the explicit finite original-moment
matrices ECC24--28 and ECC39; the two-by-two inverse uses its adjugate
and positive determinant. No coefficient or metric is replaced.
For an entirely expanded two-by-two formula, put
$J_b=\left(\begin{smallmatrix}\alpha_b&\beta_b\\
 \bar\beta_b&\delta_b\end{smallmatrix}\right)$ and
$d_b=\alpha_b\delta_b-|\beta_b|^2$. Then
\[
 a_t=\frac{\alpha_W\delta_U+\delta_W\alpha_U
                -2\operatorname{Re}(\beta_W\bar\beta_U)}{d_U},
 \qquad b_t=d_W/d_U.
 \tag{OCQ5}
\]
Positive definiteness shows that $J_WJ_U^{-1}$ is similar to the
positive Hermitian matrix $J_U^{-1/2}J_WJ_U^{-1/2}$.
Its eigenvalues are therefore the positive numbers in OCQ4, including
the repeated case $a_t^2=4b_t$. Isometric coordinates for the inverse
of OCQ3 give the squared singular matrix
$\mathcal G_{W,t}^{-1/2}\mathcal G_{U,t}
 \mathcal G_{W,t}^{-1/2}$, whose two parity blocks have these same
eigenvalues. Consequently all four inverse singular values are
\[
 \sqrt{\kappa_+},\ \sqrt{\kappa_+},\
 \sqrt{\kappa_-},\ \sqrt{\kappa_-}.
 \tag{OCQ6}
\]
Taking their ordered products proves every inverse exterior norm:
\[
 \begin{array}{c|cccc}
 j&1&2&3&4\\\hline
 \|\wedge^j\overline{\mathcal T}_t^{-1}\|&
 \sqrt{\kappa_+}&\kappa_+&\kappa_+\sqrt{\kappa_-}&b_t.
 \end{array}\tag{OCQ7}
\]
Both $J_{b,t}$ tend to their positive evaluation--derivative Grams
$J_{b,0}$, proved in ECC34. Substitution of those same matrices in
OCQ4--7 gives every finite nonzero limiting constant at $t=0$.
In particular no power of $|t|^{-1}$ occurs in this induced original
conductor inverse, although ECC37 and ECC43 compute such powers for
the raw-value and trace inverses separately.

\subsection{Exact cancellation in the original signed observations}
For $t\ne0$ the raw signed observation is $O_t$ from ECC31 and its
original covariance is $C_{b,t}=O_t\diag(H_b^{-1},H_b^{-1})O_t^*$.
The same unchanged $Z_t$ from ECC38 acts at both receivers. ECC39 is
the full congruence
\[
 Z_tC_{b,t}Z_t^*=\diag(J_{b,t},J_{b,t}),\qquad
 C_{b,t}^{-1}=Z_t^*\mathcal G_{b,t}Z_t\quad(b=U,W).
 \tag{OCQ8}
\]
Thus $Z_t$ is an isometry from each raw-value quotient metric to its
own jet quotient metric, simultaneously. The induced conductor on raw
values is also the coordinate identity by OCQ2 with $O_t$ in place
of $L_t$. Conjugation by the same $Z_t$ proves that its entire inverse
spectrum is OCQ6, not just that a leading determinant happens to cancel.
For example the exact top exterior cancellation is
\[
 \sqrt{\frac{\det C_{W,t}}{\det C_{U,t}}}
 =\frac{\det J_{W,t}}{\det J_{U,t}}=b_t.
 \tag{OCQ9}
\]
The common factors from $|s_0s_t|$ and the divided difference are
identical in numerator and denominator. Both odd signs and both even
directions survive this identity. The raw change is singular at zero;
the four-dimensional jet quotient OCQ1--3 extends through zero and
provides its specified, positive limiting metric.

\subsection{A bound by the unchanged original inverse-exterior quantity}
Let $\pi_{U,t}:\mathcal U\to\mathcal U/\ker q_{U,t}$ be the
quotient map, and $\iota_{W,t}:\mathcal W/\ker q_{W,t}\to\mathcal W$
the unique minimum-norm representative. The first is a contraction;
the second is an isometry by definition of the quotient norm. OCQ2
proves the exact inverse factorization
\[
 \overline{\mathcal T}_t^{-1}
   =\pi_{U,t}\mathcal T^{-1}\iota_{W,t}.
 \tag{OCQ10}
\]
To verify it, lift a target class, apply $\mathcal T^{-1}$ and take
its source class; OCQ2 gives that class the same four coordinates.
Changing its target representative changes the source by its full
kernel, so this verification is independent of the chosen lift.

Let $\gamma_1\ge\cdots\ge\gamma_8>0$ be the actual eight singular
values of $(T_A\oplus T_A)^{-1}$ in its original two metrics, as in
ET33--39. For any rank-less-than-$j$ operator $B$ approximating
$\mathcal T^{-1}$, the operator $\pi_U B\iota_W$ has rank less than
$j$ and approximation error at most $\|\mathcal T^{-1}-B\|$.
The least rank-less-than-$j$ error equals the $j$th singular value:
the singular-value decomposition attains it by retaining the first
$j-1$ singular summands, and any rank-less-than-$j$ map kills a vector
in the span of the first $j$ right singular vectors, giving the reverse
inequality. Applying this fact to OCQ10 proves
\[
 s_j(\overline{\mathcal T}_t^{-1})\le\gamma_j\quad(1\le j\le4),
 \qquad
 \boxed{\|\wedge^j\overline{\mathcal T}_t^{-1}\|
       \le\prod_{k=1}^j\gamma_k
       =\|\wedge^j\mathcal T^{-1}\|.}
 \tag{OCQ11}
\]
This holds at every admitted parameter, including the collision itself.
It applies directly to the actual original inverse-exterior estimate,
with the original source and target norms. It proves that this
particular collision quotient creates no additional inverse-exterior
loss. The conclusion is confined to the exact four-state quotient
and the fixed original conductor used throughout this calculation;
it does not assert an arithmetic Frobenius identification or an RH
endpoint.

\section{Visual atlas of the exact correspondence}
The figures below supplement the complete proofs. Their captions retain original source identity, exact proof locators, and the distinction between plotted quantities and schematic incidence.
\begin{figure}[p]
\centering
\includegraphics[width=\linewidth]{fable_01_four_marks.png}
\caption{The four separate tetrahedral locales and the complete marked complex columns of $K$. The displayed real view is an isometry of $\mathbf1^\perp$, with its basis specified; it does not project the complex polynomial points into an invented geometry. The original common fibre has seven affine points and eight completed factor states. Full proofs: FC1--7, FC15a, FC31--44, GF7--9, FX1--8, MO11. Original marked construction: The Clankers \cite{ES488,ES}.}
\label{fig:fable_01_four_marks}
\end{figure}
\clearpage
\begin{figure}[p]
\centering
\includegraphics[width=\linewidth]{fable_02_exact_receiver.png}
\caption{The exact original marked receiver and all of its polynomial, action and minimum-metric corrections. The original observation is $O_+K^{-1}$; the original generator is $K\diag(k\rho_l)K^{-1}$. The two positive residuals arise from different kernels. Full proofs: MO1--10 and AO1--23; original sources \cite{Observed,EightMinimum,CornerNonzero,Observe}; marked columns \cite{ES488}.}
\label{fig:fable_02_exact_receiver}
\end{figure}
\clearpage
\begin{figure}[p]
\centering
\includegraphics[width=\linewidth]{fable_03_five_escapes.png}
\caption{The complete eight-branch specialization of the actual marked arithmetic orbit MO12, on the punctured disk constructed in MO15. Four states escape in direction $e_a$ with quarter-power rate; one escapes in direction $e_w$ with third-order rate; three have the displayed finite limits. The diagram is a branch-count schematic, not a real-time trajectory plot. Full proofs and original observed norm constants: MO12--20, AO19--23. Original polynomial: \cite{ES488}.}
\label{fig:fable_03_five_escapes}
\end{figure}
\clearpage
\begin{figure}[p]
\centering
\includegraphics[width=\linewidth]{fable_04_collision_spectrum.png}
\caption{Exact coordinate magnitudes of the selected-root collision branch RC1 for $10^{-4}\le\varepsilon\le0.49$, without rescaling any coordinate. The four original-Gram singular asymptotes are proved in RC3--8 with all constants, and are not inferred from the graph. This curve has seven nearby affine states, unlike the transverse orbit of Figure~\ref{fig:fable_03_five_escapes}. Full original-conductor transfer: RC9. Branch provenance: received derivation \cite[Section 4]{WebFable}; original polynomial \cite{ES488}; original metric dictionary \cite{WCF}.}
\label{fig:fable_04_collision_spectrum}
\end{figure}
\clearpage
\begin{figure}[p]
\centering
\includegraphics[width=\linewidth]{fable_05_observation_recovery.png}
\caption{The exact finite recovery of every actually observed label, including all nonzero coefficients of both escape directions in the marked basis. Original-Gram cross terms are retained. No active column is replaced by a numerical example. At an empty active support, the complete receiver retains all four hidden labels. Full proofs: AO1--23 and MO20; original observation source \cite{Observe}; original marked columns \cite{ES488}.}
\label{fig:fable_05_observation_recovery}
\end{figure}
\clearpage
\begin{figure}[p]
\centering
\includegraphics[width=\linewidth]{fable_06_signed_loop.png}
\caption{The full signed-state boundary loop and its nonzero mixing. The varying separation belongs to the auxiliary quartet, while the receiving original conductor is held fixed. The rank-one limiting residue and rank-three finite based mixing are different evaluated matrices, both retained. Full proofs: SE1--18, SM1--19, SF1--18 and TC1--11. Original marked construction \cite{ES488}; original metric provenance \cite{WCF}.}
\label{fig:fable_06_signed_loop}
\end{figure}
\clearpage
\begin{figure}[p]
\centering
\includegraphics[width=\linewidth]{fable_07_exact_specialization.png}
\caption{The exact two-leg specialization $8\to7\to3$, with both original target paths, their nonzero discriminants before the endpoints, all escaping labels and the full fixed-Gram limits. The first leg loses only the negative infinity state from the original chart. The second loses the two signs each over the N and T roots. The complete source coordinates and proof, including the T/E cross pairing under $J$, are VA1--6. This is a concatenation of two specified paths rather than an arithmetic orbit. Original marked polynomial and points: The Clankers \cite{ES488}.}
\label{fig:fable_07_exact_specialization}
\end{figure}
\clearpage
% END VISUAL_ATLAS

% BEGIN RATIONAL_RESIDUE_VISUAL_ATLAS
\clearpage
\section{Visual atlas of the rational and arithmetic continuation}
\begin{figure}[p]
\centering
\includegraphics[width=\linewidth]{fable_08_two_forms.png}
\caption{The full eight-state rational receiver transports the complete Weil form and each of the two original direct-sum norms. The displayed upper norm is the actual target norm with $G_Y$, while RW19--20 supply the actual source norm with $G_X$ and $H_{U,\omega}$. Complete matrices, both inverses and every kernel component are RW1--20; all original source coefficients and moments are NH26--42. Original conductor \cite{WCF}; received rational test construction \cite{ResidueIntake}. No sign of the Weil form is inferred from positivity of the Gamma metric.}
\label{fig:fable_08_two_forms}
\end{figure}
\clearpage
\begin{figure}[p]
\centering
\includegraphics[width=\linewidth]{fable_09_native_height.png}
\caption{Raw eigenvalues of the complete target coefficient Gram at the original permitted choice $s=1$ and test pole $\omega=2+i\tau$. Solid curves evaluate the exact finite moments with 90-digit arithmetic, and dashed curves display the proved asymptotes with every constant. The full determinant check has relative error below $3\cdot10^{-58}$ on the plotted sample. The exact statements and all compounds, flags, source constants and conductor identities are RW7--11 and NH1--42. Original Gamma dictionary and human source identity \cite{WCF}. These numerical curves illustrate the proof; they are not a zero search.}
\label{fig:fable_09_native_height}
\end{figure}
\clearpage
\begin{figure}[p]
\centering
\includegraphics[width=\linewidth]{fable_10_residue_boundary.png}
\caption{Two sequences of original ES--Fable points approach the same marked M with the same displayed $a$-coordinate distance. AX7--8 prove that their selected arithmetic residues respectively tend to zero and diverge in absolute value. AR8--9 compute the distinct global residue at infinity. The complete map, original Jacobian and all completion factors remain unchanged. Original marked map \cite{ES488}; completed zeta and Gamma modulus \cite{DLMFXi,DLMFGamma}. The plotted values use the stated exact expressions.}
\label{fig:fable_10_residue_boundary}
\end{figure}
\clearpage
\begin{figure}[p]
\centering
\includegraphics[width=\linewidth]{fable_11_confluent_jets.png}
\caption{The complete paired-evaluation morphism through the collision of the literal quartet. Its exact determinant remains nonzero; its limit retains values and derivatives in a four-dimensional coefficient algebra. AR14--17 provide every original coefficient matrix and the full Hermite inverse, while AX1--3 retain the completion jet and both Laurent coefficients of the colliding arithmetic residues. The resultant-one factor space has no finite state over a multiple root; its limiting jet is nevertheless completely specified. Original marked factor map \cite{ES488}, completion \cite{DLMFXi}, received calculation \cite{ResidueIntake}. The drawing shows incidence, not distances.}
\label{fig:fable_11_confluent_jets}
\end{figure}
\clearpage
% END RATIONAL_RESIDUE_VISUAL_ATLAS


\clearpage
\section{Visual atlas of the complete native Weil continuation}
\begin{figure}[p]
\centering
\includegraphics[width=\linewidth]{fable_12_continuous_positivity.png}
\caption{The entire continuous real interval $|\tau|\le2.97\cdot10^{12}$
has a strictly positive complete cubic form. The four displayed intervals
are exact; no invented zero ordinates are plotted. Their selected-zero
Gram is bounded below by $b_*$, and every unverified reflected zero is
included in the possible negative tail $e_*$. The remainder $\mu_*=b_*-e_*$
is strictly positive. CP1--18 proves every factor and endpoint.
The human inputs are Platt--Trudgian \cite[Theorem 1]{PlattTrudgian}
and Hasanalizade--Shen--Wong \cite[Corollary 1.2]{HSW}.}
\end{figure}
\clearpage
\begin{figure}[p]
\centering
\includegraphics[width=\linewidth]{fable_13_native_certified_spectrum.png}
\caption{All twenty certified generalized eigenvalues at five poles are shown
in the unchanged original target norm with the permitted $s=1$ and
$M=\sqrt{2\pi}$. Each marker lies in its exact rational interval of
relative width below $10^{-60}$; four graph zeros remain explicitly
displayed off the logarithmic axis. There is no interpolation across poles.
NV1--9 proves completeness, simplicity and the cofactor eigenline formulas.
The complete coefficient identity is \cite[MP6--15]{MovableWeil}; the
finite arithmetic provider is \cite[WP24--27]{WeilCertificates}; the
interval-arithmetic provenance is Johansson \cite{JohanssonArb}.}
\end{figure}
\clearpage
\begin{figure}[p]
\centering
\includegraphics[width=\linewidth]{fable_14_exact_native_return.png}
\caption{Every original $M,N,T,E$ sign state remains in the full receiver.
The four-dimensional evaluation-graph kernel is exact throughout the
continuous interval, and the displayed section uniquely minimizes both
original norms. The conductor arrow carries the exact source representative
to its exact target; the two metrics remain separate. CP19--29 and NW1--16
prove the maps, inverses, full traces and bounds; NV7--9 gives the certified
eigenline return. Here $t_W=\operatorname{tr}H_\tau$ and
$t_U=\operatorname{tr}H_{U,\tau}$, fully evaluated in NW11--16.
The original quartet domain is $0<\delta<1/2$, $\gamma>2$.
The marked polynomial is credited to the ES reader \cite{ES488}, and
the human orthogonal-polynomial sources remain in \cite{WCF}.}
\end{figure}
\clearpage

\clearpage
\section{The completed critical fibre, residue duality and trace loss}
\begin{figure}[p]
\centering
\includegraphics[width=\linewidth]{15_residue_trace_and_critical_jets.png}
\caption{The two exact length-four blocks at the formal paired collision
retain the original upper labels M,T and lower labels N,E, each with both
nearby signs. Both full residue matrices are displayed, with the original
factor $-1/2$, original reciprocal factor $t$, and every $\gamma$-dependent
entry. Their direct sum has rank eight, while ordinary trace has rank two.
The exact comparison is multiplication by $2t^3$ followed by the residue
duality, so its six-dimensional kernel is completely identified. The
two arithmetic measurements retain both $X$ and $X'$ without division
by either. FS21--27 and RD9--17 prove the algebras, matrices and tests;
RD18--23 gives every test vector in both original conductor spaces
and their unchanged norms. This is an exact algebraic diagram, with no
sampled or assumed arithmetic zero. Classical Hankel/Bezout provenance
is Jouve and Rodriguez-Villegas \cite[Section 2]{JouveRV}; the marked
polynomial source remains the ES reader \cite{ES488}.}
\end{figure}
\clearpage

\clearpage
\section{The marked inverse pole and every original-metric cost}
\begin{figure}[p]
\centering
\includegraphics[width=\linewidth]{16_marked_inverse_pole.png}
\caption{The exact original marked receiver is isometric to the two-parity
coefficient receiver, with the full mixed source Gram and unchanged original
conductor. MR1--7 give every map, inverse and metric entry. The two chains
shown with their original coefficients span all eight dimensions; their
successive images have ranks $8,6,4,2,0$, and the phase-sensitive identity
$\mathscr L^3=-\Theta_0/(2^{16}g^6)$ identifies the final image with
the critical trace image. MR8--12 prove both chains and their original
source/target returns. RS8--19 and RS21--29 evaluate every singular constant
using all original moments and Gamma weights. MR13--18 give the complete
rational inverse remainder and finite-parameter bounds; MR19--26 give a
positive explicit threshold and uniform errors for all eight singular
values and every exterior rank. This is a symbolic diagram, not a sampled
zero set. It varies the completed quartic while holding the original
reference conductor fixed. Human provenance remains the ES reader
\cite{ES488}, the original polynomial account \cite{Tao}, the Gamma
dictionary \cite{DLMF18Source}, and the residue/Hankel/Bezout construction
\cite[Section 2]{JouveRV}.}
\end{figure}
\clearpage

\clearpage
\section{The global residue zeros and their original cubic return}
\begin{figure}[p]
\centering\includegraphics[width=\linewidth]{17_global_residue_cubic.png}
\caption{Exact symbolic map diagram, with no sampled arithmetic roots.
GB1--22 proves the global eight-section basis and its complete multiplication.
GD3--12 computes both lost infinity signs, the actual dualizing line and
the perfect state--dual evaluation. GD14--21 gives every original marked
and physical receiver, the full dual metric, all inverse singular constants
and exterior limits; GD26--28 gives the positive finite-parameter bounds.
GD22--25 identifies the cubic residue weight with the exact cancellation
in the original sign map. GDR19--21 additionally proves that this weighted
map has Jacobian determinant $a^{12}/4$ and rank one at its collapsed
boundary. The fixed original polynomial still has determinant $-2$.
The global dualizing obstruction and its tensor powers are GDR14; the
same proof explicitly retains the perfect residue on $A=-1/2$.
Human basis and ideal provenance: Wood \cite{WoodBinary}, with her
Birch--Merriman/Nakagawa attribution. Classical residue provenance:
Jouve--Rodriguez-Villegas \cite{JouveRV}; original polynomial and marking:
\cite{Tao,ES488}.}
\end{figure}
\clearpage
\section{Every state in the ES-to-paired-quartet family}
\begin{figure}[p]
\centering\includegraphics[width=\linewidth]{18_complete_es_family.png}
\caption{The left panel is the exact projection of EF1 to the complex
root coordinate for $p=5,\gamma=2$; straight segments are the displayed
linear root paths. It projects out the sign coordinate $t$ and does not
identify its two values. All four labels and eight coefficient states
are retained: the M/T length-four block stays nonreduced, while the four
N/E signed branches join in the second length-four endpoint block.
EF5--8 gives the complete projectors, coefficient matrices and inverses;
EF9--12 keeps all cofactor residue corrections. EF13--15 gives each
original affine inverse branch and its exact original-norm escape.
EF16--24 computes its lower-block inverse, all retained-metric constants,
every exterior limit and exact finite-parameter singular values.
The entire eight-state trace is singular throughout; only the explicitly
identified lower block is inverted. These are auxiliary polynomial
root trajectories, not asserted arithmetic zeta-zero paths.
Human ES/polynomial sources remain \cite{ES488,Tao}; the received local
comparison \cite[Section 5]{ESRHReceived} is extended here to the whole
eight-dimensional family.}
\end{figure}
\clearpage

\clearpage
\section{The whole signed-cover coefficient comparison}
\begin{figure}[p]
\centering\includegraphics[width=\linewidth]{19_es_total_coefficient_extension.png}
\caption{An exact symbolic diagram of ET1--39, not a plot of arithmetic
zeta zeros. The three coefficient directions retain every one of the
eight signed basis directions; they are not labelled as twenty-four
distinct geometric points. ET13--15 proves the full inclusion, its
quotient, its conductor ideal and its specialized Tor sequence.
ET16--22 gives every root and sign action, all commutation maps, and
the eight length-three residue chains. ET23--36 returns the complete
inverse through the fixed original source and target, evaluates all
finite singular values and every exterior norm with the original
moments and Gamma factors present. ET37--39 gives a finite two-sided
bound in every exterior degree. Independent proofs are ETR1--31.
Immediate coefficient source: \cite[eqs.\ (3)--(8)]{DefiningPrimeReceived}.
Original human polynomial and ES provenance: \cite{Tao,ES488};
the original Gamma dictionary is retained in FC20--24 and RS21--26.}
\end{figure}


\clearpage
\section{The defining prime and the original collision quotient}
\begin{figure}[p]
\centering\includegraphics[width=\linewidth]{20_prime_and_original_collision.png}
\caption{Exact symbolic diagram of the full arithmetic inclusion and the
actual original collision quotient. PS6--25 and ISR1--34 retain all eight
signed arithmetic directions, every unit and every inclusion entry.
PB1--8 and PES1--12 impose the original positivity constraint; the
one-divisible valuation-one stratum is permitted by the proved bound,
not asserted to occur. The integral-ideal statement assumes $p\nmid S$;
the different fractional chart and its exact twist are ISR18--23.
ECC31--43 calculates the raw signed covariance and trace losses in each
original metric. The displayed four states are the colliding block of
the complete eight-state algebra, whose complementary block is retained
in ECC11. OCQ1--11 proves that the same raw-to-jet change acts at both
receivers, gives every finite inverse singular value of the induced
original conductor, and proves its exterior bound by the unchanged
original conductor inverse. The image is not a zeta-zero trajectory.
Human sources at use: \cite{ES488,Tao,DLMF18Source,JouveRV}; immediate
received proposal: \cite{DefiningPrimeReceived}.}
\end{figure}

\begin{thebibliography}{99}
\bibitem{ES488} The Clankers, \emph{Erd\H{o}s--Straus project reader},
canonical 488-page edition, deposited source version 69,
\texttt{erdos\_straus\_project\_reader.tex}, theorem
\texttt{explicit-four-time-Jacobian-collision}, source lines
19484--19588. Source SHA-256:
\texttt{678b3f70ee7d2a6d00f6c588d7778e390118ab4743a5759612dbba6430984c78}.
The formal source uses the labels $\mathrm M,\mathrm N,\mathrm T,\mathrm E$.
The seven-point fibre and missing eighth factor-sign state are proved
here as an extension, rather than attributed to this reader.
\bibitem{ES} The Clankers, \emph{Arithmetic Shells, Finite Support Geometry,
and Lorentz Lifts: Exact Identities, Certificates, and Symmetry Constructions},
dated 4 August 2026, \texttt{es\_fable\_c123\_preprint.tex},
consulted source snapshot, 19--20 September 2026. Exact theorem/equation
labels are cited at use. Source SHA-256:
\texttt{0e71a9c146eb860a9b66a2287dca673ab675e894e545df0ede642e74a4f353c4}.
The four-Time and conductor comparison in
this note is a new calculation, not an attribution to the human authors
of the original Jacobian example.
\bibitem{WCF} Split-Zero programme, \emph{A polynomial forward bound for
the full weighted conductor}, \texttt{WEIGHTED\_CONDUCTOR\_FORWARD.tex},
WCF1--6, source intake 17 September 2026. This source credits the
Meixner--Pollaczek dictionary to the NIST DLMF Chapter 18 by T. H. Koornwinder, R. Wong, R. Koekoek, R. F. Swarttouw and W. P. Reinhardt, equations
18.19.8--9, 18.22.8 and 18.23.7; its original phases and masses are
retained here, and no new orthogonality theorem is asserted.
\bibitem{NC} Split-Zero programme, \emph{Vanishing of the original
proper-source kernel through conductor coprimality}, 19 September 2026,
NC1--5, \texttt{ORIGINAL\_CONDUCTOR\_COPRIMALITY.tex}.
\bibitem{Cumulative} Split-Zero programme, accepted 782-page source
collection, DFX4--5 and CNC1--2, consulted local source snapshot.
\bibitem{Tao} Terence Tao, \emph{A digestion of the Jacobian conjecture
counterexample}, 21 July 2026,
\url{https://terrytao.wordpress.com/2026/07/21/a-digestion-of-the-jacobian-conjecture-counterexample/}.
Equations (1) and the displayed three-point fibre record the
Alp\"oge--Fable example. The ES four-Time locale is an additional marked
construction; it must not be credited as a fourth preimage in that
three-dimensional fibre.
\bibitem{WebFable} User-supplied web-session derivation, \emph{Four labelled roots, eight oriented factor states, and the original conductor}, \texttt{Fable\_Conductor\_Continuation\_20260920/PROOF.md}, received 20 September 2026. No author name was supplied. Sections 1--2 give the proposed factor completion; Section 4 gives the selected-root collision branch and singular scales; Section 6 gives the linear symmetry and polynomial flow candidates. The original incoming file and exact source hash are retained in this edition. The exact identity with the original ES polynomial is proved here with $t=w$. Explicit regular boundary inverses, full native singular constants, and the original corner quotient receiver are further calculations here.
\bibitem{Observed} Split-Zero programme, \emph{The actual observed class and its conductor return}, \texttt{OBSERVED\_CLASS\_RETURN.tex}, 20 September 2026, OB1--3. The full source is retained in \texttt{supporting\_proofs.zip}; OB3 proves the kernel inclusion used in MO1.
\bibitem{CornerNonzero} Split-Zero programme, \emph{Original corner nonvanishing}, \texttt{ORIGINAL\_CORNER\_NONVANISHING.tex}, OC50--74, 20 September 2026. Exact source SHA-256 \texttt{4f9ea962fb0d2bada8da3e43d1d3914d84a205e8b87a7ce7c99c0c009f11acaa}. Full proof and original-source bibliography are retained in \texttt{supporting\_proofs.zip}. MO1 uses its explicitly stated period domain and original phase; no additional zero-verification result is imported.
\bibitem{EightMinimum} Split-Zero programme, \emph{The eight-state original minimum}, \texttt{EIGHT\_STATE\_ORIGINAL\_MINIMUM.tex}, EQ1--21, programme result SZ-20260920-016, 20 September 2026. Source SHA-256 \texttt{ff74f5622189770f92916a70a121935e810de3cbe2bdbf9e9a393ede853fdf91}. Full proof in \texttt{supporting\_proofs.zip}; the exact original metrics of EQ1--5 are used in MO5--8.
\bibitem{Observe} Split-Zero programme, \emph{Full original packet recovery from iterated observation}, \texttt{04\_FULL\_ORIGINAL\_OBSERVABILITY.md}, O1--3 and O10, 19 September 2026. Full source in \texttt{supporting\_proofs.zip}. AO1--23 specialize its original coefficient test and Lagrange recovery to the four marked corner idempotents, keeping the original output correlations and invisible coordinates.

\bibitem{ResidueIntake} Received AI-generated web-session calculation,
\emph{Fable arithmetic residues and full-Weil rational tests},
20 September 2026. The complete unchanged incoming source is retained as
\texttt{source\_intake/INCOMING.md} in this edition's source archive.
No author name was supplied. Equations (1)--(17) and (33)--(36)
are the specific inputs independently calculated here; AR13 corrects
the omitted completion-factor modulus after incoming equation (16).
The separately developed full-Weil coefficient and tail calculation
is not assumed here.
\bibitem{DLMFXi} T. M. Apostol, \emph{Zeta and Related Functions},
NIST Digital Library of Mathematical Functions, version 1.2.8,
15 September 2026, \url{https://dlmf.nist.gov/25.4}.
Original equation TeX 25.4.3 and 25.4.4 supplies the reflection
identity and exact Riemann completion; 25.2.12 records the canonical
product. The unchanged equation sources are included with their hashes.
The source notes for 25.4 refer to Apostol (1976), p.~260.
The classical line-one argument associated with J. Hadamard and
C.-J. de la Vall\'ee Poussin (1896) is reproduced in AX4--6;
it is not claimed as a new general zeta theorem.
\bibitem{DLMFTheta} W. P. Reinhardt and P. L. Walker,
\emph{Theta Functions}, NIST Digital Library of Mathematical Functions,
version 1.2.8, 15 September 2026,
\url{https://dlmf.nist.gov/20.2.E3} and
\url{https://dlmf.nist.gov/20.7.E32}.
The original unchanged equation TeX supplies the theta series and
modular transformation. AX5 derives the entire theta integral and
the stated numerical bound from these two formulas.
\bibitem{DLMFGamma} R. A. Askey and R. Roy,
\emph{Gamma Function}, NIST Digital Library of Mathematical Functions,
version 1.2.8, 15 September 2026,
\url{https://dlmf.nist.gov/5.4.E4}.
Original equation TeX 5.4.4 is retained unchanged. Its exact
Gamma modulus supplies the equality and exponential factor in AX7.


\bibitem{Fable029} Split-Zero programme,
\emph{The ES--Fable inverse correspondence in the original weighted conductor},
result SZ-20260920-029, 20 September 2026. The complete preceding source is
preserved in this edition. FC21--35 constructs the original conductor maps;
SE1--17 constructs the full signed states; RW1--20 and NH1--42 prove both
original native metrics, all coefficient matrices and source inverses.
The original ES and human classical sources remain cited at their uses.
\bibitem{WeilCertificates} Split-Zero programme,
\emph{Complete-Weil positivity for the cubic Gaussian source and certified
movable-pole matrices},
\texttt{WEIL\_POSITIVITY\_AND\_ARITHMETIC\_CERTIFICATES.tex},
result SZ-20260920-030, WP1--27, 20 September 2026.
Complete provider proof and certificate are preserved unchanged in
\texttt{supporting\_proofs.zip}. WP14 proves the Gaussian inequality;
WP23 proves the analytic zero-pole cubic bound; WP24--27 certify five poles.
Exact source SHA-256
\texttt{ffe3f0cd195a3a2652f3d47540cd4ce9bfd8cb730cab29223d730b145a2ba73f}.
\bibitem{MovableWeil} Split-Zero programme,
\emph{Exact movable-pole Weil coefficients and finite arithmetic certificates},
\texttt{MOVABLE\_POLE\_WEIL\_DERIVATION.tex}, MP1--38 including MP6a,
20 September 2026. Complete provider source retained unchanged;
MP6--15 supplies the completed-function coefficient identity used here.
Exact source SHA-256
\texttt{146f8d45ffb93bc837f4ef80f9db681a6098ff31ee2f9f88d3df4ca2687b9d42}.
\bibitem{PlattTrudgian} Dave Platt and Tim Trudgian,
\emph{The Riemann hypothesis is true up to $3\cdot10^{12}$},
arXiv:2004.09765v1, Theorem 1, original TeX label \texttt{bank};
Bull. London Math. Soc., DOI 10.1112/blms.12460.
\url{https://arxiv.org/abs/2004.09765v1}.
The exact verified height is $3\,000\,175\,332\,800$;
CP3 uses the strictly smaller $3\cdot10^{12}$.
\bibitem{HSW} Elchin Hasanalizade, Quanli Shen and Peng-Jie Wong,
\emph{Counting zeros of the Riemann zeta function},
arXiv:2107.06506v1, Corollary 1.2,
original TeX labels \texttt{main-thm}, \texttt{main-bound-1}.
\url{https://arxiv.org/abs/2107.06506v1}.
The counting convention and every constant are retained in CP4.
The short-interval existence and full-tail deductions CP5--18 are proved here.
\bibitem{JohanssonArb} Fredrik Johansson,
\emph{Arb: efficient arbitrary-precision midpoint-radius interval arithmetic},
IEEE Trans. Computers 66 (2017), 1281--1292,
DOI 10.1109/TC.2017.2690633, arXiv:1611.02831v1.
\url{https://arxiv.org/abs/1611.02831v1}.
The original author TeX's enclosure-arithmetic and polynomial/matrix sections
provide the implementation provenance. NV uses python-flint 0.9.0 at 768 bits;
its independent stored-ball review uses 1024 bits without shrinking input radii.


\bibitem{ExceptionalAtlas} Received AI-generated continuation,
\emph{Exceptional-map atlas: the finite signed-root completion},
\texttt{EXCEPTIONAL\_ATLAS.md}, 20 September 2026, EA1--51.
Exact received-source SHA-256
\texttt{fbe6de554732ccd911fe36e75b7d566bd5dbf28719592c8d6a5f56e177c9fafd}.
FS independently proves the completion, all charts and fibre algebras;
FD proves its full original-metric differential and phase constants.
RD adds the complete residue/trace comparison and its original conductor return.
The incoming account is preserved unchanged as received AI output;
it is not an original human source for the classical algebra used here.
\bibitem{JouveRV} Florent Jouve and Fernando Rodriguez-Villegas,
\emph{On the bilinear structure associated to Bezoutians},
Journal of Algebra 400 (2014), 161--184,
DOI 10.1016/j.jalgebra.2013.10.030;
original author source arXiv:1004.4208v3, 12 September 2013,
\url{https://arxiv.org/abs/1004.4208v3}.
Section 2 proves the relation between the transfer functional,
Hankel Gram and Bezout inverse; Section 2.2, proposition
\texttt{prop1}, and Section 2.3, lemma \texttt{Q-lemma}
and theorem \texttt{thBsharp}, are the precise source locators.
The signed-root and conductor matrices RD1--29 are derived here
with the original quartic coefficient retained.
The original gzip-compressed TeX source was read and retained intact
in the research corpus. This edition includes its versioned retrieval
link, source hashes, and the bounded Section 2 source-use ledger.


\bibitem{Fable037} Split-Zero programme,
\emph{The ES--Fable four-point map and the original weighted conductor},
complete cumulative source of result SZ-20260920-037, 20 September 2026.
FC20--24 gives the original moment inverse, NH26--34 the full source norm,
and RD1--29 the perfect residue and its marked return. These full proofs
are retained earlier in the present source, with their original citations.
\bibitem{ReceivedRE} Received AI-generated web-session calculation,
\emph{Perfect residue, degenerating trace, and the complete
inverse-exterior cost}, 20 September 2026, RE1--30.
No human author name was supplied. The unchanged geometry excerpt is
\texttt{RE1\_30\_RECEIVED\_GEOMETRY.md}; its parent source has SHA-256
\texttt{d7d07ac1a77d1f706391b91c759cf45abc7bb134c27f5e1ad56444d6a76c1418}.
All of RE1--30 was read. The present RS/MR calculations derive the full
original-source and marked constants, explicit chains and finite-parameter
errors. Unrelated analytical portions of the received parent file are not
used or claimed reviewed here.
\bibitem{DLMF18Source} T. H. Koornwinder, R. Wong, R. Koekoek,
R. F. Swarttouw and W. P. Reinhardt,
NIST Digital Library of Mathematical Functions, Chapter 18,
\url{https://dlmf.nist.gov/18}, equations 18.19.8--9,
18.22.8 and 18.23.7. The original orthogonality and generating-function
dictionary is credited in WCF1--6; RS21--26 applies it with the complete
original centers, phases, Gamma mass and conductor inverse moments.


\bibitem{WoodBinary} Melanie Matchett Wood,
\emph{Rings and ideals parametrized by binary $n$-ic forms},
Journal of the London Mathematical Society (2) 83 (2011), 208--231,
DOI 10.1112/jlms/jdq074. Original author TeX, arXiv:1007.5508v1,
\url{https://arxiv.org/abs/1007.5508v1}.
Equations \texttt{E:basis}, \texttt{E:mult}, \texttt{E:Isbasis};
propositions \texttt{P:Iinvert}, \texttt{P:invertextra};
theorem \texttt{T:basechpf}; dual modules \texttt{E:dualmap},
\texttt{C:JInvdif}. The author credits the ring basis to B. J. Birch
and J. R. Merriman and its multiplication to J. Nakagawa; those
earlier papers are not represented as independently read here.
Original source SHA-256:
\texttt{0eed67746523862ee847bf5a70c94d8bce3680dd0400a4258b187c3e60e4643e}.
Actual reading coverage is recorded in the accompanying source-use ledger.
\bibitem{RamificationReceived} Received AI-generated web continuation,
\texttt{RAMIFICATION\_CONTINUATION.txt}, 20 September 2026,
Section 2. No human author name supplied. Received source SHA-256:
\texttt{259b9f7bcf84b6885f3e88a2ba045825d02d0e9e11e7df66c84ffe3750e2eae3}.
Its basis candidate is proved with the original infinity orientation
and full coordinate maps in GB1--22. Its further classification claims
are not adopted by this citation.
\bibitem{ESRHReceived} Received AI-generated calculation,
\emph{Track II---Residue duality, the exact non-unit element,
and boundary transport}, 20 September 2026, Sections 1--7.
No human author name supplied. Received source SHA-256:
\texttt{20c5fb683e9ffcac10003420dd7df5ae2f695f7c4bc8bf866ab506683b97e7f9}.
Its Section 5 gives the local ES/RH primary-algebra correction.
EF1--25 proves the complete flat family, all missing signed branches,
all residue corrections and exact original-metric inverse costs.
\bibitem{Fable042} Split-Zero programme,
\emph{Global signed duality and the complete ES family},
result SZ-20260920-042, 20 September 2026. The full GB1--22,
GD1--28, GDR1--30 and EF1--25 proofs are included in this cumulative
source, with all earlier original-conductor definitions and citations.

\bibitem{DefiningPrimeReceived} Received AI-generated web calculation,
\emph{The ES bridge at the defining prime of an actual Erdos--Straus witness},
20 September 2026, complete supplied account, equations (1)--(47).
No human author name was supplied.
Original received SHA-256: \texttt{f97ec18809279c10caf20a0ffa59abaa9cf54ec5c5c0ab21761c2840b426d886}.
The whole account was read. Its coefficient proposal (3)--(8) is
proved and extended here in ET1--39 to the complete signed cover and
the full original-conductor inverse-exterior cost.
Its integral-prime, gluing, collision and covariance assertions are now
proved with the stated corrections in PS1--25, ISR1--34, ECC1--48,
PB1--8 and PES1--12. OCQ1--11 additionally proves the actual original
quotient-conductor comparison. Exact source reading and derivation
coverage are recorded in the source-use ledger.


\bibitem{Fable050} Split-Zero programme,
\emph{The ES--Fable inverse correspondence in the original weighted
conductor}, complete preceding source, 20 September 2026,
programme result SZ-20260920-050. Principal SHA-256:
\texttt{efb34b95a4f99b1848eea9ce7f4c285aabd49a806b074cfd51a1fe11f07ee46d}.
All of its proofs and human-source citations remain in this cumulative
edition; ECC uses FC20--32, RD1--15, RS21--25 and GD14--17.

\end{thebibliography}
\end{document}
